\documentclass[10pt,a4paper,oneside]{article}
\usepackage[utf8]{inputenc}
\usepackage[T1]{fontenc}
\usepackage{lmodern}
\usepackage[margin=18mm,top=22mm,bottom=20mm,headheight=14pt]{geometry}
\usepackage{xcolor}
\definecolor{ink}{HTML}{173C52}
\definecolor{teal}{HTML}{087C80}
\definecolor{muted}{HTML}{5C6F78}
\usepackage[most]{tcolorbox}
\usepackage{hyperref}
\hypersetup{colorlinks=true,linkcolor=teal,urlcolor=teal}
\usepackage{fancyhdr}
\usepackage{enumitem}
\usepackage{microtype}
\usepackage{needspace}
\usepackage{titlesec}
\usepackage{listings}
\usepackage{tabularx}
\usepackage{longtable,booktabs}
\setlength{\parindent}{0pt}
\setlength{\parskip}{4pt}
\setlist[enumerate]{topsep=3pt,itemsep=1.4pt,parsep=0pt}
\titleformat{\section}{\LARGE\bfseries\color{ink}}{\thesection}{0.6em}{}
\titleformat{\subsection}{\large\bfseries\color{teal}}{\thesubsection}{0.6em}{}
\pagestyle{fancy}
\fancyhf{}
\lhead{\footnotesize Informatika attestatsiya | 1-mavzu}
\rhead{\footnotesize 2026}
\cfoot{\footnotesize\thepage}
\renewcommand{\headrulewidth}{0.3pt}
\lstset{basicstyle=\ttfamily\footnotesize,breaklines=true,columns=fullflexible,frame=single,keepspaces=true}
\hbadness=10000
\hfuzz=2000pt
\vbadness=10000
\vfuzz=2000pt
\emergencystretch=3em
\sloppy
\begin{document}
\begin{titlepage}
\vspace*{26mm}
{\color{teal}\rule{\textwidth}{1.7pt}}\par
\vspace{20mm}
{\Large INFORMATIKA ATTESTATSIYASI \ 2026}\par
\vspace{13mm}
{\fontsize{26}{32}\selectfont\bfseries\color{ink}AXBOROT VA RAQAMLI\\ SAVODXONLIK ASOSLARI}\par
\vspace{10mm}
{\Large Izohli testlar va darsliklar nazariyasi}\par
\vspace{17mm}
{\Large 1 031 ta savol va izohli javob}\par
{\normalsize 2026-yil test spetsifikatsiyasi bo'yicha 1--3-savollarga oid yo'nalish}\par
\vspace{17mm}
{\large Y1: 1 027 \quad Y2: 2 \quad Y3: 2}\par
\vfill
{\color{muted}Har bir testdan keyin to'g'ri javob va unga tegishli izoh berilgan.\\
Asos: attestatsiya100\% loyihasi, Supabase bazasining joriy yozuvlari.}\par
\vspace{12mm}
{\color{teal}\rule{\textwidth}{1.7pt}}
\end{titlepage}
\section*{Kitobdan foydalanish}
\addcontentsline{toc}{section}{Kitobdan foydalanish}
Savolni o'qing, tegishli javob variantini tanlang yoki moslashtirish/ketma-ketlikni bajaring.
So'ng savolning bevosita pastidagi \textbf{To'g'ri javob} va \textbf{Izoh}ni o'qing.
Y1 -- bitta to'g'ri javob; Y2 -- chap va o'ng tomon mosligi; Y3 -- amallar ketma-ketligi.
\par\textbf{Manba:} 2026-yilgi informatika attestatsiya spetsifikatsiyasi hamda
attestatsiya100\% Supabase bazasining savollar, variantlar va javob kalitlari.
Savollar va ularning izohlari Supabase test bankidan olingan. Qo'shimcha nazariyalar manba darsliklarga asoslanadi.
\par\textbf{Qamrov:} Barcha 1 031 ta \textit{published} savol. 147 ta \textit{archived} yozuv kiritilmagan;
\textit{review} holatidagi yozuv yo'q. 3 ta Y1 savolida 5 tadan variant bor;
ular asl shaklida saqlangan. 1 ta savol matni boshqa savolda ham uchraydi
(ikki alohida DB yozuvi, mustaqil ravishda saqlandi).
\par\textbf{Eslatma:} Ma'lumotlar bazasidagi javoblar va izohlar aynan ularga tegishli savoldan
keyin berildi. Bu tuzilish sifati bo'yicha avtomatik tekshiruv;
barcha savollar alohida inson eksperti tomonidan qayta baholanganini anglatmaydi.
\tableofcontents
\clearpage
\section{Informatika, ma'lumot, axborot va bilim}

% DB question_id: 13e8d884-a408-4706-a22d-76fd0fa404c6
\needspace{5\baselineskip}
\noindent\textbf{\color{ink}1.}
\par
"Xabar" atamasiga eng aniq ta'rifni toping.
\begin{enumerate}[label=\Alph*),leftmargin=9mm]
\item axborotni qabul qiluvchiga yetkazish uchun ifodalangan mazmun
\item xabarni kanal bo'ylab tashiydigan fizik o'zgarish
\item axborotni izlash, yig'ish, saqlash, qayta ishlash, uzatish va himoyalash usullarini o'rganuvchi fan
\item kompyuter tizimining jismoniy qurilmalari majmui
\end{enumerate}
\begin{tcolorbox}[enhanced,breakable,colback=teal!4,colframe=teal!50!black,boxrule=0.5pt,arc=1.1mm,left=3mm,right=3mm,top=1.7mm,bottom=1.7mm,before skip=2mm,after skip=2mm]
\textbf{\color{teal}To'g'ri javob:} \textbf{A) axborotni qabul qiluvchiga yetkazish uchun ifodalangan mazmun}
\par\smallskip
\textbf{Izoh:} Xabar - axborotni qabul qiluvchiga yetkazish uchun ifodalangan mazmun.
\end{tcolorbox}
\vspace{2mm}\hrule height 0.2pt \vspace{3mm}

% DB question_id: 32d2399b-e532-4982-938a-1b626eb520b3
\needspace{5\baselineskip}
\noindent\textbf{\color{ink}2.}
\par
"Ma'lumot" atamasiga eng aniq ta'rifni toping.
\begin{enumerate}[label=\Alph*),leftmargin=9mm]
\item qayd etilgan belgi, qiymat yoki kuzatuv bo'lib, kontekstsiz mazmuni yetarli bo'lmasligi mumkin
\item qabul qiluvchi uchun mazmunli bo'lib, noaniqlikni kamaytiradigan ma'lumot
\item axborotni tushunish, bog'lash va vaziyatda qo'llash natijasida shakllangan tizimli tushuncha
\item o'rganilayotgan narsa, hodisa, jarayon yoki shaxs
\end{enumerate}
\begin{tcolorbox}[enhanced,breakable,colback=teal!4,colframe=teal!50!black,boxrule=0.5pt,arc=1.1mm,left=3mm,right=3mm,top=1.7mm,bottom=1.7mm,before skip=2mm,after skip=2mm]
\textbf{\color{teal}To'g'ri javob:} \textbf{A) qayd etilgan belgi, qiymat yoki kuzatuv bo'lib, kontekstsiz mazmuni yetarli bo'lmasligi mumkin}
\par\smallskip
\textbf{Izoh:} Ma'lumot - qayd etilgan belgi, qiymat yoki kuzatuv bo'lib, kontekstsiz mazmuni yetarli bo'lmasligi mumkin.
\end{tcolorbox}
\vspace{2mm}\hrule height 0.2pt \vspace{3mm}

% DB question_id: 5253ad21-2314-4b47-a137-1426ca4d88a2
\needspace{5\baselineskip}
\noindent\textbf{\color{ink}3.}
\par
"Signal" atamasiga eng aniq ta'rifni toping.
\begin{enumerate}[label=\Alph*),leftmargin=9mm]
\item axborotni izlash, yig'ish, saqlash, qayta ishlash, uzatish va himoyalash usullarini o'rganuvchi fan
\item xabarni kanal bo'ylab tashiydigan fizik o'zgarish
\item kompyuter tizimining jismoniy qurilmalari majmui
\item kompyuter bajaradigan dasturlar, operatsion tizimlar va algoritmik vositalar majmui
\end{enumerate}
\begin{tcolorbox}[enhanced,breakable,colback=teal!4,colframe=teal!50!black,boxrule=0.5pt,arc=1.1mm,left=3mm,right=3mm,top=1.7mm,bottom=1.7mm,before skip=2mm,after skip=2mm]
\textbf{\color{teal}To'g'ri javob:} \textbf{B) xabarni kanal bo'ylab tashiydigan fizik o'zgarish}
\par\smallskip
\textbf{Izoh:} Signal - xabarni kanal bo'ylab tashiydigan fizik o'zgarish.
\end{tcolorbox}
\vspace{2mm}\hrule height 0.2pt \vspace{3mm}

% DB question_id: 55ad928a-22ea-49ec-bd47-f7ecf37e8d88
\needspace{5\baselineskip}
\noindent\textbf{\color{ink}4.}
\par
"Axborot va ma'lumot farqi" atamasiga eng aniq ta'rifni toping.
\begin{enumerate}[label=\Alph*),leftmargin=9mm]
\item kodlash taqdim etish formatini, shifrlash esa maxfiylikni ko'zlashini ajratish
\item bit ikkilik birlik, bayt esa 8 bitdan iborat ekanini ajratish
\item ma'lumot qayd etilgan qiymat, axborot esa qabul qiluvchi uchun mazmunli ma'lumot ekanini ajratish
\item rastr piksel, vektor geometrik obyektlarga asoslanishini ajratish
\end{enumerate}
\begin{tcolorbox}[enhanced,breakable,colback=teal!4,colframe=teal!50!black,boxrule=0.5pt,arc=1.1mm,left=3mm,right=3mm,top=1.7mm,bottom=1.7mm,before skip=2mm,after skip=2mm]
\textbf{\color{teal}To'g'ri javob:} \textbf{C) ma'lumot qayd etilgan qiymat, axborot esa qabul qiluvchi uchun mazmunli ma'lumot ekanini ajratish}
\par\smallskip
\textbf{Izoh:} Axborot va ma'lumot farqi - ma'lumot qayd etilgan qiymat, axborot esa qabul qiluvchi uchun mazmunli ma'lumot ekanini ajratish.
\end{tcolorbox}
\vspace{2mm}\hrule height 0.2pt \vspace{3mm}

% DB question_id: 5c9e3b70-244a-409a-ab1d-df26131c08a5
\needspace{5\baselineskip}
\noindent\textbf{\color{ink}5.}
\par
"Informatika" atamasiga eng aniq ta'rifni toping.
\begin{enumerate}[label=\Alph*),leftmargin=9mm]
\item axborotni izlash, yig'ish, saqlash, qayta ishlash, uzatish va himoyalash usullarini o'rganuvchi fan
\item kompyuter tizimining jismoniy qurilmalari majmui
\item kompyuter bajaradigan dasturlar, operatsion tizimlar va algoritmik vositalar majmui
\item tizimni yaratadigan, boshqaradigan va undan foydalanadigan insonlar hamda ularning bilimlari
\end{enumerate}
\begin{tcolorbox}[enhanced,breakable,colback=teal!4,colframe=teal!50!black,boxrule=0.5pt,arc=1.1mm,left=3mm,right=3mm,top=1.7mm,bottom=1.7mm,before skip=2mm,after skip=2mm]
\textbf{\color{teal}To'g'ri javob:} \textbf{A) axborotni izlash, yig'ish, saqlash, qayta ishlash, uzatish va himoyalash usullarini o'rganuvchi fan}
\par\smallskip
\textbf{Izoh:} Informatika - axborotni izlash, yig'ish, saqlash, qayta ishlash, uzatish va himoyalash usullarini o'rganuvchi fan.
\end{tcolorbox}
\vspace{2mm}\hrule height 0.2pt \vspace{3mm}

% DB question_id: 624d2270-e382-42d7-903c-8c10bd40ec4a
\needspace{5\baselineskip}
\noindent\textbf{\color{ink}6.}
\par
"Axborot" atamasiga eng aniq ta'rifni toping.
\begin{enumerate}[label=\Alph*),leftmargin=9mm]
\item axborotni tushunish, bog'lash va vaziyatda qo'llash natijasida shakllangan tizimli tushuncha
\item qabul qiluvchi uchun mazmunli bo'lib, noaniqlikni kamaytiradigan ma'lumot
\item o'rganilayotgan narsa, hodisa, jarayon yoki shaxs
\item axborotni qabul qiluvchiga yetkazish uchun ifodalangan mazmun
\end{enumerate}
\begin{tcolorbox}[enhanced,breakable,colback=teal!4,colframe=teal!50!black,boxrule=0.5pt,arc=1.1mm,left=3mm,right=3mm,top=1.7mm,bottom=1.7mm,before skip=2mm,after skip=2mm]
\textbf{\color{teal}To'g'ri javob:} \textbf{B) qabul qiluvchi uchun mazmunli bo'lib, noaniqlikni kamaytiradigan ma'lumot}
\par\smallskip
\textbf{Izoh:} Axborot - qabul qiluvchi uchun mazmunli bo'lib, noaniqlikni kamaytiradigan ma'lumot.
\end{tcolorbox}
\vspace{2mm}\hrule height 0.2pt \vspace{3mm}

% DB question_id: 6c624b26-5c8e-4d7d-b66b-35ec14b8aea4
\needspace{5\baselineskip}
\noindent\textbf{\color{ink}7.}
\par
"Obyekt" atamasiga eng aniq ta'rifni toping.
\begin{enumerate}[label=\Alph*),leftmargin=9mm]
\item axborotni qabul qiluvchiga yetkazish uchun ifodalangan mazmun
\item xabarni kanal bo'ylab tashiydigan fizik o'zgarish
\item axborotni izlash, yig'ish, saqlash, qayta ishlash, uzatish va himoyalash usullarini o'rganuvchi fan
\item o'rganilayotgan narsa, hodisa, jarayon yoki shaxs
\end{enumerate}
\begin{tcolorbox}[enhanced,breakable,colback=teal!4,colframe=teal!50!black,boxrule=0.5pt,arc=1.1mm,left=3mm,right=3mm,top=1.7mm,bottom=1.7mm,before skip=2mm,after skip=2mm]
\textbf{\color{teal}To'g'ri javob:} \textbf{D) o'rganilayotgan narsa, hodisa, jarayon yoki shaxs}
\par\smallskip
\textbf{Izoh:} Obyekt - o'rganilayotgan narsa, hodisa, jarayon yoki shaxs.
\end{tcolorbox}
\vspace{2mm}\hrule height 0.2pt \vspace{3mm}

% DB question_id: 80b5f18b-ae83-4949-81c8-2c7e55dfbbd8
\needspace{5\baselineskip}
\noindent\textbf{\color{ink}8.}
\par
"Insoniy ta'minot" atamasiga eng aniq ta'rifni toping.
\begin{enumerate}[label=\Alph*),leftmargin=9mm]
\item qayd etilgan belgi, qiymat yoki kuzatuv bo'lib, kontekstsiz mazmuni yetarli bo'lmasligi mumkin
\item qabul qiluvchi uchun mazmunli bo'lib, noaniqlikni kamaytiradigan ma'lumot
\item axborotni tushunish, bog'lash va vaziyatda qo'llash natijasida shakllangan tizimli tushuncha
\item tizimni yaratadigan, boshqaradigan va undan foydalanadigan insonlar hamda ularning bilimlari
\end{enumerate}
\begin{tcolorbox}[enhanced,breakable,colback=teal!4,colframe=teal!50!black,boxrule=0.5pt,arc=1.1mm,left=3mm,right=3mm,top=1.7mm,bottom=1.7mm,before skip=2mm,after skip=2mm]
\textbf{\color{teal}To'g'ri javob:} \textbf{D) tizimni yaratadigan, boshqaradigan va undan foydalanadigan insonlar hamda ularning bilimlari}
\par\smallskip
\textbf{Izoh:} Insoniy ta'minot - tizimni yaratadigan, boshqaradigan va undan foydalanadigan insonlar hamda ularning bilimlari.
\end{tcolorbox}
\vspace{2mm}\hrule height 0.2pt \vspace{3mm}

% DB question_id: d02a677b-c7ea-4c5e-a970-3e8945932334
\needspace{5\baselineskip}
\noindent\textbf{\color{ink}9.}
\par
"Bilim" atamasiga eng aniq ta'rifni toping.
\begin{enumerate}[label=\Alph*),leftmargin=9mm]
\item o'rganilayotgan narsa, hodisa, jarayon yoki shaxs
\item axborotni qabul qiluvchiga yetkazish uchun ifodalangan mazmun
\item axborotni tushunish, bog'lash va vaziyatda qo'llash natijasida shakllangan tizimli tushuncha
\item xabarni kanal bo'ylab tashiydigan fizik o'zgarish
\end{enumerate}
\begin{tcolorbox}[enhanced,breakable,colback=teal!4,colframe=teal!50!black,boxrule=0.5pt,arc=1.1mm,left=3mm,right=3mm,top=1.7mm,bottom=1.7mm,before skip=2mm,after skip=2mm]
\textbf{\color{teal}To'g'ri javob:} \textbf{C) axborotni tushunish, bog'lash va vaziyatda qo'llash natijasida shakllangan tizimli tushuncha}
\par\smallskip
\textbf{Izoh:} Bilim - axborotni tushunish, bog'lash va vaziyatda qo'llash natijasida shakllangan tizimli tushuncha.
\end{tcolorbox}
\vspace{2mm}\hrule height 0.2pt \vspace{3mm}

% DB question_id: ffdb48ac-9c84-5afd-91b8-f3fe42e9ee78
\needspace{5\baselineskip}
\noindent\textbf{\color{ink}10.}
\par
Axborot tizimi ta'minotining uch asosiy tarkibini ko'rsating.
\begin{enumerate}[label=\Alph*),leftmargin=9mm]
\item Software, Hardware, Office
\item Hardware, Freeware, Shareware
\item Freeware, Shareware, Brainware
\item Software, Hardware, Brainware
\end{enumerate}
\begin{tcolorbox}[enhanced,breakable,colback=teal!4,colframe=teal!50!black,boxrule=0.5pt,arc=1.1mm,left=3mm,right=3mm,top=1.7mm,bottom=1.7mm,before skip=2mm,after skip=2mm]
\textbf{\color{teal}To'g'ri javob:} \textbf{D) Software, Hardware, Brainware}
\par\smallskip
\textbf{Izoh:} Texnik ta'minot - Hardware, dasturiy ta'minot - Software, insoniy/aqliy ta'minot - Brainware.
\end{tcolorbox}
\vspace{2mm}\hrule height 0.2pt \vspace{3mm}

% DB question_id: 1ec82c34-b4ad-4fd7-9191-3cb53cf097dd
\needspace{5\baselineskip}
\noindent\textbf{\color{ink}11.}
\par
Quyidagi ta'rif qaysi atamaga tegishli? Axborotni tushunish, bog'lash va vaziyatda qo'llash natijasida shakllangan tizimli tushuncha.
\begin{enumerate}[label=\Alph*),leftmargin=9mm]
\item Bilim
\item Obyekt
\item Xabar
\item Signal
\end{enumerate}
\begin{tcolorbox}[enhanced,breakable,colback=teal!4,colframe=teal!50!black,boxrule=0.5pt,arc=1.1mm,left=3mm,right=3mm,top=1.7mm,bottom=1.7mm,before skip=2mm,after skip=2mm]
\textbf{\color{teal}To'g'ri javob:} \textbf{A) Bilim}
\par\smallskip
\textbf{Izoh:} Ta'rif Bilim tushunchasiga tegishli.
\end{tcolorbox}
\vspace{2mm}\hrule height 0.2pt \vspace{3mm}

% DB question_id: 2a3d9713-68be-44f8-83a8-05ca510404d0
\needspace{5\baselineskip}
\noindent\textbf{\color{ink}12.}
\par
Quyidagi ta'rif qaysi atamaga tegishli? Tizimni yaratadigan, boshqaradigan va undan foydalanadigan insonlar hamda ularning bilimlari.
\begin{enumerate}[label=\Alph*),leftmargin=9mm]
\item Ma'lumot
\item Insoniy ta'minot
\item Axborot
\item Bilim
\end{enumerate}
\begin{tcolorbox}[enhanced,breakable,colback=teal!4,colframe=teal!50!black,boxrule=0.5pt,arc=1.1mm,left=3mm,right=3mm,top=1.7mm,bottom=1.7mm,before skip=2mm,after skip=2mm]
\textbf{\color{teal}To'g'ri javob:} \textbf{B) Insoniy ta'minot}
\par\smallskip
\textbf{Izoh:} Ta'rif Insoniy ta'minot tushunchasiga tegishli.
\end{tcolorbox}
\vspace{2mm}\hrule height 0.2pt \vspace{3mm}

% DB question_id: 620092b7-47c5-4dc3-8e68-48be59cf20b4
\needspace{5\baselineskip}
\noindent\textbf{\color{ink}13.}
\par
Quyidagi ta'rif qaysi atamaga tegishli? Axborotni izlash, yig'ish, saqlash, qayta ishlash, uzatish va himoyalash usullarini o'rganuvchi fan.
\begin{enumerate}[label=\Alph*),leftmargin=9mm]
\item Texnik ta'minot
\item Dasturiy ta'minot
\item Informatika
\item Insoniy ta'minot
\end{enumerate}
\begin{tcolorbox}[enhanced,breakable,colback=teal!4,colframe=teal!50!black,boxrule=0.5pt,arc=1.1mm,left=3mm,right=3mm,top=1.7mm,bottom=1.7mm,before skip=2mm,after skip=2mm]
\textbf{\color{teal}To'g'ri javob:} \textbf{C) Informatika}
\par\smallskip
\textbf{Izoh:} Ta'rif Informatika tushunchasiga tegishli.
\end{tcolorbox}
\vspace{2mm}\hrule height 0.2pt \vspace{3mm}

% DB question_id: 6e2791b8-406a-4135-b87e-15075d305468
\needspace{5\baselineskip}
\noindent\textbf{\color{ink}14.}
\par
Quyidagi ta'rif qaysi atamaga tegishli? Xabarni kanal bo'ylab tashiydigan fizik o'zgarish.
\begin{enumerate}[label=\Alph*),leftmargin=9mm]
\item Informatika
\item Texnik ta'minot
\item Dasturiy ta'minot
\item Signal
\end{enumerate}
\begin{tcolorbox}[enhanced,breakable,colback=teal!4,colframe=teal!50!black,boxrule=0.5pt,arc=1.1mm,left=3mm,right=3mm,top=1.7mm,bottom=1.7mm,before skip=2mm,after skip=2mm]
\textbf{\color{teal}To'g'ri javob:} \textbf{D) Signal}
\par\smallskip
\textbf{Izoh:} Ta'rif Signal tushunchasiga tegishli.
\end{tcolorbox}
\vspace{2mm}\hrule height 0.2pt \vspace{3mm}

% DB question_id: 816b3d37-081c-4587-8704-8b95797b5919
\needspace{5\baselineskip}
\noindent\textbf{\color{ink}15.}
\par
Quyidagi ta'rif qaysi atamaga tegishli? Qayd etilgan belgi, qiymat yoki kuzatuv bo'lib, kontekstsiz mazmuni yetarli bo'lmasligi mumkin.
\begin{enumerate}[label=\Alph*),leftmargin=9mm]
\item Axborot
\item Bilim
\item Ma'lumot
\item Obyekt
\end{enumerate}
\begin{tcolorbox}[enhanced,breakable,colback=teal!4,colframe=teal!50!black,boxrule=0.5pt,arc=1.1mm,left=3mm,right=3mm,top=1.7mm,bottom=1.7mm,before skip=2mm,after skip=2mm]
\textbf{\color{teal}To'g'ri javob:} \textbf{C) Ma'lumot}
\par\smallskip
\textbf{Izoh:} Ta'rif Ma'lumot tushunchasiga tegishli.
\end{tcolorbox}
\vspace{2mm}\hrule height 0.2pt \vspace{3mm}

% DB question_id: 9b026331-039d-4b38-8291-a0ab21726aac
\needspace{5\baselineskip}
\noindent\textbf{\color{ink}16.}
\par
Quyidagi ta'rif qaysi atamaga tegishli? Axborotni qabul qiluvchiga yetkazish uchun ifodalangan mazmun.
\begin{enumerate}[label=\Alph*),leftmargin=9mm]
\item Signal
\item Informatika
\item Xabar
\item Texnik ta'minot
\end{enumerate}
\begin{tcolorbox}[enhanced,breakable,colback=teal!4,colframe=teal!50!black,boxrule=0.5pt,arc=1.1mm,left=3mm,right=3mm,top=1.7mm,bottom=1.7mm,before skip=2mm,after skip=2mm]
\textbf{\color{teal}To'g'ri javob:} \textbf{C) Xabar}
\par\smallskip
\textbf{Izoh:} Ta'rif Xabar tushunchasiga tegishli.
\end{tcolorbox}
\vspace{2mm}\hrule height 0.2pt \vspace{3mm}

% DB question_id: b5461d9e-9239-4805-8042-df18e7f4b129
\needspace{5\baselineskip}
\noindent\textbf{\color{ink}17.}
\par
Axborotni bilimga aylantirish jarayonlari qaysi javobda to'g'ri berilgan?
\begin{enumerate}[label=\Alph*),leftmargin=9mm]
\item to'plash, saqlash, qayta ishlash, tahlil, qo'llash
\item to'plash, tarqatish, o'chirish, baholash, qo'llash
\item saqlash, nusxalash, uzatish, filtrlash
\item yig'ish, saralash, chop etish, arxivlash, qo'llash
\end{enumerate}
\begin{tcolorbox}[enhanced,breakable,colback=teal!4,colframe=teal!50!black,boxrule=0.5pt,arc=1.1mm,left=3mm,right=3mm,top=1.7mm,bottom=1.7mm,before skip=2mm,after skip=2mm]
\textbf{\color{teal}To'g'ri javob:} \textbf{A) to'plash, saqlash, qayta ishlash, tahlil, qo'llash}
\par\smallskip
\textbf{Izoh:} Manbadagi A variantda axborot to'planadi, saqlanadi, qayta ishlanadi, tahlil qilinadi va amalda qo'llanadi. Mazmuniy tekshiruv bilan mos.
\end{tcolorbox}
\vspace{2mm}\hrule height 0.2pt \vspace{3mm}

% DB question_id: c5c49e0f-6c9b-4bdd-b7fb-b68f5e54bd35
\needspace{5\baselineskip}
\noindent\textbf{\color{ink}18.}
\par
Quyidagi ta'rif qaysi atamaga tegishli? Qabul qiluvchi uchun mazmunli bo'lib, noaniqlikni kamaytiradigan ma'lumot.
\begin{enumerate}[label=\Alph*),leftmargin=9mm]
\item Bilim
\item Obyekt
\item Xabar
\item Axborot
\end{enumerate}
\begin{tcolorbox}[enhanced,breakable,colback=teal!4,colframe=teal!50!black,boxrule=0.5pt,arc=1.1mm,left=3mm,right=3mm,top=1.7mm,bottom=1.7mm,before skip=2mm,after skip=2mm]
\textbf{\color{teal}To'g'ri javob:} \textbf{D) Axborot}
\par\smallskip
\textbf{Izoh:} Ta'rif Axborot tushunchasiga tegishli.
\end{tcolorbox}
\vspace{2mm}\hrule height 0.2pt \vspace{3mm}

% DB question_id: da63a6e3-d9b4-4589-8fc9-e0cc80327680
\needspace{5\baselineskip}
\noindent\textbf{\color{ink}19.}
\par
Quyidagi ta'rif qaysi atamaga tegishli? O'rganilayotgan narsa, hodisa, jarayon yoki shaxs.
\begin{enumerate}[label=\Alph*),leftmargin=9mm]
\item Xabar
\item Obyekt
\item Signal
\item Informatika
\end{enumerate}
\begin{tcolorbox}[enhanced,breakable,colback=teal!4,colframe=teal!50!black,boxrule=0.5pt,arc=1.1mm,left=3mm,right=3mm,top=1.7mm,bottom=1.7mm,before skip=2mm,after skip=2mm]
\textbf{\color{teal}To'g'ri javob:} \textbf{B) Obyekt}
\par\smallskip
\textbf{Izoh:} Ta'rif Obyekt tushunchasiga tegishli.
\end{tcolorbox}
\vspace{2mm}\hrule height 0.2pt \vspace{3mm}

% DB question_id: de4754f3-09cb-4259-aa84-5134b7af7e95
\needspace{5\baselineskip}
\noindent\textbf{\color{ink}20.}
\par
Quyidagi ta'rif qaysi atamaga tegishli? Ma'lumot qayd etilgan qiymat, axborot esa qabul qiluvchi uchun mazmunli ma'lumot ekanini ajratish.
\begin{enumerate}[label=\Alph*),leftmargin=9mm]
\item Axborot va ma'lumot farqi
\item Kodlash va shifrlash farqi
\item Bit va bayt farqi
\item Rastr va vektor farqi
\end{enumerate}
\begin{tcolorbox}[enhanced,breakable,colback=teal!4,colframe=teal!50!black,boxrule=0.5pt,arc=1.1mm,left=3mm,right=3mm,top=1.7mm,bottom=1.7mm,before skip=2mm,after skip=2mm]
\textbf{\color{teal}To'g'ri javob:} \textbf{A) Axborot va ma'lumot farqi}
\par\smallskip
\textbf{Izoh:} Ta'rif Axborot va ma'lumot farqi tushunchasiga tegishli.
\end{tcolorbox}
\vspace{2mm}\hrule height 0.2pt \vspace{3mm}

% DB question_id: 024c0191-5975-5432-a629-527eed778ed9
\needspace{5\baselineskip}
\noindent\textbf{\color{ink}21.}
\par
Administrator foydalanuvchilar huquqlarini sozladi va xodimlarni o'qitdi. Ushbu vaziyatdagi asosiy tushunchani aniqlang.
\begin{enumerate}[label=\Alph*),leftmargin=9mm]
\item Analog va raqamli axborot
\item Insoniy ta'minot
\item Axborotlashgan jamiyat
\item Informatika
\end{enumerate}
\begin{tcolorbox}[enhanced,breakable,colback=teal!4,colframe=teal!50!black,boxrule=0.5pt,arc=1.1mm,left=3mm,right=3mm,top=1.7mm,bottom=1.7mm,before skip=2mm,after skip=2mm]
\textbf{\color{teal}To'g'ri javob:} \textbf{B) Insoniy ta'minot}
\par\smallskip
\textbf{Izoh:} Vaziyat Insoniy ta'minot tushunchasini ifodalaydi: axborot tizimini yaratadigan, boshqaradigan va undan foydalanadigan insonlar hamda ularning malakasi.
\end{tcolorbox}
\vspace{2mm}\hrule height 0.2pt \vspace{3mm}

% DB question_id: 0c5f3e6e-d46d-51c1-be5d-81984e54e9bb
\needspace{5\baselineskip}
\noindent\textbf{\color{ink}22.}
\par
Quyidagi ta'rif qaysi tushunchaga tegishli? Axborot olinadigan yoki tasvirlanadigan narsa, shaxs, jarayon yoxud tizim.
\begin{enumerate}[label=\Alph*),leftmargin=9mm]
\item Obyekt
\item Hodisa
\item Xabar
\item Signal
\end{enumerate}
\begin{tcolorbox}[enhanced,breakable,colback=teal!4,colframe=teal!50!black,boxrule=0.5pt,arc=1.1mm,left=3mm,right=3mm,top=1.7mm,bottom=1.7mm,before skip=2mm,after skip=2mm]
\textbf{\color{teal}To'g'ri javob:} \textbf{A) Obyekt}
\par\smallskip
\textbf{Izoh:} Ta'rif Obyekt tushunchasiga tegishli.
\end{tcolorbox}
\vspace{2mm}\hrule height 0.2pt \vspace{3mm}

% DB question_id: 0fa65f53-ebcf-5bcb-8cf3-bb456d9de91e
\needspace{5\baselineskip}
\noindent\textbf{\color{ink}23.}
\par
O'quvchi ob-havo ma'lumotiga qarab issiqda ko'proq suv ichish kerakligini xulosa qildi. Ushbu vaziyatdagi asosiy tushunchani aniqlang.
\begin{enumerate}[label=\Alph*),leftmargin=9mm]
\item Ma'lumot \ensuremath{\to} axborot \ensuremath{\to} bilim zanjiri
\item Bilim
\item Bilimlar bazasi
\item Obyekt
\end{enumerate}
\begin{tcolorbox}[enhanced,breakable,colback=teal!4,colframe=teal!50!black,boxrule=0.5pt,arc=1.1mm,left=3mm,right=3mm,top=1.7mm,bottom=1.7mm,before skip=2mm,after skip=2mm]
\textbf{\color{teal}To'g'ri javob:} \textbf{B) Bilim}
\par\smallskip
\textbf{Izoh:} Vaziyat Bilim tushunchasini ifodalaydi: axborotni tushunish, tajriba bilan bog'lash va undan qaror chiqarishda foydalanish qobiliyati.
\end{tcolorbox}
\vspace{2mm}\hrule height 0.2pt \vspace{3mm}

% DB question_id: 10bd96d7-8ccf-5c10-84cc-bb23471a4e37
\needspace{5\baselineskip}
\noindent\textbf{\color{ink}24.}
\par
"Xabar" atamasining eng aniq ta'rifini toping.
\begin{enumerate}[label=\Alph*),leftmargin=9mm]
\item axborotni tashuvchi fizik o'zgarish yoki belgi ketma-ketligi
\item axborotni saqlash yoki yetkazishga xizmat qiladigan moddiy yoxud fizik vosita
\item axborotni jo'natuvchidan qabul qiluvchiga yetkazuvchi mazmunli ifoda
\item ko'rish, eshitish, hid, ta'm va tana sezgisi orqali tashqi muhitdan ma'lumot olish jarayoni
\end{enumerate}
\begin{tcolorbox}[enhanced,breakable,colback=teal!4,colframe=teal!50!black,boxrule=0.5pt,arc=1.1mm,left=3mm,right=3mm,top=1.7mm,bottom=1.7mm,before skip=2mm,after skip=2mm]
\textbf{\color{teal}To'g'ri javob:} \textbf{C) axborotni jo'natuvchidan qabul qiluvchiga yetkazuvchi mazmunli ifoda}
\par\smallskip
\textbf{Izoh:} Xabar - axborotni jo'natuvchidan qabul qiluvchiga yetkazuvchi mazmunli ifoda.
\end{tcolorbox}
\vspace{2mm}\hrule height 0.2pt \vspace{3mm}

% DB question_id: 1827c88c-8bb4-44f5-8a72-e7022efcd24c
\needspace{5\baselineskip}
\noindent\textbf{\color{ink}25.}
\par
Qog'oz shakldagi hujjat skanerlanib sonli faylga aylantirildi; u elektron imzo bilan tasdiqlanib, kompyuter tizimi orqali yuborildi. Jarayonda qaysi texnologiya qo'llangan?
\begin{enumerate}[label=\Alph*),leftmargin=9mm]
\item Analog texnologiya
\item Axborot manbasi
\item Raqamli texnologiya
\item Axborot almashinuvi
\end{enumerate}
\begin{tcolorbox}[enhanced,breakable,colback=teal!4,colframe=teal!50!black,boxrule=0.5pt,arc=1.1mm,left=3mm,right=3mm,top=1.7mm,bottom=1.7mm,before skip=2mm,after skip=2mm]
\textbf{\color{teal}To'g'ri javob:} \textbf{C) Raqamli texnologiya}
\par\smallskip
\textbf{Izoh:} Skanerlash, elektron faylni qayta ishlash va elektron imzolash raqamli texnologiya amallaridir. Axborot manbasi va axborot almashinuvi esa texnologiya turining aynan nomi emas.
\end{tcolorbox}
\vspace{2mm}\hrule height 0.2pt \vspace{3mm}

% DB question_id: 19be2fda-3a39-501c-970d-99a92c7575bd
\needspace{5\baselineskip}
\noindent\textbf{\color{ink}26.}
\par
"Dasturiy ta'minot" atamasining eng aniq ta'rifini toping.
\begin{enumerate}[label=\Alph*),leftmargin=9mm]
\item kompyuter bajaradigan dasturlar va ularni boshqaruvchi ko'rsatmalar majmui
\item axborot tizimini yaratadigan, boshqaradigan va undan foydalanadigan insonlar hamda ularning malakasi
\item analog axborot uzluksiz, raqamli axborot esa diskret qiymatlar orqali ifodalanadigan ikki xil tasvirlash shakli
\item axborot yaratish, undan foydalanish va bilim almashish iqtisodiy hamda ijtimoiy taraqqiyotning asosiy omiliga aylangan jamiyat
\end{enumerate}
\begin{tcolorbox}[enhanced,breakable,colback=teal!4,colframe=teal!50!black,boxrule=0.5pt,arc=1.1mm,left=3mm,right=3mm,top=1.7mm,bottom=1.7mm,before skip=2mm,after skip=2mm]
\textbf{\color{teal}To'g'ri javob:} \textbf{A) kompyuter bajaradigan dasturlar va ularni boshqaruvchi ko'rsatmalar majmui}
\par\smallskip
\textbf{Izoh:} Dasturiy ta'minot - kompyuter bajaradigan dasturlar va ularni boshqaruvchi ko'rsatmalar majmui.
\end{tcolorbox}
\vspace{2mm}\hrule height 0.2pt \vspace{3mm}

% DB question_id: 19c7aac8-c96f-497b-811f-c76878c08b2b
\needspace{5\baselineskip}
\noindent\textbf{\color{ink}27.}
\par
Quyidagi ta'riflardan qaysi biri <<axborot>> tushunchasini <<ma'lumot>> va <<bilim>>dan aniq farqlab beradi?
\begin{enumerate}[label=\Alph*),leftmargin=9mm]
\item Axborot - bu kompyuter xotirasida saqlanadigan har qanday raqamli fayllar majmui.
\item Axborot - bu har qanday belgilar ketma-ketligi bo'lib, tushunarliligi va foydaliligi muhim emas.
\item Axborot - bu ma'lum bir manbadan qabul qiluvchi uchun mazmun kasb etadigan ma'lumotlar to'plami.
\item Axborot - bu inson ongida shakllangan, asoslangan qarashlar majmui, ya'ni tizimlashtirilgan nazariy qarashlar.
\end{enumerate}
\begin{tcolorbox}[enhanced,breakable,colback=teal!4,colframe=teal!50!black,boxrule=0.5pt,arc=1.1mm,left=3mm,right=3mm,top=1.7mm,bottom=1.7mm,before skip=2mm,after skip=2mm]
\textbf{\color{teal}To'g'ri javob:} \textbf{C) Axborot - bu ma'lum bir manbadan qabul qiluvchi uchun mazmun kasb etadigan ma'lumotlar to'plami.}
\par\smallskip
\textbf{Izoh:} Axborot qabul qiluvchi uchun ma'no kasb etgan ma'lumotdir; bilim esa qayta ishlangan va o'zlashtirilgan axborot bilan bog'liq.
\end{tcolorbox}
\vspace{2mm}\hrule height 0.2pt \vspace{3mm}

% DB question_id: 1c0a752b-ac62-59e3-bdbc-20b1388ee048
\needspace{5\baselineskip}
\noindent\textbf{\color{ink}28.}
\par
O'qituvchi kompyuter, proyektor va onlayn aloqa vositasidan birgalikda foydalandi. Ushbu vaziyatdagi asosiy tushunchani aniqlang.
\begin{enumerate}[label=\Alph*),leftmargin=9mm]
\item Raqamli texnologiya
\item AKT
\item Ma'lumot
\item Axborot
\end{enumerate}
\begin{tcolorbox}[enhanced,breakable,colback=teal!4,colframe=teal!50!black,boxrule=0.5pt,arc=1.1mm,left=3mm,right=3mm,top=1.7mm,bottom=1.7mm,before skip=2mm,after skip=2mm]
\textbf{\color{teal}To'g'ri javob:} \textbf{B) AKT}
\par\smallskip
\textbf{Izoh:} Vaziyat AKT tushunchasini ifodalaydi: axborotni yaratish, qayta ishlash va almashishda kompyuter hamda telekommunikatsiya vositalaridan foydalanish sohasi.
\end{tcolorbox}
\vspace{2mm}\hrule height 0.2pt \vspace{3mm}

% DB question_id: 22155483-6420-55bf-9c2a-37f013401cb4
\needspace{5\baselineskip}
\noindent\textbf{\color{ink}29.}
\par
"Raqamli texnologiya" atamasining eng aniq ta'rifini toping.
\begin{enumerate}[label=\Alph*),leftmargin=9mm]
\item qayd etilgan belgi, qiymat yoki kuzatuv bo'lib, kontekstsiz mazmuni yetarli bo'lmasligi mumkin
\item qabul qiluvchi uchun mazmunli bo'lib, noaniqlikni kamaytiradigan ma'lumot
\item ma'lumotni diskret sonli ko'rinishda yaratish, qayta ishlash yoki uzatishga asoslangan texnologiya
\item axborotni tushunish, tajriba bilan bog'lash va undan qaror chiqarishda foydalanish qobiliyati
\end{enumerate}
\begin{tcolorbox}[enhanced,breakable,colback=teal!4,colframe=teal!50!black,boxrule=0.5pt,arc=1.1mm,left=3mm,right=3mm,top=1.7mm,bottom=1.7mm,before skip=2mm,after skip=2mm]
\textbf{\color{teal}To'g'ri javob:} \textbf{C) ma'lumotni diskret sonli ko'rinishda yaratish, qayta ishlash yoki uzatishga asoslangan texnologiya}
\par\smallskip
\textbf{Izoh:} Raqamli texnologiya - ma'lumotni diskret sonli ko'rinishda yaratish, qayta ishlash yoki uzatishga asoslangan texnologiya.
\end{tcolorbox}
\vspace{2mm}\hrule height 0.2pt \vspace{3mm}

% DB question_id: 2832eb3b-5d9a-5c9d-8866-991e764488a7
\needspace{5\baselineskip}
\noindent\textbf{\color{ink}30.}
\par
Quyidagi ta'rif qaysi tushunchaga tegishli? Kompyuter bajaradigan dasturlar va ularni boshqaruvchi ko'rsatmalar majmui.
\begin{enumerate}[label=\Alph*),leftmargin=9mm]
\item Dasturiy ta'minot
\item Insoniy ta'minot
\item Analog va raqamli axborot
\item Axborotlashgan jamiyat
\end{enumerate}
\begin{tcolorbox}[enhanced,breakable,colback=teal!4,colframe=teal!50!black,boxrule=0.5pt,arc=1.1mm,left=3mm,right=3mm,top=1.7mm,bottom=1.7mm,before skip=2mm,after skip=2mm]
\textbf{\color{teal}To'g'ri javob:} \textbf{A) Dasturiy ta'minot}
\par\smallskip
\textbf{Izoh:} Ta'rif Dasturiy ta'minot tushunchasiga tegishli.
\end{tcolorbox}
\vspace{2mm}\hrule height 0.2pt \vspace{3mm}

% DB question_id: 350d4307-5990-5b7a-8fd2-bb03b7574f42
\needspace{5\baselineskip}
\noindent\textbf{\color{ink}31.}
\par
"Bilim" atamasining eng aniq ta'rifini toping.
\begin{enumerate}[label=\Alph*),leftmargin=9mm]
\item xom qiymatga kontekst berilib axborot, axborot tahlil qilinib amaliy bilim hosil bo'lish jarayoni
\item axborotni tushunish, tajriba bilan bog'lash va undan qaror chiqarishda foydalanish qobiliyati
\item muayyan soha bo'yicha faktlar va qoidalar tartibli saqlanadigan tuzilma
\item axborot olinadigan yoki tasvirlanadigan narsa, shaxs, jarayon yoxud tizim
\end{enumerate}
\begin{tcolorbox}[enhanced,breakable,colback=teal!4,colframe=teal!50!black,boxrule=0.5pt,arc=1.1mm,left=3mm,right=3mm,top=1.7mm,bottom=1.7mm,before skip=2mm,after skip=2mm]
\textbf{\color{teal}To'g'ri javob:} \textbf{B) axborotni tushunish, tajriba bilan bog'lash va undan qaror chiqarishda foydalanish qobiliyati}
\par\smallskip
\textbf{Izoh:} Bilim - axborotni tushunish, tajriba bilan bog'lash va undan qaror chiqarishda foydalanish qobiliyati.
\end{tcolorbox}
\vspace{2mm}\hrule height 0.2pt \vspace{3mm}

% DB question_id: 36446ac4-368b-5cd5-a18a-4903c9b55ade
\needspace{5\baselineskip}
\noindent\textbf{\color{ink}32.}
\par
Maktab kursida axborot bilan ishlash usullari, algoritmlar va raqamli vositalar birgalikda o'rganildi. Ushbu vaziyatdagi asosiy tushunchani aniqlang.
\begin{enumerate}[label=\Alph*),leftmargin=9mm]
\item Informatika
\item AKT
\item Raqamli texnologiya
\item Ma'lumot
\end{enumerate}
\begin{tcolorbox}[enhanced,breakable,colback=teal!4,colframe=teal!50!black,boxrule=0.5pt,arc=1.1mm,left=3mm,right=3mm,top=1.7mm,bottom=1.7mm,before skip=2mm,after skip=2mm]
\textbf{\color{teal}To'g'ri javob:} \textbf{A) Informatika}
\par\smallskip
\textbf{Izoh:} Vaziyat Informatika tushunchasini ifodalaydi: axborotni yig'ish, saqlash, qayta ishlash, uzatish va undan foydalanish usullarini o'rganuvchi fan.
\end{tcolorbox}
\vspace{2mm}\hrule height 0.2pt \vspace{3mm}

% DB question_id: 368f57fc-b088-46f2-ae88-a1178d32be8e
\needspace{5\baselineskip}
\noindent\textbf{\color{ink}33.}
\par
Axborot yaratish, undan foydalanish va bilim almashish iqtisodiy hamda ijtimoiy taraqqiyotning muhim omiliga aylangan. Ushbu holat jamiyat rivojining qaysi turini ifodalaydi?
\begin{enumerate}[label=\Alph*),leftmargin=9mm]
\item Sanoatlashgan jamiyat
\item Agrar yo'nalishdagi jamiyat
\item Raqamli iqtisodiyot
\item Axborotlashgan jamiyat
\end{enumerate}
\begin{tcolorbox}[enhanced,breakable,colback=teal!4,colframe=teal!50!black,boxrule=0.5pt,arc=1.1mm,left=3mm,right=3mm,top=1.7mm,bottom=1.7mm,before skip=2mm,after skip=2mm]
\textbf{\color{teal}To'g'ri javob:} \textbf{D) Axborotlashgan jamiyat}
\par\smallskip
\textbf{Izoh:} Axborotlashgan jamiyatda axborotni yaratish, undan foydalanish va bilim almashish jamiyatning yetakchi rivojlanish omillaridan bo'ladi. Raqamli iqtisodiyot asosan iqtisodiy faoliyatning raqamli tomoniga qaratilgan.
\end{tcolorbox}
\vspace{2mm}\hrule height 0.2pt \vspace{3mm}

% DB question_id: 3a438c47-6c8b-5ce0-a43d-31c2fc4fc089
\needspace{5\baselineskip}
\noindent\textbf{\color{ink}34.}
\par
Operatsion tizim va matn muharriri yangilandi. Ushbu vaziyatdagi asosiy tushunchani aniqlang.
\begin{enumerate}[label=\Alph*),leftmargin=9mm]
\item Dasturiy ta'minot
\item Insoniy ta'minot
\item Analog va raqamli axborot
\item Axborotlashgan jamiyat
\end{enumerate}
\begin{tcolorbox}[enhanced,breakable,colback=teal!4,colframe=teal!50!black,boxrule=0.5pt,arc=1.1mm,left=3mm,right=3mm,top=1.7mm,bottom=1.7mm,before skip=2mm,after skip=2mm]
\textbf{\color{teal}To'g'ri javob:} \textbf{A) Dasturiy ta'minot}
\par\smallskip
\textbf{Izoh:} Vaziyat Dasturiy ta'minot tushunchasini ifodalaydi: kompyuter bajaradigan dasturlar va ularni boshqaruvchi ko'rsatmalar majmui.
\end{tcolorbox}
\vspace{2mm}\hrule height 0.2pt \vspace{3mm}

% DB question_id: 3bb6604f-41be-57f6-983d-f84d857427fe
\needspace{5\baselineskip}
\noindent\textbf{\color{ink}35.}
\par
Telefon ekraniga "Dars 9:00 da boshlanadi" yozuvi keldi. Ushbu vaziyatdagi asosiy tushunchani aniqlang.
\begin{enumerate}[label=\Alph*),leftmargin=9mm]
\item Signal
\item Axborot tashuvchi
\item Xabar
\item Insonning axborotni qabul qilishi
\end{enumerate}
\begin{tcolorbox}[enhanced,breakable,colback=teal!4,colframe=teal!50!black,boxrule=0.5pt,arc=1.1mm,left=3mm,right=3mm,top=1.7mm,bottom=1.7mm,before skip=2mm,after skip=2mm]
\textbf{\color{teal}To'g'ri javob:} \textbf{C) Xabar}
\par\smallskip
\textbf{Izoh:} Vaziyat Xabar tushunchasini ifodalaydi: axborotni jo'natuvchidan qabul qiluvchiga yetkazuvchi mazmunli ifoda.
\end{tcolorbox}
\vspace{2mm}\hrule height 0.2pt \vspace{3mm}

% DB question_id: 3d3d1c52-e3fe-4402-ba29-30b33bedd1bc
\needspace{5\baselineskip}
\noindent\textbf{\color{ink}36.}
\par
O'quvchi <<axborotni kompyuter va kommunikatsiya vositalarida izlash, to'plash, saqlash, qayta ishlash va uzatishning nazariy hamda amaliy usullarini o'rganadigan fan>> ta'rifini o'qidi. Qaysi fan nazarda tutilgan?
\begin{enumerate}[label=\Alph*),leftmargin=9mm]
\item Informatika - axborot jarayonlari, algoritm va ma'lumotlarni o'rganuvchi fan
\item Kibernetika - tizimlarni boshqarish va qayta aloqa tamoyillarini o'rganuvchi fan
\item Telekommunikatsiya - aloqa tarmoqlari orqali signallarni uzatish sohasi
\item Kompyuter injiniringi - kompyuter apparat qismlarini yaratish va sinash sohasi
\end{enumerate}
\begin{tcolorbox}[enhanced,breakable,colback=teal!4,colframe=teal!50!black,boxrule=0.5pt,arc=1.1mm,left=3mm,right=3mm,top=1.7mm,bottom=1.7mm,before skip=2mm,after skip=2mm]
\textbf{\color{teal}To'g'ri javob:} \textbf{A) Informatika - axborot jarayonlari, algoritm va ma'lumotlarni o'rganuvchi fan}
\par\smallskip
\textbf{Izoh:} Informatika axborot, uni ifodalash va qayta ishlash jarayonlari, algoritm, axborot tizimlari va texnologiyalarini o'rganadigan fan. Elektr kabellari yoki qurilma mexanikasi alohida tor sohaga tegishli.
\end{tcolorbox}
\vspace{2mm}\hrule height 0.2pt \vspace{3mm}

% DB question_id: 443b3926-daa3-4927-bfce-a696dbe7ea46
\needspace{5\baselineskip}
\noindent\textbf{\color{ink}37.}
\par
Jadval katagida faqat 42 soni yozilgan, lekin u nimani anglatishi ko'rsatilmagan. Bu vaziyatdagi asosiy tushunchani aniqlang.
\begin{enumerate}[label=\Alph*),leftmargin=9mm]
\item Ma'lumot
\item Axborot
\item Bilim
\item Obyekt
\end{enumerate}
\begin{tcolorbox}[enhanced,breakable,colback=teal!4,colframe=teal!50!black,boxrule=0.5pt,arc=1.1mm,left=3mm,right=3mm,top=1.7mm,bottom=1.7mm,before skip=2mm,after skip=2mm]
\textbf{\color{teal}To'g'ri javob:} \textbf{A) Ma'lumot}
\par\smallskip
\textbf{Izoh:} Vaziyat Ma'lumotga misol: qayd etilgan belgi, qiymat yoki kuzatuv bo'lib, kontekstsiz mazmuni yetarli bo'lmasligi mumkin.
\end{tcolorbox}
\vspace{2mm}\hrule height 0.2pt \vspace{3mm}

% DB question_id: 44bfe025-1933-4055-b7bd-e98c6ad3c753
\needspace{5\baselineskip}
\noindent\textbf{\color{ink}38.}
\par
Administrator foydalanuvchilarga ruxsat berdi va tizim ishini nazorat qildi. Bu vaziyatdagi asosiy tushunchani aniqlang.
\begin{enumerate}[label=\Alph*),leftmargin=9mm]
\item Ma'lumot
\item Axborot
\item Bilim
\item Insoniy ta'minot
\end{enumerate}
\begin{tcolorbox}[enhanced,breakable,colback=teal!4,colframe=teal!50!black,boxrule=0.5pt,arc=1.1mm,left=3mm,right=3mm,top=1.7mm,bottom=1.7mm,before skip=2mm,after skip=2mm]
\textbf{\color{teal}To'g'ri javob:} \textbf{D) Insoniy ta'minot}
\par\smallskip
\textbf{Izoh:} Vaziyat Insoniy ta'minotga misol: tizimni yaratadigan, boshqaradigan va undan foydalanadigan insonlar hamda ularning bilimlari.
\end{tcolorbox}
\vspace{2mm}\hrule height 0.2pt \vspace{3mm}

% DB question_id: 472f0874-056c-5443-a7be-921d14ea5f0c
\needspace{5\baselineskip}
\noindent\textbf{\color{ink}39.}
\par
"Axborot" atamasining eng aniq ta'rifini toping.
\begin{enumerate}[label=\Alph*),leftmargin=9mm]
\item qabul qiluvchi uchun mazmunli bo'lib, noaniqlikni kamaytiradigan ma'lumot
\item axborotni tushunish, tajriba bilan bog'lash va undan qaror chiqarishda foydalanish qobiliyati
\item xom qiymatga kontekst berilib axborot, axborot tahlil qilinib amaliy bilim hosil bo'lish jarayoni
\item muayyan soha bo'yicha faktlar va qoidalar tartibli saqlanadigan tuzilma
\end{enumerate}
\begin{tcolorbox}[enhanced,breakable,colback=teal!4,colframe=teal!50!black,boxrule=0.5pt,arc=1.1mm,left=3mm,right=3mm,top=1.7mm,bottom=1.7mm,before skip=2mm,after skip=2mm]
\textbf{\color{teal}To'g'ri javob:} \textbf{A) qabul qiluvchi uchun mazmunli bo'lib, noaniqlikni kamaytiradigan ma'lumot}
\par\smallskip
\textbf{Izoh:} Axborot - qabul qiluvchi uchun mazmunli bo'lib, noaniqlikni kamaytiradigan ma'lumot.
\end{tcolorbox}
\vspace{2mm}\hrule height 0.2pt \vspace{3mm}

% DB question_id: 4c35dedd-001f-5ea0-8d29-bcde9a074750
\needspace{5\baselineskip}
\noindent\textbf{\color{ink}40.}
\par
"Signal" atamasining eng aniq ta'rifini toping.
\begin{enumerate}[label=\Alph*),leftmargin=9mm]
\item axborotni saqlash yoki yetkazishga xizmat qiladigan moddiy yoxud fizik vosita
\item ko'rish, eshitish, hid, ta'm va tana sezgisi orqali tashqi muhitdan ma'lumot olish jarayoni
\item kompyuter tizimida ma'lumot avval kiritilib, qayta ishlanib, natija chiqariladigan asosiy model
\item axborotni tashuvchi fizik o'zgarish yoki belgi ketma-ketligi
\end{enumerate}
\begin{tcolorbox}[enhanced,breakable,colback=teal!4,colframe=teal!50!black,boxrule=0.5pt,arc=1.1mm,left=3mm,right=3mm,top=1.7mm,bottom=1.7mm,before skip=2mm,after skip=2mm]
\textbf{\color{teal}To'g'ri javob:} \textbf{D) axborotni tashuvchi fizik o'zgarish yoki belgi ketma-ketligi}
\par\smallskip
\textbf{Izoh:} Signal - axborotni tashuvchi fizik o'zgarish yoki belgi ketma-ketligi.
\end{tcolorbox}
\vspace{2mm}\hrule height 0.2pt \vspace{3mm}

% DB question_id: 4d266d2f-aeed-5362-9cac-7df1c8d3ab1a
\needspace{5\baselineskip}
\noindent\textbf{\color{ink}41.}
\par
"Ma'lumot \ensuremath{\to} axborot \ensuremath{\to} bilim zanjiri" atamasining eng aniq ta'rifini toping.
\begin{enumerate}[label=\Alph*),leftmargin=9mm]
\item muayyan soha bo'yicha faktlar va qoidalar tartibli saqlanadigan tuzilma
\item axborot olinadigan yoki tasvirlanadigan narsa, shaxs, jarayon yoxud tizim
\item xom qiymatga kontekst berilib axborot, axborot tahlil qilinib amaliy bilim hosil bo'lish jarayoni
\item muayyan vaqt va sharoitda sodir bo'ladigan voqea yoki o'zgarish
\end{enumerate}
\begin{tcolorbox}[enhanced,breakable,colback=teal!4,colframe=teal!50!black,boxrule=0.5pt,arc=1.1mm,left=3mm,right=3mm,top=1.7mm,bottom=1.7mm,before skip=2mm,after skip=2mm]
\textbf{\color{teal}To'g'ri javob:} \textbf{C) xom qiymatga kontekst berilib axborot, axborot tahlil qilinib amaliy bilim hosil bo'lish jarayoni}
\par\smallskip
\textbf{Izoh:} Ma'lumot \ensuremath{\to} axborot \ensuremath{\to} bilim zanjiri - xom qiymatga kontekst berilib axborot, axborot tahlil qilinib amaliy bilim hosil bo'lish jarayoni.
\end{tcolorbox}
\vspace{2mm}\hrule height 0.2pt \vspace{3mm}

% DB question_id: 506ec704-62b8-51bc-8288-0ec940633bc7
\needspace{5\baselineskip}
\noindent\textbf{\color{ink}42.}
\par
Quyidagi ta'rif qaysi tushunchaga tegishli? Axborotni yig'ish, saqlash, qayta ishlash, uzatish va undan foydalanish usullarini o'rganuvchi fan.
\begin{enumerate}[label=\Alph*),leftmargin=9mm]
\item Informatika
\item AKT
\item Raqamli texnologiya
\item Ma'lumot
\end{enumerate}
\begin{tcolorbox}[enhanced,breakable,colback=teal!4,colframe=teal!50!black,boxrule=0.5pt,arc=1.1mm,left=3mm,right=3mm,top=1.7mm,bottom=1.7mm,before skip=2mm,after skip=2mm]
\textbf{\color{teal}To'g'ri javob:} \textbf{A) Informatika}
\par\smallskip
\textbf{Izoh:} Ta'rif Informatika tushunchasiga tegishli.
\end{tcolorbox}
\vspace{2mm}\hrule height 0.2pt \vspace{3mm}

% DB question_id: 52d2f764-c113-57b8-93f1-623dcd952dc1
\needspace{5\baselineskip}
\noindent\textbf{\color{ink}43.}
\par
Quyidagi ta'rif qaysi tushunchaga tegishli? Axborotni tushunish, tajriba bilan bog'lash va undan qaror chiqarishda foydalanish qobiliyati.
\begin{enumerate}[label=\Alph*),leftmargin=9mm]
\item Ma'lumot \ensuremath{\to} axborot \ensuremath{\to} bilim zanjiri
\item Bilim
\item Bilimlar bazasi
\item Obyekt
\end{enumerate}
\begin{tcolorbox}[enhanced,breakable,colback=teal!4,colframe=teal!50!black,boxrule=0.5pt,arc=1.1mm,left=3mm,right=3mm,top=1.7mm,bottom=1.7mm,before skip=2mm,after skip=2mm]
\textbf{\color{teal}To'g'ri javob:} \textbf{B) Bilim}
\par\smallskip
\textbf{Izoh:} Ta'rif Bilim tushunchasiga tegishli.
\end{tcolorbox}
\vspace{2mm}\hrule height 0.2pt \vspace{3mm}

% DB question_id: 5c36908a-d347-5b91-9c8d-aebc228daf07
\needspace{5\baselineskip}
\noindent\textbf{\color{ink}44.}
\par
Maktab kutubxonasi haqidagi ma'lumotlar yig'ildi. Ushbu vaziyatdagi asosiy tushunchani aniqlang.
\begin{enumerate}[label=\Alph*),leftmargin=9mm]
\item Obyekt
\item Hodisa
\item Xabar
\item Signal
\end{enumerate}
\begin{tcolorbox}[enhanced,breakable,colback=teal!4,colframe=teal!50!black,boxrule=0.5pt,arc=1.1mm,left=3mm,right=3mm,top=1.7mm,bottom=1.7mm,before skip=2mm,after skip=2mm]
\textbf{\color{teal}To'g'ri javob:} \textbf{A) Obyekt}
\par\smallskip
\textbf{Izoh:} Vaziyat Obyekt tushunchasini ifodalaydi: axborot olinadigan yoki tasvirlanadigan narsa, shaxs, jarayon yoxud tizim.
\end{tcolorbox}
\vspace{2mm}\hrule height 0.2pt \vspace{3mm}

% DB question_id: 5e6dec2e-a992-4e1f-babe-5e9a2b80ccb8
\needspace{5\baselineskip}
\noindent\textbf{\color{ink}45.}
\par
Do'kon terminali shtrix-kodni o'qidi, dastur mahsulot narxini hisoblab chiqdi va ekran natijani ko'rsatdi. Kompyuter faoliyatining qaysi umumiy modeli aks etgan?
\begin{enumerate}[label=\Alph*),leftmargin=9mm]
\item Saqlash-qidirish-tiklash ish modeli
\item Kiritish-kodlash-saqlash ish modeli
\item Kiritish-qayta ishlash-chiqarish modeli
\item Uzatish-qabul qilish-tasdiqlash ish modeli
\end{enumerate}
\begin{tcolorbox}[enhanced,breakable,colback=teal!4,colframe=teal!50!black,boxrule=0.5pt,arc=1.1mm,left=3mm,right=3mm,top=1.7mm,bottom=1.7mm,before skip=2mm,after skip=2mm]
\textbf{\color{teal}To'g'ri javob:} \textbf{C) Kiritish-qayta ishlash-chiqarish modeli}
\par\smallskip
\textbf{Izoh:} Shtrix-kodni o'qish - kiritish, narxni hisoblash - qayta ishlash, ekranda natijani ko'rsatish - chiqarish. Demak, kiritish-qayta ishlash-chiqarish modeli.
\end{tcolorbox}
\vspace{2mm}\hrule height 0.2pt \vspace{3mm}

% DB question_id: 5ec1a774-385a-53fc-93fa-cd5436d45248
\needspace{5\baselineskip}
\noindent\textbf{\color{ink}46.}
\par
"Ma'lumot" atamasining eng aniq ta'rifini toping.
\begin{enumerate}[label=\Alph*),leftmargin=9mm]
\item qabul qiluvchi uchun mazmunli bo'lib, noaniqlikni kamaytiradigan ma'lumot
\item axborotni tushunish, tajriba bilan bog'lash va undan qaror chiqarishda foydalanish qobiliyati
\item xom qiymatga kontekst berilib axborot, axborot tahlil qilinib amaliy bilim hosil bo'lish jarayoni
\item qayd etilgan belgi, qiymat yoki kuzatuv bo'lib, kontekstsiz mazmuni yetarli bo'lmasligi mumkin
\end{enumerate}
\begin{tcolorbox}[enhanced,breakable,colback=teal!4,colframe=teal!50!black,boxrule=0.5pt,arc=1.1mm,left=3mm,right=3mm,top=1.7mm,bottom=1.7mm,before skip=2mm,after skip=2mm]
\textbf{\color{teal}To'g'ri javob:} \textbf{D) qayd etilgan belgi, qiymat yoki kuzatuv bo'lib, kontekstsiz mazmuni yetarli bo'lmasligi mumkin}
\par\smallskip
\textbf{Izoh:} Ma'lumot - qayd etilgan belgi, qiymat yoki kuzatuv bo'lib, kontekstsiz mazmuni yetarli bo'lmasligi mumkin.
\end{tcolorbox}
\vspace{2mm}\hrule height 0.2pt \vspace{3mm}

% DB question_id: 601cb78b-c7c9-4de0-8d32-074dd55a4ab2
\needspace{5\baselineskip}
\noindent\textbf{\color{ink}47.}
\par
O'qituvchi ma'lumotni saqlash, qayta ishlash va uzatishning umumiy qonuniyatlarini o'rganmoqda. Bu vaziyatdagi asosiy tushunchani aniqlang.
\begin{enumerate}[label=\Alph*),leftmargin=9mm]
\item Informatika
\item Texnik ta'minot
\item Dasturiy ta'minot
\item Insoniy ta'minot
\end{enumerate}
\begin{tcolorbox}[enhanced,breakable,colback=teal!4,colframe=teal!50!black,boxrule=0.5pt,arc=1.1mm,left=3mm,right=3mm,top=1.7mm,bottom=1.7mm,before skip=2mm,after skip=2mm]
\textbf{\color{teal}To'g'ri javob:} \textbf{A) Informatika}
\par\smallskip
\textbf{Izoh:} Vaziyat Informatikaga misol: axborotni izlash, yig'ish, saqlash, qayta ishlash, uzatish va himoyalash usullarini o'rganuvchi fan.
\end{tcolorbox}
\vspace{2mm}\hrule height 0.2pt \vspace{3mm}

% DB question_id: 6355acc3-d592-5db6-b1e3-9448d559fdda
\needspace{5\baselineskip}
\noindent\textbf{\color{ink}48.}
\par
Matn qog'ozda bosilib, keyin arxivga qo'yildi. Ushbu vaziyatdagi asosiy tushunchani aniqlang.
\begin{enumerate}[label=\Alph*),leftmargin=9mm]
\item Axborot tashuvchi
\item Insonning axborotni qabul qilishi
\item Kiritish-qayta ishlash-chiqarish modeli
\item Texnik ta'minot
\end{enumerate}
\begin{tcolorbox}[enhanced,breakable,colback=teal!4,colframe=teal!50!black,boxrule=0.5pt,arc=1.1mm,left=3mm,right=3mm,top=1.7mm,bottom=1.7mm,before skip=2mm,after skip=2mm]
\textbf{\color{teal}To'g'ri javob:} \textbf{A) Axborot tashuvchi}
\par\smallskip
\textbf{Izoh:} Vaziyat Axborot tashuvchi tushunchasini ifodalaydi: axborotni saqlash yoki yetkazishga xizmat qiladigan moddiy yoxud fizik vosita.
\end{tcolorbox}
\vspace{2mm}\hrule height 0.2pt \vspace{3mm}

% DB question_id: 64105957-f49b-5682-b9fd-1b4b353dff50
\needspace{5\baselineskip}
\noindent\textbf{\color{ink}49.}
\par
Quyidagi ta'rif qaysi tushunchaga tegishli? Muayyan vaqt va sharoitda sodir bo'ladigan voqea yoki o'zgarish.
\begin{enumerate}[label=\Alph*),leftmargin=9mm]
\item Xabar
\item Hodisa
\item Signal
\item Axborot tashuvchi
\end{enumerate}
\begin{tcolorbox}[enhanced,breakable,colback=teal!4,colframe=teal!50!black,boxrule=0.5pt,arc=1.1mm,left=3mm,right=3mm,top=1.7mm,bottom=1.7mm,before skip=2mm,after skip=2mm]
\textbf{\color{teal}To'g'ri javob:} \textbf{B) Hodisa}
\par\smallskip
\textbf{Izoh:} Ta'rif Hodisa tushunchasiga tegishli.
\end{tcolorbox}
\vspace{2mm}\hrule height 0.2pt \vspace{3mm}

% DB question_id: 6533c3e0-c87b-55bf-87fa-9e99b5eb7795
\needspace{5\baselineskip}
\noindent\textbf{\color{ink}50.}
\par
Quyidagi ta'rif qaysi tushunchaga tegishli? Qayd etilgan belgi, qiymat yoki kuzatuv bo'lib, kontekstsiz mazmuni yetarli bo'lmasligi mumkin.
\begin{enumerate}[label=\Alph*),leftmargin=9mm]
\item Axborot
\item Bilim
\item Ma'lumot \ensuremath{\to} axborot \ensuremath{\to} bilim zanjiri
\item Ma'lumot
\end{enumerate}
\begin{tcolorbox}[enhanced,breakable,colback=teal!4,colframe=teal!50!black,boxrule=0.5pt,arc=1.1mm,left=3mm,right=3mm,top=1.7mm,bottom=1.7mm,before skip=2mm,after skip=2mm]
\textbf{\color{teal}To'g'ri javob:} \textbf{D) Ma'lumot}
\par\smallskip
\textbf{Izoh:} Ta'rif Ma'lumot tushunchasiga tegishli.
\end{tcolorbox}
\vspace{2mm}\hrule height 0.2pt \vspace{3mm}

% DB question_id: 68825207-5be9-56f3-8916-1f0ddd3ca24d
\needspace{5\baselineskip}
\noindent\textbf{\color{ink}51.}
\par
"Kiritish-qayta ishlash-chiqarish modeli" atamasining eng aniq ta'rifini toping.
\begin{enumerate}[label=\Alph*),leftmargin=9mm]
\item kompyuter tizimining barcha jismoniy qurilmalari majmui
\item kompyuter bajaradigan dasturlar va ularni boshqaruvchi ko'rsatmalar majmui
\item kompyuter tizimida ma'lumot avval kiritilib, qayta ishlanib, natija chiqariladigan asosiy model
\item axborot tizimini yaratadigan, boshqaradigan va undan foydalanadigan insonlar hamda ularning malakasi
\end{enumerate}
\begin{tcolorbox}[enhanced,breakable,colback=teal!4,colframe=teal!50!black,boxrule=0.5pt,arc=1.1mm,left=3mm,right=3mm,top=1.7mm,bottom=1.7mm,before skip=2mm,after skip=2mm]
\textbf{\color{teal}To'g'ri javob:} \textbf{C) kompyuter tizimida ma'lumot avval kiritilib, qayta ishlanib, natija chiqariladigan asosiy model}
\par\smallskip
\textbf{Izoh:} Kiritish-qayta ishlash-chiqarish modeli - kompyuter tizimida ma'lumot avval kiritilib, qayta ishlanib, natija chiqariladigan asosiy model.
\end{tcolorbox}
\vspace{2mm}\hrule height 0.2pt \vspace{3mm}

% DB question_id: 6a75c605-29b6-5c44-8bbb-d02d69de0798
\needspace{5\baselineskip}
\noindent\textbf{\color{ink}52.}
\par
"Analog va raqamli axborot" atamasining eng aniq ta'rifini toping.
\begin{enumerate}[label=\Alph*),leftmargin=9mm]
\item axborot yaratish, undan foydalanish va bilim almashish iqtisodiy hamda ijtimoiy taraqqiyotning asosiy omiliga aylangan jamiyat
\item axborotni yig'ish, saqlash, qayta ishlash, uzatish va undan foydalanish usullarini o'rganuvchi fan
\item analog axborot uzluksiz, raqamli axborot esa diskret qiymatlar orqali ifodalanadigan ikki xil tasvirlash shakli
\item axborotni yaratish, qayta ishlash va almashishda kompyuter hamda telekommunikatsiya vositalaridan foydalanish sohasi
\end{enumerate}
\begin{tcolorbox}[enhanced,breakable,colback=teal!4,colframe=teal!50!black,boxrule=0.5pt,arc=1.1mm,left=3mm,right=3mm,top=1.7mm,bottom=1.7mm,before skip=2mm,after skip=2mm]
\textbf{\color{teal}To'g'ri javob:} \textbf{C) analog axborot uzluksiz, raqamli axborot esa diskret qiymatlar orqali ifodalanadigan ikki xil tasvirlash shakli}
\par\smallskip
\textbf{Izoh:} Analog va raqamli axborot - analog axborot uzluksiz, raqamli axborot esa diskret qiymatlar orqali ifodalanadigan ikki xil tasvirlash shakli.
\end{tcolorbox}
\vspace{2mm}\hrule height 0.2pt \vspace{3mm}

% DB question_id: 6be2783c-3389-564a-910e-4f7b3b3e25d3
\needspace{5\baselineskip}
\noindent\textbf{\color{ink}53.}
\par
"Bugungi havo harorati 37,5 \ensuremath{^\circ}C" xabari foydalanuvchiga aniq mazmun berdi. Ushbu vaziyatdagi asosiy tushunchani aniqlang.
\begin{enumerate}[label=\Alph*),leftmargin=9mm]
\item Axborot
\item Bilim
\item Ma'lumot \ensuremath{\to} axborot \ensuremath{\to} bilim zanjiri
\item Bilimlar bazasi
\end{enumerate}
\begin{tcolorbox}[enhanced,breakable,colback=teal!4,colframe=teal!50!black,boxrule=0.5pt,arc=1.1mm,left=3mm,right=3mm,top=1.7mm,bottom=1.7mm,before skip=2mm,after skip=2mm]
\textbf{\color{teal}To'g'ri javob:} \textbf{A) Axborot}
\par\smallskip
\textbf{Izoh:} Vaziyat Axborot tushunchasini ifodalaydi: qabul qiluvchi uchun mazmunli bo'lib, noaniqlikni kamaytiradigan ma'lumot.
\end{tcolorbox}
\vspace{2mm}\hrule height 0.2pt \vspace{3mm}

% DB question_id: 6e455b6c-8143-4de4-a50d-db7e0d2e2fe6
\needspace{5\baselineskip}
\noindent\textbf{\color{ink}54.}
\par
Sensor qiymati jadvalga yozildi, izohlandi va keyin qurilmani o'chirish qarori chiqarildi. Ushbu vaziyatdagi asosiy tushunchani aniqlang.
\begin{enumerate}[label=\Alph*),leftmargin=9mm]
\item Ma'lumot \ensuremath{\to} kodlash \ensuremath{\to} uzatish zanjiri
\item Sensor \ensuremath{\to} aktuator \ensuremath{\to} signal zanjiri
\item Ma'lumot \ensuremath{\to} axborot \ensuremath{\to} bilim zanjiri
\item O'lchash \ensuremath{\to} filtrlash \ensuremath{\to} tasvirlash jarayoni
\end{enumerate}
\begin{tcolorbox}[enhanced,breakable,colback=teal!4,colframe=teal!50!black,boxrule=0.5pt,arc=1.1mm,left=3mm,right=3mm,top=1.7mm,bottom=1.7mm,before skip=2mm,after skip=2mm]
\textbf{\color{teal}To'g'ri javob:} \textbf{C) Ma'lumot \ensuremath{\to} axborot \ensuremath{\to} bilim zanjiri}
\par\smallskip
\textbf{Izoh:} Sensor qiymati kontekstsiz xom ma'lumotdir. Uni jadvalda izohlash axborotga, axborotdan xulosa chiqarib qurilmani boshqarish qarori esa amaliy bilimga o'tishni bildiradi. Boshqa ketma-ketliklar texnik signallarni qayta ishlash bosqichlarini tavsiflaydi.
\end{tcolorbox}
\vspace{2mm}\hrule height 0.2pt \vspace{3mm}

% DB question_id: 7724aebe-0037-5ebb-a5c1-c06c2d019844
\needspace{5\baselineskip}
\noindent\textbf{\color{ink}55.}
\par
"AKT" atamasining eng aniq ta'rifini toping.
\begin{enumerate}[label=\Alph*),leftmargin=9mm]
\item ma'lumotni diskret sonli ko'rinishda yaratish, qayta ishlash yoki uzatishga asoslangan texnologiya
\item axborotni yaratish, qayta ishlash va almashishda kompyuter hamda telekommunikatsiya vositalaridan foydalanish sohasi
\item qayd etilgan belgi, qiymat yoki kuzatuv bo'lib, kontekstsiz mazmuni yetarli bo'lmasligi mumkin
\item qabul qiluvchi uchun mazmunli bo'lib, noaniqlikni kamaytiradigan ma'lumot
\end{enumerate}
\begin{tcolorbox}[enhanced,breakable,colback=teal!4,colframe=teal!50!black,boxrule=0.5pt,arc=1.1mm,left=3mm,right=3mm,top=1.7mm,bottom=1.7mm,before skip=2mm,after skip=2mm]
\textbf{\color{teal}To'g'ri javob:} \textbf{B) axborotni yaratish, qayta ishlash va almashishda kompyuter hamda telekommunikatsiya vositalaridan foydalanish sohasi}
\par\smallskip
\textbf{Izoh:} AKT - axborotni yaratish, qayta ishlash va almashishda kompyuter hamda telekommunikatsiya vositalaridan foydalanish sohasi.
\end{tcolorbox}
\vspace{2mm}\hrule height 0.2pt \vspace{3mm}

% DB question_id: 78183ae8-951b-42cd-9d96-08e3e750b56f
\needspace{5\baselineskip}
\noindent\textbf{\color{ink}56.}
\par
Davlat xizmatlari, maktab ta'limi, savdo va kundalik muloqotda axborotni yaratish, undan foydalanish va bilim almashish jamiyat taraqqiyotining asosiy omiliga aylandi. Bu qaysi jamiyat turiga xos?
\begin{enumerate}[label=\Alph*),leftmargin=9mm]
\item Sanoatlashgan jamiyat
\item Raqamli iqtisodiyot
\item Agrar yo'nalishdagi jamiyat
\item Axborotlashgan jamiyat
\end{enumerate}
\begin{tcolorbox}[enhanced,breakable,colback=teal!4,colframe=teal!50!black,boxrule=0.5pt,arc=1.1mm,left=3mm,right=3mm,top=1.7mm,bottom=1.7mm,before skip=2mm,after skip=2mm]
\textbf{\color{teal}To'g'ri javob:} \textbf{D) Axborotlashgan jamiyat}
\par\smallskip
\textbf{Izoh:} Axborotlashgan jamiyatda axborotni yaratish, undan foydalanish va bilim almashish keng ijtimoiy-iqtisodiy faoliyatning yetakchi omiliga aylanadi. Raqamli iqtisodiyot bu jarayonning iqtisodiy qismiga e'tibor beradi.
\end{tcolorbox}
\vspace{2mm}\hrule height 0.2pt \vspace{3mm}

% DB question_id: 788862d9-f896-5dff-9e42-5d67a30a1bbf
\needspace{5\baselineskip}
\noindent\textbf{\color{ink}57.}
\par
O'quvchi diagrammani ko'rdi, ogohlantirish tovushini eshitdi va issiqlikni sezdi. Ushbu vaziyatdagi asosiy tushunchani aniqlang.
\begin{enumerate}[label=\Alph*),leftmargin=9mm]
\item Kiritish-qayta ishlash-chiqarish modeli
\item Insonning axborotni qabul qilishi
\item Texnik ta'minot
\item Dasturiy ta'minot
\end{enumerate}
\begin{tcolorbox}[enhanced,breakable,colback=teal!4,colframe=teal!50!black,boxrule=0.5pt,arc=1.1mm,left=3mm,right=3mm,top=1.7mm,bottom=1.7mm,before skip=2mm,after skip=2mm]
\textbf{\color{teal}To'g'ri javob:} \textbf{B) Insonning axborotni qabul qilishi}
\par\smallskip
\textbf{Izoh:} Vaziyat Insonning axborotni qabul qilishi tushunchasini ifodalaydi: ko'rish, eshitish, hid, ta'm va tana sezgisi orqali tashqi muhitdan ma'lumot olish jarayoni.
\end{tcolorbox}
\vspace{2mm}\hrule height 0.2pt \vspace{3mm}

% DB question_id: 7d1600a7-9555-5f98-9761-b73a7f960ab9
\needspace{5\baselineskip}
\noindent\textbf{\color{ink}58.}
\par
Quyidagi ta'rif qaysi tushunchaga tegishli? Axborotni jo'natuvchidan qabul qiluvchiga yetkazuvchi mazmunli ifoda.
\begin{enumerate}[label=\Alph*),leftmargin=9mm]
\item Signal
\item Axborot tashuvchi
\item Xabar
\item Insonning axborotni qabul qilishi
\end{enumerate}
\begin{tcolorbox}[enhanced,breakable,colback=teal!4,colframe=teal!50!black,boxrule=0.5pt,arc=1.1mm,left=3mm,right=3mm,top=1.7mm,bottom=1.7mm,before skip=2mm,after skip=2mm]
\textbf{\color{teal}To'g'ri javob:} \textbf{C) Xabar}
\par\smallskip
\textbf{Izoh:} Ta'rif Xabar tushunchasiga tegishli.
\end{tcolorbox}
\vspace{2mm}\hrule height 0.2pt \vspace{3mm}

% DB question_id: 856c8502-1734-480d-a7ad-3313f53bc54e
\needspace{5\baselineskip}
\noindent\textbf{\color{ink}59.}
\par
Operator klaviaturadan sonlarni kiritdi, dastur ularning yig'indisini hisoblab, natijani monitorga chiqardi. Bu uch bosqich kompyuterning qaysi umumiy modelini ko'rsatadi?
\begin{enumerate}[label=\Alph*),leftmargin=9mm]
\item Saqlash-qidirish-tiklash ish modeli
\item Yaratish-sinash-tarqatish ish modeli
\item Kiritish-qayta ishlash-chiqarish modeli
\item Uzatish-qabul qilish-tasdiqlash modeli
\end{enumerate}
\begin{tcolorbox}[enhanced,breakable,colback=teal!4,colframe=teal!50!black,boxrule=0.5pt,arc=1.1mm,left=3mm,right=3mm,top=1.7mm,bottom=1.7mm,before skip=2mm,after skip=2mm]
\textbf{\color{teal}To'g'ri javob:} \textbf{C) Kiritish-qayta ishlash-chiqarish modeli}
\par\smallskip
\textbf{Izoh:} Klaviaturadan ma'lumot olish kiritish, yig'indini hisoblash qayta ishlash, monitor natijasi chiqarish bosqichidir. Bular kiritish-qayta ishlash-chiqarish modelini tashkil qiladi.
\end{tcolorbox}
\vspace{2mm}\hrule height 0.2pt \vspace{3mm}

% DB question_id: 86461dba-3ae4-53d6-bd65-9a3ea521324b
\needspace{5\baselineskip}
\noindent\textbf{\color{ink}60.}
\par
"Obyekt" atamasining eng aniq ta'rifini toping.
\begin{enumerate}[label=\Alph*),leftmargin=9mm]
\item axborot olinadigan yoki tasvirlanadigan narsa, shaxs, jarayon yoxud tizim
\item muayyan vaqt va sharoitda sodir bo'ladigan voqea yoki o'zgarish
\item axborotni jo'natuvchidan qabul qiluvchiga yetkazuvchi mazmunli ifoda
\item axborotni tashuvchi fizik o'zgarish yoki belgi ketma-ketligi
\end{enumerate}
\begin{tcolorbox}[enhanced,breakable,colback=teal!4,colframe=teal!50!black,boxrule=0.5pt,arc=1.1mm,left=3mm,right=3mm,top=1.7mm,bottom=1.7mm,before skip=2mm,after skip=2mm]
\textbf{\color{teal}To'g'ri javob:} \textbf{A) axborot olinadigan yoki tasvirlanadigan narsa, shaxs, jarayon yoxud tizim}
\par\smallskip
\textbf{Izoh:} Obyekt - axborot olinadigan yoki tasvirlanadigan narsa, shaxs, jarayon yoxud tizim.
\end{tcolorbox}
\vspace{2mm}\hrule height 0.2pt \vspace{3mm}

% DB question_id: 8c546919-2b72-5f90-b066-94fcf6d54cce
\needspace{5\baselineskip}
\noindent\textbf{\color{ink}61.}
\par
Tizim blok, monitor, klaviatura va printer inventar qilindi. Ushbu vaziyatdagi asosiy tushunchani aniqlang.
\begin{enumerate}[label=\Alph*),leftmargin=9mm]
\item Dasturiy ta'minot
\item Insoniy ta'minot
\item Analog va raqamli axborot
\item Texnik ta'minot
\end{enumerate}
\begin{tcolorbox}[enhanced,breakable,colback=teal!4,colframe=teal!50!black,boxrule=0.5pt,arc=1.1mm,left=3mm,right=3mm,top=1.7mm,bottom=1.7mm,before skip=2mm,after skip=2mm]
\textbf{\color{teal}To'g'ri javob:} \textbf{D) Texnik ta'minot}
\par\smallskip
\textbf{Izoh:} Vaziyat Texnik ta'minot tushunchasini ifodalaydi: kompyuter tizimining barcha jismoniy qurilmalari majmui.
\end{tcolorbox}
\vspace{2mm}\hrule height 0.2pt \vspace{3mm}

% DB question_id: 8d3f395c-f4b4-51b4-86fe-f766576cb247
\needspace{5\baselineskip}
\noindent\textbf{\color{ink}62.}
\par
Quyidagi ta'rif qaysi tushunchaga tegishli? Kompyuter tizimining barcha jismoniy qurilmalari majmui.
\begin{enumerate}[label=\Alph*),leftmargin=9mm]
\item Dasturiy ta'minot
\item Insoniy ta'minot
\item Analog va raqamli axborot
\item Texnik ta'minot
\end{enumerate}
\begin{tcolorbox}[enhanced,breakable,colback=teal!4,colframe=teal!50!black,boxrule=0.5pt,arc=1.1mm,left=3mm,right=3mm,top=1.7mm,bottom=1.7mm,before skip=2mm,after skip=2mm]
\textbf{\color{teal}To'g'ri javob:} \textbf{D) Texnik ta'minot}
\par\smallskip
\textbf{Izoh:} Ta'rif Texnik ta'minot tushunchasiga tegishli.
\end{tcolorbox}
\vspace{2mm}\hrule height 0.2pt \vspace{3mm}

% DB question_id: 8ee5aa36-a817-5bc4-a5c2-78364f9a962e
\needspace{5\baselineskip}
\noindent\textbf{\color{ink}63.}
\par
"Insoniy ta'minot" atamasining eng aniq ta'rifini toping.
\begin{enumerate}[label=\Alph*),leftmargin=9mm]
\item analog axborot uzluksiz, raqamli axborot esa diskret qiymatlar orqali ifodalanadigan ikki xil tasvirlash shakli
\item axborot tizimini yaratadigan, boshqaradigan va undan foydalanadigan insonlar hamda ularning malakasi
\item axborot yaratish, undan foydalanish va bilim almashish iqtisodiy hamda ijtimoiy taraqqiyotning asosiy omiliga aylangan jamiyat
\item axborotni yig'ish, saqlash, qayta ishlash, uzatish va undan foydalanish usullarini o'rganuvchi fan
\end{enumerate}
\begin{tcolorbox}[enhanced,breakable,colback=teal!4,colframe=teal!50!black,boxrule=0.5pt,arc=1.1mm,left=3mm,right=3mm,top=1.7mm,bottom=1.7mm,before skip=2mm,after skip=2mm]
\textbf{\color{teal}To'g'ri javob:} \textbf{B) axborot tizimini yaratadigan, boshqaradigan va undan foydalanadigan insonlar hamda ularning malakasi}
\par\smallskip
\textbf{Izoh:} Insoniy ta'minot - axborot tizimini yaratadigan, boshqaradigan va undan foydalanadigan insonlar hamda ularning malakasi.
\end{tcolorbox}
\vspace{2mm}\hrule height 0.2pt \vspace{3mm}

% DB question_id: 8f4d1678-4268-521b-82c6-05b8d6e72117
\needspace{5\baselineskip}
\noindent\textbf{\color{ink}64.}
\par
"Axborotlashgan jamiyat" atamasining eng aniq ta'rifini toping.
\begin{enumerate}[label=\Alph*),leftmargin=9mm]
\item axborotni yig'ish, saqlash, qayta ishlash, uzatish va undan foydalanish usullarini o'rganuvchi fan
\item axborotni yaratish, qayta ishlash va almashishda kompyuter hamda telekommunikatsiya vositalaridan foydalanish sohasi
\item ma'lumotni diskret sonli ko'rinishda yaratish, qayta ishlash yoki uzatishga asoslangan texnologiya
\item axborot yaratish, undan foydalanish va bilim almashish iqtisodiy hamda ijtimoiy taraqqiyotning asosiy omiliga aylangan jamiyat
\end{enumerate}
\begin{tcolorbox}[enhanced,breakable,colback=teal!4,colframe=teal!50!black,boxrule=0.5pt,arc=1.1mm,left=3mm,right=3mm,top=1.7mm,bottom=1.7mm,before skip=2mm,after skip=2mm]
\textbf{\color{teal}To'g'ri javob:} \textbf{D) axborot yaratish, undan foydalanish va bilim almashish iqtisodiy hamda ijtimoiy taraqqiyotning asosiy omiliga aylangan jamiyat}
\par\smallskip
\textbf{Izoh:} Axborotlashgan jamiyat - axborot yaratish, undan foydalanish va bilim almashish iqtisodiy hamda ijtimoiy taraqqiyotning asosiy omiliga aylangan jamiyat.
\end{tcolorbox}
\vspace{2mm}\hrule height 0.2pt \vspace{3mm}

% DB question_id: 92267022-bc80-5d5d-92f4-b1a38f700a5a
\needspace{5\baselineskip}
\noindent\textbf{\color{ink}65.}
\par
"Bilimlar bazasi" atamasining eng aniq ta'rifini toping.
\begin{enumerate}[label=\Alph*),leftmargin=9mm]
\item axborot olinadigan yoki tasvirlanadigan narsa, shaxs, jarayon yoxud tizim
\item muayyan vaqt va sharoitda sodir bo'ladigan voqea yoki o'zgarish
\item axborotni jo'natuvchidan qabul qiluvchiga yetkazuvchi mazmunli ifoda
\item muayyan soha bo'yicha faktlar va qoidalar tartibli saqlanadigan tuzilma
\end{enumerate}
\begin{tcolorbox}[enhanced,breakable,colback=teal!4,colframe=teal!50!black,boxrule=0.5pt,arc=1.1mm,left=3mm,right=3mm,top=1.7mm,bottom=1.7mm,before skip=2mm,after skip=2mm]
\textbf{\color{teal}To'g'ri javob:} \textbf{D) muayyan soha bo'yicha faktlar va qoidalar tartibli saqlanadigan tuzilma}
\par\smallskip
\textbf{Izoh:} Bilimlar bazasi - muayyan soha bo'yicha faktlar va qoidalar tartibli saqlanadigan tuzilma.
\end{tcolorbox}
\vspace{2mm}\hrule height 0.2pt \vspace{3mm}

% DB question_id: 9865a7d7-ff13-58f7-a2ad-75e9be28235f
\needspace{5\baselineskip}
\noindent\textbf{\color{ink}66.}
\par
"Axborot tashuvchi" atamasining eng aniq ta'rifini toping.
\begin{enumerate}[label=\Alph*),leftmargin=9mm]
\item axborotni saqlash yoki yetkazishga xizmat qiladigan moddiy yoxud fizik vosita
\item ko'rish, eshitish, hid, ta'm va tana sezgisi orqali tashqi muhitdan ma'lumot olish jarayoni
\item kompyuter tizimida ma'lumot avval kiritilib, qayta ishlanib, natija chiqariladigan asosiy model
\item kompyuter tizimining barcha jismoniy qurilmalari majmui
\end{enumerate}
\begin{tcolorbox}[enhanced,breakable,colback=teal!4,colframe=teal!50!black,boxrule=0.5pt,arc=1.1mm,left=3mm,right=3mm,top=1.7mm,bottom=1.7mm,before skip=2mm,after skip=2mm]
\textbf{\color{teal}To'g'ri javob:} \textbf{A) axborotni saqlash yoki yetkazishga xizmat qiladigan moddiy yoxud fizik vosita}
\par\smallskip
\textbf{Izoh:} Axborot tashuvchi - axborotni saqlash yoki yetkazishga xizmat qiladigan moddiy yoxud fizik vosita.
\end{tcolorbox}
\vspace{2mm}\hrule height 0.2pt \vspace{3mm}

% DB question_id: 9f1c443d-1bff-56e3-8fb6-04c9d3c5c9d0
\needspace{5\baselineskip}
\noindent\textbf{\color{ink}67.}
\par
Quyidagi ta'rif qaysi tushunchaga tegishli? Qabul qiluvchi uchun mazmunli bo'lib, noaniqlikni kamaytiradigan ma'lumot.
\begin{enumerate}[label=\Alph*),leftmargin=9mm]
\item Axborot
\item Bilim
\item Ma'lumot \ensuremath{\to} axborot \ensuremath{\to} bilim zanjiri
\item Bilimlar bazasi
\end{enumerate}
\begin{tcolorbox}[enhanced,breakable,colback=teal!4,colframe=teal!50!black,boxrule=0.5pt,arc=1.1mm,left=3mm,right=3mm,top=1.7mm,bottom=1.7mm,before skip=2mm,after skip=2mm]
\textbf{\color{teal}To'g'ri javob:} \textbf{A) Axborot}
\par\smallskip
\textbf{Izoh:} Ta'rif Axborot tushunchasiga tegishli.
\end{tcolorbox}
\vspace{2mm}\hrule height 0.2pt \vspace{3mm}

% DB question_id: a6ee6c00-73cc-53ab-ae2d-7e7645ff29e4
\needspace{5\baselineskip}
\noindent\textbf{\color{ink}68.}
\par
Jadval katagida faqat "37,5" soni yozilgan, lekin u nimani anglatishi ko'rsatilmagan. Ushbu vaziyatdagi asosiy tushunchani aniqlang.
\begin{enumerate}[label=\Alph*),leftmargin=9mm]
\item Axborot
\item Bilim
\item Ma'lumot \ensuremath{\to} axborot \ensuremath{\to} bilim zanjiri
\item Ma'lumot
\end{enumerate}
\begin{tcolorbox}[enhanced,breakable,colback=teal!4,colframe=teal!50!black,boxrule=0.5pt,arc=1.1mm,left=3mm,right=3mm,top=1.7mm,bottom=1.7mm,before skip=2mm,after skip=2mm]
\textbf{\color{teal}To'g'ri javob:} \textbf{D) Ma'lumot}
\par\smallskip
\textbf{Izoh:} Vaziyat Ma'lumot tushunchasini ifodalaydi: qayd etilgan belgi, qiymat yoki kuzatuv bo'lib, kontekstsiz mazmuni yetarli bo'lmasligi mumkin.
\end{tcolorbox}
\vspace{2mm}\hrule height 0.2pt \vspace{3mm}

% DB question_id: a9edea8e-33ca-5366-8cc5-8eeb868f1baf
\needspace{5\baselineskip}
\noindent\textbf{\color{ink}69.}
\par
"Hodisa" atamasining eng aniq ta'rifini toping.
\begin{enumerate}[label=\Alph*),leftmargin=9mm]
\item axborotni jo'natuvchidan qabul qiluvchiga yetkazuvchi mazmunli ifoda
\item muayyan vaqt va sharoitda sodir bo'ladigan voqea yoki o'zgarish
\item axborotni tashuvchi fizik o'zgarish yoki belgi ketma-ketligi
\item axborotni saqlash yoki yetkazishga xizmat qiladigan moddiy yoxud fizik vosita
\end{enumerate}
\begin{tcolorbox}[enhanced,breakable,colback=teal!4,colframe=teal!50!black,boxrule=0.5pt,arc=1.1mm,left=3mm,right=3mm,top=1.7mm,bottom=1.7mm,before skip=2mm,after skip=2mm]
\textbf{\color{teal}To'g'ri javob:} \textbf{B) muayyan vaqt va sharoitda sodir bo'ladigan voqea yoki o'zgarish}
\par\smallskip
\textbf{Izoh:} Hodisa - muayyan vaqt va sharoitda sodir bo'ladigan voqea yoki o'zgarish.
\end{tcolorbox}
\vspace{2mm}\hrule height 0.2pt \vspace{3mm}

% DB question_id: ade1ccfa-cc81-4e38-a2f9-e5f4b6ce8ab6
\needspace{5\baselineskip}
\noindent\textbf{\color{ink}70.}
\par
Harorat datchigi sonli qiymat yubordi; o'lchov me'yor bilan solishtirib talqin qilindi; yakunda sovitish bo'yicha asosli qaror ishlab chiqildi. Xom qiymatdan bilimga o'tish qaysi tartibda kechdi?
\begin{enumerate}[label=\Alph*),leftmargin=9mm]
\item Bilim \ensuremath{\to} ma'lumot \ensuremath{\to} axborot zanjiri
\item Axborot \ensuremath{\to} ma'lumot \ensuremath{\to} bilim zanjiri
\item Ma'lumot \ensuremath{\to} axborot \ensuremath{\to} bilim zanjiri
\item Ma'lumot \ensuremath{\to} bilim \ensuremath{\to} axborot zanjiri
\end{enumerate}
\begin{tcolorbox}[enhanced,breakable,colback=teal!4,colframe=teal!50!black,boxrule=0.5pt,arc=1.1mm,left=3mm,right=3mm,top=1.7mm,bottom=1.7mm,before skip=2mm,after skip=2mm]
\textbf{\color{teal}To'g'ri javob:} \textbf{C) Ma'lumot \ensuremath{\to} axborot \ensuremath{\to} bilim zanjiri}
\par\smallskip
\textbf{Izoh:} Xom qiymat ma'lumot hisoblanadi. U vaziyatga bog'lab talqin qilinganda axborot, undan asoslangan qaror va umumlashgan tushuncha olinganda bilim shakllanadi.
\end{tcolorbox}
\vspace{2mm}\hrule height 0.2pt \vspace{3mm}

% DB question_id: b6fa2020-f265-5704-acd5-7c496fce4f30
\needspace{5\baselineskip}
\noindent\textbf{\color{ink}71.}
\par
Serverning soat 10:15 da o'chib qolishi qayd etildi. Ushbu vaziyatdagi asosiy tushunchani aniqlang.
\begin{enumerate}[label=\Alph*),leftmargin=9mm]
\item Xabar
\item Hodisa
\item Signal
\item Axborot tashuvchi
\end{enumerate}
\begin{tcolorbox}[enhanced,breakable,colback=teal!4,colframe=teal!50!black,boxrule=0.5pt,arc=1.1mm,left=3mm,right=3mm,top=1.7mm,bottom=1.7mm,before skip=2mm,after skip=2mm]
\textbf{\color{teal}To'g'ri javob:} \textbf{B) Hodisa}
\par\smallskip
\textbf{Izoh:} Vaziyat Hodisa tushunchasini ifodalaydi: muayyan vaqt va sharoitda sodir bo'ladigan voqea yoki o'zgarish.
\end{tcolorbox}
\vspace{2mm}\hrule height 0.2pt \vspace{3mm}

% DB question_id: b8ee7cb5-6649-5a73-b91b-e961887396fc
\needspace{5\baselineskip}
\noindent\textbf{\color{ink}72.}
\par
Quyidagi ta'rif qaysi tushunchaga tegishli? Axborotni tashuvchi fizik o'zgarish yoki belgi ketma-ketligi.
\begin{enumerate}[label=\Alph*),leftmargin=9mm]
\item Axborot tashuvchi
\item Insonning axborotni qabul qilishi
\item Kiritish-qayta ishlash-chiqarish modeli
\item Signal
\end{enumerate}
\begin{tcolorbox}[enhanced,breakable,colback=teal!4,colframe=teal!50!black,boxrule=0.5pt,arc=1.1mm,left=3mm,right=3mm,top=1.7mm,bottom=1.7mm,before skip=2mm,after skip=2mm]
\textbf{\color{teal}To'g'ri javob:} \textbf{D) Signal}
\par\smallskip
\textbf{Izoh:} Ta'rif Signal tushunchasiga tegishli.
\end{tcolorbox}
\vspace{2mm}\hrule height 0.2pt \vspace{3mm}

% DB question_id: bbe897ff-2a6f-5417-a753-3a7b09c94bd0
\needspace{5\baselineskip}
\noindent\textbf{\color{ink}73.}
\par
Quyidagi ta'rif qaysi tushunchaga tegishli? Axborotni saqlash yoki yetkazishga xizmat qiladigan moddiy yoxud fizik vosita.
\begin{enumerate}[label=\Alph*),leftmargin=9mm]
\item Axborot tashuvchi
\item Insonning axborotni qabul qilishi
\item Kiritish-qayta ishlash-chiqarish modeli
\item Texnik ta'minot
\end{enumerate}
\begin{tcolorbox}[enhanced,breakable,colback=teal!4,colframe=teal!50!black,boxrule=0.5pt,arc=1.1mm,left=3mm,right=3mm,top=1.7mm,bottom=1.7mm,before skip=2mm,after skip=2mm]
\textbf{\color{teal}To'g'ri javob:} \textbf{A) Axborot tashuvchi}
\par\smallskip
\textbf{Izoh:} Ta'rif Axborot tashuvchi tushunchasiga tegishli.
\end{tcolorbox}
\vspace{2mm}\hrule height 0.2pt \vspace{3mm}

% DB question_id: bf8cf476-8f34-4fc1-989f-c7f64f6ec73d
\needspace{5\baselineskip}
\noindent\textbf{\color{ink}74.}
\par
Mikrofonning uzluksiz elektr signali ADC yordamida sonli namunalarga o'tkazildi. Qaysi ikki axborot ifodalash shakli o'rtasidagi farq ko'rsatilgan?
\begin{enumerate}[label=\Alph*),leftmargin=9mm]
\item Matnli va grafik axborot
\item Ovozli va video axborot
\item Analog va raqamli axborot
\item Ochiq va yopiq axborot
\end{enumerate}
\begin{tcolorbox}[enhanced,breakable,colback=teal!4,colframe=teal!50!black,boxrule=0.5pt,arc=1.1mm,left=3mm,right=3mm,top=1.7mm,bottom=1.7mm,before skip=2mm,after skip=2mm]
\textbf{\color{teal}To'g'ri javob:} \textbf{C) Analog va raqamli axborot}
\par\smallskip
\textbf{Izoh:} Mikrofon hosil qilgan uzluksiz signal analog, ADCdan keyingi diskret sonli namunalar raqamli ifodadir. Farq uzluksizlik va diskretlikda.
\end{tcolorbox}
\vspace{2mm}\hrule height 0.2pt \vspace{3mm}

% DB question_id: c28a00c2-90ff-54f3-84e7-be1548be6870
\needspace{5\baselineskip}
\noindent\textbf{\color{ink}75.}
\par
Optik tolada yorug'lik impulslari ketma-ketligi uzatildi. Ushbu vaziyatdagi asosiy tushunchani aniqlang.
\begin{enumerate}[label=\Alph*),leftmargin=9mm]
\item Axborot tashuvchi
\item Insonning axborotni qabul qilishi
\item Kiritish-qayta ishlash-chiqarish modeli
\item Signal
\end{enumerate}
\begin{tcolorbox}[enhanced,breakable,colback=teal!4,colframe=teal!50!black,boxrule=0.5pt,arc=1.1mm,left=3mm,right=3mm,top=1.7mm,bottom=1.7mm,before skip=2mm,after skip=2mm]
\textbf{\color{teal}To'g'ri javob:} \textbf{D) Signal}
\par\smallskip
\textbf{Izoh:} Vaziyat Signal tushunchasini ifodalaydi: axborotni tashuvchi fizik o'zgarish yoki belgi ketma-ketligi.
\end{tcolorbox}
\vspace{2mm}\hrule height 0.2pt \vspace{3mm}

% DB question_id: c3462947-dd64-5c52-b78d-85c57518db2c
\needspace{5\baselineskip}
\noindent\textbf{\color{ink}76.}
\par
Quyidagi ta'rif qaysi tushunchaga tegishli? Axborot tizimini yaratadigan, boshqaradigan va undan foydalanadigan insonlar hamda ularning malakasi.
\begin{enumerate}[label=\Alph*),leftmargin=9mm]
\item Analog va raqamli axborot
\item Insoniy ta'minot
\item Axborotlashgan jamiyat
\item Informatika
\end{enumerate}
\begin{tcolorbox}[enhanced,breakable,colback=teal!4,colframe=teal!50!black,boxrule=0.5pt,arc=1.1mm,left=3mm,right=3mm,top=1.7mm,bottom=1.7mm,before skip=2mm,after skip=2mm]
\textbf{\color{teal}To'g'ri javob:} \textbf{B) Insoniy ta'minot}
\par\smallskip
\textbf{Izoh:} Ta'rif Insoniy ta'minot tushunchasiga tegishli.
\end{tcolorbox}
\vspace{2mm}\hrule height 0.2pt \vspace{3mm}

% DB question_id: c6a69555-20b0-572c-94df-c641c2a7a12b
\needspace{5\baselineskip}
\noindent\textbf{\color{ink}77.}
\par
Quyidagi ta'rif qaysi tushunchaga tegishli? Axborotni yaratish, qayta ishlash va almashishda kompyuter hamda telekommunikatsiya vositalaridan foydalanish sohasi.
\begin{enumerate}[label=\Alph*),leftmargin=9mm]
\item Raqamli texnologiya
\item AKT
\item Ma'lumot
\item Axborot
\end{enumerate}
\begin{tcolorbox}[enhanced,breakable,colback=teal!4,colframe=teal!50!black,boxrule=0.5pt,arc=1.1mm,left=3mm,right=3mm,top=1.7mm,bottom=1.7mm,before skip=2mm,after skip=2mm]
\textbf{\color{teal}To'g'ri javob:} \textbf{B) AKT}
\par\smallskip
\textbf{Izoh:} Ta'rif AKT tushunchasiga tegishli.
\end{tcolorbox}
\vspace{2mm}\hrule height 0.2pt \vspace{3mm}

% DB question_id: c8aa0008-2b2f-5a5c-9625-e5c482a91c12
\needspace{5\baselineskip}
\noindent\textbf{\color{ink}78.}
\par
"Informatika" atamasining eng aniq ta'rifini toping.
\begin{enumerate}[label=\Alph*),leftmargin=9mm]
\item axborotni yig'ish, saqlash, qayta ishlash, uzatish va undan foydalanish usullarini o'rganuvchi fan
\item axborotni yaratish, qayta ishlash va almashishda kompyuter hamda telekommunikatsiya vositalaridan foydalanish sohasi
\item ma'lumotni diskret sonli ko'rinishda yaratish, qayta ishlash yoki uzatishga asoslangan texnologiya
\item qayd etilgan belgi, qiymat yoki kuzatuv bo'lib, kontekstsiz mazmuni yetarli bo'lmasligi mumkin
\end{enumerate}
\begin{tcolorbox}[enhanced,breakable,colback=teal!4,colframe=teal!50!black,boxrule=0.5pt,arc=1.1mm,left=3mm,right=3mm,top=1.7mm,bottom=1.7mm,before skip=2mm,after skip=2mm]
\textbf{\color{teal}To'g'ri javob:} \textbf{A) axborotni yig'ish, saqlash, qayta ishlash, uzatish va undan foydalanish usullarini o'rganuvchi fan}
\par\smallskip
\textbf{Izoh:} Informatika - axborotni yig'ish, saqlash, qayta ishlash, uzatish va undan foydalanish usullarini o'rganuvchi fan.
\end{tcolorbox}
\vspace{2mm}\hrule height 0.2pt \vspace{3mm}

% DB question_id: c98a4753-f7b0-4850-9c35-c7ca67bd3c2a
\needspace{5\baselineskip}
\noindent\textbf{\color{ink}79.}
\par
"42" kontekstsiz, "50 dan 42" esa natijani anglatdi. Bu vaziyatdagi asosiy tushunchani aniqlang.
\begin{enumerate}[label=\Alph*),leftmargin=9mm]
\item Kodlash va shifrlash farqi
\item Bit va bayt farqi
\item Axborot va ma'lumot farqi
\item Rastr va vektor farqi
\end{enumerate}
\begin{tcolorbox}[enhanced,breakable,colback=teal!4,colframe=teal!50!black,boxrule=0.5pt,arc=1.1mm,left=3mm,right=3mm,top=1.7mm,bottom=1.7mm,before skip=2mm,after skip=2mm]
\textbf{\color{teal}To'g'ri javob:} \textbf{C) Axborot va ma'lumot farqi}
\par\smallskip
\textbf{Izoh:} Vaziyat Axborot va ma'lumot farqiga misol: ma'lumot qayd etilgan qiymat, axborot esa qabul qiluvchi uchun mazmunli ma'lumot ekanini ajratish.
\end{tcolorbox}
\vspace{2mm}\hrule height 0.2pt \vspace{3mm}

% DB question_id: cfac2302-c40a-5132-9557-40c3c58a807a
\needspace{5\baselineskip}
\noindent\textbf{\color{ink}80.}
\par
Quyidagi ta'rif qaysi tushunchaga tegishli? Ko'rish, eshitish, hid, ta'm va tana sezgisi orqali tashqi muhitdan ma'lumot olish jarayoni.
\begin{enumerate}[label=\Alph*),leftmargin=9mm]
\item Kiritish-qayta ishlash-chiqarish modeli
\item Insonning axborotni qabul qilishi
\item Texnik ta'minot
\item Dasturiy ta'minot
\end{enumerate}
\begin{tcolorbox}[enhanced,breakable,colback=teal!4,colframe=teal!50!black,boxrule=0.5pt,arc=1.1mm,left=3mm,right=3mm,top=1.7mm,bottom=1.7mm,before skip=2mm,after skip=2mm]
\textbf{\color{teal}To'g'ri javob:} \textbf{B) Insonning axborotni qabul qilishi}
\par\smallskip
\textbf{Izoh:} Ta'rif Insonning axborotni qabul qilishi tushunchasiga tegishli.
\end{tcolorbox}
\vspace{2mm}\hrule height 0.2pt \vspace{3mm}

% DB question_id: d42e3b2a-0961-424f-8b3f-f2815e254113
\needspace{5\baselineskip}
\noindent\textbf{\color{ink}81.}
\par
Sensordan olingan ma'lumot ikkilik kodda saqlanib, dastur orqali qayta ishlanadi. Ushbu ishlash usuli qaysi texnologiya turiga xos?
\begin{enumerate}[label=\Alph*),leftmargin=9mm]
\item Analog texnologiya
\item Mexanik texnologiya
\item Raqamli texnologiya
\item Gibrid texnologiya
\end{enumerate}
\begin{tcolorbox}[enhanced,breakable,colback=teal!4,colframe=teal!50!black,boxrule=0.5pt,arc=1.1mm,left=3mm,right=3mm,top=1.7mm,bottom=1.7mm,before skip=2mm,after skip=2mm]
\textbf{\color{teal}To'g'ri javob:} \textbf{C) Raqamli texnologiya}
\par\smallskip
\textbf{Izoh:} Raqamli texnologiyalar diskret sonli kodda axborotni ifodalaydi va qayta ishlaydi. Analog usulda signal uzluksiz, gibrid tizimda esa ikkala texnologiya aralash kelishi mumkin.
\end{tcolorbox}
\vspace{2mm}\hrule height 0.2pt \vspace{3mm}

% DB question_id: d6472433-4a74-59b7-87a5-b64cf759ff83
\needspace{5\baselineskip}
\noindent\textbf{\color{ink}82.}
\par
"Texnik ta'minot" atamasining eng aniq ta'rifini toping.
\begin{enumerate}[label=\Alph*),leftmargin=9mm]
\item kompyuter bajaradigan dasturlar va ularni boshqaruvchi ko'rsatmalar majmui
\item axborot tizimini yaratadigan, boshqaradigan va undan foydalanadigan insonlar hamda ularning malakasi
\item analog axborot uzluksiz, raqamli axborot esa diskret qiymatlar orqali ifodalanadigan ikki xil tasvirlash shakli
\item kompyuter tizimining barcha jismoniy qurilmalari majmui
\end{enumerate}
\begin{tcolorbox}[enhanced,breakable,colback=teal!4,colframe=teal!50!black,boxrule=0.5pt,arc=1.1mm,left=3mm,right=3mm,top=1.7mm,bottom=1.7mm,before skip=2mm,after skip=2mm]
\textbf{\color{teal}To'g'ri javob:} \textbf{D) kompyuter tizimining barcha jismoniy qurilmalari majmui}
\par\smallskip
\textbf{Izoh:} Texnik ta'minot - kompyuter tizimining barcha jismoniy qurilmalari majmui.
\end{tcolorbox}
\vspace{2mm}\hrule height 0.2pt \vspace{3mm}

% DB question_id: d8a3a5e4-f917-456c-8cae-e7a7ffa679c7
\needspace{5\baselineskip}
\noindent\textbf{\color{ink}83.}
\par
Ekspert tizimida soha faktlari va <<agar ... unda ...>> qoidalari xulosa chiqarishda ishlatilishi uchun tartibli saqlanadi. Tizimning bu tarkibiy qismi nima deb ataladi?
\begin{enumerate}[label=\Alph*),leftmargin=9mm]
\item Ma'lumotlar bazasi
\item Xulosalash mexanizmi
\item Foydalanuvchi interfeysi
\item Bilimlar bazasi
\end{enumerate}
\begin{tcolorbox}[enhanced,breakable,colback=teal!4,colframe=teal!50!black,boxrule=0.5pt,arc=1.1mm,left=3mm,right=3mm,top=1.7mm,bottom=1.7mm,before skip=2mm,after skip=2mm]
\textbf{\color{teal}To'g'ri javob:} \textbf{D) Bilimlar bazasi}
\par\smallskip
\textbf{Izoh:} Bilimlar bazasi - ekspert tizimidagi fakt va qoidalar to'plami. Xulosalash mexanizmi esa mavjud bilimlar asosida xulosa chiqaradi; interfeys foydalanuvchi bilan ishlashni ta'minlaydi.
\end{tcolorbox}
\vspace{2mm}\hrule height 0.2pt \vspace{3mm}

% DB question_id: dbda02dd-8c38-4989-b8ab-e36e7a5ee9a8
\needspace{5\baselineskip}
\noindent\textbf{\color{ink}84.}
\par
Mikrofon qabul qilgan tovush uzluksiz signal, kompyuter xotirasidagi shu tovushning namunalari esa diskret qiymatlardir. Bu taqqoslash axborotning qaysi ikki ifodalanish turiga oid?
\begin{enumerate}[label=\Alph*),leftmargin=9mm]
\item Matnli va grafikli axborot
\item Ovozli va tasvirli axborot
\item Analog va raqamli axborot
\item Ochiq va maxfiy axborot
\end{enumerate}
\begin{tcolorbox}[enhanced,breakable,colback=teal!4,colframe=teal!50!black,boxrule=0.5pt,arc=1.1mm,left=3mm,right=3mm,top=1.7mm,bottom=1.7mm,before skip=2mm,after skip=2mm]
\textbf{\color{teal}To'g'ri javob:} \textbf{C) Analog va raqamli axborot}
\par\smallskip
\textbf{Izoh:} Uzluksiz signal analog axborotni, raqamli qurilmadagi diskret namunalar esa raqamli axborotni ifodalaydi.
\end{tcolorbox}
\vspace{2mm}\hrule height 0.2pt \vspace{3mm}

% DB question_id: e63ad566-fa2a-5478-ab78-9b4bb69a935a
\needspace{5\baselineskip}
\noindent\textbf{\color{ink}85.}
\par
"Insonning axborotni qabul qilishi" atamasining eng aniq ta'rifini toping.
\begin{enumerate}[label=\Alph*),leftmargin=9mm]
\item kompyuter tizimida ma'lumot avval kiritilib, qayta ishlanib, natija chiqariladigan asosiy model
\item ko'rish, eshitish, hid, ta'm va tana sezgisi orqali tashqi muhitdan ma'lumot olish jarayoni
\item kompyuter tizimining barcha jismoniy qurilmalari majmui
\item kompyuter bajaradigan dasturlar va ularni boshqaruvchi ko'rsatmalar majmui
\end{enumerate}
\begin{tcolorbox}[enhanced,breakable,colback=teal!4,colframe=teal!50!black,boxrule=0.5pt,arc=1.1mm,left=3mm,right=3mm,top=1.7mm,bottom=1.7mm,before skip=2mm,after skip=2mm]
\textbf{\color{teal}To'g'ri javob:} \textbf{B) ko'rish, eshitish, hid, ta'm va tana sezgisi orqali tashqi muhitdan ma'lumot olish jarayoni}
\par\smallskip
\textbf{Izoh:} Insonning axborotni qabul qilishi - ko'rish, eshitish, hid, ta'm va tana sezgisi orqali tashqi muhitdan ma'lumot olish jarayoni.
\end{tcolorbox}
\vspace{2mm}\hrule height 0.2pt \vspace{3mm}

% DB question_id: f86e0af6-b808-4084-a6e4-3ae9f9f97771
\needspace{5\baselineskip}
\noindent\textbf{\color{ink}86.}
\par
Ekspert maslahat tizimida alomatlar bilan birga <<agar belgi kuzatilsa, ehtimoliy xulosa...>> kabi sohaga oid qoidalar saqlanadi. Faktlar va qoidalar saqlanadigan qism qanday ataladi?
\begin{enumerate}[label=\Alph*),leftmargin=9mm]
\item Ma'lumotlar bazasi
\item Xulosalash mexanizmi
\item Qidiruv indeksi
\item Bilimlar bazasi
\end{enumerate}
\begin{tcolorbox}[enhanced,breakable,colback=teal!4,colframe=teal!50!black,boxrule=0.5pt,arc=1.1mm,left=3mm,right=3mm,top=1.7mm,bottom=1.7mm,before skip=2mm,after skip=2mm]
\textbf{\color{teal}To'g'ri javob:} \textbf{D) Bilimlar bazasi}
\par\smallskip
\textbf{Izoh:} Bilimlar bazasi ekspert tizimida sohaga oid fakt va mantiqiy qoidalarni saqlaydi. Xulosalash mexanizmi aynan shu qoidalardan xulosa olish uchun foydalanadi.
\end{tcolorbox}
\vspace{2mm}\hrule height 0.2pt \vspace{3mm}

% DB question_id: 20228a08-150c-4922-a4dc-045417798178
\needspace{5\baselineskip}
\noindent\textbf{\color{ink}87.}
\par
Tadqiqotda maktab tarmog'ining tezligi va yuklanishi kuzatildi. Vaziyatni tahlil qilib, unga eng aniq mos tushunchani toping.
\begin{enumerate}[label=\Alph*),leftmargin=9mm]
\item Xabar
\item Signal
\item Informatika
\item Obyekt
\end{enumerate}
\begin{tcolorbox}[enhanced,breakable,colback=teal!4,colframe=teal!50!black,boxrule=0.5pt,arc=1.1mm,left=3mm,right=3mm,top=1.7mm,bottom=1.7mm,before skip=2mm,after skip=2mm]
\textbf{\color{teal}To'g'ri javob:} \textbf{D) Obyekt}
\par\smallskip
\textbf{Izoh:} Vaziyat Obyektga misol: o'rganilayotgan narsa, hodisa, jarayon yoki shaxs.
\end{tcolorbox}
\vspace{2mm}\hrule height 0.2pt \vspace{3mm}

% DB question_id: 5fda9695-5123-45c8-8f84-2728007e6be3
\needspace{5\baselineskip}
\noindent\textbf{\color{ink}88.}
\par
"50 savoldan 42 tasi to'g'ri" xabari o'quvchining natijasini tushuntirdi. Vaziyatni tahlil qilib, unga eng aniq mos tushunchani toping.
\begin{enumerate}[label=\Alph*),leftmargin=9mm]
\item Bilim
\item Axborot
\item Obyekt
\item Xabar
\end{enumerate}
\begin{tcolorbox}[enhanced,breakable,colback=teal!4,colframe=teal!50!black,boxrule=0.5pt,arc=1.1mm,left=3mm,right=3mm,top=1.7mm,bottom=1.7mm,before skip=2mm,after skip=2mm]
\textbf{\color{teal}To'g'ri javob:} \textbf{B) Axborot}
\par\smallskip
\textbf{Izoh:} Vaziyat Axborotga misol: qabul qiluvchi uchun mazmunli bo'lib, noaniqlikni kamaytiradigan ma'lumot.
\end{tcolorbox}
\vspace{2mm}\hrule height 0.2pt \vspace{3mm}

% DB question_id: 77cb5672-1006-4e65-bf66-d88f0842f47f
\needspace{5\baselineskip}
\noindent\textbf{\color{ink}89.}
\par
O'quvchi formulani yoddan aytish bilan cheklanmay, yangi masalani yechdi. Vaziyatni tahlil qilib, unga eng aniq mos tushunchani toping.
\begin{enumerate}[label=\Alph*),leftmargin=9mm]
\item Obyekt
\item Xabar
\item Bilim
\item Signal
\end{enumerate}
\begin{tcolorbox}[enhanced,breakable,colback=teal!4,colframe=teal!50!black,boxrule=0.5pt,arc=1.1mm,left=3mm,right=3mm,top=1.7mm,bottom=1.7mm,before skip=2mm,after skip=2mm]
\textbf{\color{teal}To'g'ri javob:} \textbf{C) Bilim}
\par\smallskip
\textbf{Izoh:} Vaziyat Bilimga misol: axborotni tushunish, bog'lash va vaziyatda qo'llash natijasida shakllangan tizimli tushuncha.
\end{tcolorbox}
\vspace{2mm}\hrule height 0.2pt \vspace{3mm}

% DB question_id: 880d4289-dc7f-4067-b410-392fa9ea15bb
\needspace{5\baselineskip}
\noindent\textbf{\color{ink}90.}
\par
Server foydalanuvchiga "fayl muvaffaqiyatli yuklandi" degan bildirishnoma yubordi. Yaqin tushunchalarni farqlab, aynan qaysi atama vaziyatni eng aniq ifodalashini toping.
\begin{enumerate}[label=\Alph*),leftmargin=9mm]
\item Xabar
\item Signal
\item Informatika
\item Texnik ta'minot
\end{enumerate}
\begin{tcolorbox}[enhanced,breakable,colback=teal!4,colframe=teal!50!black,boxrule=0.5pt,arc=1.1mm,left=3mm,right=3mm,top=1.7mm,bottom=1.7mm,before skip=2mm,after skip=2mm]
\textbf{\color{teal}To'g'ri javob:} \textbf{A) Xabar}
\par\smallskip
\textbf{Izoh:} Vaziyat Xabarga misol: axborotni qabul qiluvchiga yetkazish uchun ifodalangan mazmun.
\end{tcolorbox}
\vspace{2mm}\hrule height 0.2pt \vspace{3mm}

% DB question_id: dbe6b419-8a8e-4124-9bc8-c1e6665bf21a
\needspace{5\baselineskip}
\noindent\textbf{\color{ink}91.}
\par
Optik tolada yorug'lik impulslari 0 va 1 larni tashidi. Yaqin tushunchalarni farqlab, aynan qaysi atama vaziyatni eng aniq ifodalashini toping.
\begin{enumerate}[label=\Alph*),leftmargin=9mm]
\item Informatika
\item Signal
\item Texnik ta'minot
\item Dasturiy ta'minot
\end{enumerate}
\begin{tcolorbox}[enhanced,breakable,colback=teal!4,colframe=teal!50!black,boxrule=0.5pt,arc=1.1mm,left=3mm,right=3mm,top=1.7mm,bottom=1.7mm,before skip=2mm,after skip=2mm]
\textbf{\color{teal}To'g'ri javob:} \textbf{B) Signal}
\par\smallskip
\textbf{Izoh:} Vaziyat Signalga misol: xabarni kanal bo'ylab tashiydigan fizik o'zgarish.
\end{tcolorbox}
\vspace{2mm}\hrule height 0.2pt \vspace{3mm}
\section{Axborotning turlari, shakllari, xossalari va manbalari}

% DB question_id: 0da75917-b369-40d5-bc03-f55f970f2088
\needspace{5\baselineskip}
\noindent\textbf{\color{ink}92.}
\par
"Qabul qiluvchi" atamasiga eng aniq ta'rifni toping.
\begin{enumerate}[label=\Alph*),leftmargin=9mm]
\item xabarni qabul qilib, undan foydalanadigan obyekt yoki shaxs
\item hodisa yoki tadqiqotdan bevosita olingan asl dalil
\item birlamchi manbalarni tahlil qilib yoki qayta bayon qilib yaratilgan manba
\item foydalanuvchi yoki tizim yangilamaguncha o'zgarmaydigan mazmun
\end{enumerate}
\begin{tcolorbox}[enhanced,breakable,colback=teal!4,colframe=teal!50!black,boxrule=0.5pt,arc=1.1mm,left=3mm,right=3mm,top=1.7mm,bottom=1.7mm,before skip=2mm,after skip=2mm]
\textbf{\color{teal}To'g'ri javob:} \textbf{A) xabarni qabul qilib, undan foydalanadigan obyekt yoki shaxs}
\par\smallskip
\textbf{Izoh:} Qabul qiluvchi - xabarni qabul qilib, undan foydalanadigan obyekt yoki shaxs.
\end{tcolorbox}
\vspace{2mm}\hrule height 0.2pt \vspace{3mm}

% DB question_id: 16fce42a-fe68-4a98-a1d9-9b3e4aa1b373
\needspace{5\baselineskip}
\noindent\textbf{\color{ink}93.}
\par
"Dinamik axborot" atamasiga eng aniq ta'rifni toping.
\begin{enumerate}[label=\Alph*),leftmargin=9mm]
\item vaqt, hodisa yoki ma'lumotlar bazasiga qarab avtomatik yangilanadigan mazmun
\item turli manbalardan kerakli ma'lumotlarni olish va jamlash jarayoni
\item mavjud ma'lumotni tartiblash, hisoblash yoki tahlil qilish orqali yangi natija olish jarayoni
\item axborotni keyinchalik foydalanish uchun tashuvchida qayd etish jarayoni
\end{enumerate}
\begin{tcolorbox}[enhanced,breakable,colback=teal!4,colframe=teal!50!black,boxrule=0.5pt,arc=1.1mm,left=3mm,right=3mm,top=1.7mm,bottom=1.7mm,before skip=2mm,after skip=2mm]
\textbf{\color{teal}To'g'ri javob:} \textbf{A) vaqt, hodisa yoki ma'lumotlar bazasiga qarab avtomatik yangilanadigan mazmun}
\par\smallskip
\textbf{Izoh:} Dinamik axborot - vaqt, hodisa yoki ma'lumotlar bazasiga qarab avtomatik yangilanadigan mazmun.
\end{tcolorbox}
\vspace{2mm}\hrule height 0.2pt \vspace{3mm}

% DB question_id: 1c950ba8-230e-4232-9ba7-bd8f04e2ba5e
\needspace{5\baselineskip}
\noindent\textbf{\color{ink}94.}
\par
"Axborot manbai" atamasiga eng aniq ta'rifni toping.
\begin{enumerate}[label=\Alph*),leftmargin=9mm]
\item xabar manbadan qabul qiluvchiga o'tadigan muhit yoki vosita
\item xabarni qabul qilib, undan foydalanadigan obyekt yoki shaxs
\item xabarni hosil qiluvchi yoki axborot kelib chiqadigan obyekt
\item hodisa yoki tadqiqotdan bevosita olingan asl dalil
\end{enumerate}
\begin{tcolorbox}[enhanced,breakable,colback=teal!4,colframe=teal!50!black,boxrule=0.5pt,arc=1.1mm,left=3mm,right=3mm,top=1.7mm,bottom=1.7mm,before skip=2mm,after skip=2mm]
\textbf{\color{teal}To'g'ri javob:} \textbf{C) xabarni hosil qiluvchi yoki axborot kelib chiqadigan obyekt}
\par\smallskip
\textbf{Izoh:} Axborot manbai - xabarni hosil qiluvchi yoki axborot kelib chiqadigan obyekt.
\end{tcolorbox}
\vspace{2mm}\hrule height 0.2pt \vspace{3mm}

% DB question_id: 21919500-82af-4541-815b-61bea27b2596
\needspace{5\baselineskip}
\noindent\textbf{\color{ink}95.}
\par
"Vizual axborot" atamasiga eng aniq ta'rifni toping.
\begin{enumerate}[label=\Alph*),leftmargin=9mm]
\item vaqt yoki qiymat bo'yicha uzluksiz o'zgaradigan axborot
\item ko'rish a'zosi orqali qabul qilinadigan axborot
\item alohida qiymatlar yoki belgilar ketma-ketligida ifodalangan axborot
\item axborotning obyekt holatini xatosiz va kerakli aniqlik darajasida aks ettirish xossasi
\end{enumerate}
\begin{tcolorbox}[enhanced,breakable,colback=teal!4,colframe=teal!50!black,boxrule=0.5pt,arc=1.1mm,left=3mm,right=3mm,top=1.7mm,bottom=1.7mm,before skip=2mm,after skip=2mm]
\textbf{\color{teal}To'g'ri javob:} \textbf{B) ko'rish a'zosi orqali qabul qilinadigan axborot}
\par\smallskip
\textbf{Izoh:} Vizual axborot - ko'rish a'zosi orqali qabul qilinadigan axborot.
\end{tcolorbox}
\vspace{2mm}\hrule height 0.2pt \vspace{3mm}

% DB question_id: 289295e3-db3f-482b-b30d-e3f854c75d75
\needspace{5\baselineskip}
\noindent\textbf{\color{ink}96.}
\par
"Statik axborot" atamasiga eng aniq ta'rifni toping.
\begin{enumerate}[label=\Alph*),leftmargin=9mm]
\item vaqt, hodisa yoki ma'lumotlar bazasiga qarab avtomatik yangilanadigan mazmun
\item turli manbalardan kerakli ma'lumotlarni olish va jamlash jarayoni
\item mavjud ma'lumotni tartiblash, hisoblash yoki tahlil qilish orqali yangi natija olish jarayoni
\item foydalanuvchi yoki tizim yangilamaguncha o'zgarmaydigan mazmun
\end{enumerate}
\begin{tcolorbox}[enhanced,breakable,colback=teal!4,colframe=teal!50!black,boxrule=0.5pt,arc=1.1mm,left=3mm,right=3mm,top=1.7mm,bottom=1.7mm,before skip=2mm,after skip=2mm]
\textbf{\color{teal}To'g'ri javob:} \textbf{D) foydalanuvchi yoki tizim yangilamaguncha o'zgarmaydigan mazmun}
\par\smallskip
\textbf{Izoh:} Statik axborot - foydalanuvchi yoki tizim yangilamaguncha o'zgarmaydigan mazmun.
\end{tcolorbox}
\vspace{2mm}\hrule height 0.2pt \vspace{3mm}

% DB question_id: 3a55f48d-ac58-4abe-817a-b3cdeb918ff0
\needspace{5\baselineskip}
\noindent\textbf{\color{ink}97.}
\par
"Grafik axborot" atamasiga eng aniq ta'rifni toping.
\begin{enumerate}[label=\Alph*),leftmargin=9mm]
\item tovush orqali ifodalangan va eshitish bilan qabul qilinadigan axborot
\item matn, grafika, audio, video va animatsiyaning uyg'un birikmasi
\item rasm, chizma, diagramma yoki sxema ko'rinishidagi axborot
\item ko'rish a'zosi orqali qabul qilinadigan axborot
\end{enumerate}
\begin{tcolorbox}[enhanced,breakable,colback=teal!4,colframe=teal!50!black,boxrule=0.5pt,arc=1.1mm,left=3mm,right=3mm,top=1.7mm,bottom=1.7mm,before skip=2mm,after skip=2mm]
\textbf{\color{teal}To'g'ri javob:} \textbf{C) rasm, chizma, diagramma yoki sxema ko'rinishidagi axborot}
\par\smallskip
\textbf{Izoh:} Grafik axborot - rasm, chizma, diagramma yoki sxema ko'rinishidagi axborot.
\end{tcolorbox}
\vspace{2mm}\hrule height 0.2pt \vspace{3mm}

% DB question_id: 48b9190a-39ac-422f-b6b2-53d7561de156
\needspace{5\baselineskip}
\noindent\textbf{\color{ink}98.}
\par
"Diskret axborot" atamasiga eng aniq ta'rifni toping.
\begin{enumerate}[label=\Alph*),leftmargin=9mm]
\item axborotning obyekt holatini xatosiz va kerakli aniqlik darajasida aks ettirish xossasi
\item axborotning tekshirilgan dalil va ishonchli manbaga tayanganlik xossasi
\item qaror chiqarish uchun zarur bo'lgan muhim qismlar yetarli berilganlik xossasi
\item alohida qiymatlar yoki belgilar ketma-ketligida ifodalangan axborot
\end{enumerate}
\begin{tcolorbox}[enhanced,breakable,colback=teal!4,colframe=teal!50!black,boxrule=0.5pt,arc=1.1mm,left=3mm,right=3mm,top=1.7mm,bottom=1.7mm,before skip=2mm,after skip=2mm]
\textbf{\color{teal}To'g'ri javob:} \textbf{D) alohida qiymatlar yoki belgilar ketma-ketligida ifodalangan axborot}
\par\smallskip
\textbf{Izoh:} Diskret axborot - alohida qiymatlar yoki belgilar ketma-ketligida ifodalangan axborot.
\end{tcolorbox}
\vspace{2mm}\hrule height 0.2pt \vspace{3mm}

% DB question_id: 4d25b7a7-3892-44fa-af94-4c6e5cc7b55a
\needspace{5\baselineskip}
\noindent\textbf{\color{ink}99.}
\par
"Aloqa kanali" atamasiga eng aniq ta'rifni toping.
\begin{enumerate}[label=\Alph*),leftmargin=9mm]
\item xabarni qabul qilib, undan foydalanadigan obyekt yoki shaxs
\item hodisa yoki tadqiqotdan bevosita olingan asl dalil
\item birlamchi manbalarni tahlil qilib yoki qayta bayon qilib yaratilgan manba
\item xabar manbadan qabul qiluvchiga o'tadigan muhit yoki vosita
\end{enumerate}
\begin{tcolorbox}[enhanced,breakable,colback=teal!4,colframe=teal!50!black,boxrule=0.5pt,arc=1.1mm,left=3mm,right=3mm,top=1.7mm,bottom=1.7mm,before skip=2mm,after skip=2mm]
\textbf{\color{teal}To'g'ri javob:} \textbf{D) xabar manbadan qabul qiluvchiga o'tadigan muhit yoki vosita}
\par\smallskip
\textbf{Izoh:} Aloqa kanali - xabar manbadan qabul qiluvchiga o'tadigan muhit yoki vosita.
\end{tcolorbox}
\vspace{2mm}\hrule height 0.2pt \vspace{3mm}

% DB question_id: 4f3c18a5-56f3-49a8-b4e8-0b83e818c7dc
\needspace{5\baselineskip}
\noindent\textbf{\color{ink}100.}
\par
"Matnli axborot" atamasiga eng aniq ta'rifni toping.
\begin{enumerate}[label=\Alph*),leftmargin=9mm]
\item rasm, chizma, diagramma yoki sxema ko'rinishidagi axborot
\item harf, raqam va tinish belgilaridan tuzilgan yozma axborot
\item tovush orqali ifodalangan va eshitish bilan qabul qilinadigan axborot
\item matn, grafika, audio, video va animatsiyaning uyg'un birikmasi
\end{enumerate}
\begin{tcolorbox}[enhanced,breakable,colback=teal!4,colframe=teal!50!black,boxrule=0.5pt,arc=1.1mm,left=3mm,right=3mm,top=1.7mm,bottom=1.7mm,before skip=2mm,after skip=2mm]
\textbf{\color{teal}To'g'ri javob:} \textbf{B) harf, raqam va tinish belgilaridan tuzilgan yozma axborot}
\par\smallskip
\textbf{Izoh:} Matnli axborot - harf, raqam va tinish belgilaridan tuzilgan yozma axborot.
\end{tcolorbox}
\vspace{2mm}\hrule height 0.2pt \vspace{3mm}

% DB question_id: 666c9b00-6b13-4901-b86b-00173aebf309
\needspace{5\baselineskip}
\noindent\textbf{\color{ink}101.}
\par
"Birlamchi manba" atamasiga eng aniq ta'rifni toping.
\begin{enumerate}[label=\Alph*),leftmargin=9mm]
\item birlamchi manbalarni tahlil qilib yoki qayta bayon qilib yaratilgan manba
\item hodisa yoki tadqiqotdan bevosita olingan asl dalil
\item foydalanuvchi yoki tizim yangilamaguncha o'zgarmaydigan mazmun
\item vaqt, hodisa yoki ma'lumotlar bazasiga qarab avtomatik yangilanadigan mazmun
\end{enumerate}
\begin{tcolorbox}[enhanced,breakable,colback=teal!4,colframe=teal!50!black,boxrule=0.5pt,arc=1.1mm,left=3mm,right=3mm,top=1.7mm,bottom=1.7mm,before skip=2mm,after skip=2mm]
\textbf{\color{teal}To'g'ri javob:} \textbf{B) hodisa yoki tadqiqotdan bevosita olingan asl dalil}
\par\smallskip
\textbf{Izoh:} Birlamchi manba - hodisa yoki tadqiqotdan bevosita olingan asl dalil.
\end{tcolorbox}
\vspace{2mm}\hrule height 0.2pt \vspace{3mm}

% DB question_id: 802513c3-39b6-4cf3-8d6e-e83db240fba0
\needspace{5\baselineskip}
\noindent\textbf{\color{ink}102.}
\par
"Multimedia" tushunchasining "Axborotning xossalari, turlari va shakllari" bobidagi eng aniq ta'rifini toping.
\begin{enumerate}[label=\Alph*),leftmargin=9mm]
\item matn, grafika, audio, video va animatsiyaning uyg'un birikmasi
\item ko'rish a'zosi orqali qabul qilinadigan axborot
\item vaqt yoki qiymat bo'yicha uzluksiz o'zgaradigan axborot
\item alohida qiymatlar yoki belgilar ketma-ketligida ifodalangan axborot
\end{enumerate}
\begin{tcolorbox}[enhanced,breakable,colback=teal!4,colframe=teal!50!black,boxrule=0.5pt,arc=1.1mm,left=3mm,right=3mm,top=1.7mm,bottom=1.7mm,before skip=2mm,after skip=2mm]
\textbf{\color{teal}To'g'ri javob:} \textbf{A) matn, grafika, audio, video va animatsiyaning uyg'un birikmasi}
\par\smallskip
\textbf{Izoh:} Multimedia - matn, grafika, audio, video va animatsiyaning uyg'un birikmasi.
\end{tcolorbox}
\vspace{2mm}\hrule height 0.2pt \vspace{3mm}

% DB question_id: 90a23293-eaef-4976-8447-156b9654c70d
\needspace{5\baselineskip}
\noindent\textbf{\color{ink}103.}
\par
"Audioaxborot" atamasiga eng aniq ta'rifni toping.
\begin{enumerate}[label=\Alph*),leftmargin=9mm]
\item matn, grafika, audio, video va animatsiyaning uyg'un birikmasi
\item ko'rish a'zosi orqali qabul qilinadigan axborot
\item vaqt yoki qiymat bo'yicha uzluksiz o'zgaradigan axborot
\item tovush orqali ifodalangan va eshitish bilan qabul qilinadigan axborot
\end{enumerate}
\begin{tcolorbox}[enhanced,breakable,colback=teal!4,colframe=teal!50!black,boxrule=0.5pt,arc=1.1mm,left=3mm,right=3mm,top=1.7mm,bottom=1.7mm,before skip=2mm,after skip=2mm]
\textbf{\color{teal}To'g'ri javob:} \textbf{D) tovush orqali ifodalangan va eshitish bilan qabul qilinadigan axborot}
\par\smallskip
\textbf{Izoh:} Audioaxborot - tovush orqali ifodalangan va eshitish bilan qabul qilinadigan axborot.
\end{tcolorbox}
\vspace{2mm}\hrule height 0.2pt \vspace{3mm}

% DB question_id: 9dab7ab9-cbcd-45b4-b708-9ff04d7b2f30
\needspace{5\baselineskip}
\noindent\textbf{\color{ink}104.}
\par
"Ikkilamchi manba" atamasiga eng aniq ta'rifni toping.
\begin{enumerate}[label=\Alph*),leftmargin=9mm]
\item foydalanuvchi yoki tizim yangilamaguncha o'zgarmaydigan mazmun
\item vaqt, hodisa yoki ma'lumotlar bazasiga qarab avtomatik yangilanadigan mazmun
\item birlamchi manbalarni tahlil qilib yoki qayta bayon qilib yaratilgan manba
\item turli manbalardan kerakli ma'lumotlarni olish va jamlash jarayoni
\end{enumerate}
\begin{tcolorbox}[enhanced,breakable,colback=teal!4,colframe=teal!50!black,boxrule=0.5pt,arc=1.1mm,left=3mm,right=3mm,top=1.7mm,bottom=1.7mm,before skip=2mm,after skip=2mm]
\textbf{\color{teal}To'g'ri javob:} \textbf{C) birlamchi manbalarni tahlil qilib yoki qayta bayon qilib yaratilgan manba}
\par\smallskip
\textbf{Izoh:} Ikkilamchi manba - birlamchi manbalarni tahlil qilib yoki qayta bayon qilib yaratilgan manba.
\end{tcolorbox}
\vspace{2mm}\hrule height 0.2pt \vspace{3mm}

% DB question_id: f1c98612-24dc-4ef7-bdd5-c9c811874b4a
\needspace{5\baselineskip}
\noindent\textbf{\color{ink}105.}
\par
"Analog axborot" atamasiga eng aniq ta'rifni toping.
\begin{enumerate}[label=\Alph*),leftmargin=9mm]
\item alohida qiymatlar yoki belgilar ketma-ketligida ifodalangan axborot
\item axborotning obyekt holatini xatosiz va kerakli aniqlik darajasida aks ettirish xossasi
\item vaqt yoki qiymat bo'yicha uzluksiz o'zgaradigan axborot
\item axborotning tekshirilgan dalil va ishonchli manbaga tayanganlik xossasi
\end{enumerate}
\begin{tcolorbox}[enhanced,breakable,colback=teal!4,colframe=teal!50!black,boxrule=0.5pt,arc=1.1mm,left=3mm,right=3mm,top=1.7mm,bottom=1.7mm,before skip=2mm,after skip=2mm]
\textbf{\color{teal}To'g'ri javob:} \textbf{C) vaqt yoki qiymat bo'yicha uzluksiz o'zgaradigan axborot}
\par\smallskip
\textbf{Izoh:} Analog axborot - vaqt yoki qiymat bo'yicha uzluksiz o'zgaradigan axborot.
\end{tcolorbox}
\vspace{2mm}\hrule height 0.2pt \vspace{3mm}

% DB question_id: 05224e72-6b0a-4892-87a2-024284dcf28c
\needspace{5\baselineskip}
\noindent\textbf{\color{ink}106.}
\par
Quyidagi ta'rif qaysi atamaga tegishli? Hodisa yoki tadqiqotdan bevosita olingan asl dalil.
\begin{enumerate}[label=\Alph*),leftmargin=9mm]
\item Ikkilamchi manba
\item Statik axborot
\item Dinamik axborot
\item Birlamchi manba
\end{enumerate}
\begin{tcolorbox}[enhanced,breakable,colback=teal!4,colframe=teal!50!black,boxrule=0.5pt,arc=1.1mm,left=3mm,right=3mm,top=1.7mm,bottom=1.7mm,before skip=2mm,after skip=2mm]
\textbf{\color{teal}To'g'ri javob:} \textbf{D) Birlamchi manba}
\par\smallskip
\textbf{Izoh:} Ta'rif Birlamchi manba tushunchasiga tegishli.
\end{tcolorbox}
\vspace{2mm}\hrule height 0.2pt \vspace{3mm}

% DB question_id: 0fbdd8a2-242f-45de-b900-a6cff03dcbac
\needspace{5\baselineskip}
\noindent\textbf{\color{ink}107.}
\par
Quyidagi ta'rif qaysi atamaga tegishli? Vaqt, hodisa yoki ma'lumotlar bazasiga qarab avtomatik yangilanadigan mazmun.
\begin{enumerate}[label=\Alph*),leftmargin=9mm]
\item Axborotni yig'ish
\item Axborotni qayta ishlash
\item Dinamik axborot
\item Axborotni saqlash
\end{enumerate}
\begin{tcolorbox}[enhanced,breakable,colback=teal!4,colframe=teal!50!black,boxrule=0.5pt,arc=1.1mm,left=3mm,right=3mm,top=1.7mm,bottom=1.7mm,before skip=2mm,after skip=2mm]
\textbf{\color{teal}To'g'ri javob:} \textbf{C) Dinamik axborot}
\par\smallskip
\textbf{Izoh:} Ta'rif Dinamik axborot tushunchasiga tegishli.
\end{tcolorbox}
\vspace{2mm}\hrule height 0.2pt \vspace{3mm}

% DB question_id: 1222ea74-30d9-4f6f-9551-e49aabb96ee4
\needspace{5\baselineskip}
\noindent\textbf{\color{ink}108.}
\par
Quyidagi ta'rif qaysi atamaga tegishli? Xabarni hosil qiluvchi yoki axborot kelib chiqadigan obyekt.
\begin{enumerate}[label=\Alph*),leftmargin=9mm]
\item Axborot manbai
\item Aloqa kanali
\item Qabul qiluvchi
\item Birlamchi manba
\end{enumerate}
\begin{tcolorbox}[enhanced,breakable,colback=teal!4,colframe=teal!50!black,boxrule=0.5pt,arc=1.1mm,left=3mm,right=3mm,top=1.7mm,bottom=1.7mm,before skip=2mm,after skip=2mm]
\textbf{\color{teal}To'g'ri javob:} \textbf{A) Axborot manbai}
\par\smallskip
\textbf{Izoh:} Ta'rif Axborot manbai tushunchasiga tegishli.
\end{tcolorbox}
\vspace{2mm}\hrule height 0.2pt \vspace{3mm}

% DB question_id: 31dd6717-228e-4ae6-9c5e-8689ba64742d
\needspace{5\baselineskip}
\noindent\textbf{\color{ink}109.}
\par
Quyidagi ta'rif qaysi atamaga tegishli? Ko'rish a'zosi orqali qabul qilinadigan axborot.
\begin{enumerate}[label=\Alph*),leftmargin=9mm]
\item Analog axborot
\item Diskret axborot
\item Aniqlik
\item Vizual axborot
\end{enumerate}
\begin{tcolorbox}[enhanced,breakable,colback=teal!4,colframe=teal!50!black,boxrule=0.5pt,arc=1.1mm,left=3mm,right=3mm,top=1.7mm,bottom=1.7mm,before skip=2mm,after skip=2mm]
\textbf{\color{teal}To'g'ri javob:} \textbf{D) Vizual axborot}
\par\smallskip
\textbf{Izoh:} Ta'rif Vizual axborot tushunchasiga tegishli.
\end{tcolorbox}
\vspace{2mm}\hrule height 0.2pt \vspace{3mm}

% DB question_id: 36ef133d-d6de-45d4-bd5b-67fdc446a1e0
\needspace{5\baselineskip}
\noindent\textbf{\color{ink}110.}
\par
Quyidagi ta'rif qaysi atamaga tegishli? Harf, raqam va tinish belgilaridan tuzilgan yozma axborot.
\begin{enumerate}[label=\Alph*),leftmargin=9mm]
\item Grafik axborot
\item Audioaxborot
\item Multimedia
\item Matnli axborot
\end{enumerate}
\begin{tcolorbox}[enhanced,breakable,colback=teal!4,colframe=teal!50!black,boxrule=0.5pt,arc=1.1mm,left=3mm,right=3mm,top=1.7mm,bottom=1.7mm,before skip=2mm,after skip=2mm]
\textbf{\color{teal}To'g'ri javob:} \textbf{D) Matnli axborot}
\par\smallskip
\textbf{Izoh:} Ta'rif Matnli axborot tushunchasiga tegishli.
\end{tcolorbox}
\vspace{2mm}\hrule height 0.2pt \vspace{3mm}

% DB question_id: 3a286d9b-7799-4d2e-be22-34109a594c3f
\needspace{5\baselineskip}
\noindent\textbf{\color{ink}111.}
\par
Quyidagi ta'rif qaysi atamaga tegishli? Birlamchi manbalarni tahlil qilib yoki qayta bayon qilib yaratilgan manba.
\begin{enumerate}[label=\Alph*),leftmargin=9mm]
\item Ikkilamchi manba
\item Statik axborot
\item Dinamik axborot
\item Axborotni yig'ish
\end{enumerate}
\begin{tcolorbox}[enhanced,breakable,colback=teal!4,colframe=teal!50!black,boxrule=0.5pt,arc=1.1mm,left=3mm,right=3mm,top=1.7mm,bottom=1.7mm,before skip=2mm,after skip=2mm]
\textbf{\color{teal}To'g'ri javob:} \textbf{A) Ikkilamchi manba}
\par\smallskip
\textbf{Izoh:} Ta'rif Ikkilamchi manba tushunchasiga tegishli.
\end{tcolorbox}
\vspace{2mm}\hrule height 0.2pt \vspace{3mm}

% DB question_id: 88b69132-f3a3-46e8-b499-c02a68caf6b3
\needspace{5\baselineskip}
\noindent\textbf{\color{ink}112.}
\par
Quyidagi ta'rif qaysi atamaga tegishli? Matn, grafika, audio, video va animatsiyaning uyg'un birikmasi.
\begin{enumerate}[label=\Alph*),leftmargin=9mm]
\item Vizual axborot
\item Analog axborot
\item Multimedia
\item Diskret axborot
\end{enumerate}
\begin{tcolorbox}[enhanced,breakable,colback=teal!4,colframe=teal!50!black,boxrule=0.5pt,arc=1.1mm,left=3mm,right=3mm,top=1.7mm,bottom=1.7mm,before skip=2mm,after skip=2mm]
\textbf{\color{teal}To'g'ri javob:} \textbf{C) Multimedia}
\par\smallskip
\textbf{Izoh:} Ta'rif Multimedia tushunchasiga tegishli.
\end{tcolorbox}
\vspace{2mm}\hrule height 0.2pt \vspace{3mm}

% DB question_id: 963a32c0-e01d-4a56-aa9f-e6a00c7143d4
\needspace{5\baselineskip}
\noindent\textbf{\color{ink}113.}
\par
Quyidagi ta'rif qaysi atamaga tegishli? Xabar manbadan qabul qiluvchiga o'tadigan muhit yoki vosita.
\begin{enumerate}[label=\Alph*),leftmargin=9mm]
\item Qabul qiluvchi
\item Aloqa kanali
\item Birlamchi manba
\item Ikkilamchi manba
\end{enumerate}
\begin{tcolorbox}[enhanced,breakable,colback=teal!4,colframe=teal!50!black,boxrule=0.5pt,arc=1.1mm,left=3mm,right=3mm,top=1.7mm,bottom=1.7mm,before skip=2mm,after skip=2mm]
\textbf{\color{teal}To'g'ri javob:} \textbf{B) Aloqa kanali}
\par\smallskip
\textbf{Izoh:} Ta'rif Aloqa kanali tushunchasiga tegishli.
\end{tcolorbox}
\vspace{2mm}\hrule height 0.2pt \vspace{3mm}

% DB question_id: a0b8f85f-8d46-4b63-9200-36784651fbb2
\needspace{5\baselineskip}
\noindent\textbf{\color{ink}114.}
\par
Quyidagi ta'rif qaysi atamaga tegishli? Foydalanuvchi yoki tizim yangilamaguncha o'zgarmaydigan mazmun.
\begin{enumerate}[label=\Alph*),leftmargin=9mm]
\item Dinamik axborot
\item Statik axborot
\item Axborotni yig'ish
\item Axborotni qayta ishlash
\end{enumerate}
\begin{tcolorbox}[enhanced,breakable,colback=teal!4,colframe=teal!50!black,boxrule=0.5pt,arc=1.1mm,left=3mm,right=3mm,top=1.7mm,bottom=1.7mm,before skip=2mm,after skip=2mm]
\textbf{\color{teal}To'g'ri javob:} \textbf{B) Statik axborot}
\par\smallskip
\textbf{Izoh:} Ta'rif Statik axborot tushunchasiga tegishli.
\end{tcolorbox}
\vspace{2mm}\hrule height 0.2pt \vspace{3mm}

% DB question_id: dece1d00-79d5-4c92-992f-67313ec44deb
\needspace{5\baselineskip}
\noindent\textbf{\color{ink}115.}
\par
A, E, I, O harflaridan barcha 4 belgili so'zlar tuziladi; harflar bir so'z ichida takrorlanishi mumkin. So'zlar alifbo tartibida 1 dan boshlab raqamlangan. 248-o'rindagi so'zni toping.
\begin{enumerate}[label=\Alph*),leftmargin=9mm]
\item OOEO
\item OOEA
\item OOOO
\item OOEI
\end{enumerate}
\begin{tcolorbox}[enhanced,breakable,colback=teal!4,colframe=teal!50!black,boxrule=0.5pt,arc=1.1mm,left=3mm,right=3mm,top=1.7mm,bottom=1.7mm,before skip=2mm,after skip=2mm]
\textbf{\color{teal}To'g'ri javob:} \textbf{A) OOEO}
\par\smallskip
\textbf{Izoh:} 4 ta harfdan takrorlash bilan 4\textsuperscript{4} = 256 ta so'z tuziladi. 248-o'rinning noldan boshlangan indeksi 247 = 3313\textsubscript{4}. A=0, E=1, I=2, O=3 bo'lsa, 3313\textsubscript{4} indeksi OOEO so'ziga mos.
\end{tcolorbox}
\vspace{2mm}\hrule height 0.2pt \vspace{3mm}

% DB question_id: e48e8422-c7b2-4ac8-a14b-73a7a03dd368
\needspace{5\baselineskip}
\noindent\textbf{\color{ink}116.}
\par
Quyidagi ta'rif qaysi atamaga tegishli? Alohida qiymatlar yoki belgilar ketma-ketligida ifodalangan axborot.
\begin{enumerate}[label=\Alph*),leftmargin=9mm]
\item Aniqlik
\item Diskret axborot
\item Ishonchlilik
\item To'liqlik
\end{enumerate}
\begin{tcolorbox}[enhanced,breakable,colback=teal!4,colframe=teal!50!black,boxrule=0.5pt,arc=1.1mm,left=3mm,right=3mm,top=1.7mm,bottom=1.7mm,before skip=2mm,after skip=2mm]
\textbf{\color{teal}To'g'ri javob:} \textbf{B) Diskret axborot}
\par\smallskip
\textbf{Izoh:} Ta'rif Diskret axborot tushunchasiga tegishli.
\end{tcolorbox}
\vspace{2mm}\hrule height 0.2pt \vspace{3mm}

% DB question_id: e7f2ecaa-2a57-4ab5-bf70-ae6e22f99d02
\needspace{5\baselineskip}
\noindent\textbf{\color{ink}117.}
\par
Quyidagi ta'rif qaysi atamaga tegishli? Vaqt yoki qiymat bo'yicha uzluksiz o'zgaradigan axborot.
\begin{enumerate}[label=\Alph*),leftmargin=9mm]
\item Analog axborot
\item Diskret axborot
\item Aniqlik
\item Ishonchlilik
\end{enumerate}
\begin{tcolorbox}[enhanced,breakable,colback=teal!4,colframe=teal!50!black,boxrule=0.5pt,arc=1.1mm,left=3mm,right=3mm,top=1.7mm,bottom=1.7mm,before skip=2mm,after skip=2mm]
\textbf{\color{teal}To'g'ri javob:} \textbf{A) Analog axborot}
\par\smallskip
\textbf{Izoh:} Ta'rif Analog axborot tushunchasiga tegishli.
\end{tcolorbox}
\vspace{2mm}\hrule height 0.2pt \vspace{3mm}

% DB question_id: e97163ad-c7c8-4bcb-b9e8-12b347c918fd
\needspace{5\baselineskip}
\noindent\textbf{\color{ink}118.}
\par
Quyidagi ta'rif qaysi atamaga tegishli? Xabarni qabul qilib, undan foydalanadigan obyekt yoki shaxs.
\begin{enumerate}[label=\Alph*),leftmargin=9mm]
\item Birlamchi manba
\item Ikkilamchi manba
\item Qabul qiluvchi
\item Statik axborot
\end{enumerate}
\begin{tcolorbox}[enhanced,breakable,colback=teal!4,colframe=teal!50!black,boxrule=0.5pt,arc=1.1mm,left=3mm,right=3mm,top=1.7mm,bottom=1.7mm,before skip=2mm,after skip=2mm]
\textbf{\color{teal}To'g'ri javob:} \textbf{C) Qabul qiluvchi}
\par\smallskip
\textbf{Izoh:} Ta'rif Qabul qiluvchi tushunchasiga tegishli.
\end{tcolorbox}
\vspace{2mm}\hrule height 0.2pt \vspace{3mm}

% DB question_id: ef557dad-8a11-47f3-8fd2-157d5a145378
\needspace{5\baselineskip}
\noindent\textbf{\color{ink}119.}
\par
Quyidagi ta'rif qaysi atamaga tegishli? Tovush orqali ifodalangan va eshitish bilan qabul qilinadigan axborot.
\begin{enumerate}[label=\Alph*),leftmargin=9mm]
\item Multimedia
\item Audioaxborot
\item Vizual axborot
\item Analog axborot
\end{enumerate}
\begin{tcolorbox}[enhanced,breakable,colback=teal!4,colframe=teal!50!black,boxrule=0.5pt,arc=1.1mm,left=3mm,right=3mm,top=1.7mm,bottom=1.7mm,before skip=2mm,after skip=2mm]
\textbf{\color{teal}To'g'ri javob:} \textbf{B) Audioaxborot}
\par\smallskip
\textbf{Izoh:} Ta'rif Audioaxborot tushunchasiga tegishli.
\end{tcolorbox}
\vspace{2mm}\hrule height 0.2pt \vspace{3mm}

% DB question_id: fd39aa04-4abe-4fcb-8cca-304a5cb6b774
\needspace{5\baselineskip}
\noindent\textbf{\color{ink}120.}
\par
Quyidagi ta'rif qaysi atamaga tegishli? Rasm, chizma, diagramma yoki sxema ko'rinishidagi axborot.
\begin{enumerate}[label=\Alph*),leftmargin=9mm]
\item Grafik axborot
\item Audioaxborot
\item Multimedia
\item Vizual axborot
\end{enumerate}
\begin{tcolorbox}[enhanced,breakable,colback=teal!4,colframe=teal!50!black,boxrule=0.5pt,arc=1.1mm,left=3mm,right=3mm,top=1.7mm,bottom=1.7mm,before skip=2mm,after skip=2mm]
\textbf{\color{teal}To'g'ri javob:} \textbf{A) Grafik axborot}
\par\smallskip
\textbf{Izoh:} Ta'rif Grafik axborot tushunchasiga tegishli.
\end{tcolorbox}
\vspace{2mm}\hrule height 0.2pt \vspace{3mm}

% DB question_id: 01611a84-d41d-5fb4-bda6-22306aa24d42
\needspace{5\baselineskip}
\noindent\textbf{\color{ink}121.}
\par
Jadvalda o'quvchilar ballari 72, 85 va 91 ko'rinishida berildi. Ushbu vaziyatdagi asosiy tushunchani aniqlang.
\begin{enumerate}[label=\Alph*),leftmargin=9mm]
\item Grafik axborot
\item Sonli axborot
\item Audioaxborot
\item Videoaxborot
\end{enumerate}
\begin{tcolorbox}[enhanced,breakable,colback=teal!4,colframe=teal!50!black,boxrule=0.5pt,arc=1.1mm,left=3mm,right=3mm,top=1.7mm,bottom=1.7mm,before skip=2mm,after skip=2mm]
\textbf{\color{teal}To'g'ri javob:} \textbf{B) Sonli axborot}
\par\smallskip
\textbf{Izoh:} Vaziyat Sonli axborot tushunchasini ifodalaydi: miqdor va o'lchovlarni sonlar orqali ifodalovchi axborot.
\end{tcolorbox}
\vspace{2mm}\hrule height 0.2pt \vspace{3mm}

% DB question_id: 03d02590-6473-51bb-9303-b79078994a16
\needspace{5\baselineskip}
\noindent\textbf{\color{ink}122.}
\par
"Videoaxborot" atamasining eng aniq ta'rifini toping.
\begin{enumerate}[label=\Alph*),leftmargin=9mm]
\item matn, grafika, audio yoki videoning birgalikdagi ko'rinishi
\item vaqt bo'yicha almashuvchi tasvirlar ketma-ketligi ko'rinishidagi axborot
\item vaqt davomida o'zgarmaydigan yoki bitta holatni tasvirlaydigan axborot
\item vaqt davomida holati o'zgarib boradigan axborot
\end{enumerate}
\begin{tcolorbox}[enhanced,breakable,colback=teal!4,colframe=teal!50!black,boxrule=0.5pt,arc=1.1mm,left=3mm,right=3mm,top=1.7mm,bottom=1.7mm,before skip=2mm,after skip=2mm]
\textbf{\color{teal}To'g'ri javob:} \textbf{B) vaqt bo'yicha almashuvchi tasvirlar ketma-ketligi ko'rinishidagi axborot}
\par\smallskip
\textbf{Izoh:} Videoaxborot - vaqt bo'yicha almashuvchi tasvirlar ketma-ketligi ko'rinishidagi axborot.
\end{tcolorbox}
\vspace{2mm}\hrule height 0.2pt \vspace{3mm}

% DB question_id: 0c498fb3-60d7-59f7-b07c-d669e18f5697
\needspace{5\baselineskip}
\noindent\textbf{\color{ink}123.}
\par
Qaysi juftlik bir xil tasnif mezoniga tegishli?
\begin{enumerate}[label=\Alph*),leftmargin=9mm]
\item Matnli - maxfiy
\item Birlamchi - audio
\item Aniq - video
\item Statik - dinamik
\end{enumerate}
\begin{tcolorbox}[enhanced,breakable,colback=teal!4,colframe=teal!50!black,boxrule=0.5pt,arc=1.1mm,left=3mm,right=3mm,top=1.7mm,bottom=1.7mm,before skip=2mm,after skip=2mm]
\textbf{\color{teal}To'g'ri javob:} \textbf{D) Statik - dinamik}
\par\smallskip
\textbf{Izoh:} Statik va dinamik axborot vaqt bo'yicha o'zgarish mezoniga tegishli.
\end{tcolorbox}
\vspace{2mm}\hrule height 0.2pt \vspace{3mm}

% DB question_id: 0ceafe9b-858f-5dd4-ad5d-7e8619dd4d1e
\needspace{5\baselineskip}
\noindent\textbf{\color{ink}124.}
\par
Tadqiqotchi o'zi o'tkazgan tajriba natijalarini jurnalga yozdi. Ushbu vaziyatdagi asosiy tushunchani aniqlang.
\begin{enumerate}[label=\Alph*),leftmargin=9mm]
\item Ikkilamchi manba
\item Bevosita va bilvosita manba
\item Birlamchi manba
\item Matnli axborot
\end{enumerate}
\begin{tcolorbox}[enhanced,breakable,colback=teal!4,colframe=teal!50!black,boxrule=0.5pt,arc=1.1mm,left=3mm,right=3mm,top=1.7mm,bottom=1.7mm,before skip=2mm,after skip=2mm]
\textbf{\color{teal}To'g'ri javob:} \textbf{C) Birlamchi manba}
\par\smallskip
\textbf{Izoh:} Vaziyat Birlamchi manba tushunchasini ifodalaydi: hodisa yoki tadqiqotga bevosita aloqador asl dalil yoki birinchi qo'l manbasi.
\end{tcolorbox}
\vspace{2mm}\hrule height 0.2pt \vspace{3mm}

% DB question_id: 0cee520d-7408-5f33-91dd-f88620920d2c
\needspace{5\baselineskip}
\noindent\textbf{\color{ink}125.}
\par
"Diskret axborot" atamasining eng aniq ta'rifini toping.
\begin{enumerate}[label=\Alph*),leftmargin=9mm]
\item foydalanishi cheklanmagan yoki keng jamoatchilikka e'lon qilingan axborot
\item muayyan shaxsga tegishli yoki kirishi cheklangan axborot
\item alohida, sanaladigan holat yoki qiymatlar ketma-ketligida ifodalanadigan axborot
\item hodisa yoki tadqiqotga bevosita aloqador asl dalil yoki birinchi qo'l manbasi
\end{enumerate}
\begin{tcolorbox}[enhanced,breakable,colback=teal!4,colframe=teal!50!black,boxrule=0.5pt,arc=1.1mm,left=3mm,right=3mm,top=1.7mm,bottom=1.7mm,before skip=2mm,after skip=2mm]
\textbf{\color{teal}To'g'ri javob:} \textbf{C) alohida, sanaladigan holat yoki qiymatlar ketma-ketligida ifodalanadigan axborot}
\par\smallskip
\textbf{Izoh:} Diskret axborot - alohida, sanaladigan holat yoki qiymatlar ketma-ketligida ifodalanadigan axborot.
\end{tcolorbox}
\vspace{2mm}\hrule height 0.2pt \vspace{3mm}

% DB question_id: 0f315633-8d3e-5fa4-926a-fe37c138134e
\needspace{5\baselineskip}
\noindent\textbf{\color{ink}126.}
\par
Quyidagi ta'rif qaysi tushunchaga tegishli? Alohida, sanaladigan holat yoki qiymatlar ketma-ketligida ifodalanadigan axborot.
\begin{enumerate}[label=\Alph*),leftmargin=9mm]
\item Ochiq va ommaviy axborot
\item Diskret axborot
\item Shaxsiy va maxfiy axborot
\item Birlamchi manba
\end{enumerate}
\begin{tcolorbox}[enhanced,breakable,colback=teal!4,colframe=teal!50!black,boxrule=0.5pt,arc=1.1mm,left=3mm,right=3mm,top=1.7mm,bottom=1.7mm,before skip=2mm,after skip=2mm]
\textbf{\color{teal}To'g'ri javob:} \textbf{B) Diskret axborot}
\par\smallskip
\textbf{Izoh:} Ta'rif Diskret axborot tushunchasiga tegishli.
\end{tcolorbox}
\vspace{2mm}\hrule height 0.2pt \vspace{3mm}

% DB question_id: 13e390ce-4a9d-5e84-a611-b089261fc271
\needspace{5\baselineskip}
\noindent\textbf{\color{ink}127.}
\par
Shaxsning telefon raqami ommaviy katalogda uning roziligisiz e'lon qilindi. Axborotning asosiy turi qaysi?
\begin{enumerate}[label=\Alph*),leftmargin=9mm]
\item Ommaviy mulk
\item Grafik axborot
\item Shaxsiy axborot
\item Ikkilamchi axborot
\end{enumerate}
\begin{tcolorbox}[enhanced,breakable,colback=teal!4,colframe=teal!50!black,boxrule=0.5pt,arc=1.1mm,left=3mm,right=3mm,top=1.7mm,bottom=1.7mm,before skip=2mm,after skip=2mm]
\textbf{\color{teal}To'g'ri javob:} \textbf{C) Shaxsiy axborot}
\par\smallskip
\textbf{Izoh:} Telefon raqami shaxsni aniqlashga bog'liq shaxsiy ma'lumotdir.
\end{tcolorbox}
\vspace{2mm}\hrule height 0.2pt \vspace{3mm}

% DB question_id: 14ebaedb-b544-5fec-b9d2-f05a88f2123d
\needspace{5\baselineskip}
\noindent\textbf{\color{ink}128.}
\par
Quyidagi ta'rif qaysi tushunchaga tegishli? Rasm, chizma, diagramma yoki sxema ko'rinishidagi axborot.
\begin{enumerate}[label=\Alph*),leftmargin=9mm]
\item Audioaxborot
\item Videoaxborot
\item Grafik axborot
\item Kombinatsiyalashgan axborot
\end{enumerate}
\begin{tcolorbox}[enhanced,breakable,colback=teal!4,colframe=teal!50!black,boxrule=0.5pt,arc=1.1mm,left=3mm,right=3mm,top=1.7mm,bottom=1.7mm,before skip=2mm,after skip=2mm]
\textbf{\color{teal}To'g'ri javob:} \textbf{C) Grafik axborot}
\par\smallskip
\textbf{Izoh:} Ta'rif Grafik axborot tushunchasiga tegishli.
\end{tcolorbox}
\vspace{2mm}\hrule height 0.2pt \vspace{3mm}

% DB question_id: 16061a5c-8c5e-56ba-a3d4-94293fd5da26
\needspace{5\baselineskip}
\noindent\textbf{\color{ink}129.}
\par
Foydalanuvchining paroli va tibbiy ma'lumoti himoyalangan bo'limda saqlandi. Ushbu vaziyatdagi asosiy tushunchani aniqlang.
\begin{enumerate}[label=\Alph*),leftmargin=9mm]
\item Shaxsiy va maxfiy axborot
\item Birlamchi manba
\item Ikkilamchi manba
\item Bevosita va bilvosita manba
\end{enumerate}
\begin{tcolorbox}[enhanced,breakable,colback=teal!4,colframe=teal!50!black,boxrule=0.5pt,arc=1.1mm,left=3mm,right=3mm,top=1.7mm,bottom=1.7mm,before skip=2mm,after skip=2mm]
\textbf{\color{teal}To'g'ri javob:} \textbf{A) Shaxsiy va maxfiy axborot}
\par\smallskip
\textbf{Izoh:} Vaziyat Shaxsiy va maxfiy axborot tushunchasini ifodalaydi: muayyan shaxsga tegishli yoki kirishi cheklangan axborot.
\end{tcolorbox}
\vspace{2mm}\hrule height 0.2pt \vspace{3mm}

% DB question_id: 1c7442ab-bd6d-5b04-ab57-dafaa1b65c73
\needspace{5\baselineskip}
\noindent\textbf{\color{ink}130.}
\par
Jonli grafik sensor qiymatini har soniya yangiladi. Ushbu vaziyatdagi asosiy tushunchani aniqlang.
\begin{enumerate}[label=\Alph*),leftmargin=9mm]
\item Dinamik axborot
\item Uzluksiz axborot
\item Diskret axborot
\item Ochiq va ommaviy axborot
\end{enumerate}
\begin{tcolorbox}[enhanced,breakable,colback=teal!4,colframe=teal!50!black,boxrule=0.5pt,arc=1.1mm,left=3mm,right=3mm,top=1.7mm,bottom=1.7mm,before skip=2mm,after skip=2mm]
\textbf{\color{teal}To'g'ri javob:} \textbf{A) Dinamik axborot}
\par\smallskip
\textbf{Izoh:} Vaziyat Dinamik axborot tushunchasini ifodalaydi: vaqt davomida holati o'zgarib boradigan axborot.
\end{tcolorbox}
\vspace{2mm}\hrule height 0.2pt \vspace{3mm}

% DB question_id: 21099094-3bb2-5ae2-b7ca-83ee02103cf6
\needspace{5\baselineskip}
\noindent\textbf{\color{ink}131.}
\par
Elektron hujjatda izohlar va sarlavhalar yozildi. Ushbu vaziyatdagi asosiy tushunchani aniqlang.
\begin{enumerate}[label=\Alph*),leftmargin=9mm]
\item Matnli axborot
\item Sonli axborot
\item Grafik axborot
\item Audioaxborot
\end{enumerate}
\begin{tcolorbox}[enhanced,breakable,colback=teal!4,colframe=teal!50!black,boxrule=0.5pt,arc=1.1mm,left=3mm,right=3mm,top=1.7mm,bottom=1.7mm,before skip=2mm,after skip=2mm]
\textbf{\color{teal}To'g'ri javob:} \textbf{A) Matnli axborot}
\par\smallskip
\textbf{Izoh:} Vaziyat Matnli axborot tushunchasini ifodalaydi: harf, raqam va yozuv belgilaridan tuzilgan axborot shakli.
\end{tcolorbox}
\vspace{2mm}\hrule height 0.2pt \vspace{3mm}

% DB question_id: 2125285e-d6db-59b8-8788-3fefee1d2dd4
\needspace{5\baselineskip}
\noindent\textbf{\color{ink}132.}
\par
Quyidagi ta'rif qaysi tushunchaga tegishli? Harf, raqam va yozuv belgilaridan tuzilgan axborot shakli.
\begin{enumerate}[label=\Alph*),leftmargin=9mm]
\item Matnli axborot
\item Sonli axborot
\item Grafik axborot
\item Audioaxborot
\end{enumerate}
\begin{tcolorbox}[enhanced,breakable,colback=teal!4,colframe=teal!50!black,boxrule=0.5pt,arc=1.1mm,left=3mm,right=3mm,top=1.7mm,bottom=1.7mm,before skip=2mm,after skip=2mm]
\textbf{\color{teal}To'g'ri javob:} \textbf{A) Matnli axborot}
\par\smallskip
\textbf{Izoh:} Ta'rif Matnli axborot tushunchasiga tegishli.
\end{tcolorbox}
\vspace{2mm}\hrule height 0.2pt \vspace{3mm}

% DB question_id: 274577b0-ae37-5e54-a2a3-da0c6378ac97
\needspace{5\baselineskip}
\noindent\textbf{\color{ink}133.}
\par
Qaysi holatda axborot foydali, ammo qimmatli bo'lmasligi mumkin?
\begin{enumerate}[label=\Alph*),leftmargin=9mm]
\item Yagona maxfiy kalit oshkor bo'ldi
\item Kundalik ob-havo ma'lumoti oddiy reja uchun kerak, lekin noyob emas
\item Tajriba natijasi noto'g'ri
\item Xabar tushunarsiz yozildi
\end{enumerate}
\begin{tcolorbox}[enhanced,breakable,colback=teal!4,colframe=teal!50!black,boxrule=0.5pt,arc=1.1mm,left=3mm,right=3mm,top=1.7mm,bottom=1.7mm,before skip=2mm,after skip=2mm]
\textbf{\color{teal}To'g'ri javob:} \textbf{B) Kundalik ob-havo ma'lumoti oddiy reja uchun kerak, lekin noyob emas}
\par\smallskip
\textbf{Izoh:} Foydalilik amaliy ehtiyojni, qiymatlilik esa ahamiyat va noyoblikni bildiradi.
\end{tcolorbox}
\vspace{2mm}\hrule height 0.2pt \vspace{3mm}

% DB question_id: 294938c5-2b5c-530b-a943-0a83bd1da8c4
\needspace{5\baselineskip}
\noindent\textbf{\color{ink}134.}
\par
Svetofor faqat qizil, sariq yoki yashil holatlardan birini qabul qildi. Ushbu vaziyatdagi asosiy tushunchani aniqlang.
\begin{enumerate}[label=\Alph*),leftmargin=9mm]
\item Ochiq va ommaviy axborot
\item Shaxsiy va maxfiy axborot
\item Diskret axborot
\item Birlamchi manba
\end{enumerate}
\begin{tcolorbox}[enhanced,breakable,colback=teal!4,colframe=teal!50!black,boxrule=0.5pt,arc=1.1mm,left=3mm,right=3mm,top=1.7mm,bottom=1.7mm,before skip=2mm,after skip=2mm]
\textbf{\color{teal}To'g'ri javob:} \textbf{C) Diskret axborot}
\par\smallskip
\textbf{Izoh:} Vaziyat Diskret axborot tushunchasini ifodalaydi: alohida, sanaladigan holat yoki qiymatlar ketma-ketligida ifodalanadigan axborot.
\end{tcolorbox}
\vspace{2mm}\hrule height 0.2pt \vspace{3mm}

% DB question_id: 2e918591-08a2-41fe-afc7-080f66c6b3e6
\needspace{5\baselineskip}
\noindent\textbf{\color{ink}135.}
\par
Wi-Fi orqali planshetga fayl yuborildi. Bu vaziyatdagi asosiy tushunchani aniqlang.
\begin{enumerate}[label=\Alph*),leftmargin=9mm]
\item Qabul qiluvchi
\item Birlamchi manba
\item Ikkilamchi manba
\item Aloqa kanali
\end{enumerate}
\begin{tcolorbox}[enhanced,breakable,colback=teal!4,colframe=teal!50!black,boxrule=0.5pt,arc=1.1mm,left=3mm,right=3mm,top=1.7mm,bottom=1.7mm,before skip=2mm,after skip=2mm]
\textbf{\color{teal}To'g'ri javob:} \textbf{D) Aloqa kanali}
\par\smallskip
\textbf{Izoh:} Vaziyat Aloqa kanaliga misol: xabar manbadan qabul qiluvchiga o'tadigan muhit yoki vosita.
\end{tcolorbox}
\vspace{2mm}\hrule height 0.2pt \vspace{3mm}

% DB question_id: 326bad08-56af-5aa9-b18d-d78bee0d1043
\needspace{5\baselineskip}
\noindent\textbf{\color{ink}136.}
\par
"Birlamchi manba" atamasining eng aniq ta'rifini toping.
\begin{enumerate}[label=\Alph*),leftmargin=9mm]
\item birlamchi manbalarni tahlil qiluvchi, sharhlovchi yoki umumlashtiruvchi manba
\item hodisa yoki tadqiqotga bevosita aloqador asl dalil yoki birinchi qo'l manbasi
\item bevosita manba hodisani to'g'ridan-to'g'ri, bilvosita manba esa vositachi orqali yetkazadi
\item harf, raqam va yozuv belgilaridan tuzilgan axborot shakli
\end{enumerate}
\begin{tcolorbox}[enhanced,breakable,colback=teal!4,colframe=teal!50!black,boxrule=0.5pt,arc=1.1mm,left=3mm,right=3mm,top=1.7mm,bottom=1.7mm,before skip=2mm,after skip=2mm]
\textbf{\color{teal}To'g'ri javob:} \textbf{B) hodisa yoki tadqiqotga bevosita aloqador asl dalil yoki birinchi qo'l manbasi}
\par\smallskip
\textbf{Izoh:} Birlamchi manba - hodisa yoki tadqiqotga bevosita aloqador asl dalil yoki birinchi qo'l manbasi.
\end{tcolorbox}
\vspace{2mm}\hrule height 0.2pt \vspace{3mm}

% DB question_id: 38484954-677c-5ac1-8621-c27a03f9a383
\needspace{5\baselineskip}
\noindent\textbf{\color{ink}137.}
\par
"Dinamik axborot" atamasining eng aniq ta'rifini toping.
\begin{enumerate}[label=\Alph*),leftmargin=9mm]
\item vaqt davomida holati o'zgarib boradigan axborot
\item qiymati ma'lum oraliqda uzluksiz o'zgarishi mumkin bo'lgan axborot
\item alohida, sanaladigan holat yoki qiymatlar ketma-ketligida ifodalanadigan axborot
\item foydalanishi cheklanmagan yoki keng jamoatchilikka e'lon qilingan axborot
\end{enumerate}
\begin{tcolorbox}[enhanced,breakable,colback=teal!4,colframe=teal!50!black,boxrule=0.5pt,arc=1.1mm,left=3mm,right=3mm,top=1.7mm,bottom=1.7mm,before skip=2mm,after skip=2mm]
\textbf{\color{teal}To'g'ri javob:} \textbf{A) vaqt davomida holati o'zgarib boradigan axborot}
\par\smallskip
\textbf{Izoh:} Dinamik axborot - vaqt davomida holati o'zgarib boradigan axborot.
\end{tcolorbox}
\vspace{2mm}\hrule height 0.2pt \vspace{3mm}

% DB question_id: 39eb2b9e-607f-5316-84aa-bd7ef2629779
\needspace{5\baselineskip}
\noindent\textbf{\color{ink}138.}
\par
Muzey eksponatining fotosurati ko'rsatildi. Ushbu vaziyatdagi asosiy tushunchani aniqlang.
\begin{enumerate}[label=\Alph*),leftmargin=9mm]
\item Dinamik axborot
\item Uzluksiz axborot
\item Statik axborot
\item Diskret axborot
\end{enumerate}
\begin{tcolorbox}[enhanced,breakable,colback=teal!4,colframe=teal!50!black,boxrule=0.5pt,arc=1.1mm,left=3mm,right=3mm,top=1.7mm,bottom=1.7mm,before skip=2mm,after skip=2mm]
\textbf{\color{teal}To'g'ri javob:} \textbf{C) Statik axborot}
\par\smallskip
\textbf{Izoh:} Vaziyat Statik axborot tushunchasini ifodalaydi: vaqt davomida o'zgarmaydigan yoki bitta holatni tasvirlaydigan axborot.
\end{tcolorbox}
\vspace{2mm}\hrule height 0.2pt \vspace{3mm}

% DB question_id: 3b2d1d58-7cc0-4575-a5a7-0c26337b9ef3
\needspace{5\baselineskip}
\noindent\textbf{\color{ink}139.}
\par
Tajriba paytida o'lchangan qiymatlar laboratoriya daftariga yozildi. Bu vaziyatdagi asosiy tushunchani aniqlang.
\begin{enumerate}[label=\Alph*),leftmargin=9mm]
\item Ikkilamchi manba
\item Birlamchi manba
\item Statik axborot
\item Dinamik axborot
\end{enumerate}
\begin{tcolorbox}[enhanced,breakable,colback=teal!4,colframe=teal!50!black,boxrule=0.5pt,arc=1.1mm,left=3mm,right=3mm,top=1.7mm,bottom=1.7mm,before skip=2mm,after skip=2mm]
\textbf{\color{teal}To'g'ri javob:} \textbf{B) Birlamchi manba}
\par\smallskip
\textbf{Izoh:} Vaziyat Birlamchi manbaga misol: hodisa yoki tadqiqotdan bevosita olingan asl dalil.
\end{tcolorbox}
\vspace{2mm}\hrule height 0.2pt \vspace{3mm}

% DB question_id: 3d9b3b7b-7222-5732-b101-3e9f5233d361
\needspace{5\baselineskip}
\noindent\textbf{\color{ink}140.}
\par
Quyidagi ta'rif qaysi tushunchaga tegishli? Vaqt davomida holati o'zgarib boradigan axborot.
\begin{enumerate}[label=\Alph*),leftmargin=9mm]
\item Uzluksiz axborot
\item Diskret axborot
\item Ochiq va ommaviy axborot
\item Dinamik axborot
\end{enumerate}
\begin{tcolorbox}[enhanced,breakable,colback=teal!4,colframe=teal!50!black,boxrule=0.5pt,arc=1.1mm,left=3mm,right=3mm,top=1.7mm,bottom=1.7mm,before skip=2mm,after skip=2mm]
\textbf{\color{teal}To'g'ri javob:} \textbf{D) Dinamik axborot}
\par\smallskip
\textbf{Izoh:} Ta'rif Dinamik axborot tushunchasiga tegishli.
\end{tcolorbox}
\vspace{2mm}\hrule height 0.2pt \vspace{3mm}

% DB question_id: 3da93a1c-600b-56ed-8211-9f533f96380f
\needspace{5\baselineskip}
\noindent\textbf{\color{ink}141.}
\par
Quyidagi ta'rif qaysi tushunchaga tegishli? Birlamchi manbalarni tahlil qiluvchi, sharhlovchi yoki umumlashtiruvchi manba.
\begin{enumerate}[label=\Alph*),leftmargin=9mm]
\item Bevosita va bilvosita manba
\item Ikkilamchi manba
\item Matnli axborot
\item Sonli axborot
\end{enumerate}
\begin{tcolorbox}[enhanced,breakable,colback=teal!4,colframe=teal!50!black,boxrule=0.5pt,arc=1.1mm,left=3mm,right=3mm,top=1.7mm,bottom=1.7mm,before skip=2mm,after skip=2mm]
\textbf{\color{teal}To'g'ri javob:} \textbf{B) Ikkilamchi manba}
\par\smallskip
\textbf{Izoh:} Ta'rif Ikkilamchi manba tushunchasiga tegishli.
\end{tcolorbox}
\vspace{2mm}\hrule height 0.2pt \vspace{3mm}

% DB question_id: 3fc69a62-8d43-5e1a-9751-ae758a7f1cce
\needspace{5\baselineskip}
\noindent\textbf{\color{ink}142.}
\par
"Sonli axborot" atamasining eng aniq ta'rifini toping.
\begin{enumerate}[label=\Alph*),leftmargin=9mm]
\item rasm, chizma, diagramma yoki sxema ko'rinishidagi axborot
\item tovush va nutq orqali ifodalanadigan axborot
\item miqdor va o'lchovlarni sonlar orqali ifodalovchi axborot
\item vaqt bo'yicha almashuvchi tasvirlar ketma-ketligi ko'rinishidagi axborot
\end{enumerate}
\begin{tcolorbox}[enhanced,breakable,colback=teal!4,colframe=teal!50!black,boxrule=0.5pt,arc=1.1mm,left=3mm,right=3mm,top=1.7mm,bottom=1.7mm,before skip=2mm,after skip=2mm]
\textbf{\color{teal}To'g'ri javob:} \textbf{C) miqdor va o'lchovlarni sonlar orqali ifodalovchi axborot}
\par\smallskip
\textbf{Izoh:} Sonli axborot - miqdor va o'lchovlarni sonlar orqali ifodalovchi axborot.
\end{tcolorbox}
\vspace{2mm}\hrule height 0.2pt \vspace{3mm}

% DB question_id: 41e83e01-dd6b-4280-8735-53e90f2ded03
\needspace{5\baselineskip}
\noindent\textbf{\color{ink}143.}
\par
Ob-havo natijalari ustunli diagrammada tasvirlandi. Bu vaziyatdagi asosiy tushunchani aniqlang.
\begin{enumerate}[label=\Alph*),leftmargin=9mm]
\item Audioaxborot
\item Multimedia
\item Grafik axborot
\item Vizual axborot
\end{enumerate}
\begin{tcolorbox}[enhanced,breakable,colback=teal!4,colframe=teal!50!black,boxrule=0.5pt,arc=1.1mm,left=3mm,right=3mm,top=1.7mm,bottom=1.7mm,before skip=2mm,after skip=2mm]
\textbf{\color{teal}To'g'ri javob:} \textbf{C) Grafik axborot}
\par\smallskip
\textbf{Izoh:} Vaziyat Grafik axborotga misol: rasm, chizma, diagramma yoki sxema ko'rinishidagi axborot.
\end{tcolorbox}
\vspace{2mm}\hrule height 0.2pt \vspace{3mm}

% DB question_id: 4345a41f-4e47-5ae1-9ab6-2a8f87abe509
\needspace{5\baselineskip}
\noindent\textbf{\color{ink}144.}
\par
"Kombinatsiyalashgan axborot" atamasining eng aniq ta'rifini toping.
\begin{enumerate}[label=\Alph*),leftmargin=9mm]
\item vaqt davomida o'zgarmaydigan yoki bitta holatni tasvirlaydigan axborot
\item vaqt davomida holati o'zgarib boradigan axborot
\item matn, grafika, audio yoki videoning birgalikdagi ko'rinishi
\item qiymati ma'lum oraliqda uzluksiz o'zgarishi mumkin bo'lgan axborot
\end{enumerate}
\begin{tcolorbox}[enhanced,breakable,colback=teal!4,colframe=teal!50!black,boxrule=0.5pt,arc=1.1mm,left=3mm,right=3mm,top=1.7mm,bottom=1.7mm,before skip=2mm,after skip=2mm]
\textbf{\color{teal}To'g'ri javob:} \textbf{C) matn, grafika, audio yoki videoning birgalikdagi ko'rinishi}
\par\smallskip
\textbf{Izoh:} Kombinatsiyalashgan axborot - matn, grafika, audio yoki videoning birgalikdagi ko'rinishi.
\end{tcolorbox}
\vspace{2mm}\hrule height 0.2pt \vspace{3mm}

% DB question_id: 4ab9e611-e990-52b5-9f29-53a1c9c22a67
\needspace{5\baselineskip}
\noindent\textbf{\color{ink}145.}
\par
Raqamli termometr ekrani 36,6; 36,7; 36,8 kabi qiymat ko'rsatadi. Ko'rsatish shakli qanday?
\begin{enumerate}[label=\Alph*),leftmargin=9mm]
\item Uzluksiz
\item Diskret
\item Faqat analog
\item Matnsiz
\end{enumerate}
\begin{tcolorbox}[enhanced,breakable,colback=teal!4,colframe=teal!50!black,boxrule=0.5pt,arc=1.1mm,left=3mm,right=3mm,top=1.7mm,bottom=1.7mm,before skip=2mm,after skip=2mm]
\textbf{\color{teal}To'g'ri javob:} \textbf{B) Diskret}
\par\smallskip
\textbf{Izoh:} Ekran alohida sonli qiymatlarni ko'rsatadi.
\end{tcolorbox}
\vspace{2mm}\hrule height 0.2pt \vspace{3mm}

% DB question_id: 4cc8e020-8734-50fe-a533-20e614356512
\needspace{5\baselineskip}
\noindent\textbf{\color{ink}146.}
\par
"Uzluksiz axborot" atamasining eng aniq ta'rifini toping.
\begin{enumerate}[label=\Alph*),leftmargin=9mm]
\item alohida, sanaladigan holat yoki qiymatlar ketma-ketligida ifodalanadigan axborot
\item qiymati ma'lum oraliqda uzluksiz o'zgarishi mumkin bo'lgan axborot
\item foydalanishi cheklanmagan yoki keng jamoatchilikka e'lon qilingan axborot
\item muayyan shaxsga tegishli yoki kirishi cheklangan axborot
\end{enumerate}
\begin{tcolorbox}[enhanced,breakable,colback=teal!4,colframe=teal!50!black,boxrule=0.5pt,arc=1.1mm,left=3mm,right=3mm,top=1.7mm,bottom=1.7mm,before skip=2mm,after skip=2mm]
\textbf{\color{teal}To'g'ri javob:} \textbf{B) qiymati ma'lum oraliqda uzluksiz o'zgarishi mumkin bo'lgan axborot}
\par\smallskip
\textbf{Izoh:} Uzluksiz axborot - qiymati ma'lum oraliqda uzluksiz o'zgarishi mumkin bo'lgan axborot.
\end{tcolorbox}
\vspace{2mm}\hrule height 0.2pt \vspace{3mm}

% DB question_id: 557b7341-93ec-4e2c-bb75-6001d8f365a7
\needspace{5\baselineskip}
\noindent\textbf{\color{ink}147.}
\par
Voqeani shaxsan kuzatgan guvoh intervyu berdi. Intervyuning birlamchi manba ekani uning ishonchliligi haqida qanday xulosaga olib keladi?
\begin{enumerate}[label=\Alph*),leftmargin=9mm]
\item To'g'ri: bevosita olingan har qanday axborotni tekshirish shart emas
\item To'g'ri: guvoh bergan ma'lumotlar doimo tarafkashlikdan xoli bo'ladi
\item Noto'g'ri: birlamchi manba faqat boshqa manbalardan yig'ilgan sharhdir
\item Noto'g'ri: asl manbada ham xato yoki tarafkashlik bo'lishi mumkin
\end{enumerate}
\begin{tcolorbox}[enhanced,breakable,colback=teal!4,colframe=teal!50!black,boxrule=0.5pt,arc=1.1mm,left=3mm,right=3mm,top=1.7mm,bottom=1.7mm,before skip=2mm,after skip=2mm]
\textbf{\color{teal}To'g'ri javob:} \textbf{D) Noto'g'ri: asl manbada ham xato yoki tarafkashlik bo'lishi mumkin}
\par\smallskip
\textbf{Izoh:} Birlamchi manba ma'lumotni bevosita kuzatgan yoki qayd qilgan shaxs yoki yozuv bo'lishi mumkin. Biroq xotira xatosi, noto'liq tasvir yoki tarafkashlik bo'lishi mumkinligi sabab uni ham tekshirish zarur.
\end{tcolorbox}
\vspace{2mm}\hrule height 0.2pt \vspace{3mm}

% DB question_id: 560ec44d-963e-42ee-a928-03f8851315b5
\needspace{5\baselineskip}
\noindent\textbf{\color{ink}148.}
\par
Interaktiv darsda matn, rasm, ovoz va video bir sahifada ishlatildi. Bu vaziyatdagi asosiy tushunchani aniqlang.
\begin{enumerate}[label=\Alph*),leftmargin=9mm]
\item Multimedia
\item Vizual axborot
\item Analog axborot
\item Diskret axborot
\end{enumerate}
\begin{tcolorbox}[enhanced,breakable,colback=teal!4,colframe=teal!50!black,boxrule=0.5pt,arc=1.1mm,left=3mm,right=3mm,top=1.7mm,bottom=1.7mm,before skip=2mm,after skip=2mm]
\textbf{\color{teal}To'g'ri javob:} \textbf{A) Multimedia}
\par\smallskip
\textbf{Izoh:} Vaziyat Multimediaga misol: matn, grafika, audio, video va animatsiyaning uyg'un birikmasi.
\end{tcolorbox}
\vspace{2mm}\hrule height 0.2pt \vspace{3mm}

% DB question_id: 591ad681-018d-5b47-93ed-fbdce716d8fe
\needspace{5\baselineskip}
\noindent\textbf{\color{ink}149.}
\par
Quyidagi ta'rif qaysi tushunchaga tegishli? Foydalanishi cheklanmagan yoki keng jamoatchilikka e'lon qilingan axborot.
\begin{enumerate}[label=\Alph*),leftmargin=9mm]
\item Shaxsiy va maxfiy axborot
\item Birlamchi manba
\item Ochiq va ommaviy axborot
\item Ikkilamchi manba
\end{enumerate}
\begin{tcolorbox}[enhanced,breakable,colback=teal!4,colframe=teal!50!black,boxrule=0.5pt,arc=1.1mm,left=3mm,right=3mm,top=1.7mm,bottom=1.7mm,before skip=2mm,after skip=2mm]
\textbf{\color{teal}To'g'ri javob:} \textbf{C) Ochiq va ommaviy axborot}
\par\smallskip
\textbf{Izoh:} Ta'rif Ochiq va ommaviy axborot tushunchasiga tegishli.
\end{tcolorbox}
\vspace{2mm}\hrule height 0.2pt \vspace{3mm}

% DB question_id: 596e215d-e4d9-5009-b812-a87c9921fee1
\needspace{5\baselineskip}
\noindent\textbf{\color{ink}150.}
\par
Quyidagi ta'rif qaysi tushunchaga tegishli? Tovush va nutq orqali ifodalanadigan axborot.
\begin{enumerate}[label=\Alph*),leftmargin=9mm]
\item Videoaxborot
\item Kombinatsiyalashgan axborot
\item Statik axborot
\item Audioaxborot
\end{enumerate}
\begin{tcolorbox}[enhanced,breakable,colback=teal!4,colframe=teal!50!black,boxrule=0.5pt,arc=1.1mm,left=3mm,right=3mm,top=1.7mm,bottom=1.7mm,before skip=2mm,after skip=2mm]
\textbf{\color{teal}To'g'ri javob:} \textbf{D) Audioaxborot}
\par\smallskip
\textbf{Izoh:} Ta'rif Audioaxborot tushunchasiga tegishli.
\end{tcolorbox}
\vspace{2mm}\hrule height 0.2pt \vspace{3mm}

% DB question_id: 60670eb2-1187-4c7a-9c73-ea4a0385ddf7
\needspace{5\baselineskip}
\noindent\textbf{\color{ink}151.}
\par
Atamalarni ularga mos raqamli ifodalash tavsifi bilan juftlang.
\textbf{\footnotesize Chap ustun:}
\begin{enumerate}[label=\arabic*),leftmargin=9mm]
\item ASCII
\item UTF-8
\item RGB
\item PCM
\end{enumerate}
\textbf{\footnotesize O'ng ustun:}
\begin{enumerate}[label=\Alph*),leftmargin=9mm]
\item Qizil, yashil va ko'k kanallar orqali rang ifodalash
\item Asl ko'rinishida 7 bitli belgilar kodi
\item Audio namunalarini raqamli qiymatlar ketma-ketligida saqlash
\item Unicode kod nuqtalarini 1-4 baytda kodlash shakli
\end{enumerate}
\begin{tcolorbox}[enhanced,breakable,colback=teal!4,colframe=teal!50!black,boxrule=0.5pt,arc=1.1mm,left=3mm,right=3mm,top=1.7mm,bottom=1.7mm,before skip=2mm,after skip=2mm]
\textbf{\color{teal}To'g'ri javob:} \textbf{1-B, 2-D, 3-A, 4-C}
\par\smallskip
\textbf{Izoh:} ASCII - 7 bitli kod; UTF-8 - Unicode encoding; RGB - rang modeli; PCM - raqamli audio namunalarini ifodalash usuli.
\end{tcolorbox}
\vspace{2mm}\hrule height 0.2pt \vspace{3mm}

% DB question_id: 60b09194-8d8f-5b1a-8f46-aa435010e860
\needspace{5\baselineskip}
\noindent\textbf{\color{ink}152.}
\par
"Audioaxborot" atamasining eng aniq ta'rifini toping.
\begin{enumerate}[label=\Alph*),leftmargin=9mm]
\item tovush va nutq orqali ifodalanadigan axborot
\item vaqt bo'yicha almashuvchi tasvirlar ketma-ketligi ko'rinishidagi axborot
\item matn, grafika, audio yoki videoning birgalikdagi ko'rinishi
\item vaqt davomida o'zgarmaydigan yoki bitta holatni tasvirlaydigan axborot
\end{enumerate}
\begin{tcolorbox}[enhanced,breakable,colback=teal!4,colframe=teal!50!black,boxrule=0.5pt,arc=1.1mm,left=3mm,right=3mm,top=1.7mm,bottom=1.7mm,before skip=2mm,after skip=2mm]
\textbf{\color{teal}To'g'ri javob:} \textbf{A) tovush va nutq orqali ifodalanadigan axborot}
\par\smallskip
\textbf{Izoh:} Audioaxborot - tovush va nutq orqali ifodalanadigan axborot.
\end{tcolorbox}
\vspace{2mm}\hrule height 0.2pt \vspace{3mm}

% DB question_id: 643d9cc6-ba29-5a1a-92ef-8b7ed0a54ecc
\needspace{5\baselineskip}
\noindent\textbf{\color{ink}153.}
\par
"Shaxsiy va maxfiy axborot" atamasining eng aniq ta'rifini toping.
\begin{enumerate}[label=\Alph*),leftmargin=9mm]
\item muayyan shaxsga tegishli yoki kirishi cheklangan axborot
\item hodisa yoki tadqiqotga bevosita aloqador asl dalil yoki birinchi qo'l manbasi
\item birlamchi manbalarni tahlil qiluvchi, sharhlovchi yoki umumlashtiruvchi manba
\item bevosita manba hodisani to'g'ridan-to'g'ri, bilvosita manba esa vositachi orqali yetkazadi
\end{enumerate}
\begin{tcolorbox}[enhanced,breakable,colback=teal!4,colframe=teal!50!black,boxrule=0.5pt,arc=1.1mm,left=3mm,right=3mm,top=1.7mm,bottom=1.7mm,before skip=2mm,after skip=2mm]
\textbf{\color{teal}To'g'ri javob:} \textbf{A) muayyan shaxsga tegishli yoki kirishi cheklangan axborot}
\par\smallskip
\textbf{Izoh:} Shaxsiy va maxfiy axborot - muayyan shaxsga tegishli yoki kirishi cheklangan axborot.
\end{tcolorbox}
\vspace{2mm}\hrule height 0.2pt \vspace{3mm}

% DB question_id: 65bd2cb5-08f9-5a14-a8b1-973b1e322419
\needspace{5\baselineskip}
\noindent\textbf{\color{ink}154.}
\par
Kompaniya ichki savdo bazasi va davlat statistika sayti mos ravishda qanday manba?
\begin{enumerate}[label=\Alph*),leftmargin=9mm]
\item Birlamchi va statik
\item Bevosita va maxfiy
\item Ichki va tashqi
\item Ochiq va dinamik
\end{enumerate}
\begin{tcolorbox}[enhanced,breakable,colback=teal!4,colframe=teal!50!black,boxrule=0.5pt,arc=1.1mm,left=3mm,right=3mm,top=1.7mm,bottom=1.7mm,before skip=2mm,after skip=2mm]
\textbf{\color{teal}To'g'ri javob:} \textbf{C) Ichki va tashqi}
\par\smallskip
\textbf{Izoh:} Kompaniya bazasi tashkilot ichidan, davlat sayti esa tashqaridan olinadi.
\end{tcolorbox}
\vspace{2mm}\hrule height 0.2pt \vspace{3mm}

% DB question_id: 6dcede2d-1a2b-5278-b77d-3b7c188facf0
\needspace{5\baselineskip}
\noindent\textbf{\color{ink}155.}
\par
Simobli termometr ustuni harorat bilan silliq ko'tarildi. Ushbu vaziyatdagi asosiy tushunchani aniqlang.
\begin{enumerate}[label=\Alph*),leftmargin=9mm]
\item Diskret axborot
\item Uzluksiz axborot
\item Ochiq va ommaviy axborot
\item Shaxsiy va maxfiy axborot
\end{enumerate}
\begin{tcolorbox}[enhanced,breakable,colback=teal!4,colframe=teal!50!black,boxrule=0.5pt,arc=1.1mm,left=3mm,right=3mm,top=1.7mm,bottom=1.7mm,before skip=2mm,after skip=2mm]
\textbf{\color{teal}To'g'ri javob:} \textbf{B) Uzluksiz axborot}
\par\smallskip
\textbf{Izoh:} Vaziyat Uzluksiz axborot tushunchasini ifodalaydi: qiymati ma'lum oraliqda uzluksiz o'zgarishi mumkin bo'lgan axborot.
\end{tcolorbox}
\vspace{2mm}\hrule height 0.2pt \vspace{3mm}

% DB question_id: 6e16e700-785a-5443-a8fb-3e42ec0dcfe1
\needspace{5\baselineskip}
\noindent\textbf{\color{ink}156.}
\par
"Bevosita va bilvosita manba" atamasining eng aniq ta'rifini toping.
\begin{enumerate}[label=\Alph*),leftmargin=9mm]
\item harf, raqam va yozuv belgilaridan tuzilgan axborot shakli
\item miqdor va o'lchovlarni sonlar orqali ifodalovchi axborot
\item rasm, chizma, diagramma yoki sxema ko'rinishidagi axborot
\item bevosita manba hodisani to'g'ridan-to'g'ri, bilvosita manba esa vositachi orqali yetkazadi
\end{enumerate}
\begin{tcolorbox}[enhanced,breakable,colback=teal!4,colframe=teal!50!black,boxrule=0.5pt,arc=1.1mm,left=3mm,right=3mm,top=1.7mm,bottom=1.7mm,before skip=2mm,after skip=2mm]
\textbf{\color{teal}To'g'ri javob:} \textbf{D) bevosita manba hodisani to'g'ridan-to'g'ri, bilvosita manba esa vositachi orqali yetkazadi}
\par\smallskip
\textbf{Izoh:} Bevosita va bilvosita manba - bevosita manba hodisani to'g'ridan-to'g'ri, bilvosita manba esa vositachi orqali yetkazadi.
\end{tcolorbox}
\vspace{2mm}\hrule height 0.2pt \vspace{3mm}

% DB question_id: 6f5fa8e5-e75e-5c63-afe3-0e832c8ebbee
\needspace{5\baselineskip}
\noindent\textbf{\color{ink}157.}
\par
O'lchov asbobi 25,0 \ensuremath{^\circ}C o'rniga 28,0 \ensuremath{^\circ}C ko'rsatmoqda. Qaysi xossa buzilgan?
\begin{enumerate}[label=\Alph*),leftmargin=9mm]
\item Aniqlik
\item Foydalilik
\item Ommaviylik
\item Dinamiklik
\end{enumerate}
\begin{tcolorbox}[enhanced,breakable,colback=teal!4,colframe=teal!50!black,boxrule=0.5pt,arc=1.1mm,left=3mm,right=3mm,top=1.7mm,bottom=1.7mm,before skip=2mm,after skip=2mm]
\textbf{\color{teal}To'g'ri javob:} \textbf{A) Aniqlik}
\par\smallskip
\textbf{Izoh:} Ko'rsatkich haqiqiy qiymatga mos emas, demak aniqlik buzilgan.
\end{tcolorbox}
\vspace{2mm}\hrule height 0.2pt \vspace{3mm}

% DB question_id: 707b55d5-2bbb-5499-9807-79dfdbb2a7d1
\needspace{5\baselineskip}
\noindent\textbf{\color{ink}158.}
\par
"Ikkilamchi manba" atamasining eng aniq ta'rifini toping.
\begin{enumerate}[label=\Alph*),leftmargin=9mm]
\item bevosita manba hodisani to'g'ridan-to'g'ri, bilvosita manba esa vositachi orqali yetkazadi
\item harf, raqam va yozuv belgilaridan tuzilgan axborot shakli
\item birlamchi manbalarni tahlil qiluvchi, sharhlovchi yoki umumlashtiruvchi manba
\item miqdor va o'lchovlarni sonlar orqali ifodalovchi axborot
\end{enumerate}
\begin{tcolorbox}[enhanced,breakable,colback=teal!4,colframe=teal!50!black,boxrule=0.5pt,arc=1.1mm,left=3mm,right=3mm,top=1.7mm,bottom=1.7mm,before skip=2mm,after skip=2mm]
\textbf{\color{teal}To'g'ri javob:} \textbf{C) birlamchi manbalarni tahlil qiluvchi, sharhlovchi yoki umumlashtiruvchi manba}
\par\smallskip
\textbf{Izoh:} Ikkilamchi manba - birlamchi manbalarni tahlil qiluvchi, sharhlovchi yoki umumlashtiruvchi manba.
\end{tcolorbox}
\vspace{2mm}\hrule height 0.2pt \vspace{3mm}

% DB question_id: 7516d593-8c2d-4f04-96c3-60b088dacd5e
\needspace{5\baselineskip}
\noindent\textbf{\color{ink}159.}
\par
Yo'l belgisining rangi va shakli haydovchiga ogohlantirish berdi. Bu vaziyatdagi asosiy tushunchani aniqlang.
\begin{enumerate}[label=\Alph*),leftmargin=9mm]
\item Analog axborot
\item Vizual axborot
\item Diskret axborot
\item Aniqlik
\end{enumerate}
\begin{tcolorbox}[enhanced,breakable,colback=teal!4,colframe=teal!50!black,boxrule=0.5pt,arc=1.1mm,left=3mm,right=3mm,top=1.7mm,bottom=1.7mm,before skip=2mm,after skip=2mm]
\textbf{\color{teal}To'g'ri javob:} \textbf{B) Vizual axborot}
\par\smallskip
\textbf{Izoh:} Vaziyat Vizual axborotga misol: ko'rish a'zosi orqali qabul qilinadigan axborot.
\end{tcolorbox}
\vspace{2mm}\hrule height 0.2pt \vspace{3mm}

% DB question_id: 764e68d3-e814-5687-af7d-b307e662ea17
\needspace{5\baselineskip}
\noindent\textbf{\color{ink}160.}
\par
Maqola faqat muallifning shaxsiy qarashlarini dalilsiz himoya qiladi. Qaysi xossa eng zaif?
\begin{enumerate}[label=\Alph*),leftmargin=9mm]
\item Tushunarlilik
\item Diskretlik
\item Maxfiylik
\item Obyektivlik
\end{enumerate}
\begin{tcolorbox}[enhanced,breakable,colback=teal!4,colframe=teal!50!black,boxrule=0.5pt,arc=1.1mm,left=3mm,right=3mm,top=1.7mm,bottom=1.7mm,before skip=2mm,after skip=2mm]
\textbf{\color{teal}To'g'ri javob:} \textbf{D) Obyektivlik}
\par\smallskip
\textbf{Izoh:} Bir tomonlama qarash obyektivlikni pasaytiradi.
\end{tcolorbox}
\vspace{2mm}\hrule height 0.2pt \vspace{3mm}

% DB question_id: 76993e84-21cf-5e75-9822-b045bad783a6
\needspace{5\baselineskip}
\noindent\textbf{\color{ink}161.}
\par
Jonli efirdagi video bir vaqtning o'zida qaysi ikki tasnifga mos?
\begin{enumerate}[label=\Alph*),leftmargin=9mm]
\item Matnli va statik
\item Sonli va diskret
\item Grafik va maxfiy
\item Video va dinamik
\end{enumerate}
\begin{tcolorbox}[enhanced,breakable,colback=teal!4,colframe=teal!50!black,boxrule=0.5pt,arc=1.1mm,left=3mm,right=3mm,top=1.7mm,bottom=1.7mm,before skip=2mm,after skip=2mm]
\textbf{\color{teal}To'g'ri javob:} \textbf{D) Video va dinamik}
\par\smallskip
\textbf{Izoh:} U tasvirlar ketma-ketligi va vaqt bo'yicha o'zgaruvchi axborotdir.
\end{tcolorbox}
\vspace{2mm}\hrule height 0.2pt \vspace{3mm}

% DB question_id: 7860dc5d-5a27-5e86-a10d-660b662fd225
\needspace{5\baselineskip}
\noindent\textbf{\color{ink}162.}
\par
"Statik axborot" atamasining eng aniq ta'rifini toping.
\begin{enumerate}[label=\Alph*),leftmargin=9mm]
\item vaqt davomida holati o'zgarib boradigan axborot
\item qiymati ma'lum oraliqda uzluksiz o'zgarishi mumkin bo'lgan axborot
\item alohida, sanaladigan holat yoki qiymatlar ketma-ketligida ifodalanadigan axborot
\item vaqt davomida o'zgarmaydigan yoki bitta holatni tasvirlaydigan axborot
\end{enumerate}
\begin{tcolorbox}[enhanced,breakable,colback=teal!4,colframe=teal!50!black,boxrule=0.5pt,arc=1.1mm,left=3mm,right=3mm,top=1.7mm,bottom=1.7mm,before skip=2mm,after skip=2mm]
\textbf{\color{teal}To'g'ri javob:} \textbf{D) vaqt davomida o'zgarmaydigan yoki bitta holatni tasvirlaydigan axborot}
\par\smallskip
\textbf{Izoh:} Statik axborot - vaqt davomida o'zgarmaydigan yoki bitta holatni tasvirlaydigan axborot.
\end{tcolorbox}
\vspace{2mm}\hrule height 0.2pt \vspace{3mm}

% DB question_id: 78712de9-b96b-4f84-852f-19e3660f4508
\needspace{5\baselineskip}
\noindent\textbf{\color{ink}163.}
\par
50 nafar xodimga takrorlanmaydigan 0 va 1 lardan iborat qat'iy uzunlikdagi kod ajratish uchun kodning minimal uzunligi necha bit?
\begin{enumerate}[label=\Alph*),leftmargin=9mm]
\item 6
\item 8
\item 5
\item 4
\end{enumerate}
\begin{tcolorbox}[enhanced,breakable,colback=teal!4,colframe=teal!50!black,boxrule=0.5pt,arc=1.1mm,left=3mm,right=3mm,top=1.7mm,bottom=1.7mm,before skip=2mm,after skip=2mm]
\textbf{\color{teal}To'g'ri javob:} \textbf{A) 6}
\par\smallskip
\textbf{Izoh:} 2\textsuperscript{5}=32\textless{}50\ensuremath{\leq}64=2\textsuperscript{6}. 6 bit kerak.
\end{tcolorbox}
\vspace{2mm}\hrule height 0.2pt \vspace{3mm}

% DB question_id: 7b64cf3e-c7c0-42d0-96ef-166bb64fd216
\needspace{5\baselineskip}
\noindent\textbf{\color{ink}164.}
\par
Datchik xonadagi harorat qiymatini o'lchab uzatdi. Bu vaziyatdagi asosiy tushunchani aniqlang.
\begin{enumerate}[label=\Alph*),leftmargin=9mm]
\item Aloqa kanali
\item Qabul qiluvchi
\item Axborot manbai
\item Birlamchi manba
\end{enumerate}
\begin{tcolorbox}[enhanced,breakable,colback=teal!4,colframe=teal!50!black,boxrule=0.5pt,arc=1.1mm,left=3mm,right=3mm,top=1.7mm,bottom=1.7mm,before skip=2mm,after skip=2mm]
\textbf{\color{teal}To'g'ri javob:} \textbf{C) Axborot manbai}
\par\smallskip
\textbf{Izoh:} Vaziyat Axborot manbaiga misol: xabarni hosil qiluvchi yoki axborot kelib chiqadigan obyekt.
\end{tcolorbox}
\vspace{2mm}\hrule height 0.2pt \vspace{3mm}

% DB question_id: 7db1b04b-4783-438e-a1a8-c1acf9f6dc9a
\needspace{5\baselineskip}
\noindent\textbf{\color{ink}165.}
\par
Hodisa qatnashchisining bevosita yozuvi bilan boshqa kishining aynan shu voqea haqidagi ikkilamchi sharhi solishtirildi. Manbaning voqeaga yaqinligini qaysi juftlik farqlaydi?
\begin{enumerate}[label=\Alph*),leftmargin=9mm]
\item Ochiq va cheklangan manba
\item Mahalliy va masofaviy manba
\item Bevosita va bilvosita manba
\item Raqamli va qog'oz manba
\end{enumerate}
\begin{tcolorbox}[enhanced,breakable,colback=teal!4,colframe=teal!50!black,boxrule=0.5pt,arc=1.1mm,left=3mm,right=3mm,top=1.7mm,bottom=1.7mm,before skip=2mm,after skip=2mm]
\textbf{\color{teal}To'g'ri javob:} \textbf{C) Bevosita va bilvosita manba}
\par\smallskip
\textbf{Izoh:} Bevosita manba hodisani to'g'ridan-to'g'ri kuzatuvchi yoki qayd etuvchidan olinadi; bilvosita manba esa oraliq sharhlash yoki qayta bayon orqali yetadi. Manba bevosita bo'lishi uning xatosizligini kafolatlamaydi.
\end{tcolorbox}
\vspace{2mm}\hrule height 0.2pt \vspace{3mm}

% DB question_id: 808a4ea4-759c-43ba-9e66-2d229dfbc8fa
\needspace{5\baselineskip}
\noindent\textbf{\color{ink}166.}
\par
Bir slaydda matn, diagramma va ovozli izoh uyg'un holda ishlatilgan. Bir nechta axborot shakli bitta mahsulotda jamlangan bunday taqdimot turini qanday atash ma'qul?
\begin{enumerate}[label=\Alph*),leftmargin=9mm]
\item Statik tasvirli axborot
\item Kombinatsiyalashgan axborot
\item Dinamik video axboroti
\item Yakkalangan audio axborot
\end{enumerate}
\begin{tcolorbox}[enhanced,breakable,colback=teal!4,colframe=teal!50!black,boxrule=0.5pt,arc=1.1mm,left=3mm,right=3mm,top=1.7mm,bottom=1.7mm,before skip=2mm,after skip=2mm]
\textbf{\color{teal}To'g'ri javob:} \textbf{B) Kombinatsiyalashgan axborot}
\par\smallskip
\textbf{Izoh:} Bir mahsulotda matn, grafika va audio turli ifodalash shakllari birlashganda kombinatsiyalashgan axborot yuzaga keladi. Faqat bitta media shakli bunday vaziyatni to'liq ifodalamaydi.
\end{tcolorbox}
\vspace{2mm}\hrule height 0.2pt \vspace{3mm}

% DB question_id: 83b3bed7-a46d-560c-b2b0-1826f00a2091
\needspace{5\baselineskip}
\noindent\textbf{\color{ink}167.}
\par
Vazirlikning rasmiy saytida umumiy hisobot chop etildi. Ushbu vaziyatdagi asosiy tushunchani aniqlang.
\begin{enumerate}[label=\Alph*),leftmargin=9mm]
\item Shaxsiy va maxfiy axborot
\item Birlamchi manba
\item Ikkilamchi manba
\item Ochiq va ommaviy axborot
\end{enumerate}
\begin{tcolorbox}[enhanced,breakable,colback=teal!4,colframe=teal!50!black,boxrule=0.5pt,arc=1.1mm,left=3mm,right=3mm,top=1.7mm,bottom=1.7mm,before skip=2mm,after skip=2mm]
\textbf{\color{teal}To'g'ri javob:} \textbf{D) Ochiq va ommaviy axborot}
\par\smallskip
\textbf{Izoh:} Vaziyat Ochiq va ommaviy axborot tushunchasini ifodalaydi: foydalanishi cheklanmagan yoki keng jamoatchilikka e'lon qilingan axborot.
\end{tcolorbox}
\vspace{2mm}\hrule height 0.2pt \vspace{3mm}

% DB question_id: 88f14854-6524-5522-a217-fbeb95d65258
\needspace{5\baselineskip}
\noindent\textbf{\color{ink}168.}
\par
Quyidagi ta'rif qaysi tushunchaga tegishli? Vaqt bo'yicha almashuvchi tasvirlar ketma-ketligi ko'rinishidagi axborot.
\begin{enumerate}[label=\Alph*),leftmargin=9mm]
\item Videoaxborot
\item Kombinatsiyalashgan axborot
\item Statik axborot
\item Dinamik axborot
\end{enumerate}
\begin{tcolorbox}[enhanced,breakable,colback=teal!4,colframe=teal!50!black,boxrule=0.5pt,arc=1.1mm,left=3mm,right=3mm,top=1.7mm,bottom=1.7mm,before skip=2mm,after skip=2mm]
\textbf{\color{teal}To'g'ri javob:} \textbf{A) Videoaxborot}
\par\smallskip
\textbf{Izoh:} Ta'rif Videoaxborot tushunchasiga tegishli.
\end{tcolorbox}
\vspace{2mm}\hrule height 0.2pt \vspace{3mm}

% DB question_id: 8bdbc97b-69be-4630-a07d-e23cf49d48fe
\needspace{5\baselineskip}
\noindent\textbf{\color{ink}169.}
\par
Elektron xatda gaplar va ro'yxatlar yozilgan. Bu vaziyatdagi asosiy tushunchani aniqlang.
\begin{enumerate}[label=\Alph*),leftmargin=9mm]
\item Grafik axborot
\item Matnli axborot
\item Audioaxborot
\item Multimedia
\end{enumerate}
\begin{tcolorbox}[enhanced,breakable,colback=teal!4,colframe=teal!50!black,boxrule=0.5pt,arc=1.1mm,left=3mm,right=3mm,top=1.7mm,bottom=1.7mm,before skip=2mm,after skip=2mm]
\textbf{\color{teal}To'g'ri javob:} \textbf{B) Matnli axborot}
\par\smallskip
\textbf{Izoh:} Vaziyat Matnli axborotga misol: harf, raqam va tinish belgilaridan tuzilgan yozma axborot.
\end{tcolorbox}
\vspace{2mm}\hrule height 0.2pt \vspace{3mm}

% DB question_id: 8f679434-7c2b-5b4b-8833-cc03485c643a
\needspace{5\baselineskip}
\noindent\textbf{\color{ink}170.}
\par
Natijalar ustunli diagrammada tasvirlandi. Ushbu vaziyatdagi asosiy tushunchani aniqlang.
\begin{enumerate}[label=\Alph*),leftmargin=9mm]
\item Audioaxborot
\item Videoaxborot
\item Grafik axborot
\item Kombinatsiyalashgan axborot
\end{enumerate}
\begin{tcolorbox}[enhanced,breakable,colback=teal!4,colframe=teal!50!black,boxrule=0.5pt,arc=1.1mm,left=3mm,right=3mm,top=1.7mm,bottom=1.7mm,before skip=2mm,after skip=2mm]
\textbf{\color{teal}To'g'ri javob:} \textbf{C) Grafik axborot}
\par\smallskip
\textbf{Izoh:} Vaziyat Grafik axborot tushunchasini ifodalaydi: rasm, chizma, diagramma yoki sxema ko'rinishidagi axborot.
\end{tcolorbox}
\vspace{2mm}\hrule height 0.2pt \vspace{3mm}

% DB question_id: 93ec6428-37f7-5c33-a3e9-ac51b4d0f839
\needspace{5\baselineskip}
\noindent\textbf{\color{ink}171.}
\par
Quyidagi ta'rif qaysi tushunchaga tegishli? Hodisa yoki tadqiqotga bevosita aloqador asl dalil yoki birinchi qo'l manbasi.
\begin{enumerate}[label=\Alph*),leftmargin=9mm]
\item Birlamchi manba
\item Ikkilamchi manba
\item Bevosita va bilvosita manba
\item Matnli axborot
\end{enumerate}
\begin{tcolorbox}[enhanced,breakable,colback=teal!4,colframe=teal!50!black,boxrule=0.5pt,arc=1.1mm,left=3mm,right=3mm,top=1.7mm,bottom=1.7mm,before skip=2mm,after skip=2mm]
\textbf{\color{teal}To'g'ri javob:} \textbf{A) Birlamchi manba}
\par\smallskip
\textbf{Izoh:} Ta'rif Birlamchi manba tushunchasiga tegishli.
\end{tcolorbox}
\vspace{2mm}\hrule height 0.2pt \vspace{3mm}

% DB question_id: 966750d0-6156-4a02-a167-54631f410fd3
\needspace{5\baselineskip}
\noindent\textbf{\color{ink}172.}
\par
Maktab sport zalidagi harorat xavfsizmi-yo'qmi, degan qarorni qabul qilish kerak. Qaysi axborot xabari qarorga eng ishonchli asos bo'ladi?
\begin{enumerate}[label=\Alph*),leftmargin=9mm]
\item Manbasi, sanasi, o'lchov usuli va to'liq natijalari ko'rsatilgan xabar
\item Muallif va sana berilgan, biroq o'lchov usuli ko'rsatilmagan xabar
\item Natijalar chiroyli diagrammada, ammo o'lchangan kun sanasi noma'lum xabar
\item Bir nechta odam fikri keltirilgan, lekin manbasi tekshirilmagan xabar
\end{enumerate}
\begin{tcolorbox}[enhanced,breakable,colback=teal!4,colframe=teal!50!black,boxrule=0.5pt,arc=1.1mm,left=3mm,right=3mm,top=1.7mm,bottom=1.7mm,before skip=2mm,after skip=2mm]
\textbf{\color{teal}To'g'ri javob:} \textbf{A) Manbasi, sanasi, o'lchov usuli va to'liq natijalari ko'rsatilgan xabar}
\par\smallskip
\textbf{Izoh:} Qaror uchun o'lchovning kim tomonidan, qachon va qanday usulda o'tkazilgani hamda to'liq natijalari muhim. Faqat chiroyli tasvir yoki guvohlarning shaxsiy fikri aniqlik va tekshiriluvchanlikni ta'minlamaydi.
\end{tcolorbox}
\vspace{2mm}\hrule height 0.2pt \vspace{3mm}

% DB question_id: 9727036e-efc1-5b6a-9c4f-6bb4e6101802
\needspace{5\baselineskip}
\noindent\textbf{\color{ink}173.}
\par
"Ochiq va ommaviy axborot" atamasining eng aniq ta'rifini toping.
\begin{enumerate}[label=\Alph*),leftmargin=9mm]
\item muayyan shaxsga tegishli yoki kirishi cheklangan axborot
\item hodisa yoki tadqiqotga bevosita aloqador asl dalil yoki birinchi qo'l manbasi
\item birlamchi manbalarni tahlil qiluvchi, sharhlovchi yoki umumlashtiruvchi manba
\item foydalanishi cheklanmagan yoki keng jamoatchilikka e'lon qilingan axborot
\end{enumerate}
\begin{tcolorbox}[enhanced,breakable,colback=teal!4,colframe=teal!50!black,boxrule=0.5pt,arc=1.1mm,left=3mm,right=3mm,top=1.7mm,bottom=1.7mm,before skip=2mm,after skip=2mm]
\textbf{\color{teal}To'g'ri javob:} \textbf{D) foydalanishi cheklanmagan yoki keng jamoatchilikka e'lon qilingan axborot}
\par\smallskip
\textbf{Izoh:} Ochiq va ommaviy axborot - foydalanishi cheklanmagan yoki keng jamoatchilikka e'lon qilingan axborot.
\end{tcolorbox}
\vspace{2mm}\hrule height 0.2pt \vspace{3mm}

% DB question_id: 9a585b5d-caa3-57d4-a094-b612a12db264
\needspace{5\baselineskip}
\noindent\textbf{\color{ink}174.}
\par
Ovozli xabarda yig'ilish vaqti aytildi. Ushbu vaziyatdagi asosiy tushunchani aniqlang.
\begin{enumerate}[label=\Alph*),leftmargin=9mm]
\item Videoaxborot
\item Kombinatsiyalashgan axborot
\item Statik axborot
\item Audioaxborot
\end{enumerate}
\begin{tcolorbox}[enhanced,breakable,colback=teal!4,colframe=teal!50!black,boxrule=0.5pt,arc=1.1mm,left=3mm,right=3mm,top=1.7mm,bottom=1.7mm,before skip=2mm,after skip=2mm]
\textbf{\color{teal}To'g'ri javob:} \textbf{D) Audioaxborot}
\par\smallskip
\textbf{Izoh:} Vaziyat Audioaxborot tushunchasini ifodalaydi: tovush va nutq orqali ifodalanadigan axborot.
\end{tcolorbox}
\vspace{2mm}\hrule height 0.2pt \vspace{3mm}

% DB question_id: a1ef23df-afe8-5c23-8587-f0b5265110ea
\needspace{5\baselineskip}
\noindent\textbf{\color{ink}175.}
\par
Bir hisobot 2024-yil ma'lumotlariga asoslangan, 2026-yil qarori uchun ishlatilmoqda. Eng avvalo qaysi xossa tekshirilishi kerak?
\begin{enumerate}[label=\Alph*),leftmargin=9mm]
\item Grafik shakl
\item Dolzarblik
\item Maxfiylik
\item Statiklik
\end{enumerate}
\begin{tcolorbox}[enhanced,breakable,colback=teal!4,colframe=teal!50!black,boxrule=0.5pt,arc=1.1mm,left=3mm,right=3mm,top=1.7mm,bottom=1.7mm,before skip=2mm,after skip=2mm]
\textbf{\color{teal}To'g'ri javob:} \textbf{B) Dolzarblik}
\par\smallskip
\textbf{Izoh:} Qaror vaqti uchun ma'lumotning yangiligi, ya'ni dolzarbligi muhim.
\end{tcolorbox}
\vspace{2mm}\hrule height 0.2pt \vspace{3mm}

% DB question_id: a4f40975-0052-5597-892f-0f76ee0f7676
\needspace{5\baselineskip}
\noindent\textbf{\color{ink}176.}
\par
Jadvalda faqat o'rtacha ball bor, lekin o'quvchilar soni va taqsimot ko'rsatilmagan. Qaysi sifat yetishmaydi?
\begin{enumerate}[label=\Alph*),leftmargin=9mm]
\item Aniqlik
\item Ochiqlik
\item Statiklik
\item To'liqlik
\end{enumerate}
\begin{tcolorbox}[enhanced,breakable,colback=teal!4,colframe=teal!50!black,boxrule=0.5pt,arc=1.1mm,left=3mm,right=3mm,top=1.7mm,bottom=1.7mm,before skip=2mm,after skip=2mm]
\textbf{\color{teal}To'g'ri javob:} \textbf{D) To'liqlik}
\par\smallskip
\textbf{Izoh:} Qaror uchun kerakli barcha tafsilotlar berilmagan.
\end{tcolorbox}
\vspace{2mm}\hrule height 0.2pt \vspace{3mm}

% DB question_id: a623a9dc-8438-55c9-8de2-83e34aff8842
\needspace{5\baselineskip}
\noindent\textbf{\color{ink}177.}
\par
Texnik yo'riqnoma mutaxassis bo'lmagan foydalanuvchi uchun murakkab jargon bilan yozilgan. Qaysi xossa yetishmaydi?
\begin{enumerate}[label=\Alph*),leftmargin=9mm]
\item Qiymatlilik
\item Tushunarlilik
\item Aniqlik
\item Birlamchilik
\end{enumerate}
\begin{tcolorbox}[enhanced,breakable,colback=teal!4,colframe=teal!50!black,boxrule=0.5pt,arc=1.1mm,left=3mm,right=3mm,top=1.7mm,bottom=1.7mm,before skip=2mm,after skip=2mm]
\textbf{\color{teal}To'g'ri javob:} \textbf{B) Tushunarlilik}
\par\smallskip
\textbf{Izoh:} Auditoriya matnni anglay olmayapti, tushunarlilik past.
\end{tcolorbox}
\vspace{2mm}\hrule height 0.2pt \vspace{3mm}

% DB question_id: a97511e4-ed9b-5910-8603-f9ec9b6c71b8
\needspace{5\baselineskip}
\noindent\textbf{\color{ink}178.}
\par
Quyidagi ta'rif qaysi tushunchaga tegishli? Muayyan shaxsga tegishli yoki kirishi cheklangan axborot.
\begin{enumerate}[label=\Alph*),leftmargin=9mm]
\item Birlamchi manba
\item Ikkilamchi manba
\item Bevosita va bilvosita manba
\item Shaxsiy va maxfiy axborot
\end{enumerate}
\begin{tcolorbox}[enhanced,breakable,colback=teal!4,colframe=teal!50!black,boxrule=0.5pt,arc=1.1mm,left=3mm,right=3mm,top=1.7mm,bottom=1.7mm,before skip=2mm,after skip=2mm]
\textbf{\color{teal}To'g'ri javob:} \textbf{D) Shaxsiy va maxfiy axborot}
\par\smallskip
\textbf{Izoh:} Ta'rif Shaxsiy va maxfiy axborot tushunchasiga tegishli.
\end{tcolorbox}
\vspace{2mm}\hrule height 0.2pt \vspace{3mm}

% DB question_id: ade661e4-7110-5aeb-bb9a-cb932aefd38a
\needspace{5\baselineskip}
\noindent\textbf{\color{ink}179.}
\par
Quyidagi ta'rif qaysi tushunchaga tegishli? Qiymati ma'lum oraliqda uzluksiz o'zgarishi mumkin bo'lgan axborot.
\begin{enumerate}[label=\Alph*),leftmargin=9mm]
\item Uzluksiz axborot
\item Diskret axborot
\item Ochiq va ommaviy axborot
\item Shaxsiy va maxfiy axborot
\end{enumerate}
\begin{tcolorbox}[enhanced,breakable,colback=teal!4,colframe=teal!50!black,boxrule=0.5pt,arc=1.1mm,left=3mm,right=3mm,top=1.7mm,bottom=1.7mm,before skip=2mm,after skip=2mm]
\textbf{\color{teal}To'g'ri javob:} \textbf{A) Uzluksiz axborot}
\par\smallskip
\textbf{Izoh:} Ta'rif Uzluksiz axborot tushunchasiga tegishli.
\end{tcolorbox}
\vspace{2mm}\hrule height 0.2pt \vspace{3mm}

% DB question_id: be579e76-c2b2-5a0b-b799-3733b42ad973
\needspace{5\baselineskip}
\noindent\textbf{\color{ink}180.}
\par
Quyidagi ta'rif qaysi tushunchaga tegishli? Vaqt davomida o'zgarmaydigan yoki bitta holatni tasvirlaydigan axborot.
\begin{enumerate}[label=\Alph*),leftmargin=9mm]
\item Dinamik axborot
\item Uzluksiz axborot
\item Statik axborot
\item Diskret axborot
\end{enumerate}
\begin{tcolorbox}[enhanced,breakable,colback=teal!4,colframe=teal!50!black,boxrule=0.5pt,arc=1.1mm,left=3mm,right=3mm,top=1.7mm,bottom=1.7mm,before skip=2mm,after skip=2mm]
\textbf{\color{teal}To'g'ri javob:} \textbf{C) Statik axborot}
\par\smallskip
\textbf{Izoh:} Ta'rif Statik axborot tushunchasiga tegishli.
\end{tcolorbox}
\vspace{2mm}\hrule height 0.2pt \vspace{3mm}

% DB question_id: c06823f7-479e-467d-8e7e-de57a0fc38fe
\needspace{5\baselineskip}
\noindent\textbf{\color{ink}181.}
\par
O'quvchi elektron kundalikdagi bildirishnomani o'qidi. Bu vaziyatdagi asosiy tushunchani aniqlang.
\begin{enumerate}[label=\Alph*),leftmargin=9mm]
\item Qabul qiluvchi
\item Birlamchi manba
\item Ikkilamchi manba
\item Statik axborot
\end{enumerate}
\begin{tcolorbox}[enhanced,breakable,colback=teal!4,colframe=teal!50!black,boxrule=0.5pt,arc=1.1mm,left=3mm,right=3mm,top=1.7mm,bottom=1.7mm,before skip=2mm,after skip=2mm]
\textbf{\color{teal}To'g'ri javob:} \textbf{A) Qabul qiluvchi}
\par\smallskip
\textbf{Izoh:} Vaziyat Qabul qiluvchiga misol: xabarni qabul qilib, undan foydalanadigan obyekt yoki shaxs.
\end{tcolorbox}
\vspace{2mm}\hrule height 0.2pt \vspace{3mm}

% DB question_id: c4f8dc44-983a-5671-af23-a05adb5fdfe9
\needspace{5\baselineskip}
\noindent\textbf{\color{ink}182.}
\par
Quyidagi ta'rif qaysi tushunchaga tegishli? Miqdor va o'lchovlarni sonlar orqali ifodalovchi axborot.
\begin{enumerate}[label=\Alph*),leftmargin=9mm]
\item Grafik axborot
\item Sonli axborot
\item Audioaxborot
\item Videoaxborot
\end{enumerate}
\begin{tcolorbox}[enhanced,breakable,colback=teal!4,colframe=teal!50!black,boxrule=0.5pt,arc=1.1mm,left=3mm,right=3mm,top=1.7mm,bottom=1.7mm,before skip=2mm,after skip=2mm]
\textbf{\color{teal}To'g'ri javob:} \textbf{B) Sonli axborot}
\par\smallskip
\textbf{Izoh:} Ta'rif Sonli axborot tushunchasiga tegishli.
\end{tcolorbox}
\vspace{2mm}\hrule height 0.2pt \vspace{3mm}

% DB question_id: d754da1f-395c-4646-9ac5-0b8fed26f1c5
\needspace{5\baselineskip}
\noindent\textbf{\color{ink}183.}
\par
Slaydda mavzuni tushuntirish uchun matn, diagramma va ovozli izoh birgalikda qo'llangan. Bir nechta axborot shakli uyg'unlashgan bunday ko'rinish qanday ataladi?
\begin{enumerate}[label=\Alph*),leftmargin=9mm]
\item Grafikli axborot
\item Kombinatsiyalashgan axborot
\item Dinamik axborot
\item Analog axborot
\end{enumerate}
\begin{tcolorbox}[enhanced,breakable,colback=teal!4,colframe=teal!50!black,boxrule=0.5pt,arc=1.1mm,left=3mm,right=3mm,top=1.7mm,bottom=1.7mm,before skip=2mm,after skip=2mm]
\textbf{\color{teal}To'g'ri javob:} \textbf{B) Kombinatsiyalashgan axborot}
\par\smallskip
\textbf{Izoh:} Matn, grafika va audio bir mahsulotda birlashgani uchun bu kombinatsiyalashgan axborotdir. Grafik, dinamik yoki analog tushunchalar axborotning boshqa xususiyatlarini ta'riflaydi.
\end{tcolorbox}
\vspace{2mm}\hrule height 0.2pt \vspace{3mm}

% DB question_id: df76680b-9024-57bf-bbdb-44612db5c1c2
\needspace{5\baselineskip}
\noindent\textbf{\color{ink}184.}
\par
"Grafik axborot" atamasining eng aniq ta'rifini toping.
\begin{enumerate}[label=\Alph*),leftmargin=9mm]
\item tovush va nutq orqali ifodalanadigan axborot
\item vaqt bo'yicha almashuvchi tasvirlar ketma-ketligi ko'rinishidagi axborot
\item matn, grafika, audio yoki videoning birgalikdagi ko'rinishi
\item rasm, chizma, diagramma yoki sxema ko'rinishidagi axborot
\end{enumerate}
\begin{tcolorbox}[enhanced,breakable,colback=teal!4,colframe=teal!50!black,boxrule=0.5pt,arc=1.1mm,left=3mm,right=3mm,top=1.7mm,bottom=1.7mm,before skip=2mm,after skip=2mm]
\textbf{\color{teal}To'g'ri javob:} \textbf{D) rasm, chizma, diagramma yoki sxema ko'rinishidagi axborot}
\par\smallskip
\textbf{Izoh:} Grafik axborot - rasm, chizma, diagramma yoki sxema ko'rinishidagi axborot.
\end{tcolorbox}
\vspace{2mm}\hrule height 0.2pt \vspace{3mm}

% DB question_id: e5829b43-f4cf-4700-94a2-39b0f372776c
\needspace{5\baselineskip}
\noindent\textbf{\color{ink}185.}
\par
Muallif bir nechta tadqiqot natijalarini darslik bobida umumlashtirdi. Bu vaziyatdagi asosiy tushunchani aniqlang.
\begin{enumerate}[label=\Alph*),leftmargin=9mm]
\item Statik axborot
\item Dinamik axborot
\item Ikkilamchi manba
\item Axborotni yig'ish
\end{enumerate}
\begin{tcolorbox}[enhanced,breakable,colback=teal!4,colframe=teal!50!black,boxrule=0.5pt,arc=1.1mm,left=3mm,right=3mm,top=1.7mm,bottom=1.7mm,before skip=2mm,after skip=2mm]
\textbf{\color{teal}To'g'ri javob:} \textbf{C) Ikkilamchi manba}
\par\smallskip
\textbf{Izoh:} Vaziyat Ikkilamchi manbaga misol: birlamchi manbalarni tahlil qilib yoki qayta bayon qilib yaratilgan manba.
\end{tcolorbox}
\vspace{2mm}\hrule height 0.2pt \vspace{3mm}

% DB question_id: e6932086-2e33-563e-8386-65e5b7130876
\needspace{5\baselineskip}
\noindent\textbf{\color{ink}186.}
\par
Guvohning intervyusi bilan boshqa odamning shu intervyu haqidagi sharhi taqqoslandi. Ushbu vaziyatdagi asosiy tushunchani aniqlang.
\begin{enumerate}[label=\Alph*),leftmargin=9mm]
\item Bevosita va bilvosita manba
\item Matnli axborot
\item Sonli axborot
\item Grafik axborot
\end{enumerate}
\begin{tcolorbox}[enhanced,breakable,colback=teal!4,colframe=teal!50!black,boxrule=0.5pt,arc=1.1mm,left=3mm,right=3mm,top=1.7mm,bottom=1.7mm,before skip=2mm,after skip=2mm]
\textbf{\color{teal}To'g'ri javob:} \textbf{A) Bevosita va bilvosita manba}
\par\smallskip
\textbf{Izoh:} Vaziyat Bevosita va bilvosita manba tushunchasini ifodalaydi: bevosita manba hodisani to'g'ridan-to'g'ri, bilvosita manba esa vositachi orqali yetkazadi.
\end{tcolorbox}
\vspace{2mm}\hrule height 0.2pt \vspace{3mm}

% DB question_id: e935583b-1531-5d7e-989e-8a3ec28ca926
\needspace{5\baselineskip}
\noindent\textbf{\color{ink}187.}
\par
Shifokorning bemorni bevosita ko'rikdan o'tkazib yozgan qaydi qanday manba?
\begin{enumerate}[label=\Alph*),leftmargin=9mm]
\item Birlamchi va bevosita
\item Ikkilamchi va bilvosita
\item Ommaviy va tashqi
\item Statik va ikkilamchi
\end{enumerate}
\begin{tcolorbox}[enhanced,breakable,colback=teal!4,colframe=teal!50!black,boxrule=0.5pt,arc=1.1mm,left=3mm,right=3mm,top=1.7mm,bottom=1.7mm,before skip=2mm,after skip=2mm]
\textbf{\color{teal}To'g'ri javob:} \textbf{A) Birlamchi va bevosita}
\par\smallskip
\textbf{Izoh:} Shifokor hodisani bevosita kuzatgan va asl qayd yaratgan.
\end{tcolorbox}
\vspace{2mm}\hrule height 0.2pt \vspace{3mm}

% DB question_id: ee9b5ae9-3f70-587c-93e6-a9ec5d3b9c2c
\needspace{5\baselineskip}
\noindent\textbf{\color{ink}188.}
\par
Ikki sayt bir xil yangilik beradi, biroq birinchisi rasmiy hujjatga havola qiladi, ikkinchisi muallifsiz. Qaysi xossa birinchi manbada kuchliroq?
\begin{enumerate}[label=\Alph*),leftmargin=9mm]
\item Dinamiklik
\item Diskretlik
\item Ishonchlilik
\item Kombinatsiyalanganlik
\end{enumerate}
\begin{tcolorbox}[enhanced,breakable,colback=teal!4,colframe=teal!50!black,boxrule=0.5pt,arc=1.1mm,left=3mm,right=3mm,top=1.7mm,bottom=1.7mm,before skip=2mm,after skip=2mm]
\textbf{\color{teal}To'g'ri javob:} \textbf{C) Ishonchlilik}
\par\smallskip
\textbf{Izoh:} Rasmiy dalil va aniq muallif ishonchlilikni oshiradi.
\end{tcolorbox}
\vspace{2mm}\hrule height 0.2pt \vspace{3mm}

% DB question_id: ef4678af-ad1a-56c5-a709-6067475ed5af
\needspace{5\baselineskip}
\noindent\textbf{\color{ink}189.}
\par
"Matnli axborot" atamasining eng aniq ta'rifini toping.
\begin{enumerate}[label=\Alph*),leftmargin=9mm]
\item miqdor va o'lchovlarni sonlar orqali ifodalovchi axborot
\item harf, raqam va yozuv belgilaridan tuzilgan axborot shakli
\item rasm, chizma, diagramma yoki sxema ko'rinishidagi axborot
\item tovush va nutq orqali ifodalanadigan axborot
\end{enumerate}
\begin{tcolorbox}[enhanced,breakable,colback=teal!4,colframe=teal!50!black,boxrule=0.5pt,arc=1.1mm,left=3mm,right=3mm,top=1.7mm,bottom=1.7mm,before skip=2mm,after skip=2mm]
\textbf{\color{teal}To'g'ri javob:} \textbf{B) harf, raqam va yozuv belgilaridan tuzilgan axborot shakli}
\par\smallskip
\textbf{Izoh:} Matnli axborot - harf, raqam va yozuv belgilaridan tuzilgan axborot shakli.
\end{tcolorbox}
\vspace{2mm}\hrule height 0.2pt \vspace{3mm}

% DB question_id: f896e2d3-4b72-5fa1-b1d9-cd60d14e7a02
\needspace{5\baselineskip}
\noindent\textbf{\color{ink}190.}
\par
Tajriba jarayoni kamera orqali yozib olindi. Ushbu vaziyatdagi asosiy tushunchani aniqlang.
\begin{enumerate}[label=\Alph*),leftmargin=9mm]
\item Videoaxborot
\item Kombinatsiyalashgan axborot
\item Statik axborot
\item Dinamik axborot
\end{enumerate}
\begin{tcolorbox}[enhanced,breakable,colback=teal!4,colframe=teal!50!black,boxrule=0.5pt,arc=1.1mm,left=3mm,right=3mm,top=1.7mm,bottom=1.7mm,before skip=2mm,after skip=2mm]
\textbf{\color{teal}To'g'ri javob:} \textbf{A) Videoaxborot}
\par\smallskip
\textbf{Izoh:} Vaziyat Videoaxborot tushunchasini ifodalaydi: vaqt bo'yicha almashuvchi tasvirlar ketma-ketligi ko'rinishidagi axborot.
\end{tcolorbox}
\vspace{2mm}\hrule height 0.2pt \vspace{3mm}

% DB question_id: fa01eef6-bec8-5ad5-ad35-38d8df085f19
\needspace{5\baselineskip}
\noindent\textbf{\color{ink}191.}
\par
Muallif bir nechta tadqiqot asosida tahliliy maqola yozdi. Ushbu vaziyatdagi asosiy tushunchani aniqlang.
\begin{enumerate}[label=\Alph*),leftmargin=9mm]
\item Bevosita va bilvosita manba
\item Matnli axborot
\item Sonli axborot
\item Ikkilamchi manba
\end{enumerate}
\begin{tcolorbox}[enhanced,breakable,colback=teal!4,colframe=teal!50!black,boxrule=0.5pt,arc=1.1mm,left=3mm,right=3mm,top=1.7mm,bottom=1.7mm,before skip=2mm,after skip=2mm]
\textbf{\color{teal}To'g'ri javob:} \textbf{D) Ikkilamchi manba}
\par\smallskip
\textbf{Izoh:} Vaziyat Ikkilamchi manba tushunchasini ifodalaydi: birlamchi manbalarni tahlil qiluvchi, sharhlovchi yoki umumlashtiruvchi manba.
\end{tcolorbox}
\vspace{2mm}\hrule height 0.2pt \vspace{3mm}

% DB question_id: ff7dede3-17ab-477a-bc28-4d5f2e6a9735
\needspace{5\baselineskip}
\noindent\textbf{\color{ink}192.}
\par
Podkastda dars mazmuni ovozli tarzda tushuntirildi. Bu vaziyatdagi asosiy tushunchani aniqlang.
\begin{enumerate}[label=\Alph*),leftmargin=9mm]
\item Multimedia
\item Vizual axborot
\item Analog axborot
\item Audioaxborot
\end{enumerate}
\begin{tcolorbox}[enhanced,breakable,colback=teal!4,colframe=teal!50!black,boxrule=0.5pt,arc=1.1mm,left=3mm,right=3mm,top=1.7mm,bottom=1.7mm,before skip=2mm,after skip=2mm]
\textbf{\color{teal}To'g'ri javob:} \textbf{D) Audioaxborot}
\par\smallskip
\textbf{Izoh:} Vaziyat Audioaxborotga misol: tovush orqali ifodalangan va eshitish bilan qabul qilinadigan axborot.
\end{tcolorbox}
\vspace{2mm}\hrule height 0.2pt \vspace{3mm}

% DB question_id: 145b94e2-1055-4569-8059-b105389cb62d
\needspace{5\baselineskip}
\noindent\textbf{\color{ink}193.}
\par
Raqamli termometr haroratni 22, 23 yoki 24 kabi sonlarda ko'rsatdi. Vaziyatni tahlil qilib, unga eng aniq mos tushunchani toping.
\begin{enumerate}[label=\Alph*),leftmargin=9mm]
\item Aniqlik
\item Ishonchlilik
\item To'liqlik
\item Diskret axborot
\end{enumerate}
\begin{tcolorbox}[enhanced,breakable,colback=teal!4,colframe=teal!50!black,boxrule=0.5pt,arc=1.1mm,left=3mm,right=3mm,top=1.7mm,bottom=1.7mm,before skip=2mm,after skip=2mm]
\textbf{\color{teal}To'g'ri javob:} \textbf{D) Diskret axborot}
\par\smallskip
\textbf{Izoh:} Vaziyat Diskret axborotga misol: alohida qiymatlar yoki belgilar ketma-ketligida ifodalangan axborot.
\end{tcolorbox}
\vspace{2mm}\hrule height 0.2pt \vspace{3mm}

% DB question_id: 4738f1cc-d659-4a24-b1cb-08e555d82f86
\needspace{5\baselineskip}
\noindent\textbf{\color{ink}194.}
\par
PDF qo'llanmadagi jadval nashrdan keyin o'z-o'zidan yangilanmadi. Vaziyatni tahlil qilib, unga eng aniq mos tushunchani toping.
\begin{enumerate}[label=\Alph*),leftmargin=9mm]
\item Dinamik axborot
\item Axborotni yig'ish
\item Axborotni qayta ishlash
\item Statik axborot
\end{enumerate}
\begin{tcolorbox}[enhanced,breakable,colback=teal!4,colframe=teal!50!black,boxrule=0.5pt,arc=1.1mm,left=3mm,right=3mm,top=1.7mm,bottom=1.7mm,before skip=2mm,after skip=2mm]
\textbf{\color{teal}To'g'ri javob:} \textbf{D) Statik axborot}
\par\smallskip
\textbf{Izoh:} Vaziyat Statik axborotga misol: foydalanuvchi yoki tizim yangilamaguncha o'zgarmaydigan mazmun.
\end{tcolorbox}
\vspace{2mm}\hrule height 0.2pt \vspace{3mm}

% DB question_id: e6fb3d11-1e6b-4877-a93b-99950917134c
\needspace{5\baselineskip}
\noindent\textbf{\color{ink}195.}
\par
Simobli termometr ustuni harorat o'zgarishiga silliq ergashdi. Vaziyatni tahlil qilib, unga eng aniq mos tushunchani toping.
\begin{enumerate}[label=\Alph*),leftmargin=9mm]
\item Diskret axborot
\item Aniqlik
\item Analog axborot
\item Ishonchlilik
\end{enumerate}
\begin{tcolorbox}[enhanced,breakable,colback=teal!4,colframe=teal!50!black,boxrule=0.5pt,arc=1.1mm,left=3mm,right=3mm,top=1.7mm,bottom=1.7mm,before skip=2mm,after skip=2mm]
\textbf{\color{teal}To'g'ri javob:} \textbf{C) Analog axborot}
\par\smallskip
\textbf{Izoh:} Vaziyat Analog axborotga misol: vaqt yoki qiymat bo'yicha uzluksiz o'zgaradigan axborot.
\end{tcolorbox}
\vspace{2mm}\hrule height 0.2pt \vspace{3mm}

% DB question_id: ed26a1ca-f840-432c-ac56-75fa228e34a6
\needspace{5\baselineskip}
\noindent\textbf{\color{ink}196.}
\par
Valyuta kursi sahifada yangi ma'lumot kelishi bilan o'zgardi. Vaziyatni tahlil qilib, unga eng aniq mos tushunchani toping.
\begin{enumerate}[label=\Alph*),leftmargin=9mm]
\item Dinamik axborot
\item Axborotni yig'ish
\item Axborotni qayta ishlash
\item Axborotni saqlash
\end{enumerate}
\begin{tcolorbox}[enhanced,breakable,colback=teal!4,colframe=teal!50!black,boxrule=0.5pt,arc=1.1mm,left=3mm,right=3mm,top=1.7mm,bottom=1.7mm,before skip=2mm,after skip=2mm]
\textbf{\color{teal}To'g'ri javob:} \textbf{A) Dinamik axborot}
\par\smallskip
\textbf{Izoh:} Vaziyat Dinamik axborotga misol: vaqt, hodisa yoki ma'lumotlar bazasiga qarab avtomatik yangilanadigan mazmun.
\end{tcolorbox}
\vspace{2mm}\hrule height 0.2pt \vspace{3mm}
\section{Axborot jarayonlari, izlash va raqamli madaniyat}

% DB question_id: 0ebe8329-2829-49c7-958d-72deea613223
\needspace{5\baselineskip}
\noindent\textbf{\color{ink}197.}
\par
"Raqamli iz" atamasiga eng aniq ta'rifni toping.
\begin{enumerate}[label=\Alph*),leftmargin=9mm]
\item dalil yordamida tekshirish mumkin bo'lgan da'vo
\item foydalanuvchining onlayn harakatlari natijasida qoladigan ma'lumotlar majmui
\item shaxsiy baho, did yoki munosabatni ifodalovchi mulohaza
\item faktlar va qoidalar asosida chiqarilgan natija
\end{enumerate}
\begin{tcolorbox}[enhanced,breakable,colback=teal!4,colframe=teal!50!black,boxrule=0.5pt,arc=1.1mm,left=3mm,right=3mm,top=1.7mm,bottom=1.7mm,before skip=2mm,after skip=2mm]
\textbf{\color{teal}To'g'ri javob:} \textbf{B) foydalanuvchining onlayn harakatlari natijasida qoladigan ma'lumotlar majmui}
\par\smallskip
\textbf{Izoh:} Raqamli iz - foydalanuvchining onlayn harakatlari natijasida qoladigan ma'lumotlar majmui.
\end{tcolorbox}
\vspace{2mm}\hrule height 0.2pt \vspace{3mm}

% DB question_id: 108bd8f2-7baa-40c6-8b0d-9de7e84cc2f1
\needspace{5\baselineskip}
\noindent\textbf{\color{ink}198.}
\par
"Raqamli muhit" atamasiga eng aniq ta'rifni toping.
\begin{enumerate}[label=\Alph*),leftmargin=9mm]
\item axborot va bilim ishlab chiqarish hamda xizmatlarda asosiy resursga aylangan jamiyat
\item raqamli qurilmalar, dasturlar, tarmoqlar va xizmatlar o'zaro ishlaydigan makon
\item axborotni izlash, baholash, qonuniy va xavfsiz ishlatish ko'nikmalari majmui
\item raqamli vositalardan maqsadli, tanqidiy va samarali foydalanish qobiliyati
\end{enumerate}
\begin{tcolorbox}[enhanced,breakable,colback=teal!4,colframe=teal!50!black,boxrule=0.5pt,arc=1.1mm,left=3mm,right=3mm,top=1.7mm,bottom=1.7mm,before skip=2mm,after skip=2mm]
\textbf{\color{teal}To'g'ri javob:} \textbf{B) raqamli qurilmalar, dasturlar, tarmoqlar va xizmatlar o'zaro ishlaydigan makon}
\par\smallskip
\textbf{Izoh:} Raqamli muhit - raqamli qurilmalar, dasturlar, tarmoqlar va xizmatlar o'zaro ishlaydigan makon.
\end{tcolorbox}
\vspace{2mm}\hrule height 0.2pt \vspace{3mm}

% DB question_id: 13db70c4-34da-4a22-ad8c-368395dbed91
\needspace{5\baselineskip}
\noindent\textbf{\color{ink}199.}
\par
"Ishonchlilik" atamasiga eng aniq ta'rifni toping.
\begin{enumerate}[label=\Alph*),leftmargin=9mm]
\item qaror chiqarish uchun zarur bo'lgan muhim qismlar yetarli berilganlik xossasi
\item axborotning tekshirilgan dalil va ishonchli manbaga tayanganlik xossasi
\item harf, raqam va tinish belgilaridan tuzilgan yozma axborot
\item rasm, chizma, diagramma yoki sxema ko'rinishidagi axborot
\end{enumerate}
\begin{tcolorbox}[enhanced,breakable,colback=teal!4,colframe=teal!50!black,boxrule=0.5pt,arc=1.1mm,left=3mm,right=3mm,top=1.7mm,bottom=1.7mm,before skip=2mm,after skip=2mm]
\textbf{\color{teal}To'g'ri javob:} \textbf{B) axborotning tekshirilgan dalil va ishonchli manbaga tayanganlik xossasi}
\par\smallskip
\textbf{Izoh:} Ishonchlilik - axborotning tekshirilgan dalil va ishonchli manbaga tayanganlik xossasi.
\end{tcolorbox}
\vspace{2mm}\hrule height 0.2pt \vspace{3mm}

% DB question_id: 1b97d885-85f9-4e78-ae03-be7b9daa5ada
\needspace{5\baselineskip}
\noindent\textbf{\color{ink}200.}
\par
"CC NC" atamasiga eng aniq ta'rifni toping.
\begin{enumerate}[label=\Alph*),leftmargin=9mm]
\item asarning o'zgartirilgan hosila nusxalarini tarqatishni cheklaydigan shart
\item muallif, asar nomi, manba va litsenziyani imkon qadar aniq qayd etish
\item asardan tijorat maqsadida foydalanishni cheklaydigan shart
\item raqamli muloqotdagi hurmatli va mas'uliyatli xulq qoidalari
\end{enumerate}
\begin{tcolorbox}[enhanced,breakable,colback=teal!4,colframe=teal!50!black,boxrule=0.5pt,arc=1.1mm,left=3mm,right=3mm,top=1.7mm,bottom=1.7mm,before skip=2mm,after skip=2mm]
\textbf{\color{teal}To'g'ri javob:} \textbf{C) asardan tijorat maqsadida foydalanishni cheklaydigan shart}
\par\smallskip
\textbf{Izoh:} CC NC - asardan tijorat maqsadida foydalanishni cheklaydigan shart.
\end{tcolorbox}
\vspace{2mm}\hrule height 0.2pt \vspace{3mm}

% DB question_id: 2103b496-2522-4218-b8bb-cf0a59bb46ca
\needspace{5\baselineskip}
\noindent\textbf{\color{ink}201.}
\par
"Litsenziya" atamasiga eng aniq ta'rifni toping.
\begin{enumerate}[label=\Alph*),leftmargin=9mm]
\item mualliflik huquqi cheklovlari tugagan yoki qo'llanmaydigan asarlar holati
\item foydalanishga ruxsat berib, muallifni ko'rsatishni talab qiladigan Creative Commons sharti
\item asar yoki dasturdan qanday shartlarda foydalanish mumkinligini belgilovchi ruxsat
\item hosila asarni ayni yoki mos litsenziya bilan ulashishni talab qiladigan shart
\end{enumerate}
\begin{tcolorbox}[enhanced,breakable,colback=teal!4,colframe=teal!50!black,boxrule=0.5pt,arc=1.1mm,left=3mm,right=3mm,top=1.7mm,bottom=1.7mm,before skip=2mm,after skip=2mm]
\textbf{\color{teal}To'g'ri javob:} \textbf{C) asar yoki dasturdan qanday shartlarda foydalanish mumkinligini belgilovchi ruxsat}
\par\smallskip
\textbf{Izoh:} Litsenziya - asar yoki dasturdan qanday shartlarda foydalanish mumkinligini belgilovchi ruxsat.
\end{tcolorbox}
\vspace{2mm}\hrule height 0.2pt \vspace{3mm}

% DB question_id: 25e6cd73-3d20-4af3-b963-01709572e5b9
\needspace{5\baselineskip}
\noindent\textbf{\color{ink}202.}
\par
"Raqamli axborot" atamasiga eng aniq ta'rifni toping.
\begin{enumerate}[label=\Alph*),leftmargin=9mm]
\item raqamli kod, odatda bitlar ketma-ketligi orqali ifodalangan axborot
\item raqamli qurilmalar, dasturlar, tarmoqlar va xizmatlar o'zaro ishlaydigan makon
\item axborot va bilim ishlab chiqarish hamda xizmatlarda asosiy resursga aylangan jamiyat
\item axborotni izlash, baholash, qonuniy va xavfsiz ishlatish ko'nikmalari majmui
\end{enumerate}
\begin{tcolorbox}[enhanced,breakable,colback=teal!4,colframe=teal!50!black,boxrule=0.5pt,arc=1.1mm,left=3mm,right=3mm,top=1.7mm,bottom=1.7mm,before skip=2mm,after skip=2mm]
\textbf{\color{teal}To'g'ri javob:} \textbf{A) raqamli kod, odatda bitlar ketma-ketligi orqali ifodalangan axborot}
\par\smallskip
\textbf{Izoh:} Raqamli axborot - raqamli kod, odatda bitlar ketma-ketligi orqali ifodalangan axborot.
\end{tcolorbox}
\vspace{2mm}\hrule height 0.2pt \vspace{3mm}

% DB question_id: 2b8a3bba-fc1c-47fe-a016-294b228eec71
\needspace{5\baselineskip}
\noindent\textbf{\color{ink}203.}
\par
"Xulosa" atamasiga eng aniq ta'rifni toping.
\begin{enumerate}[label=\Alph*),leftmargin=9mm]
\item faktlar va qoidalar asosida chiqarilgan natija
\item axborotni yig'ish, saqlash, qayta ishlash, uzatish va taqdim etish vosita-usullari majmui
\item raqamli kod, odatda bitlar ketma-ketligi orqali ifodalangan axborot
\item raqamli qurilmalar, dasturlar, tarmoqlar va xizmatlar o'zaro ishlaydigan makon
\end{enumerate}
\begin{tcolorbox}[enhanced,breakable,colback=teal!4,colframe=teal!50!black,boxrule=0.5pt,arc=1.1mm,left=3mm,right=3mm,top=1.7mm,bottom=1.7mm,before skip=2mm,after skip=2mm]
\textbf{\color{teal}To'g'ri javob:} \textbf{A) faktlar va qoidalar asosida chiqarilgan natija}
\par\smallskip
\textbf{Izoh:} Xulosa - faktlar va qoidalar asosida chiqarilgan natija.
\end{tcolorbox}
\vspace{2mm}\hrule height 0.2pt \vspace{3mm}

% DB question_id: 30e6cfd1-c178-4efc-8783-8f8b612be7e0
\needspace{5\baselineskip}
\noindent\textbf{\color{ink}204.}
\par
"Axborotni yig'ish" atamasiga eng aniq ta'rifni toping.
\begin{enumerate}[label=\Alph*),leftmargin=9mm]
\item mavjud ma'lumotni tartiblash, hisoblash yoki tahlil qilish orqali yangi natija olish jarayoni
\item turli manbalardan kerakli ma'lumotlarni olish va jamlash jarayoni
\item axborotni keyinchalik foydalanish uchun tashuvchida qayd etish jarayoni
\item xabarni hosil qiluvchi yoki axborot kelib chiqadigan obyekt
\end{enumerate}
\begin{tcolorbox}[enhanced,breakable,colback=teal!4,colframe=teal!50!black,boxrule=0.5pt,arc=1.1mm,left=3mm,right=3mm,top=1.7mm,bottom=1.7mm,before skip=2mm,after skip=2mm]
\textbf{\color{teal}To'g'ri javob:} \textbf{B) turli manbalardan kerakli ma'lumotlarni olish va jamlash jarayoni}
\par\smallskip
\textbf{Izoh:} Axborotni yig'ish - turli manbalardan kerakli ma'lumotlarni olish va jamlash jarayoni.
\end{tcolorbox}
\vspace{2mm}\hrule height 0.2pt \vspace{3mm}

% DB question_id: 537c99b0-00e5-43ac-95a9-763e28541739
\needspace{5\baselineskip}
\noindent\textbf{\color{ink}205.}
\par
"Fakt" atamasiga eng aniq ta'rifni toping.
\begin{enumerate}[label=\Alph*),leftmargin=9mm]
\item shaxsiy baho, did yoki munosabatni ifodalovchi mulohaza
\item faktlar va qoidalar asosida chiqarilgan natija
\item dalil yordamida tekshirish mumkin bo'lgan da'vo
\item axborotni yig'ish, saqlash, qayta ishlash, uzatish va taqdim etish vosita-usullari majmui
\end{enumerate}
\begin{tcolorbox}[enhanced,breakable,colback=teal!4,colframe=teal!50!black,boxrule=0.5pt,arc=1.1mm,left=3mm,right=3mm,top=1.7mm,bottom=1.7mm,before skip=2mm,after skip=2mm]
\textbf{\color{teal}To'g'ri javob:} \textbf{C) dalil yordamida tekshirish mumkin bo'lgan da'vo}
\par\smallskip
\textbf{Izoh:} Fakt - dalil yordamida tekshirish mumkin bo'lgan da'vo.
\end{tcolorbox}
\vspace{2mm}\hrule height 0.2pt \vspace{3mm}

% DB question_id: 55f3ff56-6897-40dc-8152-78faf75db71b
\needspace{5\baselineskip}
\noindent\textbf{\color{ink}206.}
\par
"CC ND" atamasiga eng aniq ta'rifni toping.
\begin{enumerate}[label=\Alph*),leftmargin=9mm]
\item muallif, asar nomi, manba va litsenziyani imkon qadar aniq qayd etish
\item raqamli muloqotdagi hurmatli va mas'uliyatli xulq qoidalari
\item ijodiy asarning ifoda shakli va undan foydalanish huquqlarini himoya qiladigan huquqiy tizim
\item asarning o'zgartirilgan hosila nusxalarini tarqatishni cheklaydigan shart
\end{enumerate}
\begin{tcolorbox}[enhanced,breakable,colback=teal!4,colframe=teal!50!black,boxrule=0.5pt,arc=1.1mm,left=3mm,right=3mm,top=1.7mm,bottom=1.7mm,before skip=2mm,after skip=2mm]
\textbf{\color{teal}To'g'ri javob:} \textbf{D) asarning o'zgartirilgan hosila nusxalarini tarqatishni cheklaydigan shart}
\par\smallskip
\textbf{Izoh:} CC ND - asarning o'zgartirilgan hosila nusxalarini tarqatishni cheklaydigan shart.
\end{tcolorbox}
\vspace{2mm}\hrule height 0.2pt \vspace{3mm}

% DB question_id: 5b9aa286-af1c-455b-a101-c13c5a19e0e0
\needspace{5\baselineskip}
\noindent\textbf{\color{ink}207.}
\par
"Aniqlik" atamasiga eng aniq ta'rifni toping.
\begin{enumerate}[label=\Alph*),leftmargin=9mm]
\item axborotning obyekt holatini xatosiz va kerakli aniqlik darajasida aks ettirish xossasi
\item axborotning tekshirilgan dalil va ishonchli manbaga tayanganlik xossasi
\item qaror chiqarish uchun zarur bo'lgan muhim qismlar yetarli berilganlik xossasi
\item harf, raqam va tinish belgilaridan tuzilgan yozma axborot
\end{enumerate}
\begin{tcolorbox}[enhanced,breakable,colback=teal!4,colframe=teal!50!black,boxrule=0.5pt,arc=1.1mm,left=3mm,right=3mm,top=1.7mm,bottom=1.7mm,before skip=2mm,after skip=2mm]
\textbf{\color{teal}To'g'ri javob:} \textbf{A) axborotning obyekt holatini xatosiz va kerakli aniqlik darajasida aks ettirish xossasi}
\par\smallskip
\textbf{Izoh:} Aniqlik - axborotning obyekt holatini xatosiz va kerakli aniqlik darajasida aks ettirish xossasi.
\end{tcolorbox}
\vspace{2mm}\hrule height 0.2pt \vspace{3mm}

% DB question_id: 6485baba-1e89-5ec8-a979-f5c54abe8bb9
\needspace{5\baselineskip}
\noindent\textbf{\color{ink}208.}
\par
Sinf rahbari o'quvchilar haqidagi mavjud yozuvlarga ularning qiziqishlari haqidagi yangi ma'lumotlarni qo'shdi. Bu, avvalo, qaysi axborot jarayoniga misol?
\begin{enumerate}[label=\Alph*),leftmargin=9mm]
\item uzatish
\item to'plash
\item qayta ishlash
\item izlash
\end{enumerate}
\begin{tcolorbox}[enhanced,breakable,colback=teal!4,colframe=teal!50!black,boxrule=0.5pt,arc=1.1mm,left=3mm,right=3mm,top=1.7mm,bottom=1.7mm,before skip=2mm,after skip=2mm]
\textbf{\color{teal}To'g'ri javob:} \textbf{B) to'plash}
\par\smallskip
\textbf{Izoh:} O'quvchilardan yangi ma'lumot olib, yozuvlarga jamlash axborotni yig'ish/to'plash jarayonidir.
\end{tcolorbox}
\vspace{2mm}\hrule height 0.2pt \vspace{3mm}

% DB question_id: 66cae419-a986-42eb-86fe-203f707abe34
\needspace{5\baselineskip}
\noindent\textbf{\color{ink}209.}
\par
"Dezinformatsiya" atamasiga eng aniq ta'rifni toping.
\begin{enumerate}[label=\Alph*),leftmargin=9mm]
\item o'zaro bog'langan kompyuter tarmoqlarining global tizimi
\item Internet orqali olinadigan o'zaro havolalangan veb-sahifalar va resurslar xizmati
\item veb-resursni so'raydigan va mazmunini foydalanuvchiga ko'rsatadigan dastur
\item aldash maqsadida ataylab tarqatilgan noto'g'ri axborot
\end{enumerate}
\begin{tcolorbox}[enhanced,breakable,colback=teal!4,colframe=teal!50!black,boxrule=0.5pt,arc=1.1mm,left=3mm,right=3mm,top=1.7mm,bottom=1.7mm,before skip=2mm,after skip=2mm]
\textbf{\color{teal}To'g'ri javob:} \textbf{D) aldash maqsadida ataylab tarqatilgan noto'g'ri axborot}
\par\smallskip
\textbf{Izoh:} Dezinformatsiya - aldash maqsadida ataylab tarqatilgan noto'g'ri axborot.
\end{tcolorbox}
\vspace{2mm}\hrule height 0.2pt \vspace{3mm}

% DB question_id: 6cdca415-0471-481a-89e7-3819d528b85c
\needspace{5\baselineskip}
\noindent\textbf{\color{ink}210.}
\par
"Netiket" atamasiga eng aniq ta'rifni toping.
\begin{enumerate}[label=\Alph*),leftmargin=9mm]
\item ijodiy asarning ifoda shakli va undan foydalanish huquqlarini himoya qiladigan huquqiy tizim
\item raqamli muloqotdagi hurmatli va mas'uliyatli xulq qoidalari
\item boshqa shaxs g'oyasi yoki ifodasini manbasiz o'ziniki qilib ko'rsatish
\item asar yoki dasturdan qanday shartlarda foydalanish mumkinligini belgilovchi ruxsat
\end{enumerate}
\begin{tcolorbox}[enhanced,breakable,colback=teal!4,colframe=teal!50!black,boxrule=0.5pt,arc=1.1mm,left=3mm,right=3mm,top=1.7mm,bottom=1.7mm,before skip=2mm,after skip=2mm]
\textbf{\color{teal}To'g'ri javob:} \textbf{B) raqamli muloqotdagi hurmatli va mas'uliyatli xulq qoidalari}
\par\smallskip
\textbf{Izoh:} Netiket - raqamli muloqotdagi hurmatli va mas'uliyatli xulq qoidalari.
\end{tcolorbox}
\vspace{2mm}\hrule height 0.2pt \vspace{3mm}

% DB question_id: 6dbaac3d-41d0-4a24-8cdc-bc249756981d
\needspace{5\baselineskip}
\noindent\textbf{\color{ink}211.}
\par
"Axborot texnologiyasi" atamasiga eng aniq ta'rifni toping.
\begin{enumerate}[label=\Alph*),leftmargin=9mm]
\item raqamli kod, odatda bitlar ketma-ketligi orqali ifodalangan axborot
\item raqamli qurilmalar, dasturlar, tarmoqlar va xizmatlar o'zaro ishlaydigan makon
\item axborot va bilim ishlab chiqarish hamda xizmatlarda asosiy resursga aylangan jamiyat
\item axborotni yig'ish, saqlash, qayta ishlash, uzatish va taqdim etish vosita-usullari majmui
\end{enumerate}
\begin{tcolorbox}[enhanced,breakable,colback=teal!4,colframe=teal!50!black,boxrule=0.5pt,arc=1.1mm,left=3mm,right=3mm,top=1.7mm,bottom=1.7mm,before skip=2mm,after skip=2mm]
\textbf{\color{teal}To'g'ri javob:} \textbf{D) axborotni yig'ish, saqlash, qayta ishlash, uzatish va taqdim etish vosita-usullari majmui}
\par\smallskip
\textbf{Izoh:} Axborot texnologiyasi - axborotni yig'ish, saqlash, qayta ishlash, uzatish va taqdim etish vosita-usullari majmui.
\end{tcolorbox}
\vspace{2mm}\hrule height 0.2pt \vspace{3mm}

% DB question_id: 71b34d05-1195-4201-8b65-84af020dce41
\needspace{5\baselineskip}
\noindent\textbf{\color{ink}212.}
\par
"Mualliflik huquqi" atamasiga eng aniq ta'rifni toping.
\begin{enumerate}[label=\Alph*),leftmargin=9mm]
\item ijodiy asarning ifoda shakli va undan foydalanish huquqlarini himoya qiladigan huquqiy tizim
\item boshqa shaxs g'oyasi yoki ifodasini manbasiz o'ziniki qilib ko'rsatish
\item asar yoki dasturdan qanday shartlarda foydalanish mumkinligini belgilovchi ruxsat
\item mualliflik huquqi cheklovlari tugagan yoki qo'llanmaydigan asarlar holati
\end{enumerate}
\begin{tcolorbox}[enhanced,breakable,colback=teal!4,colframe=teal!50!black,boxrule=0.5pt,arc=1.1mm,left=3mm,right=3mm,top=1.7mm,bottom=1.7mm,before skip=2mm,after skip=2mm]
\textbf{\color{teal}To'g'ri javob:} \textbf{A) ijodiy asarning ifoda shakli va undan foydalanish huquqlarini himoya qiladigan huquqiy tizim}
\par\smallskip
\textbf{Izoh:} Mualliflik huquqi - ijodiy asarning ifoda shakli va undan foydalanish huquqlarini himoya qiladigan huquqiy tizim.
\end{tcolorbox}
\vspace{2mm}\hrule height 0.2pt \vspace{3mm}

% DB question_id: 7ba47682-d14f-4e7b-8d93-13cf4a3be118
\needspace{5\baselineskip}
\noindent\textbf{\color{ink}213.}
\par
"Plagiat" atamasiga eng aniq ta'rifni toping.
\begin{enumerate}[label=\Alph*),leftmargin=9mm]
\item asar yoki dasturdan qanday shartlarda foydalanish mumkinligini belgilovchi ruxsat
\item boshqa shaxs g'oyasi yoki ifodasini manbasiz o'ziniki qilib ko'rsatish
\item mualliflik huquqi cheklovlari tugagan yoki qo'llanmaydigan asarlar holati
\item foydalanishga ruxsat berib, muallifni ko'rsatishni talab qiladigan Creative Commons sharti
\end{enumerate}
\begin{tcolorbox}[enhanced,breakable,colback=teal!4,colframe=teal!50!black,boxrule=0.5pt,arc=1.1mm,left=3mm,right=3mm,top=1.7mm,bottom=1.7mm,before skip=2mm,after skip=2mm]
\textbf{\color{teal}To'g'ri javob:} \textbf{B) boshqa shaxs g'oyasi yoki ifodasini manbasiz o'ziniki qilib ko'rsatish}
\par\smallskip
\textbf{Izoh:} Plagiat - boshqa shaxs g'oyasi yoki ifodasini manbasiz o'ziniki qilib ko'rsatish.
\end{tcolorbox}
\vspace{2mm}\hrule height 0.2pt \vspace{3mm}

% DB question_id: 7db4b72d-4624-4a9f-9f8c-5ad09020a0bb
\needspace{5\baselineskip}
\noindent\textbf{\color{ink}214.}
\par
"Axborotlashgan jamiyat" atamasiga eng aniq ta'rifni toping.
\begin{enumerate}[label=\Alph*),leftmargin=9mm]
\item axborotni izlash, baholash, qonuniy va xavfsiz ishlatish ko'nikmalari majmui
\item raqamli vositalardan maqsadli, tanqidiy va samarali foydalanish qobiliyati
\item axborot va bilim ishlab chiqarish hamda xizmatlarda asosiy resursga aylangan jamiyat
\item foydalanuvchining onlayn harakatlari natijasida qoladigan ma'lumotlar majmui
\end{enumerate}
\begin{tcolorbox}[enhanced,breakable,colback=teal!4,colframe=teal!50!black,boxrule=0.5pt,arc=1.1mm,left=3mm,right=3mm,top=1.7mm,bottom=1.7mm,before skip=2mm,after skip=2mm]
\textbf{\color{teal}To'g'ri javob:} \textbf{C) axborot va bilim ishlab chiqarish hamda xizmatlarda asosiy resursga aylangan jamiyat}
\par\smallskip
\textbf{Izoh:} Axborotlashgan jamiyat - axborot va bilim ishlab chiqarish hamda xizmatlarda asosiy resursga aylangan jamiyat.
\end{tcolorbox}
\vspace{2mm}\hrule height 0.2pt \vspace{3mm}

% DB question_id: 8bbb2068-ca9e-4f99-851a-730cbd539673
\needspace{5\baselineskip}
\noindent\textbf{\color{ink}215.}
\par
"CC SA" atamasiga eng aniq ta'rifni toping.
\begin{enumerate}[label=\Alph*),leftmargin=9mm]
\item asardan tijorat maqsadida foydalanishni cheklaydigan shart
\item hosila asarni ayni yoki mos litsenziya bilan ulashishni talab qiladigan shart
\item asarning o'zgartirilgan hosila nusxalarini tarqatishni cheklaydigan shart
\item muallif, asar nomi, manba va litsenziyani imkon qadar aniq qayd etish
\end{enumerate}
\begin{tcolorbox}[enhanced,breakable,colback=teal!4,colframe=teal!50!black,boxrule=0.5pt,arc=1.1mm,left=3mm,right=3mm,top=1.7mm,bottom=1.7mm,before skip=2mm,after skip=2mm]
\textbf{\color{teal}To'g'ri javob:} \textbf{B) hosila asarni ayni yoki mos litsenziya bilan ulashishni talab qiladigan shart}
\par\smallskip
\textbf{Izoh:} CC SA - hosila asarni ayni yoki mos litsenziya bilan ulashishni talab qiladigan shart.
\end{tcolorbox}
\vspace{2mm}\hrule height 0.2pt \vspace{3mm}

% DB question_id: a7084ed2-240a-4d05-a708-22e679016adb
\needspace{5\baselineskip}
\noindent\textbf{\color{ink}216.}
\par
"To'liqlik" atamasiga eng aniq ta'rifni toping.
\begin{enumerate}[label=\Alph*),leftmargin=9mm]
\item harf, raqam va tinish belgilaridan tuzilgan yozma axborot
\item rasm, chizma, diagramma yoki sxema ko'rinishidagi axborot
\item qaror chiqarish uchun zarur bo'lgan muhim qismlar yetarli berilganlik xossasi
\item tovush orqali ifodalangan va eshitish bilan qabul qilinadigan axborot
\end{enumerate}
\begin{tcolorbox}[enhanced,breakable,colback=teal!4,colframe=teal!50!black,boxrule=0.5pt,arc=1.1mm,left=3mm,right=3mm,top=1.7mm,bottom=1.7mm,before skip=2mm,after skip=2mm]
\textbf{\color{teal}To'g'ri javob:} \textbf{C) qaror chiqarish uchun zarur bo'lgan muhim qismlar yetarli berilganlik xossasi}
\par\smallskip
\textbf{Izoh:} To'liqlik - qaror chiqarish uchun zarur bo'lgan muhim qismlar yetarli berilganlik xossasi.
\end{tcolorbox}
\vspace{2mm}\hrule height 0.2pt \vspace{3mm}

% DB question_id: b8445a6e-993c-49f7-9cb3-7fd5d94ade25
\needspace{5\baselineskip}
\noindent\textbf{\color{ink}217.}
\par
"HTTPS va ishonchlilik farqi" atamasiga eng aniq ta'rifni toping.
\begin{enumerate}[label=\Alph*),leftmargin=9mm]
\item ma'lumot qayd etilgan qiymat, axborot esa qabul qiluvchi uchun mazmunli ma'lumot ekanini ajratish
\item kodlash taqdim etish formatini, shifrlash esa maxfiylikni ko'zlashini ajratish
\item bit ikkilik birlik, bayt esa 8 bitdan iborat ekanini ajratish
\item HTTPS transportni shifrlashi, sahifa mazmunining haqiqatligini kafolatlamasligini ajratish
\end{enumerate}
\begin{tcolorbox}[enhanced,breakable,colback=teal!4,colframe=teal!50!black,boxrule=0.5pt,arc=1.1mm,left=3mm,right=3mm,top=1.7mm,bottom=1.7mm,before skip=2mm,after skip=2mm]
\textbf{\color{teal}To'g'ri javob:} \textbf{D) HTTPS transportni shifrlashi, sahifa mazmunining haqiqatligini kafolatlamasligini ajratish}
\par\smallskip
\textbf{Izoh:} HTTPS va ishonchlilik farqi - HTTPS transportni shifrlashi, sahifa mazmunining haqiqatligini kafolatlamasligini ajratish.
\end{tcolorbox}
\vspace{2mm}\hrule height 0.2pt \vspace{3mm}

% DB question_id: b9637ed4-b85b-4f41-b9ba-e7aa038116ad
\needspace{5\baselineskip}
\noindent\textbf{\color{ink}218.}
\par
"CC BY" atamasiga eng aniq ta'rifni toping.
\begin{enumerate}[label=\Alph*),leftmargin=9mm]
\item foydalanishga ruxsat berib, muallifni ko'rsatishni talab qiladigan Creative Commons sharti
\item hosila asarni ayni yoki mos litsenziya bilan ulashishni talab qiladigan shart
\item asardan tijorat maqsadida foydalanishni cheklaydigan shart
\item asarning o'zgartirilgan hosila nusxalarini tarqatishni cheklaydigan shart
\end{enumerate}
\begin{tcolorbox}[enhanced,breakable,colback=teal!4,colframe=teal!50!black,boxrule=0.5pt,arc=1.1mm,left=3mm,right=3mm,top=1.7mm,bottom=1.7mm,before skip=2mm,after skip=2mm]
\textbf{\color{teal}To'g'ri javob:} \textbf{A) foydalanishga ruxsat berib, muallifni ko'rsatishni talab qiladigan Creative Commons sharti}
\par\smallskip
\textbf{Izoh:} CC BY - foydalanishga ruxsat berib, muallifni ko'rsatishni talab qiladigan Creative Commons sharti.
\end{tcolorbox}
\vspace{2mm}\hrule height 0.2pt \vspace{3mm}

% DB question_id: c32f4b78-aa14-4917-93ee-e6f51b94c7ff
\needspace{5\baselineskip}
\noindent\textbf{\color{ink}219.}
\par
"Axborotni saqlash" atamasiga eng aniq ta'rifni toping.
\begin{enumerate}[label=\Alph*),leftmargin=9mm]
\item xabarni hosil qiluvchi yoki axborot kelib chiqadigan obyekt
\item xabar manbadan qabul qiluvchiga o'tadigan muhit yoki vosita
\item xabarni qabul qilib, undan foydalanadigan obyekt yoki shaxs
\item axborotni keyinchalik foydalanish uchun tashuvchida qayd etish jarayoni
\end{enumerate}
\begin{tcolorbox}[enhanced,breakable,colback=teal!4,colframe=teal!50!black,boxrule=0.5pt,arc=1.1mm,left=3mm,right=3mm,top=1.7mm,bottom=1.7mm,before skip=2mm,after skip=2mm]
\textbf{\color{teal}To'g'ri javob:} \textbf{D) axborotni keyinchalik foydalanish uchun tashuvchida qayd etish jarayoni}
\par\smallskip
\textbf{Izoh:} Axborotni saqlash - axborotni keyinchalik foydalanish uchun tashuvchida qayd etish jarayoni.
\end{tcolorbox}
\vspace{2mm}\hrule height 0.2pt \vspace{3mm}

% DB question_id: c5d2cd17-ef61-447c-827b-caf09d9e4a51
\needspace{5\baselineskip}
\noindent\textbf{\color{ink}220.}
\par
"Axborot madaniyati" atamasiga eng aniq ta'rifni toping.
\begin{enumerate}[label=\Alph*),leftmargin=9mm]
\item raqamli vositalardan maqsadli, tanqidiy va samarali foydalanish qobiliyati
\item foydalanuvchining onlayn harakatlari natijasida qoladigan ma'lumotlar majmui
\item dalil yordamida tekshirish mumkin bo'lgan da'vo
\item axborotni izlash, baholash, qonuniy va xavfsiz ishlatish ko'nikmalari majmui
\end{enumerate}
\begin{tcolorbox}[enhanced,breakable,colback=teal!4,colframe=teal!50!black,boxrule=0.5pt,arc=1.1mm,left=3mm,right=3mm,top=1.7mm,bottom=1.7mm,before skip=2mm,after skip=2mm]
\textbf{\color{teal}To'g'ri javob:} \textbf{D) axborotni izlash, baholash, qonuniy va xavfsiz ishlatish ko'nikmalari majmui}
\par\smallskip
\textbf{Izoh:} Axborot madaniyati - axborotni izlash, baholash, qonuniy va xavfsiz ishlatish ko'nikmalari majmui.
\end{tcolorbox}
\vspace{2mm}\hrule height 0.2pt \vspace{3mm}

% DB question_id: d8f96862-cc38-49d9-a34a-a3530eb04134
\needspace{5\baselineskip}
\noindent\textbf{\color{ink}221.}
\par
"Manba ko'rsatish" atamasiga eng aniq ta'rifni toping.
\begin{enumerate}[label=\Alph*),leftmargin=9mm]
\item muallif, asar nomi, manba va litsenziyani imkon qadar aniq qayd etish
\item raqamli muloqotdagi hurmatli va mas'uliyatli xulq qoidalari
\item ijodiy asarning ifoda shakli va undan foydalanish huquqlarini himoya qiladigan huquqiy tizim
\item boshqa shaxs g'oyasi yoki ifodasini manbasiz o'ziniki qilib ko'rsatish
\end{enumerate}
\begin{tcolorbox}[enhanced,breakable,colback=teal!4,colframe=teal!50!black,boxrule=0.5pt,arc=1.1mm,left=3mm,right=3mm,top=1.7mm,bottom=1.7mm,before skip=2mm,after skip=2mm]
\textbf{\color{teal}To'g'ri javob:} \textbf{A) muallif, asar nomi, manba va litsenziyani imkon qadar aniq qayd etish}
\par\smallskip
\textbf{Izoh:} Manba ko'rsatish - muallif, asar nomi, manba va litsenziyani imkon qadar aniq qayd etish.
\end{tcolorbox}
\vspace{2mm}\hrule height 0.2pt \vspace{3mm}

% DB question_id: da1f47ec-e0b2-4658-b743-9926986ee090
\needspace{5\baselineskip}
\noindent\textbf{\color{ink}222.}
\par
"Raqamli savodxonlik" atamasiga eng aniq ta'rifni toping.
\begin{enumerate}[label=\Alph*),leftmargin=9mm]
\item raqamli vositalardan maqsadli, tanqidiy va samarali foydalanish qobiliyati
\item foydalanuvchining onlayn harakatlari natijasida qoladigan ma'lumotlar majmui
\item dalil yordamida tekshirish mumkin bo'lgan da'vo
\item shaxsiy baho, did yoki munosabatni ifodalovchi mulohaza
\end{enumerate}
\begin{tcolorbox}[enhanced,breakable,colback=teal!4,colframe=teal!50!black,boxrule=0.5pt,arc=1.1mm,left=3mm,right=3mm,top=1.7mm,bottom=1.7mm,before skip=2mm,after skip=2mm]
\textbf{\color{teal}To'g'ri javob:} \textbf{A) raqamli vositalardan maqsadli, tanqidiy va samarali foydalanish qobiliyati}
\par\smallskip
\textbf{Izoh:} Raqamli savodxonlik - raqamli vositalardan maqsadli, tanqidiy va samarali foydalanish qobiliyati.
\end{tcolorbox}
\vspace{2mm}\hrule height 0.2pt \vspace{3mm}

% DB question_id: dd019856-35fe-589c-9fa7-d251179e4cd7
\needspace{5\baselineskip}
\noindent\textbf{\color{ink}223.}
\par
Qaysi holat raqamli kontentdan qonuniy va akademik halol foydalanishga misol bo'ladi?
\begin{enumerate}[label=\Alph*),leftmargin=9mm]
\item Muallif va manbani ko'rsatib, tahlil uchun zarur kichik parchani iqtibos qilish
\item Pulli dasturni noqonuniy nusxalab tarqatish
\item Topilgan rasmni ruxsatsiz reklamada ishlatish
\item Filmni yuklab olib do'stlarga tarqatish
\end{enumerate}
\begin{tcolorbox}[enhanced,breakable,colback=teal!4,colframe=teal!50!black,boxrule=0.5pt,arc=1.1mm,left=3mm,right=3mm,top=1.7mm,bottom=1.7mm,before skip=2mm,after skip=2mm]
\textbf{\color{teal}To'g'ri javob:} \textbf{A) Muallif va manbani ko'rsatib, tahlil uchun zarur kichik parchani iqtibos qilish}
\par\smallskip
\textbf{Izoh:} Manbasi va muallifi ko'rsatilgan, maqsadga zarur hajmdagi qisqa iqtibos qonuniy foydalanish shartlariga mos; qolganlari huquq egasining ruxsatisiz tarqatish yoki tijoriy foydalanishdir.
\end{tcolorbox}
\vspace{2mm}\hrule height 0.2pt \vspace{3mm}

% DB question_id: e08d3665-5300-4321-a505-46357a8c3842
\needspace{5\baselineskip}
\noindent\textbf{\color{ink}224.}
\par
"Axborotni qayta ishlash" atamasiga eng aniq ta'rifni toping.
\begin{enumerate}[label=\Alph*),leftmargin=9mm]
\item axborotni keyinchalik foydalanish uchun tashuvchida qayd etish jarayoni
\item xabarni hosil qiluvchi yoki axborot kelib chiqadigan obyekt
\item mavjud ma'lumotni tartiblash, hisoblash yoki tahlil qilish orqali yangi natija olish jarayoni
\item xabar manbadan qabul qiluvchiga o'tadigan muhit yoki vosita
\end{enumerate}
\begin{tcolorbox}[enhanced,breakable,colback=teal!4,colframe=teal!50!black,boxrule=0.5pt,arc=1.1mm,left=3mm,right=3mm,top=1.7mm,bottom=1.7mm,before skip=2mm,after skip=2mm]
\textbf{\color{teal}To'g'ri javob:} \textbf{C) mavjud ma'lumotni tartiblash, hisoblash yoki tahlil qilish orqali yangi natija olish jarayoni}
\par\smallskip
\textbf{Izoh:} Axborotni qayta ishlash - mavjud ma'lumotni tartiblash, hisoblash yoki tahlil qilish orqali yangi natija olish jarayoni.
\end{tcolorbox}
\vspace{2mm}\hrule height 0.2pt \vspace{3mm}

% DB question_id: e81682b2-aa91-4334-879e-0e04c3ff3ab6
\needspace{5\baselineskip}
\noindent\textbf{\color{ink}225.}
\par
"Manba ishonchliligi" atamasiga eng aniq ta'rifni toping.
\begin{enumerate}[label=\Alph*),leftmargin=9mm]
\item aldash maqsadida ataylab tarqatilgan noto'g'ri axborot
\item o'zaro bog'langan kompyuter tarmoqlarining global tizimi
\item muallif, sana, dalil, maqsad va mustaqil tasdiq orqali baholanadigan sifat
\item Internet orqali olinadigan o'zaro havolalangan veb-sahifalar va resurslar xizmati
\end{enumerate}
\begin{tcolorbox}[enhanced,breakable,colback=teal!4,colframe=teal!50!black,boxrule=0.5pt,arc=1.1mm,left=3mm,right=3mm,top=1.7mm,bottom=1.7mm,before skip=2mm,after skip=2mm]
\textbf{\color{teal}To'g'ri javob:} \textbf{C) muallif, sana, dalil, maqsad va mustaqil tasdiq orqali baholanadigan sifat}
\par\smallskip
\textbf{Izoh:} Manba ishonchliligi - muallif, sana, dalil, maqsad va mustaqil tasdiq orqali baholanadigan sifat.
\end{tcolorbox}
\vspace{2mm}\hrule height 0.2pt \vspace{3mm}

% DB question_id: f37d5f73-1312-4a8c-be80-f7c5b3aa2611
\needspace{5\baselineskip}
\noindent\textbf{\color{ink}226.}
\par
"Ommaviy mulk" atamasiga eng aniq ta'rifni toping.
\begin{enumerate}[label=\Alph*),leftmargin=9mm]
\item foydalanishga ruxsat berib, muallifni ko'rsatishni talab qiladigan Creative Commons sharti
\item hosila asarni ayni yoki mos litsenziya bilan ulashishni talab qiladigan shart
\item asardan tijorat maqsadida foydalanishni cheklaydigan shart
\item mualliflik huquqi cheklovlari tugagan yoki qo'llanmaydigan asarlar holati
\end{enumerate}
\begin{tcolorbox}[enhanced,breakable,colback=teal!4,colframe=teal!50!black,boxrule=0.5pt,arc=1.1mm,left=3mm,right=3mm,top=1.7mm,bottom=1.7mm,before skip=2mm,after skip=2mm]
\textbf{\color{teal}To'g'ri javob:} \textbf{D) mualliflik huquqi cheklovlari tugagan yoki qo'llanmaydigan asarlar holati}
\par\smallskip
\textbf{Izoh:} Ommaviy mulk - mualliflik huquqi cheklovlari tugagan yoki qo'llanmaydigan asarlar holati.
\end{tcolorbox}
\vspace{2mm}\hrule height 0.2pt \vspace{3mm}

% DB question_id: fb0ecdb9-6ee2-44aa-aa6c-219c18d1f733
\needspace{5\baselineskip}
\noindent\textbf{\color{ink}227.}
\par
"Fikr" atamasiga eng aniq ta'rifni toping.
\begin{enumerate}[label=\Alph*),leftmargin=9mm]
\item faktlar va qoidalar asosida chiqarilgan natija
\item axborotni yig'ish, saqlash, qayta ishlash, uzatish va taqdim etish vosita-usullari majmui
\item raqamli kod, odatda bitlar ketma-ketligi orqali ifodalangan axborot
\item shaxsiy baho, did yoki munosabatni ifodalovchi mulohaza
\end{enumerate}
\begin{tcolorbox}[enhanced,breakable,colback=teal!4,colframe=teal!50!black,boxrule=0.5pt,arc=1.1mm,left=3mm,right=3mm,top=1.7mm,bottom=1.7mm,before skip=2mm,after skip=2mm]
\textbf{\color{teal}To'g'ri javob:} \textbf{D) shaxsiy baho, did yoki munosabatni ifodalovchi mulohaza}
\par\smallskip
\textbf{Izoh:} Fikr - shaxsiy baho, did yoki munosabatni ifodalovchi mulohaza.
\end{tcolorbox}
\vspace{2mm}\hrule height 0.2pt \vspace{3mm}

% DB question_id: 0fd51da1-3bf5-4a89-994a-ae757db7ea76
\needspace{5\baselineskip}
\noindent\textbf{\color{ink}228.}
\par
Quyidagi ta'rif qaysi atamaga tegishli? Ijodiy asarning ifoda shakli va undan foydalanish huquqlarini himoya qiladigan huquqiy tizim.
\begin{enumerate}[label=\Alph*),leftmargin=9mm]
\item Plagiat
\item Litsenziya
\item Mualliflik huquqi
\item Ommaviy mulk
\end{enumerate}
\begin{tcolorbox}[enhanced,breakable,colback=teal!4,colframe=teal!50!black,boxrule=0.5pt,arc=1.1mm,left=3mm,right=3mm,top=1.7mm,bottom=1.7mm,before skip=2mm,after skip=2mm]
\textbf{\color{teal}To'g'ri javob:} \textbf{C) Mualliflik huquqi}
\par\smallskip
\textbf{Izoh:} Ta'rif Mualliflik huquqi tushunchasiga tegishli.
\end{tcolorbox}
\vspace{2mm}\hrule height 0.2pt \vspace{3mm}

% DB question_id: 1cf3c927-3512-4080-9e7e-444706b5f860
\needspace{5\baselineskip}
\noindent\textbf{\color{ink}229.}
\par
Quyidagi ta'rif qaysi atamaga tegishli? Mavjud ma'lumotni tartiblash, hisoblash yoki tahlil qilish orqali yangi natija olish jarayoni.
\begin{enumerate}[label=\Alph*),leftmargin=9mm]
\item Axborotni qayta ishlash
\item Axborotni saqlash
\item Axborot manbai
\item Aloqa kanali
\end{enumerate}
\begin{tcolorbox}[enhanced,breakable,colback=teal!4,colframe=teal!50!black,boxrule=0.5pt,arc=1.1mm,left=3mm,right=3mm,top=1.7mm,bottom=1.7mm,before skip=2mm,after skip=2mm]
\textbf{\color{teal}To'g'ri javob:} \textbf{A) Axborotni qayta ishlash}
\par\smallskip
\textbf{Izoh:} Ta'rif Axborotni qayta ishlash tushunchasiga tegishli.
\end{tcolorbox}
\vspace{2mm}\hrule height 0.2pt \vspace{3mm}

% DB question_id: 2a0b9329-a9ba-4135-b76f-cd02426d8d36
\needspace{5\baselineskip}
\noindent\textbf{\color{ink}230.}
\par
Quyidagi ta'rif qaysi atamaga tegishli? Axborotni izlash, baholash, qonuniy va xavfsiz ishlatish ko'nikmalari majmui.
\begin{enumerate}[label=\Alph*),leftmargin=9mm]
\item Raqamli savodxonlik
\item Axborot madaniyati
\item Raqamli iz
\item Fakt
\end{enumerate}
\begin{tcolorbox}[enhanced,breakable,colback=teal!4,colframe=teal!50!black,boxrule=0.5pt,arc=1.1mm,left=3mm,right=3mm,top=1.7mm,bottom=1.7mm,before skip=2mm,after skip=2mm]
\textbf{\color{teal}To'g'ri javob:} \textbf{B) Axborot madaniyati}
\par\smallskip
\textbf{Izoh:} Ta'rif Axborot madaniyati tushunchasiga tegishli.
\end{tcolorbox}
\vspace{2mm}\hrule height 0.2pt \vspace{3mm}

% DB question_id: 2f219ecd-fe09-426d-9f75-e157dc80d6bb
\needspace{5\baselineskip}
\noindent\textbf{\color{ink}231.}
\par
Quyidagi ta'rif qaysi atamaga tegishli? Dalil yordamida tekshirish mumkin bo'lgan da'vo.
\begin{enumerate}[label=\Alph*),leftmargin=9mm]
\item Fakt
\item Fikr
\item Xulosa
\item Axborot texnologiyasi
\end{enumerate}
\begin{tcolorbox}[enhanced,breakable,colback=teal!4,colframe=teal!50!black,boxrule=0.5pt,arc=1.1mm,left=3mm,right=3mm,top=1.7mm,bottom=1.7mm,before skip=2mm,after skip=2mm]
\textbf{\color{teal}To'g'ri javob:} \textbf{A) Fakt}
\par\smallskip
\textbf{Izoh:} Ta'rif Fakt tushunchasiga tegishli.
\end{tcolorbox}
\vspace{2mm}\hrule height 0.2pt \vspace{3mm}

% DB question_id: 332746b7-5ef0-4e0b-b6fb-4b52c73996dc
\needspace{5\baselineskip}
\noindent\textbf{\color{ink}232.}
\par
Quyidagi ta'rif qaysi atamaga tegishli? Muallif, sana, dalil, maqsad va mustaqil tasdiq orqali baholanadigan sifat.
\begin{enumerate}[label=\Alph*),leftmargin=9mm]
\item Manba ishonchliligi
\item Dezinformatsiya
\item Internet
\item WWW
\end{enumerate}
\begin{tcolorbox}[enhanced,breakable,colback=teal!4,colframe=teal!50!black,boxrule=0.5pt,arc=1.1mm,left=3mm,right=3mm,top=1.7mm,bottom=1.7mm,before skip=2mm,after skip=2mm]
\textbf{\color{teal}To'g'ri javob:} \textbf{A) Manba ishonchliligi}
\par\smallskip
\textbf{Izoh:} Ta'rif Manba ishonchliligi tushunchasiga tegishli.
\end{tcolorbox}
\vspace{2mm}\hrule height 0.2pt \vspace{3mm}

% DB question_id: 3405997e-5676-438d-a78c-883cc1021977
\needspace{5\baselineskip}
\noindent\textbf{\color{ink}233.}
\par
Quyidagi ta'rif qaysi atamaga tegishli? Asardan tijorat maqsadida foydalanishni cheklaydigan shart.
\begin{enumerate}[label=\Alph*),leftmargin=9mm]
\item CC NC
\item CC ND
\item Manba ko'rsatish
\item Netiket
\end{enumerate}
\begin{tcolorbox}[enhanced,breakable,colback=teal!4,colframe=teal!50!black,boxrule=0.5pt,arc=1.1mm,left=3mm,right=3mm,top=1.7mm,bottom=1.7mm,before skip=2mm,after skip=2mm]
\textbf{\color{teal}To'g'ri javob:} \textbf{A) CC NC}
\par\smallskip
\textbf{Izoh:} Ta'rif CC NC tushunchasiga tegishli.
\end{tcolorbox}
\vspace{2mm}\hrule height 0.2pt \vspace{3mm}

% DB question_id: 3705f22a-1540-46a5-84dc-34a771064b21
\needspace{5\baselineskip}
\noindent\textbf{\color{ink}234.}
\par
Quyidagi ta'rif qaysi atamaga tegishli? Raqamli qurilmalar, dasturlar, tarmoqlar va xizmatlar o'zaro ishlaydigan makon.
\begin{enumerate}[label=\Alph*),leftmargin=9mm]
\item Axborotlashgan jamiyat
\item Axborot madaniyati
\item Raqamli savodxonlik
\item Raqamli muhit
\end{enumerate}
\begin{tcolorbox}[enhanced,breakable,colback=teal!4,colframe=teal!50!black,boxrule=0.5pt,arc=1.1mm,left=3mm,right=3mm,top=1.7mm,bottom=1.7mm,before skip=2mm,after skip=2mm]
\textbf{\color{teal}To'g'ri javob:} \textbf{D) Raqamli muhit}
\par\smallskip
\textbf{Izoh:} Ta'rif Raqamli muhit tushunchasiga tegishli.
\end{tcolorbox}
\vspace{2mm}\hrule height 0.2pt \vspace{3mm}

% DB question_id: 39f8c4bd-1479-4345-b802-d8547f5e6e06
\needspace{5\baselineskip}
\noindent\textbf{\color{ink}235.}
\par
Quyidagi ta'rif qaysi atamaga tegishli? Boshqa shaxs g'oyasi yoki ifodasini manbasiz o'ziniki qilib ko'rsatish.
\begin{enumerate}[label=\Alph*),leftmargin=9mm]
\item Litsenziya
\item Ommaviy mulk
\item CC BY
\item Plagiat
\end{enumerate}
\begin{tcolorbox}[enhanced,breakable,colback=teal!4,colframe=teal!50!black,boxrule=0.5pt,arc=1.1mm,left=3mm,right=3mm,top=1.7mm,bottom=1.7mm,before skip=2mm,after skip=2mm]
\textbf{\color{teal}To'g'ri javob:} \textbf{D) Plagiat}
\par\smallskip
\textbf{Izoh:} Ta'rif Plagiat tushunchasiga tegishli.
\end{tcolorbox}
\vspace{2mm}\hrule height 0.2pt \vspace{3mm}

% DB question_id: 566b4d53-8ba5-454a-a5e9-8df47d0e4a09
\needspace{5\baselineskip}
\noindent\textbf{\color{ink}236.}
\par
Quyidagi ta'rif qaysi atamaga tegishli? Asarning o'zgartirilgan hosila nusxalarini tarqatishni cheklaydigan shart.
\begin{enumerate}[label=\Alph*),leftmargin=9mm]
\item Manba ko'rsatish
\item CC ND
\item Netiket
\item Mualliflik huquqi
\end{enumerate}
\begin{tcolorbox}[enhanced,breakable,colback=teal!4,colframe=teal!50!black,boxrule=0.5pt,arc=1.1mm,left=3mm,right=3mm,top=1.7mm,bottom=1.7mm,before skip=2mm,after skip=2mm]
\textbf{\color{teal}To'g'ri javob:} \textbf{B) CC ND}
\par\smallskip
\textbf{Izoh:} Ta'rif CC ND tushunchasiga tegishli.
\end{tcolorbox}
\vspace{2mm}\hrule height 0.2pt \vspace{3mm}

% DB question_id: 567e018c-4d3a-4b98-b32b-0a9d5c7e1784
\needspace{5\baselineskip}
\noindent\textbf{\color{ink}237.}
\par
Quyidagi ta'rif qaysi atamaga tegishli? Qaror chiqarish uchun zarur bo'lgan muhim qismlar yetarli berilganlik xossasi.
\begin{enumerate}[label=\Alph*),leftmargin=9mm]
\item To'liqlik
\item Matnli axborot
\item Grafik axborot
\item Audioaxborot
\end{enumerate}
\begin{tcolorbox}[enhanced,breakable,colback=teal!4,colframe=teal!50!black,boxrule=0.5pt,arc=1.1mm,left=3mm,right=3mm,top=1.7mm,bottom=1.7mm,before skip=2mm,after skip=2mm]
\textbf{\color{teal}To'g'ri javob:} \textbf{A) To'liqlik}
\par\smallskip
\textbf{Izoh:} Ta'rif To'liqlik tushunchasiga tegishli.
\end{tcolorbox}
\vspace{2mm}\hrule height 0.2pt \vspace{3mm}

% DB question_id: 59a01bda-d873-4935-a4aa-9fe0af562b66
\needspace{5\baselineskip}
\noindent\textbf{\color{ink}238.}
\par
Quyidagi ta'rif qaysi atamaga tegishli? Shaxsiy baho, did yoki munosabatni ifodalovchi mulohaza.
\begin{enumerate}[label=\Alph*),leftmargin=9mm]
\item Xulosa
\item Fikr
\item Axborot texnologiyasi
\item Raqamli axborot
\end{enumerate}
\begin{tcolorbox}[enhanced,breakable,colback=teal!4,colframe=teal!50!black,boxrule=0.5pt,arc=1.1mm,left=3mm,right=3mm,top=1.7mm,bottom=1.7mm,before skip=2mm,after skip=2mm]
\textbf{\color{teal}To'g'ri javob:} \textbf{B) Fikr}
\par\smallskip
\textbf{Izoh:} Ta'rif Fikr tushunchasiga tegishli.
\end{tcolorbox}
\vspace{2mm}\hrule height 0.2pt \vspace{3mm}

% DB question_id: 59a300a7-86e8-4a1f-aa01-d51ab8dc781d
\needspace{5\baselineskip}
\noindent\textbf{\color{ink}239.}
\par
Quyidagi ta'rif qaysi atamaga tegishli? Axborot va bilim ishlab chiqarish hamda xizmatlarda asosiy resursga aylangan jamiyat.
\begin{enumerate}[label=\Alph*),leftmargin=9mm]
\item Axborotlashgan jamiyat
\item Axborot madaniyati
\item Raqamli savodxonlik
\item Raqamli iz
\end{enumerate}
\begin{tcolorbox}[enhanced,breakable,colback=teal!4,colframe=teal!50!black,boxrule=0.5pt,arc=1.1mm,left=3mm,right=3mm,top=1.7mm,bottom=1.7mm,before skip=2mm,after skip=2mm]
\textbf{\color{teal}To'g'ri javob:} \textbf{A) Axborotlashgan jamiyat}
\par\smallskip
\textbf{Izoh:} Ta'rif Axborotlashgan jamiyat tushunchasiga tegishli.
\end{tcolorbox}
\vspace{2mm}\hrule height 0.2pt \vspace{3mm}

% DB question_id: 6a3c704a-e7df-4f9f-9e8a-7b4b835f3c36
\needspace{5\baselineskip}
\noindent\textbf{\color{ink}240.}
\par
Quyidagi ta'rif qaysi atamaga tegishli? Hosila asarni ayni yoki mos litsenziya bilan ulashishni talab qiladigan shart.
\begin{enumerate}[label=\Alph*),leftmargin=9mm]
\item CC NC
\item CC ND
\item Manba ko'rsatish
\item CC SA
\end{enumerate}
\begin{tcolorbox}[enhanced,breakable,colback=teal!4,colframe=teal!50!black,boxrule=0.5pt,arc=1.1mm,left=3mm,right=3mm,top=1.7mm,bottom=1.7mm,before skip=2mm,after skip=2mm]
\textbf{\color{teal}To'g'ri javob:} \textbf{D) CC SA}
\par\smallskip
\textbf{Izoh:} Ta'rif CC SA tushunchasiga tegishli.
\end{tcolorbox}
\vspace{2mm}\hrule height 0.2pt \vspace{3mm}

% DB question_id: 7a00dfa5-770d-46a6-853f-0ed6c990d9f1
\needspace{5\baselineskip}
\noindent\textbf{\color{ink}241.}
\par
Quyidagi ta'rif qaysi atamaga tegishli? Axborotning tekshirilgan dalil va ishonchli manbaga tayanganlik xossasi.
\begin{enumerate}[label=\Alph*),leftmargin=9mm]
\item To'liqlik
\item Matnli axborot
\item Grafik axborot
\item Ishonchlilik
\end{enumerate}
\begin{tcolorbox}[enhanced,breakable,colback=teal!4,colframe=teal!50!black,boxrule=0.5pt,arc=1.1mm,left=3mm,right=3mm,top=1.7mm,bottom=1.7mm,before skip=2mm,after skip=2mm]
\textbf{\color{teal}To'g'ri javob:} \textbf{D) Ishonchlilik}
\par\smallskip
\textbf{Izoh:} Ta'rif Ishonchlilik tushunchasiga tegishli.
\end{tcolorbox}
\vspace{2mm}\hrule height 0.2pt \vspace{3mm}

% DB question_id: 8dc8d509-186a-4384-b1b1-a085ddf86785
\needspace{5\baselineskip}
\noindent\textbf{\color{ink}242.}
\par
Quyidagi ta'rif qaysi atamaga tegishli? Raqamli vositalardan maqsadli, tanqidiy va samarali foydalanish qobiliyati.
\begin{enumerate}[label=\Alph*),leftmargin=9mm]
\item Raqamli iz
\item Fakt
\item Raqamli savodxonlik
\item Fikr
\end{enumerate}
\begin{tcolorbox}[enhanced,breakable,colback=teal!4,colframe=teal!50!black,boxrule=0.5pt,arc=1.1mm,left=3mm,right=3mm,top=1.7mm,bottom=1.7mm,before skip=2mm,after skip=2mm]
\textbf{\color{teal}To'g'ri javob:} \textbf{C) Raqamli savodxonlik}
\par\smallskip
\textbf{Izoh:} Ta'rif Raqamli savodxonlik tushunchasiga tegishli.
\end{tcolorbox}
\vspace{2mm}\hrule height 0.2pt \vspace{3mm}

% DB question_id: 8e87b72f-0468-438a-b984-ad99fffeeefe
\needspace{5\baselineskip}
\noindent\textbf{\color{ink}243.}
\par
Quyidagi ta'rif qaysi atamaga tegishli? Https transportni shifrlashi, sahifa mazmunining haqiqatligini kafolatlamasligini ajratish.
\begin{enumerate}[label=\Alph*),leftmargin=9mm]
\item Axborot va ma'lumot farqi
\item HTTPS va ishonchlilik farqi
\item Kodlash va shifrlash farqi
\item Bit va bayt farqi
\end{enumerate}
\begin{tcolorbox}[enhanced,breakable,colback=teal!4,colframe=teal!50!black,boxrule=0.5pt,arc=1.1mm,left=3mm,right=3mm,top=1.7mm,bottom=1.7mm,before skip=2mm,after skip=2mm]
\textbf{\color{teal}To'g'ri javob:} \textbf{B) HTTPS va ishonchlilik farqi}
\par\smallskip
\textbf{Izoh:} Ta'rif HTTPS va ishonchlilik farqi tushunchasiga tegishli.
\end{tcolorbox}
\vspace{2mm}\hrule height 0.2pt \vspace{3mm}

% DB question_id: 9327cf55-a8e7-4fc5-98bb-9f100a333f7d
\needspace{5\baselineskip}
\noindent\textbf{\color{ink}244.}
\par
Quyidagi ta'rif qaysi atamaga tegishli? Foydalanishga ruxsat berib, muallifni ko'rsatishni talab qiladigan creative commons sharti.
\begin{enumerate}[label=\Alph*),leftmargin=9mm]
\item CC SA
\item CC NC
\item CC BY
\item CC ND
\end{enumerate}
\begin{tcolorbox}[enhanced,breakable,colback=teal!4,colframe=teal!50!black,boxrule=0.5pt,arc=1.1mm,left=3mm,right=3mm,top=1.7mm,bottom=1.7mm,before skip=2mm,after skip=2mm]
\textbf{\color{teal}To'g'ri javob:} \textbf{C) CC BY}
\par\smallskip
\textbf{Izoh:} Ta'rif CC BY tushunchasiga tegishli.
\end{tcolorbox}
\vspace{2mm}\hrule height 0.2pt \vspace{3mm}

% DB question_id: 93942866-6526-46a7-b0c8-4a6d93f61a96
\needspace{5\baselineskip}
\noindent\textbf{\color{ink}245.}
\par
Quyidagi ta'rif qaysi atamaga tegishli? Turli manbalardan kerakli ma'lumotlarni olish va jamlash jarayoni.
\begin{enumerate}[label=\Alph*),leftmargin=9mm]
\item Axborotni qayta ishlash
\item Axborotni saqlash
\item Axborot manbai
\item Axborotni yig'ish
\end{enumerate}
\begin{tcolorbox}[enhanced,breakable,colback=teal!4,colframe=teal!50!black,boxrule=0.5pt,arc=1.1mm,left=3mm,right=3mm,top=1.7mm,bottom=1.7mm,before skip=2mm,after skip=2mm]
\textbf{\color{teal}To'g'ri javob:} \textbf{D) Axborotni yig'ish}
\par\smallskip
\textbf{Izoh:} Ta'rif Axborotni yig'ish tushunchasiga tegishli.
\end{tcolorbox}
\vspace{2mm}\hrule height 0.2pt \vspace{3mm}

% DB question_id: aef08623-5a6b-42b3-9334-0788967662ca
\needspace{5\baselineskip}
\noindent\textbf{\color{ink}246.}
\par
Quyidagi ta'rif qaysi atamaga tegishli? Faktlar va qoidalar asosida chiqarilgan natija.
\begin{enumerate}[label=\Alph*),leftmargin=9mm]
\item Axborot texnologiyasi
\item Raqamli axborot
\item Xulosa
\item Raqamli muhit
\end{enumerate}
\begin{tcolorbox}[enhanced,breakable,colback=teal!4,colframe=teal!50!black,boxrule=0.5pt,arc=1.1mm,left=3mm,right=3mm,top=1.7mm,bottom=1.7mm,before skip=2mm,after skip=2mm]
\textbf{\color{teal}To'g'ri javob:} \textbf{C) Xulosa}
\par\smallskip
\textbf{Izoh:} Ta'rif Xulosa tushunchasiga tegishli.
\end{tcolorbox}
\vspace{2mm}\hrule height 0.2pt \vspace{3mm}

% DB question_id: b3e02161-b61d-438e-bc9e-1865fb152144
\needspace{5\baselineskip}
\noindent\textbf{\color{ink}247.}
\par
Quyidagi ta'rif qaysi atamaga tegishli? Axborotni keyinchalik foydalanish uchun tashuvchida qayd etish jarayoni.
\begin{enumerate}[label=\Alph*),leftmargin=9mm]
\item Axborot manbai
\item Axborotni saqlash
\item Aloqa kanali
\item Qabul qiluvchi
\end{enumerate}
\begin{tcolorbox}[enhanced,breakable,colback=teal!4,colframe=teal!50!black,boxrule=0.5pt,arc=1.1mm,left=3mm,right=3mm,top=1.7mm,bottom=1.7mm,before skip=2mm,after skip=2mm]
\textbf{\color{teal}To'g'ri javob:} \textbf{B) Axborotni saqlash}
\par\smallskip
\textbf{Izoh:} Ta'rif Axborotni saqlash tushunchasiga tegishli.
\end{tcolorbox}
\vspace{2mm}\hrule height 0.2pt \vspace{3mm}

% DB question_id: bb4de605-f8a9-4d1f-b8e8-d055b29ab98f
\needspace{5\baselineskip}
\noindent\textbf{\color{ink}248.}
\par
Quyidagi ta'rif qaysi atamaga tegishli? Foydalanuvchining onlayn harakatlari natijasida qoladigan ma'lumotlar majmui.
\begin{enumerate}[label=\Alph*),leftmargin=9mm]
\item Fakt
\item Fikr
\item Xulosa
\item Raqamli iz
\end{enumerate}
\begin{tcolorbox}[enhanced,breakable,colback=teal!4,colframe=teal!50!black,boxrule=0.5pt,arc=1.1mm,left=3mm,right=3mm,top=1.7mm,bottom=1.7mm,before skip=2mm,after skip=2mm]
\textbf{\color{teal}To'g'ri javob:} \textbf{D) Raqamli iz}
\par\smallskip
\textbf{Izoh:} Ta'rif Raqamli iz tushunchasiga tegishli.
\end{tcolorbox}
\vspace{2mm}\hrule height 0.2pt \vspace{3mm}

% DB question_id: c1b1f09b-bd2f-4117-b3e2-8ad43cce743e
\needspace{5\baselineskip}
\noindent\textbf{\color{ink}249.}
\par
Quyidagi ta'rif qaysi atamaga tegishli? Asar yoki dasturdan qanday shartlarda foydalanish mumkinligini belgilovchi ruxsat.
\begin{enumerate}[label=\Alph*),leftmargin=9mm]
\item Litsenziya
\item Ommaviy mulk
\item CC BY
\item CC SA
\end{enumerate}
\begin{tcolorbox}[enhanced,breakable,colback=teal!4,colframe=teal!50!black,boxrule=0.5pt,arc=1.1mm,left=3mm,right=3mm,top=1.7mm,bottom=1.7mm,before skip=2mm,after skip=2mm]
\textbf{\color{teal}To'g'ri javob:} \textbf{A) Litsenziya}
\par\smallskip
\textbf{Izoh:} Ta'rif Litsenziya tushunchasiga tegishli.
\end{tcolorbox}
\vspace{2mm}\hrule height 0.2pt \vspace{3mm}

% DB question_id: c762327f-ef8b-4151-add6-6aabb48a78e9
\needspace{5\baselineskip}
\noindent\textbf{\color{ink}250.}
\par
Quyidagi ta'rif qaysi atamaga tegishli? Axborotni yig'ish, saqlash, qayta ishlash, uzatish va taqdim etish vosita-usullari majmui.
\begin{enumerate}[label=\Alph*),leftmargin=9mm]
\item Raqamli axborot
\item Axborot texnologiyasi
\item Raqamli muhit
\item Axborotlashgan jamiyat
\end{enumerate}
\begin{tcolorbox}[enhanced,breakable,colback=teal!4,colframe=teal!50!black,boxrule=0.5pt,arc=1.1mm,left=3mm,right=3mm,top=1.7mm,bottom=1.7mm,before skip=2mm,after skip=2mm]
\textbf{\color{teal}To'g'ri javob:} \textbf{B) Axborot texnologiyasi}
\par\smallskip
\textbf{Izoh:} Ta'rif Axborot texnologiyasi tushunchasiga tegishli.
\end{tcolorbox}
\vspace{2mm}\hrule height 0.2pt \vspace{3mm}

% DB question_id: cc7a7a77-7b2b-4ce2-bd69-a33a4afdb7cb
\needspace{5\baselineskip}
\noindent\textbf{\color{ink}251.}
\par
Quyidagi ta'rif qaysi atamaga tegishli? Raqamli muloqotdagi hurmatli va mas'uliyatli xulq qoidalari.
\begin{enumerate}[label=\Alph*),leftmargin=9mm]
\item Mualliflik huquqi
\item Plagiat
\item Litsenziya
\item Netiket
\end{enumerate}
\begin{tcolorbox}[enhanced,breakable,colback=teal!4,colframe=teal!50!black,boxrule=0.5pt,arc=1.1mm,left=3mm,right=3mm,top=1.7mm,bottom=1.7mm,before skip=2mm,after skip=2mm]
\textbf{\color{teal}To'g'ri javob:} \textbf{D) Netiket}
\par\smallskip
\textbf{Izoh:} Ta'rif Netiket tushunchasiga tegishli.
\end{tcolorbox}
\vspace{2mm}\hrule height 0.2pt \vspace{3mm}

% DB question_id: d35df013-35b1-4be2-a04f-590cee3fcd13
\needspace{5\baselineskip}
\noindent\textbf{\color{ink}252.}
\par
Quyidagi ta'rif qaysi atamaga tegishli? Muallif, asar nomi, manba va litsenziyani imkon qadar aniq qayd etish.
\begin{enumerate}[label=\Alph*),leftmargin=9mm]
\item Netiket
\item Mualliflik huquqi
\item Manba ko'rsatish
\item Plagiat
\end{enumerate}
\begin{tcolorbox}[enhanced,breakable,colback=teal!4,colframe=teal!50!black,boxrule=0.5pt,arc=1.1mm,left=3mm,right=3mm,top=1.7mm,bottom=1.7mm,before skip=2mm,after skip=2mm]
\textbf{\color{teal}To'g'ri javob:} \textbf{C) Manba ko'rsatish}
\par\smallskip
\textbf{Izoh:} Ta'rif Manba ko'rsatish tushunchasiga tegishli.
\end{tcolorbox}
\vspace{2mm}\hrule height 0.2pt \vspace{3mm}

% DB question_id: df5f770f-113f-46df-812e-af7915f75739
\needspace{5\baselineskip}
\noindent\textbf{\color{ink}253.}
\par
Quyidagi ta'rif qaysi atamaga tegishli? Aldash maqsadida ataylab tarqatilgan noto'g'ri axborot.
\begin{enumerate}[label=\Alph*),leftmargin=9mm]
\item Internet
\item Dezinformatsiya
\item WWW
\item Veb-brauzer
\end{enumerate}
\begin{tcolorbox}[enhanced,breakable,colback=teal!4,colframe=teal!50!black,boxrule=0.5pt,arc=1.1mm,left=3mm,right=3mm,top=1.7mm,bottom=1.7mm,before skip=2mm,after skip=2mm]
\textbf{\color{teal}To'g'ri javob:} \textbf{B) Dezinformatsiya}
\par\smallskip
\textbf{Izoh:} Ta'rif Dezinformatsiya tushunchasiga tegishli.
\end{tcolorbox}
\vspace{2mm}\hrule height 0.2pt \vspace{3mm}

% DB question_id: e0ef996f-01a4-4fd2-b07e-6eabc7a39e61
\needspace{5\baselineskip}
\noindent\textbf{\color{ink}254.}
\par
Quyidagi ta'rif qaysi atamaga tegishli? Raqamli kod, odatda bitlar ketma-ketligi orqali ifodalangan axborot.
\begin{enumerate}[label=\Alph*),leftmargin=9mm]
\item Raqamli muhit
\item Axborotlashgan jamiyat
\item Raqamli axborot
\item Axborot madaniyati
\end{enumerate}
\begin{tcolorbox}[enhanced,breakable,colback=teal!4,colframe=teal!50!black,boxrule=0.5pt,arc=1.1mm,left=3mm,right=3mm,top=1.7mm,bottom=1.7mm,before skip=2mm,after skip=2mm]
\textbf{\color{teal}To'g'ri javob:} \textbf{C) Raqamli axborot}
\par\smallskip
\textbf{Izoh:} Ta'rif Raqamli axborot tushunchasiga tegishli.
\end{tcolorbox}
\vspace{2mm}\hrule height 0.2pt \vspace{3mm}

% DB question_id: f24bdc1f-86d4-4c4e-9bc0-f9d3ea4dfa0e
\needspace{5\baselineskip}
\noindent\textbf{\color{ink}255.}
\par
Quyidagi ta'rif qaysi atamaga tegishli? Axborotning obyekt holatini xatosiz va kerakli aniqlik darajasida aks ettirish xossasi.
\begin{enumerate}[label=\Alph*),leftmargin=9mm]
\item Ishonchlilik
\item To'liqlik
\item Aniqlik
\item Matnli axborot
\end{enumerate}
\begin{tcolorbox}[enhanced,breakable,colback=teal!4,colframe=teal!50!black,boxrule=0.5pt,arc=1.1mm,left=3mm,right=3mm,top=1.7mm,bottom=1.7mm,before skip=2mm,after skip=2mm]
\textbf{\color{teal}To'g'ri javob:} \textbf{C) Aniqlik}
\par\smallskip
\textbf{Izoh:} Ta'rif Aniqlik tushunchasiga tegishli.
\end{tcolorbox}
\vspace{2mm}\hrule height 0.2pt \vspace{3mm}

% DB question_id: fb3bb82a-916b-4567-865a-1a94dbdc5ca2
\needspace{5\baselineskip}
\noindent\textbf{\color{ink}256.}
\par
Quyidagi ta'rif qaysi atamaga tegishli? Mualliflik huquqi cheklovlari tugagan yoki qo'llanmaydigan asarlar holati.
\begin{enumerate}[label=\Alph*),leftmargin=9mm]
\item CC BY
\item Ommaviy mulk
\item CC SA
\item CC NC
\end{enumerate}
\begin{tcolorbox}[enhanced,breakable,colback=teal!4,colframe=teal!50!black,boxrule=0.5pt,arc=1.1mm,left=3mm,right=3mm,top=1.7mm,bottom=1.7mm,before skip=2mm,after skip=2mm]
\textbf{\color{teal}To'g'ri javob:} \textbf{B) Ommaviy mulk}
\par\smallskip
\textbf{Izoh:} Ta'rif Ommaviy mulk tushunchasiga tegishli.
\end{tcolorbox}
\vspace{2mm}\hrule height 0.2pt \vspace{3mm}

% DB question_id: 0083f32c-ef8d-5621-953c-6100972f81eb
\needspace{5\baselineskip}
\noindent\textbf{\color{ink}257.}
\par
"Mantiqiy qidiruv operatorlari" atamasining eng aniq ta'rifini toping.
\begin{enumerate}[label=\Alph*),leftmargin=9mm]
\item muallif, sana, dalil, maqsad va boshqa manbalar bilan moslikni tekshirish jarayoni
\item axborotning kimga, nima uchun va qaysi sharoitda taqdim etilishini hisobga olish tamoyili
\item qidiruv shartlarini AND, OR va NOT kabi operatorlar bilan aniqlashtirish usuli
\item xat mavzusini aniq yozish, hurmatli uslub, to'g'ri adreslash va ilovani tekshirish qoidalari
\end{enumerate}
\begin{tcolorbox}[enhanced,breakable,colback=teal!4,colframe=teal!50!black,boxrule=0.5pt,arc=1.1mm,left=3mm,right=3mm,top=1.7mm,bottom=1.7mm,before skip=2mm,after skip=2mm]
\textbf{\color{teal}To'g'ri javob:} \textbf{C) qidiruv shartlarini AND, OR va NOT kabi operatorlar bilan aniqlashtirish usuli}
\par\smallskip
\textbf{Izoh:} Mantiqiy qidiruv operatorlari - qidiruv shartlarini AND, OR va NOT kabi operatorlar bilan aniqlashtirish usuli.
\end{tcolorbox}
\vspace{2mm}\hrule height 0.2pt \vspace{3mm}

% DB question_id: 049b9fab-999e-5da3-8272-9c7374d65c09
\needspace{5\baselineskip}
\noindent\textbf{\color{ink}258.}
\par
"Himoyalash" atamasining eng aniq ta'rifini toping.
\begin{enumerate}[label=\Alph*),leftmargin=9mm]
\item keraksiz yoki maxfiy axborotni belgilangan tartibda qayta tiklab bo'lmaydigan darajada o'chirish jarayoni
\item veb resurslarini indekslab, foydalanuvchi yozgan muhim so'zlarga mos natijalarni topib beruvchi xizmat va so'rov elementi
\item axborotni ruxsatsiz kirish, o'zgartirish yoki yo'qolishdan saqlash choralarini qo'llash jarayoni
\item so'zlar aynan berilgan tartibda uchrashini talab qiladigan qidiruv usuli
\end{enumerate}
\begin{tcolorbox}[enhanced,breakable,colback=teal!4,colframe=teal!50!black,boxrule=0.5pt,arc=1.1mm,left=3mm,right=3mm,top=1.7mm,bottom=1.7mm,before skip=2mm,after skip=2mm]
\textbf{\color{teal}To'g'ri javob:} \textbf{C) axborotni ruxsatsiz kirish, o'zgartirish yoki yo'qolishdan saqlash choralarini qo'llash jarayoni}
\par\smallskip
\textbf{Izoh:} Himoyalash - axborotni ruxsatsiz kirish, o'zgartirish yoki yo'qolishdan saqlash choralarini qo'llash jarayoni.
\end{tcolorbox}
\vspace{2mm}\hrule height 0.2pt \vspace{3mm}

% DB question_id: 051b2505-16d1-56b1-b03c-f794b321f980
\needspace{5\baselineskip}
\noindent\textbf{\color{ink}259.}
\par
Quyidagi ta'rif qaysi tushunchaga tegishli? Axborotni ruxsatsiz kirish, o'zgartirish yoki yo'qolishdan saqlash choralarini qo'llash jarayoni.
\begin{enumerate}[label=\Alph*),leftmargin=9mm]
\item Yo'q qilish
\item Himoyalash
\item Qidiruv tizimi va kalit so'z
\item Aniq ibora bilan qidirish
\end{enumerate}
\begin{tcolorbox}[enhanced,breakable,colback=teal!4,colframe=teal!50!black,boxrule=0.5pt,arc=1.1mm,left=3mm,right=3mm,top=1.7mm,bottom=1.7mm,before skip=2mm,after skip=2mm]
\textbf{\color{teal}To'g'ri javob:} \textbf{B) Himoyalash}
\par\smallskip
\textbf{Izoh:} Ta'rif Himoyalash tushunchasiga tegishli.
\end{tcolorbox}
\vspace{2mm}\hrule height 0.2pt \vspace{3mm}

% DB question_id: 06e82163-0cb0-57ff-8ae8-5ab609a1f451
\needspace{5\baselineskip}
\noindent\textbf{\color{ink}260.}
\par
"Tarqatish" atamasining eng aniq ta'rifini toping.
\begin{enumerate}[label=\Alph*),leftmargin=9mm]
\item axborotni bir nechta foydalanuvchi yoki auditoriyaga yetkazish jarayoni
\item axborotni auditoriya va maqsadga mos shaklda ko'rsatish jarayoni
\item axborotni ruxsatsiz kirish, o'zgartirish yoki yo'qolishdan saqlash choralarini qo'llash jarayoni
\item keraksiz yoki maxfiy axborotni belgilangan tartibda qayta tiklab bo'lmaydigan darajada o'chirish jarayoni
\end{enumerate}
\begin{tcolorbox}[enhanced,breakable,colback=teal!4,colframe=teal!50!black,boxrule=0.5pt,arc=1.1mm,left=3mm,right=3mm,top=1.7mm,bottom=1.7mm,before skip=2mm,after skip=2mm]
\textbf{\color{teal}To'g'ri javob:} \textbf{A) axborotni bir nechta foydalanuvchi yoki auditoriyaga yetkazish jarayoni}
\par\smallskip
\textbf{Izoh:} Tarqatish - axborotni bir nechta foydalanuvchi yoki auditoriyaga yetkazish jarayoni.
\end{tcolorbox}
\vspace{2mm}\hrule height 0.2pt \vspace{3mm}

% DB question_id: 070e926c-6441-5edd-96c8-ef4a7e20a79a
\needspace{5\baselineskip}
\noindent\textbf{\color{ink}261.}
\par
"Guruhlash" atamasining eng aniq ta'rifini toping.
\begin{enumerate}[label=\Alph*),leftmargin=9mm]
\item axborotni manbadan kanal orqali qabul qiluvchiga yetkazish va uni olish jarayoni
\item umumiy belgi bo'yicha elementlarni toifalarga ajratish jarayoni
\item axborotning mazmunini saqlagan holda uning qo'shimcha nusxasini yaratish jarayoni
\item axborotni bir nechta foydalanuvchi yoki auditoriyaga yetkazish jarayoni
\end{enumerate}
\begin{tcolorbox}[enhanced,breakable,colback=teal!4,colframe=teal!50!black,boxrule=0.5pt,arc=1.1mm,left=3mm,right=3mm,top=1.7mm,bottom=1.7mm,before skip=2mm,after skip=2mm]
\textbf{\color{teal}To'g'ri javob:} \textbf{B) umumiy belgi bo'yicha elementlarni toifalarga ajratish jarayoni}
\par\smallskip
\textbf{Izoh:} Guruhlash - umumiy belgi bo'yicha elementlarni toifalarga ajratish jarayoni.
\end{tcolorbox}
\vspace{2mm}\hrule height 0.2pt \vspace{3mm}

% DB question_id: 080734d8-23a4-4f81-ba82-9a963de0f891
\needspace{5\baselineskip}
\noindent\textbf{\color{ink}262.}
\par
Internetdan topilgan xabarda muallif, sana va tasdiqlovchi dalillar ko'rsatilmagan. Uni sinf guruhiga ulashishdan oldin eng mas'uliyatli ish qaysi?
\begin{enumerate}[label=\Alph*),leftmargin=9mm]
\item Post matnini qiziqarli sarlavha bilan qayta yozib, dalillarni tekshirmasdan sinf guruhiga darhol tarqatish
\item Ko'p izoh yozilganini yetarli dalil deb olib, muallifi noma'lum xabarni hech tekshirmasdan tarqatish
\item Boshqa noma'lum kanallardagi shunga o'xshash xabarlarni ko'rib, manba va sanani mustaqil tekshirmasdan yuborish
\item Manbani tekshirish, boshqa ishonchli manbalar bilan solishtirish va auditoriyaga mos ravishda taqdim etish
\end{enumerate}
\begin{tcolorbox}[enhanced,breakable,colback=teal!4,colframe=teal!50!black,boxrule=0.5pt,arc=1.1mm,left=3mm,right=3mm,top=1.7mm,bottom=1.7mm,before skip=2mm,after skip=2mm]
\textbf{\color{teal}To'g'ri javob:} \textbf{D) Manbani tekshirish, boshqa ishonchli manbalar bilan solishtirish va auditoriyaga mos ravishda taqdim etish}
\par\smallskip
\textbf{Izoh:} Muallif, sana va dalil bo'lmagan post mustaqil tekshirilishi, boshqa ishonchli manbalar bilan qiyoslanishi va auditoriya ehtiyojiga mos ravishda ehtiyotkorlik bilan taqdim etilishi kerak. Omma fikri yoki takroriy repost dalil emas.
\end{tcolorbox}
\vspace{2mm}\hrule height 0.2pt \vspace{3mm}

% DB question_id: 135b0bed-8db2-5157-bb4e-02556fcad1a5
\needspace{5\baselineskip}
\noindent\textbf{\color{ink}263.}
\par
Maqola muallifi va dalillari rasmiy manba bilan solishtirildi. Ushbu vaziyatdagi asosiy tushunchani aniqlang.
\begin{enumerate}[label=\Alph*),leftmargin=9mm]
\item Auditoriya, maqsad va kontekst
\item Qidiruv natijalarini baholash
\item Elektron pochta madaniyati
\item Netiket
\end{enumerate}
\begin{tcolorbox}[enhanced,breakable,colback=teal!4,colframe=teal!50!black,boxrule=0.5pt,arc=1.1mm,left=3mm,right=3mm,top=1.7mm,bottom=1.7mm,before skip=2mm,after skip=2mm]
\textbf{\color{teal}To'g'ri javob:} \textbf{B) Qidiruv natijalarini baholash}
\par\smallskip
\textbf{Izoh:} Vaziyat Qidiruv natijalarini baholash tushunchasini ifodalaydi: muallif, sana, dalil, maqsad va boshqa manbalar bilan moslikni tekshirish jarayoni.
\end{tcolorbox}
\vspace{2mm}\hrule height 0.2pt \vspace{3mm}

% DB question_id: 16d4b852-c31b-415a-b91b-6a3e3148a7eb
\needspace{5\baselineskip}
\noindent\textbf{\color{ink}264.}
\par
Rasm qayta ishlatildi va muallif nomi hamda havola berildi. Bu vaziyatdagi asosiy tushunchani aniqlang.
\begin{enumerate}[label=\Alph*),leftmargin=9mm]
\item CC BY
\item CC SA
\item CC NC
\item CC ND
\end{enumerate}
\begin{tcolorbox}[enhanced,breakable,colback=teal!4,colframe=teal!50!black,boxrule=0.5pt,arc=1.1mm,left=3mm,right=3mm,top=1.7mm,bottom=1.7mm,before skip=2mm,after skip=2mm]
\textbf{\color{teal}To'g'ri javob:} \textbf{A) CC BY}
\par\smallskip
\textbf{Izoh:} Vaziyat CC BYga misol: foydalanishga ruxsat berib, muallifni ko'rsatishni talab qiladigan Creative Commons sharti.
\end{tcolorbox}
\vspace{2mm}\hrule height 0.2pt \vspace{3mm}

% DB question_id: 1ae81b88-4044-4ad5-8077-658f640bcf1f
\needspace{5\baselineskip}
\noindent\textbf{\color{ink}265.}
\par
O'quvchi maqola parchasini muallifni ko'rsatmasdan topshirdi. Bu vaziyatdagi asosiy tushunchani aniqlang.
\begin{enumerate}[label=\Alph*),leftmargin=9mm]
\item Litsenziya
\item Plagiat
\item Ommaviy mulk
\item CC BY
\end{enumerate}
\begin{tcolorbox}[enhanced,breakable,colback=teal!4,colframe=teal!50!black,boxrule=0.5pt,arc=1.1mm,left=3mm,right=3mm,top=1.7mm,bottom=1.7mm,before skip=2mm,after skip=2mm]
\textbf{\color{teal}To'g'ri javob:} \textbf{B) Plagiat}
\par\smallskip
\textbf{Izoh:} Vaziyat Plagiatga misol: boshqa shaxs g'oyasi yoki ifodasini manbasiz o'ziniki qilib ko'rsatish.
\end{tcolorbox}
\vspace{2mm}\hrule height 0.2pt \vspace{3mm}

% DB question_id: 1dd23bba-f2d4-5c54-a87b-36439e72e5a0
\needspace{5\baselineskip}
\noindent\textbf{\color{ink}266.}
\par
Xat yuborishdan oldin mavzu, adresat va biriktirma qayta tekshirildi. Ushbu vaziyatdagi asosiy tushunchani aniqlang.
\begin{enumerate}[label=\Alph*),leftmargin=9mm]
\item Netiket
\item Raqamli iz
\item Axborotni taqdim etish shaklini tanlash
\item Elektron pochta madaniyati
\end{enumerate}
\begin{tcolorbox}[enhanced,breakable,colback=teal!4,colframe=teal!50!black,boxrule=0.5pt,arc=1.1mm,left=3mm,right=3mm,top=1.7mm,bottom=1.7mm,before skip=2mm,after skip=2mm]
\textbf{\color{teal}To'g'ri javob:} \textbf{D) Elektron pochta madaniyati}
\par\smallskip
\textbf{Izoh:} Vaziyat Elektron pochta madaniyati tushunchasini ifodalaydi: xat mavzusini aniq yozish, hurmatli uslub, to'g'ri adreslash va ilovani tekshirish qoidalari.
\end{tcolorbox}
\vspace{2mm}\hrule height 0.2pt \vspace{3mm}

% DB question_id: 1f79d530-0c68-4a16-b602-53c97ee05c5b
\needspace{5\baselineskip}
\noindent\textbf{\color{ink}267.}
\par
Asardan mualliflik ruxsatisiz foydalanish mumkin, lekin soxta mualliflik qilish mumkin emas. Bu vaziyatdagi asosiy tushunchani aniqlang.
\begin{enumerate}[label=\Alph*),leftmargin=9mm]
\item CC BY
\item CC SA
\item CC NC
\item Ommaviy mulk
\end{enumerate}
\begin{tcolorbox}[enhanced,breakable,colback=teal!4,colframe=teal!50!black,boxrule=0.5pt,arc=1.1mm,left=3mm,right=3mm,top=1.7mm,bottom=1.7mm,before skip=2mm,after skip=2mm]
\textbf{\color{teal}To'g'ri javob:} \textbf{D) Ommaviy mulk}
\par\smallskip
\textbf{Izoh:} Vaziyat Ommaviy mulkga misol: mualliflik huquqi cheklovlari tugagan yoki qo'llanmaydigan asarlar holati.
\end{tcolorbox}
\vspace{2mm}\hrule height 0.2pt \vspace{3mm}

% DB question_id: 2aa96b5e-5571-5b1e-93b6-a51a715c20a9
\needspace{5\baselineskip}
\noindent\textbf{\color{ink}268.}
\par
Quyidagi ta'rif qaysi tushunchaga tegishli? Qidiruv shartlarini and, or va not kabi operatorlar bilan aniqlashtirish usuli.
\begin{enumerate}[label=\Alph*),leftmargin=9mm]
\item Qidiruv natijalarini baholash
\item Mantiqiy qidiruv operatorlari
\item Auditoriya, maqsad va kontekst
\item Elektron pochta madaniyati
\end{enumerate}
\begin{tcolorbox}[enhanced,breakable,colback=teal!4,colframe=teal!50!black,boxrule=0.5pt,arc=1.1mm,left=3mm,right=3mm,top=1.7mm,bottom=1.7mm,before skip=2mm,after skip=2mm]
\textbf{\color{teal}To'g'ri javob:} \textbf{B) Mantiqiy qidiruv operatorlari}
\par\smallskip
\textbf{Izoh:} Ta'rif Mantiqiy qidiruv operatorlari tushunchasiga tegishli.
\end{tcolorbox}
\vspace{2mm}\hrule height 0.2pt \vspace{3mm}

% DB question_id: 2ad57c9b-56eb-5ad3-88cd-82e5ea854ddd
\needspace{5\baselineskip}
\noindent\textbf{\color{ink}269.}
\par
Rahbariyat uchun natijalar qisqa diagrammada berildi. Ushbu vaziyatdagi asosiy tushunchani aniqlang.
\begin{enumerate}[label=\Alph*),leftmargin=9mm]
\item Himoyalash
\item Yo'q qilish
\item Qidiruv tizimi va kalit so'z
\item Taqdim etish
\end{enumerate}
\begin{tcolorbox}[enhanced,breakable,colback=teal!4,colframe=teal!50!black,boxrule=0.5pt,arc=1.1mm,left=3mm,right=3mm,top=1.7mm,bottom=1.7mm,before skip=2mm,after skip=2mm]
\textbf{\color{teal}To'g'ri javob:} \textbf{D) Taqdim etish}
\par\smallskip
\textbf{Izoh:} Vaziyat Taqdim etish tushunchasini ifodalaydi: axborotni auditoriya va maqsadga mos shaklda ko'rsatish jarayoni.
\end{tcolorbox}
\vspace{2mm}\hrule height 0.2pt \vspace{3mm}

% DB question_id: 2e60b13a-9908-59d0-991f-96a00af461fe
\needspace{5\baselineskip}
\noindent\textbf{\color{ink}270.}
\par
Quyidagi ta'rif qaysi tushunchaga tegishli? Axborotni keyinchalik foydalanish uchun saqlab qolish hamda uzoq muddat tartibli jamlash jarayoni.
\begin{enumerate}[label=\Alph*),leftmargin=9mm]
\item Qayta ishlash va tahlil
\item Izlash
\item Saralash
\item Saqlash va arxivlash
\end{enumerate}
\begin{tcolorbox}[enhanced,breakable,colback=teal!4,colframe=teal!50!black,boxrule=0.5pt,arc=1.1mm,left=3mm,right=3mm,top=1.7mm,bottom=1.7mm,before skip=2mm,after skip=2mm]
\textbf{\color{teal}To'g'ri javob:} \textbf{D) Saqlash va arxivlash}
\par\smallskip
\textbf{Izoh:} Ta'rif Saqlash va arxivlash tushunchasiga tegishli.
\end{tcolorbox}
\vspace{2mm}\hrule height 0.2pt \vspace{3mm}

% DB question_id: 345dce1a-43a7-5b24-9276-40498d82d16d
\needspace{5\baselineskip}
\noindent\textbf{\color{ink}271.}
\par
Quyidagi ta'rif qaysi tushunchaga tegishli? Axborotni sezgi yoki qurilma orqali olish va uni axborot tizimiga kiritish jarayoni.
\begin{enumerate}[label=\Alph*),leftmargin=9mm]
\item Saqlash va arxivlash
\item Qayta ishlash va tahlil
\item Qabul qilish va kiritish
\item Izlash
\end{enumerate}
\begin{tcolorbox}[enhanced,breakable,colback=teal!4,colframe=teal!50!black,boxrule=0.5pt,arc=1.1mm,left=3mm,right=3mm,top=1.7mm,bottom=1.7mm,before skip=2mm,after skip=2mm]
\textbf{\color{teal}To'g'ri javob:} \textbf{C) Qabul qilish va kiritish}
\par\smallskip
\textbf{Izoh:} Ta'rif Qabul qilish va kiritish tushunchasiga tegishli.
\end{tcolorbox}
\vspace{2mm}\hrule height 0.2pt \vspace{3mm}

% DB question_id: 37eb396c-371a-5dd9-90c4-97717f775c46
\needspace{5\baselineskip}
\noindent\textbf{\color{ink}272.}
\par
"Qabul qilish va kiritish" atamasining eng aniq ta'rifini toping.
\begin{enumerate}[label=\Alph*),leftmargin=9mm]
\item axborotni keyinchalik foydalanish uchun saqlab qolish hamda uzoq muddat tartibli jamlash jarayoni
\item ma'lumotdan yangi natija olish, undagi bog'lanish va xulosalarni aniqlash jarayoni
\item kerakli axborotni manbalar yoki ma'lumotlar to'plamidan topish jarayoni
\item axborotni sezgi yoki qurilma orqali olish va uni axborot tizimiga kiritish jarayoni
\end{enumerate}
\begin{tcolorbox}[enhanced,breakable,colback=teal!4,colframe=teal!50!black,boxrule=0.5pt,arc=1.1mm,left=3mm,right=3mm,top=1.7mm,bottom=1.7mm,before skip=2mm,after skip=2mm]
\textbf{\color{teal}To'g'ri javob:} \textbf{D) axborotni sezgi yoki qurilma orqali olish va uni axborot tizimiga kiritish jarayoni}
\par\smallskip
\textbf{Izoh:} Qabul qilish va kiritish - axborotni sezgi yoki qurilma orqali olish va uni axborot tizimiga kiritish jarayoni.
\end{tcolorbox}
\vspace{2mm}\hrule height 0.2pt \vspace{3mm}

% DB question_id: 42867d0d-0b9b-5c91-b58b-a332a8dfe316
\needspace{5\baselineskip}
\noindent\textbf{\color{ink}273.}
\par
Quyidagi ta'rif qaysi tushunchaga tegishli? Axborotning mazmunini saqlagan holda uning qo'shimcha nusxasini yaratish jarayoni.
\begin{enumerate}[label=\Alph*),leftmargin=9mm]
\item Tarqatish
\item Taqdim etish
\item Nusxalash
\item Himoyalash
\end{enumerate}
\begin{tcolorbox}[enhanced,breakable,colback=teal!4,colframe=teal!50!black,boxrule=0.5pt,arc=1.1mm,left=3mm,right=3mm,top=1.7mm,bottom=1.7mm,before skip=2mm,after skip=2mm]
\textbf{\color{teal}To'g'ri javob:} \textbf{C) Nusxalash}
\par\smallskip
\textbf{Izoh:} Ta'rif Nusxalash tushunchasiga tegishli.
\end{tcolorbox}
\vspace{2mm}\hrule height 0.2pt \vspace{3mm}

% DB question_id: 43359b14-ba63-5537-9702-456bc431fa24
\needspace{5\baselineskip}
\noindent\textbf{\color{ink}274.}
\par
Quyidagi ta'rif qaysi tushunchaga tegishli? So'zlar aynan berilgan tartibda uchrashini talab qiladigan qidiruv usuli.
\begin{enumerate}[label=\Alph*),leftmargin=9mm]
\item Aniq ibora bilan qidirish
\item Mantiqiy qidiruv operatorlari
\item Qidiruv natijalarini baholash
\item Auditoriya, maqsad va kontekst
\end{enumerate}
\begin{tcolorbox}[enhanced,breakable,colback=teal!4,colframe=teal!50!black,boxrule=0.5pt,arc=1.1mm,left=3mm,right=3mm,top=1.7mm,bottom=1.7mm,before skip=2mm,after skip=2mm]
\textbf{\color{teal}To'g'ri javob:} \textbf{A) Aniq ibora bilan qidirish}
\par\smallskip
\textbf{Izoh:} Ta'rif Aniq ibora bilan qidirish tushunchasiga tegishli.
\end{tcolorbox}
\vspace{2mm}\hrule height 0.2pt \vspace{3mm}

% DB question_id: 43ee115a-e857-5f62-a92d-7cb17771d456
\needspace{5\baselineskip}
\noindent\textbf{\color{ink}275.}
\par
Quyidagi ta'rif qaysi tushunchaga tegishli? Axborotni auditoriya va maqsadga mos shaklda ko'rsatish jarayoni.
\begin{enumerate}[label=\Alph*),leftmargin=9mm]
\item Taqdim etish
\item Himoyalash
\item Yo'q qilish
\item Qidiruv tizimi va kalit so'z
\end{enumerate}
\begin{tcolorbox}[enhanced,breakable,colback=teal!4,colframe=teal!50!black,boxrule=0.5pt,arc=1.1mm,left=3mm,right=3mm,top=1.7mm,bottom=1.7mm,before skip=2mm,after skip=2mm]
\textbf{\color{teal}To'g'ri javob:} \textbf{A) Taqdim etish}
\par\smallskip
\textbf{Izoh:} Ta'rif Taqdim etish tushunchasiga tegishli.
\end{tcolorbox}
\vspace{2mm}\hrule height 0.2pt \vspace{3mm}

% DB question_id: 455f6a6f-c663-58a4-a33a-54c0582caeca
\needspace{5\baselineskip}
\noindent\textbf{\color{ink}276.}
\par
"Netiket" atamasining eng aniq ta'rifini toping.
\begin{enumerate}[label=\Alph*),leftmargin=9mm]
\item onlayn faoliyat natijasida xizmatlar va qurilmalarda qoladigan ma'lumotlar majmui
\item mazmunni auditoriya ehtiyoji va maqsadga eng mos matn, jadval, diagramma, audio yoki video shaklida berish qarori
\item raqamli muloqotda hurmat, aniqlik va xavfsizlikka rioya qilish qoidalari
\item yangi mazmun hosil qilish hamda kerakli ma'lumotlarni turli manbalardan to'plash jarayoni
\end{enumerate}
\begin{tcolorbox}[enhanced,breakable,colback=teal!4,colframe=teal!50!black,boxrule=0.5pt,arc=1.1mm,left=3mm,right=3mm,top=1.7mm,bottom=1.7mm,before skip=2mm,after skip=2mm]
\textbf{\color{teal}To'g'ri javob:} \textbf{C) raqamli muloqotda hurmat, aniqlik va xavfsizlikka rioya qilish qoidalari}
\par\smallskip
\textbf{Izoh:} Netiket - raqamli muloqotda hurmat, aniqlik va xavfsizlikka rioya qilish qoidalari.
\end{tcolorbox}
\vspace{2mm}\hrule height 0.2pt \vspace{3mm}

% DB question_id: 4c7289f8-96df-5ecb-adcc-3e1db3c1a8f8
\needspace{5\baselineskip}
\noindent\textbf{\color{ink}277.}
\par
"Qidiruv natijalarini baholash" atamasining eng aniq ta'rifini toping.
\begin{enumerate}[label=\Alph*),leftmargin=9mm]
\item axborotning kimga, nima uchun va qaysi sharoitda taqdim etilishini hisobga olish tamoyili
\item xat mavzusini aniq yozish, hurmatli uslub, to'g'ri adreslash va ilovani tekshirish qoidalari
\item raqamli muloqotda hurmat, aniqlik va xavfsizlikka rioya qilish qoidalari
\item muallif, sana, dalil, maqsad va boshqa manbalar bilan moslikni tekshirish jarayoni
\end{enumerate}
\begin{tcolorbox}[enhanced,breakable,colback=teal!4,colframe=teal!50!black,boxrule=0.5pt,arc=1.1mm,left=3mm,right=3mm,top=1.7mm,bottom=1.7mm,before skip=2mm,after skip=2mm]
\textbf{\color{teal}To'g'ri javob:} \textbf{D) muallif, sana, dalil, maqsad va boshqa manbalar bilan moslikni tekshirish jarayoni}
\par\smallskip
\textbf{Izoh:} Qidiruv natijalarini baholash - muallif, sana, dalil, maqsad va boshqa manbalar bilan moslikni tekshirish jarayoni.
\end{tcolorbox}
\vspace{2mm}\hrule height 0.2pt \vspace{3mm}

% DB question_id: 4da0ecbd-eea8-5599-8ad6-3f3ebea890e4
\needspace{5\baselineskip}
\noindent\textbf{\color{ink}278.}
\par
Quyidagi ta'rif qaysi tushunchaga tegishli? Veb resurslarini indekslab, foydalanuvchi yozgan muhim so'zlarga mos natijalarni topib beruvchi xizmat va so'rov elementi.
\begin{enumerate}[label=\Alph*),leftmargin=9mm]
\item Aniq ibora bilan qidirish
\item Mantiqiy qidiruv operatorlari
\item Qidiruv natijalarini baholash
\item Qidiruv tizimi va kalit so'z
\end{enumerate}
\begin{tcolorbox}[enhanced,breakable,colback=teal!4,colframe=teal!50!black,boxrule=0.5pt,arc=1.1mm,left=3mm,right=3mm,top=1.7mm,bottom=1.7mm,before skip=2mm,after skip=2mm]
\textbf{\color{teal}To'g'ri javob:} \textbf{D) Qidiruv tizimi va kalit so'z}
\par\smallskip
\textbf{Izoh:} Ta'rif Qidiruv tizimi va kalit so'z tushunchasiga tegishli.
\end{tcolorbox}
\vspace{2mm}\hrule height 0.2pt \vspace{3mm}

% DB question_id: 4f160e4c-99f0-4440-b37f-2b627d7c1768
\needspace{5\baselineskip}
\noindent\textbf{\color{ink}279.}
\par
Raqamli materiallardan quyidagi qaysi holatda foydalanish odatda mualliflik huquqiga mos keladi?
\begin{enumerate}[label=\Alph*),leftmargin=9mm]
\item Logotipni ruxsatsiz o'zgartirib joylash
\item Pullik kursni ruxsatsiz bepul tarqatish
\item Muallif ruxsatisiz materialni yuklab tarqatish
\item Qonuniy iqtibos doirasida manbani ko'rsatib foydalanish
\end{enumerate}
\begin{tcolorbox}[enhanced,breakable,colback=teal!4,colframe=teal!50!black,boxrule=0.5pt,arc=1.1mm,left=3mm,right=3mm,top=1.7mm,bottom=1.7mm,before skip=2mm,after skip=2mm]
\textbf{\color{teal}To'g'ri javob:} \textbf{D) Qonuniy iqtibos doirasida manbani ko'rsatib foydalanish}
\par\smallskip
\textbf{Izoh:} Iqtibos uchun qonunda belgilangan shartlar bajarilishi va manba ko'rsatilishi kerak.
\end{tcolorbox}
\vspace{2mm}\hrule height 0.2pt \vspace{3mm}

% DB question_id: 513b9193-8c05-5518-8636-e5bc8f4c4da6
\needspace{5\baselineskip}
\noindent\textbf{\color{ink}280.}
\par
"Filtrlash" atamasining eng aniq ta'rifini toping.
\begin{enumerate}[label=\Alph*),leftmargin=9mm]
\item faqat berilgan shartga mos yozuvlarni ajratib ko'rsatish jarayoni
\item umumiy belgi bo'yicha elementlarni toifalarga ajratish jarayoni
\item axborotni manbadan kanal orqali qabul qiluvchiga yetkazish va uni olish jarayoni
\item axborotning mazmunini saqlagan holda uning qo'shimcha nusxasini yaratish jarayoni
\end{enumerate}
\begin{tcolorbox}[enhanced,breakable,colback=teal!4,colframe=teal!50!black,boxrule=0.5pt,arc=1.1mm,left=3mm,right=3mm,top=1.7mm,bottom=1.7mm,before skip=2mm,after skip=2mm]
\textbf{\color{teal}To'g'ri javob:} \textbf{A) faqat berilgan shartga mos yozuvlarni ajratib ko'rsatish jarayoni}
\par\smallskip
\textbf{Izoh:} Filtrlash - faqat berilgan shartga mos yozuvlarni ajratib ko'rsatish jarayoni.
\end{tcolorbox}
\vspace{2mm}\hrule height 0.2pt \vspace{3mm}

% DB question_id: 51e483db-ffea-5641-b9d2-72031525e0f9
\needspace{5\baselineskip}
\noindent\textbf{\color{ink}281.}
\par
"Saralash" atamasining eng aniq ta'rifini toping.
\begin{enumerate}[label=\Alph*),leftmargin=9mm]
\item faqat berilgan shartga mos yozuvlarni ajratib ko'rsatish jarayoni
\item umumiy belgi bo'yicha elementlarni toifalarga ajratish jarayoni
\item axborotni manbadan kanal orqali qabul qiluvchiga yetkazish va uni olish jarayoni
\item elementlarni tanlangan mezon bo'yicha tartibga keltirish jarayoni
\end{enumerate}
\begin{tcolorbox}[enhanced,breakable,colback=teal!4,colframe=teal!50!black,boxrule=0.5pt,arc=1.1mm,left=3mm,right=3mm,top=1.7mm,bottom=1.7mm,before skip=2mm,after skip=2mm]
\textbf{\color{teal}To'g'ri javob:} \textbf{D) elementlarni tanlangan mezon bo'yicha tartibga keltirish jarayoni}
\par\smallskip
\textbf{Izoh:} Saralash - elementlarni tanlangan mezon bo'yicha tartibga keltirish jarayoni.
\end{tcolorbox}
\vspace{2mm}\hrule height 0.2pt \vspace{3mm}

% DB question_id: 54384f7d-df44-5572-8e90-f745719adfa4
\needspace{5\baselineskip}
\noindent\textbf{\color{ink}282.}
\par
"Yo'q qilish" atamasining eng aniq ta'rifini toping.
\begin{enumerate}[label=\Alph*),leftmargin=9mm]
\item veb resurslarini indekslab, foydalanuvchi yozgan muhim so'zlarga mos natijalarni topib beruvchi xizmat va so'rov elementi
\item so'zlar aynan berilgan tartibda uchrashini talab qiladigan qidiruv usuli
\item qidiruv shartlarini AND, OR va NOT kabi operatorlar bilan aniqlashtirish usuli
\item keraksiz yoki maxfiy axborotni belgilangan tartibda qayta tiklab bo'lmaydigan darajada o'chirish jarayoni
\end{enumerate}
\begin{tcolorbox}[enhanced,breakable,colback=teal!4,colframe=teal!50!black,boxrule=0.5pt,arc=1.1mm,left=3mm,right=3mm,top=1.7mm,bottom=1.7mm,before skip=2mm,after skip=2mm]
\textbf{\color{teal}To'g'ri javob:} \textbf{D) keraksiz yoki maxfiy axborotni belgilangan tartibda qayta tiklab bo'lmaydigan darajada o'chirish jarayoni}
\par\smallskip
\textbf{Izoh:} Yo'q qilish - keraksiz yoki maxfiy axborotni belgilangan tartibda qayta tiklab bo'lmaydigan darajada o'chirish jarayoni.
\end{tcolorbox}
\vspace{2mm}\hrule height 0.2pt \vspace{3mm}

% DB question_id: 578a7074-0168-5d45-aa04-0ce2430f448f
\needspace{5\baselineskip}
\noindent\textbf{\color{ink}283.}
\par
"Nusxalash" atamasining eng aniq ta'rifini toping.
\begin{enumerate}[label=\Alph*),leftmargin=9mm]
\item axborotni bir nechta foydalanuvchi yoki auditoriyaga yetkazish jarayoni
\item axborotni auditoriya va maqsadga mos shaklda ko'rsatish jarayoni
\item axborotni ruxsatsiz kirish, o'zgartirish yoki yo'qolishdan saqlash choralarini qo'llash jarayoni
\item axborotning mazmunini saqlagan holda uning qo'shimcha nusxasini yaratish jarayoni
\end{enumerate}
\begin{tcolorbox}[enhanced,breakable,colback=teal!4,colframe=teal!50!black,boxrule=0.5pt,arc=1.1mm,left=3mm,right=3mm,top=1.7mm,bottom=1.7mm,before skip=2mm,after skip=2mm]
\textbf{\color{teal}To'g'ri javob:} \textbf{D) axborotning mazmunini saqlagan holda uning qo'shimcha nusxasini yaratish jarayoni}
\par\smallskip
\textbf{Izoh:} Nusxalash - axborotning mazmunini saqlagan holda uning qo'shimcha nusxasini yaratish jarayoni.
\end{tcolorbox}
\vspace{2mm}\hrule height 0.2pt \vspace{3mm}

% DB question_id: 58c1a6d7-047d-568b-b055-b7d0f1eb7aeb
\needspace{5\baselineskip}
\noindent\textbf{\color{ink}284.}
\par
"Axborotni yaratish va yig'ish" atamasining eng aniq ta'rifini toping.
\begin{enumerate}[label=\Alph*),leftmargin=9mm]
\item axborotni sezgi yoki qurilma orqali olish va uni axborot tizimiga kiritish jarayoni
\item axborotni keyinchalik foydalanish uchun saqlab qolish hamda uzoq muddat tartibli jamlash jarayoni
\item yangi mazmun hosil qilish hamda kerakli ma'lumotlarni turli manbalardan to'plash jarayoni
\item ma'lumotdan yangi natija olish, undagi bog'lanish va xulosalarni aniqlash jarayoni
\end{enumerate}
\begin{tcolorbox}[enhanced,breakable,colback=teal!4,colframe=teal!50!black,boxrule=0.5pt,arc=1.1mm,left=3mm,right=3mm,top=1.7mm,bottom=1.7mm,before skip=2mm,after skip=2mm]
\textbf{\color{teal}To'g'ri javob:} \textbf{C) yangi mazmun hosil qilish hamda kerakli ma'lumotlarni turli manbalardan to'plash jarayoni}
\par\smallskip
\textbf{Izoh:} Axborotni yaratish va yig'ish - yangi mazmun hosil qilish hamda kerakli ma'lumotlarni turli manbalardan to'plash jarayoni.
\end{tcolorbox}
\vspace{2mm}\hrule height 0.2pt \vspace{3mm}

% DB question_id: 5ac8f89a-a59f-57ad-bd6d-6fe9dc609d4c
\needspace{5\baselineskip}
\noindent\textbf{\color{ink}285.}
\par
"Qidiruv tizimi va kalit so'z" atamasining eng aniq ta'rifini toping.
\begin{enumerate}[label=\Alph*),leftmargin=9mm]
\item veb resurslarini indekslab, foydalanuvchi yozgan muhim so'zlarga mos natijalarni topib beruvchi xizmat va so'rov elementi
\item so'zlar aynan berilgan tartibda uchrashini talab qiladigan qidiruv usuli
\item qidiruv shartlarini AND, OR va NOT kabi operatorlar bilan aniqlashtirish usuli
\item muallif, sana, dalil, maqsad va boshqa manbalar bilan moslikni tekshirish jarayoni
\end{enumerate}
\begin{tcolorbox}[enhanced,breakable,colback=teal!4,colframe=teal!50!black,boxrule=0.5pt,arc=1.1mm,left=3mm,right=3mm,top=1.7mm,bottom=1.7mm,before skip=2mm,after skip=2mm]
\textbf{\color{teal}To'g'ri javob:} \textbf{A) veb resurslarini indekslab, foydalanuvchi yozgan muhim so'zlarga mos natijalarni topib beruvchi xizmat va so'rov elementi}
\par\smallskip
\textbf{Izoh:} Qidiruv tizimi va kalit so'z - veb resurslarini indekslab, foydalanuvchi yozgan muhim so'zlarga mos natijalarni topib beruvchi xizmat va so'rov elementi.
\end{tcolorbox}
\vspace{2mm}\hrule height 0.2pt \vspace{3mm}

% DB question_id: 5f195fd4-6298-5047-9285-5f1f2ec9dcf3
\needspace{5\baselineskip}
\noindent\textbf{\color{ink}286.}
\par
Guruh xabarida haqorat ishlatilmay, mavzu aniq yozildi. Ushbu vaziyatdagi asosiy tushunchani aniqlang.
\begin{enumerate}[label=\Alph*),leftmargin=9mm]
\item Netiket
\item Raqamli iz
\item Axborotni taqdim etish shaklini tanlash
\item Axborotni yaratish va yig'ish
\end{enumerate}
\begin{tcolorbox}[enhanced,breakable,colback=teal!4,colframe=teal!50!black,boxrule=0.5pt,arc=1.1mm,left=3mm,right=3mm,top=1.7mm,bottom=1.7mm,before skip=2mm,after skip=2mm]
\textbf{\color{teal}To'g'ri javob:} \textbf{A) Netiket}
\par\smallskip
\textbf{Izoh:} Vaziyat Netiket tushunchasini ifodalaydi: raqamli muloqotda hurmat, aniqlik va xavfsizlikka rioya qilish qoidalari.
\end{tcolorbox}
\vspace{2mm}\hrule height 0.2pt \vspace{3mm}

% DB question_id: 619f9e90-f1a6-573c-872c-41e0b18845a1
\needspace{5\baselineskip}
\noindent\textbf{\color{ink}287.}
\par
"Qayta ishlash va tahlil" atamasining eng aniq ta'rifini toping.
\begin{enumerate}[label=\Alph*),leftmargin=9mm]
\item kerakli axborotni manbalar yoki ma'lumotlar to'plamidan topish jarayoni
\item ma'lumotdan yangi natija olish, undagi bog'lanish va xulosalarni aniqlash jarayoni
\item elementlarni tanlangan mezon bo'yicha tartibga keltirish jarayoni
\item faqat berilgan shartga mos yozuvlarni ajratib ko'rsatish jarayoni
\end{enumerate}
\begin{tcolorbox}[enhanced,breakable,colback=teal!4,colframe=teal!50!black,boxrule=0.5pt,arc=1.1mm,left=3mm,right=3mm,top=1.7mm,bottom=1.7mm,before skip=2mm,after skip=2mm]
\textbf{\color{teal}To'g'ri javob:} \textbf{B) ma'lumotdan yangi natija olish, undagi bog'lanish va xulosalarni aniqlash jarayoni}
\par\smallskip
\textbf{Izoh:} Qayta ishlash va tahlil - ma'lumotdan yangi natija olish, undagi bog'lanish va xulosalarni aniqlash jarayoni.
\end{tcolorbox}
\vspace{2mm}\hrule height 0.2pt \vspace{3mm}

% DB question_id: 61a4f4d6-6faa-51ab-9e9c-9b3aab4baaa8
\needspace{5\baselineskip}
\noindent\textbf{\color{ink}288.}
\par
Quyidagi ta'rif qaysi tushunchaga tegishli? Elementlarni tanlangan mezon bo'yicha tartibga keltirish jarayoni.
\begin{enumerate}[label=\Alph*),leftmargin=9mm]
\item Filtrlash
\item Guruhlash
\item Saralash
\item Uzatish va qabul qilish
\end{enumerate}
\begin{tcolorbox}[enhanced,breakable,colback=teal!4,colframe=teal!50!black,boxrule=0.5pt,arc=1.1mm,left=3mm,right=3mm,top=1.7mm,bottom=1.7mm,before skip=2mm,after skip=2mm]
\textbf{\color{teal}To'g'ri javob:} \textbf{C) Saralash}
\par\smallskip
\textbf{Izoh:} Ta'rif Saralash tushunchasiga tegishli.
\end{tcolorbox}
\vspace{2mm}\hrule height 0.2pt \vspace{3mm}

% DB question_id: 66d27c1b-39ce-55e4-9fbe-0708d94c4e06
\needspace{5\baselineskip}
\noindent\textbf{\color{ink}289.}
\par
Quyidagi ta'rif qaysi tushunchaga tegishli? Muallif, sana, dalil, maqsad va boshqa manbalar bilan moslikni tekshirish jarayoni.
\begin{enumerate}[label=\Alph*),leftmargin=9mm]
\item Auditoriya, maqsad va kontekst
\item Elektron pochta madaniyati
\item Qidiruv natijalarini baholash
\item Netiket
\end{enumerate}
\begin{tcolorbox}[enhanced,breakable,colback=teal!4,colframe=teal!50!black,boxrule=0.5pt,arc=1.1mm,left=3mm,right=3mm,top=1.7mm,bottom=1.7mm,before skip=2mm,after skip=2mm]
\textbf{\color{teal}To'g'ri javob:} \textbf{C) Qidiruv natijalarini baholash}
\par\smallskip
\textbf{Izoh:} Ta'rif Qidiruv natijalarini baholash tushunchasiga tegishli.
\end{tcolorbox}
\vspace{2mm}\hrule height 0.2pt \vspace{3mm}

% DB question_id: 6e261f6e-a931-42c3-8dfc-61186b069ade
\needspace{5\baselineskip}
\noindent\textbf{\color{ink}290.}
\par
Fotosurat muallifi nusxa olish va tarqatish shartlarini belgiladi. Bu vaziyatdagi asosiy tushunchani aniqlang.
\begin{enumerate}[label=\Alph*),leftmargin=9mm]
\item Mualliflik huquqi
\item Plagiat
\item Litsenziya
\item Ommaviy mulk
\end{enumerate}
\begin{tcolorbox}[enhanced,breakable,colback=teal!4,colframe=teal!50!black,boxrule=0.5pt,arc=1.1mm,left=3mm,right=3mm,top=1.7mm,bottom=1.7mm,before skip=2mm,after skip=2mm]
\textbf{\color{teal}To'g'ri javob:} \textbf{A) Mualliflik huquqi}
\par\smallskip
\textbf{Izoh:} Vaziyat Mualliflik huquqiga misol: ijodiy asarning ifoda shakli va undan foydalanish huquqlarini himoya qiladigan huquqiy tizim.
\end{tcolorbox}
\vspace{2mm}\hrule height 0.2pt \vspace{3mm}

% DB question_id: 718c0128-12dc-52f6-aa65-d2728a442aac
\needspace{5\baselineskip}
\noindent\textbf{\color{ink}291.}
\par
Quyidagi ta'rif qaysi tushunchaga tegishli? Keraksiz yoki maxfiy axborotni belgilangan tartibda qayta tiklab bo'lmaydigan darajada o'chirish jarayoni.
\begin{enumerate}[label=\Alph*),leftmargin=9mm]
\item Qidiruv tizimi va kalit so'z
\item Aniq ibora bilan qidirish
\item Yo'q qilish
\item Mantiqiy qidiruv operatorlari
\end{enumerate}
\begin{tcolorbox}[enhanced,breakable,colback=teal!4,colframe=teal!50!black,boxrule=0.5pt,arc=1.1mm,left=3mm,right=3mm,top=1.7mm,bottom=1.7mm,before skip=2mm,after skip=2mm]
\textbf{\color{teal}To'g'ri javob:} \textbf{C) Yo'q qilish}
\par\smallskip
\textbf{Izoh:} Ta'rif Yo'q qilish tushunchasiga tegishli.
\end{tcolorbox}
\vspace{2mm}\hrule height 0.2pt \vspace{3mm}

% DB question_id: 74a49843-5b3b-5cef-9a89-7dd0a8a21ca2
\needspace{5\baselineskip}
\noindent\textbf{\color{ink}292.}
\par
"Izlash" atamasining eng aniq ta'rifini toping.
\begin{enumerate}[label=\Alph*),leftmargin=9mm]
\item elementlarni tanlangan mezon bo'yicha tartibga keltirish jarayoni
\item faqat berilgan shartga mos yozuvlarni ajratib ko'rsatish jarayoni
\item kerakli axborotni manbalar yoki ma'lumotlar to'plamidan topish jarayoni
\item umumiy belgi bo'yicha elementlarni toifalarga ajratish jarayoni
\end{enumerate}
\begin{tcolorbox}[enhanced,breakable,colback=teal!4,colframe=teal!50!black,boxrule=0.5pt,arc=1.1mm,left=3mm,right=3mm,top=1.7mm,bottom=1.7mm,before skip=2mm,after skip=2mm]
\textbf{\color{teal}To'g'ri javob:} \textbf{C) kerakli axborotni manbalar yoki ma'lumotlar to'plamidan topish jarayoni}
\par\smallskip
\textbf{Izoh:} Izlash - kerakli axborotni manbalar yoki ma'lumotlar to'plamidan topish jarayoni.
\end{tcolorbox}
\vspace{2mm}\hrule height 0.2pt \vspace{3mm}

% DB question_id: 768a1ce6-5bfd-53e2-9168-d63a095e4583
\needspace{5\baselineskip}
\noindent\textbf{\color{ink}293.}
\par
"Taqdim etish" atamasining eng aniq ta'rifini toping.
\begin{enumerate}[label=\Alph*),leftmargin=9mm]
\item axborotni ruxsatsiz kirish, o'zgartirish yoki yo'qolishdan saqlash choralarini qo'llash jarayoni
\item axborotni auditoriya va maqsadga mos shaklda ko'rsatish jarayoni
\item keraksiz yoki maxfiy axborotni belgilangan tartibda qayta tiklab bo'lmaydigan darajada o'chirish jarayoni
\item veb resurslarini indekslab, foydalanuvchi yozgan muhim so'zlarga mos natijalarni topib beruvchi xizmat va so'rov elementi
\end{enumerate}
\begin{tcolorbox}[enhanced,breakable,colback=teal!4,colframe=teal!50!black,boxrule=0.5pt,arc=1.1mm,left=3mm,right=3mm,top=1.7mm,bottom=1.7mm,before skip=2mm,after skip=2mm]
\textbf{\color{teal}To'g'ri javob:} \textbf{B) axborotni auditoriya va maqsadga mos shaklda ko'rsatish jarayoni}
\par\smallskip
\textbf{Izoh:} Taqdim etish - axborotni auditoriya va maqsadga mos shaklda ko'rsatish jarayoni.
\end{tcolorbox}
\vspace{2mm}\hrule height 0.2pt \vspace{3mm}

% DB question_id: 77ab7f62-813b-4bfe-935f-24527b04355f
\needspace{5\baselineskip}
\noindent\textbf{\color{ink}294.}
\par
Elektron xat jo'natuvchidan pochta serveriga, u yerdan qabul qiluvchining qutisiga yetib bordi. Bu jarayonning asosiy axborot amallari qaysi?
\begin{enumerate}[label=\Alph*),leftmargin=9mm]
\item Uzatish va qabul qilish
\item Yaratish va tahrirlash jarayoni
\item Shifrlash va dekodlash jarayoni
\item Qidirish va saqlash jarayoni
\end{enumerate}
\begin{tcolorbox}[enhanced,breakable,colback=teal!4,colframe=teal!50!black,boxrule=0.5pt,arc=1.1mm,left=3mm,right=3mm,top=1.7mm,bottom=1.7mm,before skip=2mm,after skip=2mm]
\textbf{\color{teal}To'g'ri javob:} \textbf{A) Uzatish va qabul qilish}
\par\smallskip
\textbf{Izoh:} Jo'natuvchi xabarni tarmoq orqali uzatadi, qabul qiluvchi esa uni oladi. Shifrlash, qidirish yoki tahrirlash xatning har safar jo'natilishiga xos majburiy bosqich emas.
\end{tcolorbox}
\vspace{2mm}\hrule height 0.2pt \vspace{3mm}

% DB question_id: 77eb3ca4-8439-583d-bb78-d5136d1624d7
\needspace{5\baselineskip}
\noindent\textbf{\color{ink}295.}
\par
E'lon barcha sinf guruhlariga yuborildi. Ushbu vaziyatdagi asosiy tushunchani aniqlang.
\begin{enumerate}[label=\Alph*),leftmargin=9mm]
\item Taqdim etish
\item Himoyalash
\item Tarqatish
\item Yo'q qilish
\end{enumerate}
\begin{tcolorbox}[enhanced,breakable,colback=teal!4,colframe=teal!50!black,boxrule=0.5pt,arc=1.1mm,left=3mm,right=3mm,top=1.7mm,bottom=1.7mm,before skip=2mm,after skip=2mm]
\textbf{\color{teal}To'g'ri javob:} \textbf{C) Tarqatish}
\par\smallskip
\textbf{Izoh:} Vaziyat Tarqatish tushunchasini ifodalaydi: axborotni bir nechta foydalanuvchi yoki auditoriyaga yetkazish jarayoni.
\end{tcolorbox}
\vspace{2mm}\hrule height 0.2pt \vspace{3mm}

% DB question_id: 78a96e0c-9fe7-5721-8d96-c3a8ddb9c9b2
\needspace{5\baselineskip}
\noindent\textbf{\color{ink}296.}
\par
Quyidagi ta'rif qaysi tushunchaga tegishli? Xat mavzusini aniq yozish, hurmatli uslub, to'g'ri adreslash va ilovani tekshirish qoidalari.
\begin{enumerate}[label=\Alph*),leftmargin=9mm]
\item Elektron pochta madaniyati
\item Netiket
\item Raqamli iz
\item Axborotni taqdim etish shaklini tanlash
\end{enumerate}
\begin{tcolorbox}[enhanced,breakable,colback=teal!4,colframe=teal!50!black,boxrule=0.5pt,arc=1.1mm,left=3mm,right=3mm,top=1.7mm,bottom=1.7mm,before skip=2mm,after skip=2mm]
\textbf{\color{teal}To'g'ri javob:} \textbf{A) Elektron pochta madaniyati}
\par\smallskip
\textbf{Izoh:} Ta'rif Elektron pochta madaniyati tushunchasiga tegishli.
\end{tcolorbox}
\vspace{2mm}\hrule height 0.2pt \vspace{3mm}

% DB question_id: 7b2787df-4d54-5cf7-8eba-f2dcc0fde87f
\needspace{5\baselineskip}
\noindent\textbf{\color{ink}297.}
\par
Hujjatga kirish huquqi va kuchli parol o'rnatildi. Ushbu vaziyatdagi asosiy tushunchani aniqlang.
\begin{enumerate}[label=\Alph*),leftmargin=9mm]
\item Himoyalash
\item Yo'q qilish
\item Qidiruv tizimi va kalit so'z
\item Aniq ibora bilan qidirish
\end{enumerate}
\begin{tcolorbox}[enhanced,breakable,colback=teal!4,colframe=teal!50!black,boxrule=0.5pt,arc=1.1mm,left=3mm,right=3mm,top=1.7mm,bottom=1.7mm,before skip=2mm,after skip=2mm]
\textbf{\color{teal}To'g'ri javob:} \textbf{A) Himoyalash}
\par\smallskip
\textbf{Izoh:} Vaziyat Himoyalash tushunchasini ifodalaydi: axborotni ruxsatsiz kirish, o'zgartirish yoki yo'qolishdan saqlash choralarini qo'llash jarayoni.
\end{tcolorbox}
\vspace{2mm}\hrule height 0.2pt \vspace{3mm}

% DB question_id: 7b430e5b-007f-5051-ad61-d69dd5f7992a
\needspace{5\baselineskip}
\noindent\textbf{\color{ink}298.}
\par
Quyidagi ta'rif qaysi tushunchaga tegishli? Yangi mazmun hosil qilish hamda kerakli ma'lumotlarni turli manbalardan to'plash jarayoni.
\begin{enumerate}[label=\Alph*),leftmargin=9mm]
\item Qabul qilish va kiritish
\item Axborotni yaratish va yig'ish
\item Saqlash va arxivlash
\item Qayta ishlash va tahlil
\end{enumerate}
\begin{tcolorbox}[enhanced,breakable,colback=teal!4,colframe=teal!50!black,boxrule=0.5pt,arc=1.1mm,left=3mm,right=3mm,top=1.7mm,bottom=1.7mm,before skip=2mm,after skip=2mm]
\textbf{\color{teal}To'g'ri javob:} \textbf{B) Axborotni yaratish va yig'ish}
\par\smallskip
\textbf{Izoh:} Ta'rif Axborotni yaratish va yig'ish tushunchasiga tegishli.
\end{tcolorbox}
\vspace{2mm}\hrule height 0.2pt \vspace{3mm}

% DB question_id: 88fcc37c-88fa-5973-a644-e80cb1183129
\needspace{5\baselineskip}
\noindent\textbf{\color{ink}299.}
\par
Natijada ikki mavzu birga qatnashishi AND orqali talab qilindi. Ushbu vaziyatdagi asosiy tushunchani aniqlang.
\begin{enumerate}[label=\Alph*),leftmargin=9mm]
\item Mantiqiy qidiruv operatorlari
\item Qidiruv natijalarini baholash
\item Auditoriya, maqsad va kontekst
\item Elektron pochta madaniyati
\end{enumerate}
\begin{tcolorbox}[enhanced,breakable,colback=teal!4,colframe=teal!50!black,boxrule=0.5pt,arc=1.1mm,left=3mm,right=3mm,top=1.7mm,bottom=1.7mm,before skip=2mm,after skip=2mm]
\textbf{\color{teal}To'g'ri javob:} \textbf{A) Mantiqiy qidiruv operatorlari}
\par\smallskip
\textbf{Izoh:} Vaziyat Mantiqiy qidiruv operatorlari tushunchasini ifodalaydi: qidiruv shartlarini AND, OR va NOT kabi operatorlar bilan aniqlashtirish usuli.
\end{tcolorbox}
\vspace{2mm}\hrule height 0.2pt \vspace{3mm}

% DB question_id: 8f5c32fa-344b-5f76-9e3b-3971e4acb034
\needspace{5\baselineskip}
\noindent\textbf{\color{ink}300.}
\par
"Elektron pochta madaniyati" atamasining eng aniq ta'rifini toping.
\begin{enumerate}[label=\Alph*),leftmargin=9mm]
\item raqamli muloqotda hurmat, aniqlik va xavfsizlikka rioya qilish qoidalari
\item xat mavzusini aniq yozish, hurmatli uslub, to'g'ri adreslash va ilovani tekshirish qoidalari
\item onlayn faoliyat natijasida xizmatlar va qurilmalarda qoladigan ma'lumotlar majmui
\item mazmunni auditoriya ehtiyoji va maqsadga eng mos matn, jadval, diagramma, audio yoki video shaklida berish qarori
\end{enumerate}
\begin{tcolorbox}[enhanced,breakable,colback=teal!4,colframe=teal!50!black,boxrule=0.5pt,arc=1.1mm,left=3mm,right=3mm,top=1.7mm,bottom=1.7mm,before skip=2mm,after skip=2mm]
\textbf{\color{teal}To'g'ri javob:} \textbf{B) xat mavzusini aniq yozish, hurmatli uslub, to'g'ri adreslash va ilovani tekshirish qoidalari}
\par\smallskip
\textbf{Izoh:} Elektron pochta madaniyati - xat mavzusini aniq yozish, hurmatli uslub, to'g'ri adreslash va ilovani tekshirish qoidalari.
\end{tcolorbox}
\vspace{2mm}\hrule height 0.2pt \vspace{3mm}

% DB question_id: 93830872-49b4-5c8f-ae7a-543ceadc47f5
\needspace{5\baselineskip}
\noindent\textbf{\color{ink}301.}
\par
80 balldan yuqori natijalar ekranda qoldirildi. Ushbu vaziyatdagi asosiy tushunchani aniqlang.
\begin{enumerate}[label=\Alph*),leftmargin=9mm]
\item Guruhlash
\item Uzatish va qabul qilish
\item Filtrlash
\item Nusxalash
\end{enumerate}
\begin{tcolorbox}[enhanced,breakable,colback=teal!4,colframe=teal!50!black,boxrule=0.5pt,arc=1.1mm,left=3mm,right=3mm,top=1.7mm,bottom=1.7mm,before skip=2mm,after skip=2mm]
\textbf{\color{teal}To'g'ri javob:} \textbf{C) Filtrlash}
\par\smallskip
\textbf{Izoh:} Vaziyat Filtrlash tushunchasini ifodalaydi: faqat berilgan shartga mos yozuvlarni ajratib ko'rsatish jarayoni.
\end{tcolorbox}
\vspace{2mm}\hrule height 0.2pt \vspace{3mm}

% DB question_id: 94e3e26a-4ca1-5164-8ca2-2ec0db6c0483
\needspace{5\baselineskip}
\noindent\textbf{\color{ink}302.}
\par
"Axborotni taqdim etish shaklini tanlash" atamasining eng aniq ta'rifini toping.
\begin{enumerate}[label=\Alph*),leftmargin=9mm]
\item mazmunni auditoriya ehtiyoji va maqsadga eng mos matn, jadval, diagramma, audio yoki video shaklida berish qarori
\item yangi mazmun hosil qilish hamda kerakli ma'lumotlarni turli manbalardan to'plash jarayoni
\item axborotni sezgi yoki qurilma orqali olish va uni axborot tizimiga kiritish jarayoni
\item axborotni keyinchalik foydalanish uchun saqlab qolish hamda uzoq muddat tartibli jamlash jarayoni
\end{enumerate}
\begin{tcolorbox}[enhanced,breakable,colback=teal!4,colframe=teal!50!black,boxrule=0.5pt,arc=1.1mm,left=3mm,right=3mm,top=1.7mm,bottom=1.7mm,before skip=2mm,after skip=2mm]
\textbf{\color{teal}To'g'ri javob:} \textbf{A) mazmunni auditoriya ehtiyoji va maqsadga eng mos matn, jadval, diagramma, audio yoki video shaklida berish qarori}
\par\smallskip
\textbf{Izoh:} Axborotni taqdim etish shaklini tanlash - mazmunni auditoriya ehtiyoji va maqsadga eng mos matn, jadval, diagramma, audio yoki video shaklida berish qarori.
\end{tcolorbox}
\vspace{2mm}\hrule height 0.2pt \vspace{3mm}

% DB question_id: 96655822-2825-5b68-85c4-25d355041fbe
\needspace{5\baselineskip}
\noindent\textbf{\color{ink}303.}
\par
Eski maxfiy tashuvchi xavfsiz tozalandi. Ushbu vaziyatdagi asosiy tushunchani aniqlang.
\begin{enumerate}[label=\Alph*),leftmargin=9mm]
\item Qidiruv tizimi va kalit so'z
\item Yo'q qilish
\item Aniq ibora bilan qidirish
\item Mantiqiy qidiruv operatorlari
\end{enumerate}
\begin{tcolorbox}[enhanced,breakable,colback=teal!4,colframe=teal!50!black,boxrule=0.5pt,arc=1.1mm,left=3mm,right=3mm,top=1.7mm,bottom=1.7mm,before skip=2mm,after skip=2mm]
\textbf{\color{teal}To'g'ri javob:} \textbf{B) Yo'q qilish}
\par\smallskip
\textbf{Izoh:} Vaziyat Yo'q qilish tushunchasini ifodalaydi: keraksiz yoki maxfiy axborotni belgilangan tartibda qayta tiklab bo'lmaydigan darajada o'chirish jarayoni.
\end{tcolorbox}
\vspace{2mm}\hrule height 0.2pt \vspace{3mm}

% DB question_id: 9b916ae7-2c82-5a4e-bc2e-5f59f0b6779f
\needspace{5\baselineskip}
\noindent\textbf{\color{ink}304.}
\par
Ko'p sonli qiymatlar taqqoslash uchun jadval emas, ustunli diagrammada ko'rsatildi. Ushbu vaziyatdagi asosiy tushunchani aniqlang.
\begin{enumerate}[label=\Alph*),leftmargin=9mm]
\item Axborotni yaratish va yig'ish
\item Qabul qilish va kiritish
\item Axborotni taqdim etish shaklini tanlash
\item Saqlash va arxivlash
\end{enumerate}
\begin{tcolorbox}[enhanced,breakable,colback=teal!4,colframe=teal!50!black,boxrule=0.5pt,arc=1.1mm,left=3mm,right=3mm,top=1.7mm,bottom=1.7mm,before skip=2mm,after skip=2mm]
\textbf{\color{teal}To'g'ri javob:} \textbf{C) Axborotni taqdim etish shaklini tanlash}
\par\smallskip
\textbf{Izoh:} Vaziyat Axborotni taqdim etish shaklini tanlash tushunchasini ifodalaydi: mazmunni auditoriya ehtiyoji va maqsadga eng mos matn, jadval, diagramma, audio yoki video shaklida berish qarori.
\end{tcolorbox}
\vspace{2mm}\hrule height 0.2pt \vspace{3mm}

% DB question_id: 9ddbde56-eed2-5a87-9efb-1caf415c125c
\needspace{5\baselineskip}
\noindent\textbf{\color{ink}305.}
\par
Qidiruvlar, postlar va xaridlar tarixi foydalanuvchi faoliyati haqida yozuv qoldirdi. Ushbu vaziyatdagi asosiy tushunchani aniqlang.
\begin{enumerate}[label=\Alph*),leftmargin=9mm]
\item Axborotni taqdim etish shaklini tanlash
\item Raqamli iz
\item Axborotni yaratish va yig'ish
\item Qabul qilish va kiritish
\end{enumerate}
\begin{tcolorbox}[enhanced,breakable,colback=teal!4,colframe=teal!50!black,boxrule=0.5pt,arc=1.1mm,left=3mm,right=3mm,top=1.7mm,bottom=1.7mm,before skip=2mm,after skip=2mm]
\textbf{\color{teal}To'g'ri javob:} \textbf{B) Raqamli iz}
\par\smallskip
\textbf{Izoh:} Vaziyat Raqamli iz tushunchasini ifodalaydi: onlayn faoliyat natijasida xizmatlar va qurilmalarda qoladigan ma'lumotlar majmui.
\end{tcolorbox}
\vspace{2mm}\hrule height 0.2pt \vspace{3mm}

% DB question_id: a0a03864-3f39-5bde-8651-2366cbed7fc2
\needspace{5\baselineskip}
\noindent\textbf{\color{ink}306.}
\par
Quyidagi ta'rif qaysi tushunchaga tegishli? Mazmunni auditoriya ehtiyoji va maqsadga eng mos matn, jadval, diagramma, audio yoki video shaklida berish qarori.
\begin{enumerate}[label=\Alph*),leftmargin=9mm]
\item Axborotni yaratish va yig'ish
\item Qabul qilish va kiritish
\item Saqlash va arxivlash
\item Axborotni taqdim etish shaklini tanlash
\end{enumerate}
\begin{tcolorbox}[enhanced,breakable,colback=teal!4,colframe=teal!50!black,boxrule=0.5pt,arc=1.1mm,left=3mm,right=3mm,top=1.7mm,bottom=1.7mm,before skip=2mm,after skip=2mm]
\textbf{\color{teal}To'g'ri javob:} \textbf{D) Axborotni taqdim etish shaklini tanlash}
\par\smallskip
\textbf{Izoh:} Ta'rif Axborotni taqdim etish shaklini tanlash tushunchasiga tegishli.
\end{tcolorbox}
\vspace{2mm}\hrule height 0.2pt \vspace{3mm}

% DB question_id: a1450cad-83a3-565c-8297-ad4e4d859d1e
\needspace{5\baselineskip}
\noindent\textbf{\color{ink}307.}
\par
Ro'yxat familiya bo'yicha A dan Z gacha joylashtirildi. Ushbu vaziyatdagi asosiy tushunchani aniqlang.
\begin{enumerate}[label=\Alph*),leftmargin=9mm]
\item Filtrlash
\item Saralash
\item Guruhlash
\item Uzatish va qabul qilish
\end{enumerate}
\begin{tcolorbox}[enhanced,breakable,colback=teal!4,colframe=teal!50!black,boxrule=0.5pt,arc=1.1mm,left=3mm,right=3mm,top=1.7mm,bottom=1.7mm,before skip=2mm,after skip=2mm]
\textbf{\color{teal}To'g'ri javob:} \textbf{B) Saralash}
\par\smallskip
\textbf{Izoh:} Vaziyat Saralash tushunchasini ifodalaydi: elementlarni tanlangan mezon bo'yicha tartibga keltirish jarayoni.
\end{tcolorbox}
\vspace{2mm}\hrule height 0.2pt \vspace{3mm}

% DB question_id: a4a4140c-3970-4da5-ad2e-562839aa8df4
\needspace{5\baselineskip}
\noindent\textbf{\color{ink}308.}
\par
Davlat, ta'lim va biznes xizmatlarining katta qismi elektron shaklga o'tdi. Bu vaziyatdagi asosiy tushunchani aniqlang.
\begin{enumerate}[label=\Alph*),leftmargin=9mm]
\item Axborot madaniyati
\item Raqamli savodxonlik
\item Axborotlashgan jamiyat
\item Raqamli iz
\end{enumerate}
\begin{tcolorbox}[enhanced,breakable,colback=teal!4,colframe=teal!50!black,boxrule=0.5pt,arc=1.1mm,left=3mm,right=3mm,top=1.7mm,bottom=1.7mm,before skip=2mm,after skip=2mm]
\textbf{\color{teal}To'g'ri javob:} \textbf{C) Axborotlashgan jamiyat}
\par\smallskip
\textbf{Izoh:} Vaziyat Axborotlashgan jamiyatga misol: axborot va bilim ishlab chiqarish hamda xizmatlarda asosiy resursga aylangan jamiyat.
\end{tcolorbox}
\vspace{2mm}\hrule height 0.2pt \vspace{3mm}

% DB question_id: a7f57ab8-006b-4bd8-9f0c-bcbae15a89d5
\needspace{5\baselineskip}
\noindent\textbf{\color{ink}309.}
\par
O'quvchi elektron faylni pochta orqali o'qituvchiga yubordi. Xabar serverlardan o'tib, o'qituvchining qutisiga yetdi. Axborot jarayonining qaysi jufti markaziy?
\begin{enumerate}[label=\Alph*),leftmargin=9mm]
\item Ma'lumotni yaratish va saqlash
\item Uzatish va qabul qilish
\item Ma'lumotni qidirish va saralash
\item Ma'lumotni tahrirlash va chop etish
\end{enumerate}
\begin{tcolorbox}[enhanced,breakable,colback=teal!4,colframe=teal!50!black,boxrule=0.5pt,arc=1.1mm,left=3mm,right=3mm,top=1.7mm,bottom=1.7mm,before skip=2mm,after skip=2mm]
\textbf{\color{teal}To'g'ri javob:} \textbf{B) Uzatish va qabul qilish}
\par\smallskip
\textbf{Izoh:} Fayl aloqa kanali orqali jo'natilib, qabul qiluvchiga yetkazilishi - axborotni uzatish va qabul qilishdir. Yaratish, saqlash yoki chop etish alohida amallardir.
\end{tcolorbox}
\vspace{2mm}\hrule height 0.2pt \vspace{3mm}

% DB question_id: b5649dc6-4a2b-56cc-8bfb-059f7d7e951a
\needspace{5\baselineskip}
\noindent\textbf{\color{ink}310.}
\par
Quyidagi ta'rif qaysi tushunchaga tegishli? Axborotni bir nechta foydalanuvchi yoki auditoriyaga yetkazish jarayoni.
\begin{enumerate}[label=\Alph*),leftmargin=9mm]
\item Taqdim etish
\item Himoyalash
\item Yo'q qilish
\item Tarqatish
\end{enumerate}
\begin{tcolorbox}[enhanced,breakable,colback=teal!4,colframe=teal!50!black,boxrule=0.5pt,arc=1.1mm,left=3mm,right=3mm,top=1.7mm,bottom=1.7mm,before skip=2mm,after skip=2mm]
\textbf{\color{teal}To'g'ri javob:} \textbf{D) Tarqatish}
\par\smallskip
\textbf{Izoh:} Ta'rif Tarqatish tushunchasiga tegishli.
\end{tcolorbox}
\vspace{2mm}\hrule height 0.2pt \vspace{3mm}

% DB question_id: b9c1ff61-f69b-569a-8752-0a15d4bbb281
\needspace{5\baselineskip}
\noindent\textbf{\color{ink}311.}
\par
O'quvchi tajriba haqida hisobot yozdi va so'rovnoma javoblarini unga jamladi. Ushbu vaziyatdagi asosiy tushunchani aniqlang.
\begin{enumerate}[label=\Alph*),leftmargin=9mm]
\item Axborotni yaratish va yig'ish
\item Qabul qilish va kiritish
\item Saqlash va arxivlash
\item Qayta ishlash va tahlil
\end{enumerate}
\begin{tcolorbox}[enhanced,breakable,colback=teal!4,colframe=teal!50!black,boxrule=0.5pt,arc=1.1mm,left=3mm,right=3mm,top=1.7mm,bottom=1.7mm,before skip=2mm,after skip=2mm]
\textbf{\color{teal}To'g'ri javob:} \textbf{A) Axborotni yaratish va yig'ish}
\par\smallskip
\textbf{Izoh:} Vaziyat Axborotni yaratish va yig'ish tushunchasini ifodalaydi: yangi mazmun hosil qilish hamda kerakli ma'lumotlarni turli manbalardan to'plash jarayoni.
\end{tcolorbox}
\vspace{2mm}\hrule height 0.2pt \vspace{3mm}

% DB question_id: b9f1669c-71a2-5109-b991-2a0fe5b524ac
\needspace{5\baselineskip}
\noindent\textbf{\color{ink}312.}
\par
"Uzatish va qabul qilish" atamasining eng aniq ta'rifini toping.
\begin{enumerate}[label=\Alph*),leftmargin=9mm]
\item axborotning mazmunini saqlagan holda uning qo'shimcha nusxasini yaratish jarayoni
\item axborotni bir nechta foydalanuvchi yoki auditoriyaga yetkazish jarayoni
\item axborotni manbadan kanal orqali qabul qiluvchiga yetkazish va uni olish jarayoni
\item axborotni auditoriya va maqsadga mos shaklda ko'rsatish jarayoni
\end{enumerate}
\begin{tcolorbox}[enhanced,breakable,colback=teal!4,colframe=teal!50!black,boxrule=0.5pt,arc=1.1mm,left=3mm,right=3mm,top=1.7mm,bottom=1.7mm,before skip=2mm,after skip=2mm]
\textbf{\color{teal}To'g'ri javob:} \textbf{C) axborotni manbadan kanal orqali qabul qiluvchiga yetkazish va uni olish jarayoni}
\par\smallskip
\textbf{Izoh:} Uzatish va qabul qilish - axborotni manbadan kanal orqali qabul qiluvchiga yetkazish va uni olish jarayoni.
\end{tcolorbox}
\vspace{2mm}\hrule height 0.2pt \vspace{3mm}

% DB question_id: bd46862e-c9a1-595e-883f-f21266472023
\needspace{5\baselineskip}
\noindent\textbf{\color{ink}313.}
\par
Hisobotning asl fayli saqlanib, zaxira nusxasi yaratildi. Ushbu vaziyatdagi asosiy tushunchani aniqlang.
\begin{enumerate}[label=\Alph*),leftmargin=9mm]
\item Tarqatish
\item Nusxalash
\item Taqdim etish
\item Himoyalash
\end{enumerate}
\begin{tcolorbox}[enhanced,breakable,colback=teal!4,colframe=teal!50!black,boxrule=0.5pt,arc=1.1mm,left=3mm,right=3mm,top=1.7mm,bottom=1.7mm,before skip=2mm,after skip=2mm]
\textbf{\color{teal}To'g'ri javob:} \textbf{B) Nusxalash}
\par\smallskip
\textbf{Izoh:} Vaziyat Nusxalash tushunchasini ifodalaydi: axborotning mazmunini saqlagan holda uning qo'shimcha nusxasini yaratish jarayoni.
\end{tcolorbox}
\vspace{2mm}\hrule height 0.2pt \vspace{3mm}

% DB question_id: bfb7bf45-0dc6-51be-99d4-f5ab25bac0d4
\needspace{5\baselineskip}
\noindent\textbf{\color{ink}314.}
\par
Yil yakuni hujjatlari sana bo'yicha arxivga joylandi. Ushbu vaziyatdagi asosiy tushunchani aniqlang.
\begin{enumerate}[label=\Alph*),leftmargin=9mm]
\item Qayta ishlash va tahlil
\item Izlash
\item Saqlash va arxivlash
\item Saralash
\end{enumerate}
\begin{tcolorbox}[enhanced,breakable,colback=teal!4,colframe=teal!50!black,boxrule=0.5pt,arc=1.1mm,left=3mm,right=3mm,top=1.7mm,bottom=1.7mm,before skip=2mm,after skip=2mm]
\textbf{\color{teal}To'g'ri javob:} \textbf{C) Saqlash va arxivlash}
\par\smallskip
\textbf{Izoh:} Vaziyat Saqlash va arxivlash tushunchasini ifodalaydi: axborotni keyinchalik foydalanish uchun saqlab qolish hamda uzoq muddat tartibli jamlash jarayoni.
\end{tcolorbox}
\vspace{2mm}\hrule height 0.2pt \vspace{3mm}

% DB question_id: c2b18ca9-23c0-58d9-b6ca-8336f5433c96
\needspace{5\baselineskip}
\noindent\textbf{\color{ink}315.}
\par
Foydalanuvchi iborani qo'shtirnoq ichida yozdi. Ushbu vaziyatdagi asosiy tushunchani aniqlang.
\begin{enumerate}[label=\Alph*),leftmargin=9mm]
\item Mantiqiy qidiruv operatorlari
\item Qidiruv natijalarini baholash
\item Auditoriya, maqsad va kontekst
\item Aniq ibora bilan qidirish
\end{enumerate}
\begin{tcolorbox}[enhanced,breakable,colback=teal!4,colframe=teal!50!black,boxrule=0.5pt,arc=1.1mm,left=3mm,right=3mm,top=1.7mm,bottom=1.7mm,before skip=2mm,after skip=2mm]
\textbf{\color{teal}To'g'ri javob:} \textbf{D) Aniq ibora bilan qidirish}
\par\smallskip
\textbf{Izoh:} Vaziyat Aniq ibora bilan qidirish tushunchasini ifodalaydi: so'zlar aynan berilgan tartibda uchrashini talab qiladigan qidiruv usuli.
\end{tcolorbox}
\vspace{2mm}\hrule height 0.2pt \vspace{3mm}

% DB question_id: c3aa4e7a-cfb3-569b-9be4-e3a22f0b94aa
\needspace{5\baselineskip}
\noindent\textbf{\color{ink}316.}
\par
Quyidagi ta'rif qaysi tushunchaga tegishli? Onlayn faoliyat natijasida xizmatlar va qurilmalarda qoladigan ma'lumotlar majmui.
\begin{enumerate}[label=\Alph*),leftmargin=9mm]
\item Axborotni taqdim etish shaklini tanlash
\item Axborotni yaratish va yig'ish
\item Raqamli iz
\item Qabul qilish va kiritish
\end{enumerate}
\begin{tcolorbox}[enhanced,breakable,colback=teal!4,colframe=teal!50!black,boxrule=0.5pt,arc=1.1mm,left=3mm,right=3mm,top=1.7mm,bottom=1.7mm,before skip=2mm,after skip=2mm]
\textbf{\color{teal}To'g'ri javob:} \textbf{C) Raqamli iz}
\par\smallskip
\textbf{Izoh:} Ta'rif Raqamli iz tushunchasiga tegishli.
\end{tcolorbox}
\vspace{2mm}\hrule height 0.2pt \vspace{3mm}

% DB question_id: c3e0a137-d9c5-4a2d-8a1c-77b31c12f4d7
\needspace{5\baselineskip}
\noindent\textbf{\color{ink}317.}
\par
O'quvchi ma'lumotni tekshirib, manba ko'rsatib va shaxsiy ma'lumotni oshkor qilmay ulashdi. Bu vaziyatdagi asosiy tushunchani aniqlang.
\begin{enumerate}[label=\Alph*),leftmargin=9mm]
\item Raqamli savodxonlik
\item Raqamli iz
\item Fakt
\item Axborot madaniyati
\end{enumerate}
\begin{tcolorbox}[enhanced,breakable,colback=teal!4,colframe=teal!50!black,boxrule=0.5pt,arc=1.1mm,left=3mm,right=3mm,top=1.7mm,bottom=1.7mm,before skip=2mm,after skip=2mm]
\textbf{\color{teal}To'g'ri javob:} \textbf{D) Axborot madaniyati}
\par\smallskip
\textbf{Izoh:} Vaziyat Axborot madaniyatiga misol: axborotni izlash, baholash, qonuniy va xavfsiz ishlatish ko'nikmalari majmui.
\end{tcolorbox}
\vspace{2mm}\hrule height 0.2pt \vspace{3mm}

% DB question_id: ccf70852-61f7-4fe0-bb06-a79a66116a49
\needspace{5\baselineskip}
\noindent\textbf{\color{ink}318.}
\par
Elektron jurnal baholarni kiritish, hisoblash va ota-onaga ko'rsatishni birlashtirdi. Bu vaziyatdagi asosiy tushunchani aniqlang.
\begin{enumerate}[label=\Alph*),leftmargin=9mm]
\item Raqamli axborot
\item Raqamli muhit
\item Axborotlashgan jamiyat
\item Axborot texnologiyasi
\end{enumerate}
\begin{tcolorbox}[enhanced,breakable,colback=teal!4,colframe=teal!50!black,boxrule=0.5pt,arc=1.1mm,left=3mm,right=3mm,top=1.7mm,bottom=1.7mm,before skip=2mm,after skip=2mm]
\textbf{\color{teal}To'g'ri javob:} \textbf{D) Axborot texnologiyasi}
\par\smallskip
\textbf{Izoh:} Vaziyat Axborot texnologiyasiga misol: axborotni yig'ish, saqlash, qayta ishlash, uzatish va taqdim etish vosita-usullari majmui.
\end{tcolorbox}
\vspace{2mm}\hrule height 0.2pt \vspace{3mm}

% DB question_id: d09966e1-94ae-4b3d-8d8a-d21092596840
\needspace{5\baselineskip}
\noindent\textbf{\color{ink}319.}
\par
Sinflarning test ballari jadvalidan har bir sinf o'rtachasi hisoblandi, ikki oy natijalari taqqoslanib o'sish tendensiyasi aniqlandi. Asosiy jarayon qaysi?
\begin{enumerate}[label=\Alph*),leftmargin=9mm]
\item Qayta ishlash va tahlil
\item Ma'lumotni qidirish va filtrlash
\item Ma'lumotni saqlash va nusxalash
\item Ma'lumotni saralash va guruhlash
\end{enumerate}
\begin{tcolorbox}[enhanced,breakable,colback=teal!4,colframe=teal!50!black,boxrule=0.5pt,arc=1.1mm,left=3mm,right=3mm,top=1.7mm,bottom=1.7mm,before skip=2mm,after skip=2mm]
\textbf{\color{teal}To'g'ri javob:} \textbf{A) Qayta ishlash va tahlil}
\par\smallskip
\textbf{Izoh:} O'rtachani hisoblash va ikki davr o'zgarishidan xulosa chiqarish dastlabki ma'lumotdan yangi natija olish - qayta ishlash va tahlilni ifodalaydi. Filtrlash yoki saralashning o'zi bunday xulosani bermaydi.
\end{tcolorbox}
\vspace{2mm}\hrule height 0.2pt \vspace{3mm}

% DB question_id: d1cccae9-dc67-513c-8323-4dd1f04f001e
\needspace{5\baselineskip}
\noindent\textbf{\color{ink}320.}
\par
Quyidagi ta'rif qaysi tushunchaga tegishli? Kerakli axborotni manbalar yoki ma'lumotlar to'plamidan topish jarayoni.
\begin{enumerate}[label=\Alph*),leftmargin=9mm]
\item Saralash
\item Izlash
\item Filtrlash
\item Guruhlash
\end{enumerate}
\begin{tcolorbox}[enhanced,breakable,colback=teal!4,colframe=teal!50!black,boxrule=0.5pt,arc=1.1mm,left=3mm,right=3mm,top=1.7mm,bottom=1.7mm,before skip=2mm,after skip=2mm]
\textbf{\color{teal}To'g'ri javob:} \textbf{B) Izlash}
\par\smallskip
\textbf{Izoh:} Ta'rif Izlash tushunchasiga tegishli.
\end{tcolorbox}
\vspace{2mm}\hrule height 0.2pt \vspace{3mm}

% DB question_id: d2d0968a-a4fe-4cac-9bd2-42e780f6f780
\needspace{5\baselineskip}
\noindent\textbf{\color{ink}321.}
\par
Ovoz ADC yordamida ikkilik kodga aylantirildi. Bu vaziyatdagi asosiy tushunchani aniqlang.
\begin{enumerate}[label=\Alph*),leftmargin=9mm]
\item Raqamli axborot
\item Raqamli muhit
\item Axborotlashgan jamiyat
\item Axborot madaniyati
\end{enumerate}
\begin{tcolorbox}[enhanced,breakable,colback=teal!4,colframe=teal!50!black,boxrule=0.5pt,arc=1.1mm,left=3mm,right=3mm,top=1.7mm,bottom=1.7mm,before skip=2mm,after skip=2mm]
\textbf{\color{teal}To'g'ri javob:} \textbf{A) Raqamli axborot}
\par\smallskip
\textbf{Izoh:} Vaziyat Raqamli axborotga misol: raqamli kod, odatda bitlar ketma-ketligi orqali ifodalangan axborot.
\end{tcolorbox}
\vspace{2mm}\hrule height 0.2pt \vspace{3mm}

% DB question_id: d2e83993-2d91-430a-a0a3-73710ff17e79
\needspace{5\baselineskip}
\noindent\textbf{\color{ink}322.}
\par
Dastur faqat ta'lim maqsadida ishlatish sharti bilan berildi. Bu vaziyatdagi asosiy tushunchani aniqlang.
\begin{enumerate}[label=\Alph*),leftmargin=9mm]
\item Ommaviy mulk
\item CC BY
\item Litsenziya
\item CC SA
\end{enumerate}
\begin{tcolorbox}[enhanced,breakable,colback=teal!4,colframe=teal!50!black,boxrule=0.5pt,arc=1.1mm,left=3mm,right=3mm,top=1.7mm,bottom=1.7mm,before skip=2mm,after skip=2mm]
\textbf{\color{teal}To'g'ri javob:} \textbf{C) Litsenziya}
\par\smallskip
\textbf{Izoh:} Vaziyat Litsenziyaga misol: asar yoki dasturdan qanday shartlarda foydalanish mumkinligini belgilovchi ruxsat.
\end{tcolorbox}
\vspace{2mm}\hrule height 0.2pt \vspace{3mm}

% DB question_id: d39f79fc-3987-547e-8e65-c5408f805ddd
\needspace{5\baselineskip}
\noindent\textbf{\color{ink}323.}
\par
"Auditoriya, maqsad va kontekst" atamasining eng aniq ta'rifini toping.
\begin{enumerate}[label=\Alph*),leftmargin=9mm]
\item axborotning kimga, nima uchun va qaysi sharoitda taqdim etilishini hisobga olish tamoyili
\item xat mavzusini aniq yozish, hurmatli uslub, to'g'ri adreslash va ilovani tekshirish qoidalari
\item raqamli muloqotda hurmat, aniqlik va xavfsizlikka rioya qilish qoidalari
\item onlayn faoliyat natijasida xizmatlar va qurilmalarda qoladigan ma'lumotlar majmui
\end{enumerate}
\begin{tcolorbox}[enhanced,breakable,colback=teal!4,colframe=teal!50!black,boxrule=0.5pt,arc=1.1mm,left=3mm,right=3mm,top=1.7mm,bottom=1.7mm,before skip=2mm,after skip=2mm]
\textbf{\color{teal}To'g'ri javob:} \textbf{A) axborotning kimga, nima uchun va qaysi sharoitda taqdim etilishini hisobga olish tamoyili}
\par\smallskip
\textbf{Izoh:} Auditoriya, maqsad va kontekst - axborotning kimga, nima uchun va qaysi sharoitda taqdim etilishini hisobga olish tamoyili.
\end{tcolorbox}
\vspace{2mm}\hrule height 0.2pt \vspace{3mm}

% DB question_id: d736da81-b352-5a7a-8dd8-432cc14be4f9
\needspace{5\baselineskip}
\noindent\textbf{\color{ink}324.}
\par
Quyidagi ta'rif qaysi tushunchaga tegishli? Faqat berilgan shartga mos yozuvlarni ajratib ko'rsatish jarayoni.
\begin{enumerate}[label=\Alph*),leftmargin=9mm]
\item Guruhlash
\item Uzatish va qabul qilish
\item Nusxalash
\item Filtrlash
\end{enumerate}
\begin{tcolorbox}[enhanced,breakable,colback=teal!4,colframe=teal!50!black,boxrule=0.5pt,arc=1.1mm,left=3mm,right=3mm,top=1.7mm,bottom=1.7mm,before skip=2mm,after skip=2mm]
\textbf{\color{teal}To'g'ri javob:} \textbf{D) Filtrlash}
\par\smallskip
\textbf{Izoh:} Ta'rif Filtrlash tushunchasiga tegishli.
\end{tcolorbox}
\vspace{2mm}\hrule height 0.2pt \vspace{3mm}

% DB question_id: d81ef3de-5280-5a39-9556-48ac57e2308f
\needspace{5\baselineskip}
\noindent\textbf{\color{ink}325.}
\par
Ballar o'rtachasi hisoblanib, eng katta o'sish kuzatilgan sinf topildi. Ushbu vaziyatdagi asosiy tushunchani aniqlang.
\begin{enumerate}[label=\Alph*),leftmargin=9mm]
\item Izlash
\item Saralash
\item Filtrlash
\item Qayta ishlash va tahlil
\end{enumerate}
\begin{tcolorbox}[enhanced,breakable,colback=teal!4,colframe=teal!50!black,boxrule=0.5pt,arc=1.1mm,left=3mm,right=3mm,top=1.7mm,bottom=1.7mm,before skip=2mm,after skip=2mm]
\textbf{\color{teal}To'g'ri javob:} \textbf{D) Qayta ishlash va tahlil}
\par\smallskip
\textbf{Izoh:} Vaziyat Qayta ishlash va tahlil tushunchasini ifodalaydi: ma'lumotdan yangi natija olish, undagi bog'lanish va xulosalarni aniqlash jarayoni.
\end{tcolorbox}
\vspace{2mm}\hrule height 0.2pt \vspace{3mm}

% DB question_id: e3e7b9d8-4be5-5bf3-abec-4c2ba895becd
\needspace{5\baselineskip}
\noindent\textbf{\color{ink}326.}
\par
Skaner qog'ozdagi matnni kompyuterga uzatdi. Ushbu vaziyatdagi asosiy tushunchani aniqlang.
\begin{enumerate}[label=\Alph*),leftmargin=9mm]
\item Saqlash va arxivlash
\item Qabul qilish va kiritish
\item Qayta ishlash va tahlil
\item Izlash
\end{enumerate}
\begin{tcolorbox}[enhanced,breakable,colback=teal!4,colframe=teal!50!black,boxrule=0.5pt,arc=1.1mm,left=3mm,right=3mm,top=1.7mm,bottom=1.7mm,before skip=2mm,after skip=2mm]
\textbf{\color{teal}To'g'ri javob:} \textbf{B) Qabul qilish va kiritish}
\par\smallskip
\textbf{Izoh:} Vaziyat Qabul qilish va kiritish tushunchasini ifodalaydi: axborotni sezgi yoki qurilma orqali olish va uni axborot tizimiga kiritish jarayoni.
\end{tcolorbox}
\vspace{2mm}\hrule height 0.2pt \vspace{3mm}

% DB question_id: e49d19a8-7c94-5622-bed6-fc8a6bd93e14
\needspace{5\baselineskip}
\noindent\textbf{\color{ink}327.}
\par
"Aniq ibora bilan qidirish" atamasining eng aniq ta'rifini toping.
\begin{enumerate}[label=\Alph*),leftmargin=9mm]
\item qidiruv shartlarini AND, OR va NOT kabi operatorlar bilan aniqlashtirish usuli
\item so'zlar aynan berilgan tartibda uchrashini talab qiladigan qidiruv usuli
\item muallif, sana, dalil, maqsad va boshqa manbalar bilan moslikni tekshirish jarayoni
\item axborotning kimga, nima uchun va qaysi sharoitda taqdim etilishini hisobga olish tamoyili
\end{enumerate}
\begin{tcolorbox}[enhanced,breakable,colback=teal!4,colframe=teal!50!black,boxrule=0.5pt,arc=1.1mm,left=3mm,right=3mm,top=1.7mm,bottom=1.7mm,before skip=2mm,after skip=2mm]
\textbf{\color{teal}To'g'ri javob:} \textbf{B) so'zlar aynan berilgan tartibda uchrashini talab qiladigan qidiruv usuli}
\par\smallskip
\textbf{Izoh:} Aniq ibora bilan qidirish - so'zlar aynan berilgan tartibda uchrashini talab qiladigan qidiruv usuli.
\end{tcolorbox}
\vspace{2mm}\hrule height 0.2pt \vspace{3mm}

% DB question_id: e7f1ebf0-151a-53f7-8dca-963bbad45332
\needspace{5\baselineskip}
\noindent\textbf{\color{ink}328.}
\par
"Saqlash va arxivlash" atamasining eng aniq ta'rifini toping.
\begin{enumerate}[label=\Alph*),leftmargin=9mm]
\item axborotni keyinchalik foydalanish uchun saqlab qolish hamda uzoq muddat tartibli jamlash jarayoni
\item ma'lumotdan yangi natija olish, undagi bog'lanish va xulosalarni aniqlash jarayoni
\item kerakli axborotni manbalar yoki ma'lumotlar to'plamidan topish jarayoni
\item elementlarni tanlangan mezon bo'yicha tartibga keltirish jarayoni
\end{enumerate}
\begin{tcolorbox}[enhanced,breakable,colback=teal!4,colframe=teal!50!black,boxrule=0.5pt,arc=1.1mm,left=3mm,right=3mm,top=1.7mm,bottom=1.7mm,before skip=2mm,after skip=2mm]
\textbf{\color{teal}To'g'ri javob:} \textbf{A) axborotni keyinchalik foydalanish uchun saqlab qolish hamda uzoq muddat tartibli jamlash jarayoni}
\par\smallskip
\textbf{Izoh:} Saqlash va arxivlash - axborotni keyinchalik foydalanish uchun saqlab qolish hamda uzoq muddat tartibli jamlash jarayoni.
\end{tcolorbox}
\vspace{2mm}\hrule height 0.2pt \vspace{3mm}

% DB question_id: f07d97b0-bdcd-4fdf-b722-ec3a1270215b
\needspace{5\baselineskip}
\noindent\textbf{\color{ink}329.}
\par
Creative Commons qisqartmalarini mos cheklov yoki talab bilan juftlang.
\textbf{\footnotesize Chap ustun:}
\begin{enumerate}[label=\arabic*),leftmargin=9mm]
\item BY
\item SA
\item NC
\item ND
\end{enumerate}
\textbf{\footnotesize O'ng ustun:}
\begin{enumerate}[label=\Alph*),leftmargin=9mm]
\item Tijoriy foydalanishni cheklash
\item Muallif/manbani ko'rsatish talabi
\item O'zgartirilgan hosila nusxalarni tarqatishni cheklash
\item Hosila asarni xuddi shu litsenziya shartida tarqatish
\end{enumerate}
\begin{tcolorbox}[enhanced,breakable,colback=teal!4,colframe=teal!50!black,boxrule=0.5pt,arc=1.1mm,left=3mm,right=3mm,top=1.7mm,bottom=1.7mm,before skip=2mm,after skip=2mm]
\textbf{\color{teal}To'g'ri javob:} \textbf{1-B, 2-D, 3-A, 4-C}
\par\smallskip
\textbf{Izoh:} CC belgilarida BY - attribution, SA - share alike, NC - non-commercial, ND - no derivatives.
\end{tcolorbox}
\vspace{2mm}\hrule height 0.2pt \vspace{3mm}

% DB question_id: f08c8ee6-92dd-506d-958e-d77bb9364d52
\needspace{5\baselineskip}
\noindent\textbf{\color{ink}330.}
\par
Quyidagi ta'rif qaysi tushunchaga tegishli? Umumiy belgi bo'yicha elementlarni toifalarga ajratish jarayoni.
\begin{enumerate}[label=\Alph*),leftmargin=9mm]
\item Guruhlash
\item Uzatish va qabul qilish
\item Nusxalash
\item Tarqatish
\end{enumerate}
\begin{tcolorbox}[enhanced,breakable,colback=teal!4,colframe=teal!50!black,boxrule=0.5pt,arc=1.1mm,left=3mm,right=3mm,top=1.7mm,bottom=1.7mm,before skip=2mm,after skip=2mm]
\textbf{\color{teal}To'g'ri javob:} \textbf{A) Guruhlash}
\par\smallskip
\textbf{Izoh:} Ta'rif Guruhlash tushunchasiga tegishli.
\end{tcolorbox}
\vspace{2mm}\hrule height 0.2pt \vspace{3mm}

% DB question_id: f2c5dfc9-c239-5a31-a649-604ae7d23cc2
\needspace{5\baselineskip}
\noindent\textbf{\color{ink}331.}
\par
Bir xil natija o'quvchiga sodda, mutaxassisga esa batafsil shaklda berildi. Ushbu vaziyatdagi asosiy tushunchani aniqlang.
\begin{enumerate}[label=\Alph*),leftmargin=9mm]
\item Elektron pochta madaniyati
\item Netiket
\item Auditoriya, maqsad va kontekst
\item Raqamli iz
\end{enumerate}
\begin{tcolorbox}[enhanced,breakable,colback=teal!4,colframe=teal!50!black,boxrule=0.5pt,arc=1.1mm,left=3mm,right=3mm,top=1.7mm,bottom=1.7mm,before skip=2mm,after skip=2mm]
\textbf{\color{teal}To'g'ri javob:} \textbf{C) Auditoriya, maqsad va kontekst}
\par\smallskip
\textbf{Izoh:} Vaziyat Auditoriya, maqsad va kontekst tushunchasini ifodalaydi: axborotning kimga, nima uchun va qaysi sharoitda taqdim etilishini hisobga olish tamoyili.
\end{tcolorbox}
\vspace{2mm}\hrule height 0.2pt \vspace{3mm}

% DB question_id: f4ae44b1-4817-5f93-9015-4287e6c3b168
\needspace{5\baselineskip}
\noindent\textbf{\color{ink}332.}
\par
Quyidagi ta'rif qaysi tushunchaga tegishli? Axborotning kimga, nima uchun va qaysi sharoitda taqdim etilishini hisobga olish tamoyili.
\begin{enumerate}[label=\Alph*),leftmargin=9mm]
\item Elektron pochta madaniyati
\item Netiket
\item Raqamli iz
\item Auditoriya, maqsad va kontekst
\end{enumerate}
\begin{tcolorbox}[enhanced,breakable,colback=teal!4,colframe=teal!50!black,boxrule=0.5pt,arc=1.1mm,left=3mm,right=3mm,top=1.7mm,bottom=1.7mm,before skip=2mm,after skip=2mm]
\textbf{\color{teal}To'g'ri javob:} \textbf{D) Auditoriya, maqsad va kontekst}
\par\smallskip
\textbf{Izoh:} Ta'rif Auditoriya, maqsad va kontekst tushunchasiga tegishli.
\end{tcolorbox}
\vspace{2mm}\hrule height 0.2pt \vspace{3mm}

% DB question_id: f520c5ea-f7af-5a39-851f-d95199b42364
\needspace{5\baselineskip}
\noindent\textbf{\color{ink}333.}
\par
Fayllar matn, rasm va video turlariga ajratildi. Ushbu vaziyatdagi asosiy tushunchani aniqlang.
\begin{enumerate}[label=\Alph*),leftmargin=9mm]
\item Uzatish va qabul qilish
\item Nusxalash
\item Tarqatish
\item Guruhlash
\end{enumerate}
\begin{tcolorbox}[enhanced,breakable,colback=teal!4,colframe=teal!50!black,boxrule=0.5pt,arc=1.1mm,left=3mm,right=3mm,top=1.7mm,bottom=1.7mm,before skip=2mm,after skip=2mm]
\textbf{\color{teal}To'g'ri javob:} \textbf{D) Guruhlash}
\par\smallskip
\textbf{Izoh:} Vaziyat Guruhlash tushunchasini ifodalaydi: umumiy belgi bo'yicha elementlarni toifalarga ajratish jarayoni.
\end{tcolorbox}
\vspace{2mm}\hrule height 0.2pt \vspace{3mm}

% DB question_id: f58aaaec-3f26-5aa0-add9-702799dbd23b
\needspace{5\baselineskip}
\noindent\textbf{\color{ink}334.}
\par
Foydalanuvchi hujjatlar orasidan kerakli hisobotni so'rov bilan topdi. Ushbu vaziyatdagi asosiy tushunchani aniqlang.
\begin{enumerate}[label=\Alph*),leftmargin=9mm]
\item Izlash
\item Saralash
\item Filtrlash
\item Guruhlash
\end{enumerate}
\begin{tcolorbox}[enhanced,breakable,colback=teal!4,colframe=teal!50!black,boxrule=0.5pt,arc=1.1mm,left=3mm,right=3mm,top=1.7mm,bottom=1.7mm,before skip=2mm,after skip=2mm]
\textbf{\color{teal}To'g'ri javob:} \textbf{A) Izlash}
\par\smallskip
\textbf{Izoh:} Vaziyat Izlash tushunchasini ifodalaydi: kerakli axborotni manbalar yoki ma'lumotlar to'plamidan topish jarayoni.
\end{tcolorbox}
\vspace{2mm}\hrule height 0.2pt \vspace{3mm}

% DB question_id: f80a19a3-78dc-58ea-a7c9-dad3b7781aa5
\needspace{5\baselineskip}
\noindent\textbf{\color{ink}335.}
\par
"Raqamli iz" atamasining eng aniq ta'rifini toping.
\begin{enumerate}[label=\Alph*),leftmargin=9mm]
\item mazmunni auditoriya ehtiyoji va maqsadga eng mos matn, jadval, diagramma, audio yoki video shaklida berish qarori
\item yangi mazmun hosil qilish hamda kerakli ma'lumotlarni turli manbalardan to'plash jarayoni
\item axborotni sezgi yoki qurilma orqali olish va uni axborot tizimiga kiritish jarayoni
\item onlayn faoliyat natijasida xizmatlar va qurilmalarda qoladigan ma'lumotlar majmui
\end{enumerate}
\begin{tcolorbox}[enhanced,breakable,colback=teal!4,colframe=teal!50!black,boxrule=0.5pt,arc=1.1mm,left=3mm,right=3mm,top=1.7mm,bottom=1.7mm,before skip=2mm,after skip=2mm]
\textbf{\color{teal}To'g'ri javob:} \textbf{D) onlayn faoliyat natijasida xizmatlar va qurilmalarda qoladigan ma'lumotlar majmui}
\par\smallskip
\textbf{Izoh:} Raqamli iz - onlayn faoliyat natijasida xizmatlar va qurilmalarda qoladigan ma'lumotlar majmui.
\end{tcolorbox}
\vspace{2mm}\hrule height 0.2pt \vspace{3mm}

% DB question_id: f868e1ed-da21-5b35-a729-1c46fa92d46b
\needspace{5\baselineskip}
\noindent\textbf{\color{ink}336.}
\par
Foydalanuvchi qidiruv satriga "audio fayl hajmi formula" deb yozdi. Ushbu vaziyatdagi asosiy tushunchani aniqlang.
\begin{enumerate}[label=\Alph*),leftmargin=9mm]
\item Aniq ibora bilan qidirish
\item Mantiqiy qidiruv operatorlari
\item Qidiruv tizimi va kalit so'z
\item Qidiruv natijalarini baholash
\end{enumerate}
\begin{tcolorbox}[enhanced,breakable,colback=teal!4,colframe=teal!50!black,boxrule=0.5pt,arc=1.1mm,left=3mm,right=3mm,top=1.7mm,bottom=1.7mm,before skip=2mm,after skip=2mm]
\textbf{\color{teal}To'g'ri javob:} \textbf{C) Qidiruv tizimi va kalit so'z}
\par\smallskip
\textbf{Izoh:} Vaziyat Qidiruv tizimi va kalit so'z tushunchasini ifodalaydi: veb resurslarini indekslab, foydalanuvchi yozgan muhim so'zlarga mos natijalarni topib beruvchi xizmat va so'rov elementi.
\end{tcolorbox}
\vspace{2mm}\hrule height 0.2pt \vspace{3mm}

% DB question_id: f8b17255-2811-592e-be51-c4b6f523c9de
\needspace{5\baselineskip}
\noindent\textbf{\color{ink}337.}
\par
Quyidagi ta'rif qaysi tushunchaga tegishli? Raqamli muloqotda hurmat, aniqlik va xavfsizlikka rioya qilish qoidalari.
\begin{enumerate}[label=\Alph*),leftmargin=9mm]
\item Raqamli iz
\item Netiket
\item Axborotni taqdim etish shaklini tanlash
\item Axborotni yaratish va yig'ish
\end{enumerate}
\begin{tcolorbox}[enhanced,breakable,colback=teal!4,colframe=teal!50!black,boxrule=0.5pt,arc=1.1mm,left=3mm,right=3mm,top=1.7mm,bottom=1.7mm,before skip=2mm,after skip=2mm]
\textbf{\color{teal}To'g'ri javob:} \textbf{B) Netiket}
\par\smallskip
\textbf{Izoh:} Ta'rif Netiket tushunchasiga tegishli.
\end{tcolorbox}
\vspace{2mm}\hrule height 0.2pt \vspace{3mm}

% DB question_id: fd1106a2-531f-4526-ba06-ef1cd85374a9
\needspace{5\baselineskip}
\noindent\textbf{\color{ink}338.}
\par
Internetdan olingan axborot asosida ishonchli xulosa chiqarish jarayonini tartiblang.
\begin{enumerate}[label=\Alph*),leftmargin=9mm]
\item Manba, dalil, sana va muallifni baholash
\item Axborot ehtiyoji va savolni aniq belgilash
\item Tekshirilgan dalillarni solishtirib xulosa chiqarish
\item Mos manbalarni qidirib topish
\end{enumerate}
\begin{tcolorbox}[enhanced,breakable,colback=teal!4,colframe=teal!50!black,boxrule=0.5pt,arc=1.1mm,left=3mm,right=3mm,top=1.7mm,bottom=1.7mm,before skip=2mm,after skip=2mm]
\textbf{\color{teal}To'g'ri javob:} \textbf{B -- D -- A -- C}
\par\smallskip
\textbf{Izoh:} Savol \ensuremath{\to} qidiruv \ensuremath{\to} manbani baholash \ensuremath{\to} dalillardan xulosa.
\end{tcolorbox}
\vspace{2mm}\hrule height 0.2pt \vspace{3mm}

% DB question_id: fd1df548-15fe-441b-becc-eb9e356bb3fa
\needspace{5\baselineskip}
\noindent\textbf{\color{ink}339.}
\par
O'quvchi bulutli hujjat, videomuloqot va LMSdan birgalikda foydalandi. Bu vaziyatdagi asosiy tushunchani aniqlang.
\begin{enumerate}[label=\Alph*),leftmargin=9mm]
\item Axborotlashgan jamiyat
\item Raqamli muhit
\item Axborot madaniyati
\item Raqamli savodxonlik
\end{enumerate}
\begin{tcolorbox}[enhanced,breakable,colback=teal!4,colframe=teal!50!black,boxrule=0.5pt,arc=1.1mm,left=3mm,right=3mm,top=1.7mm,bottom=1.7mm,before skip=2mm,after skip=2mm]
\textbf{\color{teal}To'g'ri javob:} \textbf{B) Raqamli muhit}
\par\smallskip
\textbf{Izoh:} Vaziyat Raqamli muhitga misol: raqamli qurilmalar, dasturlar, tarmoqlar va xizmatlar o'zaro ishlaydigan makon.
\end{tcolorbox}
\vspace{2mm}\hrule height 0.2pt \vspace{3mm}

% DB question_id: 083af927-4fb0-44e8-a1df-5b335d0be358
\needspace{5\baselineskip}
\noindent\textbf{\color{ink}340.}
\par
So'rovnoma javoblari yagona jadvalga to'plandi. Vaziyatni tahlil qilib, unga eng aniq mos tushunchani toping.
\begin{enumerate}[label=\Alph*),leftmargin=9mm]
\item Axborotni qayta ishlash
\item Axborotni yig'ish
\item Axborotni saqlash
\item Axborot manbai
\end{enumerate}
\begin{tcolorbox}[enhanced,breakable,colback=teal!4,colframe=teal!50!black,boxrule=0.5pt,arc=1.1mm,left=3mm,right=3mm,top=1.7mm,bottom=1.7mm,before skip=2mm,after skip=2mm]
\textbf{\color{teal}To'g'ri javob:} \textbf{B) Axborotni yig'ish}
\par\smallskip
\textbf{Izoh:} Vaziyat Axborotni yig'ishga misol: turli manbalardan kerakli ma'lumotlarni olish va jamlash jarayoni.
\end{tcolorbox}
\vspace{2mm}\hrule height 0.2pt \vspace{3mm}

% DB question_id: 4651f092-0384-4732-9266-ddb424df0b88
\needspace{5\baselineskip}
\noindent\textbf{\color{ink}341.}
\par
Asl fayl ulashildi, ammo tahrirlangan versiya tarqatilmadi. Vaziyatni tahlil qilib, unga eng aniq mos tushunchani toping.
\begin{enumerate}[label=\Alph*),leftmargin=9mm]
\item Manba ko'rsatish
\item Netiket
\item Mualliflik huquqi
\item CC ND
\end{enumerate}
\begin{tcolorbox}[enhanced,breakable,colback=teal!4,colframe=teal!50!black,boxrule=0.5pt,arc=1.1mm,left=3mm,right=3mm,top=1.7mm,bottom=1.7mm,before skip=2mm,after skip=2mm]
\textbf{\color{teal}To'g'ri javob:} \textbf{D) CC ND}
\par\smallskip
\textbf{Izoh:} Vaziyat CC NDga misol: asarning o'zgartirilgan hosila nusxalarini tarqatishni cheklaydigan shart.
\end{tcolorbox}
\vspace{2mm}\hrule height 0.2pt \vspace{3mm}

% DB question_id: 5e3f9557-c704-4413-81fe-83d4edc5e0b9
\needspace{5\baselineskip}
\noindent\textbf{\color{ink}342.}
\par
O'lchov natijasi 23,4 \ensuremath{^\circ}C deb asbob imkoniyatiga mos yozildi. Vaziyatni tahlil qilib, unga eng aniq mos tushunchani toping.
\begin{enumerate}[label=\Alph*),leftmargin=9mm]
\item Aniqlik
\item Ishonchlilik
\item To'liqlik
\item Matnli axborot
\end{enumerate}
\begin{tcolorbox}[enhanced,breakable,colback=teal!4,colframe=teal!50!black,boxrule=0.5pt,arc=1.1mm,left=3mm,right=3mm,top=1.7mm,bottom=1.7mm,before skip=2mm,after skip=2mm]
\textbf{\color{teal}To'g'ri javob:} \textbf{A) Aniqlik}
\par\smallskip
\textbf{Izoh:} Vaziyat Aniqlikga misol: axborotning obyekt holatini xatosiz va kerakli aniqlik darajasida aks ettirish xossasi.
\end{tcolorbox}
\vspace{2mm}\hrule height 0.2pt \vspace{3mm}

% DB question_id: 619866ae-bc81-45ef-b983-75edb8232876
\needspace{5\baselineskip}
\noindent\textbf{\color{ink}343.}
\par
Material reklama daromadi olinadigan mahsulotda ruxsatsiz ishlatilmadi. Vaziyatni tahlil qilib, unga eng aniq mos tushunchani toping.
\begin{enumerate}[label=\Alph*),leftmargin=9mm]
\item CC ND
\item Manba ko'rsatish
\item CC NC
\item Netiket
\end{enumerate}
\begin{tcolorbox}[enhanced,breakable,colback=teal!4,colframe=teal!50!black,boxrule=0.5pt,arc=1.1mm,left=3mm,right=3mm,top=1.7mm,bottom=1.7mm,before skip=2mm,after skip=2mm]
\textbf{\color{teal}To'g'ri javob:} \textbf{C) CC NC}
\par\smallskip
\textbf{Izoh:} Vaziyat CC NCga misol: asardan tijorat maqsadida foydalanishni cheklaydigan shart.
\end{tcolorbox}
\vspace{2mm}\hrule height 0.2pt \vspace{3mm}

% DB question_id: 73036021-7d95-4030-a28d-acc57491f494
\needspace{5\baselineskip}
\noindent\textbf{\color{ink}344.}
\par
Fayl xossalarida uning hajmi 8 MiB ekani ko'rsatildi. Vaziyatni tahlil qilib, unga eng aniq mos tushunchani toping.
\begin{enumerate}[label=\Alph*),leftmargin=9mm]
\item Fikr
\item Xulosa
\item Fakt
\item Axborot texnologiyasi
\end{enumerate}
\begin{tcolorbox}[enhanced,breakable,colback=teal!4,colframe=teal!50!black,boxrule=0.5pt,arc=1.1mm,left=3mm,right=3mm,top=1.7mm,bottom=1.7mm,before skip=2mm,after skip=2mm]
\textbf{\color{teal}To'g'ri javob:} \textbf{C) Fakt}
\par\smallskip
\textbf{Izoh:} Vaziyat Faktga misol: dalil yordamida tekshirish mumkin bo'lgan da'vo.
\end{tcolorbox}
\vspace{2mm}\hrule height 0.2pt \vspace{3mm}

% DB question_id: 8dfe2cb0-cc86-4554-8e6a-b1eded5e834e
\needspace{5\baselineskip}
\noindent\textbf{\color{ink}345.}
\par
Saytga kirish, izoh yozish va joylashuv ruxsati tarixda qayd etildi. Vaziyatni tahlil qilib, unga eng aniq mos tushunchani toping.
\begin{enumerate}[label=\Alph*),leftmargin=9mm]
\item Fakt
\item Raqamli iz
\item Fikr
\item Xulosa
\end{enumerate}
\begin{tcolorbox}[enhanced,breakable,colback=teal!4,colframe=teal!50!black,boxrule=0.5pt,arc=1.1mm,left=3mm,right=3mm,top=1.7mm,bottom=1.7mm,before skip=2mm,after skip=2mm]
\textbf{\color{teal}To'g'ri javob:} \textbf{B) Raqamli iz}
\par\smallskip
\textbf{Izoh:} Vaziyat Raqamli izga misol: foydalanuvchining onlayn harakatlari natijasida qoladigan ma'lumotlar majmui.
\end{tcolorbox}
\vspace{2mm}\hrule height 0.2pt \vspace{3mm}

% DB question_id: 966cc545-c7d9-470d-8529-aae5025e08e7
\needspace{5\baselineskip}
\noindent\textbf{\color{ink}346.}
\par
Tahrirlangan material asl litsenziyaga mos litsenziyada tarqatildi. Vaziyatni tahlil qilib, unga eng aniq mos tushunchani toping.
\begin{enumerate}[label=\Alph*),leftmargin=9mm]
\item CC NC
\item CC SA
\item CC ND
\item Manba ko'rsatish
\end{enumerate}
\begin{tcolorbox}[enhanced,breakable,colback=teal!4,colframe=teal!50!black,boxrule=0.5pt,arc=1.1mm,left=3mm,right=3mm,top=1.7mm,bottom=1.7mm,before skip=2mm,after skip=2mm]
\textbf{\color{teal}To'g'ri javob:} \textbf{B) CC SA}
\par\smallskip
\textbf{Izoh:} Vaziyat CC SAga misol: hosila asarni ayni yoki mos litsenziya bilan ulashishni talab qiladigan shart.
\end{tcolorbox}
\vspace{2mm}\hrule height 0.2pt \vspace{3mm}

% DB question_id: b78aed39-e948-41cd-b8cd-522d9a6986dc
\needspace{5\baselineskip}
\noindent\textbf{\color{ink}347.}
\par
Foydalanuvchi faylni yaratdi, tekshirdi, xavfsiz ulashdi va natijani baholadi. Vaziyatni tahlil qilib, unga eng aniq mos tushunchani toping.
\begin{enumerate}[label=\Alph*),leftmargin=9mm]
\item Raqamli savodxonlik
\item Raqamli iz
\item Fakt
\item Fikr
\end{enumerate}
\begin{tcolorbox}[enhanced,breakable,colback=teal!4,colframe=teal!50!black,boxrule=0.5pt,arc=1.1mm,left=3mm,right=3mm,top=1.7mm,bottom=1.7mm,before skip=2mm,after skip=2mm]
\textbf{\color{teal}To'g'ri javob:} \textbf{A) Raqamli savodxonlik}
\par\smallskip
\textbf{Izoh:} Vaziyat Raqamli savodxonlikga misol: raqamli vositalardan maqsadli, tanqidiy va samarali foydalanish qobiliyati.
\end{tcolorbox}
\vspace{2mm}\hrule height 0.2pt \vspace{3mm}

% DB question_id: 43633c79-45db-4b54-9b91-e4ff44cf5fd1
\needspace{5\baselineskip}
\noindent\textbf{\color{ink}348.}
\par
Muallif soxta raqamni odamlarni chalg'itish uchun bilib turib e'lon qildi. Yaqin tushunchalarni farqlab, aynan qaysi atama vaziyatni eng aniq ifodalashini toping.
\begin{enumerate}[label=\Alph*),leftmargin=9mm]
\item Internet
\item WWW
\item Veb-brauzer
\item Dezinformatsiya
\end{enumerate}
\begin{tcolorbox}[enhanced,breakable,colback=teal!4,colframe=teal!50!black,boxrule=0.5pt,arc=1.1mm,left=3mm,right=3mm,top=1.7mm,bottom=1.7mm,before skip=2mm,after skip=2mm]
\textbf{\color{teal}To'g'ri javob:} \textbf{D) Dezinformatsiya}
\par\smallskip
\textbf{Izoh:} Vaziyat Dezinformatsiyaga misol: aldash maqsadida ataylab tarqatilgan noto'g'ri axborot.
\end{tcolorbox}
\vspace{2mm}\hrule height 0.2pt \vspace{3mm}

% DB question_id: 81a3bcb3-7148-460b-a4d5-913fed8f0ed3
\needspace{5\baselineskip}
\noindent\textbf{\color{ink}349.}
\par
Baholar yig'indisi va o'rtacha qiymati hisoblandi. Yaqin tushunchalarni farqlab, aynan qaysi atama vaziyatni eng aniq ifodalashini toping.
\begin{enumerate}[label=\Alph*),leftmargin=9mm]
\item Axborotni saqlash
\item Axborot manbai
\item Axborotni qayta ishlash
\item Aloqa kanali
\end{enumerate}
\begin{tcolorbox}[enhanced,breakable,colback=teal!4,colframe=teal!50!black,boxrule=0.5pt,arc=1.1mm,left=3mm,right=3mm,top=1.7mm,bottom=1.7mm,before skip=2mm,after skip=2mm]
\textbf{\color{teal}To'g'ri javob:} \textbf{C) Axborotni qayta ishlash}
\par\smallskip
\textbf{Izoh:} Vaziyat Axborotni qayta ishlashga misol: mavjud ma'lumotni tartiblash, hisoblash yoki tahlil qilish orqali yangi natija olish jarayoni.
\end{tcolorbox}
\vspace{2mm}\hrule height 0.2pt \vspace{3mm}

% DB question_id: 921aceed-0cfd-4f32-ac1c-1f0163f5baaf
\needspace{5\baselineskip}
\noindent\textbf{\color{ink}350.}
\par
Slaydda rasm muallifi, URL va CC litsenziyasi yozildi. Yaqin tushunchalarni farqlab, aynan qaysi atama vaziyatni eng aniq ifodalashini toping.
\begin{enumerate}[label=\Alph*),leftmargin=9mm]
\item Manba ko'rsatish
\item Netiket
\item Mualliflik huquqi
\item Plagiat
\end{enumerate}
\begin{tcolorbox}[enhanced,breakable,colback=teal!4,colframe=teal!50!black,boxrule=0.5pt,arc=1.1mm,left=3mm,right=3mm,top=1.7mm,bottom=1.7mm,before skip=2mm,after skip=2mm]
\textbf{\color{teal}To'g'ri javob:} \textbf{A) Manba ko'rsatish}
\par\smallskip
\textbf{Izoh:} Vaziyat Manba ko'rsatishga misol: muallif, asar nomi, manba va litsenziyani imkon qadar aniq qayd etish.
\end{tcolorbox}
\vspace{2mm}\hrule height 0.2pt \vspace{3mm}

% DB question_id: 92d783e9-1b52-426d-884d-3c63fca62d6b
\needspace{5\baselineskip}
\noindent\textbf{\color{ink}351.}
\par
Qulf belgili saytdagi da'vo baribir boshqa manbalar bilan tekshirildi. Yaqin tushunchalarni farqlab, aynan qaysi atama vaziyatni eng aniq ifodalashini toping.
\begin{enumerate}[label=\Alph*),leftmargin=9mm]
\item Axborot va ma'lumot farqi
\item Kodlash va shifrlash farqi
\item Bit va bayt farqi
\item HTTPS va ishonchlilik farqi
\end{enumerate}
\begin{tcolorbox}[enhanced,breakable,colback=teal!4,colframe=teal!50!black,boxrule=0.5pt,arc=1.1mm,left=3mm,right=3mm,top=1.7mm,bottom=1.7mm,before skip=2mm,after skip=2mm]
\textbf{\color{teal}To'g'ri javob:} \textbf{D) HTTPS va ishonchlilik farqi}
\par\smallskip
\textbf{Izoh:} Vaziyat HTTPS va ishonchlilik farqiga misol: HTTPS transportni shifrlashi, sahifa mazmunining haqiqatligini kafolatlamasligini ajratish.
\end{tcolorbox}
\vspace{2mm}\hrule height 0.2pt \vspace{3mm}

% DB question_id: 9633f05f-f93a-4db0-8d5c-ef0e90df0efa
\needspace{5\baselineskip}
\noindent\textbf{\color{ink}352.}
\par
Bo'sh joy 4 MiB, fayl 8 MiB bo'lgani uchun u sig'masligi aytildi. Yaqin tushunchalarni farqlab, aynan qaysi atama vaziyatni eng aniq ifodalashini toping.
\begin{enumerate}[label=\Alph*),leftmargin=9mm]
\item Xulosa
\item Axborot texnologiyasi
\item Raqamli axborot
\item Raqamli muhit
\end{enumerate}
\begin{tcolorbox}[enhanced,breakable,colback=teal!4,colframe=teal!50!black,boxrule=0.5pt,arc=1.1mm,left=3mm,right=3mm,top=1.7mm,bottom=1.7mm,before skip=2mm,after skip=2mm]
\textbf{\color{teal}To'g'ri javob:} \textbf{A) Xulosa}
\par\smallskip
\textbf{Izoh:} Vaziyat Xulosaga misol: faktlar va qoidalar asosida chiqarilgan natija.
\end{tcolorbox}
\vspace{2mm}\hrule height 0.2pt \vspace{3mm}

% DB question_id: 9d71abc0-bd3e-43a6-a37d-ad78ed54d108
\needspace{5\baselineskip}
\noindent\textbf{\color{ink}353.}
\par
Maqola muallifi va dalillari tekshirilib, boshqa manba bilan solishtirildi. Yaqin tushunchalarni farqlab, aynan qaysi atama vaziyatni eng aniq ifodalashini toping.
\begin{enumerate}[label=\Alph*),leftmargin=9mm]
\item Dezinformatsiya
\item Internet
\item Manba ishonchliligi
\item WWW
\end{enumerate}
\begin{tcolorbox}[enhanced,breakable,colback=teal!4,colframe=teal!50!black,boxrule=0.5pt,arc=1.1mm,left=3mm,right=3mm,top=1.7mm,bottom=1.7mm,before skip=2mm,after skip=2mm]
\textbf{\color{teal}To'g'ri javob:} \textbf{C) Manba ishonchliligi}
\par\smallskip
\textbf{Izoh:} Vaziyat Manba ishonchliligiga misol: muallif, sana, dalil, maqsad va mustaqil tasdiq orqali baholanadigan sifat.
\end{tcolorbox}
\vspace{2mm}\hrule height 0.2pt \vspace{3mm}

% DB question_id: ceffba25-afeb-4bbc-9a00-66feb5e57675
\needspace{5\baselineskip}
\noindent\textbf{\color{ink}354.}
\par
Jadvalda natija bor, ammo sana va o'lchov birligi yo'qligi sabab xulosa chiqarib bo'lmadi. Yaqin tushunchalarni farqlab, aynan qaysi atama vaziyatni eng aniq ifodalashini toping.
\begin{enumerate}[label=\Alph*),leftmargin=9mm]
\item Matnli axborot
\item Grafik axborot
\item To'liqlik
\item Audioaxborot
\end{enumerate}
\begin{tcolorbox}[enhanced,breakable,colback=teal!4,colframe=teal!50!black,boxrule=0.5pt,arc=1.1mm,left=3mm,right=3mm,top=1.7mm,bottom=1.7mm,before skip=2mm,after skip=2mm]
\textbf{\color{teal}To'g'ri javob:} \textbf{C) To'liqlik}
\par\smallskip
\textbf{Izoh:} Vaziyat To'liqlikga misol: qaror chiqarish uchun zarur bo'lgan muhim qismlar yetarli berilganlik xossasi.
\end{tcolorbox}
\vspace{2mm}\hrule height 0.2pt \vspace{3mm}

% DB question_id: dfff606a-8d1a-46ba-9120-c70e23125819
\needspace{5\baselineskip}
\noindent\textbf{\color{ink}355.}
\par
Foydalanuvchi haqorat qilmay, mavzuga mos va shaxsiy ma'lumotni oshkor etmay yozdi. Yaqin tushunchalarni farqlab, aynan qaysi atama vaziyatni eng aniq ifodalashini toping.
\begin{enumerate}[label=\Alph*),leftmargin=9mm]
\item Mualliflik huquqi
\item Netiket
\item Plagiat
\item Litsenziya
\end{enumerate}
\begin{tcolorbox}[enhanced,breakable,colback=teal!4,colframe=teal!50!black,boxrule=0.5pt,arc=1.1mm,left=3mm,right=3mm,top=1.7mm,bottom=1.7mm,before skip=2mm,after skip=2mm]
\textbf{\color{teal}To'g'ri javob:} \textbf{B) Netiket}
\par\smallskip
\textbf{Izoh:} Vaziyat Netiketga misol: raqamli muloqotdagi hurmatli va mas'uliyatli xulq qoidalari.
\end{tcolorbox}
\vspace{2mm}\hrule height 0.2pt \vspace{3mm}

% DB question_id: e59a81a4-45f7-4643-8728-78330396b8ec
\needspace{5\baselineskip}
\noindent\textbf{\color{ink}356.}
\par
Foydalanuvchi "bu dastur eng qulay" deb o'z qarashini bildirdi. Yaqin tushunchalarni farqlab, aynan qaysi atama vaziyatni eng aniq ifodalashini toping.
\begin{enumerate}[label=\Alph*),leftmargin=9mm]
\item Xulosa
\item Axborot texnologiyasi
\item Raqamli axborot
\item Fikr
\end{enumerate}
\begin{tcolorbox}[enhanced,breakable,colback=teal!4,colframe=teal!50!black,boxrule=0.5pt,arc=1.1mm,left=3mm,right=3mm,top=1.7mm,bottom=1.7mm,before skip=2mm,after skip=2mm]
\textbf{\color{teal}To'g'ri javob:} \textbf{D) Fikr}
\par\smallskip
\textbf{Izoh:} Vaziyat Fikrga misol: shaxsiy baho, did yoki munosabatni ifodalovchi mulohaza.
\end{tcolorbox}
\vspace{2mm}\hrule height 0.2pt \vspace{3mm}

% DB question_id: eaf608e9-1619-479b-a6ec-9a20306757be
\needspace{5\baselineskip}
\noindent\textbf{\color{ink}357.}
\par
Hujjat server diskiga yozilib, zaxira nusxasi yaratildi. Yaqin tushunchalarni farqlab, aynan qaysi atama vaziyatni eng aniq ifodalashini toping.
\begin{enumerate}[label=\Alph*),leftmargin=9mm]
\item Axborot manbai
\item Aloqa kanali
\item Qabul qiluvchi
\item Axborotni saqlash
\end{enumerate}
\begin{tcolorbox}[enhanced,breakable,colback=teal!4,colframe=teal!50!black,boxrule=0.5pt,arc=1.1mm,left=3mm,right=3mm,top=1.7mm,bottom=1.7mm,before skip=2mm,after skip=2mm]
\textbf{\color{teal}To'g'ri javob:} \textbf{D) Axborotni saqlash}
\par\smallskip
\textbf{Izoh:} Vaziyat Axborotni saqlashga misol: axborotni keyinchalik foydalanish uchun tashuvchida qayd etish jarayoni.
\end{tcolorbox}
\vspace{2mm}\hrule height 0.2pt \vspace{3mm}

% DB question_id: f2c2c5e3-3f40-4955-afd7-32ebda0c4e44
\needspace{5\baselineskip}
\noindent\textbf{\color{ink}358.}
\par
Xabar ikki mustaqil rasmiy manba bilan tasdiqlandi. Yaqin tushunchalarni farqlab, aynan qaysi atama vaziyatni eng aniq ifodalashini toping.
\begin{enumerate}[label=\Alph*),leftmargin=9mm]
\item To'liqlik
\item Ishonchlilik
\item Matnli axborot
\item Grafik axborot
\end{enumerate}
\begin{tcolorbox}[enhanced,breakable,colback=teal!4,colframe=teal!50!black,boxrule=0.5pt,arc=1.1mm,left=3mm,right=3mm,top=1.7mm,bottom=1.7mm,before skip=2mm,after skip=2mm]
\textbf{\color{teal}To'g'ri javob:} \textbf{B) Ishonchlilik}
\par\smallskip
\textbf{Izoh:} Vaziyat Ishonchlilikga misol: axborotning tekshirilgan dalil va ishonchli manbaga tayanganlik xossasi.
\end{tcolorbox}
\vspace{2mm}\hrule height 0.2pt \vspace{3mm}
\section{Belgi, kod, kodlash va axborot hajmi}

% DB question_id: 290605ba-6de5-52b5-aa8a-d7d94a196a6a
\needspace{5\baselineskip}
\noindent\textbf{\color{ink}359.}
\par
Sezar shifrida \texttt{A \ensuremath{\to} D}, \texttt{B \ensuremath{\to} E} tarzida har bir harf 3 ta o'rin oldinga surilgan. \texttt{PRWKHUERDUG} kodini deshifrlang.
\begin{enumerate}[label=\Alph*),leftmargin=9mm]
\item MOTHERBEARD
\item MOTHERBOARD
\item SPEAKERBOARD
\item SPEAKERBEARD
\end{enumerate}
\begin{tcolorbox}[enhanced,breakable,colback=teal!4,colframe=teal!50!black,boxrule=0.5pt,arc=1.1mm,left=3mm,right=3mm,top=1.7mm,bottom=1.7mm,before skip=2mm,after skip=2mm]
\textbf{\color{teal}To'g'ri javob:} \textbf{B) MOTHERBOARD}
\par\smallskip
\textbf{Izoh:} Deshifrlashda har bir harf alifboda 3 ta o'rin orqaga suriladi: \texttt{PRWKHUERDUG \ensuremath{\to} MOTHERBOARD}.
\end{tcolorbox}
\vspace{2mm}\hrule height 0.2pt \vspace{3mm}

% DB question_id: 297e4362-5d8a-4049-b6a1-5ed7ba96a3ec
\needspace{5\baselineskip}
\noindent\textbf{\color{ink}360.}
\par
"Dekodlash" atamasiga eng aniq ta'rifni toping.
\begin{enumerate}[label=\Alph*),leftmargin=9mm]
\item kodlangan axborotni kelishilgan qoida asosida qayta tiklash
\item har bir belgi bir xil sondagi kod elementi bilan ifodalanadigan kod
\item turli belgilar turli uzunlikdagi kodlar bilan ifodalanadigan kod
\item axborotni saqlash yoki uzatish uchun fayl hajmini kamaytirish jarayoni
\end{enumerate}
\begin{tcolorbox}[enhanced,breakable,colback=teal!4,colframe=teal!50!black,boxrule=0.5pt,arc=1.1mm,left=3mm,right=3mm,top=1.7mm,bottom=1.7mm,before skip=2mm,after skip=2mm]
\textbf{\color{teal}To'g'ri javob:} \textbf{A) kodlangan axborotni kelishilgan qoida asosida qayta tiklash}
\par\smallskip
\textbf{Izoh:} Dekodlash - kodlangan axborotni kelishilgan qoida asosida qayta tiklash.
\end{tcolorbox}
\vspace{2mm}\hrule height 0.2pt \vspace{3mm}

% DB question_id: 2d4b2e75-0590-5929-974f-d71a96d0cbeb
\needspace{5\baselineskip}
\noindent\textbf{\color{ink}361.}
\par
Standart ASCII jadvalida o'nlik kodi 76 bo'lgan belgi xotirada 8 bitli ko'rinishda qanday yoziladi?
\begin{enumerate}[label=\Alph*),leftmargin=9mm]
\item 01001100
\item 01001011
\item 01001101
\item 10011000
\end{enumerate}
\begin{tcolorbox}[enhanced,breakable,colback=teal!4,colframe=teal!50!black,boxrule=0.5pt,arc=1.1mm,left=3mm,right=3mm,top=1.7mm,bottom=1.7mm,before skip=2mm,after skip=2mm]
\textbf{\color{teal}To'g'ri javob:} \textbf{A) 01001100}
\par\smallskip
\textbf{Izoh:} 76\textsubscript{1}\textsubscript{0}=01001100\textsubscript{2}; bu standart ASCII dagi \texttt{L} harfi.
\end{tcolorbox}
\vspace{2mm}\hrule height 0.2pt \vspace{3mm}

% DB question_id: 3fd2b32c-c89f-482e-9f79-9a066b209818
\needspace{5\baselineskip}
\noindent\textbf{\color{ink}362.}
\par
"Kodlash va shifrlash farqi" atamasiga eng aniq ta'rifni toping.
\begin{enumerate}[label=\Alph*),leftmargin=9mm]
\item bit ikkilik birlik, bayt esa 8 bitdan iborat ekanini ajratish
\item rastr piksel, vektor geometrik obyektlarga asoslanishini ajratish
\item kodek oqimni kodlashi, konteyner esa oqimlarni bir faylda saqlashini ajratish
\item kodlash taqdim etish formatini, shifrlash esa maxfiylikni ko'zlashini ajratish
\end{enumerate}
\begin{tcolorbox}[enhanced,breakable,colback=teal!4,colframe=teal!50!black,boxrule=0.5pt,arc=1.1mm,left=3mm,right=3mm,top=1.7mm,bottom=1.7mm,before skip=2mm,after skip=2mm]
\textbf{\color{teal}To'g'ri javob:} \textbf{D) kodlash taqdim etish formatini, shifrlash esa maxfiylikni ko'zlashini ajratish}
\par\smallskip
\textbf{Izoh:} Kodlash va shifrlash farqi - kodlash taqdim etish formatini, shifrlash esa maxfiylikni ko'zlashini ajratish.
\end{tcolorbox}
\vspace{2mm}\hrule height 0.2pt \vspace{3mm}

% DB question_id: 444c940c-bee8-5f1d-982b-8bcc555f8d35
\needspace{5\baselineskip}
\noindent\textbf{\color{ink}363.}
\par
26 harfli ingliz alifbosida Sezar shifri bilan harflar 2 o'rin oldinga siljitiladi. \texttt{GOLL} so'zining shifrmatnini toping.
\begin{enumerate}[label=\Alph*),leftmargin=9mm]
\item IQNN
\item HPMM
\item JROO
\item EMJJ
\end{enumerate}
\begin{tcolorbox}[enhanced,breakable,colback=teal!4,colframe=teal!50!black,boxrule=0.5pt,arc=1.1mm,left=3mm,right=3mm,top=1.7mm,bottom=1.7mm,before skip=2mm,after skip=2mm]
\textbf{\color{teal}To'g'ri javob:} \textbf{A) IQNN}
\par\smallskip
\textbf{Izoh:} G\ensuremath{\to}I, O\ensuremath{\to}Q, L\ensuremath{\to}N, L\ensuremath{\to}N; shifrmatn IQNN.
\end{tcolorbox}
\vspace{2mm}\hrule height 0.2pt \vspace{3mm}

% DB question_id: 44f47f17-3f0e-4cda-b054-1c4c6b6d2585
\needspace{5\baselineskip}
\noindent\textbf{\color{ink}364.}
\par
"Bit" atamasiga eng aniq ta'rifni toping.
\begin{enumerate}[label=\Alph*),leftmargin=9mm]
\item sakkiz bitdan tashkil topgan axborot birligi
\item teng ehtimolli N holatdan birini kodlash uchun i = log\textsubscript{2}N munosabatini beradigan formula
\item ikki teng ehtimolli holatdan birini ajratish uchun yetarli eng kichik axborot birligi
\item kodlashda qatnashadigan turli belgilar yoki holatlar soni
\end{enumerate}
\begin{tcolorbox}[enhanced,breakable,colback=teal!4,colframe=teal!50!black,boxrule=0.5pt,arc=1.1mm,left=3mm,right=3mm,top=1.7mm,bottom=1.7mm,before skip=2mm,after skip=2mm]
\textbf{\color{teal}To'g'ri javob:} \textbf{C) ikki teng ehtimolli holatdan birini ajratish uchun yetarli eng kichik axborot birligi}
\par\smallskip
\textbf{Izoh:} Bit - ikki teng ehtimolli holatdan birini ajratish uchun yetarli eng kichik axborot birligi.
\end{tcolorbox}
\vspace{2mm}\hrule height 0.2pt \vspace{3mm}

% DB question_id: 4521a0cb-8366-413c-b57f-89c9dc5ed06d
\needspace{5\baselineskip}
\noindent\textbf{\color{ink}365.}
\par
Quyidagi ta'rif qaysi atamaga tegishli? Sakkiz bitdan tashkil topgan axborot birligi.
\begin{enumerate}[label=\Alph*),leftmargin=9mm]
\item Xartli formulasi
\item Alifbo quvvati
\item Minimal butun kod uzunligi
\item Bayt
\end{enumerate}
\begin{tcolorbox}[enhanced,breakable,colback=teal!4,colframe=teal!50!black,boxrule=0.5pt,arc=1.1mm,left=3mm,right=3mm,top=1.7mm,bottom=1.7mm,before skip=2mm,after skip=2mm]
\textbf{\color{teal}To'g'ri javob:} \textbf{D) Bayt}
\par\smallskip
\textbf{Izoh:} Ta'rif Bayt tushunchasiga tegishli.
\end{tcolorbox}
\vspace{2mm}\hrule height 0.2pt \vspace{3mm}

% DB question_id: 56d907ff-c587-46d7-a8d6-dcaed9c26939
\needspace{5\baselineskip}
\noindent\textbf{\color{ink}366.}
\par
"Doimiy uzunlikdagi kod" atamasiga eng aniq ta'rifni toping.
\begin{enumerate}[label=\Alph*),leftmargin=9mm]
\item turli belgilar turli uzunlikdagi kodlar bilan ifodalanadigan kod
\item har bir belgi bir xil sondagi kod elementi bilan ifodalanadigan kod
\item axborotni saqlash yoki uzatish uchun fayl hajmini kamaytirish jarayoni
\item mazmunni kalitsiz o'qib bo'lmaydigan ko'rinishga aylantirish jarayoni
\end{enumerate}
\begin{tcolorbox}[enhanced,breakable,colback=teal!4,colframe=teal!50!black,boxrule=0.5pt,arc=1.1mm,left=3mm,right=3mm,top=1.7mm,bottom=1.7mm,before skip=2mm,after skip=2mm]
\textbf{\color{teal}To'g'ri javob:} \textbf{B) har bir belgi bir xil sondagi kod elementi bilan ifodalanadigan kod}
\par\smallskip
\textbf{Izoh:} Doimiy uzunlikdagi kod - har bir belgi bir xil sondagi kod elementi bilan ifodalanadigan kod.
\end{tcolorbox}
\vspace{2mm}\hrule height 0.2pt \vspace{3mm}

% DB question_id: 71540b0a-0736-4b9c-8e0d-cbcec536bd09
\needspace{5\baselineskip}
\noindent\textbf{\color{ink}367.}
\par
"kB" atamasiga eng aniq ta'rifni toping.
\begin{enumerate}[label=\Alph*),leftmargin=9mm]
\item IEC standartida 1024 baytga teng ikkilik birlik
\item tashuvchida saqlanishi mumkin bo'lgan axborotning umumiy hajmi
\item ikki teng ehtimolli holatdan birini ajratish uchun yetarli eng kichik axborot birligi
\item SI standartida 1000 baytga teng o'nlik birlik
\end{enumerate}
\begin{tcolorbox}[enhanced,breakable,colback=teal!4,colframe=teal!50!black,boxrule=0.5pt,arc=1.1mm,left=3mm,right=3mm,top=1.7mm,bottom=1.7mm,before skip=2mm,after skip=2mm]
\textbf{\color{teal}To'g'ri javob:} \textbf{D) SI standartida 1000 baytga teng o'nlik birlik}
\par\smallskip
\textbf{Izoh:} kB - SI standartida 1000 baytga teng o'nlik birlik.
\end{tcolorbox}
\vspace{2mm}\hrule height 0.2pt \vspace{3mm}

% DB question_id: 7585f7ab-3517-5d5c-aae7-e034b60a50c8
\needspace{5\baselineskip}
\noindent\textbf{\color{ink}368.}
\par
IBM PC CP437 8 bitli kod sahifasidagi oxirgi kod 255\textsubscript{1}\textsubscript{0}. Uning o'n oltilik yozuvini toping.
\begin{enumerate}[label=\Alph*),leftmargin=9mm]
\item FF\textsubscript{1}\textsubscript{6}
\item 7F\textsubscript{1}\textsubscript{6}
\item FE\textsubscript{1}\textsubscript{6}
\item 100\textsubscript{1}\textsubscript{6}
\end{enumerate}
\begin{tcolorbox}[enhanced,breakable,colback=teal!4,colframe=teal!50!black,boxrule=0.5pt,arc=1.1mm,left=3mm,right=3mm,top=1.7mm,bottom=1.7mm,before skip=2mm,after skip=2mm]
\textbf{\color{teal}To'g'ri javob:} \textbf{A) FF\textsubscript{1}\textsubscript{6}}
\par\smallskip
\textbf{Izoh:} 255=15\ensuremath{\cdot}16+15; ikkala raqam ham F, shuning uchun FF\textsubscript{1}\textsubscript{6}.
\end{tcolorbox}
\vspace{2mm}\hrule height 0.2pt \vspace{3mm}

% DB question_id: 7812f249-b8ee-498d-971b-67a41067bf71
\needspace{5\baselineskip}
\noindent\textbf{\color{ink}369.}
\par
"Alifbo" atamasiga eng aniq ta'rifni toping.
\begin{enumerate}[label=\Alph*),leftmargin=9mm]
\item axborot elementlarini boshqa belgilar yoki signallar bilan moslashtirish qoidalari tizimi
\item kodlashda foydalaniladigan ruxsat etilgan belgilar to'plami
\item axborotni belgilangan qoida bo'yicha boshqa ko'rinishga o'tkazish
\item kodlangan axborotni kelishilgan qoida asosida qayta tiklash
\end{enumerate}
\begin{tcolorbox}[enhanced,breakable,colback=teal!4,colframe=teal!50!black,boxrule=0.5pt,arc=1.1mm,left=3mm,right=3mm,top=1.7mm,bottom=1.7mm,before skip=2mm,after skip=2mm]
\textbf{\color{teal}To'g'ri javob:} \textbf{B) kodlashda foydalaniladigan ruxsat etilgan belgilar to'plami}
\par\smallskip
\textbf{Izoh:} Alifbo - kodlashda foydalaniladigan ruxsat etilgan belgilar to'plami.
\end{tcolorbox}
\vspace{2mm}\hrule height 0.2pt \vspace{3mm}

% DB question_id: 7e87067f-d49e-43b8-8dea-44df2616d0cd
\needspace{5\baselineskip}
\noindent\textbf{\color{ink}370.}
\par
"Bit va bayt farqi" atamasiga eng aniq ta'rifni toping.
\begin{enumerate}[label=\Alph*),leftmargin=9mm]
\item bit ikkilik birlik, bayt esa 8 bitdan iborat ekanini ajratish
\item rastr piksel, vektor geometrik obyektlarga asoslanishini ajratish
\item kodek oqimni kodlashi, konteyner esa oqimlarni bir faylda saqlashini ajratish
\item validatsiya qoidaga moslikni, verifikatsiya aslga moslikni tekshirishini ajratish
\end{enumerate}
\begin{tcolorbox}[enhanced,breakable,colback=teal!4,colframe=teal!50!black,boxrule=0.5pt,arc=1.1mm,left=3mm,right=3mm,top=1.7mm,bottom=1.7mm,before skip=2mm,after skip=2mm]
\textbf{\color{teal}To'g'ri javob:} \textbf{A) bit ikkilik birlik, bayt esa 8 bitdan iborat ekanini ajratish}
\par\smallskip
\textbf{Izoh:} Bit va bayt farqi - bit ikkilik birlik, bayt esa 8 bitdan iborat ekanini ajratish.
\end{tcolorbox}
\vspace{2mm}\hrule height 0.2pt \vspace{3mm}

% DB question_id: 87f1046b-7b2d-4e57-8c7c-b2c5d867f792
\needspace{5\baselineskip}
\noindent\textbf{\color{ink}371.}
\par
"kB va KiB farqi" atamasiga eng aniq ta'rifni toping.
\begin{enumerate}[label=\Alph*),leftmargin=9mm]
\item baud simvol tezligi, bit/s esa axborot bitlari tezligi ekanini ajratish
\item kB 1000 bayt, KiB 1024 bayt ekanini standart belgilari bilan ajratish
\item HTTPS transportni shifrlashi, sahifa mazmunining haqiqatligini kafolatlamasligini ajratish
\item ma'lumot qayd etilgan qiymat, axborot esa qabul qiluvchi uchun mazmunli ma'lumot ekanini ajratish
\end{enumerate}
\begin{tcolorbox}[enhanced,breakable,colback=teal!4,colframe=teal!50!black,boxrule=0.5pt,arc=1.1mm,left=3mm,right=3mm,top=1.7mm,bottom=1.7mm,before skip=2mm,after skip=2mm]
\textbf{\color{teal}To'g'ri javob:} \textbf{B) kB 1000 bayt, KiB 1024 bayt ekanini standart belgilari bilan ajratish}
\par\smallskip
\textbf{Izoh:} kB va KiB farqi - kB 1000 bayt, KiB 1024 bayt ekanini standart belgilari bilan ajratish.
\end{tcolorbox}
\vspace{2mm}\hrule height 0.2pt \vspace{3mm}

% DB question_id: 898fa073-a621-52b7-9dc2-b20046da5140
\needspace{5\baselineskip}
\noindent\textbf{\color{ink}372.}
\par
Video 5 daqiqa davom etadi va sekundiga 25 kadr yoziladi. Videodagi jami kadrlar sonini toping.
\begin{enumerate}[label=\Alph*),leftmargin=9mm]
\item 7500
\item 125
\item 1500
\item 12 500
\end{enumerate}
\begin{tcolorbox}[enhanced,breakable,colback=teal!4,colframe=teal!50!black,boxrule=0.5pt,arc=1.1mm,left=3mm,right=3mm,top=1.7mm,bottom=1.7mm,before skip=2mm,after skip=2mm]
\textbf{\color{teal}To'g'ri javob:} \textbf{A) 7500}
\par\smallskip
\textbf{Izoh:} 5 daqiqa=300 sekund. 300\ensuremath{\cdot}25=7500 kadr.
\end{tcolorbox}
\vspace{2mm}\hrule height 0.2pt \vspace{3mm}

% DB question_id: 8cbdfc9e-0430-489b-88d2-c5987403c7e3
\needspace{5\baselineskip}
\noindent\textbf{\color{ink}373.}
\par
"Kod" atamasiga eng aniq ta'rifni toping.
\begin{enumerate}[label=\Alph*),leftmargin=9mm]
\item axborotni belgilangan qoida bo'yicha boshqa ko'rinishga o'tkazish
\item kodlangan axborotni kelishilgan qoida asosida qayta tiklash
\item axborot elementlarini boshqa belgilar yoki signallar bilan moslashtirish qoidalari tizimi
\item har bir belgi bir xil sondagi kod elementi bilan ifodalanadigan kod
\end{enumerate}
\begin{tcolorbox}[enhanced,breakable,colback=teal!4,colframe=teal!50!black,boxrule=0.5pt,arc=1.1mm,left=3mm,right=3mm,top=1.7mm,bottom=1.7mm,before skip=2mm,after skip=2mm]
\textbf{\color{teal}To'g'ri javob:} \textbf{C) axborot elementlarini boshqa belgilar yoki signallar bilan moslashtirish qoidalari tizimi}
\par\smallskip
\textbf{Izoh:} Kod - axborot elementlarini boshqa belgilar yoki signallar bilan moslashtirish qoidalari tizimi.
\end{tcolorbox}
\vspace{2mm}\hrule height 0.2pt \vspace{3mm}

% DB question_id: 99e148e3-b04b-43ca-96d2-c22218e66466
\needspace{5\baselineskip}
\noindent\textbf{\color{ink}374.}
\par
Quyidagi ta'rif qaysi atamaga tegishli? Ikki teng ehtimolli holatdan birini ajratish uchun yetarli eng kichik axborot birligi.
\begin{enumerate}[label=\Alph*),leftmargin=9mm]
\item Bayt
\item Xartli formulasi
\item Bit
\item Alifbo quvvati
\end{enumerate}
\begin{tcolorbox}[enhanced,breakable,colback=teal!4,colframe=teal!50!black,boxrule=0.5pt,arc=1.1mm,left=3mm,right=3mm,top=1.7mm,bottom=1.7mm,before skip=2mm,after skip=2mm]
\textbf{\color{teal}To'g'ri javob:} \textbf{C) Bit}
\par\smallskip
\textbf{Izoh:} Ta'rif Bit tushunchasiga tegishli.
\end{tcolorbox}
\vspace{2mm}\hrule height 0.2pt \vspace{3mm}

% DB question_id: 9f0d12fd-1305-52da-bf04-0167994f30cd
\needspace{5\baselineskip}
\noindent\textbf{\color{ink}375.}
\par
Standart ASCII da \texttt{A} harfining o'nlik kodi 65. Lotin bosh harflari ketma-ket kodlangan bo'lsa, \texttt{F} harfining o'nlik kodi nechaga teng?
\begin{enumerate}[label=\Alph*),leftmargin=9mm]
\item 70
\item 69
\item 71
\item 75
\end{enumerate}
\begin{tcolorbox}[enhanced,breakable,colback=teal!4,colframe=teal!50!black,boxrule=0.5pt,arc=1.1mm,left=3mm,right=3mm,top=1.7mm,bottom=1.7mm,before skip=2mm,after skip=2mm]
\textbf{\color{teal}To'g'ri javob:} \textbf{A) 70}
\par\smallskip
\textbf{Izoh:} A=65, B=66, C=67, D=68, E=69, F=70.
\end{tcolorbox}
\vspace{2mm}\hrule height 0.2pt \vspace{3mm}

% DB question_id: a5fd48cb-8f00-4df5-bce8-5e5206eb7419
\needspace{5\baselineskip}
\noindent\textbf{\color{ink}376.}
\par
"KiB" atamasiga eng aniq ta'rifni toping.
\begin{enumerate}[label=\Alph*),leftmargin=9mm]
\item IEC standartida 1024 baytga teng ikkilik birlik
\item tashuvchida saqlanishi mumkin bo'lgan axborotning umumiy hajmi
\item ikki teng ehtimolli holatdan birini ajratish uchun yetarli eng kichik axborot birligi
\item sakkiz bitdan tashkil topgan axborot birligi
\end{enumerate}
\begin{tcolorbox}[enhanced,breakable,colback=teal!4,colframe=teal!50!black,boxrule=0.5pt,arc=1.1mm,left=3mm,right=3mm,top=1.7mm,bottom=1.7mm,before skip=2mm,after skip=2mm]
\textbf{\color{teal}To'g'ri javob:} \textbf{A) IEC standartida 1024 baytga teng ikkilik birlik}
\par\smallskip
\textbf{Izoh:} KiB - IEC standartida 1024 baytga teng ikkilik birlik.
\end{tcolorbox}
\vspace{2mm}\hrule height 0.2pt \vspace{3mm}

% DB question_id: b2a0be85-d2f1-4b9a-93dd-01c9e72d5180
\needspace{5\baselineskip}
\noindent\textbf{\color{ink}377.}
\par
"Bayt" atamasiga eng aniq ta'rifni toping.
\begin{enumerate}[label=\Alph*),leftmargin=9mm]
\item teng ehtimolli N holatdan birini kodlash uchun i = log\textsubscript{2}N munosabatini beradigan formula
\item kodlashda qatnashadigan turli belgilar yoki holatlar soni
\item N holat uchun log\textsubscript{2}N yuqoriga yaxlitlangan bitlar soni
\item sakkiz bitdan tashkil topgan axborot birligi
\end{enumerate}
\begin{tcolorbox}[enhanced,breakable,colback=teal!4,colframe=teal!50!black,boxrule=0.5pt,arc=1.1mm,left=3mm,right=3mm,top=1.7mm,bottom=1.7mm,before skip=2mm,after skip=2mm]
\textbf{\color{teal}To'g'ri javob:} \textbf{D) sakkiz bitdan tashkil topgan axborot birligi}
\par\smallskip
\textbf{Izoh:} Bayt - sakkiz bitdan tashkil topgan axborot birligi.
\end{tcolorbox}
\vspace{2mm}\hrule height 0.2pt \vspace{3mm}

% DB question_id: be3c4799-0300-45de-8f46-3581fbf3fbd4
\needspace{5\baselineskip}
\noindent\textbf{\color{ink}378.}
\par
"Siqish" tushunchasining "Belgi, kod, kodlash va shifrlash" bobidagi eng aniq ta'rifini toping.
\begin{enumerate}[label=\Alph*),leftmargin=9mm]
\item mazmunni kalitsiz o'qib bo'lmaydigan ko'rinishga aylantirish jarayoni
\item ochiq va yopiq kalitlar juftidan foydalanadigan shifrlash usuli
\item obyekt, xossa yoki tushunchani ifodalash uchun qabul qilingan ishora
\item axborotni saqlash yoki uzatish uchun fayl hajmini kamaytirish jarayoni
\end{enumerate}
\begin{tcolorbox}[enhanced,breakable,colback=teal!4,colframe=teal!50!black,boxrule=0.5pt,arc=1.1mm,left=3mm,right=3mm,top=1.7mm,bottom=1.7mm,before skip=2mm,after skip=2mm]
\textbf{\color{teal}To'g'ri javob:} \textbf{D) axborotni saqlash yoki uzatish uchun fayl hajmini kamaytirish jarayoni}
\par\smallskip
\textbf{Izoh:} Siqish - axborotni saqlash yoki uzatish uchun fayl hajmini kamaytirish jarayoni.
\end{tcolorbox}
\vspace{2mm}\hrule height 0.2pt \vspace{3mm}

% DB question_id: c7712f88-3990-407a-bc9c-a94277b77e94
\needspace{5\baselineskip}
\noindent\textbf{\color{ink}379.}
\par
"Kodlash" atamasiga eng aniq ta'rifni toping.
\begin{enumerate}[label=\Alph*),leftmargin=9mm]
\item kodlangan axborotni kelishilgan qoida asosida qayta tiklash
\item har bir belgi bir xil sondagi kod elementi bilan ifodalanadigan kod
\item turli belgilar turli uzunlikdagi kodlar bilan ifodalanadigan kod
\item axborotni belgilangan qoida bo'yicha boshqa ko'rinishga o'tkazish
\end{enumerate}
\begin{tcolorbox}[enhanced,breakable,colback=teal!4,colframe=teal!50!black,boxrule=0.5pt,arc=1.1mm,left=3mm,right=3mm,top=1.7mm,bottom=1.7mm,before skip=2mm,after skip=2mm]
\textbf{\color{teal}To'g'ri javob:} \textbf{D) axborotni belgilangan qoida bo'yicha boshqa ko'rinishga o'tkazish}
\par\smallskip
\textbf{Izoh:} Kodlash - axborotni belgilangan qoida bo'yicha boshqa ko'rinishga o'tkazish.
\end{tcolorbox}
\vspace{2mm}\hrule height 0.2pt \vspace{3mm}

% DB question_id: cb0364d9-03ce-5f4a-ab21-8ae118706453
\needspace{5\baselineskip}
\noindent\textbf{\color{ink}380.}
\par
Manbadagi IBM PC CP437 kengaytirilgan 8 bitli kod sahifasida \texttt{\textcopyright{}} belgisi qaysi o'nlik kodga ega?
\begin{enumerate}[label=\Alph*),leftmargin=9mm]
\item 184
\item 169
\item 127
\item 255
\end{enumerate}
\begin{tcolorbox}[enhanced,breakable,colback=teal!4,colframe=teal!50!black,boxrule=0.5pt,arc=1.1mm,left=3mm,right=3mm,top=1.7mm,bottom=1.7mm,before skip=2mm,after skip=2mm]
\textbf{\color{teal}To'g'ri javob:} \textbf{A) 184}
\par\smallskip
\textbf{Izoh:} Berilgan jadval CP437 tartibiga mos: unda \textcopyright{} belgisi 184-kodda. Bu standart 7 bitli ASCII emas; kengaytirilgan kod sahifasidir.
\end{tcolorbox}
\vspace{2mm}\hrule height 0.2pt \vspace{3mm}

% DB question_id: cd121829-21ca-44fa-a8ce-aa279457884c
\needspace{5\baselineskip}
\noindent\textbf{\color{ink}381.}
\par
"O'zgaruvchan uzunlikdagi kod" atamasiga eng aniq ta'rifni toping.
\begin{enumerate}[label=\Alph*),leftmargin=9mm]
\item axborotni saqlash yoki uzatish uchun fayl hajmini kamaytirish jarayoni
\item mazmunni kalitsiz o'qib bo'lmaydigan ko'rinishga aylantirish jarayoni
\item turli belgilar turli uzunlikdagi kodlar bilan ifodalanadigan kod
\item ochiq va yopiq kalitlar juftidan foydalanadigan shifrlash usuli
\end{enumerate}
\begin{tcolorbox}[enhanced,breakable,colback=teal!4,colframe=teal!50!black,boxrule=0.5pt,arc=1.1mm,left=3mm,right=3mm,top=1.7mm,bottom=1.7mm,before skip=2mm,after skip=2mm]
\textbf{\color{teal}To'g'ri javob:} \textbf{C) turli belgilar turli uzunlikdagi kodlar bilan ifodalanadigan kod}
\par\smallskip
\textbf{Izoh:} O'zgaruvchan uzunlikdagi kod - turli belgilar turli uzunlikdagi kodlar bilan ifodalanadigan kod.
\end{tcolorbox}
\vspace{2mm}\hrule height 0.2pt \vspace{3mm}

% DB question_id: cdd3a0e0-618c-425b-b040-3540b414831e
\needspace{5\baselineskip}
\noindent\textbf{\color{ink}382.}
\par
"Kodek va konteyner farqi" atamasiga eng aniq ta'rifni toping.
\begin{enumerate}[label=\Alph*),leftmargin=9mm]
\item validatsiya qoidaga moslikni, verifikatsiya aslga moslikni tekshirishini ajratish
\item Internet tarmoqlar infratuzilmasi, WWW esa undagi xizmatlardan biri ekanini ajratish
\item kodek oqimni kodlashi, konteyner esa oqimlarni bir faylda saqlashini ajratish
\item kB 1000 bayt, KiB 1024 bayt ekanini standart belgilari bilan ajratish
\end{enumerate}
\begin{tcolorbox}[enhanced,breakable,colback=teal!4,colframe=teal!50!black,boxrule=0.5pt,arc=1.1mm,left=3mm,right=3mm,top=1.7mm,bottom=1.7mm,before skip=2mm,after skip=2mm]
\textbf{\color{teal}To'g'ri javob:} \textbf{C) kodek oqimni kodlashi, konteyner esa oqimlarni bir faylda saqlashini ajratish}
\par\smallskip
\textbf{Izoh:} Kodek va konteyner farqi - kodek oqimni kodlashi, konteyner esa oqimlarni bir faylda saqlashini ajratish.
\end{tcolorbox}
\vspace{2mm}\hrule height 0.2pt \vspace{3mm}

% DB question_id: e6ff54c1-d0d2-4506-8d0d-ae0798a7511c
\needspace{5\baselineskip}
\noindent\textbf{\color{ink}383.}
\par
"Xotira sig'imi" atamasiga eng aniq ta'rifni toping.
\begin{enumerate}[label=\Alph*),leftmargin=9mm]
\item ikki teng ehtimolli holatdan birini ajratish uchun yetarli eng kichik axborot birligi
\item tashuvchida saqlanishi mumkin bo'lgan axborotning umumiy hajmi
\item sakkiz bitdan tashkil topgan axborot birligi
\item teng ehtimolli N holatdan birini kodlash uchun i = log\textsubscript{2}N munosabatini beradigan formula
\end{enumerate}
\begin{tcolorbox}[enhanced,breakable,colback=teal!4,colframe=teal!50!black,boxrule=0.5pt,arc=1.1mm,left=3mm,right=3mm,top=1.7mm,bottom=1.7mm,before skip=2mm,after skip=2mm]
\textbf{\color{teal}To'g'ri javob:} \textbf{B) tashuvchida saqlanishi mumkin bo'lgan axborotning umumiy hajmi}
\par\smallskip
\textbf{Izoh:} Xotira sig'imi - tashuvchida saqlanishi mumkin bo'lgan axborotning umumiy hajmi.
\end{tcolorbox}
\vspace{2mm}\hrule height 0.2pt \vspace{3mm}

% DB question_id: fd9832ab-5f66-4fd0-9c0f-4ef6ecafeec0
\needspace{5\baselineskip}
\noindent\textbf{\color{ink}384.}
\par
"Belgi" atamasiga eng aniq ta'rifni toping.
\begin{enumerate}[label=\Alph*),leftmargin=9mm]
\item obyekt, xossa yoki tushunchani ifodalash uchun qabul qilingan ishora
\item kodlashda foydalaniladigan ruxsat etilgan belgilar to'plami
\item axborot elementlarini boshqa belgilar yoki signallar bilan moslashtirish qoidalari tizimi
\item axborotni belgilangan qoida bo'yicha boshqa ko'rinishga o'tkazish
\end{enumerate}
\begin{tcolorbox}[enhanced,breakable,colback=teal!4,colframe=teal!50!black,boxrule=0.5pt,arc=1.1mm,left=3mm,right=3mm,top=1.7mm,bottom=1.7mm,before skip=2mm,after skip=2mm]
\textbf{\color{teal}To'g'ri javob:} \textbf{A) obyekt, xossa yoki tushunchani ifodalash uchun qabul qilingan ishora}
\par\smallskip
\textbf{Izoh:} Belgi - obyekt, xossa yoki tushunchani ifodalash uchun qabul qilingan ishora.
\end{tcolorbox}
\vspace{2mm}\hrule height 0.2pt \vspace{3mm}

% DB question_id: 0277a43a-1e95-5653-8308-a7d8e1b686ae
\needspace{5\baselineskip}
\noindent\textbf{\color{ink}385.}
\par
1856 KiB axborot necha baytga teng? 1 KiB = 1024 bayt deb oling.
\begin{enumerate}[label=\Alph*),leftmargin=9mm]
\item 1900544 bayt
\item 15204352 bayt
\item 1856000 bayt
\item 237568 bayt
\end{enumerate}
\begin{tcolorbox}[enhanced,breakable,colback=teal!4,colframe=teal!50!black,boxrule=0.5pt,arc=1.1mm,left=3mm,right=3mm,top=1.7mm,bottom=1.7mm,before skip=2mm,after skip=2mm]
\textbf{\color{teal}To'g'ri javob:} \textbf{A) 1900544 bayt}
\par\smallskip
\textbf{Izoh:} 1856 \ensuremath{\times} 1024 = 1900544 bayt.
\end{tcolorbox}
\vspace{2mm}\hrule height 0.2pt \vspace{3mm}

% DB question_id: 073627af-81b4-52ae-bce1-66d6e8f5123b
\needspace{5\baselineskip}
\noindent\textbf{\color{ink}386.}
\par
17414 KiB axborot necha baytga teng? 1 KiB = 1024 bayt deb oling.
\begin{enumerate}[label=\Alph*),leftmargin=9mm]
\item 142655488 bayt
\item 17831936 bayt
\item 17414000 bayt
\item 2228992 bayt
\end{enumerate}
\begin{tcolorbox}[enhanced,breakable,colback=teal!4,colframe=teal!50!black,boxrule=0.5pt,arc=1.1mm,left=3mm,right=3mm,top=1.7mm,bottom=1.7mm,before skip=2mm,after skip=2mm]
\textbf{\color{teal}To'g'ri javob:} \textbf{B) 17831936 bayt}
\par\smallskip
\textbf{Izoh:} 17414 \ensuremath{\times} 1024 = 17831936 bayt.
\end{tcolorbox}
\vspace{2mm}\hrule height 0.2pt \vspace{3mm}

% DB question_id: 07b66eb3-e133-5aa8-b94e-9c696469eb49
\needspace{5\baselineskip}
\noindent\textbf{\color{ink}387.}
\par
42545 KiB axborot necha baytga teng? 1 KiB = 1024 bayt deb oling.
\begin{enumerate}[label=\Alph*),leftmargin=9mm]
\item 5445760 bayt
\item 43566080 bayt
\item 348528640 bayt
\item 42545000 bayt
\end{enumerate}
\begin{tcolorbox}[enhanced,breakable,colback=teal!4,colframe=teal!50!black,boxrule=0.5pt,arc=1.1mm,left=3mm,right=3mm,top=1.7mm,bottom=1.7mm,before skip=2mm,after skip=2mm]
\textbf{\color{teal}To'g'ri javob:} \textbf{B) 43566080 bayt}
\par\smallskip
\textbf{Izoh:} 42545 \ensuremath{\times} 1024 = 43566080 bayt.
\end{tcolorbox}
\vspace{2mm}\hrule height 0.2pt \vspace{3mm}

% DB question_id: 096726a2-3567-4e10-838f-0eaa0bbab1b5
\needspace{5\baselineskip}
\noindent\textbf{\color{ink}388.}
\par
Quyidagi ta'rif qaysi atamaga tegishli? Axborot elementlarini boshqa belgilar yoki signallar bilan moslashtirish qoidalari tizimi.
\begin{enumerate}[label=\Alph*),leftmargin=9mm]
\item Kod
\item Kodlash
\item Dekodlash
\item Doimiy uzunlikdagi kod
\end{enumerate}
\begin{tcolorbox}[enhanced,breakable,colback=teal!4,colframe=teal!50!black,boxrule=0.5pt,arc=1.1mm,left=3mm,right=3mm,top=1.7mm,bottom=1.7mm,before skip=2mm,after skip=2mm]
\textbf{\color{teal}To'g'ri javob:} \textbf{A) Kod}
\par\smallskip
\textbf{Izoh:} Ta'rif Kod tushunchasiga tegishli.
\end{tcolorbox}
\vspace{2mm}\hrule height 0.2pt \vspace{3mm}

% DB question_id: 0b81ea98-aaaa-59fd-8b8c-6392863b4f19
\needspace{5\baselineskip}
\noindent\textbf{\color{ink}389.}
\par
850 ta o'rindiqning har biriga takrorlanmaydigan, bir xil uzunlikdagi ikkilik kod berish uchun kodning minimal uzunligi necha bit bo'lishi kerak?
\begin{enumerate}[label=\Alph*),leftmargin=9mm]
\item 9 bit
\item 10 bit
\item 11 bit
\item 12 bit
\end{enumerate}
\begin{tcolorbox}[enhanced,breakable,colback=teal!4,colframe=teal!50!black,boxrule=0.5pt,arc=1.1mm,left=3mm,right=3mm,top=1.7mm,bottom=1.7mm,before skip=2mm,after skip=2mm]
\textbf{\color{teal}To'g'ri javob:} \textbf{B) 10 bit}
\par\smallskip
\textbf{Izoh:} 2\textasciicircum{}9=512\textless{}850\ensuremath{\leq}1024=2\textasciicircum{}10, demak kamida 10 bit kerak.
\end{tcolorbox}
\vspace{2mm}\hrule height 0.2pt \vspace{3mm}

% DB question_id: 0ff616ac-be0d-5c95-bab0-0763ba6af4bc
\needspace{5\baselineskip}
\noindent\textbf{\color{ink}390.}
\par
A qabila 32 belgili alifboda 80 belgili, B qabila 64 belgili alifboda 60 belgili xabar yozdi. Har bir alifboda belgilar bir xil minimal uzunlikdagi ikkilik kod bilan kodlansa, xabarlar hajmi farqi necha bit?
\begin{enumerate}[label=\Alph*),leftmargin=9mm]
\item 40 bit
\item 760 bit
\item 5 bit
\item 360 bit
\end{enumerate}
\begin{tcolorbox}[enhanced,breakable,colback=teal!4,colframe=teal!50!black,boxrule=0.5pt,arc=1.1mm,left=3mm,right=3mm,top=1.7mm,bottom=1.7mm,before skip=2mm,after skip=2mm]
\textbf{\color{teal}To'g'ri javob:} \textbf{A) 40 bit}
\par\smallskip
\textbf{Izoh:} A: log\textsubscript{2}32=5 bit/belgi, 80\ensuremath{\cdot}5=400 bit. B: log\textsubscript{2}64=6 bit/belgi, 60\ensuremath{\cdot}6=360 bit. Farq 400-360=40 bit.
\end{tcolorbox}
\vspace{2mm}\hrule height 0.2pt \vspace{3mm}

% DB question_id: 1018a9f9-3ac6-5dfc-a27e-b68be8efd6a1
\needspace{5\baselineskip}
\noindent\textbf{\color{ink}391.}
\par
65107 KiB axborot necha baytga teng? 1 KiB = 1024 bayt deb oling.
\begin{enumerate}[label=\Alph*),leftmargin=9mm]
\item 65107000 bayt
\item 66669568 bayt
\item 8333696 bayt
\item 533356544 bayt
\end{enumerate}
\begin{tcolorbox}[enhanced,breakable,colback=teal!4,colframe=teal!50!black,boxrule=0.5pt,arc=1.1mm,left=3mm,right=3mm,top=1.7mm,bottom=1.7mm,before skip=2mm,after skip=2mm]
\textbf{\color{teal}To'g'ri javob:} \textbf{B) 66669568 bayt}
\par\smallskip
\textbf{Izoh:} 65107 \ensuremath{\times} 1024 = 66669568 bayt.
\end{tcolorbox}
\vspace{2mm}\hrule height 0.2pt \vspace{3mm}

% DB question_id: 1216a8f9-7c27-53fa-ab22-845ff54f7b6c
\needspace{5\baselineskip}
\noindent\textbf{\color{ink}392.}
\par
55821 KiB axborot necha baytga teng? 1 KiB = 1024 bayt deb oling.
\begin{enumerate}[label=\Alph*),leftmargin=9mm]
\item 457285632 bayt
\item 57160704 bayt
\item 7145088 bayt
\item 55821000 bayt
\end{enumerate}
\begin{tcolorbox}[enhanced,breakable,colback=teal!4,colframe=teal!50!black,boxrule=0.5pt,arc=1.1mm,left=3mm,right=3mm,top=1.7mm,bottom=1.7mm,before skip=2mm,after skip=2mm]
\textbf{\color{teal}To'g'ri javob:} \textbf{B) 57160704 bayt}
\par\smallskip
\textbf{Izoh:} 55821 \ensuremath{\times} 1024 = 57160704 bayt.
\end{tcolorbox}
\vspace{2mm}\hrule height 0.2pt \vspace{3mm}

% DB question_id: 14fc9bcf-c1c9-50ea-8bed-7a215752c446
\needspace{5\baselineskip}
\noindent\textbf{\color{ink}393.}
\par
17162 KiB axborot necha baytga teng? 1 KiB = 1024 bayt deb oling.
\begin{enumerate}[label=\Alph*),leftmargin=9mm]
\item 17573888 bayt
\item 17162000 bayt
\item 140591104 bayt
\item 2196736 bayt
\end{enumerate}
\begin{tcolorbox}[enhanced,breakable,colback=teal!4,colframe=teal!50!black,boxrule=0.5pt,arc=1.1mm,left=3mm,right=3mm,top=1.7mm,bottom=1.7mm,before skip=2mm,after skip=2mm]
\textbf{\color{teal}To'g'ri javob:} \textbf{A) 17573888 bayt}
\par\smallskip
\textbf{Izoh:} 17162 \ensuremath{\times} 1024 = 17573888 bayt.
\end{tcolorbox}
\vspace{2mm}\hrule height 0.2pt \vspace{3mm}

% DB question_id: 1e5808c5-7b3e-55ad-8f90-a5aabbe59e91
\needspace{5\baselineskip}
\noindent\textbf{\color{ink}394.}
\par
57852 KiB axborot necha baytga teng? 1 KiB = 1024 bayt deb oling.
\begin{enumerate}[label=\Alph*),leftmargin=9mm]
\item 7405056 bayt
\item 59240448 bayt
\item 57852000 bayt
\item 473923584 bayt
\end{enumerate}
\begin{tcolorbox}[enhanced,breakable,colback=teal!4,colframe=teal!50!black,boxrule=0.5pt,arc=1.1mm,left=3mm,right=3mm,top=1.7mm,bottom=1.7mm,before skip=2mm,after skip=2mm]
\textbf{\color{teal}To'g'ri javob:} \textbf{B) 59240448 bayt}
\par\smallskip
\textbf{Izoh:} 57852 \ensuremath{\times} 1024 = 59240448 bayt.
\end{tcolorbox}
\vspace{2mm}\hrule height 0.2pt \vspace{3mm}

% DB question_id: 2615dce2-9e4b-5737-92f6-124e072f63fc
\needspace{5\baselineskip}
\noindent\textbf{\color{ink}395.}
\par
Harbiy bazadagi 89 ta tankning har biriga barcha tanklar uchun bir xil uzunlikdagi minimal ikkilik kod berilgan. Qurilma marradan o'tgan 72 ta tank kodini bitlar oqimi sifatida ketma-ket yozsa, qancha bayt xotira ishlatiladi?
\begin{enumerate}[label=\Alph*),leftmargin=9mm]
\item 63 bayt
\item 72 bayt
\item 504 bayt
\item 89 bayt
\end{enumerate}
\begin{tcolorbox}[enhanced,breakable,colback=teal!4,colframe=teal!50!black,boxrule=0.5pt,arc=1.1mm,left=3mm,right=3mm,top=1.7mm,bottom=1.7mm,before skip=2mm,after skip=2mm]
\textbf{\color{teal}To'g'ri javob:} \textbf{A) 63 bayt}
\par\smallskip
\textbf{Izoh:} 2\textsuperscript{6}\textless{}89\ensuremath{\leq}2\textsuperscript{7}, shuning uchun har kod 7 bit. 72\ensuremath{\cdot}7=504 bit=63 bayt.
\end{tcolorbox}
\vspace{2mm}\hrule height 0.2pt \vspace{3mm}

% DB question_id: 2d434ade-3769-5095-8b59-8aeb9b8270b7
\needspace{5\baselineskip}
\noindent\textbf{\color{ink}396.}
\par
53826 KiB axborot necha baytga teng? 1 KiB = 1024 bayt deb oling.
\begin{enumerate}[label=\Alph*),leftmargin=9mm]
\item 53826000 bayt
\item 55117824 bayt
\item 6889728 bayt
\item 440942592 bayt
\end{enumerate}
\begin{tcolorbox}[enhanced,breakable,colback=teal!4,colframe=teal!50!black,boxrule=0.5pt,arc=1.1mm,left=3mm,right=3mm,top=1.7mm,bottom=1.7mm,before skip=2mm,after skip=2mm]
\textbf{\color{teal}To'g'ri javob:} \textbf{B) 55117824 bayt}
\par\smallskip
\textbf{Izoh:} 53826 \ensuremath{\times} 1024 = 55117824 bayt.
\end{tcolorbox}
\vspace{2mm}\hrule height 0.2pt \vspace{3mm}

% DB question_id: 2f5a5727-0731-5307-a915-88e7e0c42d2e
\needspace{5\baselineskip}
\noindent\textbf{\color{ink}397.}
\par
6133 KiB axborot necha baytga teng? 1 KiB = 1024 bayt deb oling.
\begin{enumerate}[label=\Alph*),leftmargin=9mm]
\item 785024 bayt
\item 6280192 bayt
\item 6133000 bayt
\item 50241536 bayt
\end{enumerate}
\begin{tcolorbox}[enhanced,breakable,colback=teal!4,colframe=teal!50!black,boxrule=0.5pt,arc=1.1mm,left=3mm,right=3mm,top=1.7mm,bottom=1.7mm,before skip=2mm,after skip=2mm]
\textbf{\color{teal}To'g'ri javob:} \textbf{B) 6280192 bayt}
\par\smallskip
\textbf{Izoh:} 6133 \ensuremath{\times} 1024 = 6280192 bayt.
\end{tcolorbox}
\vspace{2mm}\hrule height 0.2pt \vspace{3mm}

% DB question_id: 3fe94e8d-482e-5f29-af7f-c341ebe0bb1e
\needspace{5\baselineskip}
\noindent\textbf{\color{ink}398.}
\par
8163 KiB axborot necha baytga teng? 1 KiB = 1024 bayt deb oling.
\begin{enumerate}[label=\Alph*),leftmargin=9mm]
\item 66871296 bayt
\item 8358912 bayt
\item 8163000 bayt
\item 1044864 bayt
\end{enumerate}
\begin{tcolorbox}[enhanced,breakable,colback=teal!4,colframe=teal!50!black,boxrule=0.5pt,arc=1.1mm,left=3mm,right=3mm,top=1.7mm,bottom=1.7mm,before skip=2mm,after skip=2mm]
\textbf{\color{teal}To'g'ri javob:} \textbf{B) 8358912 bayt}
\par\smallskip
\textbf{Izoh:} 8163 \ensuremath{\times} 1024 = 8358912 bayt.
\end{tcolorbox}
\vspace{2mm}\hrule height 0.2pt \vspace{3mm}

% DB question_id: 3feb1a88-7422-5956-bb32-987a52acdee4
\needspace{5\baselineskip}
\noindent\textbf{\color{ink}399.}
\par
26 harfli ingliz alifbosida Sezar shifri qo'llanadi: har bir harf alifboda 4 o'rin oldinga siljitiladi va \texttt{Z} dan keyin \texttt{A} ga o'tiladi. \texttt{SINFDOSHLAR} so'zining shifrmatnini toping.
\begin{enumerate}[label=\Alph*),leftmargin=9mm]
\item WMRJHSWLPEV
\item VLQIGRVKODU
\item XNSKITXMQFW
\item UKPHFQUJMBT
\end{enumerate}
\begin{tcolorbox}[enhanced,breakable,colback=teal!4,colframe=teal!50!black,boxrule=0.5pt,arc=1.1mm,left=3mm,right=3mm,top=1.7mm,bottom=1.7mm,before skip=2mm,after skip=2mm]
\textbf{\color{teal}To'g'ri javob:} \textbf{A) WMRJHSWLPEV}
\par\smallskip
\textbf{Izoh:} S\ensuremath{\to}W, I\ensuremath{\to}M, N\ensuremath{\to}R, F\ensuremath{\to}J, D\ensuremath{\to}H, O\ensuremath{\to}S, S\ensuremath{\to}W, H\ensuremath{\to}L, L\ensuremath{\to}P, A\ensuremath{\to}E, R\ensuremath{\to}V. Natija: WMRJHSWLPEV.
\end{tcolorbox}
\vspace{2mm}\hrule height 0.2pt \vspace{3mm}

% DB question_id: 404016a4-3931-5ab9-acba-5bb62c01f971
\needspace{5\baselineskip}
\noindent\textbf{\color{ink}400.}
\par
\texttt{CHEVROLET} so'zidagi turli harflar alifbo tartibida joylashtirilib, 0 dan boshlab bir xil minimal uzunlikdagi ikkilik kodlar berildi. \texttt{R} harfining kodi qaysi?
\begin{enumerate}[label=\Alph*),leftmargin=9mm]
\item 101
\item 100
\item 011
\item 110
\end{enumerate}
\begin{tcolorbox}[enhanced,breakable,colback=teal!4,colframe=teal!50!black,boxrule=0.5pt,arc=1.1mm,left=3mm,right=3mm,top=1.7mm,bottom=1.7mm,before skip=2mm,after skip=2mm]
\textbf{\color{teal}To'g'ri javob:} \textbf{A) 101}
\par\smallskip
\textbf{Izoh:} Turli harflar C,E,H,L,O,R,T,V - jami 8 ta, demak 3 bit yetadi. Kodlar 000 dan boshlab berilganda R oltinchi harf, ya'ni indeksi 5 va kodi 101\textsubscript{2}.
\end{tcolorbox}
\vspace{2mm}\hrule height 0.2pt \vspace{3mm}

% DB question_id: 462c06e4-156f-440f-b1a8-9d0e6e63be38
\needspace{5\baselineskip}
\noindent\textbf{\color{ink}401.}
\par
Quyidagi ta'rif qaysi atamaga tegishli? Kodlangan axborotni kelishilgan qoida asosida qayta tiklash.
\begin{enumerate}[label=\Alph*),leftmargin=9mm]
\item Doimiy uzunlikdagi kod
\item O'zgaruvchan uzunlikdagi kod
\item Dekodlash
\item Siqish
\end{enumerate}
\begin{tcolorbox}[enhanced,breakable,colback=teal!4,colframe=teal!50!black,boxrule=0.5pt,arc=1.1mm,left=3mm,right=3mm,top=1.7mm,bottom=1.7mm,before skip=2mm,after skip=2mm]
\textbf{\color{teal}To'g'ri javob:} \textbf{C) Dekodlash}
\par\smallskip
\textbf{Izoh:} Ta'rif Dekodlash tushunchasiga tegishli.
\end{tcolorbox}
\vspace{2mm}\hrule height 0.2pt \vspace{3mm}

% DB question_id: 46c5dc37-8aab-58a7-a58c-bae256ed98b8
\needspace{5\baselineskip}
\noindent\textbf{\color{ink}402.}
\par
Bir test natijasining JSON fayli o'rtacha 16 KiB. Bir kunda 40 960 ta natija yuborilsa, jami hajm necha MiB bo'ladi? 1 MiB=1024 KiB.
\begin{enumerate}[label=\Alph*),leftmargin=9mm]
\item 640 MiB
\item 655,36 MiB
\item 40 MiB
\item 6 400 MiB
\end{enumerate}
\begin{tcolorbox}[enhanced,breakable,colback=teal!4,colframe=teal!50!black,boxrule=0.5pt,arc=1.1mm,left=3mm,right=3mm,top=1.7mm,bottom=1.7mm,before skip=2mm,after skip=2mm]
\textbf{\color{teal}To'g'ri javob:} \textbf{A) 640 MiB}
\par\smallskip
\textbf{Izoh:} 16\ensuremath{\cdot}40 960=655 360 KiB; 655 360/1024=640 MiB.
\end{tcolorbox}
\vspace{2mm}\hrule height 0.2pt \vspace{3mm}

% DB question_id: 47ca9051-d941-5ada-b509-bebecdcfd085
\needspace{5\baselineskip}
\noindent\textbf{\color{ink}403.}
\par
2,5 GiB fayl arxivlanganda hajmi 25\% ga kamaydi. Hosil bo'lgan fayl 16 Mibit/s tezlikda uzatilsa, uzatish qancha sekund davom etadi? 1 GiB=1024 MiB.
\begin{enumerate}[label=\Alph*),leftmargin=9mm]
\item 960 s
\item 1 280 s
\item 120 s
\item 768 s
\end{enumerate}
\begin{tcolorbox}[enhanced,breakable,colback=teal!4,colframe=teal!50!black,boxrule=0.5pt,arc=1.1mm,left=3mm,right=3mm,top=1.7mm,bottom=1.7mm,before skip=2mm,after skip=2mm]
\textbf{\color{teal}To'g'ri javob:} \textbf{A) 960 s}
\par\smallskip
\textbf{Izoh:} Arxiv hajmi 2,5\ensuremath{\cdot}0,75=1,875 GiB=1920 MiB=15 360 Mibit. 15 360/16=960 s.
\end{tcolorbox}
\vspace{2mm}\hrule height 0.2pt \vspace{3mm}

% DB question_id: 49110132-648b-4376-b477-7ee1e4d5b9b7
\needspace{5\baselineskip}
\noindent\textbf{\color{ink}404.}
\par
8 ta ketma-ket bit bitta guruh sifatida saqlandi. Bu vaziyatdagi asosiy tushunchani aniqlang.
\begin{enumerate}[label=\Alph*),leftmargin=9mm]
\item Xartli formulasi
\item Alifbo quvvati
\item Minimal butun kod uzunligi
\item Bayt
\end{enumerate}
\begin{tcolorbox}[enhanced,breakable,colback=teal!4,colframe=teal!50!black,boxrule=0.5pt,arc=1.1mm,left=3mm,right=3mm,top=1.7mm,bottom=1.7mm,before skip=2mm,after skip=2mm]
\textbf{\color{teal}To'g'ri javob:} \textbf{D) Bayt}
\par\smallskip
\textbf{Izoh:} Vaziyat Baytga misol: sakkiz bitdan tashkil topgan axborot birligi.
\end{tcolorbox}
\vspace{2mm}\hrule height 0.2pt \vspace{3mm}

% DB question_id: 4def309b-9fba-52f1-9ef5-3390820eddb6
\needspace{5\baselineskip}
\noindent\textbf{\color{ink}405.}
\par
21978 KiB axborot necha baytga teng? 1 KiB = 1024 bayt deb oling.
\begin{enumerate}[label=\Alph*),leftmargin=9mm]
\item 21978000 bayt
\item 22505472 bayt
\item 180043776 bayt
\item 2813184 bayt
\end{enumerate}
\begin{tcolorbox}[enhanced,breakable,colback=teal!4,colframe=teal!50!black,boxrule=0.5pt,arc=1.1mm,left=3mm,right=3mm,top=1.7mm,bottom=1.7mm,before skip=2mm,after skip=2mm]
\textbf{\color{teal}To'g'ri javob:} \textbf{B) 22505472 bayt}
\par\smallskip
\textbf{Izoh:} 21978 \ensuremath{\times} 1024 = 22505472 bayt.
\end{tcolorbox}
\vspace{2mm}\hrule height 0.2pt \vspace{3mm}

% DB question_id: 50c05e2f-8517-5a2a-8be2-f1f96361df8f
\needspace{5\baselineskip}
\noindent\textbf{\color{ink}406.}
\par
26 harfli ingliz alifbosida Sezar shifri qo'llanadi: harf 4 o'rin oldinga suriladi va Z dan keyin A ga o'tiladi. \texttt{COMPUTER} so'zining shifrmatni qaysi?
\begin{enumerate}[label=\Alph*),leftmargin=9mm]
\item GSQTYXIV
\item GSQTYVIX
\item DPNQVUFS
\item YKILQPAN
\end{enumerate}
\begin{tcolorbox}[enhanced,breakable,colback=teal!4,colframe=teal!50!black,boxrule=0.5pt,arc=1.1mm,left=3mm,right=3mm,top=1.7mm,bottom=1.7mm,before skip=2mm,after skip=2mm]
\textbf{\color{teal}To'g'ri javob:} \textbf{A) GSQTYXIV}
\par\smallskip
\textbf{Izoh:} C\ensuremath{\to}G, O\ensuremath{\to}S, M\ensuremath{\to}Q, P\ensuremath{\to}T, U\ensuremath{\to}Y, T\ensuremath{\to}X, E\ensuremath{\to}I, R\ensuremath{\to}V; natija GSQTYXIV.
\end{tcolorbox}
\vspace{2mm}\hrule height 0.2pt \vspace{3mm}

% DB question_id: 5b2e1f0c-3bfd-562e-a95b-263e89eebf48
\needspace{5\baselineskip}
\noindent\textbf{\color{ink}407.}
\par
44540 KiB axborot necha baytga teng? 1 KiB = 1024 bayt deb oling.
\begin{enumerate}[label=\Alph*),leftmargin=9mm]
\item 364871680 bayt
\item 5701120 bayt
\item 44540000 bayt
\item 45608960 bayt
\end{enumerate}
\begin{tcolorbox}[enhanced,breakable,colback=teal!4,colframe=teal!50!black,boxrule=0.5pt,arc=1.1mm,left=3mm,right=3mm,top=1.7mm,bottom=1.7mm,before skip=2mm,after skip=2mm]
\textbf{\color{teal}To'g'ri javob:} \textbf{D) 45608960 bayt}
\par\smallskip
\textbf{Izoh:} 44540 \ensuremath{\times} 1024 = 45608960 bayt.
\end{tcolorbox}
\vspace{2mm}\hrule height 0.2pt \vspace{3mm}

% DB question_id: 5b937adf-8cec-5473-9184-6f5b6a851ba8
\needspace{5\baselineskip}
\noindent\textbf{\color{ink}408.}
\par
5881 KiB axborot necha baytga teng? 1 KiB = 1024 bayt deb oling.
\begin{enumerate}[label=\Alph*),leftmargin=9mm]
\item 6022144 bayt
\item 5881000 bayt
\item 48177152 bayt
\item 752768 bayt
\end{enumerate}
\begin{tcolorbox}[enhanced,breakable,colback=teal!4,colframe=teal!50!black,boxrule=0.5pt,arc=1.1mm,left=3mm,right=3mm,top=1.7mm,bottom=1.7mm,before skip=2mm,after skip=2mm]
\textbf{\color{teal}To'g'ri javob:} \textbf{A) 6022144 bayt}
\par\smallskip
\textbf{Izoh:} 5881 \ensuremath{\times} 1024 = 6022144 bayt.
\end{tcolorbox}
\vspace{2mm}\hrule height 0.2pt \vspace{3mm}

% DB question_id: 5fb4c80d-9a11-54c4-967b-9d66aece2eef
\needspace{5\baselineskip}
\noindent\textbf{\color{ink}409.}
\par
26 harfli ingliz alifbosida Sezar shifri bilan harflar 5 o'rin oldinga siljitiladi. \texttt{KITOBLAR} so'zining shifrmatnini toping.
\begin{enumerate}[label=\Alph*),leftmargin=9mm]
\item PNYTGQFW
\item OMXSFPEV
\item QOZUHRGX
\item FDOJWGVN
\end{enumerate}
\begin{tcolorbox}[enhanced,breakable,colback=teal!4,colframe=teal!50!black,boxrule=0.5pt,arc=1.1mm,left=3mm,right=3mm,top=1.7mm,bottom=1.7mm,before skip=2mm,after skip=2mm]
\textbf{\color{teal}To'g'ri javob:} \textbf{A) PNYTGQFW}
\par\smallskip
\textbf{Izoh:} K\ensuremath{\to}P, I\ensuremath{\to}N, T\ensuremath{\to}Y, O\ensuremath{\to}T, B\ensuremath{\to}G, L\ensuremath{\to}Q, A\ensuremath{\to}F, R\ensuremath{\to}W. Natija PNYTGQFW.
\end{tcolorbox}
\vspace{2mm}\hrule height 0.2pt \vspace{3mm}

% DB question_id: 601c6357-aebc-4fbd-8bd8-b0d8137a2aea
\needspace{5\baselineskip}
\noindent\textbf{\color{ink}410.}
\par
Quyidagi ta'rif qaysi atamaga tegishli? Iec standartida 1024 baytga teng ikkilik birlik.
\begin{enumerate}[label=\Alph*),leftmargin=9mm]
\item KiB
\item Xotira sig'imi
\item Bit
\item Bayt
\end{enumerate}
\begin{tcolorbox}[enhanced,breakable,colback=teal!4,colframe=teal!50!black,boxrule=0.5pt,arc=1.1mm,left=3mm,right=3mm,top=1.7mm,bottom=1.7mm,before skip=2mm,after skip=2mm]
\textbf{\color{teal}To'g'ri javob:} \textbf{A) KiB}
\par\smallskip
\textbf{Izoh:} Ta'rif KiB tushunchasiga tegishli.
\end{tcolorbox}
\vspace{2mm}\hrule height 0.2pt \vspace{3mm}

% DB question_id: 60fc09dd-f327-514b-864b-94ed6ff35c52
\needspace{5\baselineskip}
\noindent\textbf{\color{ink}411.}
\par
46571 KiB axborot necha baytga teng? 1 KiB = 1024 bayt deb oling.
\begin{enumerate}[label=\Alph*),leftmargin=9mm]
\item 381509632 bayt
\item 47688704 bayt
\item 5961088 bayt
\item 46571000 bayt
\end{enumerate}
\begin{tcolorbox}[enhanced,breakable,colback=teal!4,colframe=teal!50!black,boxrule=0.5pt,arc=1.1mm,left=3mm,right=3mm,top=1.7mm,bottom=1.7mm,before skip=2mm,after skip=2mm]
\textbf{\color{teal}To'g'ri javob:} \textbf{B) 47688704 bayt}
\par\smallskip
\textbf{Izoh:} 46571 \ensuremath{\times} 1024 = 47688704 bayt.
\end{tcolorbox}
\vspace{2mm}\hrule height 0.2pt \vspace{3mm}

% DB question_id: 61e59577-48ba-4b27-b5d3-1174005bc5a1
\needspace{5\baselineskip}
\noindent\textbf{\color{ink}412.}
\par
26 harfli lotin alifbosida Sezar shifri bilan har bir harf 4 pozitsiya oldinga suriladi. \texttt{BILIM} so'zining shifrlangan ko'rinishini toping.
\begin{enumerate}[label=\Alph*),leftmargin=9mm]
\item FMPQM
\item FMPMQ
\item ELOLP
\item ELOLPQ
\end{enumerate}
\begin{tcolorbox}[enhanced,breakable,colback=teal!4,colframe=teal!50!black,boxrule=0.5pt,arc=1.1mm,left=3mm,right=3mm,top=1.7mm,bottom=1.7mm,before skip=2mm,after skip=2mm]
\textbf{\color{teal}To'g'ri javob:} \textbf{B) FMPMQ}
\par\smallskip
\textbf{Izoh:} B\ensuremath{\to}F, I\ensuremath{\to}M, L\ensuremath{\to}P, I\ensuremath{\to}M, M\ensuremath{\to}Q. Natija FMPMQ.
\end{tcolorbox}
\vspace{2mm}\hrule height 0.2pt \vspace{3mm}

% DB question_id: 626833eb-4969-47ab-b964-efdb70c2f34f
\needspace{5\baselineskip}
\noindent\textbf{\color{ink}413.}
\par
Quyidagi ta'rif qaysi atamaga tegishli? Turli belgilar turli uzunlikdagi kodlar bilan ifodalanadigan kod.
\begin{enumerate}[label=\Alph*),leftmargin=9mm]
\item O'zgaruvchan uzunlikdagi kod
\item Siqish
\item Shifrlash
\item Asimmetrik shifrlash
\end{enumerate}
\begin{tcolorbox}[enhanced,breakable,colback=teal!4,colframe=teal!50!black,boxrule=0.5pt,arc=1.1mm,left=3mm,right=3mm,top=1.7mm,bottom=1.7mm,before skip=2mm,after skip=2mm]
\textbf{\color{teal}To'g'ri javob:} \textbf{A) O'zgaruvchan uzunlikdagi kod}
\par\smallskip
\textbf{Izoh:} Ta'rif O'zgaruvchan uzunlikdagi kod tushunchasiga tegishli.
\end{tcolorbox}
\vspace{2mm}\hrule height 0.2pt \vspace{3mm}

% DB question_id: 633b3872-a13b-5b97-b79c-8e92cc0cdd88
\needspace{5\baselineskip}
\noindent\textbf{\color{ink}414.}
\par
\texttt{OLIMPIADA} so'zi 3\ensuremath{\times}3 jadvalga satrma-satr yozildi, so'ng ustunma-ustun yuqoridan pastga o'qildi. Hosil bo'lgan shifrmatnni toping.
\begin{enumerate}[label=\Alph*),leftmargin=9mm]
\item OMALPDIIA
\item OILMIPDAA
\item ADAPMIILO
\item OLIMPIADA
\end{enumerate}
\begin{tcolorbox}[enhanced,breakable,colback=teal!4,colframe=teal!50!black,boxrule=0.5pt,arc=1.1mm,left=3mm,right=3mm,top=1.7mm,bottom=1.7mm,before skip=2mm,after skip=2mm]
\textbf{\color{teal}To'g'ri javob:} \textbf{A) OMALPDIIA}
\par\smallskip
\textbf{Izoh:} Jadval satrlari OLI / MPI / ADA. Ustunlar bo'yicha OMA, LPD, IIA o'qiladi: OMALPDIIA.
\end{tcolorbox}
\vspace{2mm}\hrule height 0.2pt \vspace{3mm}

% DB question_id: 6510b752-e677-5109-9283-68ffcfc8d539
\needspace{5\baselineskip}
\noindent\textbf{\color{ink}415.}
\par
JShShIR tuzilishida birinchi raqam jins va tug'ilgan asr indeksini, keyingi 6 raqam esa \texttt{KKOOYY} formatidagi tug'ilgan sanani bildiradi. 24.08.1992 da tug'ilgan ayolning JShShIRidagi dastlabki 7 raqam qaysi?
\begin{enumerate}[label=\Alph*),leftmargin=9mm]
\item 4240892
\item 3240892
\item 4082492
\item 4920824
\end{enumerate}
\begin{tcolorbox}[enhanced,breakable,colback=teal!4,colframe=teal!50!black,boxrule=0.5pt,arc=1.1mm,left=3mm,right=3mm,top=1.7mm,bottom=1.7mm,before skip=2mm,after skip=2mm]
\textbf{\color{teal}To'g'ri javob:} \textbf{A) 4240892}
\par\smallskip
\textbf{Izoh:} 1900-1999 yillarda tug'ilgan ayol indeksi 4; sana KKOOYY=240892. Dastlabki 7 raqam 4|240892=4240892. To'liq 14 raqamni hudud, tartib va nazorat raqamisiz aniqlab bo'lmaydi.
\end{tcolorbox}
\vspace{2mm}\hrule height 0.2pt \vspace{3mm}

% DB question_id: 67b0a4c2-fb95-5bd6-b716-5c7429420a6d
\needspace{5\baselineskip}
\noindent\textbf{\color{ink}416.}
\par
60134 KiB axborot necha baytga teng? 1 KiB = 1024 bayt deb oling.
\begin{enumerate}[label=\Alph*),leftmargin=9mm]
\item 61577216 bayt
\item 60134000 bayt
\item 7697152 bayt
\item 492617728 bayt
\end{enumerate}
\begin{tcolorbox}[enhanced,breakable,colback=teal!4,colframe=teal!50!black,boxrule=0.5pt,arc=1.1mm,left=3mm,right=3mm,top=1.7mm,bottom=1.7mm,before skip=2mm,after skip=2mm]
\textbf{\color{teal}To'g'ri javob:} \textbf{A) 61577216 bayt}
\par\smallskip
\textbf{Izoh:} 60134 \ensuremath{\times} 1024 = 61577216 bayt.
\end{tcolorbox}
\vspace{2mm}\hrule height 0.2pt \vspace{3mm}

% DB question_id: 72e7a4c0-10cb-4a6b-a900-83a0d1eaed9f
\needspace{5\baselineskip}
\noindent\textbf{\color{ink}417.}
\par
Quyidagi ta'rif qaysi atamaga tegishli? Si standartida 1000 baytga teng o'nlik birlik.
\begin{enumerate}[label=\Alph*),leftmargin=9mm]
\item KiB
\item Xotira sig'imi
\item Bit
\item kB
\end{enumerate}
\begin{tcolorbox}[enhanced,breakable,colback=teal!4,colframe=teal!50!black,boxrule=0.5pt,arc=1.1mm,left=3mm,right=3mm,top=1.7mm,bottom=1.7mm,before skip=2mm,after skip=2mm]
\textbf{\color{teal}To'g'ri javob:} \textbf{D) kB}
\par\smallskip
\textbf{Izoh:} Ta'rif kB tushunchasiga tegishli.
\end{tcolorbox}
\vspace{2mm}\hrule height 0.2pt \vspace{3mm}

% DB question_id: 8048696d-8f0c-4c94-8cb5-859cb19f726d
\needspace{5\baselineskip}
\noindent\textbf{\color{ink}418.}
\par
7-bit ASCII kodlari M=77, A=65, K=75, T=84, B=66 (o'nlik) bo'lsa, \texttt{MAKTAB} so'zining belgilar kodlarini ketma-ket yozib hosil qilingan ikkilik satrni toping.
\begin{enumerate}[label=\Alph*),leftmargin=9mm]
\item 011111111110010101101110011011110111011010
\item 010101110110100101101110011001000110111101
\item 100001100110010101101110011011110111011001
\item 100110110000011001011101010010000011000010
\end{enumerate}
\begin{tcolorbox}[enhanced,breakable,colback=teal!4,colframe=teal!50!black,boxrule=0.5pt,arc=1.1mm,left=3mm,right=3mm,top=1.7mm,bottom=1.7mm,before skip=2mm,after skip=2mm]
\textbf{\color{teal}To'g'ri javob:} \textbf{D) 100110110000011001011101010010000011000010}
\par\smallskip
\textbf{Izoh:} M=1001101, A=1000001, K=1001011, T=1010100, A=1000001, B=1000010. Ketma-ket birlashtirish D variantini beradi.
\end{tcolorbox}
\vspace{2mm}\hrule height 0.2pt \vspace{3mm}

% DB question_id: 81a99e1d-544c-4be9-bf76-935f6f62aba2
\needspace{5\baselineskip}
\noindent\textbf{\color{ink}419.}
\par
Quyidagi ta'rif qaysi atamaga tegishli? Axborotni belgilangan qoida bo'yicha boshqa ko'rinishga o'tkazish.
\begin{enumerate}[label=\Alph*),leftmargin=9mm]
\item Dekodlash
\item Kodlash
\item Doimiy uzunlikdagi kod
\item O'zgaruvchan uzunlikdagi kod
\end{enumerate}
\begin{tcolorbox}[enhanced,breakable,colback=teal!4,colframe=teal!50!black,boxrule=0.5pt,arc=1.1mm,left=3mm,right=3mm,top=1.7mm,bottom=1.7mm,before skip=2mm,after skip=2mm]
\textbf{\color{teal}To'g'ri javob:} \textbf{B) Kodlash}
\par\smallskip
\textbf{Izoh:} Ta'rif Kodlash tushunchasiga tegishli.
\end{tcolorbox}
\vspace{2mm}\hrule height 0.2pt \vspace{3mm}

% DB question_id: 81e5e3f7-5483-5014-8d42-21b9c64698f3
\needspace{5\baselineskip}
\noindent\textbf{\color{ink}420.}
\par
IBM PC CP437 kengaytirilgan kod sahifasida \texttt{10001101\textsubscript{2}} kodi qaysi belgiga mos?
\begin{enumerate}[label=\Alph*),leftmargin=9mm]
\item i
\item \textcopyright{}
\item A
\item C
\end{enumerate}
\begin{tcolorbox}[enhanced,breakable,colback=teal!4,colframe=teal!50!black,boxrule=0.5pt,arc=1.1mm,left=3mm,right=3mm,top=1.7mm,bottom=1.7mm,before skip=2mm,after skip=2mm]
\textbf{\color{teal}To'g'ri javob:} \textbf{A) i}
\par\smallskip
\textbf{Izoh:} 10001101\textsubscript{2}=141\textsubscript{1}\textsubscript{0}. Manbadagi CP437 jadvalida 141-kod \texttt{i} belgisiga mos.
\end{tcolorbox}
\vspace{2mm}\hrule height 0.2pt \vspace{3mm}

% DB question_id: 83984570-2478-51aa-b394-d8d93b478c83
\needspace{5\baselineskip}
\noindent\textbf{\color{ink}421.}
\par
56109 KiB axborot necha baytga teng? 1 KiB = 1024 bayt deb oling.
\begin{enumerate}[label=\Alph*),leftmargin=9mm]
\item 57455616 bayt
\item 7181952 bayt
\item 56109000 bayt
\item 459644928 bayt
\end{enumerate}
\begin{tcolorbox}[enhanced,breakable,colback=teal!4,colframe=teal!50!black,boxrule=0.5pt,arc=1.1mm,left=3mm,right=3mm,top=1.7mm,bottom=1.7mm,before skip=2mm,after skip=2mm]
\textbf{\color{teal}To'g'ri javob:} \textbf{A) 57455616 bayt}
\par\smallskip
\textbf{Izoh:} 56109 \ensuremath{\times} 1024 = 57455616 bayt.
\end{tcolorbox}
\vspace{2mm}\hrule height 0.2pt \vspace{3mm}

% DB question_id: 88c2e557-cd9d-5f98-9db4-8cbfac5bfc1d
\needspace{5\baselineskip}
\noindent\textbf{\color{ink}422.}
\par
Ko'k piksel \texttt{00}, oq piksel \texttt{01} bilan kodlanadi. Ikki qatorli 3\ensuremath{\times}2 rastrning qator kodlari mos ravishda \texttt{010100} va \texttt{010001}. Qaysi ranglar joylashuvi to'g'ri?
\begin{enumerate}[label=\Alph*),leftmargin=9mm]
\item oq-oq-ko'k / oq-ko'k-oq
\item ko'k-oq-ko'k / oq-ko'k-oq
\item oq-ko'k-oq / ko'k-oq-ko'k
\item ko'k-ko'k-oq / oq-ko'k-oq
\end{enumerate}
\begin{tcolorbox}[enhanced,breakable,colback=teal!4,colframe=teal!50!black,boxrule=0.5pt,arc=1.1mm,left=3mm,right=3mm,top=1.7mm,bottom=1.7mm,before skip=2mm,after skip=2mm]
\textbf{\color{teal}To'g'ri javob:} \textbf{A) oq-oq-ko'k / oq-ko'k-oq}
\par\smallskip
\textbf{Izoh:} Kodlar ikki bitdan ajratiladi: 01|01|00 = oq-oq-ko'k; 01|00|01 = oq-ko'k-oq.
\end{tcolorbox}
\vspace{2mm}\hrule height 0.2pt \vspace{3mm}

% DB question_id: 890cca66-09a8-5a17-bf00-2120ae8bc53f
\needspace{5\baselineskip}
\noindent\textbf{\color{ink}423.}
\par
17701 KiB axborot necha baytga teng? 1 KiB = 1024 bayt deb oling.
\begin{enumerate}[label=\Alph*),leftmargin=9mm]
\item 2265728 bayt
\item 145006592 bayt
\item 17701000 bayt
\item 18125824 bayt
\end{enumerate}
\begin{tcolorbox}[enhanced,breakable,colback=teal!4,colframe=teal!50!black,boxrule=0.5pt,arc=1.1mm,left=3mm,right=3mm,top=1.7mm,bottom=1.7mm,before skip=2mm,after skip=2mm]
\textbf{\color{teal}To'g'ri javob:} \textbf{D) 18125824 bayt}
\par\smallskip
\textbf{Izoh:} 17701 \ensuremath{\times} 1024 = 18125824 bayt.
\end{tcolorbox}
\vspace{2mm}\hrule height 0.2pt \vspace{3mm}

% DB question_id: 8d087248-e510-45a6-b25b-b80e5e96aec9
\needspace{5\baselineskip}
\noindent\textbf{\color{ink}424.}
\par
\texttt{10111111110011111110001101} xabarini quyidagi prefiks kod jadvali bo'yicha dekodlang: A=00, B=01, C=1011, D=1100, E=1101, F=1110, G=1111.
\begin{enumerate}[label=\Alph*),leftmargin=9mm]
\item GCDGFAE
\item CGDGFAE
\item EAFGDGC
\item CGDGFEA
\end{enumerate}
\begin{tcolorbox}[enhanced,breakable,colback=teal!4,colframe=teal!50!black,boxrule=0.5pt,arc=1.1mm,left=3mm,right=3mm,top=1.7mm,bottom=1.7mm,before skip=2mm,after skip=2mm]
\textbf{\color{teal}To'g'ri javob:} \textbf{B) CGDGFAE}
\par\smallskip
\textbf{Izoh:} Bitlar 1011|1111|1100|1111|1110|00|1101 tarzida ajraladi; bu C|G|D|G|F|A|E ni beradi.
\end{tcolorbox}
\vspace{2mm}\hrule height 0.2pt \vspace{3mm}

% DB question_id: 93650bfb-6977-5a48-940e-af0403c3e443
\needspace{5\baselineskip}
\noindent\textbf{\color{ink}425.}
\par
3599 KiB axborot necha baytga teng? 1 KiB = 1024 bayt deb oling.
\begin{enumerate}[label=\Alph*),leftmargin=9mm]
\item 3599000 bayt
\item 460672 bayt
\item 3685376 bayt
\item 29483008 bayt
\end{enumerate}
\begin{tcolorbox}[enhanced,breakable,colback=teal!4,colframe=teal!50!black,boxrule=0.5pt,arc=1.1mm,left=3mm,right=3mm,top=1.7mm,bottom=1.7mm,before skip=2mm,after skip=2mm]
\textbf{\color{teal}To'g'ri javob:} \textbf{C) 3685376 bayt}
\par\smallskip
\textbf{Izoh:} 3599 \ensuremath{\times} 1024 = 3685376 bayt.
\end{tcolorbox}
\vspace{2mm}\hrule height 0.2pt \vspace{3mm}

% DB question_id: 96fac222-5c45-5b5f-a879-29f975d9b939
\needspace{5\baselineskip}
\noindent\textbf{\color{ink}426.}
\par
33007 KiB axborot necha baytga teng? 1 KiB = 1024 bayt deb oling.
\begin{enumerate}[label=\Alph*),leftmargin=9mm]
\item 4224896 bayt
\item 33007000 bayt
\item 270393344 bayt
\item 33799168 bayt
\end{enumerate}
\begin{tcolorbox}[enhanced,breakable,colback=teal!4,colframe=teal!50!black,boxrule=0.5pt,arc=1.1mm,left=3mm,right=3mm,top=1.7mm,bottom=1.7mm,before skip=2mm,after skip=2mm]
\textbf{\color{teal}To'g'ri javob:} \textbf{D) 33799168 bayt}
\par\smallskip
\textbf{Izoh:} 33007 \ensuremath{\times} 1024 = 33799168 bayt.
\end{tcolorbox}
\vspace{2mm}\hrule height 0.2pt \vspace{3mm}

% DB question_id: 978dffd6-b346-5b84-9156-7b7e7a04a847
\needspace{5\baselineskip}
\noindent\textbf{\color{ink}427.}
\par
19696 KiB axborot necha baytga teng? 1 KiB = 1024 bayt deb oling.
\begin{enumerate}[label=\Alph*),leftmargin=9mm]
\item 161349632 bayt
\item 19696000 bayt
\item 20168704 bayt
\item 2521088 bayt
\end{enumerate}
\begin{tcolorbox}[enhanced,breakable,colback=teal!4,colframe=teal!50!black,boxrule=0.5pt,arc=1.1mm,left=3mm,right=3mm,top=1.7mm,bottom=1.7mm,before skip=2mm,after skip=2mm]
\textbf{\color{teal}To'g'ri javob:} \textbf{C) 20168704 bayt}
\par\smallskip
\textbf{Izoh:} 19696 \ensuremath{\times} 1024 = 20168704 bayt.
\end{tcolorbox}
\vspace{2mm}\hrule height 0.2pt \vspace{3mm}

% DB question_id: 988482c9-880a-564f-bed9-eae9d69f9265
\needspace{5\baselineskip}
\noindent\textbf{\color{ink}428.}
\par
58104 KiB axborot necha baytga teng? 1 KiB = 1024 bayt deb oling.
\begin{enumerate}[label=\Alph*),leftmargin=9mm]
\item 475987968 bayt
\item 7437312 bayt
\item 59498496 bayt
\item 58104000 bayt
\end{enumerate}
\begin{tcolorbox}[enhanced,breakable,colback=teal!4,colframe=teal!50!black,boxrule=0.5pt,arc=1.1mm,left=3mm,right=3mm,top=1.7mm,bottom=1.7mm,before skip=2mm,after skip=2mm]
\textbf{\color{teal}To'g'ri javob:} \textbf{C) 59498496 bayt}
\par\smallskip
\textbf{Izoh:} 58104 \ensuremath{\times} 1024 = 59498496 bayt.
\end{tcolorbox}
\vspace{2mm}\hrule height 0.2pt \vspace{3mm}

% DB question_id: 99d44254-fc26-4692-a706-6cff354c8cb2
\needspace{5\baselineskip}
\noindent\textbf{\color{ink}429.}
\par
Quyidagi ta'rif qaysi atamaga tegishli? Har bir belgi bir xil sondagi kod elementi bilan ifodalanadigan kod.
\begin{enumerate}[label=\Alph*),leftmargin=9mm]
\item O'zgaruvchan uzunlikdagi kod
\item Siqish
\item Shifrlash
\item Doimiy uzunlikdagi kod
\end{enumerate}
\begin{tcolorbox}[enhanced,breakable,colback=teal!4,colframe=teal!50!black,boxrule=0.5pt,arc=1.1mm,left=3mm,right=3mm,top=1.7mm,bottom=1.7mm,before skip=2mm,after skip=2mm]
\textbf{\color{teal}To'g'ri javob:} \textbf{D) Doimiy uzunlikdagi kod}
\par\smallskip
\textbf{Izoh:} Ta'rif Doimiy uzunlikdagi kod tushunchasiga tegishli.
\end{tcolorbox}
\vspace{2mm}\hrule height 0.2pt \vspace{3mm}

% DB question_id: a0a6b621-d799-55bf-9d39-22ba6dbb2c73
\needspace{5\baselineskip}
\noindent\textbf{\color{ink}430.}
\par
Standart ASCII da \texttt{MAKTAB} so'zining har bir harfini 8 bitli ikkilik kod bilan yozilgan to'g'ri ketma-ketlikni toping.
\begin{enumerate}[label=\Alph*),leftmargin=9mm]
\item 01001101 01000001 01001011 01010100 01000001 01000010
\item 01001100 01000001 01001011 01010100 01000001 01000010
\item 01001101 01000010 01001011 01010100 01000001 01000010
\item 01001101 01000001 01001010 01010100 01000001 01000010
\end{enumerate}
\begin{tcolorbox}[enhanced,breakable,colback=teal!4,colframe=teal!50!black,boxrule=0.5pt,arc=1.1mm,left=3mm,right=3mm,top=1.7mm,bottom=1.7mm,before skip=2mm,after skip=2mm]
\textbf{\color{teal}To'g'ri javob:} \textbf{A) 01001101 01000001 01001011 01010100 01000001 01000010}
\par\smallskip
\textbf{Izoh:} M=77=01001101, A=65=01000001, K=75=01001011, T=84=01010100, A=65, B=66=01000010.
\end{tcolorbox}
\vspace{2mm}\hrule height 0.2pt \vspace{3mm}

% DB question_id: a2a61e82-bf9b-4dad-8cc3-34ca14ec8c0b
\needspace{5\baselineskip}
\noindent\textbf{\color{ink}431.}
\par
Quyidagi ta'rif qaysi atamaga tegishli? Axborotni saqlash yoki uzatish uchun fayl hajmini kamaytirish jarayoni.
\begin{enumerate}[label=\Alph*),leftmargin=9mm]
\item Shifrlash
\item Siqish
\item Asimmetrik shifrlash
\item Belgi
\end{enumerate}
\begin{tcolorbox}[enhanced,breakable,colback=teal!4,colframe=teal!50!black,boxrule=0.5pt,arc=1.1mm,left=3mm,right=3mm,top=1.7mm,bottom=1.7mm,before skip=2mm,after skip=2mm]
\textbf{\color{teal}To'g'ri javob:} \textbf{B) Siqish}
\par\smallskip
\textbf{Izoh:} Ta'rif Siqish tushunchasiga tegishli.
\end{tcolorbox}
\vspace{2mm}\hrule height 0.2pt \vspace{3mm}

% DB question_id: a65919ac-62fc-532b-b2b6-2bbb3d535376
\needspace{5\baselineskip}
\noindent\textbf{\color{ink}432.}
\par
2000 xil mahsulotning har biriga yagona, bir xil uzunlikdagi ikkilik ID beriladi. Bitta ID uchun eng kamida necha bit kerak?
\begin{enumerate}[label=\Alph*),leftmargin=9mm]
\item 11 bit
\item 10 bit
\item 12 bit
\item 2000 bit
\end{enumerate}
\begin{tcolorbox}[enhanced,breakable,colback=teal!4,colframe=teal!50!black,boxrule=0.5pt,arc=1.1mm,left=3mm,right=3mm,top=1.7mm,bottom=1.7mm,before skip=2mm,after skip=2mm]
\textbf{\color{teal}To'g'ri javob:} \textbf{A) 11 bit}
\par\smallskip
\textbf{Izoh:} 2\textsuperscript{1}\textsuperscript{0}=1024\textless{}2000\ensuremath{\leq}2048=2\textsuperscript{1}\textsuperscript{1}, shuning uchun bitta ID uchun kamida 11 bit kerak.
\end{tcolorbox}
\vspace{2mm}\hrule height 0.2pt \vspace{3mm}

% DB question_id: a6f688cd-442c-4c67-ba8b-45e12f139013
\needspace{5\baselineskip}
\noindent\textbf{\color{ink}433.}
\par
Kalit faqat yoqilgan yoki o'chirilgan holatdan birini oldi. Bu vaziyatdagi asosiy tushunchani aniqlang.
\begin{enumerate}[label=\Alph*),leftmargin=9mm]
\item Bayt
\item Xartli formulasi
\item Bit
\item Alifbo quvvati
\end{enumerate}
\begin{tcolorbox}[enhanced,breakable,colback=teal!4,colframe=teal!50!black,boxrule=0.5pt,arc=1.1mm,left=3mm,right=3mm,top=1.7mm,bottom=1.7mm,before skip=2mm,after skip=2mm]
\textbf{\color{teal}To'g'ri javob:} \textbf{C) Bit}
\par\smallskip
\textbf{Izoh:} Vaziyat Bitga misol: ikki teng ehtimolli holatdan birini ajratish uchun yetarli eng kichik axborot birligi.
\end{tcolorbox}
\vspace{2mm}\hrule height 0.2pt \vspace{3mm}

% DB question_id: aaa1222b-35bf-55a7-86e4-683ffa9b2d8d
\needspace{5\baselineskip}
\noindent\textbf{\color{ink}434.}
\par
JShShIR tuzilishida birinchi raqam jins va tug'ilgan asr indeksini, keyingi 6 raqam esa \texttt{KKOOYY} formatidagi tug'ilgan sanani bildiradi. 19.02.1992 da tug'ilgan erkakning JShShIRidagi dastlabki 7 raqam qaysi?
\begin{enumerate}[label=\Alph*),leftmargin=9mm]
\item 3190292
\item 4190292
\item 3021992
\item 3920219
\end{enumerate}
\begin{tcolorbox}[enhanced,breakable,colback=teal!4,colframe=teal!50!black,boxrule=0.5pt,arc=1.1mm,left=3mm,right=3mm,top=1.7mm,bottom=1.7mm,before skip=2mm,after skip=2mm]
\textbf{\color{teal}To'g'ri javob:} \textbf{A) 3190292}
\par\smallskip
\textbf{Izoh:} 1900-1999 yillarda tug'ilgan erkak indeksi 3; sana KKOOYY=190292. Dastlabki 7 raqam 3|190292=3190292. To'liq JShShIR uchun yana hudud, tartib va nazorat raqamlari kerak.
\end{tcolorbox}
\vspace{2mm}\hrule height 0.2pt \vspace{3mm}

% DB question_id: ad628ede-b81b-505c-8389-af299ecc274b
\needspace{5\baselineskip}
\noindent\textbf{\color{ink}435.}
\par
Bir kanalli siqilmagan audio 48 000 Hz, 32 bit/namuna bilan 4 daqiqa yozildi. Fayl hajmini eng yaqin butun MiBga yaxlitlang. 1 MiB=2\textsuperscript{2}\textsuperscript{0} bayt.
\begin{enumerate}[label=\Alph*),leftmargin=9mm]
\item 44 MiB
\item 87 MiB
\item 125 MiB
\item 175 MiB
\end{enumerate}
\begin{tcolorbox}[enhanced,breakable,colback=teal!4,colframe=teal!50!black,boxrule=0.5pt,arc=1.1mm,left=3mm,right=3mm,top=1.7mm,bottom=1.7mm,before skip=2mm,after skip=2mm]
\textbf{\color{teal}To'g'ri javob:} \textbf{A) 44 MiB}
\par\smallskip
\textbf{Izoh:} Hajm 48 000\ensuremath{\cdot}32\ensuremath{\cdot}240/8=46 080 000 bayt. Bu 46 080 000/2\textsuperscript{2}\textsuperscript{0}\ensuremath{\approx}43,95 MiB, ya'ni 44 MiB.
\end{tcolorbox}
\vspace{2mm}\hrule height 0.2pt \vspace{3mm}

% DB question_id: af7f702a-53fa-5689-bd45-5d38099eb64b
\needspace{5\baselineskip}
\noindent\textbf{\color{ink}436.}
\par
28982 KiB axborot necha baytga teng? 1 KiB = 1024 bayt deb oling.
\begin{enumerate}[label=\Alph*),leftmargin=9mm]
\item 3709696 bayt
\item 237420544 bayt
\item 28982000 bayt
\item 29677568 bayt
\end{enumerate}
\begin{tcolorbox}[enhanced,breakable,colback=teal!4,colframe=teal!50!black,boxrule=0.5pt,arc=1.1mm,left=3mm,right=3mm,top=1.7mm,bottom=1.7mm,before skip=2mm,after skip=2mm]
\textbf{\color{teal}To'g'ri javob:} \textbf{D) 29677568 bayt}
\par\smallskip
\textbf{Izoh:} 28982 \ensuremath{\times} 1024 = 29677568 bayt.
\end{tcolorbox}
\vspace{2mm}\hrule height 0.2pt \vspace{3mm}

% DB question_id: afb4f27a-9b91-5aab-bb48-d79ea4884b52
\needspace{5\baselineskip}
\noindent\textbf{\color{ink}437.}
\par
Pasport seriyasi 26 harfli lotin alifbosidan olingan 2 ta belgi, raqami esa 10 raqamli alifbodan olingan 7 ta belgidan iborat. Har qismda belgilar bir xil eng kam bitlar bilan kodlansa, seriya, raqam va jami uchun mos ravishda necha bit kerak?
\begin{enumerate}[label=\Alph*),leftmargin=9mm]
\item 10 bit; 28 bit; jami 38 bit
\item 12 bit; 28 bit; jami 40 bit
\item 10 bit; 21 bit; jami 31 bit
\item 52 bit; 70 bit; jami 122 bit
\end{enumerate}
\begin{tcolorbox}[enhanced,breakable,colback=teal!4,colframe=teal!50!black,boxrule=0.5pt,arc=1.1mm,left=3mm,right=3mm,top=1.7mm,bottom=1.7mm,before skip=2mm,after skip=2mm]
\textbf{\color{teal}To'g'ri javob:} \textbf{A) 10 bit; 28 bit; jami 38 bit}
\par\smallskip
\textbf{Izoh:} 26 harf uchun 5 bit: 2\ensuremath{\cdot}5=10 bit. 10 raqam uchun 4 bit: 7\ensuremath{\cdot}4=28 bit. Jami 38 bit.
\end{tcolorbox}
\vspace{2mm}\hrule height 0.2pt \vspace{3mm}

% DB question_id: b0ede91a-cfda-4267-be76-a20507917b6d
\needspace{5\baselineskip}
\noindent\textbf{\color{ink}438.}
\par
Quyidagi ta'rif qaysi atamaga tegishli? Bit ikkilik birlik, bayt esa 8 bitdan iborat ekanini ajratish.
\begin{enumerate}[label=\Alph*),leftmargin=9mm]
\item Rastr va vektor farqi
\item Kodek va konteyner farqi
\item Bit va bayt farqi
\item Validatsiya va verifikatsiya farqi
\end{enumerate}
\begin{tcolorbox}[enhanced,breakable,colback=teal!4,colframe=teal!50!black,boxrule=0.5pt,arc=1.1mm,left=3mm,right=3mm,top=1.7mm,bottom=1.7mm,before skip=2mm,after skip=2mm]
\textbf{\color{teal}To'g'ri javob:} \textbf{C) Bit va bayt farqi}
\par\smallskip
\textbf{Izoh:} Ta'rif Bit va bayt farqi tushunchasiga tegishli.
\end{tcolorbox}
\vspace{2mm}\hrule height 0.2pt \vspace{3mm}

% DB question_id: b44808e9-b3d1-5f55-8fff-0349ed94ab2c
\needspace{5\baselineskip}
\noindent\textbf{\color{ink}439.}
\par
30977 KiB axborot necha baytga teng? 1 KiB = 1024 bayt deb oling.
\begin{enumerate}[label=\Alph*),leftmargin=9mm]
\item 253763584 bayt
\item 30977000 bayt
\item 3965056 bayt
\item 31720448 bayt
\end{enumerate}
\begin{tcolorbox}[enhanced,breakable,colback=teal!4,colframe=teal!50!black,boxrule=0.5pt,arc=1.1mm,left=3mm,right=3mm,top=1.7mm,bottom=1.7mm,before skip=2mm,after skip=2mm]
\textbf{\color{teal}To'g'ri javob:} \textbf{D) 31720448 bayt}
\par\smallskip
\textbf{Izoh:} 30977 \ensuremath{\times} 1024 = 31720448 bayt.
\end{tcolorbox}
\vspace{2mm}\hrule height 0.2pt \vspace{3mm}

% DB question_id: b705688a-80f1-5aa0-ba95-9a5a73c527b0
\needspace{5\baselineskip}
\noindent\textbf{\color{ink}440.}
\par
15419 KiB axborot necha baytga teng? 1 KiB = 1024 bayt deb oling.
\begin{enumerate}[label=\Alph*),leftmargin=9mm]
\item 15419000 bayt
\item 15789056 bayt
\item 126312448 bayt
\item 1973632 bayt
\end{enumerate}
\begin{tcolorbox}[enhanced,breakable,colback=teal!4,colframe=teal!50!black,boxrule=0.5pt,arc=1.1mm,left=3mm,right=3mm,top=1.7mm,bottom=1.7mm,before skip=2mm,after skip=2mm]
\textbf{\color{teal}To'g'ri javob:} \textbf{B) 15789056 bayt}
\par\smallskip
\textbf{Izoh:} 15419 \ensuremath{\times} 1024 = 15789056 bayt.
\end{tcolorbox}
\vspace{2mm}\hrule height 0.2pt \vspace{3mm}

% DB question_id: b7756b73-ab85-54f3-a50d-c21d6aa682df
\needspace{5\baselineskip}
\noindent\textbf{\color{ink}441.}
\par
3851 KiB axborot necha baytga teng? 1 KiB = 1024 bayt deb oling.
\begin{enumerate}[label=\Alph*),leftmargin=9mm]
\item 492928 bayt
\item 3851000 bayt
\item 3943424 bayt
\item 31547392 bayt
\end{enumerate}
\begin{tcolorbox}[enhanced,breakable,colback=teal!4,colframe=teal!50!black,boxrule=0.5pt,arc=1.1mm,left=3mm,right=3mm,top=1.7mm,bottom=1.7mm,before skip=2mm,after skip=2mm]
\textbf{\color{teal}To'g'ri javob:} \textbf{C) 3943424 bayt}
\par\smallskip
\textbf{Izoh:} 3851 \ensuremath{\times} 1024 = 3943424 bayt.
\end{tcolorbox}
\vspace{2mm}\hrule height 0.2pt \vspace{3mm}

% DB question_id: bb35956d-a9b0-5ec7-beb3-5fb3ce066fd2
\needspace{5\baselineskip}
\noindent\textbf{\color{ink}442.}
\par
35541 KiB axborot necha baytga teng? 1 KiB = 1024 bayt deb oling.
\begin{enumerate}[label=\Alph*),leftmargin=9mm]
\item 291151872 bayt
\item 35541000 bayt
\item 4549248 bayt
\item 36393984 bayt
\end{enumerate}
\begin{tcolorbox}[enhanced,breakable,colback=teal!4,colframe=teal!50!black,boxrule=0.5pt,arc=1.1mm,left=3mm,right=3mm,top=1.7mm,bottom=1.7mm,before skip=2mm,after skip=2mm]
\textbf{\color{teal}To'g'ri javob:} \textbf{D) 36393984 bayt}
\par\smallskip
\textbf{Izoh:} 35541 \ensuremath{\times} 1024 = 36393984 bayt.
\end{tcolorbox}
\vspace{2mm}\hrule height 0.2pt \vspace{3mm}

% DB question_id: bed3e5ce-e56c-4c86-99b9-144207f1f221
\needspace{5\baselineskip}
\noindent\textbf{\color{ink}443.}
\par
Quyidagi ta'rif qaysi atamaga tegishli? Kb 1000 bayt, kib 1024 bayt ekanini standart belgilari bilan ajratish.
\begin{enumerate}[label=\Alph*),leftmargin=9mm]
\item Baud va bit/s farqi
\item HTTPS va ishonchlilik farqi
\item Axborot va ma'lumot farqi
\item kB va KiB farqi
\end{enumerate}
\begin{tcolorbox}[enhanced,breakable,colback=teal!4,colframe=teal!50!black,boxrule=0.5pt,arc=1.1mm,left=3mm,right=3mm,top=1.7mm,bottom=1.7mm,before skip=2mm,after skip=2mm]
\textbf{\color{teal}To'g'ri javob:} \textbf{D) kB va KiB farqi}
\par\smallskip
\textbf{Izoh:} Ta'rif kB va KiB farqi tushunchasiga tegishli.
\end{tcolorbox}
\vspace{2mm}\hrule height 0.2pt \vspace{3mm}

% DB question_id: c051b25e-a9be-5221-bfc5-7feb205135a0
\needspace{5\baselineskip}
\noindent\textbf{\color{ink}444.}
\par
13137 KiB axborot necha baytga teng? 1 KiB = 1024 bayt deb oling.
\begin{enumerate}[label=\Alph*),leftmargin=9mm]
\item 107618304 bayt
\item 13137000 bayt
\item 13452288 bayt
\item 1681536 bayt
\end{enumerate}
\begin{tcolorbox}[enhanced,breakable,colback=teal!4,colframe=teal!50!black,boxrule=0.5pt,arc=1.1mm,left=3mm,right=3mm,top=1.7mm,bottom=1.7mm,before skip=2mm,after skip=2mm]
\textbf{\color{teal}To'g'ri javob:} \textbf{C) 13452288 bayt}
\par\smallskip
\textbf{Izoh:} 13137 \ensuremath{\times} 1024 = 13452288 bayt.
\end{tcolorbox}
\vspace{2mm}\hrule height 0.2pt \vspace{3mm}

% DB question_id: c925d0a0-6566-4b31-8002-6a9f4372dbf7
\needspace{5\baselineskip}
\noindent\textbf{\color{ink}445.}
\par
Quyidagi ta'rif qaysi atamaga tegishli? Obyekt, xossa yoki tushunchani ifodalash uchun qabul qilingan ishora.
\begin{enumerate}[label=\Alph*),leftmargin=9mm]
\item Alifbo
\item Kod
\item Belgi
\item Kodlash
\end{enumerate}
\begin{tcolorbox}[enhanced,breakable,colback=teal!4,colframe=teal!50!black,boxrule=0.5pt,arc=1.1mm,left=3mm,right=3mm,top=1.7mm,bottom=1.7mm,before skip=2mm,after skip=2mm]
\textbf{\color{teal}To'g'ri javob:} \textbf{C) Belgi}
\par\smallskip
\textbf{Izoh:} Ta'rif Belgi tushunchasiga tegishli.
\end{tcolorbox}
\vspace{2mm}\hrule height 0.2pt \vspace{3mm}

% DB question_id: c9489009-afc1-5c4a-a45a-c1ff318dbcf6
\needspace{5\baselineskip}
\noindent\textbf{\color{ink}446.}
\par
19444 KiB axborot necha baytga teng? 1 KiB = 1024 bayt deb oling.
\begin{enumerate}[label=\Alph*),leftmargin=9mm]
\item 159285248 bayt
\item 2488832 bayt
\item 19444000 bayt
\item 19910656 bayt
\end{enumerate}
\begin{tcolorbox}[enhanced,breakable,colback=teal!4,colframe=teal!50!black,boxrule=0.5pt,arc=1.1mm,left=3mm,right=3mm,top=1.7mm,bottom=1.7mm,before skip=2mm,after skip=2mm]
\textbf{\color{teal}To'g'ri javob:} \textbf{D) 19910656 bayt}
\par\smallskip
\textbf{Izoh:} 19444 \ensuremath{\times} 1024 = 19910656 bayt.
\end{tcolorbox}
\vspace{2mm}\hrule height 0.2pt \vspace{3mm}

% DB question_id: c99c4c9b-bce7-4a81-afbf-5abe67a64a71
\needspace{5\baselineskip}
\noindent\textbf{\color{ink}447.}
\par
Quyidagi ta'rif qaysi atamaga tegishli? Kodek oqimni kodlashi, konteyner esa oqimlarni bir faylda saqlashini ajratish.
\begin{enumerate}[label=\Alph*),leftmargin=9mm]
\item Kodek va konteyner farqi
\item Validatsiya va verifikatsiya farqi
\item Internet va WWW farqi
\item kB va KiB farqi
\end{enumerate}
\begin{tcolorbox}[enhanced,breakable,colback=teal!4,colframe=teal!50!black,boxrule=0.5pt,arc=1.1mm,left=3mm,right=3mm,top=1.7mm,bottom=1.7mm,before skip=2mm,after skip=2mm]
\textbf{\color{teal}To'g'ri javob:} \textbf{A) Kodek va konteyner farqi}
\par\smallskip
\textbf{Izoh:} Ta'rif Kodek va konteyner farqi tushunchasiga tegishli.
\end{tcolorbox}
\vspace{2mm}\hrule height 0.2pt \vspace{3mm}

% DB question_id: cfdc4248-2226-5546-8946-835fc65378f9
\needspace{5\baselineskip}
\noindent\textbf{\color{ink}448.}
\par
Rang kodlari: sariq=\texttt{110}, qora=\texttt{000}, ko'k=\texttt{001}, havorang=\texttt{011}, oq=\texttt{111}, qizil=\texttt{100}. Yuqori qatori sariq-qora-ko'k, pastki qatori havorang-oq-qizil bo'lgan 3\ensuremath{\times}2 rastrning kodi qaysi?
\begin{enumerate}[label=\Alph*),leftmargin=9mm]
\item 110 000 001 / 011 111 100
\item 001 111 100 / 110 111 001
\item 101 100 001 / 000 110 010
\item 100 101 111 / 011 111 110
\end{enumerate}
\begin{tcolorbox}[enhanced,breakable,colback=teal!4,colframe=teal!50!black,boxrule=0.5pt,arc=1.1mm,left=3mm,right=3mm,top=1.7mm,bottom=1.7mm,before skip=2mm,after skip=2mm]
\textbf{\color{teal}To'g'ri javob:} \textbf{A) 110 000 001 / 011 111 100}
\par\smallskip
\textbf{Izoh:} Ranglarni berilgan kodlarga to'g'ridan-to'g'ri almashtirsak yuqori qator 110 000 001, pastki qator 011 111 100 bo'ladi.
\end{tcolorbox}
\vspace{2mm}\hrule height 0.2pt \vspace{3mm}

% DB question_id: d47cdf8f-f0df-43ec-8294-1358b184fc6d
\needspace{5\baselineskip}
\noindent\textbf{\color{ink}449.}
\par
Quyidagi ta'rif qaysi atamaga tegishli? Kodlashda foydalaniladigan ruxsat etilgan belgilar to'plami.
\begin{enumerate}[label=\Alph*),leftmargin=9mm]
\item Kod
\item Kodlash
\item Dekodlash
\item Alifbo
\end{enumerate}
\begin{tcolorbox}[enhanced,breakable,colback=teal!4,colframe=teal!50!black,boxrule=0.5pt,arc=1.1mm,left=3mm,right=3mm,top=1.7mm,bottom=1.7mm,before skip=2mm,after skip=2mm]
\textbf{\color{teal}To'g'ri javob:} \textbf{D) Alifbo}
\par\smallskip
\textbf{Izoh:} Ta'rif Alifbo tushunchasiga tegishli.
\end{tcolorbox}
\vspace{2mm}\hrule height 0.2pt \vspace{3mm}

% DB question_id: d689d338-e80c-44d9-b615-8c40d660c73f
\needspace{5\baselineskip}
\noindent\textbf{\color{ink}450.}
\par
Quyidagi ta'rif qaysi atamaga tegishli? Kodlash taqdim etish formatini, shifrlash esa maxfiylikni ko'zlashini ajratish.
\begin{enumerate}[label=\Alph*),leftmargin=9mm]
\item Bit va bayt farqi
\item Kodlash va shifrlash farqi
\item Rastr va vektor farqi
\item Kodek va konteyner farqi
\end{enumerate}
\begin{tcolorbox}[enhanced,breakable,colback=teal!4,colframe=teal!50!black,boxrule=0.5pt,arc=1.1mm,left=3mm,right=3mm,top=1.7mm,bottom=1.7mm,before skip=2mm,after skip=2mm]
\textbf{\color{teal}To'g'ri javob:} \textbf{B) Kodlash va shifrlash farqi}
\par\smallskip
\textbf{Izoh:} Ta'rif Kodlash va shifrlash farqi tushunchasiga tegishli.
\end{tcolorbox}
\vspace{2mm}\hrule height 0.2pt \vspace{3mm}

% DB question_id: d7b624f4-46b2-504c-b341-aa8b9124f2a0
\needspace{5\baselineskip}
\noindent\textbf{\color{ink}451.}
\par
21726 KiB axborot necha baytga teng? 1 KiB = 1024 bayt deb oling.
\begin{enumerate}[label=\Alph*),leftmargin=9mm]
\item 2780928 bayt
\item 21726000 bayt
\item 177979392 bayt
\item 22247424 bayt
\end{enumerate}
\begin{tcolorbox}[enhanced,breakable,colback=teal!4,colframe=teal!50!black,boxrule=0.5pt,arc=1.1mm,left=3mm,right=3mm,top=1.7mm,bottom=1.7mm,before skip=2mm,after skip=2mm]
\textbf{\color{teal}To'g'ri javob:} \textbf{D) 22247424 bayt}
\par\smallskip
\textbf{Izoh:} 21726 \ensuremath{\times} 1024 = 22247424 bayt.
\end{tcolorbox}
\vspace{2mm}\hrule height 0.2pt \vspace{3mm}

% DB question_id: e321e600-a1b7-4393-b0dd-411db42cc791
\needspace{5\baselineskip}
\noindent\textbf{\color{ink}452.}
\par
Quyidagi ta'rif qaysi atamaga tegishli? Tashuvchida saqlanishi mumkin bo'lgan axborotning umumiy hajmi.
\begin{enumerate}[label=\Alph*),leftmargin=9mm]
\item Bit
\item Xotira sig'imi
\item Bayt
\item Xartli formulasi
\end{enumerate}
\begin{tcolorbox}[enhanced,breakable,colback=teal!4,colframe=teal!50!black,boxrule=0.5pt,arc=1.1mm,left=3mm,right=3mm,top=1.7mm,bottom=1.7mm,before skip=2mm,after skip=2mm]
\textbf{\color{teal}To'g'ri javob:} \textbf{B) Xotira sig'imi}
\par\smallskip
\textbf{Izoh:} Ta'rif Xotira sig'imi tushunchasiga tegishli.
\end{tcolorbox}
\vspace{2mm}\hrule height 0.2pt \vspace{3mm}

% DB question_id: e7231d0a-8bb7-57c2-95e6-1005a858b03c
\needspace{5\baselineskip}
\noindent\textbf{\color{ink}453.}
\par
160 sahifali kitobning har sahifasida 48 satr, har satrida 64 belgi bor. Har bir belgi 1 bayt bo'lsa, kitob 21 soat 20 daqiqada o'qib chiqildi. O'rtacha o'qish tezligi necha bayt/sekund?
\begin{enumerate}[label=\Alph*),leftmargin=9mm]
\item 6,4 bayt/s
\item 64 bayt/s
\item 3,2 bayt/s
\item 8 bayt/s
\end{enumerate}
\begin{tcolorbox}[enhanced,breakable,colback=teal!4,colframe=teal!50!black,boxrule=0.5pt,arc=1.1mm,left=3mm,right=3mm,top=1.7mm,bottom=1.7mm,before skip=2mm,after skip=2mm]
\textbf{\color{teal}To'g'ri javob:} \textbf{A) 6,4 bayt/s}
\par\smallskip
\textbf{Izoh:} Hajm 160\ensuremath{\cdot}48\ensuremath{\cdot}64=491 520 bayt. Vaqt 21\ensuremath{\cdot}3600+20\ensuremath{\cdot}60=76 800 s. Tezlik 491 520/76 800=6,4 bayt/s.
\end{tcolorbox}
\vspace{2mm}\hrule height 0.2pt \vspace{3mm}

% DB question_id: ec62293c-8e04-5a70-b0f0-36908a0b6b36
\needspace{5\baselineskip}
\noindent\textbf{\color{ink}454.}
\par
26 harfli ingliz alifbosida Sezar shifri bilan harflar 2 o'rin oldinga siljitiladi; apostrof o'zgarmaydi. \texttt{O'QUVCHI} matnining shifrmatnini toping.
\begin{enumerate}[label=\Alph*),leftmargin=9mm]
\item Q'SWXEJK
\item P'RVWDIJ
\item S'UYZGLM
\item M'OSTAFG
\end{enumerate}
\begin{tcolorbox}[enhanced,breakable,colback=teal!4,colframe=teal!50!black,boxrule=0.5pt,arc=1.1mm,left=3mm,right=3mm,top=1.7mm,bottom=1.7mm,before skip=2mm,after skip=2mm]
\textbf{\color{teal}To'g'ri javob:} \textbf{A) Q'SWXEJK}
\par\smallskip
\textbf{Izoh:} O\ensuremath{\to}Q, Q\ensuremath{\to}S, U\ensuremath{\to}W, V\ensuremath{\to}X, C\ensuremath{\to}E, H\ensuremath{\to}J, I\ensuremath{\to}K; apostrof saqlanadi. Natija Q'SWXEJK.
\end{tcolorbox}
\vspace{2mm}\hrule height 0.2pt \vspace{3mm}

% DB question_id: f3e1a911-906d-5dbb-b9f5-5a596bf7f075
\needspace{5\baselineskip}
\noindent\textbf{\color{ink}455.}
\par
4138 KiB axborot necha baytga teng? 1 KiB = 1024 bayt deb oling.
\begin{enumerate}[label=\Alph*),leftmargin=9mm]
\item 4138000 bayt
\item 4237312 bayt
\item 33898496 bayt
\item 529664 bayt
\end{enumerate}
\begin{tcolorbox}[enhanced,breakable,colback=teal!4,colframe=teal!50!black,boxrule=0.5pt,arc=1.1mm,left=3mm,right=3mm,top=1.7mm,bottom=1.7mm,before skip=2mm,after skip=2mm]
\textbf{\color{teal}To'g'ri javob:} \textbf{B) 4237312 bayt}
\par\smallskip
\textbf{Izoh:} 4138 \ensuremath{\times} 1024 = 4237312 bayt.
\end{tcolorbox}
\vspace{2mm}\hrule height 0.2pt \vspace{3mm}

% DB question_id: 02e46735-fa65-54b0-ba60-5e3206108fc1
\needspace{5\baselineskip}
\noindent\textbf{\color{ink}456.}
\par
Yo'qotishli siqish qachon ma'qul?
\begin{enumerate}[label=\Alph*),leftmargin=9mm]
\item Matnning biror belgisini yo'qotib bo'lmaydigan holatda
\item Parolni saqlashda
\item Qabul qilinadigan sifat yo'qotishi evaziga media hajmini keskin kamaytirishda
\item Aniq hisob jadvalida
\end{enumerate}
\begin{tcolorbox}[enhanced,breakable,colback=teal!4,colframe=teal!50!black,boxrule=0.5pt,arc=1.1mm,left=3mm,right=3mm,top=1.7mm,bottom=1.7mm,before skip=2mm,after skip=2mm]
\textbf{\color{teal}To'g'ri javob:} \textbf{C) Qabul qilinadigan sifat yo'qotishi evaziga media hajmini keskin kamaytirishda}
\par\smallskip
\textbf{Izoh:} Lossy ko'proq tasvir, audio va videoda qo'llanadi.
\end{tcolorbox}
\vspace{2mm}\hrule height 0.2pt \vspace{3mm}

% DB question_id: 050c891c-0798-4d27-bcd2-0c14fbf0a5f4
\needspace{5\baselineskip}
\noindent\textbf{\color{ink}457.}
\par
5 bayt necha bitga teng?
\begin{enumerate}[label=\Alph*),leftmargin=9mm]
\item 13 bit
\item 32 bit
\item 40 bit
\item 48 bit
\end{enumerate}
\begin{tcolorbox}[enhanced,breakable,colback=teal!4,colframe=teal!50!black,boxrule=0.5pt,arc=1.1mm,left=3mm,right=3mm,top=1.7mm,bottom=1.7mm,before skip=2mm,after skip=2mm]
\textbf{\color{teal}To'g'ri javob:} \textbf{C) 40 bit}
\par\smallskip
\textbf{Izoh:} 5 \ensuremath{\times} 8 = 40 bit.
\end{tcolorbox}
\vspace{2mm}\hrule height 0.2pt \vspace{3mm}

% DB question_id: 052ce3ce-6aca-5cf7-8e2f-9db1aa4f9817
\needspace{5\baselineskip}
\noindent\textbf{\color{ink}458.}
\par
Kitob 500 sahifa, har sahifada 35 satr, har satrda 60 belgi. Har belgi 1 bayt. Matn hajmini B da toping.
\begin{enumerate}[label=\Alph*),leftmargin=9mm]
\item 1 050 000 B
\item 8 400 000 B
\item 131 250 B
\item 1 050 003 B
\end{enumerate}
\begin{tcolorbox}[enhanced,breakable,colback=teal!4,colframe=teal!50!black,boxrule=0.5pt,arc=1.1mm,left=3mm,right=3mm,top=1.7mm,bottom=1.7mm,before skip=2mm,after skip=2mm]
\textbf{\color{teal}To'g'ri javob:} \textbf{A) 1 050 000 B}
\par\smallskip
\textbf{Izoh:} 500 \ensuremath{\times} 35 \ensuremath{\times} 60 \ensuremath{\times} 1 = 1 050 000 B = 1 050 000 B.
\end{tcolorbox}
\vspace{2mm}\hrule height 0.2pt \vspace{3mm}

% DB question_id: 05916aa2-c1fb-4182-a542-abffbc21fe6f
\needspace{5\baselineskip}
\noindent\textbf{\color{ink}459.}
\par
Mashina so'zi nimaga bog'liq?
\begin{enumerate}[label=\Alph*),leftmargin=9mm]
\item Protsessor bir amal bilan qayta ishlaydigan bitlar soniga
\item Protsessor takt generatorining bir soniyadagi impulslar chastotasiga
\item Operativ xotiraning umumiy saqlash hajmi nechta baytligiga
\item Protsessorda mavjud bo'lgan mustaqil hisoblash yadrolarining soniga
\end{enumerate}
\begin{tcolorbox}[enhanced,breakable,colback=teal!4,colframe=teal!50!black,boxrule=0.5pt,arc=1.1mm,left=3mm,right=3mm,top=1.7mm,bottom=1.7mm,before skip=2mm,after skip=2mm]
\textbf{\color{teal}To'g'ri javob:} \textbf{A) Protsessor bir amal bilan qayta ishlaydigan bitlar soniga}
\par\smallskip
\textbf{Izoh:} Mashina so'zi odatda protsessor bir amalda qayta ishlay oladigan ma'lumot uzunligi (razryadlilik) bilan bog'liq. Takt chastotasi protsessor tezligini, operativ xotira hajmi saqlash imkoniyatini, yadrolar soni esa parallel hisoblash resurslarini bildiradi.
\end{tcolorbox}
\vspace{2mm}\hrule height 0.2pt \vspace{3mm}

% DB question_id: 0a1b9146-8c64-4a82-8c3b-98c06ed0530c
\needspace{5\baselineskip}
\noindent\textbf{\color{ink}460.}
\par
3 MiB necha KiB?
\begin{enumerate}[label=\Alph*),leftmargin=9mm]
\item 3000 KiB
\item 3144 KiB
\item 4096 KiB
\item 3072 KiB
\end{enumerate}
\begin{tcolorbox}[enhanced,breakable,colback=teal!4,colframe=teal!50!black,boxrule=0.5pt,arc=1.1mm,left=3mm,right=3mm,top=1.7mm,bottom=1.7mm,before skip=2mm,after skip=2mm]
\textbf{\color{teal}To'g'ri javob:} \textbf{D) 3072 KiB}
\par\smallskip
\textbf{Izoh:} 1 MiB = 1024 KiB; 3 \ensuremath{\times} 1024 = 3072.
\end{tcolorbox}
\vspace{2mm}\hrule height 0.2pt \vspace{3mm}

% DB question_id: 0ae910a9-d0a3-5781-bb95-c1c9d4751fd7
\needspace{5\baselineskip}
\noindent\textbf{\color{ink}461.}
\par
200 ta turli belgini doimiy uzunlikdagi ikkilik kod bilan kodlash uchun minimal kod uzunligi qancha?
\begin{enumerate}[label=\Alph*),leftmargin=9mm]
\item 8 bit
\item 7 bit
\item 200 bit
\item 9 bit
\end{enumerate}
\begin{tcolorbox}[enhanced,breakable,colback=teal!4,colframe=teal!50!black,boxrule=0.5pt,arc=1.1mm,left=3mm,right=3mm,top=1.7mm,bottom=1.7mm,before skip=2mm,after skip=2mm]
\textbf{\color{teal}To'g'ri javob:} \textbf{A) 8 bit}
\par\smallskip
\textbf{Izoh:} 2\textsuperscript{n} \ensuremath{\geq} 200; minimal n = 8.
\end{tcolorbox}
\vspace{2mm}\hrule height 0.2pt \vspace{3mm}

% DB question_id: 0c315aa6-3dfb-46ac-b97b-0b449451072f
\needspace{5\baselineskip}
\noindent\textbf{\color{ink}462.}
\par
Maktabdagi 448 o'quvchining har biriga 448 kishining barchasini ajrata oladigan bir xil uzunlikdagi minimal ikkilik ID berilgan. Shu kod bilan 1/4 qismi - 112 o'quvchi IDlarini saqlash uchun necha bayt kerak?
\begin{enumerate}[label=\Alph*),leftmargin=9mm]
\item 112 bayt
\item 126 bayt
\item 144 bayt
\item 504 bayt
\end{enumerate}
\begin{tcolorbox}[enhanced,breakable,colback=teal!4,colframe=teal!50!black,boxrule=0.5pt,arc=1.1mm,left=3mm,right=3mm,top=1.7mm,bottom=1.7mm,before skip=2mm,after skip=2mm]
\textbf{\color{teal}To'g'ri javob:} \textbf{B) 126 bayt}
\par\smallskip
\textbf{Izoh:} 448 turli kodni ajratish uchun 2\textsuperscript{8}=256 yetmaydi, 2\textsuperscript{9}=512 yetadi: bir ID 9 bit. 448/4=112 kishi; 112\ensuremath{\times}9=1008 bit, 1008/8=126 bayt.
\end{tcolorbox}
\vspace{2mm}\hrule height 0.2pt \vspace{3mm}

% DB question_id: 0d063802-d46c-48a6-b0a5-10cda7126273
\needspace{5\baselineskip}
\noindent\textbf{\color{ink}463.}
\par
Svetoforning qizil chirog'i to'xtash talabini bildirdi. Bu vaziyatdagi asosiy tushunchani aniqlang.
\begin{enumerate}[label=\Alph*),leftmargin=9mm]
\item Belgi
\item Alifbo
\item Kod
\item Kodlash
\end{enumerate}
\begin{tcolorbox}[enhanced,breakable,colback=teal!4,colframe=teal!50!black,boxrule=0.5pt,arc=1.1mm,left=3mm,right=3mm,top=1.7mm,bottom=1.7mm,before skip=2mm,after skip=2mm]
\textbf{\color{teal}To'g'ri javob:} \textbf{A) Belgi}
\par\smallskip
\textbf{Izoh:} Vaziyat Belgiga misol: obyekt, xossa yoki tushunchani ifodalash uchun qabul qilingan ishora.
\end{tcolorbox}
\vspace{2mm}\hrule height 0.2pt \vspace{3mm}

% DB question_id: 0d11cc2a-dd3c-4f20-be5e-545759201a3b
\needspace{5\baselineskip}
\noindent\textbf{\color{ink}464.}
\par
215\textsubscript{8} \ensuremath{\leq} x \ensuremath{\leq} 93\textsubscript{1}\textsubscript{6} oraliqdagi JUFT butun sonlarni ikkilikka o'tkazing. Ikkilik yozuvida ketma-ket <<00>> bo'lgan sonlarni chiqarib tashlang. Nechta son qoladi va qolganlar yig'indisi o'nlikda nechaga teng?
\begin{enumerate}[label=\Alph*),leftmargin=9mm]
\item 0 ta son, yig'indi 0
\item 1 ta son, yig'indi 142
\item 2 ta son, yig'indi 288
\item 3 ta son, yig'indi 432
\end{enumerate}
\begin{tcolorbox}[enhanced,breakable,colback=teal!4,colframe=teal!50!black,boxrule=0.5pt,arc=1.1mm,left=3mm,right=3mm,top=1.7mm,bottom=1.7mm,before skip=2mm,after skip=2mm]
\textbf{\color{teal}To'g'ri javob:} \textbf{A) 0 ta son, yig'indi 0}
\par\smallskip
\textbf{Izoh:} 215\textsubscript{8}=141\textsubscript{1}\textsubscript{0} va 93\textsubscript{1}\textsubscript{6}=147\textsubscript{1}\textsubscript{0}. Oraliqdagi juftlar 142=10001110\textsubscript{2}, 144=10010000\textsubscript{2}, 146=10010010\textsubscript{2}. Uchalasi ham ketma-ket 00 ni o'z ichiga oladi; hammasi o'chiriladi va bo'sh yig'indi 0.
\end{tcolorbox}
\vspace{2mm}\hrule height 0.2pt \vspace{3mm}

% DB question_id: 107ebd40-f788-561b-9049-f37aadfa208e
\needspace{5\baselineskip}
\noindent\textbf{\color{ink}465.}
\par
Tasvir 10\ensuremath{\times}15 sm o'lchamda, 150 dpi aniqlikda va 8 rangli palitrada siqilmasdan saqlanadi. 1 dyuym=2,5 sm, sarlavha hisobga olinmasa, hajm necha bayt?
\begin{enumerate}[label=\Alph*),leftmargin=9mm]
\item 202 500 bayt
\item 540 000 bayt
\item 67 500 bayt
\item 1 620 000 bayt
\end{enumerate}
\begin{tcolorbox}[enhanced,breakable,colback=teal!4,colframe=teal!50!black,boxrule=0.5pt,arc=1.1mm,left=3mm,right=3mm,top=1.7mm,bottom=1.7mm,before skip=2mm,after skip=2mm]
\textbf{\color{teal}To'g'ri javob:} \textbf{A) 202 500 bayt}
\par\smallskip
\textbf{Izoh:} 10/2,5\ensuremath{\cdot}150=600 va 15/2,5\ensuremath{\cdot}150=900 piksel. 8 rang uchun 3 bit/piksel. Hajm 600\ensuremath{\cdot}900\ensuremath{\cdot}3/8=202 500 bayt.
\end{tcolorbox}
\vspace{2mm}\hrule height 0.2pt \vspace{3mm}

% DB question_id: 1133c5ee-bda6-45e3-8b8f-43228e02771e
\needspace{5\baselineskip}
\noindent\textbf{\color{ink}466.}
\par
Stereo 22 000 Hz, 24 bit/kanal audio 120 soniya siqilmasdan yozildi. Hajmni eng yaqin butun MiBga yaxlitlang (1 MiB=2\textsuperscript{2}\textsuperscript{0} bayt).
\begin{enumerate}[label=\Alph*),leftmargin=9mm]
\item 12 MiB
\item 14 MiB
\item 15 MiB
\item 20 MiB
\end{enumerate}
\begin{tcolorbox}[enhanced,breakable,colback=teal!4,colframe=teal!50!black,boxrule=0.5pt,arc=1.1mm,left=3mm,right=3mm,top=1.7mm,bottom=1.7mm,before skip=2mm,after skip=2mm]
\textbf{\color{teal}To'g'ri javob:} \textbf{C) 15 MiB}
\par\smallskip
\textbf{Izoh:} Hajm 22000\ensuremath{\times}24\ensuremath{\times}2\ensuremath{\times}120/8=15840000 bayt. 15840000/1048576\ensuremath{\approx}15,106 MiB, yaqin butun son 15 MiB.
\end{tcolorbox}
\vspace{2mm}\hrule height 0.2pt \vspace{3mm}

% DB question_id: 11a2dcda-9ab4-4fce-8d0a-67c1330add8e
\needspace{5\baselineskip}
\noindent\textbf{\color{ink}467.}
\par
6 GiB hajmli flesh-xotira to'liq bo'sh. Doimiy yozish tezligi 512 MiB/s bo'lsa, boshqa kechikishlarsiz to'ldirish necha sekund davom etadi?
\begin{enumerate}[label=\Alph*),leftmargin=9mm]
\item 6 soniya
\item 10 soniya
\item 12 soniya
\item 24 soniya
\end{enumerate}
\begin{tcolorbox}[enhanced,breakable,colback=teal!4,colframe=teal!50!black,boxrule=0.5pt,arc=1.1mm,left=3mm,right=3mm,top=1.7mm,bottom=1.7mm,before skip=2mm,after skip=2mm]
\textbf{\color{teal}To'g'ri javob:} \textbf{C) 12 soniya}
\par\smallskip
\textbf{Izoh:} 6 GiB=6\ensuremath{\times}1024=6144 MiB. 6144 MiB / 512 MiB/s=12 sekund. Birliklar ikkalasi binary GiB/MiB bo'lgani uchun muvofiq hisoblandi.
\end{tcolorbox}
\vspace{2mm}\hrule height 0.2pt \vspace{3mm}

% DB question_id: 15de5a09-74c5-5d4b-9abb-b0a0e0204c87
\needspace{5\baselineskip}
\noindent\textbf{\color{ink}468.}
\par
3 ta turli belgini doimiy uzunlikdagi ikkilik kod bilan kodlash uchun minimal kod uzunligi qancha?
\begin{enumerate}[label=\Alph*),leftmargin=9mm]
\item 1 bit
\item 2 bit
\item 3 bit
\item 5 bit
\end{enumerate}
\begin{tcolorbox}[enhanced,breakable,colback=teal!4,colframe=teal!50!black,boxrule=0.5pt,arc=1.1mm,left=3mm,right=3mm,top=1.7mm,bottom=1.7mm,before skip=2mm,after skip=2mm]
\textbf{\color{teal}To'g'ri javob:} \textbf{B) 2 bit}
\par\smallskip
\textbf{Izoh:} 2\textsuperscript{n} \ensuremath{\geq} 3; minimal n = 2.
\end{tcolorbox}
\vspace{2mm}\hrule height 0.2pt \vspace{3mm}

% DB question_id: 160379d1-e16c-5507-a439-b60136ca8b49
\needspace{5\baselineskip}
\noindent\textbf{\color{ink}469.}
\par
750 MiB axborot 50 s da uzatildi. O'rtacha tezlikni toping.
\begin{enumerate}[label=\Alph*),leftmargin=9mm]
\item 1,88 MiB/s
\item 120 MiB/s
\item 15 MiB/s
\item 900 MiB/s
\end{enumerate}
\begin{tcolorbox}[enhanced,breakable,colback=teal!4,colframe=teal!50!black,boxrule=0.5pt,arc=1.1mm,left=3mm,right=3mm,top=1.7mm,bottom=1.7mm,before skip=2mm,after skip=2mm]
\textbf{\color{teal}To'g'ri javob:} \textbf{C) 15 MiB/s}
\par\smallskip
\textbf{Izoh:} Vaqt sekundga o'tkazilib, v = I/t formula bilan hisoblanadi. Natija: 15 MiB/s.
\end{tcolorbox}
\vspace{2mm}\hrule height 0.2pt \vspace{3mm}

% DB question_id: 1614106b-425c-5462-909c-df0877b215a8
\needspace{5\baselineskip}
\noindent\textbf{\color{ink}470.}
\par
100 ta turli belgini doimiy uzunlikdagi ikkilik kod bilan kodlash uchun minimal kod uzunligi qancha?
\begin{enumerate}[label=\Alph*),leftmargin=9mm]
\item 6 bit
\item 100 bit
\item 8 bit
\item 7 bit
\end{enumerate}
\begin{tcolorbox}[enhanced,breakable,colback=teal!4,colframe=teal!50!black,boxrule=0.5pt,arc=1.1mm,left=3mm,right=3mm,top=1.7mm,bottom=1.7mm,before skip=2mm,after skip=2mm]
\textbf{\color{teal}To'g'ri javob:} \textbf{D) 7 bit}
\par\smallskip
\textbf{Izoh:} 2\textsuperscript{n} \ensuremath{\geq} 100; minimal n = 7.
\end{tcolorbox}
\vspace{2mm}\hrule height 0.2pt \vspace{3mm}

% DB question_id: 17170d02-bc17-5b11-84cf-d657fff3c07e
\needspace{5\baselineskip}
\noindent\textbf{\color{ink}471.}
\par
b va B yozuvlari nimani bildiradi?
\begin{enumerate}[label=\Alph*),leftmargin=9mm]
\item b - bayt, B - bit
\item b - bit, B - bayt
\item Ikkalasi bayt
\item Ikkalasi bit
\end{enumerate}
\begin{tcolorbox}[enhanced,breakable,colback=teal!4,colframe=teal!50!black,boxrule=0.5pt,arc=1.1mm,left=3mm,right=3mm,top=1.7mm,bottom=1.7mm,before skip=2mm,after skip=2mm]
\textbf{\color{teal}To'g'ri javob:} \textbf{B) b - bit, B - bayt}
\par\smallskip
\textbf{Izoh:} Harf registri birlikni o'zgartiradi.
\end{tcolorbox}
\vspace{2mm}\hrule height 0.2pt \vspace{3mm}

% DB question_id: 19f15cc2-1c4c-5f41-aabe-a2b23d5a13d7
\needspace{5\baselineskip}
\noindent\textbf{\color{ink}472.}
\par
250 ta belgi, har biri 3 baytdan kodlangan. Xabar hajmi qancha?
\begin{enumerate}[label=\Alph*),leftmargin=9mm]
\item 250 B
\item 750 B
\item 6 000 B
\item 93,75 B
\end{enumerate}
\begin{tcolorbox}[enhanced,breakable,colback=teal!4,colframe=teal!50!black,boxrule=0.5pt,arc=1.1mm,left=3mm,right=3mm,top=1.7mm,bottom=1.7mm,before skip=2mm,after skip=2mm]
\textbf{\color{teal}To'g'ri javob:} \textbf{B) 750 B}
\par\smallskip
\textbf{Izoh:} Hajm = 250 \ensuremath{\times} 3 = 750 B.
\end{tcolorbox}
\vspace{2mm}\hrule height 0.2pt \vspace{3mm}

% DB question_id: 1a91d252-8416-4aea-9fcd-227918aa9497
\needspace{5\baselineskip}
\noindent\textbf{\color{ink}473.}
\par
Ikki xonali xy\textsubscript{9} soni 47\textsubscript{8} ga teng. Bu yerda x va y - to'qqizlik sanoq sistemasining raqamlari. x\ensuremath{\cdot}y ko'paytmasini ikkilik sanoq sistemasida toping.
\begin{enumerate}[label=\Alph*),leftmargin=9mm]
\item 1001\textsubscript{2}
\item 1010\textsubscript{2}
\item 1100\textsubscript{2}
\item 1101\textsubscript{2}
\end{enumerate}
\begin{tcolorbox}[enhanced,breakable,colback=teal!4,colframe=teal!50!black,boxrule=0.5pt,arc=1.1mm,left=3mm,right=3mm,top=1.7mm,bottom=1.7mm,before skip=2mm,after skip=2mm]
\textbf{\color{teal}To'g'ri javob:} \textbf{C) 1100\textsubscript{2}}
\par\smallskip
\textbf{Izoh:} 47\textsubscript{8}=4\ensuremath{\cdot}8+7=39\textsubscript{1}\textsubscript{0}. 9x+y=39 va 0\ensuremath{\leq}x,y\ensuremath{\leq}8, demak x=4, y=3. Ularning ko'paytmasi 12\textsubscript{1}\textsubscript{0}=1100\textsubscript{2}.
\end{tcolorbox}
\vspace{2mm}\hrule height 0.2pt \vspace{3mm}

% DB question_id: 1b18f5f0-0f7c-5400-beb6-d95689e2d378
\needspace{5\baselineskip}
\noindent\textbf{\color{ink}474.}
\par
Vijiner shifrida ingliz alifbosi \texttt{A=0,...,Z=25} tartibida olinadi, ochiq matn va takrorlangan kalit modullari qo'shilib 26 bo'yicha qoldiq olinadi. \texttt{ATTESTATSIYA} matnini \texttt{KITOB} kaliti bilan shifrlang.
\begin{enumerate}[label=\Alph*),leftmargin=9mm]
\item KBMSTDIMGJII
\item LCNTUEJNHKJJ
\item KBLRSCILFHHI
\item QGEPAQDLQGEP
\end{enumerate}
\begin{tcolorbox}[enhanced,breakable,colback=teal!4,colframe=teal!50!black,boxrule=0.5pt,arc=1.1mm,left=3mm,right=3mm,top=1.7mm,bottom=1.7mm,before skip=2mm,after skip=2mm]
\textbf{\color{teal}To'g'ri javob:} \textbf{A) KBMSTDIMGJII}
\par\smallskip
\textbf{Izoh:} Kalit KITOBKITOBKI ko'rinishida takrorlanadi. Harflar modullari 26 bo'yicha qo'shilganda KBMSTDIMGJII hosil bo'ladi.
\end{tcolorbox}
\vspace{2mm}\hrule height 0.2pt \vspace{3mm}

% DB question_id: 1c1d98f7-0e92-42ab-97de-b2a93fb36b1f
\needspace{5\baselineskip}
\noindent\textbf{\color{ink}475.}
\par
Internet ulanishida kanalning reklama qilingan maksimal tezligidan ko'ra amalda foydali ma'lumotning vaqt birligida qancha uzatilishi o'lchandi. Throughput atamasi nimani anglatadi?
\begin{enumerate}[label=\Alph*),leftmargin=9mm]
\item Amalda erishilgan foydali uzatish tezligini
\item Tarmoq kanalining maksimal nazariy uzatish sig'imini
\item Uzatish davomida bir paketning kechikish vaqtini
\item Aloqa kanalidagi ma'lumot yo'qotilishining umumiy ulushini
\end{enumerate}
\begin{tcolorbox}[enhanced,breakable,colback=teal!4,colframe=teal!50!black,boxrule=0.5pt,arc=1.1mm,left=3mm,right=3mm,top=1.7mm,bottom=1.7mm,before skip=2mm,after skip=2mm]
\textbf{\color{teal}To'g'ri javob:} \textbf{A) Amalda erishilgan foydali uzatish tezligini}
\par\smallskip
\textbf{Izoh:} Throughput amalda foydali yetkazib berilgan ma'lumot miqdorining vaqtga nisbatidir. Nazariy kanal sig'imi bandwidth, kechikish latency va paket yo'qotish ulushi boshqa ko'rsatkichlardir.
\end{tcolorbox}
\vspace{2mm}\hrule height 0.2pt \vspace{3mm}

% DB question_id: 1ef6977e-808c-4f74-845d-d12a506afe34
\needspace{5\baselineskip}
\noindent\textbf{\color{ink}476.}
\par
64 bit qiymat 8 baytga aylantirildi. Bu vaziyatdagi asosiy tushunchani aniqlang.
\begin{enumerate}[label=\Alph*),leftmargin=9mm]
\item Bit va bayt farqi
\item Rastr va vektor farqi
\item Kodek va konteyner farqi
\item Validatsiya va verifikatsiya farqi
\end{enumerate}
\begin{tcolorbox}[enhanced,breakable,colback=teal!4,colframe=teal!50!black,boxrule=0.5pt,arc=1.1mm,left=3mm,right=3mm,top=1.7mm,bottom=1.7mm,before skip=2mm,after skip=2mm]
\textbf{\color{teal}To'g'ri javob:} \textbf{A) Bit va bayt farqi}
\par\smallskip
\textbf{Izoh:} Vaziyat Bit va bayt farqiga misol: bit ikkilik birlik, bayt esa 8 bitdan iborat ekanini ajratish.
\end{tcolorbox}
\vspace{2mm}\hrule height 0.2pt \vspace{3mm}

% DB question_id: 24b7eed9-d71f-5fe1-9dfa-2b70283836b9
\needspace{5\baselineskip}
\noindent\textbf{\color{ink}477.}
\par
Kanalning nazariy tezligi 80 Mbit/s, amaliy throughput 52 Mbit/s. Samaradorlik necha foiz?
\begin{enumerate}[label=\Alph*),leftmargin=9mm]
\item 0,65 \%
\item 35 \%
\item 8,12 \%
\item 65 \%
\end{enumerate}
\begin{tcolorbox}[enhanced,breakable,colback=teal!4,colframe=teal!50!black,boxrule=0.5pt,arc=1.1mm,left=3mm,right=3mm,top=1.7mm,bottom=1.7mm,before skip=2mm,after skip=2mm]
\textbf{\color{teal}To'g'ri javob:} \textbf{D) 65 \%}
\par\smallskip
\textbf{Izoh:} \ensuremath{\eta} = 52/80 \ensuremath{\times} 100\% = 65\%.
\end{tcolorbox}
\vspace{2mm}\hrule height 0.2pt \vspace{3mm}

% DB question_id: 25b1b026-e443-4d57-843e-f05162be9f39
\needspace{5\baselineskip}
\noindent\textbf{\color{ink}478.}
\par
16 bit/belgi kodlashda matndan bitta viloyat nomi bilan unga yondosh vergul va bitta bo'sh joy o'chirildi; hajm 26 baytga kamaydi. \texttt{Qashqadaryo}, \texttt{Surxandaryo}, \texttt{Farg'ona}, \texttt{Xorazm} orasida 11 harfli ikki nom bo'lsa, matnda birinchi uchraganini tanlang. Qaysi nom o'chirilgan?
\begin{enumerate}[label=\Alph*),leftmargin=9mm]
\item Surxandaryo
\item Farg'ona
\item Qashqadaryo
\item Xorazm
\end{enumerate}
\begin{tcolorbox}[enhanced,breakable,colback=teal!4,colframe=teal!50!black,boxrule=0.5pt,arc=1.1mm,left=3mm,right=3mm,top=1.7mm,bottom=1.7mm,before skip=2mm,after skip=2mm]
\textbf{\color{teal}To'g'ri javob:} \textbf{C) Qashqadaryo}
\par\smallskip
\textbf{Izoh:} 26 bayt=208 bit, har belgi 16 bit \ensuremath{\Rightarrow} 13 belgi o'chgan. Vergul va bo'sh joy 2 belgi, nom 11 ta harf. Qashqadaryo va Surxandaryo 11 harfli; shartdagi <<birinchi uchragani>> qoidasi bo'yicha Qashqadaryo tanlanadi.
\end{tcolorbox}
\vspace{2mm}\hrule height 0.2pt \vspace{3mm}

% DB question_id: 26fcf900-9b92-5949-9d08-c936041c87d9
\needspace{5\baselineskip}
\noindent\textbf{\color{ink}479.}
\par
Asl hajm 900 MB, siqilgan hajm 300 MB. Siqish koeffitsiyenti qancha?
\begin{enumerate}[label=\Alph*),leftmargin=9mm]
\item 897 marta
\item 900 marta
\item 1,5 marta
\item 3 marta
\end{enumerate}
\begin{tcolorbox}[enhanced,breakable,colback=teal!4,colframe=teal!50!black,boxrule=0.5pt,arc=1.1mm,left=3mm,right=3mm,top=1.7mm,bottom=1.7mm,before skip=2mm,after skip=2mm]
\textbf{\color{teal}To'g'ri javob:} \textbf{D) 3 marta}
\par\smallskip
\textbf{Izoh:} 900/300 = 3 marta.
\end{tcolorbox}
\vspace{2mm}\hrule height 0.2pt \vspace{3mm}

% DB question_id: 27873e4d-47fb-50e0-b476-cde82614cb6f
\needspace{5\baselineskip}
\noindent\textbf{\color{ink}480.}
\par
Jadval katagi \texttt{(satr; ustun)} bilan berilgan:
\medskip\par
| |0|1|2|3|4|5|6|7|8|9|
\par
|---|---|---|---|---|---|---|---|---|---|---|
\par
|0|a|x|b|o|r|o|t| |t|e|
\par
|1|x|n|o|l|o|g|i|y|a|l|
\par
|2|a|r|i| |h|a|q|i|d|a|
\par
|3|m|a|q|o|l|a|l|a|r| |
\medskip\par
\texttt{(0;3), (3;0), (3;7), (2;8)} kataklaridan olingan so'z Sezar usulida 3 o'rin oldinga siljitildi. Natijani toping.
\begin{enumerate}[label=\Alph*),leftmargin=9mm]
\item RPDG
\item OMAD
\item RPDO
\item QOCF
\end{enumerate}
\begin{tcolorbox}[enhanced,breakable,colback=teal!4,colframe=teal!50!black,boxrule=0.5pt,arc=1.1mm,left=3mm,right=3mm,top=1.7mm,bottom=1.7mm,before skip=2mm,after skip=2mm]
\textbf{\color{teal}To'g'ri javob:} \textbf{A) RPDG}
\par\smallskip
\textbf{Izoh:} Koordinatalar O, M, A, D harflarini beradi: OMAD. Har birini 3 o'rin oldinga sursak O\ensuremath{\to}R, M\ensuremath{\to}P, A\ensuremath{\to}D, D\ensuremath{\to}G; natija RPDG.
\end{tcolorbox}
\vspace{2mm}\hrule height 0.2pt \vspace{3mm}

% DB question_id: 284cc577-8f80-421c-8e14-dd01d53282d4
\needspace{5\baselineskip}
\noindent\textbf{\color{ink}481.}
\par
Bir daqiqa davom etgan siqilmagan stereo audio 48 000 Hz, 16 bit/kanal bilan yozildi. Fayl hajmi nechchi KiB? 1 KiB=1024 bayt, sarlavha hisobga olinmaydi.
\begin{enumerate}[label=\Alph*),leftmargin=9mm]
\item 11 250 KiB
\item 11 520 KiB
\item 5 625 KiB
\item 22 500 KiB
\end{enumerate}
\begin{tcolorbox}[enhanced,breakable,colback=teal!4,colframe=teal!50!black,boxrule=0.5pt,arc=1.1mm,left=3mm,right=3mm,top=1.7mm,bottom=1.7mm,before skip=2mm,after skip=2mm]
\textbf{\color{teal}To'g'ri javob:} \textbf{A) 11 250 KiB}
\par\smallskip
\textbf{Izoh:} Bir sekunddagi hajm 48000\ensuremath{\times}16\ensuremath{\times}2/8=192000 bayt. 60 sekundda 11520000 bayt. 1024 ga bo'lganda 11250 KiB.
\end{tcolorbox}
\vspace{2mm}\hrule height 0.2pt \vspace{3mm}

% DB question_id: 28656acf-7cb1-5381-b144-52ec4f7651ee
\needspace{5\baselineskip}
\noindent\textbf{\color{ink}482.}
\par
Ikkilik kod qaysi alifbodan foydalanadi?
\begin{enumerate}[label=\Alph*),leftmargin=9mm]
\item 0-9
\item 0 va 1
\item A-F
\item Nuqta va vergul
\end{enumerate}
\begin{tcolorbox}[enhanced,breakable,colback=teal!4,colframe=teal!50!black,boxrule=0.5pt,arc=1.1mm,left=3mm,right=3mm,top=1.7mm,bottom=1.7mm,before skip=2mm,after skip=2mm]
\textbf{\color{teal}To'g'ri javob:} \textbf{B) 0 va 1}
\par\smallskip
\textbf{Izoh:} Ikkilik alifbo ikki belgidan iborat.
\end{tcolorbox}
\vspace{2mm}\hrule height 0.2pt \vspace{3mm}

% DB question_id: 28ee8950-2d01-56b7-bcee-4c9f88f88cd1
\needspace{5\baselineskip}
\noindent\textbf{\color{ink}483.}
\par
768 KiB nechta MiB ga teng? Savolda yozilgan SI/IEC birliklariga aynan amal qiling.
\begin{enumerate}[label=\Alph*),leftmargin=9mm]
\item 0,094 MiB
\item 6 MiB
\item 0,75 MiB
\item 768 MiB
\end{enumerate}
\begin{tcolorbox}[enhanced,breakable,colback=teal!4,colframe=teal!50!black,boxrule=0.5pt,arc=1.1mm,left=3mm,right=3mm,top=1.7mm,bottom=1.7mm,before skip=2mm,after skip=2mm]
\textbf{\color{teal}To'g'ri javob:} \textbf{C) 0,75 MiB}
\par\smallskip
\textbf{Izoh:} Birliklar orasidagi ko'paytirish yoki bo'lish natijasi 0,75 MiB.
\end{tcolorbox}
\vspace{2mm}\hrule height 0.2pt \vspace{3mm}

% DB question_id: 28f83a8d-7276-5c85-ae0f-2abaf47f073f
\needspace{5\baselineskip}
\noindent\textbf{\color{ink}484.}
\par
O'zgaruvchan uzunlikdagi kodning afzalligi nimada?
\begin{enumerate}[label=\Alph*),leftmargin=9mm]
\item Har doim dekodlab bo'lmaydi
\item Faqat 8 bit ishlatadi
\item Rang chuqurligini oshiradi
\item Ko'p uchraydigan belgilarni qisqa kodlab hajmni kamaytirishi mumkin
\end{enumerate}
\begin{tcolorbox}[enhanced,breakable,colback=teal!4,colframe=teal!50!black,boxrule=0.5pt,arc=1.1mm,left=3mm,right=3mm,top=1.7mm,bottom=1.7mm,before skip=2mm,after skip=2mm]
\textbf{\color{teal}To'g'ri javob:} \textbf{D) Ko'p uchraydigan belgilarni qisqa kodlab hajmni kamaytirishi mumkin}
\par\smallskip
\textbf{Izoh:} Chastotaga mos kod uzunligi siqishni ta'minlaydi.
\end{tcolorbox}
\vspace{2mm}\hrule height 0.2pt \vspace{3mm}

% DB question_id: 2d874e74-7e66-53fc-8a4e-4a1863ff50d1
\needspace{5\baselineskip}
\noindent\textbf{\color{ink}485.}
\par
1024 MiB fayl 8:1 siqish koeffitsiyenti bilan siqildi. Yangi hajm qancha?
\begin{enumerate}[label=\Alph*),leftmargin=9mm]
\item 896 MiB
\item 128 MiB
\item 1 024 MiB
\item 64 MiB
\end{enumerate}
\begin{tcolorbox}[enhanced,breakable,colback=teal!4,colframe=teal!50!black,boxrule=0.5pt,arc=1.1mm,left=3mm,right=3mm,top=1.7mm,bottom=1.7mm,before skip=2mm,after skip=2mm]
\textbf{\color{teal}To'g'ri javob:} \textbf{B) 128 MiB}
\par\smallskip
\textbf{Izoh:} 1024/8 = 128 MiB.
\end{tcolorbox}
\vspace{2mm}\hrule height 0.2pt \vspace{3mm}

% DB question_id: 2daf5ee2-a59e-5190-b5cd-51916a0bec96
\needspace{5\baselineskip}
\noindent\textbf{\color{ink}486.}
\par
Kitob 64 sahifa, har sahifada 48 satr, har satrda 64 belgi. Har belgi 1 bayt. Matn hajmini KiB da toping.
\begin{enumerate}[label=\Alph*),leftmargin=9mm]
\item 196 608 KiB
\item 192 KiB
\item 1 572 864 KiB
\item 196,608 KiB
\end{enumerate}
\begin{tcolorbox}[enhanced,breakable,colback=teal!4,colframe=teal!50!black,boxrule=0.5pt,arc=1.1mm,left=3mm,right=3mm,top=1.7mm,bottom=1.7mm,before skip=2mm,after skip=2mm]
\textbf{\color{teal}To'g'ri javob:} \textbf{B) 192 KiB}
\par\smallskip
\textbf{Izoh:} 64 \ensuremath{\times} 48 \ensuremath{\times} 64 \ensuremath{\times} 1 = 196 608 B = 192 KiB.
\end{tcolorbox}
\vspace{2mm}\hrule height 0.2pt \vspace{3mm}

% DB question_id: 2e9b2113-476f-5ac9-acbc-c2aa4bcc1a63
\needspace{5\baselineskip}
\noindent\textbf{\color{ink}487.}
\par
Bayt qancha?
\begin{enumerate}[label=\Alph*),leftmargin=9mm]
\item 8 bit
\item 4 bit
\item 16 bit
\item 1000 bit
\end{enumerate}
\begin{tcolorbox}[enhanced,breakable,colback=teal!4,colframe=teal!50!black,boxrule=0.5pt,arc=1.1mm,left=3mm,right=3mm,top=1.7mm,bottom=1.7mm,before skip=2mm,after skip=2mm]
\textbf{\color{teal}To'g'ri javob:} \textbf{A) 8 bit}
\par\smallskip
\textbf{Izoh:} 1 bayt = 8 bit.
\end{tcolorbox}
\vspace{2mm}\hrule height 0.2pt \vspace{3mm}

% DB question_id: 2e9c6eb6-1050-4679-8703-beda89396b4d
\needspace{5\baselineskip}
\noindent\textbf{\color{ink}488.}
\par
960 MB faylning 25\%i 6 MB/s, qolgan 720 MBning 1/4 qismi 18 MB/s tezlikda uzatildi. Oxirgi 540 MB uchun avvalgi ikki qismning jami vaqtiga teng vaqt sarflandi. Butun faylning o'rtacha uzatish tezligi nechchi MB/s?
\begin{enumerate}[label=\Alph*),leftmargin=9mm]
\item 6 MB/s
\item 9,6 MB/s
\item 12 MB/s
\item 19,2 MB/s
\end{enumerate}
\begin{tcolorbox}[enhanced,breakable,colback=teal!4,colframe=teal!50!black,boxrule=0.5pt,arc=1.1mm,left=3mm,right=3mm,top=1.7mm,bottom=1.7mm,before skip=2mm,after skip=2mm]
\textbf{\color{teal}To'g'ri javob:} \textbf{B) 9,6 MB/s}
\par\smallskip
\textbf{Izoh:} Birinchi qism 240/6=40 soniya, ikkinchi 180/18=10 soniya. Qolgan 540 MB ham 40+10=50 soniyada uzatiladi; jami 100 soniya. O'rtacha 960/100=9,6 MB/s.
\end{tcolorbox}
\vspace{2mm}\hrule height 0.2pt \vspace{3mm}

% DB question_id: 320cbae5-7fcc-5de0-8371-8b7961cb91a3
\needspace{5\baselineskip}
\noindent\textbf{\color{ink}489.}
\par
Kitob 220 sahifa, har sahifada 28 satr, har satrda 68 belgi. Har belgi 1 bayt. Matn hajmini KiB da toping.
\begin{enumerate}[label=\Alph*),leftmargin=9mm]
\item 418 880 KiB
\item 409,062 KiB
\item 3 351 040 KiB
\item 418,88 KiB
\end{enumerate}
\begin{tcolorbox}[enhanced,breakable,colback=teal!4,colframe=teal!50!black,boxrule=0.5pt,arc=1.1mm,left=3mm,right=3mm,top=1.7mm,bottom=1.7mm,before skip=2mm,after skip=2mm]
\textbf{\color{teal}To'g'ri javob:} \textbf{B) 409,062 KiB}
\par\smallskip
\textbf{Izoh:} 220 \ensuremath{\times} 28 \ensuremath{\times} 68 \ensuremath{\times} 1 = 418 880 B = 409,062 KiB.
\end{tcolorbox}
\vspace{2mm}\hrule height 0.2pt \vspace{3mm}

% DB question_id: 3451a189-57c3-554c-b59c-ea8b76d61e20
\needspace{5\baselineskip}
\noindent\textbf{\color{ink}490.}
\par
Kitob 50 sahifa, har sahifada 40 satr, har satrda 80 belgi. Har belgi 2 bayt. Matn hajmini KiB da toping.
\begin{enumerate}[label=\Alph*),leftmargin=9mm]
\item 160 000 KiB
\item 2 560 000 KiB
\item 312,5 KiB
\item 320 KiB
\end{enumerate}
\begin{tcolorbox}[enhanced,breakable,colback=teal!4,colframe=teal!50!black,boxrule=0.5pt,arc=1.1mm,left=3mm,right=3mm,top=1.7mm,bottom=1.7mm,before skip=2mm,after skip=2mm]
\textbf{\color{teal}To'g'ri javob:} \textbf{C) 312,5 KiB}
\par\smallskip
\textbf{Izoh:} 50 \ensuremath{\times} 40 \ensuremath{\times} 80 \ensuremath{\times} 2 = 320 000 B = 312,5 KiB.
\end{tcolorbox}
\vspace{2mm}\hrule height 0.2pt \vspace{3mm}

% DB question_id: 3863d52e-59b9-5c72-a90b-8d69202fdb74
\needspace{5\baselineskip}
\noindent\textbf{\color{ink}491.}
\par
0,5 GB nechta MB ga teng? Savolda yozilgan SI/IEC birliklariga aynan amal qiling.
\begin{enumerate}[label=\Alph*),leftmargin=9mm]
\item 500 MB
\item 62,5 MB
\item 4 000 MB
\item 512 000 MB
\end{enumerate}
\begin{tcolorbox}[enhanced,breakable,colback=teal!4,colframe=teal!50!black,boxrule=0.5pt,arc=1.1mm,left=3mm,right=3mm,top=1.7mm,bottom=1.7mm,before skip=2mm,after skip=2mm]
\textbf{\color{teal}To'g'ri javob:} \textbf{A) 500 MB}
\par\smallskip
\textbf{Izoh:} Birliklar orasidagi ko'paytirish yoki bo'lish natijasi 500 MB.
\end{tcolorbox}
\vspace{2mm}\hrule height 0.2pt \vspace{3mm}

% DB question_id: 39394f73-9841-5872-8633-3f964e2ca826
\needspace{5\baselineskip}
\noindent\textbf{\color{ink}492.}
\par
Kanalning nazariy tezligi 200 Mbit/s, amaliy throughput 150 Mbit/s. Samaradorlik necha foiz?
\begin{enumerate}[label=\Alph*),leftmargin=9mm]
\item 0,75 \%
\item 25 \%
\item 9,38 \%
\item 75 \%
\end{enumerate}
\begin{tcolorbox}[enhanced,breakable,colback=teal!4,colframe=teal!50!black,boxrule=0.5pt,arc=1.1mm,left=3mm,right=3mm,top=1.7mm,bottom=1.7mm,before skip=2mm,after skip=2mm]
\textbf{\color{teal}To'g'ri javob:} \textbf{D) 75 \%}
\par\smallskip
\textbf{Izoh:} \ensuremath{\eta} = 150/200 \ensuremath{\times} 100\% = 75\%.
\end{tcolorbox}
\vspace{2mm}\hrule height 0.2pt \vspace{3mm}

% DB question_id: 3af71df6-e919-4238-8df4-94890bdac4b4
\needspace{5\baselineskip}
\noindent\textbf{\color{ink}493.}
\par
Ikkilik yozuvda faqat 0 va 1 belgilaridan foydalanildi. Bu vaziyatdagi asosiy tushunchani aniqlang.
\begin{enumerate}[label=\Alph*),leftmargin=9mm]
\item Kod
\item Alifbo
\item Kodlash
\item Dekodlash
\end{enumerate}
\begin{tcolorbox}[enhanced,breakable,colback=teal!4,colframe=teal!50!black,boxrule=0.5pt,arc=1.1mm,left=3mm,right=3mm,top=1.7mm,bottom=1.7mm,before skip=2mm,after skip=2mm]
\textbf{\color{teal}To'g'ri javob:} \textbf{B) Alifbo}
\par\smallskip
\textbf{Izoh:} Vaziyat Alifboga misol: kodlashda foydalaniladigan ruxsat etilgan belgilar to'plami.
\end{tcolorbox}
\vspace{2mm}\hrule height 0.2pt \vspace{3mm}

% DB question_id: 3bf240cd-140f-4572-93a0-a61924ba7f9b
\needspace{5\baselineskip}
\noindent\textbf{\color{ink}494.}
\par
Kod jadvali nima?
\begin{enumerate}[label=\Alph*),leftmargin=9mm]
\item Belgilar uchun mos ko'rinishlar va shrift tasvirlari ro'yxati
\item Ma'lumot fayllarining nomlari va diskdagi manzillari ro'yxati
\item Belgilar va ularga mos kodlar ro'yxati
\item Kompyuter tarmog'idagi qurilmalar va ularning IP manzillari ro'yxati
\end{enumerate}
\begin{tcolorbox}[enhanced,breakable,colback=teal!4,colframe=teal!50!black,boxrule=0.5pt,arc=1.1mm,left=3mm,right=3mm,top=1.7mm,bottom=1.7mm,before skip=2mm,after skip=2mm]
\textbf{\color{teal}To'g'ri javob:} \textbf{C) Belgilar va ularga mos kodlar ro'yxati}
\par\smallskip
\textbf{Izoh:} Kod jadvali belgilarni ularga mos sonli yoki ikkilik kodlar bilan bog'laydi. Shriftlar belgining ko'rinishini, fayl kataloglari saqlangan obyektlarning joyini, tarmoq jadvallari esa qurilma manzillarini aks ettiradi.
\end{tcolorbox}
\vspace{2mm}\hrule height 0.2pt \vspace{3mm}

% DB question_id: 3ed47751-4392-5060-aa5a-909d1c47005e
\needspace{5\baselineskip}
\noindent\textbf{\color{ink}495.}
\par
640 kbit nechta kB ga teng? Savolda yozilgan SI/IEC birliklariga aynan amal qiling.
\begin{enumerate}[label=\Alph*),leftmargin=9mm]
\item 10 kB
\item 80 kB
\item 640 kB
\item 81 920 kB
\end{enumerate}
\begin{tcolorbox}[enhanced,breakable,colback=teal!4,colframe=teal!50!black,boxrule=0.5pt,arc=1.1mm,left=3mm,right=3mm,top=1.7mm,bottom=1.7mm,before skip=2mm,after skip=2mm]
\textbf{\color{teal}To'g'ri javob:} \textbf{B) 80 kB}
\par\smallskip
\textbf{Izoh:} Birliklar orasidagi ko'paytirish yoki bo'lish natijasi 80 kB.
\end{tcolorbox}
\vspace{2mm}\hrule height 0.2pt \vspace{3mm}

% DB question_id: 423fd42d-4c7c-5e05-b888-9346be74dd27
\needspace{5\baselineskip}
\noindent\textbf{\color{ink}496.}
\par
1200 MB fayl hajmi 35\% ga kamaydi. Siqilgan hajm qancha?
\begin{enumerate}[label=\Alph*),leftmargin=9mm]
\item 420 MB
\item 780 MB
\item 27 300 MB
\item 65 MB
\end{enumerate}
\begin{tcolorbox}[enhanced,breakable,colback=teal!4,colframe=teal!50!black,boxrule=0.5pt,arc=1.1mm,left=3mm,right=3mm,top=1.7mm,bottom=1.7mm,before skip=2mm,after skip=2mm]
\textbf{\color{teal}To'g'ri javob:} \textbf{B) 780 MB}
\par\smallskip
\textbf{Izoh:} 1200 \ensuremath{\times} (1-35/100) = 780 MB.
\end{tcolorbox}
\vspace{2mm}\hrule height 0.2pt \vspace{3mm}

% DB question_id: 447567c6-fafa-56aa-a6bc-b6af60fd32b3
\needspace{5\baselineskip}
\noindent\textbf{\color{ink}497.}
\par
1 024 Mbit nechta MB ga teng? Savolda yozilgan SI/IEC birliklariga aynan amal qiling.
\begin{enumerate}[label=\Alph*),leftmargin=9mm]
\item 16 MB
\item 1 024 MB
\item 131 072 MB
\item 128 MB
\end{enumerate}
\begin{tcolorbox}[enhanced,breakable,colback=teal!4,colframe=teal!50!black,boxrule=0.5pt,arc=1.1mm,left=3mm,right=3mm,top=1.7mm,bottom=1.7mm,before skip=2mm,after skip=2mm]
\textbf{\color{teal}To'g'ri javob:} \textbf{D) 128 MB}
\par\smallskip
\textbf{Izoh:} Birliklar orasidagi ko'paytirish yoki bo'lish natijasi 128 MB.
\end{tcolorbox}
\vspace{2mm}\hrule height 0.2pt \vspace{3mm}

% DB question_id: 4552f2f7-d6aa-597d-97c3-39cb8775d1b7
\needspace{5\baselineskip}
\noindent\textbf{\color{ink}498.}
\par
80 MB hajmli fayl 20 Mbit/s tezlikda uzatilsa, vaqt qancha?
\begin{enumerate}[label=\Alph*),leftmargin=9mm]
\item 4 s
\item 256 s
\item 40 s
\item 32 s
\end{enumerate}
\begin{tcolorbox}[enhanced,breakable,colback=teal!4,colframe=teal!50!black,boxrule=0.5pt,arc=1.1mm,left=3mm,right=3mm,top=1.7mm,bottom=1.7mm,before skip=2mm,after skip=2mm]
\textbf{\color{teal}To'g'ri javob:} \textbf{D) 32 s}
\par\smallskip
\textbf{Izoh:} Hajm bitga o'tkazilib, t = I/v formula bilan hisoblanadi. Natija: 32 s.
\end{tcolorbox}
\vspace{2mm}\hrule height 0.2pt \vspace{3mm}

% DB question_id: 46dafd34-f747-527b-8d20-e9d5fc7a9c4b
\needspace{5\baselineskip}
\noindent\textbf{\color{ink}499.}
\par
Kitob 100 sahifa, har sahifada 30 satr, har satrda 60 belgi. Har belgi 1 bayt. Matn hajmini B da toping.
\begin{enumerate}[label=\Alph*),leftmargin=9mm]
\item 1 440 000 B
\item 22 500 B
\item 180 000 B
\item 180 003 B
\end{enumerate}
\begin{tcolorbox}[enhanced,breakable,colback=teal!4,colframe=teal!50!black,boxrule=0.5pt,arc=1.1mm,left=3mm,right=3mm,top=1.7mm,bottom=1.7mm,before skip=2mm,after skip=2mm]
\textbf{\color{teal}To'g'ri javob:} \textbf{C) 180 000 B}
\par\smallskip
\textbf{Izoh:} 100 \ensuremath{\times} 30 \ensuremath{\times} 60 \ensuremath{\times} 1 = 180 000 B = 180 000 B.
\end{tcolorbox}
\vspace{2mm}\hrule height 0.2pt \vspace{3mm}

% DB question_id: 47d9d7f2-a801-5c56-a8d8-a74c5bffcde2
\needspace{5\baselineskip}
\noindent\textbf{\color{ink}500.}
\par
2 GiB nechta MiB ga teng? Savolda yozilgan SI/IEC birliklariga aynan amal qiling.
\begin{enumerate}[label=\Alph*),leftmargin=9mm]
\item 256 MiB
\item 16 384 MiB
\item 2 048 MiB
\item 2 MiB
\end{enumerate}
\begin{tcolorbox}[enhanced,breakable,colback=teal!4,colframe=teal!50!black,boxrule=0.5pt,arc=1.1mm,left=3mm,right=3mm,top=1.7mm,bottom=1.7mm,before skip=2mm,after skip=2mm]
\textbf{\color{teal}To'g'ri javob:} \textbf{C) 2 048 MiB}
\par\smallskip
\textbf{Izoh:} Birliklar orasidagi ko'paytirish yoki bo'lish natijasi 2 048 MiB.
\end{tcolorbox}
\vspace{2mm}\hrule height 0.2pt \vspace{3mm}

% DB question_id: 48669f4d-98d9-5f48-b4f5-ed432e08235c
\needspace{5\baselineskip}
\noindent\textbf{\color{ink}501.}
\par
7 MiB nechta B ga teng? Savolda yozilgan SI/IEC birliklariga aynan amal qiling.
\begin{enumerate}[label=\Alph*),leftmargin=9mm]
\item 7 340 032 B
\item 917 504 B
\item 58 720 256 B
\item 7 168 B
\end{enumerate}
\begin{tcolorbox}[enhanced,breakable,colback=teal!4,colframe=teal!50!black,boxrule=0.5pt,arc=1.1mm,left=3mm,right=3mm,top=1.7mm,bottom=1.7mm,before skip=2mm,after skip=2mm]
\textbf{\color{teal}To'g'ri javob:} \textbf{A) 7 340 032 B}
\par\smallskip
\textbf{Izoh:} Birliklar orasidagi ko'paytirish yoki bo'lish natijasi 7 340 032 B.
\end{tcolorbox}
\vspace{2mm}\hrule height 0.2pt \vspace{3mm}

% DB question_id: 4dc2c0e8-fa0c-5a4e-b16b-ba20dad31d1f
\needspace{5\baselineskip}
\noindent\textbf{\color{ink}502.}
\par
Kanalning nazariy tezligi 1000 Mbit/s, amaliy throughput 820 Mbit/s. Samaradorlik necha foiz?
\begin{enumerate}[label=\Alph*),leftmargin=9mm]
\item 82 \%
\item 0,82 \%
\item 18 \%
\item 10,25 \%
\end{enumerate}
\begin{tcolorbox}[enhanced,breakable,colback=teal!4,colframe=teal!50!black,boxrule=0.5pt,arc=1.1mm,left=3mm,right=3mm,top=1.7mm,bottom=1.7mm,before skip=2mm,after skip=2mm]
\textbf{\color{teal}To'g'ri javob:} \textbf{A) 82 \%}
\par\smallskip
\textbf{Izoh:} \ensuremath{\eta} = 820/1000 \ensuremath{\times} 100\% = 82\%.
\end{tcolorbox}
\vspace{2mm}\hrule height 0.2pt \vspace{3mm}

% DB question_id: 4dea5e24-a2fd-5028-8f1e-e91e6bed58d3
\needspace{5\baselineskip}
\noindent\textbf{\color{ink}503.}
\par
2 097 152 B nechta MiB ga teng? Savolda yozilgan SI/IEC birliklariga aynan amal qiling.
\begin{enumerate}[label=\Alph*),leftmargin=9mm]
\item 16 MiB
\item 0,002 MiB
\item 2 048 MiB
\item 2 MiB
\end{enumerate}
\begin{tcolorbox}[enhanced,breakable,colback=teal!4,colframe=teal!50!black,boxrule=0.5pt,arc=1.1mm,left=3mm,right=3mm,top=1.7mm,bottom=1.7mm,before skip=2mm,after skip=2mm]
\textbf{\color{teal}To'g'ri javob:} \textbf{D) 2 MiB}
\par\smallskip
\textbf{Izoh:} Birliklar orasidagi ko'paytirish yoki bo'lish natijasi 2 MiB.
\end{tcolorbox}
\vspace{2mm}\hrule height 0.2pt \vspace{3mm}

% DB question_id: 4edae311-01b1-5af3-85a6-0f5d1cfe14b0
\needspace{5\baselineskip}
\noindent\textbf{\color{ink}504.}
\par
Teng ehtimolli 64 ta holatdan birini ifodalash uchun Xartli formulasi bo'yicha qancha axborot kerak?
\begin{enumerate}[label=\Alph*),leftmargin=9mm]
\item 64 bit
\item 7 bit
\item 6 bit
\item 5 bit
\end{enumerate}
\begin{tcolorbox}[enhanced,breakable,colback=teal!4,colframe=teal!50!black,boxrule=0.5pt,arc=1.1mm,left=3mm,right=3mm,top=1.7mm,bottom=1.7mm,before skip=2mm,after skip=2mm]
\textbf{\color{teal}To'g'ri javob:} \textbf{C) 6 bit}
\par\smallskip
\textbf{Izoh:} I = log\textsubscript{2}64 = 6 bit.
\end{tcolorbox}
\vspace{2mm}\hrule height 0.2pt \vspace{3mm}

% DB question_id: 5099def2-fccd-5b2f-9f45-549b0dfaef73
\needspace{5\baselineskip}
\noindent\textbf{\color{ink}505.}
\par
1 Gbit/s tezlik 30 s davom etsa, qancha axborot uzatiladi?
\begin{enumerate}[label=\Alph*),leftmargin=9mm]
\item 0,47 GB
\item 3,75 GB
\item 30 GB
\item 0,06 GB
\end{enumerate}
\begin{tcolorbox}[enhanced,breakable,colback=teal!4,colframe=teal!50!black,boxrule=0.5pt,arc=1.1mm,left=3mm,right=3mm,top=1.7mm,bottom=1.7mm,before skip=2mm,after skip=2mm]
\textbf{\color{teal}To'g'ri javob:} \textbf{B) 3,75 GB}
\par\smallskip
\textbf{Izoh:} I = v \ensuremath{\times} t; vaqt sekundga va bit/bayt birligi mos holga keltiriladi. Natija: 3,75 GB.
\end{tcolorbox}
\vspace{2mm}\hrule height 0.2pt \vspace{3mm}

% DB question_id: 50af1ee9-3249-5e12-81a9-f08c5daad0a2
\needspace{5\baselineskip}
\noindent\textbf{\color{ink}506.}
\par
Nibble qancha?
\begin{enumerate}[label=\Alph*),leftmargin=9mm]
\item 8 bit
\item 16 bit
\item 1024 bit
\item 4 bit
\end{enumerate}
\begin{tcolorbox}[enhanced,breakable,colback=teal!4,colframe=teal!50!black,boxrule=0.5pt,arc=1.1mm,left=3mm,right=3mm,top=1.7mm,bottom=1.7mm,before skip=2mm,after skip=2mm]
\textbf{\color{teal}To'g'ri javob:} \textbf{D) 4 bit}
\par\smallskip
\textbf{Izoh:} Nibble - yarim bayt, ya'ni 4 bit.
\end{tcolorbox}
\vspace{2mm}\hrule height 0.2pt \vspace{3mm}

% DB question_id: 54e1f273-f9d2-4169-be49-26b5c2060a24
\needspace{5\baselineskip}
\noindent\textbf{\color{ink}507.}
\par
Kodlash nima?
\begin{enumerate}[label=\Alph*),leftmargin=9mm]
\item Axborotni maxfiy kalit yordamida begonalardan yashirish jarayoni
\item Axborot nusxalarini bir xotiradan ikkinchisiga ko'chirish jarayoni
\item Axborot faylining haqiqiyligini elektron imzo bilan tasdiqlash jarayoni
\item Axborotni qoidaga ko'ra boshqa ko'rinishga o'tkazish
\end{enumerate}
\begin{tcolorbox}[enhanced,breakable,colback=teal!4,colframe=teal!50!black,boxrule=0.5pt,arc=1.1mm,left=3mm,right=3mm,top=1.7mm,bottom=1.7mm,before skip=2mm,after skip=2mm]
\textbf{\color{teal}To'g'ri javob:} \textbf{D) Axborotni qoidaga ko'ra boshqa ko'rinishga o'tkazish}
\par\smallskip
\textbf{Izoh:} Kodlash ma'lumotni belgilangan qoidaga ko'ra boshqa tasvirlash shakliga o'tkazishdir. Shifrlash kodlashdan farqli ravishda maxfiylikni ta'minlaydi; nusxalash joylashuvni o'zgartiradi, elektron imzo esa yaxlitlik va mualliflikni tekshirishga xizmat qiladi.
\end{tcolorbox}
\vspace{2mm}\hrule height 0.2pt \vspace{3mm}

% DB question_id: 574f6469-0608-5ddb-9caa-efcdec8b6b47
\needspace{5\baselineskip}
\noindent\textbf{\color{ink}508.}
\par
100 MB fayl siqilgach 40\% asl hajm qoldi. Yangi hajm qancha?
\begin{enumerate}[label=\Alph*),leftmargin=9mm]
\item 60 MB
\item 40 MB
\item 1 600 MB
\item 20 MB
\end{enumerate}
\begin{tcolorbox}[enhanced,breakable,colback=teal!4,colframe=teal!50!black,boxrule=0.5pt,arc=1.1mm,left=3mm,right=3mm,top=1.7mm,bottom=1.7mm,before skip=2mm,after skip=2mm]
\textbf{\color{teal}To'g'ri javob:} \textbf{B) 40 MB}
\par\smallskip
\textbf{Izoh:} 100 \ensuremath{\times} 40/100 = 40 MB.
\end{tcolorbox}
\vspace{2mm}\hrule height 0.2pt \vspace{3mm}

% DB question_id: 5b079077-9859-4e81-9fab-90667aa65f51
\needspace{5\baselineskip}
\noindent\textbf{\color{ink}509.}
\par
Xartli formulasi qachon qo'llanadi?
\begin{enumerate}[label=\Alph*),leftmargin=9mm]
\item Turli ehtimolli holatlarning o'rtacha axborot miqdorini hisoblash uchun
\item Tarmoq orqali uzatilayotgan signaldagi shovqin va yo'qotishni baholash uchun
\item Xatolarni tuzatuvchi kodda ortiqcha nazorat bitlari sonini belgilash uchun
\item Teng ehtimolli N ta holatdan birini tanlash axborotini topishda
\end{enumerate}
\begin{tcolorbox}[enhanced,breakable,colback=teal!4,colframe=teal!50!black,boxrule=0.5pt,arc=1.1mm,left=3mm,right=3mm,top=1.7mm,bottom=1.7mm,before skip=2mm,after skip=2mm]
\textbf{\color{teal}To'g'ri javob:} \textbf{D) Teng ehtimolli N ta holatdan birini tanlash axborotini topishda}
\par\smallskip
\textbf{Izoh:} Xartli formulasida I = log\textsubscript{2}N va N ta mumkin bo'lgan holat teng ehtimolli deb qaraladi. Ehtimollari turlicha bo'lgan holatlarning o'rtacha axboroti Shennon entropiyasi bilan tavsiflanadi. Shovqin yoki nazorat bitlarini hisoblash boshqa masalalardir.
\end{tcolorbox}
\vspace{2mm}\hrule height 0.2pt \vspace{3mm}

% DB question_id: 5cfb7923-5ab1-524a-848b-b738af90c35b
\needspace{5\baselineskip}
\noindent\textbf{\color{ink}510.}
\par
Asl hajm 450 MB, siqilgan hajm 180 MB. Siqish koeffitsiyenti qancha?
\begin{enumerate}[label=\Alph*),leftmargin=9mm]
\item 447,5 marta
\item 450 marta
\item 2,5 marta
\item 1,25 marta
\end{enumerate}
\begin{tcolorbox}[enhanced,breakable,colback=teal!4,colframe=teal!50!black,boxrule=0.5pt,arc=1.1mm,left=3mm,right=3mm,top=1.7mm,bottom=1.7mm,before skip=2mm,after skip=2mm]
\textbf{\color{teal}To'g'ri javob:} \textbf{C) 2,5 marta}
\par\smallskip
\textbf{Izoh:} 450/180 = 2,5 marta.
\end{tcolorbox}
\vspace{2mm}\hrule height 0.2pt \vspace{3mm}

% DB question_id: 5d75ecb4-9bff-57cf-a7bd-0f11bd3b52cb
\needspace{5\baselineskip}
\noindent\textbf{\color{ink}511.}
\par
Vijiner shifrida ingliz alifbosi \texttt{A=0,...,Z=25} tartibida olinadi va kalit takrorlanadi. \texttt{BILIMDON} matnini \texttt{PYTHON} kaliti bilan shifrlang.
\begin{enumerate}[label=\Alph*),leftmargin=9mm]
\item QGEPAQDL
\item RHFQBREL
\item QGEOAPDK
\item PYTHONPY
\end{enumerate}
\begin{tcolorbox}[enhanced,breakable,colback=teal!4,colframe=teal!50!black,boxrule=0.5pt,arc=1.1mm,left=3mm,right=3mm,top=1.7mm,bottom=1.7mm,before skip=2mm,after skip=2mm]
\textbf{\color{teal}To'g'ri javob:} \textbf{A) QGEPAQDL}
\par\smallskip
\textbf{Izoh:} Takrorlangan kalit PYTHONPY. Modullar yig'indisi 26 bo'yicha olinsa QGEPAQDL chiqadi.
\end{tcolorbox}
\vspace{2mm}\hrule height 0.2pt \vspace{3mm}

% DB question_id: 5f02a7dc-f8b0-5a56-9999-6825cb82f865
\needspace{5\baselineskip}
\noindent\textbf{\color{ink}512.}
\par
1024 ta belgi, har biri 2 baytdan kodlangan. Xabar hajmi qancha?
\begin{enumerate}[label=\Alph*),leftmargin=9mm]
\item 1 024 KiB
\item 16 384 KiB
\item 256 KiB
\item 2 KiB
\end{enumerate}
\begin{tcolorbox}[enhanced,breakable,colback=teal!4,colframe=teal!50!black,boxrule=0.5pt,arc=1.1mm,left=3mm,right=3mm,top=1.7mm,bottom=1.7mm,before skip=2mm,after skip=2mm]
\textbf{\color{teal}To'g'ri javob:} \textbf{D) 2 KiB}
\par\smallskip
\textbf{Izoh:} Hajm = 1024 \ensuremath{\times} 2 = 2048 B = 2 KiB.
\end{tcolorbox}
\vspace{2mm}\hrule height 0.2pt \vspace{3mm}

% DB question_id: 6285684e-b001-5ad7-8016-d36e43b94921
\needspace{5\baselineskip}
\noindent\textbf{\color{ink}513.}
\par
1200 ta belgi, har biri 1 baytdan kodlangan. Xabar hajmi qancha?
\begin{enumerate}[label=\Alph*),leftmargin=9mm]
\item 9 600 B
\item 150 B
\item 1 200 B
\item 1 203 B
\end{enumerate}
\begin{tcolorbox}[enhanced,breakable,colback=teal!4,colframe=teal!50!black,boxrule=0.5pt,arc=1.1mm,left=3mm,right=3mm,top=1.7mm,bottom=1.7mm,before skip=2mm,after skip=2mm]
\textbf{\color{teal}To'g'ri javob:} \textbf{C) 1 200 B}
\par\smallskip
\textbf{Izoh:} Hajm = 1200 \ensuremath{\times} 1 = 1200 B.
\end{tcolorbox}
\vspace{2mm}\hrule height 0.2pt \vspace{3mm}

% DB question_id: 649b861f-2246-46bf-a0d3-149c28c2d66a
\needspace{5\baselineskip}
\noindent\textbf{\color{ink}514.}
\par
Har bir harfga ma'lum bitlar ketma-ketligi biriktirildi. Bu vaziyatdagi asosiy tushunchani aniqlang.
\begin{enumerate}[label=\Alph*),leftmargin=9mm]
\item Kodlash
\item Dekodlash
\item Kod
\item Doimiy uzunlikdagi kod
\end{enumerate}
\begin{tcolorbox}[enhanced,breakable,colback=teal!4,colframe=teal!50!black,boxrule=0.5pt,arc=1.1mm,left=3mm,right=3mm,top=1.7mm,bottom=1.7mm,before skip=2mm,after skip=2mm]
\textbf{\color{teal}To'g'ri javob:} \textbf{C) Kod}
\par\smallskip
\textbf{Izoh:} Vaziyat Kodga misol: axborot elementlarini boshqa belgilar yoki signallar bilan moslashtirish qoidalari tizimi.
\end{tcolorbox}
\vspace{2mm}\hrule height 0.2pt \vspace{3mm}

% DB question_id: 673d12f6-1181-4515-bee9-9399bdf5472d
\needspace{5\baselineskip}
\noindent\textbf{\color{ink}515.}
\par
MP4 ichidagi H.264 video dekodlandi. Bu vaziyatdagi asosiy tushunchani aniqlang.
\begin{enumerate}[label=\Alph*),leftmargin=9mm]
\item Validatsiya va verifikatsiya farqi
\item Internet va WWW farqi
\item Kodek va konteyner farqi
\item kB va KiB farqi
\end{enumerate}
\begin{tcolorbox}[enhanced,breakable,colback=teal!4,colframe=teal!50!black,boxrule=0.5pt,arc=1.1mm,left=3mm,right=3mm,top=1.7mm,bottom=1.7mm,before skip=2mm,after skip=2mm]
\textbf{\color{teal}To'g'ri javob:} \textbf{C) Kodek va konteyner farqi}
\par\smallskip
\textbf{Izoh:} Vaziyat Kodek va konteyner farqiga misol: kodek oqimni kodlashi, konteyner esa oqimlarni bir faylda saqlashini ajratish.
\end{tcolorbox}
\vspace{2mm}\hrule height 0.2pt \vspace{3mm}

% DB question_id: 696b3ce6-d7e6-417b-a2b8-7600211a23e6
\needspace{5\baselineskip}
\noindent\textbf{\color{ink}516.}
\par
Pasportda 2 ta lotin bosh harfi (26 imkoniyat) va 7 ta o'nlik raqam bor. Har bir harf va har bir raqam alohida bir xil uzunlikdagi minimal ikkilik kodda saqlansa, jami nechta bit kerak?
\begin{enumerate}[label=\Alph*),leftmargin=9mm]
\item 32 bit
\item 38 bit
\item 40 bit
\item 76 bit
\end{enumerate}
\begin{tcolorbox}[enhanced,breakable,colback=teal!4,colframe=teal!50!black,boxrule=0.5pt,arc=1.1mm,left=3mm,right=3mm,top=1.7mm,bottom=1.7mm,before skip=2mm,after skip=2mm]
\textbf{\color{teal}To'g'ri javob:} \textbf{B) 38 bit}
\par\smallskip
\textbf{Izoh:} 26 harf uchun ceil(log\textsubscript{2}26)=5 bit, 0...9 raqamlari uchun ceil(log\textsubscript{2}10)=4 bit. Ikki harf 10 bit, yetti raqam 28 bit, jami 38 bit.
\end{tcolorbox}
\vspace{2mm}\hrule height 0.2pt \vspace{3mm}

% DB question_id: 6d2bd337-c409-4e4e-a1da-714c53892db0
\needspace{5\baselineskip}
\noindent\textbf{\color{ink}517.}
\par
Kodlashda har bir harf yoki xabar oldindan kelishilgan qoida bilan boshqa belgilar ketma-ketligida ifodalanadi. Aynan kod atamasining ta'rifi qaysi?
\begin{enumerate}[label=\Alph*),leftmargin=9mm]
\item Belgilar uchun belgilangan barcha kodlarni jamlovchi jadval
\item Belgiga yoki xabarga mos qoida asosidagi ifoda
\item Xabarni boshqa qurilmaga uzatish uchun tanlangan kanal
\item Kodlar ichidagi takroriy belgilarni qisqartirishga xizmat qiladigan usul
\end{enumerate}
\begin{tcolorbox}[enhanced,breakable,colback=teal!4,colframe=teal!50!black,boxrule=0.5pt,arc=1.1mm,left=3mm,right=3mm,top=1.7mm,bottom=1.7mm,before skip=2mm,after skip=2mm]
\textbf{\color{teal}To'g'ri javob:} \textbf{B) Belgiga yoki xabarga mos qoida asosidagi ifoda}
\par\smallskip
\textbf{Izoh:} Kod - belgiga yoki xabarga qoida asosida moslangan ifoda. Kod jadvali ko'plab belgilar va kodlar mosligini ko'rsatadi; uzatish kanali yoki siqish algoritmi boshqa tushunchadir.
\end{tcolorbox}
\vspace{2mm}\hrule height 0.2pt \vspace{3mm}

% DB question_id: 6dfc0e05-7343-4dd2-ba58-57b6aa1df8ee
\needspace{5\baselineskip}
\noindent\textbf{\color{ink}518.}
\par
Har lampochka qizil yoki yashil holatdan bittasini oladi. Qurilma kamida 200 turli signal kombinatsiyasi berishi uchun eng kamida nechta mustaqil lampochka kerak?
\begin{enumerate}[label=\Alph*),leftmargin=9mm]
\item 7
\item 8
\item 9
\item 10
\end{enumerate}
\begin{tcolorbox}[enhanced,breakable,colback=teal!4,colframe=teal!50!black,boxrule=0.5pt,arc=1.1mm,left=3mm,right=3mm,top=1.7mm,bottom=1.7mm,before skip=2mm,after skip=2mm]
\textbf{\color{teal}To'g'ri javob:} \textbf{B) 8}
\par\smallskip
\textbf{Izoh:} n lampochkaning holatlari 2\textsuperscript{n} ta. 2\textsuperscript{7}=128 yetmaydi, 2\textsuperscript{8}=256 esa kamida 200 ta turli holatni beradi. Minimum n=8.
\end{tcolorbox}
\vspace{2mm}\hrule height 0.2pt \vspace{3mm}

% DB question_id: 72d810bc-a764-5557-89ec-febca6351962
\needspace{5\baselineskip}
\noindent\textbf{\color{ink}519.}
\par
240 MB fayl 4:1 siqish koeffitsiyenti bilan siqildi. Yangi hajm qancha?
\begin{enumerate}[label=\Alph*),leftmargin=9mm]
\item 180 MB
\item 240 MB
\item 60 MB
\item 30 MB
\end{enumerate}
\begin{tcolorbox}[enhanced,breakable,colback=teal!4,colframe=teal!50!black,boxrule=0.5pt,arc=1.1mm,left=3mm,right=3mm,top=1.7mm,bottom=1.7mm,before skip=2mm,after skip=2mm]
\textbf{\color{teal}To'g'ri javob:} \textbf{C) 60 MB}
\par\smallskip
\textbf{Izoh:} 240/4 = 60 MB.
\end{tcolorbox}
\vspace{2mm}\hrule height 0.2pt \vspace{3mm}

% DB question_id: 73af6fd6-4464-5771-b77f-fdd5c236f8b1
\needspace{5\baselineskip}
\noindent\textbf{\color{ink}520.}
\par
2,5 MB nechta Mbit ga teng? Savolda yozilgan SI/IEC birliklariga aynan amal qiling.
\begin{enumerate}[label=\Alph*),leftmargin=9mm]
\item 2,5 Mbit
\item 160 Mbit
\item 20 Mbit
\item 20 480 Mbit
\end{enumerate}
\begin{tcolorbox}[enhanced,breakable,colback=teal!4,colframe=teal!50!black,boxrule=0.5pt,arc=1.1mm,left=3mm,right=3mm,top=1.7mm,bottom=1.7mm,before skip=2mm,after skip=2mm]
\textbf{\color{teal}To'g'ri javob:} \textbf{C) 20 Mbit}
\par\smallskip
\textbf{Izoh:} Birliklar orasidagi ko'paytirish yoki bo'lish natijasi 20 Mbit.
\end{tcolorbox}
\vspace{2mm}\hrule height 0.2pt \vspace{3mm}

% DB question_id: 759b631a-294a-59b3-95df-1fbdea5134f6
\needspace{5\baselineskip}
\noindent\textbf{\color{ink}521.}
\par
600 MiB hajmli fayl 40 Mbit/s tezlikda uzatilsa, vaqt qancha?
\begin{enumerate}[label=\Alph*),leftmargin=9mm]
\item 15,73 s
\item 1006,64 s
\item 125,83 s
\item 157,29 s
\end{enumerate}
\begin{tcolorbox}[enhanced,breakable,colback=teal!4,colframe=teal!50!black,boxrule=0.5pt,arc=1.1mm,left=3mm,right=3mm,top=1.7mm,bottom=1.7mm,before skip=2mm,after skip=2mm]
\textbf{\color{teal}To'g'ri javob:} \textbf{C) 125,83 s}
\par\smallskip
\textbf{Izoh:} Hajm bitga o'tkazilib, t = I/v formula bilan hisoblanadi. Natija: 125.83 s.
\end{tcolorbox}
\vspace{2mm}\hrule height 0.2pt \vspace{3mm}

% DB question_id: 75ae48c6-1315-5457-af78-61209df0a335
\needspace{5\baselineskip}
\noindent\textbf{\color{ink}522.}
\par
Teng ehtimolli 8 ta holatdan birini ifodalash uchun Xartli formulasi bo'yicha qancha axborot kerak?
\begin{enumerate}[label=\Alph*),leftmargin=9mm]
\item 3 bit
\item 8 bit
\item 4 bit
\item 2 bit
\end{enumerate}
\begin{tcolorbox}[enhanced,breakable,colback=teal!4,colframe=teal!50!black,boxrule=0.5pt,arc=1.1mm,left=3mm,right=3mm,top=1.7mm,bottom=1.7mm,before skip=2mm,after skip=2mm]
\textbf{\color{teal}To'g'ri javob:} \textbf{A) 3 bit}
\par\smallskip
\textbf{Izoh:} I = log\textsubscript{2}8 = 3 bit.
\end{tcolorbox}
\vspace{2mm}\hrule height 0.2pt \vspace{3mm}

% DB question_id: 76a3ddb8-f837-5900-a348-997fa53e24f7
\needspace{5\baselineskip}
\noindent\textbf{\color{ink}523.}
\par
2 048 B nechta KiB ga teng? Savolda yozilgan SI/IEC birliklariga aynan amal qiling.
\begin{enumerate}[label=\Alph*),leftmargin=9mm]
\item 16 KiB
\item 0,002 KiB
\item 2 048 KiB
\item 2 KiB
\end{enumerate}
\begin{tcolorbox}[enhanced,breakable,colback=teal!4,colframe=teal!50!black,boxrule=0.5pt,arc=1.1mm,left=3mm,right=3mm,top=1.7mm,bottom=1.7mm,before skip=2mm,after skip=2mm]
\textbf{\color{teal}To'g'ri javob:} \textbf{D) 2 KiB}
\par\smallskip
\textbf{Izoh:} Birliklar orasidagi ko'paytirish yoki bo'lish natijasi 2 KiB.
\end{tcolorbox}
\vspace{2mm}\hrule height 0.2pt \vspace{3mm}

% DB question_id: 76c8491e-e378-5ee1-98b8-5fb604801a23
\needspace{5\baselineskip}
\noindent\textbf{\color{ink}524.}
\par
32 GB nechta Gbit ga teng? Savolda yozilgan SI/IEC birliklariga aynan amal qiling.
\begin{enumerate}[label=\Alph*),leftmargin=9mm]
\item 32 Gbit
\item 2 048 Gbit
\item 262 144 Gbit
\item 256 Gbit
\end{enumerate}
\begin{tcolorbox}[enhanced,breakable,colback=teal!4,colframe=teal!50!black,boxrule=0.5pt,arc=1.1mm,left=3mm,right=3mm,top=1.7mm,bottom=1.7mm,before skip=2mm,after skip=2mm]
\textbf{\color{teal}To'g'ri javob:} \textbf{D) 256 Gbit}
\par\smallskip
\textbf{Izoh:} Birliklar orasidagi ko'paytirish yoki bo'lish natijasi 256 Gbit.
\end{tcolorbox}
\vspace{2mm}\hrule height 0.2pt \vspace{3mm}

% DB question_id: 7790ef49-8773-5e1d-bc94-07a05c95af2d
\needspace{5\baselineskip}
\noindent\textbf{\color{ink}525.}
\par
33 ta turli belgini doimiy uzunlikdagi ikkilik kod bilan kodlash uchun minimal kod uzunligi qancha?
\begin{enumerate}[label=\Alph*),leftmargin=9mm]
\item 5 bit
\item 6 bit
\item 33 bit
\item 7 bit
\end{enumerate}
\begin{tcolorbox}[enhanced,breakable,colback=teal!4,colframe=teal!50!black,boxrule=0.5pt,arc=1.1mm,left=3mm,right=3mm,top=1.7mm,bottom=1.7mm,before skip=2mm,after skip=2mm]
\textbf{\color{teal}To'g'ri javob:} \textbf{B) 6 bit}
\par\smallskip
\textbf{Izoh:} 2\textsuperscript{n} \ensuremath{\geq} 33; minimal n = 6.
\end{tcolorbox}
\vspace{2mm}\hrule height 0.2pt \vspace{3mm}

% DB question_id: 7990a98e-7970-4f2c-9a04-9e5fcfeea89e
\needspace{5\baselineskip}
\noindent\textbf{\color{ink}526.}
\par
Tarmoq qurilmasi ishlab chiqaruvchisi hajmni o'nlik prefiksda ko'rsatdi. Vaziyatni tahlil qilib, unga eng aniq mos tushunchani toping.
\begin{enumerate}[label=\Alph*),leftmargin=9mm]
\item KiB
\item kB
\item Xotira sig'imi
\item Bit
\end{enumerate}
\begin{tcolorbox}[enhanced,breakable,colback=teal!4,colframe=teal!50!black,boxrule=0.5pt,arc=1.1mm,left=3mm,right=3mm,top=1.7mm,bottom=1.7mm,before skip=2mm,after skip=2mm]
\textbf{\color{teal}To'g'ri javob:} \textbf{B) kB}
\par\smallskip
\textbf{Izoh:} Vaziyat kBga misol: SI standartida 1000 baytga teng o'nlik birlik.
\end{tcolorbox}
\vspace{2mm}\hrule height 0.2pt \vspace{3mm}

% DB question_id: 79f55735-9685-514f-8fd5-3ccadbe73289
\needspace{5\baselineskip}
\noindent\textbf{\color{ink}527.}
\par
20 Mbit/s tezlik 90 s davom etsa, qancha axborot uzatiladi?
\begin{enumerate}[label=\Alph*),leftmargin=9mm]
\item 28,12 MB
\item 1 800 MB
\item 3,75 MB
\item 225 MB
\end{enumerate}
\begin{tcolorbox}[enhanced,breakable,colback=teal!4,colframe=teal!50!black,boxrule=0.5pt,arc=1.1mm,left=3mm,right=3mm,top=1.7mm,bottom=1.7mm,before skip=2mm,after skip=2mm]
\textbf{\color{teal}To'g'ri javob:} \textbf{D) 225 MB}
\par\smallskip
\textbf{Izoh:} I = v \ensuremath{\times} t; vaqt sekundga va bit/bayt birligi mos holga keltiriladi. Natija: 225 MB.
\end{tcolorbox}
\vspace{2mm}\hrule height 0.2pt \vspace{3mm}

% DB question_id: 7a7482b7-86a3-410b-86d1-48ad9fcc0e39
\needspace{5\baselineskip}
\noindent\textbf{\color{ink}528.}
\par
Qabul qilingan bitlardan ekranda asl matn hosil qilindi. Bu vaziyatdagi asosiy tushunchani aniqlang.
\begin{enumerate}[label=\Alph*),leftmargin=9mm]
\item Dekodlash
\item Doimiy uzunlikdagi kod
\item O'zgaruvchan uzunlikdagi kod
\item Siqish
\end{enumerate}
\begin{tcolorbox}[enhanced,breakable,colback=teal!4,colframe=teal!50!black,boxrule=0.5pt,arc=1.1mm,left=3mm,right=3mm,top=1.7mm,bottom=1.7mm,before skip=2mm,after skip=2mm]
\textbf{\color{teal}To'g'ri javob:} \textbf{A) Dekodlash}
\par\smallskip
\textbf{Izoh:} Vaziyat Dekodlashga misol: kodlangan axborotni kelishilgan qoida asosida qayta tiklash.
\end{tcolorbox}
\vspace{2mm}\hrule height 0.2pt \vspace{3mm}

% DB question_id: 7d192ceb-eccc-5f3c-a8f2-02b4054117a3
\needspace{5\baselineskip}
\noindent\textbf{\color{ink}529.}
\par
Kitob 200 sahifa, har sahifada 40 satr, har satrda 75 belgi. Har belgi 2 bayt. Matn hajmini MiB da toping.
\begin{enumerate}[label=\Alph*),leftmargin=9mm]
\item 600 000 MiB
\item 9 600 000 MiB
\item 1,144 MiB
\item 1 200 MiB
\end{enumerate}
\begin{tcolorbox}[enhanced,breakable,colback=teal!4,colframe=teal!50!black,boxrule=0.5pt,arc=1.1mm,left=3mm,right=3mm,top=1.7mm,bottom=1.7mm,before skip=2mm,after skip=2mm]
\textbf{\color{teal}To'g'ri javob:} \textbf{C) 1,144 MiB}
\par\smallskip
\textbf{Izoh:} 200 \ensuremath{\times} 40 \ensuremath{\times} 75 \ensuremath{\times} 2 = 1 200 000 B = 1,144 MiB.
\end{tcolorbox}
\vspace{2mm}\hrule height 0.2pt \vspace{3mm}

% DB question_id: 7e7eaca1-2008-525f-aa7a-f8bb91745ddb
\needspace{5\baselineskip}
\noindent\textbf{\color{ink}530.}
\par
3 GiB hajmli fayl 24 MiB/s tezlikda uzatilsa, vaqt qancha?
\begin{enumerate}[label=\Alph*),leftmargin=9mm]
\item 128 s
\item 16 s
\item 1 024 s
\item 160 s
\end{enumerate}
\begin{tcolorbox}[enhanced,breakable,colback=teal!4,colframe=teal!50!black,boxrule=0.5pt,arc=1.1mm,left=3mm,right=3mm,top=1.7mm,bottom=1.7mm,before skip=2mm,after skip=2mm]
\textbf{\color{teal}To'g'ri javob:} \textbf{A) 128 s}
\par\smallskip
\textbf{Izoh:} Hajm bitga o'tkazilib, t = I/v formula bilan hisoblanadi. Natija: 128 s.
\end{tcolorbox}
\vspace{2mm}\hrule height 0.2pt \vspace{3mm}

% DB question_id: 8264cb99-c8ff-59a1-a0a0-fea0b89dfa07
\needspace{5\baselineskip}
\noindent\textbf{\color{ink}531.}
\par
Teng ehtimolli 4 ta holatdan birini ifodalash uchun Xartli formulasi bo'yicha qancha axborot kerak?
\begin{enumerate}[label=\Alph*),leftmargin=9mm]
\item 4 bit
\item 3 bit
\item 1 bit
\item 2 bit
\end{enumerate}
\begin{tcolorbox}[enhanced,breakable,colback=teal!4,colframe=teal!50!black,boxrule=0.5pt,arc=1.1mm,left=3mm,right=3mm,top=1.7mm,bottom=1.7mm,before skip=2mm,after skip=2mm]
\textbf{\color{teal}To'g'ri javob:} \textbf{D) 2 bit}
\par\smallskip
\textbf{Izoh:} I = log\textsubscript{2}4 = 2 bit.
\end{tcolorbox}
\vspace{2mm}\hrule height 0.2pt \vspace{3mm}

% DB question_id: 86bfd308-3aa4-5e73-bf79-d9c3cbc424b4
\needspace{5\baselineskip}
\noindent\textbf{\color{ink}532.}
\par
4 MB nechta B ga teng? Savolda yozilgan SI/IEC birliklariga aynan amal qiling.
\begin{enumerate}[label=\Alph*),leftmargin=9mm]
\item 500 000 B
\item 32 000 000 B
\item 3906,25 B
\item 4 000 000 B
\end{enumerate}
\begin{tcolorbox}[enhanced,breakable,colback=teal!4,colframe=teal!50!black,boxrule=0.5pt,arc=1.1mm,left=3mm,right=3mm,top=1.7mm,bottom=1.7mm,before skip=2mm,after skip=2mm]
\textbf{\color{teal}To'g'ri javob:} \textbf{D) 4 000 000 B}
\par\smallskip
\textbf{Izoh:} Birliklar orasidagi ko'paytirish yoki bo'lish natijasi 4 000 000 B.
\end{tcolorbox}
\vspace{2mm}\hrule height 0.2pt \vspace{3mm}

% DB question_id: 8b6796c3-549c-5fe2-9f59-5f7ca24fb1b5
\needspace{5\baselineskip}
\noindent\textbf{\color{ink}533.}
\par
2 GB axborot 100 s da uzatildi. O'rtacha tezlikni toping.
\begin{enumerate}[label=\Alph*),leftmargin=9mm]
\item 20 Mbit/s
\item 1 280 Mbit/s
\item 9 600 Mbit/s
\item 160 Mbit/s
\end{enumerate}
\begin{tcolorbox}[enhanced,breakable,colback=teal!4,colframe=teal!50!black,boxrule=0.5pt,arc=1.1mm,left=3mm,right=3mm,top=1.7mm,bottom=1.7mm,before skip=2mm,after skip=2mm]
\textbf{\color{teal}To'g'ri javob:} \textbf{D) 160 Mbit/s}
\par\smallskip
\textbf{Izoh:} Vaqt sekundga o'tkazilib, v = I/t formula bilan hisoblanadi. Natija: 160 Mbit/s.
\end{tcolorbox}
\vspace{2mm}\hrule height 0.2pt \vspace{3mm}

% DB question_id: 914e4f90-5468-4297-be24-8af29c724050
\needspace{5\baselineskip}
\noindent\textbf{\color{ink}534.}
\par
Kitobda 128 sahifa, har sahifada 32 qator, har qatorda 60 belgi bor. Har belgi 8 bit, xizmat metama'lumoti yo'q. 512 bit/sekund tezlikda uzatish necha daqiqa davom etadi?
\begin{enumerate}[label=\Alph*),leftmargin=9mm]
\item 60 daqiqa
\item 62 daqiqa
\item 64 daqiqa
\item 66 daqiqa
\end{enumerate}
\begin{tcolorbox}[enhanced,breakable,colback=teal!4,colframe=teal!50!black,boxrule=0.5pt,arc=1.1mm,left=3mm,right=3mm,top=1.7mm,bottom=1.7mm,before skip=2mm,after skip=2mm]
\textbf{\color{teal}To'g'ri javob:} \textbf{C) 64 daqiqa}
\par\smallskip
\textbf{Izoh:} Belgi soni 128\ensuremath{\times}32\ensuremath{\times}60=245760, bitlar soni 245760\ensuremath{\times}8=1966080. 512 bit/s bilan vaqt 3840 sekund; 3840/60=64 daqiqa. Asl PDFdagi 3600s va 60min bir xil noto'g'ri variantlar bo'lib, yangi versiyada takrorlar yo'q.
\end{tcolorbox}
\vspace{2mm}\hrule height 0.2pt \vspace{3mm}

% DB question_id: 968a2623-a879-5865-b09e-9b799b26bc69
\needspace{5\baselineskip}
\noindent\textbf{\color{ink}535.}
\par
Yo'qotishsiz siqish xususiyati qaysi?
\begin{enumerate}[label=\Alph*),leftmargin=9mm]
\item Ayrim ma'lumot ataylab tashlanadi
\item Asl ma'lumot to'liq tiklanadi
\item Fayl doimo kattalashadi
\item Kodlash bajarilmaydi
\end{enumerate}
\begin{tcolorbox}[enhanced,breakable,colback=teal!4,colframe=teal!50!black,boxrule=0.5pt,arc=1.1mm,left=3mm,right=3mm,top=1.7mm,bottom=1.7mm,before skip=2mm,after skip=2mm]
\textbf{\color{teal}To'g'ri javob:} \textbf{B) Asl ma'lumot to'liq tiklanadi}
\par\smallskip
\textbf{Izoh:} Lossless usulda bitlar to'liq tiklanadi.
\end{tcolorbox}
\vspace{2mm}\hrule height 0.2pt \vspace{3mm}

% DB question_id: 9745fc4a-4e64-5240-8470-0c4a00befd43
\needspace{5\baselineskip}
\noindent\textbf{\color{ink}536.}
\par
5 KiB nechta B ga teng? Savolda yozilgan SI/IEC birliklariga aynan amal qiling.
\begin{enumerate}[label=\Alph*),leftmargin=9mm]
\item 5 120 B
\item 640 B
\item 40 960 B
\item 5 B
\end{enumerate}
\begin{tcolorbox}[enhanced,breakable,colback=teal!4,colframe=teal!50!black,boxrule=0.5pt,arc=1.1mm,left=3mm,right=3mm,top=1.7mm,bottom=1.7mm,before skip=2mm,after skip=2mm]
\textbf{\color{teal}To'g'ri javob:} \textbf{A) 5 120 B}
\par\smallskip
\textbf{Izoh:} Birliklar orasidagi ko'paytirish yoki bo'lish natijasi 5 120 B.
\end{tcolorbox}
\vspace{2mm}\hrule height 0.2pt \vspace{3mm}

% DB question_id: 98fc692b-3380-5d37-b066-e7cca9eefe88
\needspace{5\baselineskip}
\noindent\textbf{\color{ink}537.}
\par
3 MiB nechta KiB ga teng? Savolda yozilgan SI/IEC birliklariga aynan amal qiling.
\begin{enumerate}[label=\Alph*),leftmargin=9mm]
\item 384 KiB
\item 3 072 KiB
\item 24 576 KiB
\item 3 KiB
\end{enumerate}
\begin{tcolorbox}[enhanced,breakable,colback=teal!4,colframe=teal!50!black,boxrule=0.5pt,arc=1.1mm,left=3mm,right=3mm,top=1.7mm,bottom=1.7mm,before skip=2mm,after skip=2mm]
\textbf{\color{teal}To'g'ri javob:} \textbf{B) 3 072 KiB}
\par\smallskip
\textbf{Izoh:} Birliklar orasidagi ko'paytirish yoki bo'lish natijasi 3 072 KiB.
\end{tcolorbox}
\vspace{2mm}\hrule height 0.2pt \vspace{3mm}

% DB question_id: 9bde5319-2e7a-5b12-bf23-689a970c684f
\needspace{5\baselineskip}
\noindent\textbf{\color{ink}538.}
\par
16 777 216 bit nechta MiB ga teng? Savolda yozilgan SI/IEC birliklariga aynan amal qiling.
\begin{enumerate}[label=\Alph*),leftmargin=9mm]
\item 2 MiB
\item 16 MiB
\item 0 MiB
\item 2 048 MiB
\end{enumerate}
\begin{tcolorbox}[enhanced,breakable,colback=teal!4,colframe=teal!50!black,boxrule=0.5pt,arc=1.1mm,left=3mm,right=3mm,top=1.7mm,bottom=1.7mm,before skip=2mm,after skip=2mm]
\textbf{\color{teal}To'g'ri javob:} \textbf{A) 2 MiB}
\par\smallskip
\textbf{Izoh:} Birliklar orasidagi ko'paytirish yoki bo'lish natijasi 2 MiB.
\end{tcolorbox}
\vspace{2mm}\hrule height 0.2pt \vspace{3mm}

% DB question_id: a2710e86-0d70-5874-89bd-769a2daf8aa3
\needspace{5\baselineskip}
\noindent\textbf{\color{ink}539.}
\par
Aloqa kanalida A, B, C, D, E, F, G va H harflari uzatiladi. Fano (prefiks) shartini qanoatlantiruvchi ikkilik kodda \texttt{A=000}, \texttt{B=001}, \texttt{C=0101}, \texttt{D=0100}, \texttt{E=011}, \texttt{F=101}. G va H uchun kod so'zlarining mumkin bo'lgan eng kichik umumiy uzunligini toping.
\begin{enumerate}[label=\Alph*),leftmargin=9mm]
\item 6
\item 4
\item 5
\item 3
\end{enumerate}
\begin{tcolorbox}[enhanced,breakable,colback=teal!4,colframe=teal!50!black,boxrule=0.5pt,arc=1.1mm,left=3mm,right=3mm,top=1.7mm,bottom=1.7mm,before skip=2mm,after skip=2mm]
\textbf{\color{teal}To'g'ri javob:} \textbf{C) 5}
\par\smallskip
\textbf{Izoh:} 0 bilan boshlanuvchi mavjud shoxlar berilgan kodlar bilan band. 1 shoxida \texttt{101} mavjud; qolgan ikki harfga, masalan, \texttt{11} va \texttt{100} kodlarini berish mumkin. Hech biri boshqasining prefiksi emas. Umumiy uzunlik 2+3=5; bundan qisqa ikki kodni joylashtirib bo'lmaydi.
\end{tcolorbox}
\vspace{2mm}\hrule height 0.2pt \vspace{3mm}

% DB question_id: a44dac7c-23a5-5c22-8489-4342f05c00fb
\needspace{5\baselineskip}
\noindent\textbf{\color{ink}540.}
\par
5 000 000 B nechta MB ga teng? Savolda yozilgan SI/IEC birliklariga aynan amal qiling.
\begin{enumerate}[label=\Alph*),leftmargin=9mm]
\item 5 MB
\item 0,625 MB
\item 40 MB
\item 5 120 MB
\end{enumerate}
\begin{tcolorbox}[enhanced,breakable,colback=teal!4,colframe=teal!50!black,boxrule=0.5pt,arc=1.1mm,left=3mm,right=3mm,top=1.7mm,bottom=1.7mm,before skip=2mm,after skip=2mm]
\textbf{\color{teal}To'g'ri javob:} \textbf{A) 5 MB}
\par\smallskip
\textbf{Izoh:} Birliklar orasidagi ko'paytirish yoki bo'lish natijasi 5 MB.
\end{tcolorbox}
\vspace{2mm}\hrule height 0.2pt \vspace{3mm}

% DB question_id: a82a3d54-973f-55b2-bfbe-2719d3528703
\needspace{5\baselineskip}
\noindent\textbf{\color{ink}541.}
\par
2.4 GB fayl hajmi 25\% ga kamaydi. Siqilgan hajm qancha?
\begin{enumerate}[label=\Alph*),leftmargin=9mm]
\item 0,6 GB
\item 45 GB
\item 1,8 GB
\item 75 GB
\end{enumerate}
\begin{tcolorbox}[enhanced,breakable,colback=teal!4,colframe=teal!50!black,boxrule=0.5pt,arc=1.1mm,left=3mm,right=3mm,top=1.7mm,bottom=1.7mm,before skip=2mm,after skip=2mm]
\textbf{\color{teal}To'g'ri javob:} \textbf{C) 1,8 GB}
\par\smallskip
\textbf{Izoh:} 2.4 \ensuremath{\times} (1-25/100) = 1,8 GB.
\end{tcolorbox}
\vspace{2mm}\hrule height 0.2pt \vspace{3mm}

% DB question_id: a8aef085-3e89-5b30-ad7d-af767a052c63
\needspace{5\baselineskip}
\noindent\textbf{\color{ink}542.}
\par
Kanalning nazariy tezligi 100 Mbit/s, amaliy throughput 72 Mbit/s. Samaradorlik necha foiz?
\begin{enumerate}[label=\Alph*),leftmargin=9mm]
\item 0,72 \%
\item 28 \%
\item 72 \%
\item 9 \%
\end{enumerate}
\begin{tcolorbox}[enhanced,breakable,colback=teal!4,colframe=teal!50!black,boxrule=0.5pt,arc=1.1mm,left=3mm,right=3mm,top=1.7mm,bottom=1.7mm,before skip=2mm,after skip=2mm]
\textbf{\color{teal}To'g'ri javob:} \textbf{C) 72 \%}
\par\smallskip
\textbf{Izoh:} \ensuremath{\eta} = 72/100 \ensuremath{\times} 100\% = 72\%.
\end{tcolorbox}
\vspace{2mm}\hrule height 0.2pt \vspace{3mm}

% DB question_id: ad117b10-0654-5369-b8fd-4804277eb0ac
\needspace{5\baselineskip}
\noindent\textbf{\color{ink}543.}
\par
96 MB fayl 3:1 siqish koeffitsiyenti bilan siqildi. Yangi hajm qancha?
\begin{enumerate}[label=\Alph*),leftmargin=9mm]
\item 64 MB
\item 96 MB
\item 32 MB
\item 16 MB
\end{enumerate}
\begin{tcolorbox}[enhanced,breakable,colback=teal!4,colframe=teal!50!black,boxrule=0.5pt,arc=1.1mm,left=3mm,right=3mm,top=1.7mm,bottom=1.7mm,before skip=2mm,after skip=2mm]
\textbf{\color{teal}To'g'ri javob:} \textbf{C) 32 MB}
\par\smallskip
\textbf{Izoh:} 96/3 = 32 MB.
\end{tcolorbox}
\vspace{2mm}\hrule height 0.2pt \vspace{3mm}

% DB question_id: ae6c07a7-0910-4524-9db6-c186c8a729ee
\needspace{5\baselineskip}
\noindent\textbf{\color{ink}544.}
\par
O'quvchi qoidaga ko'ra kodlangan xabarni shu qoida yordamida tushunarli boshlang'ich mazmunga qaytardi. Bajarilgan amalni eng aniq tavsiflang.
\begin{enumerate}[label=\Alph*),leftmargin=9mm]
\item Kodlangan axborotni asl mazmuniga qaytarish
\item Ochiq axborotni maxfiy shaklga o'tkazish jarayoni
\item Axborotni kamroq bitda saqlash uchun siqish
\item Tarmoqdagi xabarni ikkinchi qurilmaga qayta yuborish
\end{enumerate}
\begin{tcolorbox}[enhanced,breakable,colback=teal!4,colframe=teal!50!black,boxrule=0.5pt,arc=1.1mm,left=3mm,right=3mm,top=1.7mm,bottom=1.7mm,before skip=2mm,after skip=2mm]
\textbf{\color{teal}To'g'ri javob:} \textbf{A) Kodlangan axborotni asl mazmuniga qaytarish}
\par\smallskip
\textbf{Izoh:} Dekodlash kodlangan xabarni belgilangan moslik yordamida qayta talqin qilib mazmunini tiklashdir. Siqish, shifrlash yoki qayta uzatish bundan boshqa jarayonlar.
\end{tcolorbox}
\vspace{2mm}\hrule height 0.2pt \vspace{3mm}

% DB question_id: ae90c93f-c204-4934-8502-c0bf11a64329
\needspace{5\baselineskip}
\noindent\textbf{\color{ink}545.}
\par
A qabila alifbosida 32 belgi, B qabilada 64 belgi. A 80, B 60 belgidan iborat fixed-length kodli xabar yubordi. Ikkala xabar bit hajmlari yig'indisi nechaga teng?
\begin{enumerate}[label=\Alph*),leftmargin=9mm]
\item 760 bit
\item 400 bit
\item 360 bit
\item 780 bit
\end{enumerate}
\begin{tcolorbox}[enhanced,breakable,colback=teal!4,colframe=teal!50!black,boxrule=0.5pt,arc=1.1mm,left=3mm,right=3mm,top=1.7mm,bottom=1.7mm,before skip=2mm,after skip=2mm]
\textbf{\color{teal}To'g'ri javob:} \textbf{A) 760 bit}
\par\smallskip
\textbf{Izoh:} 32 turli belgi 5 bit, 64 turli belgi 6 bit talab qiladi. 80\ensuremath{\times}5=400 bit va 60\ensuremath{\times}6=360 bit; jami 760 bit.
\end{tcolorbox}
\vspace{2mm}\hrule height 0.2pt \vspace{3mm}

% DB question_id: b265868e-bf27-59b4-8553-c2b5b45136c3
\needspace{5\baselineskip}
\noindent\textbf{\color{ink}546.}
\par
Teng ehtimolli 2 ta holatdan birini ifodalash uchun Xartli formulasi bo'yicha qancha axborot kerak?
\begin{enumerate}[label=\Alph*),leftmargin=9mm]
\item 2 bit
\item 3 bit
\item 1 bit
\item 4 bit
\end{enumerate}
\begin{tcolorbox}[enhanced,breakable,colback=teal!4,colframe=teal!50!black,boxrule=0.5pt,arc=1.1mm,left=3mm,right=3mm,top=1.7mm,bottom=1.7mm,before skip=2mm,after skip=2mm]
\textbf{\color{teal}To'g'ri javob:} \textbf{C) 1 bit}
\par\smallskip
\textbf{Izoh:} I = log\textsubscript{2}2 = 1 bit.
\end{tcolorbox}
\vspace{2mm}\hrule height 0.2pt \vspace{3mm}

% DB question_id: b4d730ca-9c36-4fab-abc2-7318815d33e5
\needspace{5\baselineskip}
\noindent\textbf{\color{ink}547.}
\par
720 MB faylning 1/3 qismi 120 sekundda, yana 25\%i 90 sekundda yuklandi. Qolgan qismi 24 Mbit/s=3 MB/s tezlikda olindi. Butun yuklashning o'rtacha tezligini eng yaqin butun MB/sga yaxlitlang.
\begin{enumerate}[label=\Alph*),leftmargin=9mm]
\item 1 MB/s
\item 2 MB/s
\item 3 MB/s
\item 4 MB/s
\end{enumerate}
\begin{tcolorbox}[enhanced,breakable,colback=teal!4,colframe=teal!50!black,boxrule=0.5pt,arc=1.1mm,left=3mm,right=3mm,top=1.7mm,bottom=1.7mm,before skip=2mm,after skip=2mm]
\textbf{\color{teal}To'g'ri javob:} \textbf{B) 2 MB/s}
\par\smallskip
\textbf{Izoh:} Dastlab 240 MB/120s=2 MB/s, keyin 180 MB/90s=2 MB/s. Qolgan 300 MB 3 MB/s bilan 100s oladi. Jami 720 MB/310s\ensuremath{\approx}2,323 MB/s, eng yaqin butun 2.
\end{tcolorbox}
\vspace{2mm}\hrule height 0.2pt \vspace{3mm}

% DB question_id: b59c91eb-a630-51d3-91d5-668534aebe79
\needspace{5\baselineskip}
\noindent\textbf{\color{ink}548.}
\par
Teng ehtimolli 128 ta holatdan birini ifodalash uchun Xartli formulasi bo'yicha qancha axborot kerak?
\begin{enumerate}[label=\Alph*),leftmargin=9mm]
\item 128 bit
\item 8 bit
\item 6 bit
\item 7 bit
\end{enumerate}
\begin{tcolorbox}[enhanced,breakable,colback=teal!4,colframe=teal!50!black,boxrule=0.5pt,arc=1.1mm,left=3mm,right=3mm,top=1.7mm,bottom=1.7mm,before skip=2mm,after skip=2mm]
\textbf{\color{teal}To'g'ri javob:} \textbf{D) 7 bit}
\par\smallskip
\textbf{Izoh:} I = log\textsubscript{2}128 = 7 bit.
\end{tcolorbox}
\vspace{2mm}\hrule height 0.2pt \vspace{3mm}

% DB question_id: b77de3c2-2151-4974-8ab1-a121f65b1ffd
\needspace{5\baselineskip}
\noindent\textbf{\color{ink}549.}
\par
UTF-8 uchun jadval yetarli, maxfiy xabar uchun esa kalit kerak bo'ldi. Bu vaziyatdagi asosiy tushunchani aniqlang.
\begin{enumerate}[label=\Alph*),leftmargin=9mm]
\item Bit va bayt farqi
\item Rastr va vektor farqi
\item Kodek va konteyner farqi
\item Kodlash va shifrlash farqi
\end{enumerate}
\begin{tcolorbox}[enhanced,breakable,colback=teal!4,colframe=teal!50!black,boxrule=0.5pt,arc=1.1mm,left=3mm,right=3mm,top=1.7mm,bottom=1.7mm,before skip=2mm,after skip=2mm]
\textbf{\color{teal}To'g'ri javob:} \textbf{D) Kodlash va shifrlash farqi}
\par\smallskip
\textbf{Izoh:} Vaziyat Kodlash va shifrlash farqiga misol: kodlash taqdim etish formatini, shifrlash esa maxfiylikni ko'zlashini ajratish.
\end{tcolorbox}
\vspace{2mm}\hrule height 0.2pt \vspace{3mm}

% DB question_id: b84b914e-422a-53f6-965c-43799f665231
\needspace{5\baselineskip}
\noindent\textbf{\color{ink}550.}
\par
7500 ta belgi, har biri 8 baytdan kodlangan. Xabar hajmi qancha?
\begin{enumerate}[label=\Alph*),leftmargin=9mm]
\item 480 000 bit
\item 7 500 bit
\item 480 002 bit
\item 480 003 bit
\end{enumerate}
\begin{tcolorbox}[enhanced,breakable,colback=teal!4,colframe=teal!50!black,boxrule=0.5pt,arc=1.1mm,left=3mm,right=3mm,top=1.7mm,bottom=1.7mm,before skip=2mm,after skip=2mm]
\textbf{\color{teal}To'g'ri javob:} \textbf{A) 480 000 bit}
\par\smallskip
\textbf{Izoh:} Hajm = 7500 \ensuremath{\times} 8 = 60000 B = 480 000 bit.
\end{tcolorbox}
\vspace{2mm}\hrule height 0.2pt \vspace{3mm}

% DB question_id: b9b3494c-74bc-5687-8f58-6a390eb184af
\needspace{5\baselineskip}
\noindent\textbf{\color{ink}551.}
\par
Kitob 300 sahifa, har sahifada 25 satr, har satrda 70 belgi. Har belgi 1 bayt. Matn hajmini MiB da toping.
\begin{enumerate}[label=\Alph*),leftmargin=9mm]
\item 525 000 MiB
\item 0,501 MiB
\item 4 200 000 MiB
\item 525 MiB
\end{enumerate}
\begin{tcolorbox}[enhanced,breakable,colback=teal!4,colframe=teal!50!black,boxrule=0.5pt,arc=1.1mm,left=3mm,right=3mm,top=1.7mm,bottom=1.7mm,before skip=2mm,after skip=2mm]
\textbf{\color{teal}To'g'ri javob:} \textbf{B) 0,501 MiB}
\par\smallskip
\textbf{Izoh:} 300 \ensuremath{\times} 25 \ensuremath{\times} 70 \ensuremath{\times} 1 = 525 000 B = 0,501 MiB.
\end{tcolorbox}
\vspace{2mm}\hrule height 0.2pt \vspace{3mm}

% DB question_id: bf2a3a76-a5a6-4a71-aa25-2bf7df679663
\needspace{5\baselineskip}
\noindent\textbf{\color{ink}552.}
\par
Dasturga kiritilishi mumkin bo'lgan barcha ruxsat etilgan harf, raqam va maxsus ramzlar majmui berilgan. Axborotni kodlash nazariyasida alifbo qanday aniqlanadi?
\begin{enumerate}[label=\Alph*),leftmargin=9mm]
\item Kodlashda ishlatiladigan belgilar to'plami
\item Har bir belgiga mos kodlar ro'yxati
\item Alohida belgini ifodalovchi bitlar ketma-ketligi
\item Kodlangan matnni o'qiydigan dasturiy modul
\end{enumerate}
\begin{tcolorbox}[enhanced,breakable,colback=teal!4,colframe=teal!50!black,boxrule=0.5pt,arc=1.1mm,left=3mm,right=3mm,top=1.7mm,bottom=1.7mm,before skip=2mm,after skip=2mm]
\textbf{\color{teal}To'g'ri javob:} \textbf{A) Kodlashda ishlatiladigan belgilar to'plami}
\par\smallskip
\textbf{Izoh:} Alifbo kodlashda ruxsat etilgan belgilar to'plamidir. Kod jadvali belgilar bilan kodlar mosligini, kod so'zi esa alohida belgi yoki xabarning ifodasini bildiradi.
\end{tcolorbox}
\vspace{2mm}\hrule height 0.2pt \vspace{3mm}

% DB question_id: bf5ac969-c772-5be1-a5dd-3dcddb40f550
\needspace{5\baselineskip}
\noindent\textbf{\color{ink}553.}
\par
Teng ehtimolli 256 ta holatdan birini ifodalash uchun Xartli formulasi bo'yicha qancha axborot kerak?
\begin{enumerate}[label=\Alph*),leftmargin=9mm]
\item 8 bit
\item 256 bit
\item 9 bit
\item 7 bit
\end{enumerate}
\begin{tcolorbox}[enhanced,breakable,colback=teal!4,colframe=teal!50!black,boxrule=0.5pt,arc=1.1mm,left=3mm,right=3mm,top=1.7mm,bottom=1.7mm,before skip=2mm,after skip=2mm]
\textbf{\color{teal}To'g'ri javob:} \textbf{A) 8 bit}
\par\smallskip
\textbf{Izoh:} I = log\textsubscript{2}256 = 8 bit.
\end{tcolorbox}
\vspace{2mm}\hrule height 0.2pt \vspace{3mm}

% DB question_id: c5018fad-3a89-5520-8c9c-25bdfa6d0aef
\needspace{5\baselineskip}
\noindent\textbf{\color{ink}554.}
\par
80 Mbit/s tezlik 2 min davom etsa, qancha axborot uzatiladi?
\begin{enumerate}[label=\Alph*),leftmargin=9mm]
\item 1 200 MB
\item 0,12 200 MB
\item 8 200 MB
\item 0,02 200 MB
\end{enumerate}
\begin{tcolorbox}[enhanced,breakable,colback=teal!4,colframe=teal!50!black,boxrule=0.5pt,arc=1.1mm,left=3mm,right=3mm,top=1.7mm,bottom=1.7mm,before skip=2mm,after skip=2mm]
\textbf{\color{teal}To'g'ri javob:} \textbf{A) 1 200 MB}
\par\smallskip
\textbf{Izoh:} I = v \ensuremath{\times} t; vaqt sekundga va bit/bayt birligi mos holga keltiriladi. Natija: 1 200 MB.
\end{tcolorbox}
\vspace{2mm}\hrule height 0.2pt \vspace{3mm}

% DB question_id: c5826bed-5a23-57cb-807a-4d510c8e4198
\needspace{5\baselineskip}
\noindent\textbf{\color{ink}555.}
\par
256 MiB hajmli fayl 8 MiB/s tezlikda uzatilsa, vaqt qancha?
\begin{enumerate}[label=\Alph*),leftmargin=9mm]
\item 32 s
\item 4 s
\item 256 s
\item 40 s
\end{enumerate}
\begin{tcolorbox}[enhanced,breakable,colback=teal!4,colframe=teal!50!black,boxrule=0.5pt,arc=1.1mm,left=3mm,right=3mm,top=1.7mm,bottom=1.7mm,before skip=2mm,after skip=2mm]
\textbf{\color{teal}To'g'ri javob:} \textbf{A) 32 s}
\par\smallskip
\textbf{Izoh:} Hajm bitga o'tkazilib, t = I/v formula bilan hisoblanadi. Natija: 32 s.
\end{tcolorbox}
\vspace{2mm}\hrule height 0.2pt \vspace{3mm}

% DB question_id: c74c38dd-5533-5b98-8b42-928cb8cb015f
\needspace{5\baselineskip}
\noindent\textbf{\color{ink}556.}
\par
500 MB fayl hajmi 70\% ga kamaydi. Siqilgan hajm qancha?
\begin{enumerate}[label=\Alph*),leftmargin=9mm]
\item 350 MB
\item 10 500 MB
\item 30 MB
\item 150 MB
\end{enumerate}
\begin{tcolorbox}[enhanced,breakable,colback=teal!4,colframe=teal!50!black,boxrule=0.5pt,arc=1.1mm,left=3mm,right=3mm,top=1.7mm,bottom=1.7mm,before skip=2mm,after skip=2mm]
\textbf{\color{teal}To'g'ri javob:} \textbf{D) 150 MB}
\par\smallskip
\textbf{Izoh:} 500 \ensuremath{\times} (1-70/100) = 150 MB.
\end{tcolorbox}
\vspace{2mm}\hrule height 0.2pt \vspace{3mm}

% DB question_id: c9814e9b-d611-5196-97d6-e91f2affa988
\needspace{5\baselineskip}
\noindent\textbf{\color{ink}557.}
\par
Uzatish tezligi nimani bildiradi?
\begin{enumerate}[label=\Alph*),leftmargin=9mm]
\item Fayldagi belgilar sonini
\item Vaqt birligida uzatilgan axborot miqdorini
\item Kanal uzunligini
\item Kod jadvali hajmini
\end{enumerate}
\begin{tcolorbox}[enhanced,breakable,colback=teal!4,colframe=teal!50!black,boxrule=0.5pt,arc=1.1mm,left=3mm,right=3mm,top=1.7mm,bottom=1.7mm,before skip=2mm,after skip=2mm]
\textbf{\color{teal}To'g'ri javob:} \textbf{B) Vaqt birligida uzatilgan axborot miqdorini}
\par\smallskip
\textbf{Izoh:} Tezlik I/t nisbatidir.
\end{tcolorbox}
\vspace{2mm}\hrule height 0.2pt \vspace{3mm}

% DB question_id: cc5acabf-08bf-4134-9fc1-040e5163fba3
\needspace{5\baselineskip}
\noindent\textbf{\color{ink}558.}
\par
25\ensuremath{\times}30 sm rasm 300 dpi bilan skanerlandi. Har piksel RGB uchun 3 bayt, siqish yo'q. 1 dyuym=2,54 sm, 1 MB=10\textsuperscript{6} bayt bo'lsa, hajm taxminan nechchi MB?
\begin{enumerate}[label=\Alph*),leftmargin=9mm]
\item 21 MB
\item 31 MB
\item 42 MB
\item 64 MB
\end{enumerate}
\begin{tcolorbox}[enhanced,breakable,colback=teal!4,colframe=teal!50!black,boxrule=0.5pt,arc=1.1mm,left=3mm,right=3mm,top=1.7mm,bottom=1.7mm,before skip=2mm,after skip=2mm]
\textbf{\color{teal}To'g'ri javob:} \textbf{B) 31 MB}
\par\smallskip
\textbf{Izoh:} Gorizontal piksel \ensuremath{\approx}25/2,54\ensuremath{\times}300, vertikal \ensuremath{\approx}30/2,54\ensuremath{\times}300. Jami \ensuremath{\approx}10,46 mln piksel, 3 baytdan \ensuremath{\approx}31,39 mln bayt, yaqin butun MB 31.
\end{tcolorbox}
\vspace{2mm}\hrule height 0.2pt \vspace{3mm}

% DB question_id: cd374583-f0f9-5154-b71a-c02c1eb8c148
\needspace{5\baselineskip}
\noindent\textbf{\color{ink}559.}
\par
4 096 bit nechta B ga teng? Savolda yozilgan SI/IEC birliklariga aynan amal qiling.
\begin{enumerate}[label=\Alph*),leftmargin=9mm]
\item 512 B
\item 64 B
\item 4 096 B
\item 524 288 B
\end{enumerate}
\begin{tcolorbox}[enhanced,breakable,colback=teal!4,colframe=teal!50!black,boxrule=0.5pt,arc=1.1mm,left=3mm,right=3mm,top=1.7mm,bottom=1.7mm,before skip=2mm,after skip=2mm]
\textbf{\color{teal}To'g'ri javob:} \textbf{A) 512 B}
\par\smallskip
\textbf{Izoh:} Birliklar orasidagi ko'paytirish yoki bo'lish natijasi 512 B.
\end{tcolorbox}
\vspace{2mm}\hrule height 0.2pt \vspace{3mm}

% DB question_id: d0cd28b0-1a6c-4c67-a191-b8a28bc08f08
\needspace{5\baselineskip}
\noindent\textbf{\color{ink}560.}
\par
39 ta ASCII belgidan iborat satr 8 bit/belgi shaklida va shu belgilar 16 bit/belgi shaklida saqlandi. Ikki hajm orasidagi farq necha bayt?
\begin{enumerate}[label=\Alph*),leftmargin=9mm]
\item 18
\item 29
\item 39
\item 78
\end{enumerate}
\begin{tcolorbox}[enhanced,breakable,colback=teal!4,colframe=teal!50!black,boxrule=0.5pt,arc=1.1mm,left=3mm,right=3mm,top=1.7mm,bottom=1.7mm,before skip=2mm,after skip=2mm]
\textbf{\color{teal}To'g'ri javob:} \textbf{C) 39}
\par\smallskip
\textbf{Izoh:} Bir xil 39 belgi 8 bitdan 16 bitga o'tganda har biriga 8 bit=1 bayt ortadi. Jami 39\ensuremath{\times}1=39 bayt farq. Shart qat'iy fixed-width kodlash deb belgilangan.
\end{tcolorbox}
\vspace{2mm}\hrule height 0.2pt \vspace{3mm}

% DB question_id: d31eda2f-1f86-5ed6-a538-bba5eac9f59b
\needspace{5\baselineskip}
\noindent\textbf{\color{ink}561.}
\par
3000 ta belgi, har biri 4 baytdan kodlangan. Xabar hajmi qancha?
\begin{enumerate}[label=\Alph*),leftmargin=9mm]
\item 3 000 B
\item 96 000 B
\item 12 000 B
\item 1 500 B
\end{enumerate}
\begin{tcolorbox}[enhanced,breakable,colback=teal!4,colframe=teal!50!black,boxrule=0.5pt,arc=1.1mm,left=3mm,right=3mm,top=1.7mm,bottom=1.7mm,before skip=2mm,after skip=2mm]
\textbf{\color{teal}To'g'ri javob:} \textbf{C) 12 000 B}
\par\smallskip
\textbf{Izoh:} Hajm = 3000 \ensuremath{\times} 4 = 12000 B.
\end{tcolorbox}
\vspace{2mm}\hrule height 0.2pt \vspace{3mm}

% DB question_id: d3225260-4a5b-55e2-885c-346161edda0e
\needspace{5\baselineskip}
\noindent\textbf{\color{ink}562.}
\par
9 ta turli belgini doimiy uzunlikdagi ikkilik kod bilan kodlash uchun minimal kod uzunligi qancha?
\begin{enumerate}[label=\Alph*),leftmargin=9mm]
\item 3 bit
\item 9 bit
\item 5 bit
\item 4 bit
\end{enumerate}
\begin{tcolorbox}[enhanced,breakable,colback=teal!4,colframe=teal!50!black,boxrule=0.5pt,arc=1.1mm,left=3mm,right=3mm,top=1.7mm,bottom=1.7mm,before skip=2mm,after skip=2mm]
\textbf{\color{teal}To'g'ri javob:} \textbf{D) 4 bit}
\par\smallskip
\textbf{Izoh:} 2\textsuperscript{n} \ensuremath{\geq} 9; minimal n = 4.
\end{tcolorbox}
\vspace{2mm}\hrule height 0.2pt \vspace{3mm}

% DB question_id: d4a2855a-8dcd-54bd-b7ee-a88ee05453c5
\needspace{5\baselineskip}
\noindent\textbf{\color{ink}563.}
\par
300 MB axborot 5 min da uzatildi. O'rtacha tezlikni toping.
\begin{enumerate}[label=\Alph*),leftmargin=9mm]
\item 1 Mbit/s
\item 8 Mbit/s
\item 64 Mbit/s
\item 480 Mbit/s
\end{enumerate}
\begin{tcolorbox}[enhanced,breakable,colback=teal!4,colframe=teal!50!black,boxrule=0.5pt,arc=1.1mm,left=3mm,right=3mm,top=1.7mm,bottom=1.7mm,before skip=2mm,after skip=2mm]
\textbf{\color{teal}To'g'ri javob:} \textbf{B) 8 Mbit/s}
\par\smallskip
\textbf{Izoh:} Vaqt sekundga o'tkazilib, v = I/t formula bilan hisoblanadi. Natija: 8 Mbit/s.
\end{tcolorbox}
\vspace{2mm}\hrule height 0.2pt \vspace{3mm}

% DB question_id: d5cae575-7363-46ac-81b6-9158d715afeb
\needspace{5\baselineskip}
\noindent\textbf{\color{ink}564.}
\par
64 000 Hz, 24-bit, 2 kanalli siqilmagan audio fayl hajmi 120 MiB. Audioyozuv taxminan necha daqiqa davom etgan? Javobni 5 daqiqaga karrali eng yaqin songa yaxlitlang (1 MiB=2\textsuperscript{2}\textsuperscript{0} bayt).
\begin{enumerate}[label=\Alph*),leftmargin=9mm]
\item 5 daqiqa
\item 10 daqiqa
\item 15 daqiqa
\item 20 daqiqa
\end{enumerate}
\begin{tcolorbox}[enhanced,breakable,colback=teal!4,colframe=teal!50!black,boxrule=0.5pt,arc=1.1mm,left=3mm,right=3mm,top=1.7mm,bottom=1.7mm,before skip=2mm,after skip=2mm]
\textbf{\color{teal}To'g'ri javob:} \textbf{A) 5 daqiqa}
\par\smallskip
\textbf{Izoh:} Audio tezligi 64000\ensuremath{\times}24\ensuremath{\times}2/8=384000 bayt/s. Hajm 120\ensuremath{\times}1048576=125829120 bayt. Vaqt \ensuremath{\approx}327,68s\ensuremath{\approx}5,46 daqiqa, 5 daqiqaga karrali eng yaqin qiymat 5 daqiqa.
\end{tcolorbox}
\vspace{2mm}\hrule height 0.2pt \vspace{3mm}

% DB question_id: d704bafe-4c20-594a-8760-89a3cc5f9e2a
\needspace{5\baselineskip}
\noindent\textbf{\color{ink}565.}
\par
Siqilmagan stereo audio 32 000 Hz va 32 bit/kanal bilan yozildi. Fayl hajmi 60 MiB bo'lsa, davomiylikni eng yaqin butun daqiqaga yaxlitlang. 1 MiB=2\textsuperscript{2}\textsuperscript{0} bayt.
\begin{enumerate}[label=\Alph*),leftmargin=9mm]
\item 4 daqiqa
\item 2 daqiqa
\item 8 daqiqa
\item 5 daqiqa
\end{enumerate}
\begin{tcolorbox}[enhanced,breakable,colback=teal!4,colframe=teal!50!black,boxrule=0.5pt,arc=1.1mm,left=3mm,right=3mm,top=1.7mm,bottom=1.7mm,before skip=2mm,after skip=2mm]
\textbf{\color{teal}To'g'ri javob:} \textbf{A) 4 daqiqa}
\par\smallskip
\textbf{Izoh:} Oqim 32 000\ensuremath{\cdot}32\ensuremath{\cdot}2/8=256 000 bayt/s. Vaqt 60\ensuremath{\cdot}2\textsuperscript{2}\textsuperscript{0}/256 000=245,76 s=4,096 daqiqa; eng yaqin butun 4.
\end{tcolorbox}
\vspace{2mm}\hrule height 0.2pt \vspace{3mm}

% DB question_id: d710bce1-ffda-5411-97f9-3cab00e9dca3
\needspace{5\baselineskip}
\noindent\textbf{\color{ink}566.}
\par
26 harfli ingliz alifbosi (\texttt{a=0}) bo'yicha Vijiner usulida \texttt{respublika} so'zini takrorlanuvchi \texttt{partiya} kaliti bilan shifrlang.
\begin{enumerate}[label=\Alph*),leftmargin=9mm]
\item gejiczlxkr
\item gejiczexkr
\item gejcizexkr
\item gjeiczexkr
\end{enumerate}
\begin{tcolorbox}[enhanced,breakable,colback=teal!4,colframe=teal!50!black,boxrule=0.5pt,arc=1.1mm,left=3mm,right=3mm,top=1.7mm,bottom=1.7mm,before skip=2mm,after skip=2mm]
\textbf{\color{teal}To'g'ri javob:} \textbf{A) gejiczlxkr}
\par\smallskip
\textbf{Izoh:} Kalit \texttt{partiyapar} bo'ladi. Mod 26 yig'ish natijasi: r+p=g, e+a=e, s+r=j, p+t=i, u+i=c, b+y=z, l+a=l, i+p=x, k+a=k, a+r=r.
\end{tcolorbox}
\vspace{2mm}\hrule height 0.2pt \vspace{3mm}

% DB question_id: d9f0da82-fef3-5f85-9272-afd227498a29
\needspace{5\baselineskip}
\noindent\textbf{\color{ink}567.}
\par
800 MB fayl hajmi 25\% ga kamaydi. Siqilgan hajm qancha?
\begin{enumerate}[label=\Alph*),leftmargin=9mm]
\item 600 MB
\item 200 MB
\item 15 000 MB
\item 75 MB
\end{enumerate}
\begin{tcolorbox}[enhanced,breakable,colback=teal!4,colframe=teal!50!black,boxrule=0.5pt,arc=1.1mm,left=3mm,right=3mm,top=1.7mm,bottom=1.7mm,before skip=2mm,after skip=2mm]
\textbf{\color{teal}To'g'ri javob:} \textbf{A) 600 MB}
\par\smallskip
\textbf{Izoh:} 800 \ensuremath{\times} (1-25/100) = 600 MB.
\end{tcolorbox}
\vspace{2mm}\hrule height 0.2pt \vspace{3mm}

% DB question_id: dae95796-12f4-4b15-873d-033e962e5385
\needspace{5\baselineskip}
\noindent\textbf{\color{ink}568.}
\par
Matnni kodlashda 'A' harfi, '7' raqami va '?' belgisi alohida element sifatida qaraladi. Shu elementlardan biri <<belgi>> deb nimaga aytiladi?
\begin{enumerate}[label=\Alph*),leftmargin=9mm]
\item Faqat raqamlardan tashkil topgan ketma-ketlik
\item Belgilar to'plamini tartibga keltiruvchi qoida
\item Raqamlar saqlanadigan xotiraning aniq hajmi
\item Axborotni ifodalovchi alohida ramz
\end{enumerate}
\begin{tcolorbox}[enhanced,breakable,colback=teal!4,colframe=teal!50!black,boxrule=0.5pt,arc=1.1mm,left=3mm,right=3mm,top=1.7mm,bottom=1.7mm,before skip=2mm,after skip=2mm]
\textbf{\color{teal}To'g'ri javob:} \textbf{D) Axborotni ifodalovchi alohida ramz}
\par\smallskip
\textbf{Izoh:} Belgi - axborotni ifodalashda ishlatiladigan bitta alohida ramz: harf, raqam, tinish belgisi yoki boshqa simvol. Belgilar to'plami esa alifbo deb ataladi.
\end{tcolorbox}
\vspace{2mm}\hrule height 0.2pt \vspace{3mm}

% DB question_id: daf6adf1-e587-52a4-a467-e01b90d14624
\needspace{5\baselineskip}
\noindent\textbf{\color{ink}569.}
\par
Kanalning nazariy tezligi 250 Mbit/s, amaliy throughput 175 Mbit/s. Samaradorlik necha foiz?
\begin{enumerate}[label=\Alph*),leftmargin=9mm]
\item 0,7 \%
\item 70 \%
\item 30 \%
\item 8,75 \%
\end{enumerate}
\begin{tcolorbox}[enhanced,breakable,colback=teal!4,colframe=teal!50!black,boxrule=0.5pt,arc=1.1mm,left=3mm,right=3mm,top=1.7mm,bottom=1.7mm,before skip=2mm,after skip=2mm]
\textbf{\color{teal}To'g'ri javob:} \textbf{B) 70 \%}
\par\smallskip
\textbf{Izoh:} \ensuremath{\eta} = 175/250 \ensuremath{\times} 100\% = 70\%.
\end{tcolorbox}
\vspace{2mm}\hrule height 0.2pt \vspace{3mm}

% DB question_id: db279ae2-82ef-5ec9-bd33-2fa2885a773f
\needspace{5\baselineskip}
\noindent\textbf{\color{ink}570.}
\par
512 MiB axborot 256 s da uzatildi. O'rtacha tezlikni toping.
\begin{enumerate}[label=\Alph*),leftmargin=9mm]
\item 0,25 MiB/s
\item 16 MiB/s
\item 2 MiB/s
\item 120 MiB/s
\end{enumerate}
\begin{tcolorbox}[enhanced,breakable,colback=teal!4,colframe=teal!50!black,boxrule=0.5pt,arc=1.1mm,left=3mm,right=3mm,top=1.7mm,bottom=1.7mm,before skip=2mm,after skip=2mm]
\textbf{\color{teal}To'g'ri javob:} \textbf{C) 2 MiB/s}
\par\smallskip
\textbf{Izoh:} Vaqt sekundga o'tkazilib, v = I/t formula bilan hisoblanadi. Natija: 2 MiB/s.
\end{tcolorbox}
\vspace{2mm}\hrule height 0.2pt \vspace{3mm}

% DB question_id: db829f4c-9ede-5bbe-aa18-fe8fefe118b4
\needspace{5\baselineskip}
\noindent\textbf{\color{ink}571.}
\par
Teng ehtimolli 32 ta holatdan birini ifodalash uchun Xartli formulasi bo'yicha qancha axborot kerak?
\begin{enumerate}[label=\Alph*),leftmargin=9mm]
\item 32 bit
\item 5 bit
\item 6 bit
\item 4 bit
\end{enumerate}
\begin{tcolorbox}[enhanced,breakable,colback=teal!4,colframe=teal!50!black,boxrule=0.5pt,arc=1.1mm,left=3mm,right=3mm,top=1.7mm,bottom=1.7mm,before skip=2mm,after skip=2mm]
\textbf{\color{teal}To'g'ri javob:} \textbf{B) 5 bit}
\par\smallskip
\textbf{Izoh:} I = log\textsubscript{2}32 = 5 bit.
\end{tcolorbox}
\vspace{2mm}\hrule height 0.2pt \vspace{3mm}

% DB question_id: dffbfc1a-345a-4a6c-b69d-416e05e7ebcf
\needspace{5\baselineskip}
\noindent\textbf{\color{ink}572.}
\par
Server diski 500 GiB, 30\%i band. Har kuni 10 xodimning har biri 256 MiB yangi fayl yuboradi. Boshqa o'zgarishlar bo'lmasa, qolgan joy necha to'liq ish kunida to'ladi? 1 GiB=1024 MiB.
\begin{enumerate}[label=\Alph*),leftmargin=9mm]
\item 120
\item 130
\item 140
\item 150
\end{enumerate}
\begin{tcolorbox}[enhanced,breakable,colback=teal!4,colframe=teal!50!black,boxrule=0.5pt,arc=1.1mm,left=3mm,right=3mm,top=1.7mm,bottom=1.7mm,before skip=2mm,after skip=2mm]
\textbf{\color{teal}To'g'ri javob:} \textbf{C) 140}
\par\smallskip
\textbf{Izoh:} Bo'sh joy 500\ensuremath{\times}0,70=350 GiB=358400 MiB. Bir kunda 10\ensuremath{\times}256=2560 MiB yig'iladi. 358400/2560=140 kun.
\end{tcolorbox}
\vspace{2mm}\hrule height 0.2pt \vspace{3mm}

% DB question_id: e12c1cdc-f339-5393-95f2-d4a5f32aa925
\needspace{5\baselineskip}
\noindent\textbf{\color{ink}573.}
\par
Qaysi tenglik IEC standartiga mos?
\begin{enumerate}[label=\Alph*),leftmargin=9mm]
\item 1 kB = 1024 B
\item 1 MiB = 1000 KiB
\item 1 KiB = 1024 B
\item 1 MB = 1024 B
\end{enumerate}
\begin{tcolorbox}[enhanced,breakable,colback=teal!4,colframe=teal!50!black,boxrule=0.5pt,arc=1.1mm,left=3mm,right=3mm,top=1.7mm,bottom=1.7mm,before skip=2mm,after skip=2mm]
\textbf{\color{teal}To'g'ri javob:} \textbf{C) 1 KiB = 1024 B}
\par\smallskip
\textbf{Izoh:} IEC ikkilik prefiksida 1 KiB = 2\textsuperscript{1}\textsuperscript{0} bayt.
\end{tcolorbox}
\vspace{2mm}\hrule height 0.2pt \vspace{3mm}

% DB question_id: e2e6bda4-7690-5120-a78f-86c375cfc9f2
\needspace{5\baselineskip}
\noindent\textbf{\color{ink}574.}
\par
Doimiy uzunlikdagi kodning belgisi qaysi?
\begin{enumerate}[label=\Alph*),leftmargin=9mm]
\item Har kod turli uzunlikda
\item Kod faqat 1 bit
\item Har bir belgi bir xil uzunlikdagi kod oladi
\item Faqat matn uchun
\end{enumerate}
\begin{tcolorbox}[enhanced,breakable,colback=teal!4,colframe=teal!50!black,boxrule=0.5pt,arc=1.1mm,left=3mm,right=3mm,top=1.7mm,bottom=1.7mm,before skip=2mm,after skip=2mm]
\textbf{\color{teal}To'g'ri javob:} \textbf{C) Har bir belgi bir xil uzunlikdagi kod oladi}
\par\smallskip
\textbf{Izoh:} Barcha kod so'zlari teng uzunlikda bo'ladi.
\end{tcolorbox}
\vspace{2mm}\hrule height 0.2pt \vspace{3mm}

% DB question_id: e3f6e6f6-3ad0-4141-97f7-ff0e3cef98fd
\needspace{5\baselineskip}
\noindent\textbf{\color{ink}575.}
\par
Matndagi harflar UTF-8 baytlariga aylantirildi. Bu vaziyatdagi asosiy tushunchani aniqlang.
\begin{enumerate}[label=\Alph*),leftmargin=9mm]
\item Dekodlash
\item Doimiy uzunlikdagi kod
\item O'zgaruvchan uzunlikdagi kod
\item Kodlash
\end{enumerate}
\begin{tcolorbox}[enhanced,breakable,colback=teal!4,colframe=teal!50!black,boxrule=0.5pt,arc=1.1mm,left=3mm,right=3mm,top=1.7mm,bottom=1.7mm,before skip=2mm,after skip=2mm]
\textbf{\color{teal}To'g'ri javob:} \textbf{D) Kodlash}
\par\smallskip
\textbf{Izoh:} Vaziyat Kodlashga misol: axborotni belgilangan qoida bo'yicha boshqa ko'rinishga o'tkazish.
\end{tcolorbox}
\vspace{2mm}\hrule height 0.2pt \vspace{3mm}

% DB question_id: e4e33fe3-a8e2-524c-92cc-dcd4b6973fa6
\needspace{5\baselineskip}
\noindent\textbf{\color{ink}576.}
\par
Teng ehtimolli 16 ta holatdan birini ifodalash uchun Xartli formulasi bo'yicha qancha axborot kerak?
\begin{enumerate}[label=\Alph*),leftmargin=9mm]
\item 16 bit
\item 4 bit
\item 5 bit
\item 3 bit
\end{enumerate}
\begin{tcolorbox}[enhanced,breakable,colback=teal!4,colframe=teal!50!black,boxrule=0.5pt,arc=1.1mm,left=3mm,right=3mm,top=1.7mm,bottom=1.7mm,before skip=2mm,after skip=2mm]
\textbf{\color{teal}To'g'ri javob:} \textbf{B) 4 bit}
\par\smallskip
\textbf{Izoh:} I = log\textsubscript{2}16 = 4 bit.
\end{tcolorbox}
\vspace{2mm}\hrule height 0.2pt \vspace{3mm}

% DB question_id: e73af66f-181e-57ac-b2ce-4dac25feee58
\needspace{5\baselineskip}
\noindent\textbf{\color{ink}577.}
\par
315\ensuremath{\times}3072 piksel tasvir uchun, sarlavhasiz va siqilmasdan, 735 KiB xotira ajratilgan. Har piksel bir xil butun sondagi bit bilan kodlansa, ajratilgan hajmdan oshmasdan eng ko'pi bilan nechta rang kodlash mumkin? 1 KiB=1024 bayt.
\begin{enumerate}[label=\Alph*),leftmargin=9mm]
\item 64
\item 32
\item 128
\item 256
\end{enumerate}
\begin{tcolorbox}[enhanced,breakable,colback=teal!4,colframe=teal!50!black,boxrule=0.5pt,arc=1.1mm,left=3mm,right=3mm,top=1.7mm,bottom=1.7mm,before skip=2mm,after skip=2mm]
\textbf{\color{teal}To'g'ri javob:} \textbf{A) 64}
\par\smallskip
\textbf{Izoh:} Ajratilgan bitlar 735\ensuremath{\cdot}1024\ensuremath{\cdot}8=6 021 120. Piksel soni 315\ensuremath{\cdot}3072=967 680. Bir pikselga ko'pi bilan \ensuremath{\lfloor}6 021 120/967 680\ensuremath{\rfloor}=6 bit; ranglar soni 2\textsuperscript{6}=64.
\end{tcolorbox}
\vspace{2mm}\hrule height 0.2pt \vspace{3mm}

% DB question_id: e77af2bd-c98b-5a0e-a15b-81d466cf53d3
\needspace{5\baselineskip}
\noindent\textbf{\color{ink}578.}
\par
850 ta belgi, har biri 2 baytdan kodlangan. Xabar hajmi qancha?
\begin{enumerate}[label=\Alph*),leftmargin=9mm]
\item 850 B
\item 13 600 B
\item 212,5 B
\item 1 700 B
\end{enumerate}
\begin{tcolorbox}[enhanced,breakable,colback=teal!4,colframe=teal!50!black,boxrule=0.5pt,arc=1.1mm,left=3mm,right=3mm,top=1.7mm,bottom=1.7mm,before skip=2mm,after skip=2mm]
\textbf{\color{teal}To'g'ri javob:} \textbf{D) 1 700 B}
\par\smallskip
\textbf{Izoh:} Hajm = 850 \ensuremath{\times} 2 = 1700 B.
\end{tcolorbox}
\vspace{2mm}\hrule height 0.2pt \vspace{3mm}

% DB question_id: e7958818-3bac-531e-b4bc-421622c01b18
\needspace{5\baselineskip}
\noindent\textbf{\color{ink}579.}
\par
Kitob 120 sahifa, har sahifada 32 satr, har satrda 64 belgi. Har belgi 2 bayt. Matn hajmini KiB da toping.
\begin{enumerate}[label=\Alph*),leftmargin=9mm]
\item 245 760 KiB
\item 3 932 160 KiB
\item 491,52 KiB
\item 480 KiB
\end{enumerate}
\begin{tcolorbox}[enhanced,breakable,colback=teal!4,colframe=teal!50!black,boxrule=0.5pt,arc=1.1mm,left=3mm,right=3mm,top=1.7mm,bottom=1.7mm,before skip=2mm,after skip=2mm]
\textbf{\color{teal}To'g'ri javob:} \textbf{D) 480 KiB}
\par\smallskip
\textbf{Izoh:} 120 \ensuremath{\times} 32 \ensuremath{\times} 64 \ensuremath{\times} 2 = 491 520 B = 480 KiB.
\end{tcolorbox}
\vspace{2mm}\hrule height 0.2pt \vspace{3mm}

% DB question_id: e9633b58-241d-5530-ae53-e3f9f3505e66
\needspace{5\baselineskip}
\noindent\textbf{\color{ink}580.}
\par
4096 ta belgi, har biri 1 baytdan kodlangan. Xabar hajmi qancha?
\begin{enumerate}[label=\Alph*),leftmargin=9mm]
\item 4 KiB
\item 4 096 KiB
\item 32 768 KiB
\item 512 KiB
\end{enumerate}
\begin{tcolorbox}[enhanced,breakable,colback=teal!4,colframe=teal!50!black,boxrule=0.5pt,arc=1.1mm,left=3mm,right=3mm,top=1.7mm,bottom=1.7mm,before skip=2mm,after skip=2mm]
\textbf{\color{teal}To'g'ri javob:} \textbf{A) 4 KiB}
\par\smallskip
\textbf{Izoh:} Hajm = 4096 \ensuremath{\times} 1 = 4096 B = 4 KiB.
\end{tcolorbox}
\vspace{2mm}\hrule height 0.2pt \vspace{3mm}

% DB question_id: ecebf3fd-fafd-55a0-a014-fd52c74a80eb
\needspace{5\baselineskip}
\noindent\textbf{\color{ink}581.}
\par
12 MiB/s tezlik 75 s davom etsa, qancha axborot uzatiladi?
\begin{enumerate}[label=\Alph*),leftmargin=9mm]
\item 900 MiB
\item 112,5 MiB
\item 7 200 MiB
\item 15 MiB
\end{enumerate}
\begin{tcolorbox}[enhanced,breakable,colback=teal!4,colframe=teal!50!black,boxrule=0.5pt,arc=1.1mm,left=3mm,right=3mm,top=1.7mm,bottom=1.7mm,before skip=2mm,after skip=2mm]
\textbf{\color{teal}To'g'ri javob:} \textbf{A) 900 MiB}
\par\smallskip
\textbf{Izoh:} I = v \ensuremath{\times} t; vaqt sekundga va bit/bayt birligi mos holga keltiriladi. Natija: 900 MiB.
\end{tcolorbox}
\vspace{2mm}\hrule height 0.2pt \vspace{3mm}

% DB question_id: f0e73c85-ff71-5821-9a74-0cf0753e1774
\needspace{5\baselineskip}
\noindent\textbf{\color{ink}582.}
\par
1,5 GB hajmli fayl 50 Mbit/s tezlikda uzatilsa, vaqt qancha?
\begin{enumerate}[label=\Alph*),leftmargin=9mm]
\item 30 s
\item 1 920 s
\item 300 s
\item 240 s
\end{enumerate}
\begin{tcolorbox}[enhanced,breakable,colback=teal!4,colframe=teal!50!black,boxrule=0.5pt,arc=1.1mm,left=3mm,right=3mm,top=1.7mm,bottom=1.7mm,before skip=2mm,after skip=2mm]
\textbf{\color{teal}To'g'ri javob:} \textbf{D) 240 s}
\par\smallskip
\textbf{Izoh:} Hajm bitga o'tkazilib, t = I/v formula bilan hisoblanadi. Natija: 240 s.
\end{tcolorbox}
\vspace{2mm}\hrule height 0.2pt \vspace{3mm}

% DB question_id: f21e12d8-7fe6-4ad7-9179-b15d0bfa51a8
\needspace{5\baselineskip}
\noindent\textbf{\color{ink}583.}
\par
256\ensuremath{\times}128 piksel rangli tasvir siqilmasdan aniq 40 KiB joy egalladi. Har pikselga bir xil butun sonli bit ajratilgan bo'lsa, nechta turli rangni kodlash mumkin?
\begin{enumerate}[label=\Alph*),leftmargin=9mm]
\item 128
\item 256
\item 512
\item 1024
\end{enumerate}
\begin{tcolorbox}[enhanced,breakable,colback=teal!4,colframe=teal!50!black,boxrule=0.5pt,arc=1.1mm,left=3mm,right=3mm,top=1.7mm,bottom=1.7mm,before skip=2mm,after skip=2mm]
\textbf{\color{teal}To'g'ri javob:} \textbf{D) 1024}
\par\smallskip
\textbf{Izoh:} Piksel soni 256\ensuremath{\times}128=32768. Xotira 40\ensuremath{\times}1024\ensuremath{\times}8=327680 bit; bit/piksel=10. 10 bit bilan 2\textsuperscript{1}\textsuperscript{0}=1024 xil rang kodlanadi.
\end{tcolorbox}
\vspace{2mm}\hrule height 0.2pt \vspace{3mm}

% DB question_id: f30cb70d-56c5-500b-a6e5-8d7fa636b1d2
\needspace{5\baselineskip}
\noindent\textbf{\color{ink}584.}
\par
1,25 MB nechta Mbit ga teng? Savolda yozilgan SI/IEC birliklariga aynan amal qiling.
\begin{enumerate}[label=\Alph*),leftmargin=9mm]
\item 1,25 Mbit
\item 10 Mbit
\item 80 Mbit
\item 10 240 Mbit
\end{enumerate}
\begin{tcolorbox}[enhanced,breakable,colback=teal!4,colframe=teal!50!black,boxrule=0.5pt,arc=1.1mm,left=3mm,right=3mm,top=1.7mm,bottom=1.7mm,before skip=2mm,after skip=2mm]
\textbf{\color{teal}To'g'ri javob:} \textbf{B) 10 Mbit}
\par\smallskip
\textbf{Izoh:} Birliklar orasidagi ko'paytirish yoki bo'lish natijasi 10 Mbit.
\end{tcolorbox}
\vspace{2mm}\hrule height 0.2pt \vspace{3mm}

% DB question_id: f3a7f637-3fa0-538b-917b-b8967a2a0c68
\needspace{5\baselineskip}
\noindent\textbf{\color{ink}585.}
\par
900 MB hajmli fayl 15 MB/s tezlikda uzatilsa, vaqt qancha?
\begin{enumerate}[label=\Alph*),leftmargin=9mm]
\item 7,5 s
\item 60 s
\item 480 s
\item 75 s
\end{enumerate}
\begin{tcolorbox}[enhanced,breakable,colback=teal!4,colframe=teal!50!black,boxrule=0.5pt,arc=1.1mm,left=3mm,right=3mm,top=1.7mm,bottom=1.7mm,before skip=2mm,after skip=2mm]
\textbf{\color{teal}To'g'ri javob:} \textbf{B) 60 s}
\par\smallskip
\textbf{Izoh:} Hajm bitga o'tkazilib, t = I/v formula bilan hisoblanadi. Natija: 60 s.
\end{tcolorbox}
\vspace{2mm}\hrule height 0.2pt \vspace{3mm}

% DB question_id: f6237883-bbee-4aab-b940-608cdbba9672
\needspace{5\baselineskip}
\noindent\textbf{\color{ink}586.}
\par
160\ensuremath{\times}256 piksel tasvir xotirada 40 KiB joy egallagan. Har piksel mustaqil, bir xil uzunlikda va siqilmasdan kodlangan bo'lsa, maksimal nechta rang mavjud?
\begin{enumerate}[label=\Alph*),leftmargin=9mm]
\item 64
\item 128
\item 256
\item 512
\end{enumerate}
\begin{tcolorbox}[enhanced,breakable,colback=teal!4,colframe=teal!50!black,boxrule=0.5pt,arc=1.1mm,left=3mm,right=3mm,top=1.7mm,bottom=1.7mm,before skip=2mm,after skip=2mm]
\textbf{\color{teal}To'g'ri javob:} \textbf{C) 256}
\par\smallskip
\textbf{Izoh:} 160\ensuremath{\times}256=40960 piksel. 40 KiB=40960 bayt, demak 1 bayt=8 bit/piksel. Har piksel 2\textsuperscript{8}=256 xil rangni ifodalay oladi.
\end{tcolorbox}
\vspace{2mm}\hrule height 0.2pt \vspace{3mm}

% DB question_id: f7d004d3-d39e-5ca9-a665-06ae862db59a
\needspace{5\baselineskip}
\noindent\textbf{\color{ink}587.}
\par
17 ta turli belgini doimiy uzunlikdagi ikkilik kod bilan kodlash uchun minimal kod uzunligi qancha?
\begin{enumerate}[label=\Alph*),leftmargin=9mm]
\item 5 bit
\item 4 bit
\item 17 bit
\item 6 bit
\end{enumerate}
\begin{tcolorbox}[enhanced,breakable,colback=teal!4,colframe=teal!50!black,boxrule=0.5pt,arc=1.1mm,left=3mm,right=3mm,top=1.7mm,bottom=1.7mm,before skip=2mm,after skip=2mm]
\textbf{\color{teal}To'g'ri javob:} \textbf{A) 5 bit}
\par\smallskip
\textbf{Izoh:} 2\textsuperscript{n} \ensuremath{\geq} 17; minimal n = 5.
\end{tcolorbox}
\vspace{2mm}\hrule height 0.2pt \vspace{3mm}

% DB question_id: f8241875-305b-48bb-a14e-1b35e2137117
\needspace{5\baselineskip}
\noindent\textbf{\color{ink}588.}
\par
Bit nima?
\begin{enumerate}[label=\Alph*),leftmargin=9mm]
\item 8 ta bitdan tashkil topadigan xotira hajmi birligi
\item 1024 ta baytdan tashkil topadigan xotira hajmi birligi
\item 0 yoki 1 qiymatini oladigan eng kichik axborot birligi
\item Protsessor bir amalda qayta ishlaydigan bitlar guruhi
\end{enumerate}
\begin{tcolorbox}[enhanced,breakable,colback=teal!4,colframe=teal!50!black,boxrule=0.5pt,arc=1.1mm,left=3mm,right=3mm,top=1.7mm,bottom=1.7mm,before skip=2mm,after skip=2mm]
\textbf{\color{teal}To'g'ri javob:} \textbf{C) 0 yoki 1 qiymatini oladigan eng kichik axborot birligi}
\par\smallskip
\textbf{Izoh:} Bit - ikkilik sanoqdagi 0 yoki 1 holatini ifodalovchi axborot birligi. Sakkiz bit baytni, 1024 bayt esa kibibaytni ifodalaydi; protsessor bir amalda ishlov beradigan bitlar guruhi mashina so'zi deyiladi. Bular bit ta'rifiga teng emas.
\end{tcolorbox}
\vspace{2mm}\hrule height 0.2pt \vspace{3mm}

% DB question_id: f904b46f-2198-5470-8030-bfcfb7f5be8e
\needspace{5\baselineskip}
\noindent\textbf{\color{ink}589.}
\par
65 536 B nechta KiB ga teng? Savolda yozilgan SI/IEC birliklariga aynan amal qiling.
\begin{enumerate}[label=\Alph*),leftmargin=9mm]
\item 512 KiB
\item 64 KiB
\item 0,062 KiB
\item 65 536 KiB
\end{enumerate}
\begin{tcolorbox}[enhanced,breakable,colback=teal!4,colframe=teal!50!black,boxrule=0.5pt,arc=1.1mm,left=3mm,right=3mm,top=1.7mm,bottom=1.7mm,before skip=2mm,after skip=2mm]
\textbf{\color{teal}To'g'ri javob:} \textbf{B) 64 KiB}
\par\smallskip
\textbf{Izoh:} Birliklar orasidagi ko'paytirish yoki bo'lish natijasi 64 KiB.
\end{tcolorbox}
\vspace{2mm}\hrule height 0.2pt \vspace{3mm}

% DB question_id: fba11cc8-3fcc-585c-9f70-6e597a69f551
\needspace{5\baselineskip}
\noindent\textbf{\color{ink}590.}
\par
Kanalning nazariy tezligi 500 Mbit/s, amaliy throughput 300 Mbit/s. Samaradorlik necha foiz?
\begin{enumerate}[label=\Alph*),leftmargin=9mm]
\item 60 \%
\item 0,6 \%
\item 40 \%
\item 7,5 \%
\end{enumerate}
\begin{tcolorbox}[enhanced,breakable,colback=teal!4,colframe=teal!50!black,boxrule=0.5pt,arc=1.1mm,left=3mm,right=3mm,top=1.7mm,bottom=1.7mm,before skip=2mm,after skip=2mm]
\textbf{\color{teal}To'g'ri javob:} \textbf{A) 60 \%}
\par\smallskip
\textbf{Izoh:} \ensuremath{\eta} = 300/500 \ensuremath{\times} 100\% = 60\%.
\end{tcolorbox}
\vspace{2mm}\hrule height 0.2pt \vspace{3mm}

% DB question_id: fbed1d2f-1819-578d-a581-72bea616fae6
\needspace{5\baselineskip}
\noindent\textbf{\color{ink}591.}
\par
5 ta turli belgini doimiy uzunlikdagi ikkilik kod bilan kodlash uchun minimal kod uzunligi qancha?
\begin{enumerate}[label=\Alph*),leftmargin=9mm]
\item 2 bit
\item 5 bit
\item 3 bit
\item 4 bit
\end{enumerate}
\begin{tcolorbox}[enhanced,breakable,colback=teal!4,colframe=teal!50!black,boxrule=0.5pt,arc=1.1mm,left=3mm,right=3mm,top=1.7mm,bottom=1.7mm,before skip=2mm,after skip=2mm]
\textbf{\color{teal}To'g'ri javob:} \textbf{C) 3 bit}
\par\smallskip
\textbf{Izoh:} 2\textsuperscript{n} \ensuremath{\geq} 5; minimal n = 3.
\end{tcolorbox}
\vspace{2mm}\hrule height 0.2pt \vspace{3mm}

% DB question_id: fd08fe70-278f-5ef8-9bde-c4948a93cfe2
\needspace{5\baselineskip}
\noindent\textbf{\color{ink}592.}
\par
1,5 GiB nechta MiB ga teng? Savolda yozilgan SI/IEC birliklariga aynan amal qiling.
\begin{enumerate}[label=\Alph*),leftmargin=9mm]
\item 192 MiB
\item 1 536 MiB
\item 12 288 MiB
\item 1,5 MiB
\end{enumerate}
\begin{tcolorbox}[enhanced,breakable,colback=teal!4,colframe=teal!50!black,boxrule=0.5pt,arc=1.1mm,left=3mm,right=3mm,top=1.7mm,bottom=1.7mm,before skip=2mm,after skip=2mm]
\textbf{\color{teal}To'g'ri javob:} \textbf{B) 1 536 MiB}
\par\smallskip
\textbf{Izoh:} Birliklar orasidagi ko'paytirish yoki bo'lish natijasi 1 536 MiB.
\end{tcolorbox}
\vspace{2mm}\hrule height 0.2pt \vspace{3mm}

% DB question_id: ff7dcb16-5e4c-5caa-a55d-424c8c51e307
\needspace{5\baselineskip}
\noindent\textbf{\color{ink}593.}
\par
3 TiB nechta GiB ga teng? Savolda yozilgan SI/IEC birliklariga aynan amal qiling.
\begin{enumerate}[label=\Alph*),leftmargin=9mm]
\item 384 GiB
\item 24 576 GiB
\item 3 072 GiB
\item 3 GiB
\end{enumerate}
\begin{tcolorbox}[enhanced,breakable,colback=teal!4,colframe=teal!50!black,boxrule=0.5pt,arc=1.1mm,left=3mm,right=3mm,top=1.7mm,bottom=1.7mm,before skip=2mm,after skip=2mm]
\textbf{\color{teal}To'g'ri javob:} \textbf{C) 3 072 GiB}
\par\smallskip
\textbf{Izoh:} Birliklar orasidagi ko'paytirish yoki bo'lish natijasi 3 072 GiB.
\end{tcolorbox}
\vspace{2mm}\hrule height 0.2pt \vspace{3mm}

% DB question_id: 4c06b5e9-1014-48e8-ac24-d5006016bb9e
\needspace{5\baselineskip}
\noindent\textbf{\color{ink}594.}
\par
Barcha belgilar aynan 8 bitli bloklarda saqlandi. Vaziyatni tahlil qilib, unga eng aniq mos tushunchani toping.
\begin{enumerate}[label=\Alph*),leftmargin=9mm]
\item O'zgaruvchan uzunlikdagi kod
\item Doimiy uzunlikdagi kod
\item Siqish
\item Shifrlash
\end{enumerate}
\begin{tcolorbox}[enhanced,breakable,colback=teal!4,colframe=teal!50!black,boxrule=0.5pt,arc=1.1mm,left=3mm,right=3mm,top=1.7mm,bottom=1.7mm,before skip=2mm,after skip=2mm]
\textbf{\color{teal}To'g'ri javob:} \textbf{B) Doimiy uzunlikdagi kod}
\par\smallskip
\textbf{Izoh:} Vaziyat Doimiy uzunlikdagi kodga misol: har bir belgi bir xil sondagi kod elementi bilan ifodalanadigan kod.
\end{tcolorbox}
\vspace{2mm}\hrule height 0.2pt \vspace{3mm}

% DB question_id: 671b69fc-343c-4bb5-aa60-433104a488a6
\needspace{5\baselineskip}
\noindent\textbf{\color{ink}595.}
\par
Rasm yuborishdan oldin kichikroq faylga aylantirildi. Vaziyatni tahlil qilib, unga eng aniq mos tushunchani toping.
\begin{enumerate}[label=\Alph*),leftmargin=9mm]
\item Shifrlash
\item Asimmetrik shifrlash
\item Belgi
\item Siqish
\end{enumerate}
\begin{tcolorbox}[enhanced,breakable,colback=teal!4,colframe=teal!50!black,boxrule=0.5pt,arc=1.1mm,left=3mm,right=3mm,top=1.7mm,bottom=1.7mm,before skip=2mm,after skip=2mm]
\textbf{\color{teal}To'g'ri javob:} \textbf{D) Siqish}
\par\smallskip
\textbf{Izoh:} Vaziyat Siqishga misol: axborotni saqlash yoki uzatish uchun fayl hajmini kamaytirish jarayoni.
\end{tcolorbox}
\vspace{2mm}\hrule height 0.2pt \vspace{3mm}

% DB question_id: b89ff055-e1bd-4b14-84d7-6430e20cadff
\needspace{5\baselineskip}
\noindent\textbf{\color{ink}596.}
\par
Disk ishlab chiqaruvchisi kB, xotira hisobida KiB yozildi. Vaziyatni tahlil qilib, unga eng aniq mos tushunchani toping.
\begin{enumerate}[label=\Alph*),leftmargin=9mm]
\item Baud va bit/s farqi
\item kB va KiB farqi
\item HTTPS va ishonchlilik farqi
\item Axborot va ma'lumot farqi
\end{enumerate}
\begin{tcolorbox}[enhanced,breakable,colback=teal!4,colframe=teal!50!black,boxrule=0.5pt,arc=1.1mm,left=3mm,right=3mm,top=1.7mm,bottom=1.7mm,before skip=2mm,after skip=2mm]
\textbf{\color{teal}To'g'ri javob:} \textbf{B) kB va KiB farqi}
\par\smallskip
\textbf{Izoh:} Vaziyat kB va KiB farqiga misol: kB 1000 bayt, KiB 1024 bayt ekanini standart belgilari bilan ajratish.
\end{tcolorbox}
\vspace{2mm}\hrule height 0.2pt \vspace{3mm}

% DB question_id: c3d44598-5dc0-4fa2-8982-b78e824b089a
\needspace{5\baselineskip}
\noindent\textbf{\color{ink}597.}
\par
Operativ xotira hisobida 2\textsuperscript{1}\textsuperscript{0} baytlik birlik ishlatildi. Yaqin tushunchalarni farqlab, aynan qaysi atama vaziyatni eng aniq ifodalashini toping.
\begin{enumerate}[label=\Alph*),leftmargin=9mm]
\item Xotira sig'imi
\item KiB
\item Bit
\item Bayt
\end{enumerate}
\begin{tcolorbox}[enhanced,breakable,colback=teal!4,colframe=teal!50!black,boxrule=0.5pt,arc=1.1mm,left=3mm,right=3mm,top=1.7mm,bottom=1.7mm,before skip=2mm,after skip=2mm]
\textbf{\color{teal}To'g'ri javob:} \textbf{B) KiB}
\par\smallskip
\textbf{Izoh:} Vaziyat KiBga misol: IEC standartida 1024 baytga teng ikkilik birlik.
\end{tcolorbox}
\vspace{2mm}\hrule height 0.2pt \vspace{3mm}

% DB question_id: e0c33234-cb9d-4d7a-bb13-59cc40035036
\needspace{5\baselineskip}
\noindent\textbf{\color{ink}598.}
\par
Ko'p uchraydigan belgiga qisqa, kam uchraydiganga uzun kod berildi. Vaziyatni tahlil qilib, unga eng aniq mos tushunchani toping.
\begin{enumerate}[label=\Alph*),leftmargin=9mm]
\item Siqish
\item Shifrlash
\item O'zgaruvchan uzunlikdagi kod
\item Asimmetrik shifrlash
\end{enumerate}
\begin{tcolorbox}[enhanced,breakable,colback=teal!4,colframe=teal!50!black,boxrule=0.5pt,arc=1.1mm,left=3mm,right=3mm,top=1.7mm,bottom=1.7mm,before skip=2mm,after skip=2mm]
\textbf{\color{teal}To'g'ri javob:} \textbf{C) O'zgaruvchan uzunlikdagi kod}
\par\smallskip
\textbf{Izoh:} Vaziyat O'zgaruvchan uzunlikdagi kodga misol: turli belgilar turli uzunlikdagi kodlar bilan ifodalanadigan kod.
\end{tcolorbox}
\vspace{2mm}\hrule height 0.2pt \vspace{3mm}

% DB question_id: d7f61325-0273-4226-a95f-9e43abbaa093
\needspace{5\baselineskip}
\noindent\textbf{\color{ink}599.}
\par
Qaysi tenglik standart jihatdan aniq yozilgan?
\begin{enumerate}[label=\Alph*),leftmargin=9mm]
\item 1 kB = 1024 B
\item 1 KiB = 1000 B
\item 1 MB = 1024 kB har doim
\item 1 KiB = 1024 B
\end{enumerate}
\begin{tcolorbox}[enhanced,breakable,colback=teal!4,colframe=teal!50!black,boxrule=0.5pt,arc=1.1mm,left=3mm,right=3mm,top=1.7mm,bottom=1.7mm,before skip=2mm,after skip=2mm]
\textbf{\color{teal}To'g'ri javob:} \textbf{D) 1 KiB = 1024 B}
\par\smallskip
\textbf{Izoh:} IEC bo'yicha KiB = 1024 B; kB esa SI bo'yicha 1000 B.
\end{tcolorbox}
\vspace{2mm}\hrule height 0.2pt \vspace{3mm}
\section{Matn, grafika, audio va videoning raqamli tasvirlanishi}

% DB question_id: 0d9dc44a-b9ec-4aa7-9634-3b44ad597c6c
\needspace{5\baselineskip}
\noindent\textbf{\color{ink}600.}
\par
Quyidagi ta'rif qaysi atamaga tegishli? Matndagi har bir belgi yoki kod nuqtasini raqamli kod bilan ifodalash jarayoni.
\begin{enumerate}[label=\Alph*),leftmargin=9mm]
\item ASCII
\item Kengaytirilgan ASCII
\item Unicode
\item Belgini kodlash
\end{enumerate}
\begin{tcolorbox}[enhanced,breakable,colback=teal!4,colframe=teal!50!black,boxrule=0.5pt,arc=1.1mm,left=3mm,right=3mm,top=1.7mm,bottom=1.7mm,before skip=2mm,after skip=2mm]
\textbf{\color{teal}To'g'ri javob:} \textbf{D) Belgini kodlash}
\par\smallskip
\textbf{Izoh:} Ta'rif Belgini kodlash tushunchasiga tegishli.
\end{tcolorbox}
\vspace{2mm}\hrule height 0.2pt \vspace{3mm}

% DB question_id: 101739ed-5313-4edd-b00a-9bf3511bd7b3
\needspace{5\baselineskip}
\noindent\textbf{\color{ink}601.}
\par
"Siqish" tushunchasining "Fayl formatlari, kodeklar va siqish" bobidagi eng aniq ta'rifini toping.
\begin{enumerate}[label=\Alph*),leftmargin=9mm]
\item dekodlanganda asl ma'lumotni bitma-bit tiklaydigan siqish
\item axborotni kamroq bit yoki bayt bilan ifodalash jarayoni
\item hajmni kamaytirish uchun ayrim ma'lumotni qaytmas tarzda olib tashlaydigan siqish
\item ketma-ket takrorlangan qiymatni qiymat va takrorlar soni bilan ifodalash usuli
\end{enumerate}
\begin{tcolorbox}[enhanced,breakable,colback=teal!4,colframe=teal!50!black,boxrule=0.5pt,arc=1.1mm,left=3mm,right=3mm,top=1.7mm,bottom=1.7mm,before skip=2mm,after skip=2mm]
\textbf{\color{teal}To'g'ri javob:} \textbf{B) axborotni kamroq bit yoki bayt bilan ifodalash jarayoni}
\par\smallskip
\textbf{Izoh:} Siqish - axborotni kamroq bit yoki bayt bilan ifodalash jarayoni.
\end{tcolorbox}
\vspace{2mm}\hrule height 0.2pt \vspace{3mm}

% DB question_id: 1316dd43-d6ad-4d21-b313-34cf5521e216
\needspace{5\baselineskip}
\noindent\textbf{\color{ink}602.}
\par
"Matn metama'lumoti" atamasiga eng aniq ta'rifni toping.
\begin{enumerate}[label=\Alph*),leftmargin=9mm]
\item matndagi har bir belgi yoki kod nuqtasini raqamli kod bilan ifodalash jarayoni
\item asl ko'rinishida 7 bitli va 128 kod holatiga ega belgilar standarti
\item faylning kodlash turi, tuzilishi yoki xizmat ma'lumotlarini tavsiflovchi qo'shimcha ma'lumot
\item 8 bitli kod sahifalariga beriladigan umumiy nom bo'lib, yagona universal jadval emas
\end{enumerate}
\begin{tcolorbox}[enhanced,breakable,colback=teal!4,colframe=teal!50!black,boxrule=0.5pt,arc=1.1mm,left=3mm,right=3mm,top=1.7mm,bottom=1.7mm,before skip=2mm,after skip=2mm]
\textbf{\color{teal}To'g'ri javob:} \textbf{C) faylning kodlash turi, tuzilishi yoki xizmat ma'lumotlarini tavsiflovchi qo'shimcha ma'lumot}
\par\smallskip
\textbf{Izoh:} Matn metama'lumoti - faylning kodlash turi, tuzilishi yoki xizmat ma'lumotlarini tavsiflovchi qo'shimcha ma'lumot.
\end{tcolorbox}
\vspace{2mm}\hrule height 0.2pt \vspace{3mm}

% DB question_id: 13328e56-99d9-479a-a778-611a3042d534
\needspace{5\baselineskip}
\noindent\textbf{\color{ink}603.}
\par
"Kanal" atamasiga eng aniq ta'rifni toping.
\begin{enumerate}[label=\Alph*),leftmargin=9mm]
\item namunalar qiymatini odatda siqilmagan holda ketma-ket saqlash usuli
\item bir soniyada uzatiladigan yoki saqlanadigan audio bitlari soni
\item mustaqil audio signal yo'li
\item vaqt va amplituda bo'yicha uzluksiz o'zgaradigan tovush signali
\end{enumerate}
\begin{tcolorbox}[enhanced,breakable,colback=teal!4,colframe=teal!50!black,boxrule=0.5pt,arc=1.1mm,left=3mm,right=3mm,top=1.7mm,bottom=1.7mm,before skip=2mm,after skip=2mm]
\textbf{\color{teal}To'g'ri javob:} \textbf{C) mustaqil audio signal yo'li}
\par\smallskip
\textbf{Izoh:} Kanal - mustaqil audio signal yo'li.
\end{tcolorbox}
\vspace{2mm}\hrule height 0.2pt \vspace{3mm}

% DB question_id: 1454a1f0-db8f-43ce-af72-699b4258ed35
\needspace{5\baselineskip}
\noindent\textbf{\color{ink}604.}
\par
"ASCII" atamasiga eng aniq ta'rifni toping.
\begin{enumerate}[label=\Alph*),leftmargin=9mm]
\item asl ko'rinishida 7 bitli va 128 kod holatiga ega belgilar standarti
\item 8 bitli kod sahifalariga beriladigan umumiy nom bo'lib, yagona universal jadval emas
\item turli yozuv tizimlari belgilariga yagona kod nuqtalari ajratadigan standart
\item Unicode tizimida belgiga berilgan abstrakt raqam
\end{enumerate}
\begin{tcolorbox}[enhanced,breakable,colback=teal!4,colframe=teal!50!black,boxrule=0.5pt,arc=1.1mm,left=3mm,right=3mm,top=1.7mm,bottom=1.7mm,before skip=2mm,after skip=2mm]
\textbf{\color{teal}To'g'ri javob:} \textbf{A) asl ko'rinishida 7 bitli va 128 kod holatiga ega belgilar standarti}
\par\smallskip
\textbf{Izoh:} ASCII - asl ko'rinishida 7 bitli va 128 kod holatiga ega belgilar standarti.
\end{tcolorbox}
\vspace{2mm}\hrule height 0.2pt \vspace{3mm}

% DB question_id: 16f1abe4-0f08-4692-930d-ac1d56390d9b
\needspace{5\baselineskip}
\noindent\textbf{\color{ink}605.}
\par
"Media oqimi" atamasiga eng aniq ta'rifni toping.
\begin{enumerate}[label=\Alph*),leftmargin=9mm]
\item vaqt bo'yicha ketma-ket uzatiladigan audio, video yoki boshqa media ma'lumoti
\item bir nechta media turini yagona mazmun yoki interaktiv tizimda birlashtirish
\item videodagi bitta alohida tasvir
\item bir soniyada ko'rsatiladigan kadrlar soni
\end{enumerate}
\begin{tcolorbox}[enhanced,breakable,colback=teal!4,colframe=teal!50!black,boxrule=0.5pt,arc=1.1mm,left=3mm,right=3mm,top=1.7mm,bottom=1.7mm,before skip=2mm,after skip=2mm]
\textbf{\color{teal}To'g'ri javob:} \textbf{A) vaqt bo'yicha ketma-ket uzatiladigan audio, video yoki boshqa media ma'lumoti}
\par\smallskip
\textbf{Izoh:} Media oqimi - vaqt bo'yicha ketma-ket uzatiladigan audio, video yoki boshqa media ma'lumoti.
\end{tcolorbox}
\vspace{2mm}\hrule height 0.2pt \vspace{3mm}

% DB question_id: 180cfcb7-1024-4674-a067-05d70316d7d7
\needspace{5\baselineskip}
\noindent\textbf{\color{ink}606.}
\par
"Unicode" atamasiga eng aniq ta'rifni toping.
\begin{enumerate}[label=\Alph*),leftmargin=9mm]
\item Unicode tizimida belgiga berilgan abstrakt raqam
\item Unicode kod nuqtalarini 1-4 baytli ketma-ketliklarda kodlaydigan shakl
\item turli yozuv tizimlari belgilariga yagona kod nuqtalari ajratadigan standart
\item Unicode kod nuqtalarini 2 yoki 4 baytli birliklarda kodlaydigan shakl
\end{enumerate}
\begin{tcolorbox}[enhanced,breakable,colback=teal!4,colframe=teal!50!black,boxrule=0.5pt,arc=1.1mm,left=3mm,right=3mm,top=1.7mm,bottom=1.7mm,before skip=2mm,after skip=2mm]
\textbf{\color{teal}To'g'ri javob:} \textbf{C) turli yozuv tizimlari belgilariga yagona kod nuqtalari ajratadigan standart}
\par\smallskip
\textbf{Izoh:} Unicode - turli yozuv tizimlari belgilariga yagona kod nuqtalari ajratadigan standart.
\end{tcolorbox}
\vspace{2mm}\hrule height 0.2pt \vspace{3mm}

% DB question_id: 25e8954d-8433-473c-9a10-fa5b9324a935
\needspace{5\baselineskip}
\noindent\textbf{\color{ink}607.}
\par
"Audio bit tezligi" atamasiga eng aniq ta'rifni toping.
\begin{enumerate}[label=\Alph*),leftmargin=9mm]
\item bir soniyada uzatiladigan yoki saqlanadigan audio bitlari soni
\item vaqt va amplituda bo'yicha uzluksiz o'zgaradigan tovush signali
\item analog signalni raqamli namunalarga aylantiradigan qurilma yoki jarayon
\item bir soniyada olingan audio namunalar soni
\end{enumerate}
\begin{tcolorbox}[enhanced,breakable,colback=teal!4,colframe=teal!50!black,boxrule=0.5pt,arc=1.1mm,left=3mm,right=3mm,top=1.7mm,bottom=1.7mm,before skip=2mm,after skip=2mm]
\textbf{\color{teal}To'g'ri javob:} \textbf{A) bir soniyada uzatiladigan yoki saqlanadigan audio bitlari soni}
\par\smallskip
\textbf{Izoh:} Audio bit tezligi - bir soniyada uzatiladigan yoki saqlanadigan audio bitlari soni.
\end{tcolorbox}
\vspace{2mm}\hrule height 0.2pt \vspace{3mm}

% DB question_id: 348c944f-2b1c-4a5d-87a2-03e296d6447a
\needspace{5\baselineskip}
\noindent\textbf{\color{ink}608.}
\par
"Kodek" tushunchasining "Fayl formatlari, kodeklar va siqish" bobidagi eng aniq ta'rifini toping.
\begin{enumerate}[label=\Alph*),leftmargin=9mm]
\item media oqimini bir kodek yoki parametrdan boshqasiga qayta kodlash
\item ma'lumot yo'qolsa tiklash uchun alohida joyda saqlangan nusxa
\item media oqimini kodlash va dekodlash qoidalari yoki dasturiy moduli
\item fayldagi ma'lumotning tuzilishi va uni talqin qilish qoidalari
\end{enumerate}
\begin{tcolorbox}[enhanced,breakable,colback=teal!4,colframe=teal!50!black,boxrule=0.5pt,arc=1.1mm,left=3mm,right=3mm,top=1.7mm,bottom=1.7mm,before skip=2mm,after skip=2mm]
\textbf{\color{teal}To'g'ri javob:} \textbf{C) media oqimini kodlash va dekodlash qoidalari yoki dasturiy moduli}
\par\smallskip
\textbf{Izoh:} Kodek - media oqimini kodlash va dekodlash qoidalari yoki dasturiy moduli.
\end{tcolorbox}
\vspace{2mm}\hrule height 0.2pt \vspace{3mm}

% DB question_id: 35a6f0ae-2be8-47e4-9271-0352c8fdf431
\needspace{5\baselineskip}
\noindent\textbf{\color{ink}609.}
\par
"Konteyner" atamasiga eng aniq ta'rifni toping.
\begin{enumerate}[label=\Alph*),leftmargin=9mm]
\item vaqt bo'yicha ketma-ket uzatiladigan audio, video yoki boshqa media ma'lumoti
\item bir nechta media turini yagona mazmun yoki interaktiv tizimda birlashtirish
\item videodagi bitta alohida tasvir
\item video, audio, subtitr va metama'lumot oqimlarini bir faylda birlashtiradigan tuzilma
\end{enumerate}
\begin{tcolorbox}[enhanced,breakable,colback=teal!4,colframe=teal!50!black,boxrule=0.5pt,arc=1.1mm,left=3mm,right=3mm,top=1.7mm,bottom=1.7mm,before skip=2mm,after skip=2mm]
\textbf{\color{teal}To'g'ri javob:} \textbf{D) video, audio, subtitr va metama'lumot oqimlarini bir faylda birlashtiradigan tuzilma}
\par\smallskip
\textbf{Izoh:} Konteyner - video, audio, subtitr va metama'lumot oqimlarini bir faylda birlashtiradigan tuzilma.
\end{tcolorbox}
\vspace{2mm}\hrule height 0.2pt \vspace{3mm}

% DB question_id: 3709763c-c304-4a17-a56c-9ac14e2b9d10
\needspace{5\baselineskip}
\noindent\textbf{\color{ink}610.}
\par
"UTF-8" atamasiga eng aniq ta'rifni toping.
\begin{enumerate}[label=\Alph*),leftmargin=9mm]
\item Unicode kod nuqtalarini 1-4 baytli ketma-ketliklarda kodlaydigan shakl
\item Unicode kod nuqtalarini 2 yoki 4 baytli birliklarda kodlaydigan shakl
\item faylning kodlash turi, tuzilishi yoki xizmat ma'lumotlarini tavsiflovchi qo'shimcha ma'lumot
\item matndagi har bir belgi yoki kod nuqtasini raqamli kod bilan ifodalash jarayoni
\end{enumerate}
\begin{tcolorbox}[enhanced,breakable,colback=teal!4,colframe=teal!50!black,boxrule=0.5pt,arc=1.1mm,left=3mm,right=3mm,top=1.7mm,bottom=1.7mm,before skip=2mm,after skip=2mm]
\textbf{\color{teal}To'g'ri javob:} \textbf{A) Unicode kod nuqtalarini 1-4 baytli ketma-ketliklarda kodlaydigan shakl}
\par\smallskip
\textbf{Izoh:} UTF-8 - Unicode kod nuqtalarini 1-4 baytli ketma-ketliklarda kodlaydigan shakl.
\end{tcolorbox}
\vspace{2mm}\hrule height 0.2pt \vspace{3mm}

% DB question_id: 40faeebd-0177-52b0-a6af-ac550ea00de0
\needspace{5\baselineskip}
\noindent\textbf{\color{ink}611.}
\par
Qaysi holatda yo'qotishsiz siqish zarur: fayl tiklanganda dastlabki bitlar aynan qayta olinishi kerak?
\begin{enumerate}[label=\Alph*),leftmargin=9mm]
\item Faqat rang tanlashda
\item Hujjat, matn, dastur va sonli ma'lumotlarda
\item Faqat audio va videoda
\item Faqat fotosuratni ko'rsatishda
\end{enumerate}
\begin{tcolorbox}[enhanced,breakable,colback=teal!4,colframe=teal!50!black,boxrule=0.5pt,arc=1.1mm,left=3mm,right=3mm,top=1.7mm,bottom=1.7mm,before skip=2mm,after skip=2mm]
\textbf{\color{teal}To'g'ri javob:} \textbf{B) Hujjat, matn, dastur va sonli ma'lumotlarda}
\par\smallskip
\textbf{Izoh:} Matn, dastur, hujjat va sonli ma'lumotlarda bitta bitning o'zgarishi ham mazmunni buzishi mumkin; shu sabab aynan tiklanadigan yo'qotishsiz siqish qo'llanadi.
\end{tcolorbox}
\vspace{2mm}\hrule height 0.2pt \vspace{3mm}

% DB question_id: 442fbbaa-2f7c-4a0b-a848-5b707d3649f6
\needspace{5\baselineskip}
\noindent\textbf{\color{ink}612.}
\par
Quyidagi ta'rif qaysi atamaga tegishli? Bir soniyada ko'rsatiladigan kadrlar soni.
\begin{enumerate}[label=\Alph*),leftmargin=9mm]
\item Video aniqligi
\item Video bit tezligi
\item Kodek
\item FPS
\end{enumerate}
\begin{tcolorbox}[enhanced,breakable,colback=teal!4,colframe=teal!50!black,boxrule=0.5pt,arc=1.1mm,left=3mm,right=3mm,top=1.7mm,bottom=1.7mm,before skip=2mm,after skip=2mm]
\textbf{\color{teal}To'g'ri javob:} \textbf{D) FPS}
\par\smallskip
\textbf{Izoh:} Ta'rif FPS tushunchasiga tegishli.
\end{tcolorbox}
\vspace{2mm}\hrule height 0.2pt \vspace{3mm}

% DB question_id: 457845bb-7b03-4547-a74a-c2d32ea3475c
\needspace{5\baselineskip}
\noindent\textbf{\color{ink}613.}
\par
"Bit chuqurligi" atamasiga eng aniq ta'rifni toping.
\begin{enumerate}[label=\Alph*),leftmargin=9mm]
\item mustaqil audio signal yo'li
\item bitta audio namunasini kodlash uchun ajratilgan bitlar soni
\item namunalar qiymatini odatda siqilmagan holda ketma-ket saqlash usuli
\item bir soniyada uzatiladigan yoki saqlanadigan audio bitlari soni
\end{enumerate}
\begin{tcolorbox}[enhanced,breakable,colback=teal!4,colframe=teal!50!black,boxrule=0.5pt,arc=1.1mm,left=3mm,right=3mm,top=1.7mm,bottom=1.7mm,before skip=2mm,after skip=2mm]
\textbf{\color{teal}To'g'ri javob:} \textbf{B) bitta audio namunasini kodlash uchun ajratilgan bitlar soni}
\par\smallskip
\textbf{Izoh:} Bit chuqurligi - bitta audio namunasini kodlash uchun ajratilgan bitlar soni.
\end{tcolorbox}
\vspace{2mm}\hrule height 0.2pt \vspace{3mm}

% DB question_id: 489a56f5-4646-494e-a9ba-900031802f57
\needspace{5\baselineskip}
\noindent\textbf{\color{ink}614.}
\par
"ADC" atamasiga eng aniq ta'rifni toping.
\begin{enumerate}[label=\Alph*),leftmargin=9mm]
\item bir soniyada olingan audio namunalar soni
\item namuna amplitudasini ruxsat etilgan sonli darajalardan biriga moslash jarayoni
\item analog signalni raqamli namunalarga aylantiradigan qurilma yoki jarayon
\item bitta audio namunasini kodlash uchun ajratilgan bitlar soni
\end{enumerate}
\begin{tcolorbox}[enhanced,breakable,colback=teal!4,colframe=teal!50!black,boxrule=0.5pt,arc=1.1mm,left=3mm,right=3mm,top=1.7mm,bottom=1.7mm,before skip=2mm,after skip=2mm]
\textbf{\color{teal}To'g'ri javob:} \textbf{C) analog signalni raqamli namunalarga aylantiradigan qurilma yoki jarayon}
\par\smallskip
\textbf{Izoh:} ADC - analog signalni raqamli namunalarga aylantiradigan qurilma yoki jarayon.
\end{tcolorbox}
\vspace{2mm}\hrule height 0.2pt \vspace{3mm}

% DB question_id: 5b26430f-ab0f-4114-9b18-3a6ae7b1815d
\needspace{5\baselineskip}
\noindent\textbf{\color{ink}615.}
\par
"Kvantlash" atamasiga eng aniq ta'rifni toping.
\begin{enumerate}[label=\Alph*),leftmargin=9mm]
\item namuna amplitudasini ruxsat etilgan sonli darajalardan biriga moslash jarayoni
\item bitta audio namunasini kodlash uchun ajratilgan bitlar soni
\item mustaqil audio signal yo'li
\item namunalar qiymatini odatda siqilmagan holda ketma-ket saqlash usuli
\end{enumerate}
\begin{tcolorbox}[enhanced,breakable,colback=teal!4,colframe=teal!50!black,boxrule=0.5pt,arc=1.1mm,left=3mm,right=3mm,top=1.7mm,bottom=1.7mm,before skip=2mm,after skip=2mm]
\textbf{\color{teal}To'g'ri javob:} \textbf{A) namuna amplitudasini ruxsat etilgan sonli darajalardan biriga moslash jarayoni}
\par\smallskip
\textbf{Izoh:} Kvantlash - namuna amplitudasini ruxsat etilgan sonli darajalardan biriga moslash jarayoni.
\end{tcolorbox}
\vspace{2mm}\hrule height 0.2pt \vspace{3mm}

% DB question_id: 5f1e5130-7ae9-464a-9b14-0099aca417a4
\needspace{5\baselineskip}
\noindent\textbf{\color{ink}616.}
\par
"Fayl kengaytmasi" atamasiga eng aniq ta'rifni toping.
\begin{enumerate}[label=\Alph*),leftmargin=9mm]
\item fayl nomining formatni taxmin qilishga yordam beradigan oxirgi qismi
\item axborotni kamroq bit yoki bayt bilan ifodalash jarayoni
\item dekodlanganda asl ma'lumotni bitma-bit tiklaydigan siqish
\item hajmni kamaytirish uchun ayrim ma'lumotni qaytmas tarzda olib tashlaydigan siqish
\end{enumerate}
\begin{tcolorbox}[enhanced,breakable,colback=teal!4,colframe=teal!50!black,boxrule=0.5pt,arc=1.1mm,left=3mm,right=3mm,top=1.7mm,bottom=1.7mm,before skip=2mm,after skip=2mm]
\textbf{\color{teal}To'g'ri javob:} \textbf{A) fayl nomining formatni taxmin qilishga yordam beradigan oxirgi qismi}
\par\smallskip
\textbf{Izoh:} Fayl kengaytmasi - fayl nomining formatni taxmin qilishga yordam beradigan oxirgi qismi.
\end{tcolorbox}
\vspace{2mm}\hrule height 0.2pt \vspace{3mm}

% DB question_id: 6a7afb21-e982-4be7-84b3-b607bd54459f
\needspace{5\baselineskip}
\noindent\textbf{\color{ink}617.}
\par
"Video bit tezligi" atamasiga eng aniq ta'rifni toping.
\begin{enumerate}[label=\Alph*),leftmargin=9mm]
\item media oqimini kodlaydigan va dekodlaydigan algoritm yoki dasturiy modul
\item video oqimining bir soniyada uzatadigan bitlari soni
\item video, audio, subtitr va metama'lumot oqimlarini bir faylda birlashtiradigan tuzilma
\item vaqt bo'yicha ketma-ket uzatiladigan audio, video yoki boshqa media ma'lumoti
\end{enumerate}
\begin{tcolorbox}[enhanced,breakable,colback=teal!4,colframe=teal!50!black,boxrule=0.5pt,arc=1.1mm,left=3mm,right=3mm,top=1.7mm,bottom=1.7mm,before skip=2mm,after skip=2mm]
\textbf{\color{teal}To'g'ri javob:} \textbf{B) video oqimining bir soniyada uzatadigan bitlari soni}
\par\smallskip
\textbf{Izoh:} Video bit tezligi - video oqimining bir soniyada uzatadigan bitlari soni.
\end{tcolorbox}
\vspace{2mm}\hrule height 0.2pt \vspace{3mm}

% DB question_id: 7631e783-2b88-4551-a07f-ce802e97c466
\needspace{5\baselineskip}
\noindent\textbf{\color{ink}618.}
\par
"UTF-16" atamasiga eng aniq ta'rifni toping.
\begin{enumerate}[label=\Alph*),leftmargin=9mm]
\item faylning kodlash turi, tuzilishi yoki xizmat ma'lumotlarini tavsiflovchi qo'shimcha ma'lumot
\item Unicode kod nuqtalarini 2 yoki 4 baytli birliklarda kodlaydigan shakl
\item matndagi har bir belgi yoki kod nuqtasini raqamli kod bilan ifodalash jarayoni
\item asl ko'rinishida 7 bitli va 128 kod holatiga ega belgilar standarti
\end{enumerate}
\begin{tcolorbox}[enhanced,breakable,colback=teal!4,colframe=teal!50!black,boxrule=0.5pt,arc=1.1mm,left=3mm,right=3mm,top=1.7mm,bottom=1.7mm,before skip=2mm,after skip=2mm]
\textbf{\color{teal}To'g'ri javob:} \textbf{B) Unicode kod nuqtalarini 2 yoki 4 baytli birliklarda kodlaydigan shakl}
\par\smallskip
\textbf{Izoh:} UTF-16 - Unicode kod nuqtalarini 2 yoki 4 baytli birliklarda kodlaydigan shakl.
\end{tcolorbox}
\vspace{2mm}\hrule height 0.2pt \vspace{3mm}

% DB question_id: 7e49d87b-a6a8-40e8-96b4-6216a745cd08
\needspace{5\baselineskip}
\noindent\textbf{\color{ink}619.}
\par
"Namuna chastotasi" atamasiga eng aniq ta'rifni toping.
\begin{enumerate}[label=\Alph*),leftmargin=9mm]
\item namuna amplitudasini ruxsat etilgan sonli darajalardan biriga moslash jarayoni
\item bitta audio namunasini kodlash uchun ajratilgan bitlar soni
\item mustaqil audio signal yo'li
\item bir soniyada olingan audio namunalar soni
\end{enumerate}
\begin{tcolorbox}[enhanced,breakable,colback=teal!4,colframe=teal!50!black,boxrule=0.5pt,arc=1.1mm,left=3mm,right=3mm,top=1.7mm,bottom=1.7mm,before skip=2mm,after skip=2mm]
\textbf{\color{teal}To'g'ri javob:} \textbf{D) bir soniyada olingan audio namunalar soni}
\par\smallskip
\textbf{Izoh:} Namuna chastotasi - bir soniyada olingan audio namunalar soni.
\end{tcolorbox}
\vspace{2mm}\hrule height 0.2pt \vspace{3mm}

% DB question_id: 83d46683-ce19-4109-b649-e5950ae12653
\needspace{5\baselineskip}
\noindent\textbf{\color{ink}620.}
\par
"FPS" atamasiga eng aniq ta'rifni toping.
\begin{enumerate}[label=\Alph*),leftmargin=9mm]
\item har bir kadrning gorizontal va vertikal piksel o'lchami
\item video oqimining bir soniyada uzatadigan bitlari soni
\item media oqimini kodlaydigan va dekodlaydigan algoritm yoki dasturiy modul
\item bir soniyada ko'rsatiladigan kadrlar soni
\end{enumerate}
\begin{tcolorbox}[enhanced,breakable,colback=teal!4,colframe=teal!50!black,boxrule=0.5pt,arc=1.1mm,left=3mm,right=3mm,top=1.7mm,bottom=1.7mm,before skip=2mm,after skip=2mm]
\textbf{\color{teal}To'g'ri javob:} \textbf{D) bir soniyada ko'rsatiladigan kadrlar soni}
\par\smallskip
\textbf{Izoh:} FPS - bir soniyada ko'rsatiladigan kadrlar soni.
\end{tcolorbox}
\vspace{2mm}\hrule height 0.2pt \vspace{3mm}

% DB question_id: 877fe53a-36ec-4866-b9fd-4a0321693f42
\needspace{5\baselineskip}
\noindent\textbf{\color{ink}621.}
\par
"Multimedia" tushunchasining "Videoaxborot va multimedia" bobidagi eng aniq ta'rifini toping.
\begin{enumerate}[label=\Alph*),leftmargin=9mm]
\item videodagi bitta alohida tasvir
\item bir nechta media turini yagona mazmun yoki interaktiv tizimda birlashtirish
\item bir soniyada ko'rsatiladigan kadrlar soni
\item har bir kadrning gorizontal va vertikal piksel o'lchami
\end{enumerate}
\begin{tcolorbox}[enhanced,breakable,colback=teal!4,colframe=teal!50!black,boxrule=0.5pt,arc=1.1mm,left=3mm,right=3mm,top=1.7mm,bottom=1.7mm,before skip=2mm,after skip=2mm]
\textbf{\color{teal}To'g'ri javob:} \textbf{B) bir nechta media turini yagona mazmun yoki interaktiv tizimda birlashtirish}
\par\smallskip
\textbf{Izoh:} Multimedia - bir nechta media turini yagona mazmun yoki interaktiv tizimda birlashtirish.
\end{tcolorbox}
\vspace{2mm}\hrule height 0.2pt \vspace{3mm}

% DB question_id: 87d70bdb-c290-4fbe-b913-366f2da0ce85
\needspace{5\baselineskip}
\noindent\textbf{\color{ink}622.}
\par
Quyidagi ta'rif qaysi atamaga tegishli? Videodagi bitta alohida tasvir.
\begin{enumerate}[label=\Alph*),leftmargin=9mm]
\item FPS
\item Video aniqligi
\item Kadr
\item Video bit tezligi
\end{enumerate}
\begin{tcolorbox}[enhanced,breakable,colback=teal!4,colframe=teal!50!black,boxrule=0.5pt,arc=1.1mm,left=3mm,right=3mm,top=1.7mm,bottom=1.7mm,before skip=2mm,after skip=2mm]
\textbf{\color{teal}To'g'ri javob:} \textbf{C) Kadr}
\par\smallskip
\textbf{Izoh:} Ta'rif Kadr tushunchasiga tegishli.
\end{tcolorbox}
\vspace{2mm}\hrule height 0.2pt \vspace{3mm}

% DB question_id: 8bba87ac-bae5-5ced-9b76-1c3662d0570a
\needspace{5\baselineskip}
\noindent\textbf{\color{ink}623.}
\par
Rang chuqurligi 3 bit bo'lsa, palitrada nazariy jihatdan ko'pi bilan nechta rang kodlanadi?
\begin{enumerate}[label=\Alph*),leftmargin=9mm]
\item 6
\item 3
\item 1
\item 4
\item 8
\end{enumerate}
\begin{tcolorbox}[enhanced,breakable,colback=teal!4,colframe=teal!50!black,boxrule=0.5pt,arc=1.1mm,left=3mm,right=3mm,top=1.7mm,bottom=1.7mm,before skip=2mm,after skip=2mm]
\textbf{\color{teal}To'g'ri javob:} \textbf{E) 8}
\par\smallskip
\textbf{Izoh:} n bit bilan 2\textasciicircum{}n holat kodlanadi. 2\textasciicircum{}3=8 rang.
\end{tcolorbox}
\vspace{2mm}\hrule height 0.2pt \vspace{3mm}

% DB question_id: 8db68e8d-6ab8-48ef-bada-e9ab0f64c752
\needspace{5\baselineskip}
\noindent\textbf{\color{ink}624.}
\par
"Video aniqligi" atamasiga eng aniq ta'rifni toping.
\begin{enumerate}[label=\Alph*),leftmargin=9mm]
\item har bir kadrning gorizontal va vertikal piksel o'lchami
\item video oqimining bir soniyada uzatadigan bitlari soni
\item media oqimini kodlaydigan va dekodlaydigan algoritm yoki dasturiy modul
\item video, audio, subtitr va metama'lumot oqimlarini bir faylda birlashtiradigan tuzilma
\end{enumerate}
\begin{tcolorbox}[enhanced,breakable,colback=teal!4,colframe=teal!50!black,boxrule=0.5pt,arc=1.1mm,left=3mm,right=3mm,top=1.7mm,bottom=1.7mm,before skip=2mm,after skip=2mm]
\textbf{\color{teal}To'g'ri javob:} \textbf{A) har bir kadrning gorizontal va vertikal piksel o'lchami}
\par\smallskip
\textbf{Izoh:} Video aniqligi - har bir kadrning gorizontal va vertikal piksel o'lchami.
\end{tcolorbox}
\vspace{2mm}\hrule height 0.2pt \vspace{3mm}

% DB question_id: 8f1fa8c4-3287-46a4-b4be-951b0acc6dd3
\needspace{5\baselineskip}
\noindent\textbf{\color{ink}625.}
\par
"Yo'qotishsiz siqish" atamasiga eng aniq ta'rifni toping.
\begin{enumerate}[label=\Alph*),leftmargin=9mm]
\item hajmni kamaytirish uchun ayrim ma'lumotni qaytmas tarzda olib tashlaydigan siqish
\item ketma-ket takrorlangan qiymatni qiymat va takrorlar soni bilan ifodalash usuli
\item dekodlanganda asl ma'lumotni bitma-bit tiklaydigan siqish
\item asl hajmning siqilgan hajmga yoki kamayish foiziga nisbati
\end{enumerate}
\begin{tcolorbox}[enhanced,breakable,colback=teal!4,colframe=teal!50!black,boxrule=0.5pt,arc=1.1mm,left=3mm,right=3mm,top=1.7mm,bottom=1.7mm,before skip=2mm,after skip=2mm]
\textbf{\color{teal}To'g'ri javob:} \textbf{C) dekodlanganda asl ma'lumotni bitma-bit tiklaydigan siqish}
\par\smallskip
\textbf{Izoh:} Yo'qotishsiz siqish - dekodlanganda asl ma'lumotni bitma-bit tiklaydigan siqish.
\end{tcolorbox}
\vspace{2mm}\hrule height 0.2pt \vspace{3mm}

% DB question_id: a3dab8c3-d074-400d-8a41-03519f86965a
\needspace{5\baselineskip}
\noindent\textbf{\color{ink}626.}
\par
Quyidagi ta'rif qaysi atamaga tegishli? Analog signalni raqamli namunalarga aylantiradigan qurilma yoki jarayon.
\begin{enumerate}[label=\Alph*),leftmargin=9mm]
\item Namuna chastotasi
\item Kvantlash
\item ADC
\item Bit chuqurligi
\end{enumerate}
\begin{tcolorbox}[enhanced,breakable,colback=teal!4,colframe=teal!50!black,boxrule=0.5pt,arc=1.1mm,left=3mm,right=3mm,top=1.7mm,bottom=1.7mm,before skip=2mm,after skip=2mm]
\textbf{\color{teal}To'g'ri javob:} \textbf{C) ADC}
\par\smallskip
\textbf{Izoh:} Ta'rif ADC tushunchasiga tegishli.
\end{tcolorbox}
\vspace{2mm}\hrule height 0.2pt \vspace{3mm}

% DB question_id: ae997b57-6b60-45f0-9405-9a4efd12884c
\needspace{5\baselineskip}
\noindent\textbf{\color{ink}627.}
\par
"Kod nuqtasi" atamasiga eng aniq ta'rifni toping.
\begin{enumerate}[label=\Alph*),leftmargin=9mm]
\item Unicode kod nuqtalarini 1-4 baytli ketma-ketliklarda kodlaydigan shakl
\item Unicode kod nuqtalarini 2 yoki 4 baytli birliklarda kodlaydigan shakl
\item faylning kodlash turi, tuzilishi yoki xizmat ma'lumotlarini tavsiflovchi qo'shimcha ma'lumot
\item Unicode tizimida belgiga berilgan abstrakt raqam
\end{enumerate}
\begin{tcolorbox}[enhanced,breakable,colback=teal!4,colframe=teal!50!black,boxrule=0.5pt,arc=1.1mm,left=3mm,right=3mm,top=1.7mm,bottom=1.7mm,before skip=2mm,after skip=2mm]
\textbf{\color{teal}To'g'ri javob:} \textbf{D) Unicode tizimida belgiga berilgan abstrakt raqam}
\par\smallskip
\textbf{Izoh:} Kod nuqtasi - Unicode tizimida belgiga berilgan abstrakt raqam.
\end{tcolorbox}
\vspace{2mm}\hrule height 0.2pt \vspace{3mm}

% DB question_id: b3e46574-10c0-42f4-8759-f8c5eb3eb4ff
\needspace{5\baselineskip}
\noindent\textbf{\color{ink}628.}
\par
"PCM" atamasiga eng aniq ta'rifni toping.
\begin{enumerate}[label=\Alph*),leftmargin=9mm]
\item bir soniyada uzatiladigan yoki saqlanadigan audio bitlari soni
\item vaqt va amplituda bo'yicha uzluksiz o'zgaradigan tovush signali
\item analog signalni raqamli namunalarga aylantiradigan qurilma yoki jarayon
\item namunalar qiymatini odatda siqilmagan holda ketma-ket saqlash usuli
\end{enumerate}
\begin{tcolorbox}[enhanced,breakable,colback=teal!4,colframe=teal!50!black,boxrule=0.5pt,arc=1.1mm,left=3mm,right=3mm,top=1.7mm,bottom=1.7mm,before skip=2mm,after skip=2mm]
\textbf{\color{teal}To'g'ri javob:} \textbf{D) namunalar qiymatini odatda siqilmagan holda ketma-ket saqlash usuli}
\par\smallskip
\textbf{Izoh:} PCM - namunalar qiymatini odatda siqilmagan holda ketma-ket saqlash usuli.
\end{tcolorbox}
\vspace{2mm}\hrule height 0.2pt \vspace{3mm}

% DB question_id: b4769dee-bf9a-4bb8-aad1-63bfff66aba3
\needspace{5\baselineskip}
\noindent\textbf{\color{ink}629.}
\par
"Analog tovush" atamasiga eng aniq ta'rifni toping.
\begin{enumerate}[label=\Alph*),leftmargin=9mm]
\item analog signalni raqamli namunalarga aylantiradigan qurilma yoki jarayon
\item vaqt va amplituda bo'yicha uzluksiz o'zgaradigan tovush signali
\item bir soniyada olingan audio namunalar soni
\item namuna amplitudasini ruxsat etilgan sonli darajalardan biriga moslash jarayoni
\end{enumerate}
\begin{tcolorbox}[enhanced,breakable,colback=teal!4,colframe=teal!50!black,boxrule=0.5pt,arc=1.1mm,left=3mm,right=3mm,top=1.7mm,bottom=1.7mm,before skip=2mm,after skip=2mm]
\textbf{\color{teal}To'g'ri javob:} \textbf{B) vaqt va amplituda bo'yicha uzluksiz o'zgaradigan tovush signali}
\par\smallskip
\textbf{Izoh:} Analog tovush - vaqt va amplituda bo'yicha uzluksiz o'zgaradigan tovush signali.
\end{tcolorbox}
\vspace{2mm}\hrule height 0.2pt \vspace{3mm}

% DB question_id: baea72fa-4dab-4d88-83c3-ece94b941685
\needspace{5\baselineskip}
\noindent\textbf{\color{ink}630.}
\par
"Belgini kodlash" atamasiga eng aniq ta'rifni toping.
\begin{enumerate}[label=\Alph*),leftmargin=9mm]
\item asl ko'rinishida 7 bitli va 128 kod holatiga ega belgilar standarti
\item 8 bitli kod sahifalariga beriladigan umumiy nom bo'lib, yagona universal jadval emas
\item turli yozuv tizimlari belgilariga yagona kod nuqtalari ajratadigan standart
\item matndagi har bir belgi yoki kod nuqtasini raqamli kod bilan ifodalash jarayoni
\end{enumerate}
\begin{tcolorbox}[enhanced,breakable,colback=teal!4,colframe=teal!50!black,boxrule=0.5pt,arc=1.1mm,left=3mm,right=3mm,top=1.7mm,bottom=1.7mm,before skip=2mm,after skip=2mm]
\textbf{\color{teal}To'g'ri javob:} \textbf{D) matndagi har bir belgi yoki kod nuqtasini raqamli kod bilan ifodalash jarayoni}
\par\smallskip
\textbf{Izoh:} Belgini kodlash - matndagi har bir belgi yoki kod nuqtasini raqamli kod bilan ifodalash jarayoni.
\end{tcolorbox}
\vspace{2mm}\hrule height 0.2pt \vspace{3mm}

% DB question_id: bc02fd46-cfa4-479f-9edd-d723b5f06448
\needspace{5\baselineskip}
\noindent\textbf{\color{ink}631.}
\par
Quyidagi ta'rif qaysi atamaga tegishli? Vaqt va amplituda bo'yicha uzluksiz o'zgaradigan tovush signali.
\begin{enumerate}[label=\Alph*),leftmargin=9mm]
\item ADC
\item Analog tovush
\item Namuna chastotasi
\item Kvantlash
\end{enumerate}
\begin{tcolorbox}[enhanced,breakable,colback=teal!4,colframe=teal!50!black,boxrule=0.5pt,arc=1.1mm,left=3mm,right=3mm,top=1.7mm,bottom=1.7mm,before skip=2mm,after skip=2mm]
\textbf{\color{teal}To'g'ri javob:} \textbf{B) Analog tovush}
\par\smallskip
\textbf{Izoh:} Ta'rif Analog tovush tushunchasiga tegishli.
\end{tcolorbox}
\vspace{2mm}\hrule height 0.2pt \vspace{3mm}

% DB question_id: c92422e5-6570-44c7-b850-49245de7c62b
\needspace{5\baselineskip}
\noindent\textbf{\color{ink}632.}
\par
"Fayl formati" atamasiga eng aniq ta'rifni toping.
\begin{enumerate}[label=\Alph*),leftmargin=9mm]
\item fayl nomining formatni taxmin qilishga yordam beradigan oxirgi qismi
\item axborotni kamroq bit yoki bayt bilan ifodalash jarayoni
\item dekodlanganda asl ma'lumotni bitma-bit tiklaydigan siqish
\item fayldagi ma'lumotning tuzilishi va uni talqin qilish qoidalari
\end{enumerate}
\begin{tcolorbox}[enhanced,breakable,colback=teal!4,colframe=teal!50!black,boxrule=0.5pt,arc=1.1mm,left=3mm,right=3mm,top=1.7mm,bottom=1.7mm,before skip=2mm,after skip=2mm]
\textbf{\color{teal}To'g'ri javob:} \textbf{D) fayldagi ma'lumotning tuzilishi va uni talqin qilish qoidalari}
\par\smallskip
\textbf{Izoh:} Fayl formati - fayldagi ma'lumotning tuzilishi va uni talqin qilish qoidalari.
\end{tcolorbox}
\vspace{2mm}\hrule height 0.2pt \vspace{3mm}

% DB question_id: cbf274ff-5bfa-4e87-bb66-78c9d67a1e0e
\needspace{5\baselineskip}
\noindent\textbf{\color{ink}633.}
\par
Quyidagi ta'rif qaysi atamaga tegishli? Asl ko'rinishida 7 bitli va 128 kod holatiga ega belgilar standarti.
\begin{enumerate}[label=\Alph*),leftmargin=9mm]
\item ASCII
\item Kengaytirilgan ASCII
\item Unicode
\item Kod nuqtasi
\end{enumerate}
\begin{tcolorbox}[enhanced,breakable,colback=teal!4,colframe=teal!50!black,boxrule=0.5pt,arc=1.1mm,left=3mm,right=3mm,top=1.7mm,bottom=1.7mm,before skip=2mm,after skip=2mm]
\textbf{\color{teal}To'g'ri javob:} \textbf{A) ASCII}
\par\smallskip
\textbf{Izoh:} Ta'rif ASCII tushunchasiga tegishli.
\end{tcolorbox}
\vspace{2mm}\hrule height 0.2pt \vspace{3mm}

% DB question_id: cc659189-a849-4928-ab70-603171ae182a
\needspace{5\baselineskip}
\noindent\textbf{\color{ink}634.}
\par
"Kadr" atamasiga eng aniq ta'rifni toping.
\begin{enumerate}[label=\Alph*),leftmargin=9mm]
\item bir soniyada ko'rsatiladigan kadrlar soni
\item har bir kadrning gorizontal va vertikal piksel o'lchami
\item videodagi bitta alohida tasvir
\item video oqimining bir soniyada uzatadigan bitlari soni
\end{enumerate}
\begin{tcolorbox}[enhanced,breakable,colback=teal!4,colframe=teal!50!black,boxrule=0.5pt,arc=1.1mm,left=3mm,right=3mm,top=1.7mm,bottom=1.7mm,before skip=2mm,after skip=2mm]
\textbf{\color{teal}To'g'ri javob:} \textbf{C) videodagi bitta alohida tasvir}
\par\smallskip
\textbf{Izoh:} Kadr - videodagi bitta alohida tasvir.
\end{tcolorbox}
\vspace{2mm}\hrule height 0.2pt \vspace{3mm}

% DB question_id: dc215a88-5959-49f0-9cbf-aae74085dafd
\needspace{5\baselineskip}
\noindent\textbf{\color{ink}635.}
\par
"Transkodlash" atamasiga eng aniq ta'rifni toping.
\begin{enumerate}[label=\Alph*),leftmargin=9mm]
\item ma'lumot yo'qolsa tiklash uchun alohida joyda saqlangan nusxa
\item fayldagi ma'lumotning tuzilishi va uni talqin qilish qoidalari
\item fayl nomining formatni taxmin qilishga yordam beradigan oxirgi qismi
\item media oqimini bir kodek yoki parametrdan boshqasiga qayta kodlash
\end{enumerate}
\begin{tcolorbox}[enhanced,breakable,colback=teal!4,colframe=teal!50!black,boxrule=0.5pt,arc=1.1mm,left=3mm,right=3mm,top=1.7mm,bottom=1.7mm,before skip=2mm,after skip=2mm]
\textbf{\color{teal}To'g'ri javob:} \textbf{D) media oqimini bir kodek yoki parametrdan boshqasiga qayta kodlash}
\par\smallskip
\textbf{Izoh:} Transkodlash - media oqimini bir kodek yoki parametrdan boshqasiga qayta kodlash.
\end{tcolorbox}
\vspace{2mm}\hrule height 0.2pt \vspace{3mm}

% DB question_id: ebdd8a07-df5f-4503-a14f-eca5f5cc5783
\needspace{5\baselineskip}
\noindent\textbf{\color{ink}636.}
\par
"RLE" atamasiga eng aniq ta'rifni toping.
\begin{enumerate}[label=\Alph*),leftmargin=9mm]
\item ketma-ket takrorlangan qiymatni qiymat va takrorlar soni bilan ifodalash usuli
\item asl hajmning siqilgan hajmga yoki kamayish foiziga nisbati
\item media oqimini kodlash va dekodlash qoidalari yoki dasturiy moduli
\item media oqimini bir kodek yoki parametrdan boshqasiga qayta kodlash
\end{enumerate}
\begin{tcolorbox}[enhanced,breakable,colback=teal!4,colframe=teal!50!black,boxrule=0.5pt,arc=1.1mm,left=3mm,right=3mm,top=1.7mm,bottom=1.7mm,before skip=2mm,after skip=2mm]
\textbf{\color{teal}To'g'ri javob:} \textbf{A) ketma-ket takrorlangan qiymatni qiymat va takrorlar soni bilan ifodalash usuli}
\par\smallskip
\textbf{Izoh:} RLE - ketma-ket takrorlangan qiymatni qiymat va takrorlar soni bilan ifodalash usuli.
\end{tcolorbox}
\vspace{2mm}\hrule height 0.2pt \vspace{3mm}

% DB question_id: f1fea3fa-5902-486e-bf22-ec0fb70aa0df
\needspace{5\baselineskip}
\noindent\textbf{\color{ink}637.}
\par
"Yo'qotishli siqish" atamasiga eng aniq ta'rifni toping.
\begin{enumerate}[label=\Alph*),leftmargin=9mm]
\item ketma-ket takrorlangan qiymatni qiymat va takrorlar soni bilan ifodalash usuli
\item asl hajmning siqilgan hajmga yoki kamayish foiziga nisbati
\item media oqimini kodlash va dekodlash qoidalari yoki dasturiy moduli
\item hajmni kamaytirish uchun ayrim ma'lumotni qaytmas tarzda olib tashlaydigan siqish
\end{enumerate}
\begin{tcolorbox}[enhanced,breakable,colback=teal!4,colframe=teal!50!black,boxrule=0.5pt,arc=1.1mm,left=3mm,right=3mm,top=1.7mm,bottom=1.7mm,before skip=2mm,after skip=2mm]
\textbf{\color{teal}To'g'ri javob:} \textbf{D) hajmni kamaytirish uchun ayrim ma'lumotni qaytmas tarzda olib tashlaydigan siqish}
\par\smallskip
\textbf{Izoh:} Yo'qotishli siqish - hajmni kamaytirish uchun ayrim ma'lumotni qaytmas tarzda olib tashlaydigan siqish.
\end{tcolorbox}
\vspace{2mm}\hrule height 0.2pt \vspace{3mm}

% DB question_id: fdceb853-4ec0-4e8e-9ff8-4d401a642882
\needspace{5\baselineskip}
\noindent\textbf{\color{ink}638.}
\par
"Kodek" tushunchasining "Videoaxborot va multimedia" bobidagi eng aniq ta'rifini toping.
\begin{enumerate}[label=\Alph*),leftmargin=9mm]
\item video, audio, subtitr va metama'lumot oqimlarini bir faylda birlashtiradigan tuzilma
\item vaqt bo'yicha ketma-ket uzatiladigan audio, video yoki boshqa media ma'lumoti
\item media oqimini kodlaydigan va dekodlaydigan algoritm yoki dasturiy modul
\item bir nechta media turini yagona mazmun yoki interaktiv tizimda birlashtirish
\end{enumerate}
\begin{tcolorbox}[enhanced,breakable,colback=teal!4,colframe=teal!50!black,boxrule=0.5pt,arc=1.1mm,left=3mm,right=3mm,top=1.7mm,bottom=1.7mm,before skip=2mm,after skip=2mm]
\textbf{\color{teal}To'g'ri javob:} \textbf{C) media oqimini kodlaydigan va dekodlaydigan algoritm yoki dasturiy modul}
\par\smallskip
\textbf{Izoh:} Kodek - media oqimini kodlaydigan va dekodlaydigan algoritm yoki dasturiy modul.
\end{tcolorbox}
\vspace{2mm}\hrule height 0.2pt \vspace{3mm}

% DB question_id: fec7f2de-7223-408d-9bbd-0ca63e9a07e1
\needspace{5\baselineskip}
\noindent\textbf{\color{ink}639.}
\par
"Kengaytirilgan ASCII" atamasiga eng aniq ta'rifni toping.
\begin{enumerate}[label=\Alph*),leftmargin=9mm]
\item turli yozuv tizimlari belgilariga yagona kod nuqtalari ajratadigan standart
\item 8 bitli kod sahifalariga beriladigan umumiy nom bo'lib, yagona universal jadval emas
\item Unicode tizimida belgiga berilgan abstrakt raqam
\item Unicode kod nuqtalarini 1-4 baytli ketma-ketliklarda kodlaydigan shakl
\end{enumerate}
\begin{tcolorbox}[enhanced,breakable,colback=teal!4,colframe=teal!50!black,boxrule=0.5pt,arc=1.1mm,left=3mm,right=3mm,top=1.7mm,bottom=1.7mm,before skip=2mm,after skip=2mm]
\textbf{\color{teal}To'g'ri javob:} \textbf{B) 8 bitli kod sahifalariga beriladigan umumiy nom bo'lib, yagona universal jadval emas}
\par\smallskip
\textbf{Izoh:} Kengaytirilgan ASCII - 8 bitli kod sahifalariga beriladigan umumiy nom bo'lib, yagona universal jadval emas.
\end{tcolorbox}
\vspace{2mm}\hrule height 0.2pt \vspace{3mm}

% DB question_id: 1c02c066-d385-48fc-ac81-04a849766e83
\needspace{5\baselineskip}
\noindent\textbf{\color{ink}640.}
\par
Kadr 1920 \ensuremath{\times} 1080 piksel sifatida berildi. Bu vaziyatdagi asosiy tushunchani aniqlang.
\begin{enumerate}[label=\Alph*),leftmargin=9mm]
\item Video aniqligi
\item Video bit tezligi
\item Kodek
\item Konteyner
\end{enumerate}
\begin{tcolorbox}[enhanced,breakable,colback=teal!4,colframe=teal!50!black,boxrule=0.5pt,arc=1.1mm,left=3mm,right=3mm,top=1.7mm,bottom=1.7mm,before skip=2mm,after skip=2mm]
\textbf{\color{teal}To'g'ri javob:} \textbf{A) Video aniqligi}
\par\smallskip
\textbf{Izoh:} Vaziyat Video aniqligiga misol: har bir kadrning gorizontal va vertikal piksel o'lchami.
\end{tcolorbox}
\vspace{2mm}\hrule height 0.2pt \vspace{3mm}

% DB question_id: 1de06fa2-e7ea-41f9-914a-478941184a75
\needspace{5\baselineskip}
\noindent\textbf{\color{ink}641.}
\par
Quyidagi ta'rif qaysi atamaga tegishli? Unicode kod nuqtalarini 2 yoki 4 baytli birliklarda kodlaydigan shakl.
\begin{enumerate}[label=\Alph*),leftmargin=9mm]
\item Matn metama'lumoti
\item UTF-16
\item Belgini kodlash
\item ASCII
\end{enumerate}
\begin{tcolorbox}[enhanced,breakable,colback=teal!4,colframe=teal!50!black,boxrule=0.5pt,arc=1.1mm,left=3mm,right=3mm,top=1.7mm,bottom=1.7mm,before skip=2mm,after skip=2mm]
\textbf{\color{teal}To'g'ri javob:} \textbf{B) UTF-16}
\par\smallskip
\textbf{Izoh:} Ta'rif UTF-16 tushunchasiga tegishli.
\end{tcolorbox}
\vspace{2mm}\hrule height 0.2pt \vspace{3mm}

% DB question_id: 1eeb468c-c342-48cf-a6a1-92dfed775f6a
\needspace{5\baselineskip}
\noindent\textbf{\color{ink}642.}
\par
Bir baytli 0-255 kodlar mahalliy belgilar uchun turlicha talqin qilindi. Bu vaziyatdagi asosiy tushunchani aniqlang.
\begin{enumerate}[label=\Alph*),leftmargin=9mm]
\item Unicode
\item Kengaytirilgan ASCII
\item Kod nuqtasi
\item UTF-8
\end{enumerate}
\begin{tcolorbox}[enhanced,breakable,colback=teal!4,colframe=teal!50!black,boxrule=0.5pt,arc=1.1mm,left=3mm,right=3mm,top=1.7mm,bottom=1.7mm,before skip=2mm,after skip=2mm]
\textbf{\color{teal}To'g'ri javob:} \textbf{B) Kengaytirilgan ASCII}
\par\smallskip
\textbf{Izoh:} Vaziyat Kengaytirilgan ASCIIga misol: 8 bitli kod sahifalariga beriladigan umumiy nom bo'lib, yagona universal jadval emas.
\end{tcolorbox}
\vspace{2mm}\hrule height 0.2pt \vspace{3mm}

% DB question_id: 2386d02d-e1ca-4b05-bc62-0b56421ff2e4
\needspace{5\baselineskip}
\noindent\textbf{\color{ink}643.}
\par
Quyidagi ta'rif qaysi atamaga tegishli? Har bir kadrning gorizontal va vertikal piksel o'lchami.
\begin{enumerate}[label=\Alph*),leftmargin=9mm]
\item Video aniqligi
\item Video bit tezligi
\item Kodek
\item Konteyner
\end{enumerate}
\begin{tcolorbox}[enhanced,breakable,colback=teal!4,colframe=teal!50!black,boxrule=0.5pt,arc=1.1mm,left=3mm,right=3mm,top=1.7mm,bottom=1.7mm,before skip=2mm,after skip=2mm]
\textbf{\color{teal}To'g'ri javob:} \textbf{A) Video aniqligi}
\par\smallskip
\textbf{Izoh:} Ta'rif Video aniqligi tushunchasiga tegishli.
\end{tcolorbox}
\vspace{2mm}\hrule height 0.2pt \vspace{3mm}

% DB question_id: 29238013-fbca-49fe-b673-9e9720ce791a
\needspace{5\baselineskip}
\noindent\textbf{\color{ink}644.}
\par
Quyidagi ta'rif qaysi atamaga tegishli? Hajmni kamaytirish uchun ayrim ma'lumotni qaytmas tarzda olib tashlaydigan siqish.
\begin{enumerate}[label=\Alph*),leftmargin=9mm]
\item RLE
\item Yo'qotishli siqish
\item Siqish koeffitsiyenti
\item Kodek
\end{enumerate}
\begin{tcolorbox}[enhanced,breakable,colback=teal!4,colframe=teal!50!black,boxrule=0.5pt,arc=1.1mm,left=3mm,right=3mm,top=1.7mm,bottom=1.7mm,before skip=2mm,after skip=2mm]
\textbf{\color{teal}To'g'ri javob:} \textbf{B) Yo'qotishli siqish}
\par\smallskip
\textbf{Izoh:} Ta'rif Yo'qotishli siqish tushunchasiga tegishli.
\end{tcolorbox}
\vspace{2mm}\hrule height 0.2pt \vspace{3mm}

% DB question_id: 29c2b725-ed9a-4b40-a827-17cfeec5e337
\needspace{5\baselineskip}
\noindent\textbf{\color{ink}645.}
\par
44,1 kHz yozuvda har soniyada 44 100 namuna olindi. Bu vaziyatdagi asosiy tushunchani aniqlang.
\begin{enumerate}[label=\Alph*),leftmargin=9mm]
\item Kvantlash
\item Bit chuqurligi
\item Kanal
\item Namuna chastotasi
\end{enumerate}
\begin{tcolorbox}[enhanced,breakable,colback=teal!4,colframe=teal!50!black,boxrule=0.5pt,arc=1.1mm,left=3mm,right=3mm,top=1.7mm,bottom=1.7mm,before skip=2mm,after skip=2mm]
\textbf{\color{teal}To'g'ri javob:} \textbf{D) Namuna chastotasi}
\par\smallskip
\textbf{Izoh:} Vaziyat Namuna chastotasiga misol: bir soniyada olingan audio namunalar soni.
\end{tcolorbox}
\vspace{2mm}\hrule height 0.2pt \vspace{3mm}

% DB question_id: 2c15fd0a-8475-49af-8a3f-2dc698c52630
\needspace{5\baselineskip}
\noindent\textbf{\color{ink}646.}
\par
Quyidagi ta'rif qaysi atamaga tegishli? Media oqimini kodlaydigan va dekodlaydigan algoritm yoki dasturiy modul.
\begin{enumerate}[label=\Alph*),leftmargin=9mm]
\item Konteyner
\item Media oqimi
\item Kodek
\item Multimedia
\end{enumerate}
\begin{tcolorbox}[enhanced,breakable,colback=teal!4,colframe=teal!50!black,boxrule=0.5pt,arc=1.1mm,left=3mm,right=3mm,top=1.7mm,bottom=1.7mm,before skip=2mm,after skip=2mm]
\textbf{\color{teal}To'g'ri javob:} \textbf{C) Kodek}
\par\smallskip
\textbf{Izoh:} Ta'rif Kodek tushunchasiga tegishli.
\end{tcolorbox}
\vspace{2mm}\hrule height 0.2pt \vspace{3mm}

% DB question_id: 2e92ea89-2e91-41a0-8141-a2946f622fd2
\needspace{5\baselineskip}
\noindent\textbf{\color{ink}647.}
\par
Quyidagi ta'rif qaysi atamaga tegishli? Fayl nomining formatni taxmin qilishga yordam beradigan oxirgi qismi.
\begin{enumerate}[label=\Alph*),leftmargin=9mm]
\item Siqish
\item Yo'qotishsiz siqish
\item Fayl kengaytmasi
\item Yo'qotishli siqish
\end{enumerate}
\begin{tcolorbox}[enhanced,breakable,colback=teal!4,colframe=teal!50!black,boxrule=0.5pt,arc=1.1mm,left=3mm,right=3mm,top=1.7mm,bottom=1.7mm,before skip=2mm,after skip=2mm]
\textbf{\color{teal}To'g'ri javob:} \textbf{C) Fayl kengaytmasi}
\par\smallskip
\textbf{Izoh:} Ta'rif Fayl kengaytmasi tushunchasiga tegishli.
\end{tcolorbox}
\vspace{2mm}\hrule height 0.2pt \vspace{3mm}

% DB question_id: 4a27c090-8d40-472f-b1e4-a1113be5237d
\needspace{5\baselineskip}
\noindent\textbf{\color{ink}648.}
\par
Quyidagi ta'rif qaysi atamaga tegishli? Axborotni kamroq bit yoki bayt bilan ifodalash jarayoni.
\begin{enumerate}[label=\Alph*),leftmargin=9mm]
\item Yo'qotishsiz siqish
\item Yo'qotishli siqish
\item RLE
\item Siqish
\end{enumerate}
\begin{tcolorbox}[enhanced,breakable,colback=teal!4,colframe=teal!50!black,boxrule=0.5pt,arc=1.1mm,left=3mm,right=3mm,top=1.7mm,bottom=1.7mm,before skip=2mm,after skip=2mm]
\textbf{\color{teal}To'g'ri javob:} \textbf{D) Siqish}
\par\smallskip
\textbf{Izoh:} Ta'rif Siqish tushunchasiga tegishli.
\end{tcolorbox}
\vspace{2mm}\hrule height 0.2pt \vspace{3mm}

% DB question_id: 51935cc7-aaf4-4d75-9584-3a1a529b0244
\needspace{5\baselineskip}
\noindent\textbf{\color{ink}649.}
\par
Oqim 8 Mbit/s tezlikda kodlandi. Bu vaziyatdagi asosiy tushunchani aniqlang.
\begin{enumerate}[label=\Alph*),leftmargin=9mm]
\item Kodek
\item Video bit tezligi
\item Konteyner
\item Media oqimi
\end{enumerate}
\begin{tcolorbox}[enhanced,breakable,colback=teal!4,colframe=teal!50!black,boxrule=0.5pt,arc=1.1mm,left=3mm,right=3mm,top=1.7mm,bottom=1.7mm,before skip=2mm,after skip=2mm]
\textbf{\color{teal}To'g'ri javob:} \textbf{B) Video bit tezligi}
\par\smallskip
\textbf{Izoh:} Vaziyat Video bit tezligiga misol: video oqimining bir soniyada uzatadigan bitlari soni.
\end{tcolorbox}
\vspace{2mm}\hrule height 0.2pt \vspace{3mm}

% DB question_id: 53220a6d-52e8-4bd8-840b-c77601ed7bf6
\needspace{5\baselineskip}
\noindent\textbf{\color{ink}650.}
\par
Quyidagi ta'rif qaysi atamaga tegishli? Media oqimini kodlash va dekodlash qoidalari yoki dasturiy moduli.
\begin{enumerate}[label=\Alph*),leftmargin=9mm]
\item Kodek
\item Transkodlash
\item Zaxira nusxasi
\item Fayl formati
\end{enumerate}
\begin{tcolorbox}[enhanced,breakable,colback=teal!4,colframe=teal!50!black,boxrule=0.5pt,arc=1.1mm,left=3mm,right=3mm,top=1.7mm,bottom=1.7mm,before skip=2mm,after skip=2mm]
\textbf{\color{teal}To'g'ri javob:} \textbf{A) Kodek}
\par\smallskip
\textbf{Izoh:} Ta'rif Kodek tushunchasiga tegishli.
\end{tcolorbox}
\vspace{2mm}\hrule height 0.2pt \vspace{3mm}

% DB question_id: 56145d50-9db3-4df1-9e3d-25dd4a8c3658
\needspace{5\baselineskip}
\noindent\textbf{\color{ink}651.}
\par
Quyidagi ta'rif qaysi atamaga tegishli? Unicode tizimida belgiga berilgan abstrakt raqam.
\begin{enumerate}[label=\Alph*),leftmargin=9mm]
\item UTF-8
\item UTF-16
\item Matn metama'lumoti
\item Kod nuqtasi
\end{enumerate}
\begin{tcolorbox}[enhanced,breakable,colback=teal!4,colframe=teal!50!black,boxrule=0.5pt,arc=1.1mm,left=3mm,right=3mm,top=1.7mm,bottom=1.7mm,before skip=2mm,after skip=2mm]
\textbf{\color{teal}To'g'ri javob:} \textbf{D) Kod nuqtasi}
\par\smallskip
\textbf{Izoh:} Ta'rif Kod nuqtasi tushunchasiga tegishli.
\end{tcolorbox}
\vspace{2mm}\hrule height 0.2pt \vspace{3mm}

% DB question_id: 5789d8f1-4c57-4b62-89d4-8d7d9e89d11b
\needspace{5\baselineskip}
\noindent\textbf{\color{ink}652.}
\par
Mikrofon membranasi havo tebranishiga uzluksiz ergashdi. Bu vaziyatdagi asosiy tushunchani aniqlang.
\begin{enumerate}[label=\Alph*),leftmargin=9mm]
\item ADC
\item Analog tovush
\item Namuna chastotasi
\item Kvantlash
\end{enumerate}
\begin{tcolorbox}[enhanced,breakable,colback=teal!4,colframe=teal!50!black,boxrule=0.5pt,arc=1.1mm,left=3mm,right=3mm,top=1.7mm,bottom=1.7mm,before skip=2mm,after skip=2mm]
\textbf{\color{teal}To'g'ri javob:} \textbf{B) Analog tovush}
\par\smallskip
\textbf{Izoh:} Vaziyat Analog tovushga misol: vaqt va amplituda bo'yicha uzluksiz o'zgaradigan tovush signali.
\end{tcolorbox}
\vspace{2mm}\hrule height 0.2pt \vspace{3mm}

% DB question_id: 5c3e0a6e-0cbb-419a-a9b8-9e20cf65c087
\needspace{5\baselineskip}
\noindent\textbf{\color{ink}653.}
\par
Quyidagi ta'rif qaysi atamaga tegishli? 8 bitli kod sahifalariga beriladigan umumiy nom bo'lib, yagona universal jadval emas.
\begin{enumerate}[label=\Alph*),leftmargin=9mm]
\item Unicode
\item Kengaytirilgan ASCII
\item Kod nuqtasi
\item UTF-8
\end{enumerate}
\begin{tcolorbox}[enhanced,breakable,colback=teal!4,colframe=teal!50!black,boxrule=0.5pt,arc=1.1mm,left=3mm,right=3mm,top=1.7mm,bottom=1.7mm,before skip=2mm,after skip=2mm]
\textbf{\color{teal}To'g'ri javob:} \textbf{B) Kengaytirilgan ASCII}
\par\smallskip
\textbf{Izoh:} Ta'rif Kengaytirilgan ASCII tushunchasiga tegishli.
\end{tcolorbox}
\vspace{2mm}\hrule height 0.2pt \vspace{3mm}

% DB question_id: 602c448f-d58f-465b-a6e7-f80ce021c665
\needspace{5\baselineskip}
\noindent\textbf{\color{ink}654.}
\par
Quyidagi ta'rif qaysi atamaga tegishli? Turli yozuv tizimlari belgilariga yagona kod nuqtalari ajratadigan standart.
\begin{enumerate}[label=\Alph*),leftmargin=9mm]
\item Kod nuqtasi
\item UTF-8
\item Unicode
\item UTF-16
\end{enumerate}
\begin{tcolorbox}[enhanced,breakable,colback=teal!4,colframe=teal!50!black,boxrule=0.5pt,arc=1.1mm,left=3mm,right=3mm,top=1.7mm,bottom=1.7mm,before skip=2mm,after skip=2mm]
\textbf{\color{teal}To'g'ri javob:} \textbf{C) Unicode}
\par\smallskip
\textbf{Izoh:} Ta'rif Unicode tushunchasiga tegishli.
\end{tcolorbox}
\vspace{2mm}\hrule height 0.2pt \vspace{3mm}

% DB question_id: 65252fb3-9f23-4061-87bc-6f6658b5d8f8
\needspace{5\baselineskip}
\noindent\textbf{\color{ink}655.}
\par
Quyidagi ta'rif qaysi atamaga tegishli? Mustaqil audio signal yo'li.
\begin{enumerate}[label=\Alph*),leftmargin=9mm]
\item PCM
\item Audio bit tezligi
\item Kanal
\item Analog tovush
\end{enumerate}
\begin{tcolorbox}[enhanced,breakable,colback=teal!4,colframe=teal!50!black,boxrule=0.5pt,arc=1.1mm,left=3mm,right=3mm,top=1.7mm,bottom=1.7mm,before skip=2mm,after skip=2mm]
\textbf{\color{teal}To'g'ri javob:} \textbf{C) Kanal}
\par\smallskip
\textbf{Izoh:} Ta'rif Kanal tushunchasiga tegishli.
\end{tcolorbox}
\vspace{2mm}\hrule height 0.2pt \vspace{3mm}

% DB question_id: 694c90da-8f6b-4d28-85de-6e777c17a2cf
\needspace{5\baselineskip}
\noindent\textbf{\color{ink}656.}
\par
Quyidagi ta'rif qaysi atamaga tegishli? Bir soniyada olingan audio namunalar soni.
\begin{enumerate}[label=\Alph*),leftmargin=9mm]
\item Kvantlash
\item Bit chuqurligi
\item Kanal
\item Namuna chastotasi
\end{enumerate}
\begin{tcolorbox}[enhanced,breakable,colback=teal!4,colframe=teal!50!black,boxrule=0.5pt,arc=1.1mm,left=3mm,right=3mm,top=1.7mm,bottom=1.7mm,before skip=2mm,after skip=2mm]
\textbf{\color{teal}To'g'ri javob:} \textbf{D) Namuna chastotasi}
\par\smallskip
\textbf{Izoh:} Ta'rif Namuna chastotasi tushunchasiga tegishli.
\end{tcolorbox}
\vspace{2mm}\hrule height 0.2pt \vspace{3mm}

% DB question_id: 6a675bc6-5496-4070-bf32-045ebf14ca7f
\needspace{5\baselineskip}
\noindent\textbf{\color{ink}657.}
\par
Quyidagi ta'rif qaysi atamaga tegishli? Faylning kodlash turi, tuzilishi yoki xizmat ma'lumotlarini tavsiflovchi qo'shimcha ma'lumot.
\begin{enumerate}[label=\Alph*),leftmargin=9mm]
\item Belgini kodlash
\item ASCII
\item Matn metama'lumoti
\item Kengaytirilgan ASCII
\end{enumerate}
\begin{tcolorbox}[enhanced,breakable,colback=teal!4,colframe=teal!50!black,boxrule=0.5pt,arc=1.1mm,left=3mm,right=3mm,top=1.7mm,bottom=1.7mm,before skip=2mm,after skip=2mm]
\textbf{\color{teal}To'g'ri javob:} \textbf{C) Matn metama'lumoti}
\par\smallskip
\textbf{Izoh:} Ta'rif Matn metama'lumoti tushunchasiga tegishli.
\end{tcolorbox}
\vspace{2mm}\hrule height 0.2pt \vspace{3mm}

% DB question_id: 6f5b1773-859a-42ee-be82-1ee7249f94b9
\needspace{5\baselineskip}
\noindent\textbf{\color{ink}658.}
\par
Quyidagi ta'rif qaysi atamaga tegishli? Ketma-ket takrorlangan qiymatni qiymat va takrorlar soni bilan ifodalash usuli.
\begin{enumerate}[label=\Alph*),leftmargin=9mm]
\item Siqish koeffitsiyenti
\item Kodek
\item RLE
\item Transkodlash
\end{enumerate}
\begin{tcolorbox}[enhanced,breakable,colback=teal!4,colframe=teal!50!black,boxrule=0.5pt,arc=1.1mm,left=3mm,right=3mm,top=1.7mm,bottom=1.7mm,before skip=2mm,after skip=2mm]
\textbf{\color{teal}To'g'ri javob:} \textbf{C) RLE}
\par\smallskip
\textbf{Izoh:} Ta'rif RLE tushunchasiga tegishli.
\end{tcolorbox}
\vspace{2mm}\hrule height 0.2pt \vspace{3mm}

% DB question_id: 71d251ab-014a-4995-a620-dc118b671532
\needspace{5\baselineskip}
\noindent\textbf{\color{ink}659.}
\par
Quyidagi ta'rif qaysi atamaga tegishli? Vaqt bo'yicha ketma-ket uzatiladigan audio, video yoki boshqa media ma'lumoti.
\begin{enumerate}[label=\Alph*),leftmargin=9mm]
\item Media oqimi
\item Multimedia
\item Kadr
\item FPS
\end{enumerate}
\begin{tcolorbox}[enhanced,breakable,colback=teal!4,colframe=teal!50!black,boxrule=0.5pt,arc=1.1mm,left=3mm,right=3mm,top=1.7mm,bottom=1.7mm,before skip=2mm,after skip=2mm]
\textbf{\color{teal}To'g'ri javob:} \textbf{A) Media oqimi}
\par\smallskip
\textbf{Izoh:} Ta'rif Media oqimi tushunchasiga tegishli.
\end{tcolorbox}
\vspace{2mm}\hrule height 0.2pt \vspace{3mm}

% DB question_id: 732568c0-a381-4984-b6af-a463e9c1058a
\needspace{5\baselineskip}
\noindent\textbf{\color{ink}660.}
\par
Quyidagi ta'rif qaysi atamaga tegishli? Video oqimining bir soniyada uzatadigan bitlari soni.
\begin{enumerate}[label=\Alph*),leftmargin=9mm]
\item Kodek
\item Video bit tezligi
\item Konteyner
\item Media oqimi
\end{enumerate}
\begin{tcolorbox}[enhanced,breakable,colback=teal!4,colframe=teal!50!black,boxrule=0.5pt,arc=1.1mm,left=3mm,right=3mm,top=1.7mm,bottom=1.7mm,before skip=2mm,after skip=2mm]
\textbf{\color{teal}To'g'ri javob:} \textbf{B) Video bit tezligi}
\par\smallskip
\textbf{Izoh:} Ta'rif Video bit tezligi tushunchasiga tegishli.
\end{tcolorbox}
\vspace{2mm}\hrule height 0.2pt \vspace{3mm}

% DB question_id: 7dabfa77-a8b5-4978-b781-f15f84a555b2
\needspace{5\baselineskip}
\noindent\textbf{\color{ink}661.}
\par
Quyidagi ta'rif qaysi atamaga tegishli? Namuna amplitudasini ruxsat etilgan sonli darajalardan biriga moslash jarayoni.
\begin{enumerate}[label=\Alph*),leftmargin=9mm]
\item Kvantlash
\item Bit chuqurligi
\item Kanal
\item PCM
\end{enumerate}
\begin{tcolorbox}[enhanced,breakable,colback=teal!4,colframe=teal!50!black,boxrule=0.5pt,arc=1.1mm,left=3mm,right=3mm,top=1.7mm,bottom=1.7mm,before skip=2mm,after skip=2mm]
\textbf{\color{teal}To'g'ri javob:} \textbf{A) Kvantlash}
\par\smallskip
\textbf{Izoh:} Ta'rif Kvantlash tushunchasiga tegishli.
\end{tcolorbox}
\vspace{2mm}\hrule height 0.2pt \vspace{3mm}

% DB question_id: 8c5194d4-d425-49fe-a7e6-230c1e222b8d
\needspace{5\baselineskip}
\noindent\textbf{\color{ink}662.}
\par
Quyidagi ta'rif qaysi atamaga tegishli? Fayldagi ma'lumotning tuzilishi va uni talqin qilish qoidalari.
\begin{enumerate}[label=\Alph*),leftmargin=9mm]
\item Fayl kengaytmasi
\item Fayl formati
\item Siqish
\item Yo'qotishsiz siqish
\end{enumerate}
\begin{tcolorbox}[enhanced,breakable,colback=teal!4,colframe=teal!50!black,boxrule=0.5pt,arc=1.1mm,left=3mm,right=3mm,top=1.7mm,bottom=1.7mm,before skip=2mm,after skip=2mm]
\textbf{\color{teal}To'g'ri javob:} \textbf{B) Fayl formati}
\par\smallskip
\textbf{Izoh:} Ta'rif Fayl formati tushunchasiga tegishli.
\end{tcolorbox}
\vspace{2mm}\hrule height 0.2pt \vspace{3mm}

% DB question_id: 989ade2b-c203-467a-a83c-f3af717d97bd
\needspace{5\baselineskip}
\noindent\textbf{\color{ink}663.}
\par
Videoning ma'lum vaqt nuqtasidagi tasviri muzlatib ko'rsatildi. Bu vaziyatdagi asosiy tushunchani aniqlang.
\begin{enumerate}[label=\Alph*),leftmargin=9mm]
\item FPS
\item Video aniqligi
\item Kadr
\item Video bit tezligi
\end{enumerate}
\begin{tcolorbox}[enhanced,breakable,colback=teal!4,colframe=teal!50!black,boxrule=0.5pt,arc=1.1mm,left=3mm,right=3mm,top=1.7mm,bottom=1.7mm,before skip=2mm,after skip=2mm]
\textbf{\color{teal}To'g'ri javob:} \textbf{C) Kadr}
\par\smallskip
\textbf{Izoh:} Vaziyat Kadrga misol: videodagi bitta alohida tasvir.
\end{tcolorbox}
\vspace{2mm}\hrule height 0.2pt \vspace{3mm}

% DB question_id: 9bc6db37-b120-45b2-b0a1-fdbc0fb3bca3
\needspace{5\baselineskip}
\noindent\textbf{\color{ink}664.}
\par
Klaviaturadan kiritilgan A harfi baytlar ketma-ketligiga aylantirildi. Bu vaziyatdagi asosiy tushunchani aniqlang.
\begin{enumerate}[label=\Alph*),leftmargin=9mm]
\item ASCII
\item Kengaytirilgan ASCII
\item Unicode
\item Belgini kodlash
\end{enumerate}
\begin{tcolorbox}[enhanced,breakable,colback=teal!4,colframe=teal!50!black,boxrule=0.5pt,arc=1.1mm,left=3mm,right=3mm,top=1.7mm,bottom=1.7mm,before skip=2mm,after skip=2mm]
\textbf{\color{teal}To'g'ri javob:} \textbf{D) Belgini kodlash}
\par\smallskip
\textbf{Izoh:} Vaziyat Belgini kodlashga misol: matndagi har bir belgi yoki kod nuqtasini raqamli kod bilan ifodalash jarayoni.
\end{tcolorbox}
\vspace{2mm}\hrule height 0.2pt \vspace{3mm}

% DB question_id: 9d4e6d63-92b0-4da2-817f-1475ec455f76
\needspace{5\baselineskip}
\noindent\textbf{\color{ink}665.}
\par
Lotin, kirill va boshqa yozuvlar bir umumiy kod fazosida ifodalandi. Bu vaziyatdagi asosiy tushunchani aniqlang.
\begin{enumerate}[label=\Alph*),leftmargin=9mm]
\item Kod nuqtasi
\item UTF-8
\item Unicode
\item UTF-16
\end{enumerate}
\begin{tcolorbox}[enhanced,breakable,colback=teal!4,colframe=teal!50!black,boxrule=0.5pt,arc=1.1mm,left=3mm,right=3mm,top=1.7mm,bottom=1.7mm,before skip=2mm,after skip=2mm]
\textbf{\color{teal}To'g'ri javob:} \textbf{C) Unicode}
\par\smallskip
\textbf{Izoh:} Vaziyat Unicodega misol: turli yozuv tizimlari belgilariga yagona kod nuqtalari ajratadigan standart.
\end{tcolorbox}
\vspace{2mm}\hrule height 0.2pt \vspace{3mm}

% DB question_id: ad92d0b3-d52a-4f35-81e1-331086365a1a
\needspace{5\baselineskip}
\noindent\textbf{\color{ink}666.}
\par
Quyidagi ta'rif qaysi atamaga tegishli? Bitta audio namunasini kodlash uchun ajratilgan bitlar soni.
\begin{enumerate}[label=\Alph*),leftmargin=9mm]
\item Kanal
\item Bit chuqurligi
\item PCM
\item Audio bit tezligi
\end{enumerate}
\begin{tcolorbox}[enhanced,breakable,colback=teal!4,colframe=teal!50!black,boxrule=0.5pt,arc=1.1mm,left=3mm,right=3mm,top=1.7mm,bottom=1.7mm,before skip=2mm,after skip=2mm]
\textbf{\color{teal}To'g'ri javob:} \textbf{B) Bit chuqurligi}
\par\smallskip
\textbf{Izoh:} Ta'rif Bit chuqurligi tushunchasiga tegishli.
\end{tcolorbox}
\vspace{2mm}\hrule height 0.2pt \vspace{3mm}

% DB question_id: b862a08e-5f5f-4099-b6f0-ae243b2fc42a
\needspace{5\baselineskip}
\noindent\textbf{\color{ink}667.}
\par
Lotin harflari va boshqaruv belgilarining 0-127 kodlari ishlatildi. Bu vaziyatdagi asosiy tushunchani aniqlang.
\begin{enumerate}[label=\Alph*),leftmargin=9mm]
\item ASCII
\item Kengaytirilgan ASCII
\item Unicode
\item Kod nuqtasi
\end{enumerate}
\begin{tcolorbox}[enhanced,breakable,colback=teal!4,colframe=teal!50!black,boxrule=0.5pt,arc=1.1mm,left=3mm,right=3mm,top=1.7mm,bottom=1.7mm,before skip=2mm,after skip=2mm]
\textbf{\color{teal}To'g'ri javob:} \textbf{A) ASCII}
\par\smallskip
\textbf{Izoh:} Vaziyat ASCIIga misol: asl ko'rinishida 7 bitli va 128 kod holatiga ega belgilar standarti.
\end{tcolorbox}
\vspace{2mm}\hrule height 0.2pt \vspace{3mm}

% DB question_id: c090fa3f-0ae8-49dd-a1ca-4bfdd3183ee8
\needspace{5\baselineskip}
\noindent\textbf{\color{ink}668.}
\par
Quyidagi ta'rif qaysi atamaga tegishli? Unicode kod nuqtalarini 1-4 baytli ketma-ketliklarda kodlaydigan shakl.
\begin{enumerate}[label=\Alph*),leftmargin=9mm]
\item UTF-8
\item UTF-16
\item Matn metama'lumoti
\item Belgini kodlash
\end{enumerate}
\begin{tcolorbox}[enhanced,breakable,colback=teal!4,colframe=teal!50!black,boxrule=0.5pt,arc=1.1mm,left=3mm,right=3mm,top=1.7mm,bottom=1.7mm,before skip=2mm,after skip=2mm]
\textbf{\color{teal}To'g'ri javob:} \textbf{A) UTF-8}
\par\smallskip
\textbf{Izoh:} Ta'rif UTF-8 tushunchasiga tegishli.
\end{tcolorbox}
\vspace{2mm}\hrule height 0.2pt \vspace{3mm}

% DB question_id: c2d5cd43-7270-4ede-89b7-bc2fd0b4d168
\needspace{5\baselineskip}
\noindent\textbf{\color{ink}669.}
\par
Quyidagi ta'rif qaysi atamaga tegishli? Bir soniyada uzatiladigan yoki saqlanadigan audio bitlari soni.
\begin{enumerate}[label=\Alph*),leftmargin=9mm]
\item Audio bit tezligi
\item Analog tovush
\item ADC
\item Namuna chastotasi
\end{enumerate}
\begin{tcolorbox}[enhanced,breakable,colback=teal!4,colframe=teal!50!black,boxrule=0.5pt,arc=1.1mm,left=3mm,right=3mm,top=1.7mm,bottom=1.7mm,before skip=2mm,after skip=2mm]
\textbf{\color{teal}To'g'ri javob:} \textbf{A) Audio bit tezligi}
\par\smallskip
\textbf{Izoh:} Ta'rif Audio bit tezligi tushunchasiga tegishli.
\end{tcolorbox}
\vspace{2mm}\hrule height 0.2pt \vspace{3mm}

% DB question_id: d1ba0317-872a-4dbd-8281-4a8184db38ff
\needspace{5\baselineskip}
\noindent\textbf{\color{ink}670.}
\par
Quyidagi ta'rif qaysi atamaga tegishli? Media oqimini bir kodek yoki parametrdan boshqasiga qayta kodlash.
\begin{enumerate}[label=\Alph*),leftmargin=9mm]
\item Zaxira nusxasi
\item Transkodlash
\item Fayl formati
\item Fayl kengaytmasi
\end{enumerate}
\begin{tcolorbox}[enhanced,breakable,colback=teal!4,colframe=teal!50!black,boxrule=0.5pt,arc=1.1mm,left=3mm,right=3mm,top=1.7mm,bottom=1.7mm,before skip=2mm,after skip=2mm]
\textbf{\color{teal}To'g'ri javob:} \textbf{B) Transkodlash}
\par\smallskip
\textbf{Izoh:} Ta'rif Transkodlash tushunchasiga tegishli.
\end{tcolorbox}
\vspace{2mm}\hrule height 0.2pt \vspace{3mm}

% DB question_id: d7f8d93c-e24f-4225-9eef-a292f210283e
\needspace{5\baselineskip}
\noindent\textbf{\color{ink}671.}
\par
Quyidagi ta'rif qaysi atamaga tegishli? Bir nechta media turini yagona mazmun yoki interaktiv tizimda birlashtirish.
\begin{enumerate}[label=\Alph*),leftmargin=9mm]
\item Kadr
\item Multimedia
\item FPS
\item Video aniqligi
\end{enumerate}
\begin{tcolorbox}[enhanced,breakable,colback=teal!4,colframe=teal!50!black,boxrule=0.5pt,arc=1.1mm,left=3mm,right=3mm,top=1.7mm,bottom=1.7mm,before skip=2mm,after skip=2mm]
\textbf{\color{teal}To'g'ri javob:} \textbf{B) Multimedia}
\par\smallskip
\textbf{Izoh:} Ta'rif Multimedia tushunchasiga tegishli.
\end{tcolorbox}
\vspace{2mm}\hrule height 0.2pt \vspace{3mm}

% DB question_id: dbf9b524-dd9e-48c5-926f-5921a4de57ad
\needspace{5\baselineskip}
\noindent\textbf{\color{ink}672.}
\par
30 FPS videoda har soniyada 30 alohida tasvir namoyish etildi. Bu vaziyatdagi asosiy tushunchani aniqlang.
\begin{enumerate}[label=\Alph*),leftmargin=9mm]
\item Video aniqligi
\item Video bit tezligi
\item Kodek
\item FPS
\end{enumerate}
\begin{tcolorbox}[enhanced,breakable,colback=teal!4,colframe=teal!50!black,boxrule=0.5pt,arc=1.1mm,left=3mm,right=3mm,top=1.7mm,bottom=1.7mm,before skip=2mm,after skip=2mm]
\textbf{\color{teal}To'g'ri javob:} \textbf{D) FPS}
\par\smallskip
\textbf{Izoh:} Vaziyat FPSga misol: bir soniyada ko'rsatiladigan kadrlar soni.
\end{tcolorbox}
\vspace{2mm}\hrule height 0.2pt \vspace{3mm}

% DB question_id: dde7a1d6-f535-4862-9c44-1b329b89a87b
\needspace{5\baselineskip}
\noindent\textbf{\color{ink}673.}
\par
Quyidagi ta'rif qaysi atamaga tegishli? Video, audio, subtitr va metama'lumot oqimlarini bir faylda birlashtiradigan tuzilma.
\begin{enumerate}[label=\Alph*),leftmargin=9mm]
\item Media oqimi
\item Multimedia
\item Kadr
\item Konteyner
\end{enumerate}
\begin{tcolorbox}[enhanced,breakable,colback=teal!4,colframe=teal!50!black,boxrule=0.5pt,arc=1.1mm,left=3mm,right=3mm,top=1.7mm,bottom=1.7mm,before skip=2mm,after skip=2mm]
\textbf{\color{teal}To'g'ri javob:} \textbf{D) Konteyner}
\par\smallskip
\textbf{Izoh:} Ta'rif Konteyner tushunchasiga tegishli.
\end{tcolorbox}
\vspace{2mm}\hrule height 0.2pt \vspace{3mm}

% DB question_id: dfddeb59-2278-4a37-9c90-6755978744c7
\needspace{5\baselineskip}
\noindent\textbf{\color{ink}674.}
\par
Quyidagi ta'rif qaysi atamaga tegishli? Dekodlanganda asl ma'lumotni bitma-bit tiklaydigan siqish.
\begin{enumerate}[label=\Alph*),leftmargin=9mm]
\item Yo'qotishsiz siqish
\item Yo'qotishli siqish
\item RLE
\item Siqish koeffitsiyenti
\end{enumerate}
\begin{tcolorbox}[enhanced,breakable,colback=teal!4,colframe=teal!50!black,boxrule=0.5pt,arc=1.1mm,left=3mm,right=3mm,top=1.7mm,bottom=1.7mm,before skip=2mm,after skip=2mm]
\textbf{\color{teal}To'g'ri javob:} \textbf{A) Yo'qotishsiz siqish}
\par\smallskip
\textbf{Izoh:} Ta'rif Yo'qotishsiz siqish tushunchasiga tegishli.
\end{tcolorbox}
\vspace{2mm}\hrule height 0.2pt \vspace{3mm}

% DB question_id: e684f2d9-520d-4be5-a37b-ba425ecac515
\needspace{5\baselineskip}
\noindent\textbf{\color{ink}675.}
\par
Mikrofon signali kompyuterga kirishda sonli qiymatlarga o'tkazildi. Bu vaziyatdagi asosiy tushunchani aniqlang.
\begin{enumerate}[label=\Alph*),leftmargin=9mm]
\item Namuna chastotasi
\item Kvantlash
\item ADC
\item Bit chuqurligi
\end{enumerate}
\begin{tcolorbox}[enhanced,breakable,colback=teal!4,colframe=teal!50!black,boxrule=0.5pt,arc=1.1mm,left=3mm,right=3mm,top=1.7mm,bottom=1.7mm,before skip=2mm,after skip=2mm]
\textbf{\color{teal}To'g'ri javob:} \textbf{C) ADC}
\par\smallskip
\textbf{Izoh:} Vaziyat ADCga misol: analog signalni raqamli namunalarga aylantiradigan qurilma yoki jarayon.
\end{tcolorbox}
\vspace{2mm}\hrule height 0.2pt \vspace{3mm}

% DB question_id: eb4660e1-8b26-475b-8a1b-829b75550aa4
\needspace{5\baselineskip}
\noindent\textbf{\color{ink}676.}
\par
Quyidagi ta'rif qaysi atamaga tegishli? Namunalar qiymatini odatda siqilmagan holda ketma-ket saqlash usuli.
\begin{enumerate}[label=\Alph*),leftmargin=9mm]
\item Audio bit tezligi
\item Analog tovush
\item ADC
\item PCM
\end{enumerate}
\begin{tcolorbox}[enhanced,breakable,colback=teal!4,colframe=teal!50!black,boxrule=0.5pt,arc=1.1mm,left=3mm,right=3mm,top=1.7mm,bottom=1.7mm,before skip=2mm,after skip=2mm]
\textbf{\color{teal}To'g'ri javob:} \textbf{D) PCM}
\par\smallskip
\textbf{Izoh:} Ta'rif PCM tushunchasiga tegishli.
\end{tcolorbox}
\vspace{2mm}\hrule height 0.2pt \vspace{3mm}

% DB question_id: f0c2b680-93fc-4bca-b610-87650baf2a06
\needspace{5\baselineskip}
\noindent\textbf{\color{ink}677.}
\par
O'lchangan amplituda eng yaqin raqamli darajaga yaxlitlandi. Bu vaziyatdagi asosiy tushunchani aniqlang.
\begin{enumerate}[label=\Alph*),leftmargin=9mm]
\item Kvantlash
\item Bit chuqurligi
\item Kanal
\item PCM
\end{enumerate}
\begin{tcolorbox}[enhanced,breakable,colback=teal!4,colframe=teal!50!black,boxrule=0.5pt,arc=1.1mm,left=3mm,right=3mm,top=1.7mm,bottom=1.7mm,before skip=2mm,after skip=2mm]
\textbf{\color{teal}To'g'ri javob:} \textbf{A) Kvantlash}
\par\smallskip
\textbf{Izoh:} Vaziyat Kvantlashga misol: namuna amplitudasini ruxsat etilgan sonli darajalardan biriga moslash jarayoni.
\end{tcolorbox}
\vspace{2mm}\hrule height 0.2pt \vspace{3mm}

% DB question_id: 0009e8fd-8746-5553-8284-1258c61d4cfe
\needspace{5\baselineskip}
\noindent\textbf{\color{ink}678.}
\par
Video 12 Mbit/s, audio 128 kbit/s, davomiylik 180 s. Umumiy media hajmi qancha? MB = 10\textsuperscript{6} B.
\begin{enumerate}[label=\Alph*),leftmargin=9mm]
\item 270 MB
\item 272,88 MB
\item 2,88 MB
\item 2183,04 MB
\end{enumerate}
\begin{tcolorbox}[enhanced,breakable,colback=teal!4,colframe=teal!50!black,boxrule=0.5pt,arc=1.1mm,left=3mm,right=3mm,top=1.7mm,bottom=1.7mm,before skip=2mm,after skip=2mm]
\textbf{\color{teal}To'g'ri javob:} \textbf{B) 272,88 MB}
\par\smallskip
\textbf{Izoh:} Jami bitrate = 12.128 Mbit/s; hajm = jami\ensuremath{\times}180/8 = 272,88 MB.
\end{tcolorbox}
\vspace{2mm}\hrule height 0.2pt \vspace{3mm}

% DB question_id: 00b9c574-1878-5e41-be95-7d8fd1d4f66e
\needspace{5\baselineskip}
\noindent\textbf{\color{ink}679.}
\par
Kamida 5 xil rangni kodlash uchun minimal rang chuqurligi qancha?
\begin{enumerate}[label=\Alph*),leftmargin=9mm]
\item 5 bit
\item 2 bit
\item 3 bit
\item 4 bit
\end{enumerate}
\begin{tcolorbox}[enhanced,breakable,colback=teal!4,colframe=teal!50!black,boxrule=0.5pt,arc=1.1mm,left=3mm,right=3mm,top=1.7mm,bottom=1.7mm,before skip=2mm,after skip=2mm]
\textbf{\color{teal}To'g'ri javob:} \textbf{C) 3 bit}
\par\smallskip
\textbf{Izoh:} 2\textsuperscript{n} \ensuremath{\geq} 5; minimal n = 3.
\end{tcolorbox}
\vspace{2mm}\hrule height 0.2pt \vspace{3mm}

% DB question_id: 0120aba4-f4b9-59fa-a0e8-08b57aa757d8
\needspace{5\baselineskip}
\noindent\textbf{\color{ink}680.}
\par
Rezolyutsiya nimani bildiradi?
\begin{enumerate}[label=\Alph*),leftmargin=9mm]
\item Tasvirning gorizontal va vertikal piksel o'lchamini yoki zichligini
\item Faqat fayl nomini
\item Audio davomiyligini
\item Kod jadvalini
\end{enumerate}
\begin{tcolorbox}[enhanced,breakable,colback=teal!4,colframe=teal!50!black,boxrule=0.5pt,arc=1.1mm,left=3mm,right=3mm,top=1.7mm,bottom=1.7mm,before skip=2mm,after skip=2mm]
\textbf{\color{teal}To'g'ri javob:} \textbf{A) Tasvirning gorizontal va vertikal piksel o'lchamini yoki zichligini}
\par\smallskip
\textbf{Izoh:} Rezolyutsiya tasvir tafsilotlari bilan bog'liq.
\end{tcolorbox}
\vspace{2mm}\hrule height 0.2pt \vspace{3mm}

% DB question_id: 027890c2-672b-5d61-98ea-063217071d5b
\needspace{5\baselineskip}
\noindent\textbf{\color{ink}681.}
\par
44100 Hz, 16 bit, 2 kanalli PCM audioning xom bitrate'i qancha?
\begin{enumerate}[label=\Alph*),leftmargin=9mm]
\item 705 600 kbit/s
\item 176 400 kbit/s
\item 1411,2 kbit/s
\item 2 822 400 kbit/s
\end{enumerate}
\begin{tcolorbox}[enhanced,breakable,colback=teal!4,colframe=teal!50!black,boxrule=0.5pt,arc=1.1mm,left=3mm,right=3mm,top=1.7mm,bottom=1.7mm,before skip=2mm,after skip=2mm]
\textbf{\color{teal}To'g'ri javob:} \textbf{C) 1411,2 kbit/s}
\par\smallskip
\textbf{Izoh:} 44100\ensuremath{\times}16\ensuremath{\times}2 = 1411200 bit/s = 1411,2 kbit/s.
\end{tcolorbox}
\vspace{2mm}\hrule height 0.2pt \vspace{3mm}

% DB question_id: 0400b7dd-120d-5f91-84f8-34c802c40b8f
\needspace{5\baselineskip}
\noindent\textbf{\color{ink}682.}
\par
4000\ensuremath{\times}3000 pikselli, 24 bit rang chuqurligidagi siqilmagan rastr tasvir hajmini MiB da toping.
\begin{enumerate}[label=\Alph*),leftmargin=9mm]
\item 12 000 000 MiB
\item 288 000 000 MiB
\item 37,332 MiB
\item 34,332 MiB
\end{enumerate}
\begin{tcolorbox}[enhanced,breakable,colback=teal!4,colframe=teal!50!black,boxrule=0.5pt,arc=1.1mm,left=3mm,right=3mm,top=1.7mm,bottom=1.7mm,before skip=2mm,after skip=2mm]
\textbf{\color{teal}To'g'ri javob:} \textbf{D) 34,332 MiB}
\par\smallskip
\textbf{Izoh:} Hajm = 4000\ensuremath{\times}3000\ensuremath{\times}24/8 = 36 000 000 B = 34,332 MiB.
\end{tcolorbox}
\vspace{2mm}\hrule height 0.2pt \vspace{3mm}

% DB question_id: 040c87c1-79fb-5364-b3ae-1a1a05fbcf7f
\needspace{5\baselineskip}
\noindent\textbf{\color{ink}683.}
\par
1280\ensuremath{\times}720, 24 bit, 30 FPS, 2 soniyalik siqilmagan video hajmini GiB da toping.
\begin{enumerate}[label=\Alph*),leftmargin=9mm]
\item 19,775 GiB
\item 0,154 GiB
\item 2,637 GiB
\item 4,635 GiB
\end{enumerate}
\begin{tcolorbox}[enhanced,breakable,colback=teal!4,colframe=teal!50!black,boxrule=0.5pt,arc=1.1mm,left=3mm,right=3mm,top=1.7mm,bottom=1.7mm,before skip=2mm,after skip=2mm]
\textbf{\color{teal}To'g'ri javob:} \textbf{B) 0,154 GiB}
\par\smallskip
\textbf{Izoh:} Hajm = W\ensuremath{\times}H\ensuremath{\times}d/8\ensuremath{\times}FPS\ensuremath{\times}t = 0,154 GiB.
\end{tcolorbox}
\vspace{2mm}\hrule height 0.2pt \vspace{3mm}

% DB question_id: 04edd2d6-30be-55c4-95b1-63d1fddb9b43
\needspace{5\baselineskip}
\noindent\textbf{\color{ink}684.}
\par
Matnda 800 ta belgi bor, har biri 1 bayt bilan kodlangan. Hajmni B da toping.
\begin{enumerate}[label=\Alph*),leftmargin=9mm]
\item 800 B
\item 6 400 B
\item 100 B
\item 803 B
\end{enumerate}
\begin{tcolorbox}[enhanced,breakable,colback=teal!4,colframe=teal!50!black,boxrule=0.5pt,arc=1.1mm,left=3mm,right=3mm,top=1.7mm,bottom=1.7mm,before skip=2mm,after skip=2mm]
\textbf{\color{teal}To'g'ri javob:} \textbf{A) 800 B}
\par\smallskip
\textbf{Izoh:} 800 \ensuremath{\times} 1 = 800 B = 800 B.
\end{tcolorbox}
\vspace{2mm}\hrule height 0.2pt \vspace{3mm}

% DB question_id: 05525246-b05b-591f-964b-dfcfc0d7cc2b
\needspace{5\baselineskip}
\noindent\textbf{\color{ink}685.}
\par
Tasvir kengligi 5 cm, zichligi 300 PPI. 1 dyuym = 2,54 cm. Piksel kengligi taxminan qancha?
\begin{enumerate}[label=\Alph*),leftmargin=9mm]
\item 591 piksel
\item 1 500 piksel
\item 300 piksel
\item 296 piksel
\end{enumerate}
\begin{tcolorbox}[enhanced,breakable,colback=teal!4,colframe=teal!50!black,boxrule=0.5pt,arc=1.1mm,left=3mm,right=3mm,top=1.7mm,bottom=1.7mm,before skip=2mm,after skip=2mm]
\textbf{\color{teal}To'g'ri javob:} \textbf{A) 591 piksel}
\par\smallskip
\textbf{Izoh:} 5/2,54\ensuremath{\times}300 \ensuremath{\approx} 591 piksel.
\end{tcolorbox}
\vspace{2mm}\hrule height 0.2pt \vspace{3mm}

% DB question_id: 05dd64d3-72e9-5c5a-b873-381827c5e4d8
\needspace{5\baselineskip}
\noindent\textbf{\color{ink}686.}
\par
Matnda 4096 ta bir baytli va 512 ta 2 baytli belgi bor. Umumiy hajm qancha?
\begin{enumerate}[label=\Alph*),leftmargin=9mm]
\item 4 608 KiB
\item 40 960 KiB
\item 640 KiB
\item 5 KiB
\end{enumerate}
\begin{tcolorbox}[enhanced,breakable,colback=teal!4,colframe=teal!50!black,boxrule=0.5pt,arc=1.1mm,left=3mm,right=3mm,top=1.7mm,bottom=1.7mm,before skip=2mm,after skip=2mm]
\textbf{\color{teal}To'g'ri javob:} \textbf{D) 5 KiB}
\par\smallskip
\textbf{Izoh:} 4096\ensuremath{\times}1 + 512\ensuremath{\times}2 = 5120 B = 5 KiB.
\end{tcolorbox}
\vspace{2mm}\hrule height 0.2pt \vspace{3mm}

% DB question_id: 0762586b-299f-4156-99eb-3c69c2500f38
\needspace{5\baselineskip}
\noindent\textbf{\color{ink}687.}
\par
Raqamli ovoz yozib olishda namunalash chastotasi 44,1 kHz deb ko'rsatilgan. Bu miqdor fizik jihatdan nimaning sonini bildiradi?
\begin{enumerate}[label=\Alph*),leftmargin=9mm]
\item Bir soniyada uzatiladigan audio paketlar soni
\item Bitta namunani saqlash uchun ajratilgan bitlar soni
\item Bir soniyada ijro etiladigan video kadrlar soni
\item Bir soniyada olingan audio namunalar soni
\end{enumerate}
\begin{tcolorbox}[enhanced,breakable,colback=teal!4,colframe=teal!50!black,boxrule=0.5pt,arc=1.1mm,left=3mm,right=3mm,top=1.7mm,bottom=1.7mm,before skip=2mm,after skip=2mm]
\textbf{\color{teal}To'g'ri javob:} \textbf{D) Bir soniyada olingan audio namunalar soni}
\par\smallskip
\textbf{Izoh:} Namunalash chastotasi bir soniyada analog tovush signalidan olingan namunalar sonini bildiradi. 44,1 kHz qiymati soniyasiga 44 100 namuna demakdir; audio bit chuqurligi esa boshqa o'lchamdir.
\end{tcolorbox}
\vspace{2mm}\hrule height 0.2pt \vspace{3mm}

% DB question_id: 08176775-daee-5446-807b-41ba0958e1f0
\needspace{5\baselineskip}
\noindent\textbf{\color{ink}688.}
\par
8000 Hz, 8 bit, 1 kanal, 10 soniyalik siqilmagan audio hajmini B da toping.
\begin{enumerate}[label=\Alph*),leftmargin=9mm]
\item 640 000 B
\item 80 000 B
\item 160 000 B
\item 8 000 B
\end{enumerate}
\begin{tcolorbox}[enhanced,breakable,colback=teal!4,colframe=teal!50!black,boxrule=0.5pt,arc=1.1mm,left=3mm,right=3mm,top=1.7mm,bottom=1.7mm,before skip=2mm,after skip=2mm]
\textbf{\color{teal}To'g'ri javob:} \textbf{B) 80 000 B}
\par\smallskip
\textbf{Izoh:} Hajm = 8000\ensuremath{\times}8\ensuremath{\times}1\ensuremath{\times}10/8 = 80 000 B = 80 000 B.
\end{tcolorbox}
\vspace{2mm}\hrule height 0.2pt \vspace{3mm}

% DB question_id: 0ab0fe9a-fbbf-5094-b563-7d0e609de155
\needspace{5\baselineskip}
\noindent\textbf{\color{ink}689.}
\par
Loyihada 120 ta 10 MB rasm va 8 ta 1040 MB audio bor. Umumiy hajm qancha?
\begin{enumerate}[label=\Alph*),leftmargin=9mm]
\item 128 MB
\item 9 520 MB
\item 1 050 MB
\item 76 160 MB
\end{enumerate}
\begin{tcolorbox}[enhanced,breakable,colback=teal!4,colframe=teal!50!black,boxrule=0.5pt,arc=1.1mm,left=3mm,right=3mm,top=1.7mm,bottom=1.7mm,before skip=2mm,after skip=2mm]
\textbf{\color{teal}To'g'ri javob:} \textbf{B) 9 520 MB}
\par\smallskip
\textbf{Izoh:} 120\ensuremath{\times}10 + 8\ensuremath{\times}1040 = 9520 MB.
\end{tcolorbox}
\vspace{2mm}\hrule height 0.2pt \vspace{3mm}

% DB question_id: 0c6a4655-9719-543a-99d2-3b39ebed4d0f
\needspace{5\baselineskip}
\noindent\textbf{\color{ink}690.}
\par
Kamida 257 xil rangni kodlash uchun minimal rang chuqurligi qancha?
\begin{enumerate}[label=\Alph*),leftmargin=9mm]
\item 257 bit
\item 9 bit
\item 8 bit
\item 10 bit
\end{enumerate}
\begin{tcolorbox}[enhanced,breakable,colback=teal!4,colframe=teal!50!black,boxrule=0.5pt,arc=1.1mm,left=3mm,right=3mm,top=1.7mm,bottom=1.7mm,before skip=2mm,after skip=2mm]
\textbf{\color{teal}To'g'ri javob:} \textbf{B) 9 bit}
\par\smallskip
\textbf{Izoh:} 2\textsuperscript{n} \ensuremath{\geq} 257; minimal n = 9.
\end{tcolorbox}
\vspace{2mm}\hrule height 0.2pt \vspace{3mm}

% DB question_id: 0cb7e589-50f3-59ce-9afe-fb101190f147
\needspace{5\baselineskip}
\noindent\textbf{\color{ink}691.}
\par
1 bit rang chuqurligi nazariy jihatdan nechta rangni ifodalaydi?
\begin{enumerate}[label=\Alph*),leftmargin=9mm]
\item 1 rang
\item 4 rang
\item 5 rang
\item 2 rang
\end{enumerate}
\begin{tcolorbox}[enhanced,breakable,colback=teal!4,colframe=teal!50!black,boxrule=0.5pt,arc=1.1mm,left=3mm,right=3mm,top=1.7mm,bottom=1.7mm,before skip=2mm,after skip=2mm]
\textbf{\color{teal}To'g'ri javob:} \textbf{D) 2 rang}
\par\smallskip
\textbf{Izoh:} Ranglar soni = 2\textasciicircum{}1 = 2.
\end{tcolorbox}
\vspace{2mm}\hrule height 0.2pt \vspace{3mm}

% DB question_id: 0e6c7ca1-a000-5916-bd90-f6be5fe920f4
\needspace{5\baselineskip}
\noindent\textbf{\color{ink}692.}
\par
10 bit rang chuqurligi nazariy jihatdan nechta rangni ifodalaydi?
\begin{enumerate}[label=\Alph*),leftmargin=9mm]
\item 10 rang
\item 1 024 rang
\item 20 rang
\item 512 rang
\end{enumerate}
\begin{tcolorbox}[enhanced,breakable,colback=teal!4,colframe=teal!50!black,boxrule=0.5pt,arc=1.1mm,left=3mm,right=3mm,top=1.7mm,bottom=1.7mm,before skip=2mm,after skip=2mm]
\textbf{\color{teal}To'g'ri javob:} \textbf{B) 1 024 rang}
\par\smallskip
\textbf{Izoh:} Ranglar soni = 2\textasciicircum{}10 = 1 024.
\end{tcolorbox}
\vspace{2mm}\hrule height 0.2pt \vspace{3mm}

% DB question_id: 0ec942c7-da7d-4889-835b-218575103c80
\needspace{5\baselineskip}
\noindent\textbf{\color{ink}693.}
\par
Fayl formati nima?
\begin{enumerate}[label=\Alph*),leftmargin=9mm]
\item Fayl kengaytmasini ko'rsatib, uni ma'lum dastur bilan bog'lash qoidasi
\item Faylni tashuvchi qurilmada joylashtirish uchun ajratilgan sektorlar ro'yxati
\item Fayl mazmunini internet orqali uzatish uchun ishlatiladigan tarmoq protokoli
\item Ma'lumotning faylda qanday tuzilishi va talqin qilinishini belgilovchi standart
\end{enumerate}
\begin{tcolorbox}[enhanced,breakable,colback=teal!4,colframe=teal!50!black,boxrule=0.5pt,arc=1.1mm,left=3mm,right=3mm,top=1.7mm,bottom=1.7mm,before skip=2mm,after skip=2mm]
\textbf{\color{teal}To'g'ri javob:} \textbf{D) Ma'lumotning faylda qanday tuzilishi va talqin qilinishini belgilovchi standart}
\par\smallskip
\textbf{Izoh:} Fayl formati ma'lumotlar baytlari qanday tuzilishi va dastur tomonidan qanday talqin qilinishini belgilaydi. Kengaytma ko'pincha formatga ishora qiladi, lekin uning to'liq texnik ta'rifi emas; disk sektorlarining joylashuvi va uzatish protokoli ham boshqa tushunchalar.
\end{tcolorbox}
\vspace{2mm}\hrule height 0.2pt \vspace{3mm}

% DB question_id: 14203fdc-f623-53ba-8cb9-c9154b9503ee
\needspace{5\baselineskip}
\noindent\textbf{\color{ink}694.}
\par
Stereo audio nimani anglatadi?
\begin{enumerate}[label=\Alph*),leftmargin=9mm]
\item Ikki mustaqil audio kanal
\item Bitta kanal
\item To'rtta kanal
\item Siqilmagan video
\end{enumerate}
\begin{tcolorbox}[enhanced,breakable,colback=teal!4,colframe=teal!50!black,boxrule=0.5pt,arc=1.1mm,left=3mm,right=3mm,top=1.7mm,bottom=1.7mm,before skip=2mm,after skip=2mm]
\textbf{\color{teal}To'g'ri javob:} \textbf{A) Ikki mustaqil audio kanal}
\par\smallskip
\textbf{Izoh:} Stereo chap va o'ng kanaldan iborat.
\end{tcolorbox}
\vspace{2mm}\hrule height 0.2pt \vspace{3mm}

% DB question_id: 175bf11a-917b-5282-8458-0d1d6046694e
\needspace{5\baselineskip}
\noindent\textbf{\color{ink}695.}
\par
800\ensuremath{\times}600 pikselli, 16 bit rang chuqurligidagi siqilmagan rastr tasvir hajmini KiB da toping.
\begin{enumerate}[label=\Alph*),leftmargin=9mm]
\item 480 000 KiB
\item 7 680 000 KiB
\item 937,5 KiB
\item 940,5 KiB
\end{enumerate}
\begin{tcolorbox}[enhanced,breakable,colback=teal!4,colframe=teal!50!black,boxrule=0.5pt,arc=1.1mm,left=3mm,right=3mm,top=1.7mm,bottom=1.7mm,before skip=2mm,after skip=2mm]
\textbf{\color{teal}To'g'ri javob:} \textbf{C) 937,5 KiB}
\par\smallskip
\textbf{Izoh:} Hajm = 800\ensuremath{\times}600\ensuremath{\times}16/8 = 960 000 B = 937,5 KiB.
\end{tcolorbox}
\vspace{2mm}\hrule height 0.2pt \vspace{3mm}

% DB question_id: 17802a85-c3d0-4d08-bcdb-91123d0cc931
\needspace{5\baselineskip}
\noindent\textbf{\color{ink}696.}
\par
Dastur baytlarni PNG qoidalari asosida rasm sifatida o'qidi. Bu vaziyatdagi asosiy tushunchani aniqlang.
\begin{enumerate}[label=\Alph*),leftmargin=9mm]
\item Fayl kengaytmasi
\item Siqish
\item Yo'qotishsiz siqish
\item Fayl formati
\end{enumerate}
\begin{tcolorbox}[enhanced,breakable,colback=teal!4,colframe=teal!50!black,boxrule=0.5pt,arc=1.1mm,left=3mm,right=3mm,top=1.7mm,bottom=1.7mm,before skip=2mm,after skip=2mm]
\textbf{\color{teal}To'g'ri javob:} \textbf{D) Fayl formati}
\par\smallskip
\textbf{Izoh:} Vaziyat Fayl formatiga misol: fayldagi ma'lumotning tuzilishi va uni talqin qilish qoidalari.
\end{tcolorbox}
\vspace{2mm}\hrule height 0.2pt \vspace{3mm}

% DB question_id: 1874fd8a-8a90-585f-9843-36c1fffabd81
\needspace{5\baselineskip}
\noindent\textbf{\color{ink}697.}
\par
128 kbit/s bitratedagi audio 180 soniya davom etadi. Hajm qancha? MB = 10\textsuperscript{6} B.
\begin{enumerate}[label=\Alph*),leftmargin=9mm]
\item 23,04 MB
\item 2,88 MB
\item 0,36 MB
\item 5,88 MB
\end{enumerate}
\begin{tcolorbox}[enhanced,breakable,colback=teal!4,colframe=teal!50!black,boxrule=0.5pt,arc=1.1mm,left=3mm,right=3mm,top=1.7mm,bottom=1.7mm,before skip=2mm,after skip=2mm]
\textbf{\color{teal}To'g'ri javob:} \textbf{B) 2,88 MB}
\par\smallskip
\textbf{Izoh:} 128000\ensuremath{\times}180/8 = 2,88 MB.
\end{tcolorbox}
\vspace{2mm}\hrule height 0.2pt \vspace{3mm}

% DB question_id: 18fbd1c0-132a-5ffa-bf13-9ddf8e607c63
\needspace{5\baselineskip}
\noindent\textbf{\color{ink}698.}
\par
Video kadr nima?
\begin{enumerate}[label=\Alph*),leftmargin=9mm]
\item Audio namuna
\item Videodagi bitta alohida tasvir
\item Kod nuqtasi
\item Fayl papkasi
\end{enumerate}
\begin{tcolorbox}[enhanced,breakable,colback=teal!4,colframe=teal!50!black,boxrule=0.5pt,arc=1.1mm,left=3mm,right=3mm,top=1.7mm,bottom=1.7mm,before skip=2mm,after skip=2mm]
\textbf{\color{teal}To'g'ri javob:} \textbf{B) Videodagi bitta alohida tasvir}
\par\smallskip
\textbf{Izoh:} Kadrlar tez almashib harakat taassurotini beradi.
\end{tcolorbox}
\vspace{2mm}\hrule height 0.2pt \vspace{3mm}

% DB question_id: 1960762b-e9c9-5470-9d1f-2cb8884bb013
\needspace{5\baselineskip}
\noindent\textbf{\color{ink}699.}
\par
25 FPS videoda 60 soniyada nechta kadr bo'ladi?
\begin{enumerate}[label=\Alph*),leftmargin=9mm]
\item 85 kadr
\item 90 000 kadr
\item 1 500 kadr
\item 750 kadr
\end{enumerate}
\begin{tcolorbox}[enhanced,breakable,colback=teal!4,colframe=teal!50!black,boxrule=0.5pt,arc=1.1mm,left=3mm,right=3mm,top=1.7mm,bottom=1.7mm,before skip=2mm,after skip=2mm]
\textbf{\color{teal}To'g'ri javob:} \textbf{C) 1 500 kadr}
\par\smallskip
\textbf{Izoh:} Kadrlar soni = 25\ensuremath{\times}60 = 1500.
\end{tcolorbox}
\vspace{2mm}\hrule height 0.2pt \vspace{3mm}

% DB question_id: 1a292f6c-484d-5463-832e-bf54cb64e9dc
\needspace{5\baselineskip}
\noindent\textbf{\color{ink}700.}
\par
Tasvir kengligi 10 cm, zichligi 254 PPI. 1 dyuym = 2,54 cm. Piksel kengligi taxminan qancha?
\begin{enumerate}[label=\Alph*),leftmargin=9mm]
\item 2 540 piksel
\item 1 000 piksel
\item 254 piksel
\item 500 piksel
\end{enumerate}
\begin{tcolorbox}[enhanced,breakable,colback=teal!4,colframe=teal!50!black,boxrule=0.5pt,arc=1.1mm,left=3mm,right=3mm,top=1.7mm,bottom=1.7mm,before skip=2mm,after skip=2mm]
\textbf{\color{teal}To'g'ri javob:} \textbf{B) 1 000 piksel}
\par\smallskip
\textbf{Izoh:} 10/2,54\ensuremath{\times}254 \ensuremath{\approx} 1000 piksel.
\end{tcolorbox}
\vspace{2mm}\hrule height 0.2pt \vspace{3mm}

% DB question_id: 242d3a5f-31ff-55dc-90e5-7908da0079fb
\needspace{5\baselineskip}
\noindent\textbf{\color{ink}701.}
\par
Video 4 Mbit/s, audio 96 kbit/s, davomiylik 900 s. Umumiy media hajmi qancha? MB = 10\textsuperscript{6} B.
\begin{enumerate}[label=\Alph*),leftmargin=9mm]
\item 450 MB
\item 460,8 MB
\item 10,8 MB
\item 3686,4 MB
\end{enumerate}
\begin{tcolorbox}[enhanced,breakable,colback=teal!4,colframe=teal!50!black,boxrule=0.5pt,arc=1.1mm,left=3mm,right=3mm,top=1.7mm,bottom=1.7mm,before skip=2mm,after skip=2mm]
\textbf{\color{teal}To'g'ri javob:} \textbf{B) 460,8 MB}
\par\smallskip
\textbf{Izoh:} Jami bitrate = 4.096 Mbit/s; hajm = jami\ensuremath{\times}900/8 = 460,8 MB.
\end{tcolorbox}
\vspace{2mm}\hrule height 0.2pt \vspace{3mm}

% DB question_id: 2c978642-7010-5e39-a49b-9472635dffcf
\needspace{5\baselineskip}
\noindent\textbf{\color{ink}702.}
\par
Video bitrate'i 8 Mbit/s, davomiyligi 120 s. Hajm qancha? MB = 10\textsuperscript{6} B.
\begin{enumerate}[label=\Alph*),leftmargin=9mm]
\item 960 MB
\item 120 MB
\item 15 MB
\item 123 MB
\end{enumerate}
\begin{tcolorbox}[enhanced,breakable,colback=teal!4,colframe=teal!50!black,boxrule=0.5pt,arc=1.1mm,left=3mm,right=3mm,top=1.7mm,bottom=1.7mm,before skip=2mm,after skip=2mm]
\textbf{\color{teal}To'g'ri javob:} \textbf{B) 120 MB}
\par\smallskip
\textbf{Izoh:} 8\ensuremath{\times}120/8 = 120 MB.
\end{tcolorbox}
\vspace{2mm}\hrule height 0.2pt \vspace{3mm}

% DB question_id: 3196973f-5a1f-53f3-8e98-11f57aa6ca20
\needspace{5\baselineskip}
\noindent\textbf{\color{ink}703.}
\par
192 kbit/s bitratedagi audio 300 soniya davom etadi. Hajm qancha? MB = 10\textsuperscript{6} B.
\begin{enumerate}[label=\Alph*),leftmargin=9mm]
\item 57,6 MB
\item 0,9 MB
\item 7,2 MB
\item 10,2 MB
\end{enumerate}
\begin{tcolorbox}[enhanced,breakable,colback=teal!4,colframe=teal!50!black,boxrule=0.5pt,arc=1.1mm,left=3mm,right=3mm,top=1.7mm,bottom=1.7mm,before skip=2mm,after skip=2mm]
\textbf{\color{teal}To'g'ri javob:} \textbf{C) 7,2 MB}
\par\smallskip
\textbf{Izoh:} 192000\ensuremath{\times}300/8 = 7,2 MB.
\end{tcolorbox}
\vspace{2mm}\hrule height 0.2pt \vspace{3mm}

% DB question_id: 3387d6a4-010c-5024-8b98-08bb030a11fd
\needspace{5\baselineskip}
\noindent\textbf{\color{ink}704.}
\par
Video 25 Mbit/s, audio 192 kbit/s, davomiylik 120 s. Umumiy media hajmi qancha? MB = 10\textsuperscript{6} B.
\begin{enumerate}[label=\Alph*),leftmargin=9mm]
\item 377,88 MB
\item 375 MB
\item 2,88 MB
\item 3023,04 MB
\end{enumerate}
\begin{tcolorbox}[enhanced,breakable,colback=teal!4,colframe=teal!50!black,boxrule=0.5pt,arc=1.1mm,left=3mm,right=3mm,top=1.7mm,bottom=1.7mm,before skip=2mm,after skip=2mm]
\textbf{\color{teal}To'g'ri javob:} \textbf{A) 377,88 MB}
\par\smallskip
\textbf{Izoh:} Jami bitrate = 25.192 Mbit/s; hajm = jami\ensuremath{\times}120/8 = 377,88 MB.
\end{tcolorbox}
\vspace{2mm}\hrule height 0.2pt \vspace{3mm}

% DB question_id: 34633a45-36de-51ed-a89e-de8eafb177c1
\needspace{5\baselineskip}
\noindent\textbf{\color{ink}705.}
\par
16000 Hz, 16 bit, 1 kanal, 45 soniyalik siqilmagan audio hajmini MiB da toping.
\begin{enumerate}[label=\Alph*),leftmargin=9mm]
\item 11 520 000 MiB
\item 2 880 000 MiB
\item 32 000 MiB
\item 1,373 MiB
\end{enumerate}
\begin{tcolorbox}[enhanced,breakable,colback=teal!4,colframe=teal!50!black,boxrule=0.5pt,arc=1.1mm,left=3mm,right=3mm,top=1.7mm,bottom=1.7mm,before skip=2mm,after skip=2mm]
\textbf{\color{teal}To'g'ri javob:} \textbf{D) 1,373 MiB}
\par\smallskip
\textbf{Izoh:} Hajm = 16000\ensuremath{\times}16\ensuremath{\times}1\ensuremath{\times}45/8 = 1 440 000 B = 1,373 MiB.
\end{tcolorbox}
\vspace{2mm}\hrule height 0.2pt \vspace{3mm}

% DB question_id: 359a7d68-a15f-57f0-8f71-1ba7abb70590
\needspace{5\baselineskip}
\noindent\textbf{\color{ink}706.}
\par
Loyihada 250 ta 12 MB rasm va 3 ta 786 MB audio bor. Umumiy hajm qancha?
\begin{enumerate}[label=\Alph*),leftmargin=9mm]
\item 253 MB
\item 5 358 MB
\item 798 MB
\item 42 864 MB
\end{enumerate}
\begin{tcolorbox}[enhanced,breakable,colback=teal!4,colframe=teal!50!black,boxrule=0.5pt,arc=1.1mm,left=3mm,right=3mm,top=1.7mm,bottom=1.7mm,before skip=2mm,after skip=2mm]
\textbf{\color{teal}To'g'ri javob:} \textbf{B) 5 358 MB}
\par\smallskip
\textbf{Izoh:} 250\ensuremath{\times}12 + 3\ensuremath{\times}786 = 5358 MB.
\end{tcolorbox}
\vspace{2mm}\hrule height 0.2pt \vspace{3mm}

% DB question_id: 36789265-34b5-5214-b45e-9f9e4de1cea1
\needspace{5\baselineskip}
\noindent\textbf{\color{ink}707.}
\par
Kamida 3 xil rangni kodlash uchun minimal rang chuqurligi qancha?
\begin{enumerate}[label=\Alph*),leftmargin=9mm]
\item 3 bit
\item 2 bit
\item 1 bit
\item 5 bit
\end{enumerate}
\begin{tcolorbox}[enhanced,breakable,colback=teal!4,colframe=teal!50!black,boxrule=0.5pt,arc=1.1mm,left=3mm,right=3mm,top=1.7mm,bottom=1.7mm,before skip=2mm,after skip=2mm]
\textbf{\color{teal}To'g'ri javob:} \textbf{B) 2 bit}
\par\smallskip
\textbf{Izoh:} 2\textsuperscript{n} \ensuremath{\geq} 3; minimal n = 2.
\end{tcolorbox}
\vspace{2mm}\hrule height 0.2pt \vspace{3mm}

% DB question_id: 38eaa923-5292-489a-9c88-c039126470c4
\needspace{5\baselineskip}
\noindent\textbf{\color{ink}708.}
\par
Raqamli audio yozish tizimida namunalash chastotasi o'zgarmadi, lekin har bir amplituda namunasining kodlash aniqligi 16 bitdan 24 bitga oshirildi. Audio bit chuqurligi nimani bildiradi?
\begin{enumerate}[label=\Alph*),leftmargin=9mm]
\item Bitta namuna amplitudasini kodlash uchun bitlar soni
\item Bir soniyada analog signaldan olinadigan audio namunalar soni
\item Bir soniyada uzatilgan raqamli audio oqimining umumiy hajmi
\item Bitta ovoz yozuvida ishlatiladigan mustaqil kanallarning umumiy soni
\end{enumerate}
\begin{tcolorbox}[enhanced,breakable,colback=teal!4,colframe=teal!50!black,boxrule=0.5pt,arc=1.1mm,left=3mm,right=3mm,top=1.7mm,bottom=1.7mm,before skip=2mm,after skip=2mm]
\textbf{\color{teal}To'g'ri javob:} \textbf{A) Bitta namuna amplitudasini kodlash uchun bitlar soni}
\par\smallskip
\textbf{Izoh:} Audio bit chuqurligi bitta amplituda namunasi uchun ajratilgan bitlar soni. Namunalash chastotasi soniyadagi namunalar, bitrate esa vaqt birligidagi bitlar miqdoridir.
\end{tcolorbox}
\vspace{2mm}\hrule height 0.2pt \vspace{3mm}

% DB question_id: 3aae43a7-6e8b-4a96-843d-328e00d14ffb
\needspace{5\baselineskip}
\noindent\textbf{\color{ink}709.}
\par
H.264 oqimi ijro paytida dekodlandi. Vaziyatni tahlil qilib, unga eng aniq mos tushunchani toping.
\begin{enumerate}[label=\Alph*),leftmargin=9mm]
\item Kodek
\item Konteyner
\item Media oqimi
\item Multimedia
\end{enumerate}
\begin{tcolorbox}[enhanced,breakable,colback=teal!4,colframe=teal!50!black,boxrule=0.5pt,arc=1.1mm,left=3mm,right=3mm,top=1.7mm,bottom=1.7mm,before skip=2mm,after skip=2mm]
\textbf{\color{teal}To'g'ri javob:} \textbf{A) Kodek}
\par\smallskip
\textbf{Izoh:} Vaziyat Kodekga misol: media oqimini kodlaydigan va dekodlaydigan algoritm yoki dasturiy modul.
\end{tcolorbox}
\vspace{2mm}\hrule height 0.2pt \vspace{3mm}

% DB question_id: 3bd0e1b7-cc41-5e81-9697-3021dd3e4a68
\needspace{5\baselineskip}
\noindent\textbf{\color{ink}710.}
\par
11025 Hz, 8 bit, 1 kanal, 120 soniyalik siqilmagan audio hajmini MiB da toping.
\begin{enumerate}[label=\Alph*),leftmargin=9mm]
\item 1,262 MiB
\item 10 584 000 MiB
\item 2 646 000 MiB
\item 11 025 MiB
\end{enumerate}
\begin{tcolorbox}[enhanced,breakable,colback=teal!4,colframe=teal!50!black,boxrule=0.5pt,arc=1.1mm,left=3mm,right=3mm,top=1.7mm,bottom=1.7mm,before skip=2mm,after skip=2mm]
\textbf{\color{teal}To'g'ri javob:} \textbf{A) 1,262 MiB}
\par\smallskip
\textbf{Izoh:} Hajm = 11025\ensuremath{\times}8\ensuremath{\times}1\ensuremath{\times}120/8 = 1 323 000 B = 1,262 MiB.
\end{tcolorbox}
\vspace{2mm}\hrule height 0.2pt \vspace{3mm}

% DB question_id: 3bf9dfe2-e888-5f6b-bacd-895d2a827f8b
\needspace{5\baselineskip}
\noindent\textbf{\color{ink}711.}
\par
Video 5 Mbit/s, audio 320 kbit/s, davomiylik 600 s. Umumiy media hajmi qancha? MB = 10\textsuperscript{6} B.
\begin{enumerate}[label=\Alph*),leftmargin=9mm]
\item 375 MB
\item 24 MB
\item 399 MB
\item 3 192 MB
\end{enumerate}
\begin{tcolorbox}[enhanced,breakable,colback=teal!4,colframe=teal!50!black,boxrule=0.5pt,arc=1.1mm,left=3mm,right=3mm,top=1.7mm,bottom=1.7mm,before skip=2mm,after skip=2mm]
\textbf{\color{teal}To'g'ri javob:} \textbf{C) 399 MB}
\par\smallskip
\textbf{Izoh:} Jami bitrate = 5.32 Mbit/s; hajm = jami\ensuremath{\times}600/8 = 399 MB.
\end{tcolorbox}
\vspace{2mm}\hrule height 0.2pt \vspace{3mm}

% DB question_id: 3d5ab807-a44f-4df5-9501-b55c6af55ac9
\needspace{5\baselineskip}
\noindent\textbf{\color{ink}712.}
\par
Multimedia nima?
\begin{enumerate}[label=\Alph*),leftmargin=9mm]
\item Bir nechta media turining yagona mahsulotda uyg'unlashuvi
\item Bitta raqamli video oqimining faylda siqib saqlanishi
\item Turli media fayllarini bitta formatga avtomatik o'tkazish jarayoni
\item Bir turdagi media ma'lumotlarini tarmoq orqali oqimli uzatish texnologiyasi
\end{enumerate}
\begin{tcolorbox}[enhanced,breakable,colback=teal!4,colframe=teal!50!black,boxrule=0.5pt,arc=1.1mm,left=3mm,right=3mm,top=1.7mm,bottom=1.7mm,before skip=2mm,after skip=2mm]
\textbf{\color{teal}To'g'ri javob:} \textbf{A) Bir nechta media turining yagona mahsulotda uyg'unlashuvi}
\par\smallskip
\textbf{Izoh:} Multimedia matn, tasvir, ovoz, video yoki animatsiya kabi bir necha media turini bir mahsulotda birlashtiradi. Video kodlash, fayl formatini konvertatsiya qilish va oqimli uzatish alohida texnik jarayonlar bo'lib, multimedianing to'liq ta'rifi emas.
\end{tcolorbox}
\vspace{2mm}\hrule height 0.2pt \vspace{3mm}

% DB question_id: 3e67ed3f-dd7e-5173-bf16-c65307232ec3
\needspace{5\baselineskip}
\noindent\textbf{\color{ink}713.}
\par
Vektor tasvir nima?
\begin{enumerate}[label=\Alph*),leftmargin=9mm]
\item Faqat piksellar ro'yxati
\item Ovoz namunasi
\item Video konteyner
\item Geometrik obyektlar va matematik tavsiflar asosidagi tasvir
\end{enumerate}
\begin{tcolorbox}[enhanced,breakable,colback=teal!4,colframe=teal!50!black,boxrule=0.5pt,arc=1.1mm,left=3mm,right=3mm,top=1.7mm,bottom=1.7mm,before skip=2mm,after skip=2mm]
\textbf{\color{teal}To'g'ri javob:} \textbf{D) Geometrik obyektlar va matematik tavsiflar asosidagi tasvir}
\par\smallskip
\textbf{Izoh:} Vektor obyektlar masshtablanganda sifatini yaxshi saqlaydi.
\end{tcolorbox}
\vspace{2mm}\hrule height 0.2pt \vspace{3mm}

% DB question_id: 4087d4a7-9a12-5ea3-837c-12e727e68390
\needspace{5\baselineskip}
\noindent\textbf{\color{ink}714.}
\par
Matnda 1500 ta belgi bor, har biri 2 bayt bilan kodlangan. Hajmni kB da toping.
\begin{enumerate}[label=\Alph*),leftmargin=9mm]
\item 1 500 kB
\item 3 kB
\item 24 000 kB
\item 375 kB
\end{enumerate}
\begin{tcolorbox}[enhanced,breakable,colback=teal!4,colframe=teal!50!black,boxrule=0.5pt,arc=1.1mm,left=3mm,right=3mm,top=1.7mm,bottom=1.7mm,before skip=2mm,after skip=2mm]
\textbf{\color{teal}To'g'ri javob:} \textbf{B) 3 kB}
\par\smallskip
\textbf{Izoh:} 1500 \ensuremath{\times} 2 = 3000 B = 3 kB.
\end{tcolorbox}
\vspace{2mm}\hrule height 0.2pt \vspace{3mm}

% DB question_id: 43d347b6-2321-5d37-b13c-44c6f74667cc
\needspace{5\baselineskip}
\noindent\textbf{\color{ink}715.}
\par
Matnda 1200 ta belgi bor, har biri 2 bayt bilan kodlangan. Hajmni B da toping.
\begin{enumerate}[label=\Alph*),leftmargin=9mm]
\item 1 200 B
\item 2 400 B
\item 19 200 B
\item 300 B
\end{enumerate}
\begin{tcolorbox}[enhanced,breakable,colback=teal!4,colframe=teal!50!black,boxrule=0.5pt,arc=1.1mm,left=3mm,right=3mm,top=1.7mm,bottom=1.7mm,before skip=2mm,after skip=2mm]
\textbf{\color{teal}To'g'ri javob:} \textbf{B) 2 400 B}
\par\smallskip
\textbf{Izoh:} 1200 \ensuremath{\times} 2 = 2400 B = 2 400 B.
\end{tcolorbox}
\vspace{2mm}\hrule height 0.2pt \vspace{3mm}

% DB question_id: 442ec4b1-4daa-55ac-8757-26c76a4de57b
\needspace{5\baselineskip}
\noindent\textbf{\color{ink}716.}
\par
Tasvir kengligi 2,54 cm, zichligi 600 PPI. 1 dyuym = 2,54 cm. Piksel kengligi taxminan qancha?
\begin{enumerate}[label=\Alph*),leftmargin=9mm]
\item 1 524 piksel
\item 300 piksel
\item 600 piksel
\item 603 piksel
\end{enumerate}
\begin{tcolorbox}[enhanced,breakable,colback=teal!4,colframe=teal!50!black,boxrule=0.5pt,arc=1.1mm,left=3mm,right=3mm,top=1.7mm,bottom=1.7mm,before skip=2mm,after skip=2mm]
\textbf{\color{teal}To'g'ri javob:} \textbf{C) 600 piksel}
\par\smallskip
\textbf{Izoh:} 2,54/2,54\ensuremath{\times}600 \ensuremath{\approx} 600 piksel.
\end{tcolorbox}
\vspace{2mm}\hrule height 0.2pt \vspace{3mm}

% DB question_id: 45f5b954-24e9-54c1-86c3-f61ca145a1e8
\needspace{5\baselineskip}
\noindent\textbf{\color{ink}717.}
\par
Loyihada 500 ta 6 MB rasm va 4 ta 1524 MB audio bor. Umumiy hajm qancha?
\begin{enumerate}[label=\Alph*),leftmargin=9mm]
\item 9 096 MB
\item 504 MB
\item 1 530 MB
\item 72 768 MB
\end{enumerate}
\begin{tcolorbox}[enhanced,breakable,colback=teal!4,colframe=teal!50!black,boxrule=0.5pt,arc=1.1mm,left=3mm,right=3mm,top=1.7mm,bottom=1.7mm,before skip=2mm,after skip=2mm]
\textbf{\color{teal}To'g'ri javob:} \textbf{A) 9 096 MB}
\par\smallskip
\textbf{Izoh:} 500\ensuremath{\times}6 + 4\ensuremath{\times}1524 = 9096 MB.
\end{tcolorbox}
\vspace{2mm}\hrule height 0.2pt \vspace{3mm}

% DB question_id: 479c809a-a4cc-5689-b2e0-4c8627c39b8b
\needspace{5\baselineskip}
\noindent\textbf{\color{ink}718.}
\par
300\ensuremath{\times}300 pikselli, 1 bit rang chuqurligidagi siqilmagan rastr tasvir hajmini B da toping.
\begin{enumerate}[label=\Alph*),leftmargin=9mm]
\item 11 250 B
\item 90 000 B
\item 11 252 B
\item 11 253 B
\end{enumerate}
\begin{tcolorbox}[enhanced,breakable,colback=teal!4,colframe=teal!50!black,boxrule=0.5pt,arc=1.1mm,left=3mm,right=3mm,top=1.7mm,bottom=1.7mm,before skip=2mm,after skip=2mm]
\textbf{\color{teal}To'g'ri javob:} \textbf{A) 11 250 B}
\par\smallskip
\textbf{Izoh:} Hajm = 300\ensuremath{\times}300\ensuremath{\times}1/8 = 11 250 B = 11 250 B.
\end{tcolorbox}
\vspace{2mm}\hrule height 0.2pt \vspace{3mm}

% DB question_id: 4f4f0066-617a-5070-add2-91e67c478632
\needspace{5\baselineskip}
\noindent\textbf{\color{ink}719.}
\par
22000 Hz, 8 bit, 1 kanal, 60 soniyalik siqilmagan audio hajmini MB da toping.
\begin{enumerate}[label=\Alph*),leftmargin=9mm]
\item 10 560 000 MB
\item 1,32 MB
\item 2 640 000 MB
\item 22 000 MB
\end{enumerate}
\begin{tcolorbox}[enhanced,breakable,colback=teal!4,colframe=teal!50!black,boxrule=0.5pt,arc=1.1mm,left=3mm,right=3mm,top=1.7mm,bottom=1.7mm,before skip=2mm,after skip=2mm]
\textbf{\color{teal}To'g'ri javob:} \textbf{B) 1,32 MB}
\par\smallskip
\textbf{Izoh:} Hajm = 22000\ensuremath{\times}8\ensuremath{\times}1\ensuremath{\times}60/8 = 1 320 000 B = 1,32 MB.
\end{tcolorbox}
\vspace{2mm}\hrule height 0.2pt \vspace{3mm}

% DB question_id: 517de5ee-bfae-46fb-9448-f9df785b733e
\needspace{5\baselineskip}
\noindent\textbf{\color{ink}720.}
\par
16 bitli yozuv har namuna uchun 65 536 darajani ifodaladi. Vaziyatni tahlil qilib, unga eng aniq mos tushunchani toping.
\begin{enumerate}[label=\Alph*),leftmargin=9mm]
\item Kanal
\item PCM
\item Audio bit tezligi
\item Bit chuqurligi
\end{enumerate}
\begin{tcolorbox}[enhanced,breakable,colback=teal!4,colframe=teal!50!black,boxrule=0.5pt,arc=1.1mm,left=3mm,right=3mm,top=1.7mm,bottom=1.7mm,before skip=2mm,after skip=2mm]
\textbf{\color{teal}To'g'ri javob:} \textbf{D) Bit chuqurligi}
\par\smallskip
\textbf{Izoh:} Vaziyat Bit chuqurligiga misol: bitta audio namunasini kodlash uchun ajratilgan bitlar soni.
\end{tcolorbox}
\vspace{2mm}\hrule height 0.2pt \vspace{3mm}

% DB question_id: 5578cd95-65e6-5373-a737-953decd243a2
\needspace{5\baselineskip}
\noindent\textbf{\color{ink}721.}
\par
Mono audio stereo qilinsa, boshqa parametrlar o'zgarmasa hajm qanday o'zgaradi?
\begin{enumerate}[label=\Alph*),leftmargin=9mm]
\item 4 marta oshadi
\item O'zgarmaydi
\item Yarimlashadi
\item 2 marta oshadi
\end{enumerate}
\begin{tcolorbox}[enhanced,breakable,colback=teal!4,colframe=teal!50!black,boxrule=0.5pt,arc=1.1mm,left=3mm,right=3mm,top=1.7mm,bottom=1.7mm,before skip=2mm,after skip=2mm]
\textbf{\color{teal}To'g'ri javob:} \textbf{D) 2 marta oshadi}
\par\smallskip
\textbf{Izoh:} Kanal soni 1 dan 2 ga oshadi.
\end{tcolorbox}
\vspace{2mm}\hrule height 0.2pt \vspace{3mm}

% DB question_id: 5638d3c1-8fe9-547f-beb9-f8be7cd1bfa3
\needspace{5\baselineskip}
\noindent\textbf{\color{ink}722.}
\par
24 bit rang chuqurligi nazariy jihatdan nechta rangni ifodalaydi?
\begin{enumerate}[label=\Alph*),leftmargin=9mm]
\item 16 777 216 rang
\item 24 rang
\item 48 rang
\item 8 388 608 rang
\end{enumerate}
\begin{tcolorbox}[enhanced,breakable,colback=teal!4,colframe=teal!50!black,boxrule=0.5pt,arc=1.1mm,left=3mm,right=3mm,top=1.7mm,bottom=1.7mm,before skip=2mm,after skip=2mm]
\textbf{\color{teal}To'g'ri javob:} \textbf{A) 16 777 216 rang}
\par\smallskip
\textbf{Izoh:} Ranglar soni = 2\textasciicircum{}24 = 16 777 216.
\end{tcolorbox}
\vspace{2mm}\hrule height 0.2pt \vspace{3mm}

% DB question_id: 591b4907-159d-5207-98ec-8029a4296d0d
\needspace{5\baselineskip}
\noindent\textbf{\color{ink}723.}
\par
FPS nima?
\begin{enumerate}[label=\Alph*),leftmargin=9mm]
\item Pikselga bitlar soni
\item Audio kanallar soni
\item Bir soniyadagi video kadrlar soni
\item Fayl siqish foizi
\end{enumerate}
\begin{tcolorbox}[enhanced,breakable,colback=teal!4,colframe=teal!50!black,boxrule=0.5pt,arc=1.1mm,left=3mm,right=3mm,top=1.7mm,bottom=1.7mm,before skip=2mm,after skip=2mm]
\textbf{\color{teal}To'g'ri javob:} \textbf{C) Bir soniyadagi video kadrlar soni}
\par\smallskip
\textbf{Izoh:} FPS frame per second degani.
\end{tcolorbox}
\vspace{2mm}\hrule height 0.2pt \vspace{3mm}

% DB question_id: 5b34608f-1bc9-5278-bbea-7bb2c600399c
\needspace{5\baselineskip}
\noindent\textbf{\color{ink}724.}
\par
Kamida 100 xil rangni kodlash uchun minimal rang chuqurligi qancha?
\begin{enumerate}[label=\Alph*),leftmargin=9mm]
\item 7 bit
\item 100 bit
\item 6 bit
\item 8 bit
\end{enumerate}
\begin{tcolorbox}[enhanced,breakable,colback=teal!4,colframe=teal!50!black,boxrule=0.5pt,arc=1.1mm,left=3mm,right=3mm,top=1.7mm,bottom=1.7mm,before skip=2mm,after skip=2mm]
\textbf{\color{teal}To'g'ri javob:} \textbf{A) 7 bit}
\par\smallskip
\textbf{Izoh:} 2\textsuperscript{n} \ensuremath{\geq} 100; minimal n = 7.
\end{tcolorbox}
\vspace{2mm}\hrule height 0.2pt \vspace{3mm}

% DB question_id: 5d4ae150-bc85-4050-8958-04e82edeb981
\needspace{5\baselineskip}
\noindent\textbf{\color{ink}725.}
\par
250\ensuremath{\times}300 pikselli har bir tasvir uchun 40 KiB dan oshmaydigan xotira ajratilgan. Rasm siqilmagan, barcha pikselga bir xil butun bit uzunligi ishlatiladi va xizmat baytlari yo'q. Eng ko'pi bilan nechta rang kodlash mumkin?
\begin{enumerate}[label=\Alph*),leftmargin=9mm]
\item 8 ta
\item 16 ta
\item 32 ta
\item 64 ta
\end{enumerate}
\begin{tcolorbox}[enhanced,breakable,colback=teal!4,colframe=teal!50!black,boxrule=0.5pt,arc=1.1mm,left=3mm,right=3mm,top=1.7mm,bottom=1.7mm,before skip=2mm,after skip=2mm]
\textbf{\color{teal}To'g'ri javob:} \textbf{B) 16 ta}
\par\smallskip
\textbf{Izoh:} Piksel soni 250\ensuremath{\times}300=75000. Xotira 40\ensuremath{\times}1024\ensuremath{\times}8=327680 bit; bitta pikselga ko'pi bilan floor(327680/75000)=4 bit ajratish mumkin. 4 bit 2\textsuperscript{4}=16 rang beradi; 5 bit/piksel 375000 bit bo'lib limitdan oshadi.
\end{tcolorbox}
\vspace{2mm}\hrule height 0.2pt \vspace{3mm}

% DB question_id: 5e4b7073-9edc-44e1-9716-2cd627023106
\needspace{5\baselineskip}
\noindent\textbf{\color{ink}726.}
\par
Konteyner nima?
\begin{enumerate}[label=\Alph*),leftmargin=9mm]
\item Media oqimidagi ortiqcha ma'lumot hajmini kamaytiradigan siqish algoritmi
\item Videodagi har bir kadrni tashkil qiluvchi eng kichik tasvir elementi
\item Audio va video oqimlarini qayta kodlashga xizmat qiluvchi dasturiy vosita
\item Video, audio, subtitr va metama'lumot oqimlarini bitta faylda birlashtiruvchi tuzilma
\end{enumerate}
\begin{tcolorbox}[enhanced,breakable,colback=teal!4,colframe=teal!50!black,boxrule=0.5pt,arc=1.1mm,left=3mm,right=3mm,top=1.7mm,bottom=1.7mm,before skip=2mm,after skip=2mm]
\textbf{\color{teal}To'g'ri javob:} \textbf{D) Video, audio, subtitr va metama'lumot oqimlarini bitta faylda birlashtiruvchi tuzilma}
\par\smallskip
\textbf{Izoh:} Media konteyneri, masalan MP4 yoki MKV, turli oqimlarni va metama'lumotni bitta faylga joylashtiradi. Siqish algoritmi kodekka tegishli; piksel tasvir elementi; transkoder esa oqimning kodlash formatini o'zgartiradi.
\end{tcolorbox}
\vspace{2mm}\hrule height 0.2pt \vspace{3mm}

% DB question_id: 5f15c4f7-efb1-53be-8058-4bf3c774d339
\needspace{5\baselineskip}
\noindent\textbf{\color{ink}727.}
\par
Video bitrate'i 25 Mbit/s, davomiyligi 60 s. Hajm qancha? MB = 10\textsuperscript{6} B.
\begin{enumerate}[label=\Alph*),leftmargin=9mm]
\item 1 500 MB
\item 23,44 MB
\item 187,5 MB
\item 190,5 MB
\end{enumerate}
\begin{tcolorbox}[enhanced,breakable,colback=teal!4,colframe=teal!50!black,boxrule=0.5pt,arc=1.1mm,left=3mm,right=3mm,top=1.7mm,bottom=1.7mm,before skip=2mm,after skip=2mm]
\textbf{\color{teal}To'g'ri javob:} \textbf{C) 187,5 MB}
\par\smallskip
\textbf{Izoh:} 25\ensuremath{\times}60/8 = 187,5 MB.
\end{tcolorbox}
\vspace{2mm}\hrule height 0.2pt \vspace{3mm}

% DB question_id: 5f5f06e6-8381-58bb-855f-17a8439c2063
\needspace{5\baselineskip}
\noindent\textbf{\color{ink}728.}
\par
Kodek va konteyner farqi qaysi?
\begin{enumerate}[label=\Alph*),leftmargin=9mm]
\item Ikkalasi bir xil
\item Konteyner rang modeli
\item Kodek oqimni kodlaydi/siqadi, konteyner esa oqimlarni bir faylga joylaydi
\item Kodek faqat fayl nomi
\end{enumerate}
\begin{tcolorbox}[enhanced,breakable,colback=teal!4,colframe=teal!50!black,boxrule=0.5pt,arc=1.1mm,left=3mm,right=3mm,top=1.7mm,bottom=1.7mm,before skip=2mm,after skip=2mm]
\textbf{\color{teal}To'g'ri javob:} \textbf{C) Kodek oqimni kodlaydi/siqadi, konteyner esa oqimlarni bir faylga joylaydi}
\par\smallskip
\textbf{Izoh:} Masalan H.264 kodek, MP4 konteyner bo'lishi mumkin.
\end{tcolorbox}
\vspace{2mm}\hrule height 0.2pt \vspace{3mm}

% DB question_id: 619c2fb5-3318-582b-ae6e-1c6470ab831d
\needspace{5\baselineskip}
\noindent\textbf{\color{ink}729.}
\par
Matnda 4096 ta belgi bor, har biri 1 bayt bilan kodlangan. Hajmni KiB da toping.
\begin{enumerate}[label=\Alph*),leftmargin=9mm]
\item 4 096 KiB
\item 32 768 KiB
\item 4 KiB
\item 512 KiB
\end{enumerate}
\begin{tcolorbox}[enhanced,breakable,colback=teal!4,colframe=teal!50!black,boxrule=0.5pt,arc=1.1mm,left=3mm,right=3mm,top=1.7mm,bottom=1.7mm,before skip=2mm,after skip=2mm]
\textbf{\color{teal}To'g'ri javob:} \textbf{C) 4 KiB}
\par\smallskip
\textbf{Izoh:} 4096 \ensuremath{\times} 1 = 4096 B = 4 KiB.
\end{tcolorbox}
\vspace{2mm}\hrule height 0.2pt \vspace{3mm}

% DB question_id: 62980a6d-92ba-5503-b1a0-ee2b781f81ef
\needspace{5\baselineskip}
\noindent\textbf{\color{ink}730.}
\par
Loyihada 60 ta 25 MB rasm va 2 ta 170 MB audio bor. Umumiy hajm qancha?
\begin{enumerate}[label=\Alph*),leftmargin=9mm]
\item 62 MB
\item 195 MB
\item 14 720 MB
\item 1 840 MB
\end{enumerate}
\begin{tcolorbox}[enhanced,breakable,colback=teal!4,colframe=teal!50!black,boxrule=0.5pt,arc=1.1mm,left=3mm,right=3mm,top=1.7mm,bottom=1.7mm,before skip=2mm,after skip=2mm]
\textbf{\color{teal}To'g'ri javob:} \textbf{D) 1 840 MB}
\par\smallskip
\textbf{Izoh:} 60\ensuremath{\times}25 + 2\ensuremath{\times}170 = 1840 MB.
\end{tcolorbox}
\vspace{2mm}\hrule height 0.2pt \vspace{3mm}

% DB question_id: 63465b5c-e012-5c31-9ecf-d1675509cd15
\needspace{5\baselineskip}
\noindent\textbf{\color{ink}731.}
\par
Matnda 1000 ta bir baytli va 250 ta 2 baytli belgi bor. Umumiy hajm qancha?
\begin{enumerate}[label=\Alph*),leftmargin=9mm]
\item 1 250 B
\item 12 000 B
\item 187,5 B
\item 1 500 B
\end{enumerate}
\begin{tcolorbox}[enhanced,breakable,colback=teal!4,colframe=teal!50!black,boxrule=0.5pt,arc=1.1mm,left=3mm,right=3mm,top=1.7mm,bottom=1.7mm,before skip=2mm,after skip=2mm]
\textbf{\color{teal}To'g'ri javob:} \textbf{D) 1 500 B}
\par\smallskip
\textbf{Izoh:} 1000\ensuremath{\times}1 + 250\ensuremath{\times}2 = 1500 B = 1 500 B.
\end{tcolorbox}
\vspace{2mm}\hrule height 0.2pt \vspace{3mm}

% DB question_id: 63d98048-3b82-59f9-9d3f-396d32dd02fd
\needspace{5\baselineskip}
\noindent\textbf{\color{ink}732.}
\par
Bit chuqurligi 16 dan 24 ga oshsa hajm necha marta o'zgaradi?
\begin{enumerate}[label=\Alph*),leftmargin=9mm]
\item 1,5 marta oshadi
\item 2 marta oshadi
\item 8 marta oshadi
\item O'zgarmaydi
\end{enumerate}
\begin{tcolorbox}[enhanced,breakable,colback=teal!4,colframe=teal!50!black,boxrule=0.5pt,arc=1.1mm,left=3mm,right=3mm,top=1.7mm,bottom=1.7mm,before skip=2mm,after skip=2mm]
\textbf{\color{teal}To'g'ri javob:} \textbf{A) 1,5 marta oshadi}
\par\smallskip
\textbf{Izoh:} 24/16 = 1,5.
\end{tcolorbox}
\vspace{2mm}\hrule height 0.2pt \vspace{3mm}

% DB question_id: 662845f0-ed49-5c05-a5df-baa93799dce9
\needspace{5\baselineskip}
\noindent\textbf{\color{ink}733.}
\par
Analog tovush nima?
\begin{enumerate}[label=\Alph*),leftmargin=9mm]
\item Faqat 0 va 1
\item Kadrlar ketma-ketligi
\item Vaqt va amplituda bo'yicha uzluksiz signal
\item Matnli kod
\end{enumerate}
\begin{tcolorbox}[enhanced,breakable,colback=teal!4,colframe=teal!50!black,boxrule=0.5pt,arc=1.1mm,left=3mm,right=3mm,top=1.7mm,bottom=1.7mm,before skip=2mm,after skip=2mm]
\textbf{\color{teal}To'g'ri javob:} \textbf{C) Vaqt va amplituda bo'yicha uzluksiz signal}
\par\smallskip
\textbf{Izoh:} Haqiqiy tovush to'lqini uzluksizdir.
\end{tcolorbox}
\vspace{2mm}\hrule height 0.2pt \vspace{3mm}

% DB question_id: 6c731185-d5f2-525b-aaa0-e9623807574e
\needspace{5\baselineskip}
\noindent\textbf{\color{ink}734.}
\par
ASCII asl standartining asosiy xususiyati qaysi?
\begin{enumerate}[label=\Alph*),leftmargin=9mm]
\item 7 bitli 128 kod holati
\item Har belgi 4 bayt
\item Faqat rasm kodlaydi
\item Barcha yozuvlarni to'liq qamraydi
\end{enumerate}
\begin{tcolorbox}[enhanced,breakable,colback=teal!4,colframe=teal!50!black,boxrule=0.5pt,arc=1.1mm,left=3mm,right=3mm,top=1.7mm,bottom=1.7mm,before skip=2mm,after skip=2mm]
\textbf{\color{teal}To'g'ri javob:} \textbf{A) 7 bitli 128 kod holati}
\par\smallskip
\textbf{Izoh:} Asl ASCII 7 bitli va 128 holatli.
\end{tcolorbox}
\vspace{2mm}\hrule height 0.2pt \vspace{3mm}

% DB question_id: 6e276355-1760-519c-8bbc-700fd853e3f5
\needspace{5\baselineskip}
\noindent\textbf{\color{ink}735.}
\par
Video bitrate'i 50 Mbit/s, davomiyligi 30 s. Hajm qancha? MB = 10\textsuperscript{6} B.
\begin{enumerate}[label=\Alph*),leftmargin=9mm]
\item 1 500 MB
\item 23,44 MB
\item 190,5 MB
\item 187,5 MB
\end{enumerate}
\begin{tcolorbox}[enhanced,breakable,colback=teal!4,colframe=teal!50!black,boxrule=0.5pt,arc=1.1mm,left=3mm,right=3mm,top=1.7mm,bottom=1.7mm,before skip=2mm,after skip=2mm]
\textbf{\color{teal}To'g'ri javob:} \textbf{D) 187,5 MB}
\par\smallskip
\textbf{Izoh:} 50\ensuremath{\times}30/8 = 187,5 MB.
\end{tcolorbox}
\vspace{2mm}\hrule height 0.2pt \vspace{3mm}

% DB question_id: 7ae7b975-84c2-56c7-a65b-63ceececc955
\needspace{5\baselineskip}
\noindent\textbf{\color{ink}736.}
\par
44100 Hz, 24 bit, 2 kanal, 60 soniyalik siqilmagan audio hajmini MiB da toping.
\begin{enumerate}[label=\Alph*),leftmargin=9mm]
\item 127 008 000 MiB
\item 15,141 MiB
\item 7 938 000 MiB
\item 264 600 MiB
\end{enumerate}
\begin{tcolorbox}[enhanced,breakable,colback=teal!4,colframe=teal!50!black,boxrule=0.5pt,arc=1.1mm,left=3mm,right=3mm,top=1.7mm,bottom=1.7mm,before skip=2mm,after skip=2mm]
\textbf{\color{teal}To'g'ri javob:} \textbf{B) 15,141 MiB}
\par\smallskip
\textbf{Izoh:} Hajm = 44100\ensuremath{\times}24\ensuremath{\times}2\ensuremath{\times}60/8 = 15 876 000 B = 15,141 MiB.
\end{tcolorbox}
\vspace{2mm}\hrule height 0.2pt \vspace{3mm}

% DB question_id: 7bcb4faf-4a56-543c-b4f2-fe1040b17868
\needspace{5\baselineskip}
\noindent\textbf{\color{ink}737.}
\par
1024\ensuremath{\times}768 pikselli, 8 bit rang chuqurligidagi siqilmagan rastr tasvir hajmini KiB da toping.
\begin{enumerate}[label=\Alph*),leftmargin=9mm]
\item 768 KiB
\item 786 432 KiB
\item 6 291 456 KiB
\item 771 KiB
\end{enumerate}
\begin{tcolorbox}[enhanced,breakable,colback=teal!4,colframe=teal!50!black,boxrule=0.5pt,arc=1.1mm,left=3mm,right=3mm,top=1.7mm,bottom=1.7mm,before skip=2mm,after skip=2mm]
\textbf{\color{teal}To'g'ri javob:} \textbf{A) 768 KiB}
\par\smallskip
\textbf{Izoh:} Hajm = 1024\ensuremath{\times}768\ensuremath{\times}8/8 = 786 432 B = 768 KiB.
\end{tcolorbox}
\vspace{2mm}\hrule height 0.2pt \vspace{3mm}

% DB question_id: 7c16457e-e9d1-51f8-b270-d39ed861acaa
\needspace{5\baselineskip}
\noindent\textbf{\color{ink}738.}
\par
Matnda 900 ta bir baytli va 100 ta 2 baytli belgi bor. Umumiy hajm qancha?
\begin{enumerate}[label=\Alph*),leftmargin=9mm]
\item 1 000 B
\item 8 800 B
\item 1 100 B
\item 137,5 B
\end{enumerate}
\begin{tcolorbox}[enhanced,breakable,colback=teal!4,colframe=teal!50!black,boxrule=0.5pt,arc=1.1mm,left=3mm,right=3mm,top=1.7mm,bottom=1.7mm,before skip=2mm,after skip=2mm]
\textbf{\color{teal}To'g'ri javob:} \textbf{C) 1 100 B}
\par\smallskip
\textbf{Izoh:} 900\ensuremath{\times}1 + 100\ensuremath{\times}2 = 1100 B = 1 100 B.
\end{tcolorbox}
\vspace{2mm}\hrule height 0.2pt \vspace{3mm}

% DB question_id: 7cf117e5-4bee-5aa2-b944-7eb1371da211
\needspace{5\baselineskip}
\noindent\textbf{\color{ink}739.}
\par
Kodek nima?
\begin{enumerate}[label=\Alph*),leftmargin=9mm]
\item Faylni o'rab turuvchi konteyner
\item Faqat kengaytma
\item Audio yoki videoni kodlaydigan va dekodlaydigan usul yoki dastur
\item Rang modeli
\end{enumerate}
\begin{tcolorbox}[enhanced,breakable,colback=teal!4,colframe=teal!50!black,boxrule=0.5pt,arc=1.1mm,left=3mm,right=3mm,top=1.7mm,bottom=1.7mm,before skip=2mm,after skip=2mm]
\textbf{\color{teal}To'g'ri javob:} \textbf{C) Audio yoki videoni kodlaydigan va dekodlaydigan usul yoki dastur}
\par\smallskip
\textbf{Izoh:} Kodek compression/decompression vazifasini bajaradi.
\end{tcolorbox}
\vspace{2mm}\hrule height 0.2pt \vspace{3mm}

% DB question_id: 7ed881dd-0305-501f-a1b0-afbbe67b1818
\needspace{5\baselineskip}
\noindent\textbf{\color{ink}740.}
\par
Matnda 1024 ta belgi bor, har biri 2 bayt bilan kodlangan. Hajmni KiB da toping.
\begin{enumerate}[label=\Alph*),leftmargin=9mm]
\item 1 024 KiB
\item 16 384 KiB
\item 256 KiB
\item 2 KiB
\end{enumerate}
\begin{tcolorbox}[enhanced,breakable,colback=teal!4,colframe=teal!50!black,boxrule=0.5pt,arc=1.1mm,left=3mm,right=3mm,top=1.7mm,bottom=1.7mm,before skip=2mm,after skip=2mm]
\textbf{\color{teal}To'g'ri javob:} \textbf{D) 2 KiB}
\par\smallskip
\textbf{Izoh:} 1024 \ensuremath{\times} 2 = 2048 B = 2 KiB.
\end{tcolorbox}
\vspace{2mm}\hrule height 0.2pt \vspace{3mm}

% DB question_id: 7f6d8ba5-a4b9-4fc9-a8b5-6054f8dced33
\needspace{5\baselineskip}
\noindent\textbf{\color{ink}741.}
\par
Rang chuqurligi nima?
\begin{enumerate}[label=\Alph*),leftmargin=9mm]
\item Rastr tasvirning gorizontal yo'nalishidagi umumiy piksellar soni
\item Bitta piksel rangini kodlash uchun ajratilgan bitlar soni
\item Videoda bir soniyada namoyish etiladigan kadrlar umumiy soni
\item Bitta pikselning ekranda egallaydigan fizik maydon o'lchami
\end{enumerate}
\begin{tcolorbox}[enhanced,breakable,colback=teal!4,colframe=teal!50!black,boxrule=0.5pt,arc=1.1mm,left=3mm,right=3mm,top=1.7mm,bottom=1.7mm,before skip=2mm,after skip=2mm]
\textbf{\color{teal}To'g'ri javob:} \textbf{B) Bitta piksel rangini kodlash uchun ajratilgan bitlar soni}
\par\smallskip
\textbf{Izoh:} Rang chuqurligi har bir pikselning rangini kodlashga ajratilgan bitlar sonidir; nazariy ranglar soni odatda 2\textasciicircum{}n ga bog'liq. Tasvir eni piksel soni, kadrlar chastotasi vaqt birligidagi kadrlar soni, piksel maydoni esa boshqa o'lchamdir.
\end{tcolorbox}
\vspace{2mm}\hrule height 0.2pt \vspace{3mm}

% DB question_id: 809f22d2-ffde-5712-bdcf-e98519a16f35
\needspace{5\baselineskip}
\noindent\textbf{\color{ink}742.}
\par
3840\ensuremath{\times}2160 pikselli, 24 bit rang chuqurligidagi siqilmagan rastr tasvir hajmini MiB da toping.
\begin{enumerate}[label=\Alph*),leftmargin=9mm]
\item 8 294 400 MiB
\item 199 065 600 MiB
\item 23,73 MiB
\item 26,73 MiB
\end{enumerate}
\begin{tcolorbox}[enhanced,breakable,colback=teal!4,colframe=teal!50!black,boxrule=0.5pt,arc=1.1mm,left=3mm,right=3mm,top=1.7mm,bottom=1.7mm,before skip=2mm,after skip=2mm]
\textbf{\color{teal}To'g'ri javob:} \textbf{C) 23,73 MiB}
\par\smallskip
\textbf{Izoh:} Hajm = 3840\ensuremath{\times}2160\ensuremath{\times}24/8 = 24 883 200 B = 23,73 MiB.
\end{tcolorbox}
\vspace{2mm}\hrule height 0.2pt \vspace{3mm}

% DB question_id: 81636e54-620c-5206-859c-c128fe834d58
\needspace{5\baselineskip}
\noindent\textbf{\color{ink}743.}
\par
Matnda 6000 ta belgi bor, har biri 1 bayt bilan kodlangan. Hajmni KiB da toping.
\begin{enumerate}[label=\Alph*),leftmargin=9mm]
\item 5,859 KiB
\item 6 000 KiB
\item 48 000 KiB
\item 750 KiB
\end{enumerate}
\begin{tcolorbox}[enhanced,breakable,colback=teal!4,colframe=teal!50!black,boxrule=0.5pt,arc=1.1mm,left=3mm,right=3mm,top=1.7mm,bottom=1.7mm,before skip=2mm,after skip=2mm]
\textbf{\color{teal}To'g'ri javob:} \textbf{A) 5,859 KiB}
\par\smallskip
\textbf{Izoh:} 6000 \ensuremath{\times} 1 = 6000 B = 5,859 KiB.
\end{tcolorbox}
\vspace{2mm}\hrule height 0.2pt \vspace{3mm}

% DB question_id: 8ab38998-209d-5443-8ae6-782c4124f589
\needspace{5\baselineskip}
\noindent\textbf{\color{ink}744.}
\par
8192 baytli matnda har belgi 1 bayt bo'lsa, nechta belgi bor?
\begin{enumerate}[label=\Alph*),leftmargin=9mm]
\item 8 192 belgi
\item 65 536 belgi
\item 4 096 belgi
\item 8 195 belgi
\end{enumerate}
\begin{tcolorbox}[enhanced,breakable,colback=teal!4,colframe=teal!50!black,boxrule=0.5pt,arc=1.1mm,left=3mm,right=3mm,top=1.7mm,bottom=1.7mm,before skip=2mm,after skip=2mm]
\textbf{\color{teal}To'g'ri javob:} \textbf{A) 8 192 belgi}
\par\smallskip
\textbf{Izoh:} Belgilar soni = 8192/1 = 8192.
\end{tcolorbox}
\vspace{2mm}\hrule height 0.2pt \vspace{3mm}

% DB question_id: 90e4af91-a7a4-51c6-9493-3d48f773ef10
\needspace{5\baselineskip}
\noindent\textbf{\color{ink}745.}
\par
1920\ensuremath{\times}1080, 24 bit, 25 FPS, 1 soniyalik siqilmagan video hajmini MiB da toping.
\begin{enumerate}[label=\Alph*),leftmargin=9mm]
\item 18,539 MiB
\item 5,933 MiB
\item 148,315 MiB
\item 3707,886 MiB
\end{enumerate}
\begin{tcolorbox}[enhanced,breakable,colback=teal!4,colframe=teal!50!black,boxrule=0.5pt,arc=1.1mm,left=3mm,right=3mm,top=1.7mm,bottom=1.7mm,before skip=2mm,after skip=2mm]
\textbf{\color{teal}To'g'ri javob:} \textbf{C) 148,315 MiB}
\par\smallskip
\textbf{Izoh:} Hajm = W\ensuremath{\times}H\ensuremath{\times}d/8\ensuremath{\times}FPS\ensuremath{\times}t = 148,315 MiB.
\end{tcolorbox}
\vspace{2mm}\hrule height 0.2pt \vspace{3mm}

% DB question_id: 90e9e9fc-ab4c-5032-b328-b39b97efd382
\needspace{5\baselineskip}
\noindent\textbf{\color{ink}746.}
\par
Video 8 Mbit/s, audio 192 kbit/s, davomiylik 300 s. Umumiy media hajmi qancha? MB = 10\textsuperscript{6} B.
\begin{enumerate}[label=\Alph*),leftmargin=9mm]
\item 307,2 MB
\item 300 MB
\item 7,2 MB
\item 2457,6 MB
\end{enumerate}
\begin{tcolorbox}[enhanced,breakable,colback=teal!4,colframe=teal!50!black,boxrule=0.5pt,arc=1.1mm,left=3mm,right=3mm,top=1.7mm,bottom=1.7mm,before skip=2mm,after skip=2mm]
\textbf{\color{teal}To'g'ri javob:} \textbf{A) 307,2 MB}
\par\smallskip
\textbf{Izoh:} Jami bitrate = 8.192 Mbit/s; hajm = jami\ensuremath{\times}300/8 = 307,2 MB.
\end{tcolorbox}
\vspace{2mm}\hrule height 0.2pt \vspace{3mm}

% DB question_id: 93ca5c39-69e2-5aff-8f29-197451e05f07
\needspace{5\baselineskip}
\noindent\textbf{\color{ink}747.}
\par
Video bitrate'i 4 Mbit/s, davomiyligi 900 s. Hajm qancha? MB = 10\textsuperscript{6} B.
\begin{enumerate}[label=\Alph*),leftmargin=9mm]
\item 3 600 MB
\item 56,25 MB
\item 453 MB
\item 450 MB
\end{enumerate}
\begin{tcolorbox}[enhanced,breakable,colback=teal!4,colframe=teal!50!black,boxrule=0.5pt,arc=1.1mm,left=3mm,right=3mm,top=1.7mm,bottom=1.7mm,before skip=2mm,after skip=2mm]
\textbf{\color{teal}To'g'ri javob:} \textbf{D) 450 MB}
\par\smallskip
\textbf{Izoh:} 4\ensuremath{\times}900/8 = 450 MB.
\end{tcolorbox}
\vspace{2mm}\hrule height 0.2pt \vspace{3mm}

% DB question_id: 9549f7bf-1ab9-5053-9133-6cbdb8daa518
\needspace{5\baselineskip}
\noindent\textbf{\color{ink}748.}
\par
320\ensuremath{\times}240, 24 bit, 25 FPS, 2 soniyalik siqilmagan video hajmini MiB da toping.
\begin{enumerate}[label=\Alph*),leftmargin=9mm]
\item 10,986 MiB
\item 1,373 MiB
\item 0,22 MiB
\item 274,658 MiB
\end{enumerate}
\begin{tcolorbox}[enhanced,breakable,colback=teal!4,colframe=teal!50!black,boxrule=0.5pt,arc=1.1mm,left=3mm,right=3mm,top=1.7mm,bottom=1.7mm,before skip=2mm,after skip=2mm]
\textbf{\color{teal}To'g'ri javob:} \textbf{A) 10,986 MiB}
\par\smallskip
\textbf{Izoh:} Hajm = W\ensuremath{\times}H\ensuremath{\times}d/8\ensuremath{\times}FPS\ensuremath{\times}t = 10,986 MiB.
\end{tcolorbox}
\vspace{2mm}\hrule height 0.2pt \vspace{3mm}

% DB question_id: 95ef34b0-5797-4bcc-ab68-3059f7204995
\needspace{5\baselineskip}
\noindent\textbf{\color{ink}749.}
\par
Fayl uzatishdan oldin hajmi kichraytirildi. Bu vaziyatdagi asosiy tushunchani aniqlang.
\begin{enumerate}[label=\Alph*),leftmargin=9mm]
\item Yo'qotishsiz siqish
\item Siqish
\item Yo'qotishli siqish
\item RLE
\end{enumerate}
\begin{tcolorbox}[enhanced,breakable,colback=teal!4,colframe=teal!50!black,boxrule=0.5pt,arc=1.1mm,left=3mm,right=3mm,top=1.7mm,bottom=1.7mm,before skip=2mm,after skip=2mm]
\textbf{\color{teal}To'g'ri javob:} \textbf{B) Siqish}
\par\smallskip
\textbf{Izoh:} Vaziyat Siqishga misol: axborotni kamroq bit yoki bayt bilan ifodalash jarayoni.
\end{tcolorbox}
\vspace{2mm}\hrule height 0.2pt \vspace{3mm}

% DB question_id: 979f0c7f-05fd-5fc4-a32a-2b9d1ae20110
\needspace{5\baselineskip}
\noindent\textbf{\color{ink}750.}
\par
UTF-8ning xususiyati qaysi?
\begin{enumerate}[label=\Alph*),leftmargin=9mm]
\item Har doim 2 bayt
\item Faqat 7 bit
\item Belgilarni 1-4 baytli ketma-ketliklarda kodlaydi
\item Faqat videoda
\end{enumerate}
\begin{tcolorbox}[enhanced,breakable,colback=teal!4,colframe=teal!50!black,boxrule=0.5pt,arc=1.1mm,left=3mm,right=3mm,top=1.7mm,bottom=1.7mm,before skip=2mm,after skip=2mm]
\textbf{\color{teal}To'g'ri javob:} \textbf{C) Belgilarni 1-4 baytli ketma-ketliklarda kodlaydi}
\par\smallskip
\textbf{Izoh:} UTF-8 o'zgaruvchan uzunlikdagi Unicode kodlash shaklidir.
\end{tcolorbox}
\vspace{2mm}\hrule height 0.2pt \vspace{3mm}

% DB question_id: 9861745c-7455-56ba-bf8c-94270750a897
\needspace{5\baselineskip}
\noindent\textbf{\color{ink}751.}
\par
Loyihada 100 ta 8 MB rasm va 5 ta 900 MB audio bor. Umumiy hajm qancha?
\begin{enumerate}[label=\Alph*),leftmargin=9mm]
\item 105 MB
\item 908 MB
\item 5 300 MB
\item 42 400 MB
\end{enumerate}
\begin{tcolorbox}[enhanced,breakable,colback=teal!4,colframe=teal!50!black,boxrule=0.5pt,arc=1.1mm,left=3mm,right=3mm,top=1.7mm,bottom=1.7mm,before skip=2mm,after skip=2mm]
\textbf{\color{teal}To'g'ri javob:} \textbf{C) 5 300 MB}
\par\smallskip
\textbf{Izoh:} 100\ensuremath{\times}8 + 5\ensuremath{\times}900 = 5300 MB.
\end{tcolorbox}
\vspace{2mm}\hrule height 0.2pt \vspace{3mm}

% DB question_id: 99f2993a-d425-4ca8-b7fb-f5099b892e46
\needspace{5\baselineskip}
\noindent\textbf{\color{ink}752.}
\par
Fotosurat qayta saqlanganda sezilmaydigan detallar yo'qotildi. Bu vaziyatdagi asosiy tushunchani aniqlang.
\begin{enumerate}[label=\Alph*),leftmargin=9mm]
\item RLE
\item Siqish koeffitsiyenti
\item Kodek
\item Yo'qotishli siqish
\end{enumerate}
\begin{tcolorbox}[enhanced,breakable,colback=teal!4,colframe=teal!50!black,boxrule=0.5pt,arc=1.1mm,left=3mm,right=3mm,top=1.7mm,bottom=1.7mm,before skip=2mm,after skip=2mm]
\textbf{\color{teal}To'g'ri javob:} \textbf{D) Yo'qotishli siqish}
\par\smallskip
\textbf{Izoh:} Vaziyat Yo'qotishli siqishga misol: hajmni kamaytirish uchun ayrim ma'lumotni qaytmas tarzda olib tashlaydigan siqish.
\end{tcolorbox}
\vspace{2mm}\hrule height 0.2pt \vspace{3mm}

% DB question_id: 9cac11d1-e53c-5130-85b0-7e72b5b92162
\needspace{5\baselineskip}
\noindent\textbf{\color{ink}753.}
\par
Rastr tasvir nima?
\begin{enumerate}[label=\Alph*),leftmargin=9mm]
\item Faqat formulalardan tuzilgan
\item Audio oqim
\item Piksellar to'ridan tuzilgan tasvir
\item Matn kodi
\end{enumerate}
\begin{tcolorbox}[enhanced,breakable,colback=teal!4,colframe=teal!50!black,boxrule=0.5pt,arc=1.1mm,left=3mm,right=3mm,top=1.7mm,bottom=1.7mm,before skip=2mm,after skip=2mm]
\textbf{\color{teal}To'g'ri javob:} \textbf{C) Piksellar to'ridan tuzilgan tasvir}
\par\smallskip
\textbf{Izoh:} Rastr tasvirda har piksel rang qiymatiga ega.
\end{tcolorbox}
\vspace{2mm}\hrule height 0.2pt \vspace{3mm}

% DB question_id: 9e99f6f7-fa9a-52b0-ba59-775fe0a64b6b
\needspace{5\baselineskip}
\noindent\textbf{\color{ink}754.}
\par
60 FPS videoda 10 soniyada nechta kadr bo'ladi?
\begin{enumerate}[label=\Alph*),leftmargin=9mm]
\item 600 kadr
\item 70 kadr
\item 36 000 kadr
\item 300 kadr
\end{enumerate}
\begin{tcolorbox}[enhanced,breakable,colback=teal!4,colframe=teal!50!black,boxrule=0.5pt,arc=1.1mm,left=3mm,right=3mm,top=1.7mm,bottom=1.7mm,before skip=2mm,after skip=2mm]
\textbf{\color{teal}To'g'ri javob:} \textbf{A) 600 kadr}
\par\smallskip
\textbf{Izoh:} Kadrlar soni = 60\ensuremath{\times}10 = 600.
\end{tcolorbox}
\vspace{2mm}\hrule height 0.2pt \vspace{3mm}

% DB question_id: 9ebd96f8-2edd-41bd-915f-5c2550205382
\needspace{5\baselineskip}
\noindent\textbf{\color{ink}755.}
\par
Stereo yozuv chap va o'ng tovush yo'lini alohida saqladi. Vaziyatni tahlil qilib, unga eng aniq mos tushunchani toping.
\begin{enumerate}[label=\Alph*),leftmargin=9mm]
\item Kanal
\item PCM
\item Audio bit tezligi
\item Analog tovush
\end{enumerate}
\begin{tcolorbox}[enhanced,breakable,colback=teal!4,colframe=teal!50!black,boxrule=0.5pt,arc=1.1mm,left=3mm,right=3mm,top=1.7mm,bottom=1.7mm,before skip=2mm,after skip=2mm]
\textbf{\color{teal}To'g'ri javob:} \textbf{A) Kanal}
\par\smallskip
\textbf{Izoh:} Vaziyat Kanalga misol: mustaqil audio signal yo'li.
\end{tcolorbox}
\vspace{2mm}\hrule height 0.2pt \vspace{3mm}

% DB question_id: a07fbfb2-e145-544b-9f11-472c50735096
\needspace{5\baselineskip}
\noindent\textbf{\color{ink}756.}
\par
CMYK modeli qayerda ko'proq ishlatiladi?
\begin{enumerate}[label=\Alph*),leftmargin=9mm]
\item Monitor pikselida
\item Audio kanalda
\item Rangli bosmada
\item ASCII jadvalida
\end{enumerate}
\begin{tcolorbox}[enhanced,breakable,colback=teal!4,colframe=teal!50!black,boxrule=0.5pt,arc=1.1mm,left=3mm,right=3mm,top=1.7mm,bottom=1.7mm,before skip=2mm,after skip=2mm]
\textbf{\color{teal}To'g'ri javob:} \textbf{C) Rangli bosmada}
\par\smallskip
\textbf{Izoh:} CMYK bo'yoqlarni ayrish modelidir.
\end{tcolorbox}
\vspace{2mm}\hrule height 0.2pt \vspace{3mm}

% DB question_id: a22f7301-1232-5dcb-b678-f40de43004e2
\needspace{5\baselineskip}
\noindent\textbf{\color{ink}757.}
\par
12 bit rang chuqurligi nazariy jihatdan nechta rangni ifodalaydi?
\begin{enumerate}[label=\Alph*),leftmargin=9mm]
\item 12 rang
\item 24 rang
\item 4 096 rang
\item 2 048 rang
\end{enumerate}
\begin{tcolorbox}[enhanced,breakable,colback=teal!4,colframe=teal!50!black,boxrule=0.5pt,arc=1.1mm,left=3mm,right=3mm,top=1.7mm,bottom=1.7mm,before skip=2mm,after skip=2mm]
\textbf{\color{teal}To'g'ri javob:} \textbf{C) 4 096 rang}
\par\smallskip
\textbf{Izoh:} Ranglar soni = 2\textasciicircum{}12 = 4 096.
\end{tcolorbox}
\vspace{2mm}\hrule height 0.2pt \vspace{3mm}

% DB question_id: a3d848fa-b535-5650-ba40-d6ae0a40877e
\needspace{5\baselineskip}
\noindent\textbf{\color{ink}758.}
\par
48000 Hz, 24 bit, 2 kanalli PCM audioning xom bitrate'i qancha?
\begin{enumerate}[label=\Alph*),leftmargin=9mm]
\item 2 304 kbit/s
\item 1 152 000 kbit/s
\item 288 000 kbit/s
\item 4 608 000 kbit/s
\end{enumerate}
\begin{tcolorbox}[enhanced,breakable,colback=teal!4,colframe=teal!50!black,boxrule=0.5pt,arc=1.1mm,left=3mm,right=3mm,top=1.7mm,bottom=1.7mm,before skip=2mm,after skip=2mm]
\textbf{\color{teal}To'g'ri javob:} \textbf{A) 2 304 kbit/s}
\par\smallskip
\textbf{Izoh:} 48000\ensuremath{\times}24\ensuremath{\times}2 = 2304000 bit/s = 2 304 kbit/s.
\end{tcolorbox}
\vspace{2mm}\hrule height 0.2pt \vspace{3mm}

% DB question_id: a42b246a-1d6d-57f0-b336-6f042f66911b
\needspace{5\baselineskip}
\noindent\textbf{\color{ink}759.}
\par
48000 Hz, 16 bit, 1 kanalli PCM audioning xom bitrate'i qancha?
\begin{enumerate}[label=\Alph*),leftmargin=9mm]
\item 768 000 kbit/s
\item 96 000 kbit/s
\item 1 536 000 kbit/s
\item 768 kbit/s
\end{enumerate}
\begin{tcolorbox}[enhanced,breakable,colback=teal!4,colframe=teal!50!black,boxrule=0.5pt,arc=1.1mm,left=3mm,right=3mm,top=1.7mm,bottom=1.7mm,before skip=2mm,after skip=2mm]
\textbf{\color{teal}To'g'ri javob:} \textbf{D) 768 kbit/s}
\par\smallskip
\textbf{Izoh:} 48000\ensuremath{\times}16\ensuremath{\times}1 = 768000 bit/s = 768 kbit/s.
\end{tcolorbox}
\vspace{2mm}\hrule height 0.2pt \vspace{3mm}

% DB question_id: a46f2ea6-fa94-5b1e-9ebb-63edefb21291
\needspace{5\baselineskip}
\noindent\textbf{\color{ink}760.}
\par
96000 Hz, 24 bit, 2 kanal, 5 soniyalik siqilmagan audio hajmini MiB da toping.
\begin{enumerate}[label=\Alph*),leftmargin=9mm]
\item 2,747 MiB
\item 23 040 000 MiB
\item 1 440 000 MiB
\item 576 000 MiB
\end{enumerate}
\begin{tcolorbox}[enhanced,breakable,colback=teal!4,colframe=teal!50!black,boxrule=0.5pt,arc=1.1mm,left=3mm,right=3mm,top=1.7mm,bottom=1.7mm,before skip=2mm,after skip=2mm]
\textbf{\color{teal}To'g'ri javob:} \textbf{A) 2,747 MiB}
\par\smallskip
\textbf{Izoh:} Hajm = 96000\ensuremath{\times}24\ensuremath{\times}2\ensuremath{\times}5/8 = 2 880 000 B = 2,747 MiB.
\end{tcolorbox}
\vspace{2mm}\hrule height 0.2pt \vspace{3mm}

% DB question_id: b28d4861-8fc3-5dcc-b993-0104093441d5
\needspace{5\baselineskip}
\noindent\textbf{\color{ink}761.}
\par
RGB modeli qayerda tabiiy qo'llanadi?
\begin{enumerate}[label=\Alph*),leftmargin=9mm]
\item Faqat bosma siyohda
\item Yorug'lik chiqaruvchi ekranlarda
\item Audio yozuvda
\item Matn kodlashda
\end{enumerate}
\begin{tcolorbox}[enhanced,breakable,colback=teal!4,colframe=teal!50!black,boxrule=0.5pt,arc=1.1mm,left=3mm,right=3mm,top=1.7mm,bottom=1.7mm,before skip=2mm,after skip=2mm]
\textbf{\color{teal}To'g'ri javob:} \textbf{B) Yorug'lik chiqaruvchi ekranlarda}
\par\smallskip
\textbf{Izoh:} RGB qo'shiluvchi yorug'lik modelidir.
\end{tcolorbox}
\vspace{2mm}\hrule height 0.2pt \vspace{3mm}

% DB question_id: b36da9f3-83f6-5672-91d6-ea86eac92262
\needspace{5\baselineskip}
\noindent\textbf{\color{ink}762.}
\par
Video bitrate'i 12 Mbit/s, davomiyligi 300 s. Hajm qancha? MB = 10\textsuperscript{6} B.
\begin{enumerate}[label=\Alph*),leftmargin=9mm]
\item 3 600 MB
\item 56,25 MB
\item 450 MB
\item 453 MB
\end{enumerate}
\begin{tcolorbox}[enhanced,breakable,colback=teal!4,colframe=teal!50!black,boxrule=0.5pt,arc=1.1mm,left=3mm,right=3mm,top=1.7mm,bottom=1.7mm,before skip=2mm,after skip=2mm]
\textbf{\color{teal}To'g'ri javob:} \textbf{C) 450 MB}
\par\smallskip
\textbf{Izoh:} 12\ensuremath{\times}300/8 = 450 MB.
\end{tcolorbox}
\vspace{2mm}\hrule height 0.2pt \vspace{3mm}

% DB question_id: b47b184d-a1b8-57a4-b2a1-76d0c72456a1
\needspace{5\baselineskip}
\noindent\textbf{\color{ink}763.}
\par
Matnda 2000 ta bir baytli va 500 ta 3 baytli belgi bor. Umumiy hajm qancha?
\begin{enumerate}[label=\Alph*),leftmargin=9mm]
\item 3 500 B
\item 2 500 B
\item 28 000 B
\item 437,5 B
\end{enumerate}
\begin{tcolorbox}[enhanced,breakable,colback=teal!4,colframe=teal!50!black,boxrule=0.5pt,arc=1.1mm,left=3mm,right=3mm,top=1.7mm,bottom=1.7mm,before skip=2mm,after skip=2mm]
\textbf{\color{teal}To'g'ri javob:} \textbf{A) 3 500 B}
\par\smallskip
\textbf{Izoh:} 2000\ensuremath{\times}1 + 500\ensuremath{\times}3 = 3500 B = 3 500 B.
\end{tcolorbox}
\vspace{2mm}\hrule height 0.2pt \vspace{3mm}

% DB question_id: b5b42306-0b6b-554f-84ab-ddfee4c57a2a
\needspace{5\baselineskip}
\noindent\textbf{\color{ink}764.}
\par
320 kbit/s bitratedagi audio 240 soniya davom etadi. Hajm qancha? MB = 10\textsuperscript{6} B.
\begin{enumerate}[label=\Alph*),leftmargin=9mm]
\item 76,8 MB
\item 1,2 MB
\item 12,6 MB
\item 9,6 MB
\end{enumerate}
\begin{tcolorbox}[enhanced,breakable,colback=teal!4,colframe=teal!50!black,boxrule=0.5pt,arc=1.1mm,left=3mm,right=3mm,top=1.7mm,bottom=1.7mm,before skip=2mm,after skip=2mm]
\textbf{\color{teal}To'g'ri javob:} \textbf{D) 9,6 MB}
\par\smallskip
\textbf{Izoh:} 320000\ensuremath{\times}240/8 = 9,6 MB.
\end{tcolorbox}
\vspace{2mm}\hrule height 0.2pt \vspace{3mm}

% DB question_id: b6b8c762-6860-565b-ba55-aaa1e4a0f4fb
\needspace{5\baselineskip}
\noindent\textbf{\color{ink}765.}
\par
Matnda 1200 ta bir baytli va 300 ta 4 baytli belgi bor. Umumiy hajm qancha?
\begin{enumerate}[label=\Alph*),leftmargin=9mm]
\item 1 500 B
\item 2 400 B
\item 19 200 B
\item 300 B
\end{enumerate}
\begin{tcolorbox}[enhanced,breakable,colback=teal!4,colframe=teal!50!black,boxrule=0.5pt,arc=1.1mm,left=3mm,right=3mm,top=1.7mm,bottom=1.7mm,before skip=2mm,after skip=2mm]
\textbf{\color{teal}To'g'ri javob:} \textbf{B) 2 400 B}
\par\smallskip
\textbf{Izoh:} 1200\ensuremath{\times}1 + 300\ensuremath{\times}4 = 2400 B = 2 400 B.
\end{tcolorbox}
\vspace{2mm}\hrule height 0.2pt \vspace{3mm}

% DB question_id: bad14812-f1d3-5016-80eb-5b6987beda23
\needspace{5\baselineskip}
\noindent\textbf{\color{ink}766.}
\par
Kod nuqtasi nima?
\begin{enumerate}[label=\Alph*),leftmargin=9mm]
\item Faylning baytlar soni
\item Unicode tizimida belgiga berilgan abstrakt raqam
\item Piksel manzili
\item Audio namuna
\end{enumerate}
\begin{tcolorbox}[enhanced,breakable,colback=teal!4,colframe=teal!50!black,boxrule=0.5pt,arc=1.1mm,left=3mm,right=3mm,top=1.7mm,bottom=1.7mm,before skip=2mm,after skip=2mm]
\textbf{\color{teal}To'g'ri javob:} \textbf{B) Unicode tizimida belgiga berilgan abstrakt raqam}
\par\smallskip
\textbf{Izoh:} Kod nuqtasi kodlash shaklidan mustaqil belgilar identifikatoridir.
\end{tcolorbox}
\vspace{2mm}\hrule height 0.2pt \vspace{3mm}

% DB question_id: bbfd2a95-4001-5b74-a2ba-1a3d08611ce8
\needspace{5\baselineskip}
\noindent\textbf{\color{ink}767.}
\par
4 bit rang chuqurligi nazariy jihatdan nechta rangni ifodalaydi?
\begin{enumerate}[label=\Alph*),leftmargin=9mm]
\item 4 rang
\item 16 rang
\item 8 rang
\item 19 rang
\end{enumerate}
\begin{tcolorbox}[enhanced,breakable,colback=teal!4,colframe=teal!50!black,boxrule=0.5pt,arc=1.1mm,left=3mm,right=3mm,top=1.7mm,bottom=1.7mm,before skip=2mm,after skip=2mm]
\textbf{\color{teal}To'g'ri javob:} \textbf{B) 16 rang}
\par\smallskip
\textbf{Izoh:} Ranglar soni = 2\textasciicircum{}4 = 16.
\end{tcolorbox}
\vspace{2mm}\hrule height 0.2pt \vspace{3mm}

% DB question_id: bdf71bec-6574-5b31-9f20-fd413346484a
\needspace{5\baselineskip}
\noindent\textbf{\color{ink}768.}
\par
8 bit rang chuqurligi nazariy jihatdan nechta rangni ifodalaydi?
\begin{enumerate}[label=\Alph*),leftmargin=9mm]
\item 8 rang
\item 16 rang
\item 256 rang
\item 128 rang
\end{enumerate}
\begin{tcolorbox}[enhanced,breakable,colback=teal!4,colframe=teal!50!black,boxrule=0.5pt,arc=1.1mm,left=3mm,right=3mm,top=1.7mm,bottom=1.7mm,before skip=2mm,after skip=2mm]
\textbf{\color{teal}To'g'ri javob:} \textbf{C) 256 rang}
\par\smallskip
\textbf{Izoh:} Ranglar soni = 2\textasciicircum{}8 = 256.
\end{tcolorbox}
\vspace{2mm}\hrule height 0.2pt \vspace{3mm}

% DB question_id: c0937ecf-c007-5982-b79e-b1d61accd3df
\needspace{5\baselineskip}
\noindent\textbf{\color{ink}769.}
\par
48000 Hz, 24 bit, 2 kanal, 30 soniyalik siqilmagan audio hajmini MiB da toping.
\begin{enumerate}[label=\Alph*),leftmargin=9mm]
\item 69 120 000 MiB
\item 4 320 000 MiB
\item 288 000 MiB
\item 8,24 MiB
\end{enumerate}
\begin{tcolorbox}[enhanced,breakable,colback=teal!4,colframe=teal!50!black,boxrule=0.5pt,arc=1.1mm,left=3mm,right=3mm,top=1.7mm,bottom=1.7mm,before skip=2mm,after skip=2mm]
\textbf{\color{teal}To'g'ri javob:} \textbf{D) 8,24 MiB}
\par\smallskip
\textbf{Izoh:} Hajm = 48000\ensuremath{\times}24\ensuremath{\times}2\ensuremath{\times}30/8 = 8 640 000 B = 8,24 MiB.
\end{tcolorbox}
\vspace{2mm}\hrule height 0.2pt \vspace{3mm}

% DB question_id: c208d7d8-1d1e-5bc2-8a4c-eac9a6e2f65a
\needspace{5\baselineskip}
\noindent\textbf{\color{ink}770.}
\par
1280\ensuremath{\times}720 pikselli, 32 bit rang chuqurligidagi siqilmagan rastr tasvir hajmini MiB da toping.
\begin{enumerate}[label=\Alph*),leftmargin=9mm]
\item 921 600 MiB
\item 29 491 200 MiB
\item 6,516 MiB
\item 3,516 MiB
\end{enumerate}
\begin{tcolorbox}[enhanced,breakable,colback=teal!4,colframe=teal!50!black,boxrule=0.5pt,arc=1.1mm,left=3mm,right=3mm,top=1.7mm,bottom=1.7mm,before skip=2mm,after skip=2mm]
\textbf{\color{teal}To'g'ri javob:} \textbf{D) 3,516 MiB}
\par\smallskip
\textbf{Izoh:} Hajm = 1280\ensuremath{\times}720\ensuremath{\times}32/8 = 3 686 400 B = 3,516 MiB.
\end{tcolorbox}
\vspace{2mm}\hrule height 0.2pt \vspace{3mm}

% DB question_id: c3b7261a-f28c-5ca0-b196-fd59fc8bde62
\needspace{5\baselineskip}
\noindent\textbf{\color{ink}771.}
\par
256 kbit/s bitratedagi audio 600 soniya davom etadi. Hajm qancha? MB = 10\textsuperscript{6} B.
\begin{enumerate}[label=\Alph*),leftmargin=9mm]
\item 19,2 MB
\item 153,6 MB
\item 2,4 MB
\item 22,2 MB
\end{enumerate}
\begin{tcolorbox}[enhanced,breakable,colback=teal!4,colframe=teal!50!black,boxrule=0.5pt,arc=1.1mm,left=3mm,right=3mm,top=1.7mm,bottom=1.7mm,before skip=2mm,after skip=2mm]
\textbf{\color{teal}To'g'ri javob:} \textbf{A) 19,2 MB}
\par\smallskip
\textbf{Izoh:} 256000\ensuremath{\times}600/8 = 19,2 MB.
\end{tcolorbox}
\vspace{2mm}\hrule height 0.2pt \vspace{3mm}

% DB question_id: c44438d1-1d41-5820-bce7-56e02104b370
\needspace{5\baselineskip}
\noindent\textbf{\color{ink}772.}
\par
24 FPS videoda 90 soniyada nechta kadr bo'ladi?
\begin{enumerate}[label=\Alph*),leftmargin=9mm]
\item 114 kadr
\item 2 160 kadr
\item 129 600 kadr
\item 1 080 kadr
\end{enumerate}
\begin{tcolorbox}[enhanced,breakable,colback=teal!4,colframe=teal!50!black,boxrule=0.5pt,arc=1.1mm,left=3mm,right=3mm,top=1.7mm,bottom=1.7mm,before skip=2mm,after skip=2mm]
\textbf{\color{teal}To'g'ri javob:} \textbf{B) 2 160 kadr}
\par\smallskip
\textbf{Izoh:} Kadrlar soni = 24\ensuremath{\times}90 = 2160.
\end{tcolorbox}
\vspace{2mm}\hrule height 0.2pt \vspace{3mm}

% DB question_id: c4fe6349-8380-45f3-a933-6c3e74c4c81c
\needspace{5\baselineskip}
\noindent\textbf{\color{ink}773.}
\par
Hujjat nomi .pdf bilan tugadi, ammo mazmuni baribir tekshirildi. Bu vaziyatdagi asosiy tushunchani aniqlang.
\begin{enumerate}[label=\Alph*),leftmargin=9mm]
\item Fayl kengaytmasi
\item Siqish
\item Yo'qotishsiz siqish
\item Yo'qotishli siqish
\end{enumerate}
\begin{tcolorbox}[enhanced,breakable,colback=teal!4,colframe=teal!50!black,boxrule=0.5pt,arc=1.1mm,left=3mm,right=3mm,top=1.7mm,bottom=1.7mm,before skip=2mm,after skip=2mm]
\textbf{\color{teal}To'g'ri javob:} \textbf{A) Fayl kengaytmasi}
\par\smallskip
\textbf{Izoh:} Vaziyat Fayl kengaytmasiga misol: fayl nomining formatni taxmin qilishga yordam beradigan oxirgi qismi.
\end{tcolorbox}
\vspace{2mm}\hrule height 0.2pt \vspace{3mm}

% DB question_id: ca7bb483-ab67-4cd0-a703-7f67334d823c
\needspace{5\baselineskip}
\noindent\textbf{\color{ink}774.}
\par
Dastur kodi arxivdan hech bir belgi o'zgarmasdan qaytdi. Bu vaziyatdagi asosiy tushunchani aniqlang.
\begin{enumerate}[label=\Alph*),leftmargin=9mm]
\item Yo'qotishli siqish
\item RLE
\item Yo'qotishsiz siqish
\item Siqish koeffitsiyenti
\end{enumerate}
\begin{tcolorbox}[enhanced,breakable,colback=teal!4,colframe=teal!50!black,boxrule=0.5pt,arc=1.1mm,left=3mm,right=3mm,top=1.7mm,bottom=1.7mm,before skip=2mm,after skip=2mm]
\textbf{\color{teal}To'g'ri javob:} \textbf{C) Yo'qotishsiz siqish}
\par\smallskip
\textbf{Izoh:} Vaziyat Yo'qotishsiz siqishga misol: dekodlanganda asl ma'lumotni bitma-bit tiklaydigan siqish.
\end{tcolorbox}
\vspace{2mm}\hrule height 0.2pt \vspace{3mm}

% DB question_id: ccc072fb-2821-50a3-9897-baad4e5cf8b1
\needspace{5\baselineskip}
\noindent\textbf{\color{ink}775.}
\par
96 kbit/s bitratedagi audio 900 soniya davom etadi. Hajm qancha? MB = 10\textsuperscript{6} B.
\begin{enumerate}[label=\Alph*),leftmargin=9mm]
\item 86,4 MB
\item 10,8 MB
\item 1,35 MB
\item 13,8 MB
\end{enumerate}
\begin{tcolorbox}[enhanced,breakable,colback=teal!4,colframe=teal!50!black,boxrule=0.5pt,arc=1.1mm,left=3mm,right=3mm,top=1.7mm,bottom=1.7mm,before skip=2mm,after skip=2mm]
\textbf{\color{teal}To'g'ri javob:} \textbf{B) 10,8 MB}
\par\smallskip
\textbf{Izoh:} 96000\ensuremath{\times}900/8 = 10,8 MB.
\end{tcolorbox}
\vspace{2mm}\hrule height 0.2pt \vspace{3mm}

% DB question_id: cd65300a-dde6-5d73-9e57-19eebd90cc14
\needspace{5\baselineskip}
\noindent\textbf{\color{ink}776.}
\par
640\ensuremath{\times}480 pikselli, 24 bit rang chuqurligidagi siqilmagan rastr tasvir hajmini B da toping.
\begin{enumerate}[label=\Alph*),leftmargin=9mm]
\item 921 600 B
\item 307 200 B
\item 7 372 800 B
\item 921 603 B
\end{enumerate}
\begin{tcolorbox}[enhanced,breakable,colback=teal!4,colframe=teal!50!black,boxrule=0.5pt,arc=1.1mm,left=3mm,right=3mm,top=1.7mm,bottom=1.7mm,before skip=2mm,after skip=2mm]
\textbf{\color{teal}To'g'ri javob:} \textbf{A) 921 600 B}
\par\smallskip
\textbf{Izoh:} Hajm = 640\ensuremath{\times}480\ensuremath{\times}24/8 = 921 600 B = 921 600 B.
\end{tcolorbox}
\vspace{2mm}\hrule height 0.2pt \vspace{3mm}

% DB question_id: cda9944c-c410-5379-8834-55f97682a7d8
\needspace{5\baselineskip}
\noindent\textbf{\color{ink}777.}
\par
100\ensuremath{\times}200 pikselli, 8 bit rang chuqurligidagi siqilmagan rastr tasvir hajmini B da toping.
\begin{enumerate}[label=\Alph*),leftmargin=9mm]
\item 160 000 B
\item 20 002 B
\item 20 003 B
\item 20 000 B
\end{enumerate}
\begin{tcolorbox}[enhanced,breakable,colback=teal!4,colframe=teal!50!black,boxrule=0.5pt,arc=1.1mm,left=3mm,right=3mm,top=1.7mm,bottom=1.7mm,before skip=2mm,after skip=2mm]
\textbf{\color{teal}To'g'ri javob:} \textbf{D) 20 000 B}
\par\smallskip
\textbf{Izoh:} Hajm = 100\ensuremath{\times}200\ensuremath{\times}8/8 = 20 000 B = 20 000 B.
\end{tcolorbox}
\vspace{2mm}\hrule height 0.2pt \vspace{3mm}

% DB question_id: cf977475-ab69-587a-b3dc-c38be2d723d2
\needspace{5\baselineskip}
\noindent\textbf{\color{ink}778.}
\par
Bitrate nima?
\begin{enumerate}[label=\Alph*),leftmargin=9mm]
\item Fayldagi belgilar soni
\item Bir soniyada saqlanadigan yoki uzatiladigan bitlar miqdori
\item Rasm eni
\item Kod nuqtasi
\end{enumerate}
\begin{tcolorbox}[enhanced,breakable,colback=teal!4,colframe=teal!50!black,boxrule=0.5pt,arc=1.1mm,left=3mm,right=3mm,top=1.7mm,bottom=1.7mm,before skip=2mm,after skip=2mm]
\textbf{\color{teal}To'g'ri javob:} \textbf{B) Bir soniyada saqlanadigan yoki uzatiladigan bitlar miqdori}
\par\smallskip
\textbf{Izoh:} Bitrate media oqim hajmini aniqlaydi.
\end{tcolorbox}
\vspace{2mm}\hrule height 0.2pt \vspace{3mm}

% DB question_id: cf9ce1ee-d6e4-5fde-8856-26934b32e61a
\needspace{5\baselineskip}
\noindent\textbf{\color{ink}779.}
\par
30 FPS videoda 120 soniyada nechta kadr bo'ladi?
\begin{enumerate}[label=\Alph*),leftmargin=9mm]
\item 150 kadr
\item 216 000 kadr
\item 1 800 kadr
\item 3 600 kadr
\end{enumerate}
\begin{tcolorbox}[enhanced,breakable,colback=teal!4,colframe=teal!50!black,boxrule=0.5pt,arc=1.1mm,left=3mm,right=3mm,top=1.7mm,bottom=1.7mm,before skip=2mm,after skip=2mm]
\textbf{\color{teal}To'g'ri javob:} \textbf{D) 3 600 kadr}
\par\smallskip
\textbf{Izoh:} Kadrlar soni = 30\ensuremath{\times}120 = 3600.
\end{tcolorbox}
\vspace{2mm}\hrule height 0.2pt \vspace{3mm}

% DB question_id: d38d2bca-006d-4388-8da1-0347bec01a1c
\needspace{5\baselineskip}
\noindent\textbf{\color{ink}780.}
\par
MP4 faylida video va ikkita audio yo'li saqlandi. Vaziyatni tahlil qilib, unga eng aniq mos tushunchani toping.
\begin{enumerate}[label=\Alph*),leftmargin=9mm]
\item Media oqimi
\item Konteyner
\item Multimedia
\item Kadr
\end{enumerate}
\begin{tcolorbox}[enhanced,breakable,colback=teal!4,colframe=teal!50!black,boxrule=0.5pt,arc=1.1mm,left=3mm,right=3mm,top=1.7mm,bottom=1.7mm,before skip=2mm,after skip=2mm]
\textbf{\color{teal}To'g'ri javob:} \textbf{B) Konteyner}
\par\smallskip
\textbf{Izoh:} Vaziyat Konteynerga misol: video, audio, subtitr va metama'lumot oqimlarini bir faylda birlashtiradigan tuzilma.
\end{tcolorbox}
\vspace{2mm}\hrule height 0.2pt \vspace{3mm}

% DB question_id: d3bef342-9e3e-594c-bc3b-cc61c817696b
\needspace{5\baselineskip}
\noindent\textbf{\color{ink}781.}
\par
16 bit rang chuqurligi nazariy jihatdan nechta rangni ifodalaydi?
\begin{enumerate}[label=\Alph*),leftmargin=9mm]
\item 16 rang
\item 32 rang
\item 32 768 rang
\item 65 536 rang
\end{enumerate}
\begin{tcolorbox}[enhanced,breakable,colback=teal!4,colframe=teal!50!black,boxrule=0.5pt,arc=1.1mm,left=3mm,right=3mm,top=1.7mm,bottom=1.7mm,before skip=2mm,after skip=2mm]
\textbf{\color{teal}To'g'ri javob:} \textbf{D) 65 536 rang}
\par\smallskip
\textbf{Izoh:} Ranglar soni = 2\textasciicircum{}16 = 65 536.
\end{tcolorbox}
\vspace{2mm}\hrule height 0.2pt \vspace{3mm}

% DB question_id: d40258b3-70ee-5bbc-8264-c44054cd5f50
\needspace{5\baselineskip}
\noindent\textbf{\color{ink}782.}
\par
Video 20 Mbit/s, audio 256 kbit/s, davomiylik 90 s. Umumiy media hajmi qancha? MB = 10\textsuperscript{6} B.
\begin{enumerate}[label=\Alph*),leftmargin=9mm]
\item 225 MB
\item 2,88 MB
\item 1823,04 MB
\item 227,88 MB
\end{enumerate}
\begin{tcolorbox}[enhanced,breakable,colback=teal!4,colframe=teal!50!black,boxrule=0.5pt,arc=1.1mm,left=3mm,right=3mm,top=1.7mm,bottom=1.7mm,before skip=2mm,after skip=2mm]
\textbf{\color{teal}To'g'ri javob:} \textbf{D) 227,88 MB}
\par\smallskip
\textbf{Izoh:} Jami bitrate = 20.256 Mbit/s; hajm = jami\ensuremath{\times}90/8 = 227,88 MB.
\end{tcolorbox}
\vspace{2mm}\hrule height 0.2pt \vspace{3mm}

% DB question_id: d4636aa1-dcf9-592a-996f-df10d5d2fdf2
\needspace{5\baselineskip}
\noindent\textbf{\color{ink}783.}
\par
Matnda 2500 ta belgi bor, har biri 4 bayt bilan kodlangan. Hajmni KiB da toping.
\begin{enumerate}[label=\Alph*),leftmargin=9mm]
\item 2 500 KiB
\item 80 000 KiB
\item 9,766 KiB
\item 1 250 KiB
\end{enumerate}
\begin{tcolorbox}[enhanced,breakable,colback=teal!4,colframe=teal!50!black,boxrule=0.5pt,arc=1.1mm,left=3mm,right=3mm,top=1.7mm,bottom=1.7mm,before skip=2mm,after skip=2mm]
\textbf{\color{teal}To'g'ri javob:} \textbf{C) 9,766 KiB}
\par\smallskip
\textbf{Izoh:} 2500 \ensuremath{\times} 4 = 10000 B = 9,766 KiB.
\end{tcolorbox}
\vspace{2mm}\hrule height 0.2pt \vspace{3mm}

% DB question_id: d4a72e86-690d-45d9-84a3-286ddf3d57e2
\needspace{5\baselineskip}
\noindent\textbf{\color{ink}784.}
\par
Rangli tasvirni kodlashda RGB tizimida har bir rang kanali 8 bit bilan kodlansa, bitta piksel nechta bit yoki nechta bayt bilan ifodalanadi hamda nechta xil rang hosil qilish mumkin?
\begin{enumerate}[label=\Alph*),leftmargin=9mm]
\item 32 bit, 4 bayt, 4 294 967 296 rang
\item 16 bit, 2 bayt, 65 536 rang
\item 24 bit, 3 bayt, 16 777 216 rang
\item 24 bit, 4 bayt, 16 777 216 rang
\end{enumerate}
\begin{tcolorbox}[enhanced,breakable,colback=teal!4,colframe=teal!50!black,boxrule=0.5pt,arc=1.1mm,left=3mm,right=3mm,top=1.7mm,bottom=1.7mm,before skip=2mm,after skip=2mm]
\textbf{\color{teal}To'g'ri javob:} \textbf{C) 24 bit, 3 bayt, 16 777 216 rang}
\par\smallskip
\textbf{Izoh:} 3 ta kanal \ensuremath{\times} 8 bit = 24 bit = 3 bayt; ranglar soni 2\textsuperscript{2}\textsuperscript{4} = 16 777 216.
\end{tcolorbox}
\vspace{2mm}\hrule height 0.2pt \vspace{3mm}

% DB question_id: ddd0689b-28a0-4048-8f64-238ab383ed9f
\needspace{5\baselineskip}
\noindent\textbf{\color{ink}785.}
\par
Unicode nima?
\begin{enumerate}[label=\Alph*),leftmargin=9mm]
\item Belgilarni sakkiz bitli kodlar bilan cheklab saqlaydigan jadval tizimi
\item Turli yozuv tizimlari belgilariga yagona kod nuqtalari ajratadigan standart
\item Har bir matn belgisini bitta baytda saqlashni belgilaydigan kodlash
\item Belgilar ko'rinishini shriftning grafik shakllari orqali aniqlaydigan format
\end{enumerate}
\begin{tcolorbox}[enhanced,breakable,colback=teal!4,colframe=teal!50!black,boxrule=0.5pt,arc=1.1mm,left=3mm,right=3mm,top=1.7mm,bottom=1.7mm,before skip=2mm,after skip=2mm]
\textbf{\color{teal}To'g'ri javob:} \textbf{B) Turli yozuv tizimlari belgilariga yagona kod nuqtalari ajratadigan standart}
\par\smallskip
\textbf{Izoh:} Unicode turli yozuv tizimlaridagi belgilarga noyob kod nuqtalari ajratadi. Kod nuqtasining mavjudligi uni xotirada albatta bir bayt egallaydi degani emas: UTF-8 o'zgaruvchan uzunlikda kodlaydi. Shrift esa belgi shaklini, kod nuqtasini emas, tasvirlaydi.
\end{tcolorbox}
\vspace{2mm}\hrule height 0.2pt \vspace{3mm}

% DB question_id: de23129d-5031-53eb-b21c-794b3123ef6c
\needspace{5\baselineskip}
\noindent\textbf{\color{ink}786.}
\par
4096 baytli matnda har belgi 2 bayt bo'lsa, nechta belgi bor?
\begin{enumerate}[label=\Alph*),leftmargin=9mm]
\item 4 096 belgi
\item 16 384 belgi
\item 2 048 belgi
\item 1 024 belgi
\end{enumerate}
\begin{tcolorbox}[enhanced,breakable,colback=teal!4,colframe=teal!50!black,boxrule=0.5pt,arc=1.1mm,left=3mm,right=3mm,top=1.7mm,bottom=1.7mm,before skip=2mm,after skip=2mm]
\textbf{\color{teal}To'g'ri javob:} \textbf{C) 2 048 belgi}
\par\smallskip
\textbf{Izoh:} Belgilar soni = 4096/2 = 2048.
\end{tcolorbox}
\vspace{2mm}\hrule height 0.2pt \vspace{3mm}

% DB question_id: e1134e44-c226-4b8e-8175-bb3c3751e374
\needspace{5\baselineskip}
\noindent\textbf{\color{ink}787.}
\par
HSV/HSL modelining qulayligi nimada?
\begin{enumerate}[label=\Alph*),leftmargin=9mm]
\item Rangni qizil, yashil va ko'k kanallar bo'yicha aralashtirib hisoblash
\item Tasvirdagi piksellar sonini kenglik va balandlik bo'yicha o'zgartirish
\item Bosma ranglarni siyoh miqdori va qog'oz turiga moslab ajratish
\item Rang tusini, to'yinganlikni va yorug'likni intuitiv boshqarish
\end{enumerate}
\begin{tcolorbox}[enhanced,breakable,colback=teal!4,colframe=teal!50!black,boxrule=0.5pt,arc=1.1mm,left=3mm,right=3mm,top=1.7mm,bottom=1.7mm,before skip=2mm,after skip=2mm]
\textbf{\color{teal}To'g'ri javob:} \textbf{D) Rang tusini, to'yinganlikni va yorug'likni intuitiv boshqarish}
\par\smallskip
\textbf{Izoh:} HSV/HSL modellari rang tusi (Hue), to'yinganlik (Saturation) hamda yorug'lik yoki qiymat parametrini alohida boshqarishga qulay. RGB rang kanallarining qo'shilishini, CMYK esa bosma bo'yoqlarini tavsiflaydi. Piksel o'lchamini o'zgartirish rang modelining vazifasi emas.
\end{tcolorbox}
\vspace{2mm}\hrule height 0.2pt \vspace{3mm}

% DB question_id: e6db2608-b89c-5a0f-9466-f9f2da8e075c
\needspace{5\baselineskip}
\noindent\textbf{\color{ink}788.}
\par
Sampling rate yarimlashib, kanal soni 2 marta oshsa, hajm qanday o'zgaradi?
\begin{enumerate}[label=\Alph*),leftmargin=9mm]
\item 2 marta oshadi
\item O'zgarmaydi
\item Yarimlashadi
\item 4 marta oshadi
\end{enumerate}
\begin{tcolorbox}[enhanced,breakable,colback=teal!4,colframe=teal!50!black,boxrule=0.5pt,arc=1.1mm,left=3mm,right=3mm,top=1.7mm,bottom=1.7mm,before skip=2mm,after skip=2mm]
\textbf{\color{teal}To'g'ri javob:} \textbf{B) O'zgarmaydi}
\par\smallskip
\textbf{Izoh:} 0,5\ensuremath{\times}2 = 1.
\end{tcolorbox}
\vspace{2mm}\hrule height 0.2pt \vspace{3mm}

% DB question_id: e7b9b3df-7609-4e6a-85a6-b1bf1c7a941b
\needspace{5\baselineskip}
\noindent\textbf{\color{ink}789.}
\par
Oddiy lotin harfi bir bayt, ayrim boshqa belgilar ko'proq bayt egalladi. Vaziyatni tahlil qilib, unga eng aniq mos tushunchani toping.
\begin{enumerate}[label=\Alph*),leftmargin=9mm]
\item UTF-16
\item Matn metama'lumoti
\item UTF-8
\item Belgini kodlash
\end{enumerate}
\begin{tcolorbox}[enhanced,breakable,colback=teal!4,colframe=teal!50!black,boxrule=0.5pt,arc=1.1mm,left=3mm,right=3mm,top=1.7mm,bottom=1.7mm,before skip=2mm,after skip=2mm]
\textbf{\color{teal}To'g'ri javob:} \textbf{C) UTF-8}
\par\smallskip
\textbf{Izoh:} Vaziyat UTF-8ga misol: Unicode kod nuqtalarini 1-4 baytli ketma-ketliklarda kodlaydigan shakl.
\end{tcolorbox}
\vspace{2mm}\hrule height 0.2pt \vspace{3mm}

% DB question_id: e7f0f485-4eda-5e11-a7ae-7d9cad1d8a5d
\needspace{5\baselineskip}
\noindent\textbf{\color{ink}790.}
\par
640\ensuremath{\times}480, 24 bit, 30 FPS, 1 soniyalik siqilmagan video hajmini MiB da toping.
\begin{enumerate}[label=\Alph*),leftmargin=9mm]
\item 26,367 MiB
\item 3,296 MiB
\item 0,879 MiB
\item 791,016 MiB
\end{enumerate}
\begin{tcolorbox}[enhanced,breakable,colback=teal!4,colframe=teal!50!black,boxrule=0.5pt,arc=1.1mm,left=3mm,right=3mm,top=1.7mm,bottom=1.7mm,before skip=2mm,after skip=2mm]
\textbf{\color{teal}To'g'ri javob:} \textbf{A) 26,367 MiB}
\par\smallskip
\textbf{Izoh:} Hajm = W\ensuremath{\times}H\ensuremath{\times}d/8\ensuremath{\times}FPS\ensuremath{\times}t = 26,367 MiB.
\end{tcolorbox}
\vspace{2mm}\hrule height 0.2pt \vspace{3mm}

% DB question_id: ea382a64-3e2e-5155-bc8e-6afc68f90b79
\needspace{5\baselineskip}
\noindent\textbf{\color{ink}791.}
\par
44100 Hz, 16 bit, 2 kanal, 10 soniyalik siqilmagan audio hajmini MiB da toping.
\begin{enumerate}[label=\Alph*),leftmargin=9mm]
\item 14 112 000 MiB
\item 882 000 MiB
\item 1,682 MiB
\item 176 400 MiB
\end{enumerate}
\begin{tcolorbox}[enhanced,breakable,colback=teal!4,colframe=teal!50!black,boxrule=0.5pt,arc=1.1mm,left=3mm,right=3mm,top=1.7mm,bottom=1.7mm,before skip=2mm,after skip=2mm]
\textbf{\color{teal}To'g'ri javob:} \textbf{C) 1,682 MiB}
\par\smallskip
\textbf{Izoh:} Hajm = 44100\ensuremath{\times}16\ensuremath{\times}2\ensuremath{\times}10/8 = 1 764 000 B = 1,682 MiB.
\end{tcolorbox}
\vspace{2mm}\hrule height 0.2pt \vspace{3mm}

% DB question_id: ea71da76-715e-4390-a5ce-2c21197cc496
\needspace{5\baselineskip}
\noindent\textbf{\color{ink}792.}
\par
U+0041 qiymati A belgisini ifodaladi. Vaziyatni tahlil qilib, unga eng aniq mos tushunchani toping.
\begin{enumerate}[label=\Alph*),leftmargin=9mm]
\item UTF-8
\item Kod nuqtasi
\item UTF-16
\item Matn metama'lumoti
\end{enumerate}
\begin{tcolorbox}[enhanced,breakable,colback=teal!4,colframe=teal!50!black,boxrule=0.5pt,arc=1.1mm,left=3mm,right=3mm,top=1.7mm,bottom=1.7mm,before skip=2mm,after skip=2mm]
\textbf{\color{teal}To'g'ri javob:} \textbf{B) Kod nuqtasi}
\par\smallskip
\textbf{Izoh:} Vaziyat Kod nuqtasiga misol: Unicode tizimida belgiga berilgan abstrakt raqam.
\end{tcolorbox}
\vspace{2mm}\hrule height 0.2pt \vspace{3mm}

% DB question_id: eabdd7f7-5c37-5a0f-9e85-4f405f301b76
\needspace{5\baselineskip}
\noindent\textbf{\color{ink}793.}
\par
UTF-16ning xususiyati qaysi?
\begin{enumerate}[label=\Alph*),leftmargin=9mm]
\item Har doim 1 bayt
\item Faqat ASCII
\item Rang modeli
\item Kod nuqtalarini odatda 2 yoki 4 baytda ifodalaydi
\end{enumerate}
\begin{tcolorbox}[enhanced,breakable,colback=teal!4,colframe=teal!50!black,boxrule=0.5pt,arc=1.1mm,left=3mm,right=3mm,top=1.7mm,bottom=1.7mm,before skip=2mm,after skip=2mm]
\textbf{\color{teal}To'g'ri javob:} \textbf{D) Kod nuqtalarini odatda 2 yoki 4 baytda ifodalaydi}
\par\smallskip
\textbf{Izoh:} UTF-16 kod birliklari 16 bitli.
\end{tcolorbox}
\vspace{2mm}\hrule height 0.2pt \vspace{3mm}

% DB question_id: ed7b11a8-5629-5ac8-8089-6691aba86e83
\needspace{5\baselineskip}
\noindent\textbf{\color{ink}794.}
\par
12000 baytli matnda har belgi 4 bayt bo'lsa, nechta belgi bor?
\begin{enumerate}[label=\Alph*),leftmargin=9mm]
\item 12 000 belgi
\item 24 000 belgi
\item 1 500 belgi
\item 3 000 belgi
\end{enumerate}
\begin{tcolorbox}[enhanced,breakable,colback=teal!4,colframe=teal!50!black,boxrule=0.5pt,arc=1.1mm,left=3mm,right=3mm,top=1.7mm,bottom=1.7mm,before skip=2mm,after skip=2mm]
\textbf{\color{teal}To'g'ri javob:} \textbf{D) 3 000 belgi}
\par\smallskip
\textbf{Izoh:} Belgilar soni = 12000/4 = 3000.
\end{tcolorbox}
\vspace{2mm}\hrule height 0.2pt \vspace{3mm}

% DB question_id: eeb23af2-0d0d-4b46-aaf8-cb450f88c53a
\needspace{5\baselineskip}
\noindent\textbf{\color{ink}795.}
\par
Piksel nima?
\begin{enumerate}[label=\Alph*),leftmargin=9mm]
\item Rastr tasvirning eng kichik rangli elementi
\item Tasvirning har bir nuqtasiga ajratilgan rang kodlash bitlari miqdori
\item Rastr rasmning eni va bo'yini belgilovchi piksellar umumiy soni
\item Vektor tasvirda egri chiziqni boshqaradigan asosiy tayanch koordinata nuqtasi
\end{enumerate}
\begin{tcolorbox}[enhanced,breakable,colback=teal!4,colframe=teal!50!black,boxrule=0.5pt,arc=1.1mm,left=3mm,right=3mm,top=1.7mm,bottom=1.7mm,before skip=2mm,after skip=2mm]
\textbf{\color{teal}To'g'ri javob:} \textbf{A) Rastr tasvirning eng kichik rangli elementi}
\par\smallskip
\textbf{Izoh:} Piksel rastr tasvirning alohida rang qiymatiga ega eng kichik elementidir. Bir pikselga ajratilgan bitlar soni rang chuqurligini, tasvir o'lchamlari rezolyutsiyani, vektor chiziqlarining tayanch nuqtalari esa boshqa tasvirlash modelini bildiradi.
\end{tcolorbox}
\vspace{2mm}\hrule height 0.2pt \vspace{3mm}

% DB question_id: f01abfb2-cdac-59a2-96c4-cfedbcc1a72a
\needspace{5\baselineskip}
\noindent\textbf{\color{ink}796.}
\par
Namuna chastotasi 2 marta oshsa, boshqa parametrlar o'zgarmasa hajm qanday o'zgaradi?
\begin{enumerate}[label=\Alph*),leftmargin=9mm]
\item 4 marta oshadi
\item O'zgarmaydi
\item 2 marta oshadi
\item 2 marta kamayadi
\end{enumerate}
\begin{tcolorbox}[enhanced,breakable,colback=teal!4,colframe=teal!50!black,boxrule=0.5pt,arc=1.1mm,left=3mm,right=3mm,top=1.7mm,bottom=1.7mm,before skip=2mm,after skip=2mm]
\textbf{\color{teal}To'g'ri javob:} \textbf{C) 2 marta oshadi}
\par\smallskip
\textbf{Izoh:} Hajm sampling ratega to'g'ri proporsional.
\end{tcolorbox}
\vspace{2mm}\hrule height 0.2pt \vspace{3mm}

% DB question_id: f1400ee5-4385-5d34-92f6-42ac6dbf5292
\needspace{5\baselineskip}
\noindent\textbf{\color{ink}797.}
\par
Kamida 17 xil rangni kodlash uchun minimal rang chuqurligi qancha?
\begin{enumerate}[label=\Alph*),leftmargin=9mm]
\item 17 bit
\item 4 bit
\item 6 bit
\item 5 bit
\end{enumerate}
\begin{tcolorbox}[enhanced,breakable,colback=teal!4,colframe=teal!50!black,boxrule=0.5pt,arc=1.1mm,left=3mm,right=3mm,top=1.7mm,bottom=1.7mm,before skip=2mm,after skip=2mm]
\textbf{\color{teal}To'g'ri javob:} \textbf{D) 5 bit}
\par\smallskip
\textbf{Izoh:} 2\textsuperscript{n} \ensuremath{\geq} 17; minimal n = 5.
\end{tcolorbox}
\vspace{2mm}\hrule height 0.2pt \vspace{3mm}

% DB question_id: f34486a6-52e3-5ef0-a6a6-3909cd37f128
\needspace{5\baselineskip}
\noindent\textbf{\color{ink}798.}
\par
800\ensuremath{\times}600, 16 bit, 20 FPS, 5 soniyalik siqilmagan video hajmini MiB da toping.
\begin{enumerate}[label=\Alph*),leftmargin=9mm]
\item 11,444 MiB
\item 0,916 MiB
\item 1831,055 MiB
\item 91,553 MiB
\end{enumerate}
\begin{tcolorbox}[enhanced,breakable,colback=teal!4,colframe=teal!50!black,boxrule=0.5pt,arc=1.1mm,left=3mm,right=3mm,top=1.7mm,bottom=1.7mm,before skip=2mm,after skip=2mm]
\textbf{\color{teal}To'g'ri javob:} \textbf{D) 91,553 MiB}
\par\smallskip
\textbf{Izoh:} Hajm = W\ensuremath{\times}H\ensuremath{\times}d/8\ensuremath{\times}FPS\ensuremath{\times}t = 91,553 MiB.
\end{tcolorbox}
\vspace{2mm}\hrule height 0.2pt \vspace{3mm}

% DB question_id: f4714ec5-2bc6-5fcf-8d2a-dabe7e0d0489
\needspace{5\baselineskip}
\noindent\textbf{\color{ink}799.}
\par
2 bit rang chuqurligi nazariy jihatdan nechta rangni ifodalaydi?
\begin{enumerate}[label=\Alph*),leftmargin=9mm]
\item 4 rang
\item 2 rang
\item 6 rang
\item 7 rang
\end{enumerate}
\begin{tcolorbox}[enhanced,breakable,colback=teal!4,colframe=teal!50!black,boxrule=0.5pt,arc=1.1mm,left=3mm,right=3mm,top=1.7mm,bottom=1.7mm,before skip=2mm,after skip=2mm]
\textbf{\color{teal}To'g'ri javob:} \textbf{A) 4 rang}
\par\smallskip
\textbf{Izoh:} Ranglar soni = 2\textasciicircum{}2 = 4.
\end{tcolorbox}
\vspace{2mm}\hrule height 0.2pt \vspace{3mm}

% DB question_id: fb94fff6-ed5f-542e-a4c9-0b4609bcce28
\needspace{5\baselineskip}
\noindent\textbf{\color{ink}800.}
\par
1920\ensuremath{\times}1080 pikselli, 24 bit rang chuqurligidagi siqilmagan rastr tasvir hajmini MiB da toping.
\begin{enumerate}[label=\Alph*),leftmargin=9mm]
\item 2 073 600 MiB
\item 5,933 MiB
\item 49 766 400 MiB
\item 8,933 MiB
\end{enumerate}
\begin{tcolorbox}[enhanced,breakable,colback=teal!4,colframe=teal!50!black,boxrule=0.5pt,arc=1.1mm,left=3mm,right=3mm,top=1.7mm,bottom=1.7mm,before skip=2mm,after skip=2mm]
\textbf{\color{teal}To'g'ri javob:} \textbf{B) 5,933 MiB}
\par\smallskip
\textbf{Izoh:} Hajm = 1920\ensuremath{\times}1080\ensuremath{\times}24/8 = 6 220 800 B = 5,933 MiB.
\end{tcolorbox}
\vspace{2mm}\hrule height 0.2pt \vspace{3mm}

% DB question_id: 07f2b8b6-2bda-4992-abdb-82c98b7f1017
\needspace{5\baselineskip}
\noindent\textbf{\color{ink}801.}
\par
Audio AAC algoritmi bilan kodlandi. Vaziyatni tahlil qilib, unga eng aniq mos tushunchani toping.
\begin{enumerate}[label=\Alph*),leftmargin=9mm]
\item Transkodlash
\item Zaxira nusxasi
\item Kodek
\item Fayl formati
\end{enumerate}
\begin{tcolorbox}[enhanced,breakable,colback=teal!4,colframe=teal!50!black,boxrule=0.5pt,arc=1.1mm,left=3mm,right=3mm,top=1.7mm,bottom=1.7mm,before skip=2mm,after skip=2mm]
\textbf{\color{teal}To'g'ri javob:} \textbf{C) Kodek}
\par\smallskip
\textbf{Izoh:} Vaziyat Kodekga misol: media oqimini kodlash va dekodlash qoidalari yoki dasturiy moduli.
\end{tcolorbox}
\vspace{2mm}\hrule height 0.2pt \vspace{3mm}

% DB question_id: 5610e704-3eea-47fa-99a7-0c1e6c7ccb18
\needspace{5\baselineskip}
\noindent\textbf{\color{ink}802.}
\par
Asl ASCII nechta kod holatiga ega?
\begin{enumerate}[label=\Alph*),leftmargin=9mm]
\item 64
\item 128
\item 256
\item 65536
\end{enumerate}
\begin{tcolorbox}[enhanced,breakable,colback=teal!4,colframe=teal!50!black,boxrule=0.5pt,arc=1.1mm,left=3mm,right=3mm,top=1.7mm,bottom=1.7mm,before skip=2mm,after skip=2mm]
\textbf{\color{teal}To'g'ri javob:} \textbf{B) 128}
\par\smallskip
\textbf{Izoh:} ASCII 7 bitli: 2\textsuperscript{7} = 128.
\end{tcolorbox}
\vspace{2mm}\hrule height 0.2pt \vspace{3mm}

% DB question_id: 5b497a08-5f5e-4328-ab64-abd776da2d15
\needspace{5\baselineskip}
\noindent\textbf{\color{ink}803.}
\par
Jonli dars videosi yuklanish bilan bir vaqtda ijro etildi. Yaqin tushunchalarni farqlab, aynan qaysi atama vaziyatni eng aniq ifodalashini toping.
\begin{enumerate}[label=\Alph*),leftmargin=9mm]
\item Multimedia
\item Media oqimi
\item Kadr
\item FPS
\end{enumerate}
\begin{tcolorbox}[enhanced,breakable,colback=teal!4,colframe=teal!50!black,boxrule=0.5pt,arc=1.1mm,left=3mm,right=3mm,top=1.7mm,bottom=1.7mm,before skip=2mm,after skip=2mm]
\textbf{\color{teal}To'g'ri javob:} \textbf{B) Media oqimi}
\par\smallskip
\textbf{Izoh:} Vaziyat Media oqimiga misol: vaqt bo'yicha ketma-ket uzatiladigan audio, video yoki boshqa media ma'lumoti.
\end{tcolorbox}
\vspace{2mm}\hrule height 0.2pt \vspace{3mm}

% DB question_id: 5dca3231-f5dd-4ba5-940c-55c255416f1a
\needspace{5\baselineskip}
\noindent\textbf{\color{ink}804.}
\par
Asosiy tekislikdagi belgi 2 bayt, ayrim belgilar surrogat juftida saqlandi. Yaqin tushunchalarni farqlab, aynan qaysi atama vaziyatni eng aniq ifodalashini toping.
\begin{enumerate}[label=\Alph*),leftmargin=9mm]
\item Matn metama'lumoti
\item Belgini kodlash
\item UTF-16
\item ASCII
\end{enumerate}
\begin{tcolorbox}[enhanced,breakable,colback=teal!4,colframe=teal!50!black,boxrule=0.5pt,arc=1.1mm,left=3mm,right=3mm,top=1.7mm,bottom=1.7mm,before skip=2mm,after skip=2mm]
\textbf{\color{teal}To'g'ri javob:} \textbf{C) UTF-16}
\par\smallskip
\textbf{Izoh:} Vaziyat UTF-16ga misol: Unicode kod nuqtalarini 2 yoki 4 baytli birliklarda kodlaydigan shakl.
\end{tcolorbox}
\vspace{2mm}\hrule height 0.2pt \vspace{3mm}

% DB question_id: 63e4435f-6a6e-4511-b3fa-47ebf81f374c
\needspace{5\baselineskip}
\noindent\textbf{\color{ink}805.}
\par
AAAAAA ketma-ketligi A va 6 soni bilan yozildi. Vaziyatni tahlil qilib, unga eng aniq mos tushunchani toping.
\begin{enumerate}[label=\Alph*),leftmargin=9mm]
\item RLE
\item Siqish koeffitsiyenti
\item Kodek
\item Transkodlash
\end{enumerate}
\begin{tcolorbox}[enhanced,breakable,colback=teal!4,colframe=teal!50!black,boxrule=0.5pt,arc=1.1mm,left=3mm,right=3mm,top=1.7mm,bottom=1.7mm,before skip=2mm,after skip=2mm]
\textbf{\color{teal}To'g'ri javob:} \textbf{A) RLE}
\par\smallskip
\textbf{Izoh:} Vaziyat RLEga misol: ketma-ket takrorlangan qiymatni qiymat va takrorlar soni bilan ifodalash usuli.
\end{tcolorbox}
\vspace{2mm}\hrule height 0.2pt \vspace{3mm}

% DB question_id: 7b3571f5-1a3d-4567-bd21-bf691a1d9162
\needspace{5\baselineskip}
\noindent\textbf{\color{ink}806.}
\par
1000 ta ko'rinadigan belgi bor, ammo fayl 1000 baytdan katta. Eng asosli sabab qaysi?
\begin{enumerate}[label=\Alph*),leftmargin=9mm]
\item UTF-8da ayrim belgilar ko'p bayt va metama'lumot bor
\item Har belgi doim 1 bit
\item Unicode kodlamaydi
\item Probel hajm egallamaydi
\end{enumerate}
\begin{tcolorbox}[enhanced,breakable,colback=teal!4,colframe=teal!50!black,boxrule=0.5pt,arc=1.1mm,left=3mm,right=3mm,top=1.7mm,bottom=1.7mm,before skip=2mm,after skip=2mm]
\textbf{\color{teal}To'g'ri javob:} \textbf{A) UTF-8da ayrim belgilar ko'p bayt va metama'lumot bor}
\par\smallskip
\textbf{Izoh:} UTF-8 o'zgaruvchan uzunlikda; xizmat ma'lumotlari ham qo'shilishi mumkin.
\end{tcolorbox}
\vspace{2mm}\hrule height 0.2pt \vspace{3mm}

% DB question_id: 7de5606a-7008-47e7-adc5-c58e4af3e705
\needspace{5\baselineskip}
\noindent\textbf{\color{ink}807.}
\par
MP4 fayli ochilmayapti, lekin kengaytma to'g'ri. Eng asosli sabab qaysi?
\begin{enumerate}[label=\Alph*),leftmargin=9mm]
\item kerakli kodek qo'llab-quvvatlanmaydi
\item piksel mavjud emas
\item konteyner audio saqlamaydi
\item FPS faqat 1 bo'lishi kerak
\end{enumerate}
\begin{tcolorbox}[enhanced,breakable,colback=teal!4,colframe=teal!50!black,boxrule=0.5pt,arc=1.1mm,left=3mm,right=3mm,top=1.7mm,bottom=1.7mm,before skip=2mm,after skip=2mm]
\textbf{\color{teal}To'g'ri javob:} \textbf{A) kerakli kodek qo'llab-quvvatlanmaydi}
\par\smallskip
\textbf{Izoh:} Konteyner tanilsa ham ichki oqim kodeki qo'llanmasligi mumkin.
\end{tcolorbox}
\vspace{2mm}\hrule height 0.2pt \vspace{3mm}

% DB question_id: a1741b87-51fd-48d2-b71a-77147d7acda8
\needspace{5\baselineskip}
\noindent\textbf{\color{ink}808.}
\par
WAV faylda namuna qiymatlari bevosita yozildi. Yaqin tushunchalarni farqlab, aynan qaysi atama vaziyatni eng aniq ifodalashini toping.
\begin{enumerate}[label=\Alph*),leftmargin=9mm]
\item PCM
\item Audio bit tezligi
\item Analog tovush
\item ADC
\end{enumerate}
\begin{tcolorbox}[enhanced,breakable,colback=teal!4,colframe=teal!50!black,boxrule=0.5pt,arc=1.1mm,left=3mm,right=3mm,top=1.7mm,bottom=1.7mm,before skip=2mm,after skip=2mm]
\textbf{\color{teal}To'g'ri javob:} \textbf{A) PCM}
\par\smallskip
\textbf{Izoh:} Vaziyat PCMga misol: namunalar qiymatini odatda siqilmagan holda ketma-ket saqlash usuli.
\end{tcolorbox}
\vspace{2mm}\hrule height 0.2pt \vspace{3mm}

% DB question_id: 2ab2552e-7ea9-451a-92fc-e91e1513fb63
\needspace{5\baselineskip}
\noindent\textbf{\color{ink}809.}
\par
H.264 video AV1 formatidagi oqimga qayta kodlandi. Yaqin tushunchalarni farqlab, aynan qaysi atama vaziyatni eng aniq ifodalashini toping.
\begin{enumerate}[label=\Alph*),leftmargin=9mm]
\item Zaxira nusxasi
\item Fayl formati
\item Fayl kengaytmasi
\item Transkodlash
\end{enumerate}
\begin{tcolorbox}[enhanced,breakable,colback=teal!4,colframe=teal!50!black,boxrule=0.5pt,arc=1.1mm,left=3mm,right=3mm,top=1.7mm,bottom=1.7mm,before skip=2mm,after skip=2mm]
\textbf{\color{teal}To'g'ri javob:} \textbf{D) Transkodlash}
\par\smallskip
\textbf{Izoh:} Vaziyat Transkodlashga misol: media oqimini bir kodek yoki parametrdan boshqasiga qayta kodlash.
\end{tcolorbox}
\vspace{2mm}\hrule height 0.2pt \vspace{3mm}

% DB question_id: 53ba68e2-4346-418c-9fd2-3ba1a9450229
\needspace{5\baselineskip}
\noindent\textbf{\color{ink}810.}
\par
Taqdimot matn, animatsiya, ovoz va videodan foydalandi. Yaqin tushunchalarni farqlab, aynan qaysi atama vaziyatni eng aniq ifodalashini toping.
\begin{enumerate}[label=\Alph*),leftmargin=9mm]
\item Kadr
\item FPS
\item Multimedia
\item Video aniqligi
\end{enumerate}
\begin{tcolorbox}[enhanced,breakable,colback=teal!4,colframe=teal!50!black,boxrule=0.5pt,arc=1.1mm,left=3mm,right=3mm,top=1.7mm,bottom=1.7mm,before skip=2mm,after skip=2mm]
\textbf{\color{teal}To'g'ri javob:} \textbf{C) Multimedia}
\par\smallskip
\textbf{Izoh:} Vaziyat Multimediaga misol: bir nechta media turini yagona mazmun yoki interaktiv tizimda birlashtirish.
\end{tcolorbox}
\vspace{2mm}\hrule height 0.2pt \vspace{3mm}

% DB question_id: b6f5a385-c646-42d4-ac88-a479067762af
\needspace{5\baselineskip}
\noindent\textbf{\color{ink}811.}
\par
Qaysi jumla aniq?
\begin{enumerate}[label=\Alph*),leftmargin=9mm]
\item UTF-8 Unicode'ning kodlash shaklidir
\item Unicode har bir belgini doim 2 baytda saqlaydi
\item ASCII barcha yozuvlarni qamraydi
\item Kod nuqtasi fayldagi baytlar soniga doim teng
\end{enumerate}
\begin{tcolorbox}[enhanced,breakable,colback=teal!4,colframe=teal!50!black,boxrule=0.5pt,arc=1.1mm,left=3mm,right=3mm,top=1.7mm,bottom=1.7mm,before skip=2mm,after skip=2mm]
\textbf{\color{teal}To'g'ri javob:} \textbf{A) UTF-8 Unicode'ning kodlash shaklidir}
\par\smallskip
\textbf{Izoh:} Unicode kod nuqtalarini belgilaydi, UTF-8 esa ularni baytlarga kodlaydi.
\end{tcolorbox}
\vspace{2mm}\hrule height 0.2pt \vspace{3mm}

% DB question_id: dad82698-c9cb-4c92-8979-e25d621dcb7f
\needspace{5\baselineskip}
\noindent\textbf{\color{ink}812.}
\par
Ko'rinadigan belgilar hajmiga qo'shimcha ravishda sarlavha va kodlash belgisi yozildi. Yaqin tushunchalarni farqlab, aynan qaysi atama vaziyatni eng aniq ifodalashini toping.
\begin{enumerate}[label=\Alph*),leftmargin=9mm]
\item Belgini kodlash
\item ASCII
\item Kengaytirilgan ASCII
\item Matn metama'lumoti
\end{enumerate}
\begin{tcolorbox}[enhanced,breakable,colback=teal!4,colframe=teal!50!black,boxrule=0.5pt,arc=1.1mm,left=3mm,right=3mm,top=1.7mm,bottom=1.7mm,before skip=2mm,after skip=2mm]
\textbf{\color{teal}To'g'ri javob:} \textbf{D) Matn metama'lumoti}
\par\smallskip
\textbf{Izoh:} Vaziyat Matn metama'lumotiga misol: faylning kodlash turi, tuzilishi yoki xizmat ma'lumotlarini tavsiflovchi qo'shimcha ma'lumot.
\end{tcolorbox}
\vspace{2mm}\hrule height 0.2pt \vspace{3mm}
\section{Axborot hajmini hisoblashning yagona strategiyasi}

% DB question_id: 0025fecc-506a-42d0-af20-ff01ae5a0c62
\needspace{5\baselineskip}
\noindent\textbf{\color{ink}813.}
\par
"Xartli formulasi" atamasiga eng aniq ta'rifni toping.
\begin{enumerate}[label=\Alph*),leftmargin=9mm]
\item teng ehtimolli N holatdan birini kodlash uchun i = log\textsubscript{2}N munosabatini beradigan formula
\item kodlashda qatnashadigan turli belgilar yoki holatlar soni
\item N holat uchun log\textsubscript{2}N yuqoriga yaxlitlangan bitlar soni
\item SI standartida 1000 baytga teng o'nlik birlik
\end{enumerate}
\begin{tcolorbox}[enhanced,breakable,colback=teal!4,colframe=teal!50!black,boxrule=0.5pt,arc=1.1mm,left=3mm,right=3mm,top=1.7mm,bottom=1.7mm,before skip=2mm,after skip=2mm]
\textbf{\color{teal}To'g'ri javob:} \textbf{A) teng ehtimolli N holatdan birini kodlash uchun i = log\textsubscript{2}N munosabatini beradigan formula}
\par\smallskip
\textbf{Izoh:} Xartli formulasi - teng ehtimolli N holatdan birini kodlash uchun i = log\textsubscript{2}N munosabatini beradigan formula.
\end{tcolorbox}
\vspace{2mm}\hrule height 0.2pt \vspace{3mm}

% DB question_id: 0f0af8d7-8e85-4983-b4dc-1a3186137d15
\needspace{5\baselineskip}
\noindent\textbf{\color{ink}814.}
\par
"Formula tanlash" atamasiga eng aniq ta'rifni toping.
\begin{enumerate}[label=\Alph*),leftmargin=9mm]
\item berilgan axborot turiga mos hajm formulasini aniqlash bosqichi
\item hisoblashdan oldin bit-bayt, soniya-daqiqa va prefikslarni bir xil tizimga keltirish
\item siqish va xizmat ma'lumotlarisiz nazariy axborot hajmi
\item siqish algoritmi qo'llangandan keyingi fayl yoki oqim hajmi
\end{enumerate}
\begin{tcolorbox}[enhanced,breakable,colback=teal!4,colframe=teal!50!black,boxrule=0.5pt,arc=1.1mm,left=3mm,right=3mm,top=1.7mm,bottom=1.7mm,before skip=2mm,after skip=2mm]
\textbf{\color{teal}To'g'ri javob:} \textbf{A) berilgan axborot turiga mos hajm formulasini aniqlash bosqichi}
\par\smallskip
\textbf{Izoh:} Formula tanlash - berilgan axborot turiga mos hajm formulasini aniqlash bosqichi.
\end{tcolorbox}
\vspace{2mm}\hrule height 0.2pt \vspace{3mm}

% DB question_id: 109c3a9c-e461-4278-80d6-62e9809dca9d
\needspace{5\baselineskip}
\noindent\textbf{\color{ink}815.}
\par
Quyidagi ta'rif qaysi atamaga tegishli? Berilgan axborot turiga mos hajm formulasini aniqlash bosqichi.
\begin{enumerate}[label=\Alph*),leftmargin=9mm]
\item Formula tanlash
\item Birliklarni normallashtirish
\item Xom hajm
\item Siqilgan hajm
\end{enumerate}
\begin{tcolorbox}[enhanced,breakable,colback=teal!4,colframe=teal!50!black,boxrule=0.5pt,arc=1.1mm,left=3mm,right=3mm,top=1.7mm,bottom=1.7mm,before skip=2mm,after skip=2mm]
\textbf{\color{teal}To'g'ri javob:} \textbf{A) Formula tanlash}
\par\smallskip
\textbf{Izoh:} Ta'rif Formula tanlash tushunchasiga tegishli.
\end{tcolorbox}
\vspace{2mm}\hrule height 0.2pt \vspace{3mm}

% DB question_id: 11fc6fa4-72a3-47d0-b192-b06d648c6efb
\needspace{5\baselineskip}
\noindent\textbf{\color{ink}816.}
\par
"Siqish koeffitsiyenti" atamasiga eng aniq ta'rifni toping.
\begin{enumerate}[label=\Alph*),leftmargin=9mm]
\item media oqimini kodlash va dekodlash qoidalari yoki dasturiy moduli
\item asl hajmning siqilgan hajmga yoki kamayish foiziga nisbati
\item media oqimini bir kodek yoki parametrdan boshqasiga qayta kodlash
\item ma'lumot yo'qolsa tiklash uchun alohida joyda saqlangan nusxa
\end{enumerate}
\begin{tcolorbox}[enhanced,breakable,colback=teal!4,colframe=teal!50!black,boxrule=0.5pt,arc=1.1mm,left=3mm,right=3mm,top=1.7mm,bottom=1.7mm,before skip=2mm,after skip=2mm]
\textbf{\color{teal}To'g'ri javob:} \textbf{B) asl hajmning siqilgan hajmga yoki kamayish foiziga nisbati}
\par\smallskip
\textbf{Izoh:} Siqish koeffitsiyenti - asl hajmning siqilgan hajmga yoki kamayish foiziga nisbati.
\end{tcolorbox}
\vspace{2mm}\hrule height 0.2pt \vspace{3mm}

% DB question_id: 3270a5de-4b81-4870-89af-89bce140a577
\needspace{5\baselineskip}
\noindent\textbf{\color{ink}817.}
\par
Siqilmagan bitta rastr kadr 16 000 bayt. Uning o'lchami 320\ensuremath{\times}400 piksel. Har piksel uchun bir xil bit chuqurligi ishlatilsa, ko'pi bilan nechta rang kodlash mumkin?
\begin{enumerate}[label=\Alph*),leftmargin=9mm]
\item 1
\item 2
\item 4
\item 8
\end{enumerate}
\begin{tcolorbox}[enhanced,breakable,colback=teal!4,colframe=teal!50!black,boxrule=0.5pt,arc=1.1mm,left=3mm,right=3mm,top=1.7mm,bottom=1.7mm,before skip=2mm,after skip=2mm]
\textbf{\color{teal}To'g'ri javob:} \textbf{B) 2}
\par\smallskip
\textbf{Izoh:} 16 000 bayt=128 000 bit va 320\ensuremath{\cdot}400=128 000 piksel. Demak 1 bit/piksel, ranglar soni 2\textsuperscript{1}=2.
\end{tcolorbox}
\vspace{2mm}\hrule height 0.2pt \vspace{3mm}

% DB question_id: 3f666bf5-073d-4971-a9ed-4b754f72c714
\needspace{5\baselineskip}
\noindent\textbf{\color{ink}818.}
\par
"Siqilgan hajm" atamasiga eng aniq ta'rifni toping.
\begin{enumerate}[label=\Alph*),leftmargin=9mm]
\item har bir mustaqil kanal hajmini umumiy formulaga kiritish
\item FPS va davomiylik ko'paytmasi orqali topiladigan video tasvirlari soni
\item hajm va boshqa parametrlar ma'lum bo'lganda noma'lum parametrni formuladan ajratib topish
\item siqish algoritmi qo'llangandan keyingi fayl yoki oqim hajmi
\end{enumerate}
\begin{tcolorbox}[enhanced,breakable,colback=teal!4,colframe=teal!50!black,boxrule=0.5pt,arc=1.1mm,left=3mm,right=3mm,top=1.7mm,bottom=1.7mm,before skip=2mm,after skip=2mm]
\textbf{\color{teal}To'g'ri javob:} \textbf{D) siqish algoritmi qo'llangandan keyingi fayl yoki oqim hajmi}
\par\smallskip
\textbf{Izoh:} Siqilgan hajm - siqish algoritmi qo'llangandan keyingi fayl yoki oqim hajmi.
\end{tcolorbox}
\vspace{2mm}\hrule height 0.2pt \vspace{3mm}

% DB question_id: 4fb66bd8-44dc-4e70-a5f2-9100df175468
\needspace{5\baselineskip}
\noindent\textbf{\color{ink}819.}
\par
"Xom hajm" atamasiga eng aniq ta'rifni toping.
\begin{enumerate}[label=\Alph*),leftmargin=9mm]
\item siqish algoritmi qo'llangandan keyingi fayl yoki oqim hajmi
\item har bir mustaqil kanal hajmini umumiy formulaga kiritish
\item siqish va xizmat ma'lumotlarisiz nazariy axborot hajmi
\item FPS va davomiylik ko'paytmasi orqali topiladigan video tasvirlari soni
\end{enumerate}
\begin{tcolorbox}[enhanced,breakable,colback=teal!4,colframe=teal!50!black,boxrule=0.5pt,arc=1.1mm,left=3mm,right=3mm,top=1.7mm,bottom=1.7mm,before skip=2mm,after skip=2mm]
\textbf{\color{teal}To'g'ri javob:} \textbf{C) siqish va xizmat ma'lumotlarisiz nazariy axborot hajmi}
\par\smallskip
\textbf{Izoh:} Xom hajm - siqish va xizmat ma'lumotlarisiz nazariy axborot hajmi.
\end{tcolorbox}
\vspace{2mm}\hrule height 0.2pt \vspace{3mm}

% DB question_id: 568de923-1191-4470-b40e-73a66e133b28
\needspace{5\baselineskip}
\noindent\textbf{\color{ink}820.}
\par
"Birliklarni normallashtirish" atamasiga eng aniq ta'rifni toping.
\begin{enumerate}[label=\Alph*),leftmargin=9mm]
\item siqish va xizmat ma'lumotlarisiz nazariy axborot hajmi
\item hisoblashdan oldin bit-bayt, soniya-daqiqa va prefikslarni bir xil tizimga keltirish
\item siqish algoritmi qo'llangandan keyingi fayl yoki oqim hajmi
\item har bir mustaqil kanal hajmini umumiy formulaga kiritish
\end{enumerate}
\begin{tcolorbox}[enhanced,breakable,colback=teal!4,colframe=teal!50!black,boxrule=0.5pt,arc=1.1mm,left=3mm,right=3mm,top=1.7mm,bottom=1.7mm,before skip=2mm,after skip=2mm]
\textbf{\color{teal}To'g'ri javob:} \textbf{B) hisoblashdan oldin bit-bayt, soniya-daqiqa va prefikslarni bir xil tizimga keltirish}
\par\smallskip
\textbf{Izoh:} Birliklarni normallashtirish - hisoblashdan oldin bit-bayt, soniya-daqiqa va prefikslarni bir xil tizimga keltirish.
\end{tcolorbox}
\vspace{2mm}\hrule height 0.2pt \vspace{3mm}

% DB question_id: 604f66c0-c5b0-5cf7-ac79-35793a766636
\needspace{5\baselineskip}
\noindent\textbf{\color{ink}821.}
\par
Axborot hajmi 1 048 576 bit. Bu hajm necha KiB? \texttt{1 KiB = 1024 bayt}.
\begin{enumerate}[label=\Alph*),leftmargin=9mm]
\item 128 KiB
\item 1024 KiB
\item 131 072 KiB
\item 8192 KiB
\end{enumerate}
\begin{tcolorbox}[enhanced,breakable,colback=teal!4,colframe=teal!50!black,boxrule=0.5pt,arc=1.1mm,left=3mm,right=3mm,top=1.7mm,bottom=1.7mm,before skip=2mm,after skip=2mm]
\textbf{\color{teal}To'g'ri javob:} \textbf{A) 128 KiB}
\par\smallskip
\textbf{Izoh:} 1 048 576 bit \ensuremath{\div} 8 = 131 072 bayt; 131 072 \ensuremath{\div} 1024 = 128 KiB.
\end{tcolorbox}
\vspace{2mm}\hrule height 0.2pt \vspace{3mm}

% DB question_id: 9641ac72-1347-4ffd-9c2f-ee41c1efe2c5
\needspace{5\baselineskip}
\noindent\textbf{\color{ink}822.}
\par
Quyidagi ta'rif qaysi atamaga tegishli? Hisoblashdan oldin bit-bayt, soniya-daqiqa va prefikslarni bir xil tizimga keltirish.
\begin{enumerate}[label=\Alph*),leftmargin=9mm]
\item Xom hajm
\item Birliklarni normallashtirish
\item Siqilgan hajm
\item Ko'p kanalli hisob
\end{enumerate}
\begin{tcolorbox}[enhanced,breakable,colback=teal!4,colframe=teal!50!black,boxrule=0.5pt,arc=1.1mm,left=3mm,right=3mm,top=1.7mm,bottom=1.7mm,before skip=2mm,after skip=2mm]
\textbf{\color{teal}To'g'ri javob:} \textbf{B) Birliklarni normallashtirish}
\par\smallskip
\textbf{Izoh:} Ta'rif Birliklarni normallashtirish tushunchasiga tegishli.
\end{tcolorbox}
\vspace{2mm}\hrule height 0.2pt \vspace{3mm}

% DB question_id: 9d3c42d2-e1eb-47bc-8040-357cfa195f98
\needspace{5\baselineskip}
\noindent\textbf{\color{ink}823.}
\par
"Ko'p kanalli hisob" atamasiga eng aniq ta'rifni toping.
\begin{enumerate}[label=\Alph*),leftmargin=9mm]
\item har bir mustaqil kanal hajmini umumiy formulaga kiritish
\item FPS va davomiylik ko'paytmasi orqali topiladigan video tasvirlari soni
\item hajm va boshqa parametrlar ma'lum bo'lganda noma'lum parametrni formuladan ajratib topish
\item natijaning birlik, tartib va shartga mosligini baholash bosqichi
\end{enumerate}
\begin{tcolorbox}[enhanced,breakable,colback=teal!4,colframe=teal!50!black,boxrule=0.5pt,arc=1.1mm,left=3mm,right=3mm,top=1.7mm,bottom=1.7mm,before skip=2mm,after skip=2mm]
\textbf{\color{teal}To'g'ri javob:} \textbf{A) har bir mustaqil kanal hajmini umumiy formulaga kiritish}
\par\smallskip
\textbf{Izoh:} Ko'p kanalli hisob - har bir mustaqil kanal hajmini umumiy formulaga kiritish.
\end{tcolorbox}
\vspace{2mm}\hrule height 0.2pt \vspace{3mm}

% DB question_id: a04640e5-e99c-45ae-a72b-57c8a833f0d3
\needspace{5\baselineskip}
\noindent\textbf{\color{ink}824.}
\par
"Kadrlar soni" atamasiga eng aniq ta'rifni toping.
\begin{enumerate}[label=\Alph*),leftmargin=9mm]
\item hajm va boshqa parametrlar ma'lum bo'lganda noma'lum parametrni formuladan ajratib topish
\item FPS va davomiylik ko'paytmasi orqali topiladigan video tasvirlari soni
\item natijaning birlik, tartib va shartga mosligini baholash bosqichi
\item berilgan axborot turiga mos hajm formulasini aniqlash bosqichi
\end{enumerate}
\begin{tcolorbox}[enhanced,breakable,colback=teal!4,colframe=teal!50!black,boxrule=0.5pt,arc=1.1mm,left=3mm,right=3mm,top=1.7mm,bottom=1.7mm,before skip=2mm,after skip=2mm]
\textbf{\color{teal}To'g'ri javob:} \textbf{B) FPS va davomiylik ko'paytmasi orqali topiladigan video tasvirlari soni}
\par\smallskip
\textbf{Izoh:} Kadrlar soni - FPS va davomiylik ko'paytmasi orqali topiladigan video tasvirlari soni.
\end{tcolorbox}
\vspace{2mm}\hrule height 0.2pt \vspace{3mm}

% DB question_id: a56d5012-3758-42e8-8f1d-4801b2aa3439
\needspace{5\baselineskip}
\noindent\textbf{\color{ink}825.}
\par
"Alifbo quvvati" atamasiga eng aniq ta'rifni toping.
\begin{enumerate}[label=\Alph*),leftmargin=9mm]
\item N holat uchun log\textsubscript{2}N yuqoriga yaxlitlangan bitlar soni
\item kodlashda qatnashadigan turli belgilar yoki holatlar soni
\item SI standartida 1000 baytga teng o'nlik birlik
\item IEC standartida 1024 baytga teng ikkilik birlik
\end{enumerate}
\begin{tcolorbox}[enhanced,breakable,colback=teal!4,colframe=teal!50!black,boxrule=0.5pt,arc=1.1mm,left=3mm,right=3mm,top=1.7mm,bottom=1.7mm,before skip=2mm,after skip=2mm]
\textbf{\color{teal}To'g'ri javob:} \textbf{B) kodlashda qatnashadigan turli belgilar yoki holatlar soni}
\par\smallskip
\textbf{Izoh:} Alifbo quvvati - kodlashda qatnashadigan turli belgilar yoki holatlar soni.
\end{tcolorbox}
\vspace{2mm}\hrule height 0.2pt \vspace{3mm}

% DB question_id: c173124c-8e83-4e5c-8774-a71c653d216b
\needspace{5\baselineskip}
\noindent\textbf{\color{ink}826.}
\par
"Teskari masala" atamasiga eng aniq ta'rifni toping.
\begin{enumerate}[label=\Alph*),leftmargin=9mm]
\item natijaning birlik, tartib va shartga mosligini baholash bosqichi
\item berilgan axborot turiga mos hajm formulasini aniqlash bosqichi
\item hajm va boshqa parametrlar ma'lum bo'lganda noma'lum parametrni formuladan ajratib topish
\item hisoblashdan oldin bit-bayt, soniya-daqiqa va prefikslarni bir xil tizimga keltirish
\end{enumerate}
\begin{tcolorbox}[enhanced,breakable,colback=teal!4,colframe=teal!50!black,boxrule=0.5pt,arc=1.1mm,left=3mm,right=3mm,top=1.7mm,bottom=1.7mm,before skip=2mm,after skip=2mm]
\textbf{\color{teal}To'g'ri javob:} \textbf{C) hajm va boshqa parametrlar ma'lum bo'lganda noma'lum parametrni formuladan ajratib topish}
\par\smallskip
\textbf{Izoh:} Teskari masala - hajm va boshqa parametrlar ma'lum bo'lganda noma'lum parametrni formuladan ajratib topish.
\end{tcolorbox}
\vspace{2mm}\hrule height 0.2pt \vspace{3mm}

% DB question_id: cb1a62b0-6b5f-5931-a6f5-98be4501aa8d
\needspace{5\baselineskip}
\noindent\textbf{\color{ink}827.}
\par
\texttt{Informatika} so'zi 11 ta belgidan iborat. Har bir belgi 1 baytda kodlansa, so'z 11 264 marta yozilgan fayl necha KiB bo'ladi? \texttt{1 KiB = 1024 bayt}.
\begin{enumerate}[label=\Alph*),leftmargin=9mm]
\item 110 KiB
\item 121 KiB
\item 11 264 KiB
\item 11 KiB
\end{enumerate}
\begin{tcolorbox}[enhanced,breakable,colback=teal!4,colframe=teal!50!black,boxrule=0.5pt,arc=1.1mm,left=3mm,right=3mm,top=1.7mm,bottom=1.7mm,before skip=2mm,after skip=2mm]
\textbf{\color{teal}To'g'ri javob:} \textbf{B) 121 KiB}
\par\smallskip
\textbf{Izoh:} 11\ensuremath{\cdot}11 264=123 904 bayt; 123 904\ensuremath{\div}1024=121 KiB.
\end{tcolorbox}
\vspace{2mm}\hrule height 0.2pt \vspace{3mm}

% DB question_id: d8b53a68-8d65-4432-a247-559a38469b4f
\needspace{5\baselineskip}
\noindent\textbf{\color{ink}828.}
\par
"Minimal butun kod uzunligi" atamasiga eng aniq ta'rifni toping.
\begin{enumerate}[label=\Alph*),leftmargin=9mm]
\item SI standartida 1000 baytga teng o'nlik birlik
\item IEC standartida 1024 baytga teng ikkilik birlik
\item N holat uchun log\textsubscript{2}N yuqoriga yaxlitlangan bitlar soni
\item tashuvchida saqlanishi mumkin bo'lgan axborotning umumiy hajmi
\end{enumerate}
\begin{tcolorbox}[enhanced,breakable,colback=teal!4,colframe=teal!50!black,boxrule=0.5pt,arc=1.1mm,left=3mm,right=3mm,top=1.7mm,bottom=1.7mm,before skip=2mm,after skip=2mm]
\textbf{\color{teal}To'g'ri javob:} \textbf{C) N holat uchun log\textsubscript{2}N yuqoriga yaxlitlangan bitlar soni}
\par\smallskip
\textbf{Izoh:} Minimal butun kod uzunligi - N holat uchun log\textsubscript{2}N yuqoriga yaxlitlangan bitlar soni.
\end{tcolorbox}
\vspace{2mm}\hrule height 0.2pt \vspace{3mm}

% DB question_id: e80f9c6c-a915-4997-8f56-67b3a36a9d83
\needspace{5\baselineskip}
\noindent\textbf{\color{ink}829.}
\par
"Mantiqiy tekshiruv" atamasiga eng aniq ta'rifni toping.
\begin{enumerate}[label=\Alph*),leftmargin=9mm]
\item berilgan axborot turiga mos hajm formulasini aniqlash bosqichi
\item hisoblashdan oldin bit-bayt, soniya-daqiqa va prefikslarni bir xil tizimga keltirish
\item siqish va xizmat ma'lumotlarisiz nazariy axborot hajmi
\item natijaning birlik, tartib va shartga mosligini baholash bosqichi
\end{enumerate}
\begin{tcolorbox}[enhanced,breakable,colback=teal!4,colframe=teal!50!black,boxrule=0.5pt,arc=1.1mm,left=3mm,right=3mm,top=1.7mm,bottom=1.7mm,before skip=2mm,after skip=2mm]
\textbf{\color{teal}To'g'ri javob:} \textbf{D) natijaning birlik, tartib va shartga mosligini baholash bosqichi}
\par\smallskip
\textbf{Izoh:} Mantiqiy tekshiruv - natijaning birlik, tartib va shartga mosligini baholash bosqichi.
\end{tcolorbox}
\vspace{2mm}\hrule height 0.2pt \vspace{3mm}

% DB question_id: 0b6a649d-1579-5550-bf8b-3c425af9dc46
\needspace{5\baselineskip}
\noindent\textbf{\color{ink}830.}
\par
Mono audio 8 bitli kvantlash va 25 000 Hz chastotada 4 minut yozildi. Siqilmagan hajm necha KiB?
\begin{enumerate}[label=\Alph*),leftmargin=9mm]
\item 5859,375 KiB
\item 6000 KiB
\item 46 875 KiB
\item 750 KiB
\end{enumerate}
\begin{tcolorbox}[enhanced,breakable,colback=teal!4,colframe=teal!50!black,boxrule=0.5pt,arc=1.1mm,left=3mm,right=3mm,top=1.7mm,bottom=1.7mm,before skip=2mm,after skip=2mm]
\textbf{\color{teal}To'g'ri javob:} \textbf{A) 5859,375 KiB}
\par\smallskip
\textbf{Izoh:} 8\ensuremath{\cdot}25 000\ensuremath{\cdot}240=48 000 000 bit=6 000 000 bayt=5859,375 KiB.
\end{tcolorbox}
\vspace{2mm}\hrule height 0.2pt \vspace{3mm}

% DB question_id: 1ad3d01d-7e8b-5dd3-bd57-f399551ffd7d
\needspace{5\baselineskip}
\noindent\textbf{\color{ink}831.}
\par
1024\ensuremath{\times}768 pikselli siqilmagan tasvirning har pikseli 32 768 rangdan birini qabul qiladi. Kerakli eng kam videoxotira hajmini KiB da toping.
\begin{enumerate}[label=\Alph*),leftmargin=9mm]
\item 1440 KiB
\item 960 KiB
\item 1536 KiB
\item 11 796 480 KiB
\end{enumerate}
\begin{tcolorbox}[enhanced,breakable,colback=teal!4,colframe=teal!50!black,boxrule=0.5pt,arc=1.1mm,left=3mm,right=3mm,top=1.7mm,bottom=1.7mm,before skip=2mm,after skip=2mm]
\textbf{\color{teal}To'g'ri javob:} \textbf{A) 1440 KiB}
\par\smallskip
\textbf{Izoh:} 32 768=2\textsuperscript{1}\textsuperscript{5}, demak 15 bit/piksel. 1024\ensuremath{\cdot}768\ensuremath{\cdot}15/(8\ensuremath{\cdot}1024)=1440 KiB.
\end{tcolorbox}
\vspace{2mm}\hrule height 0.2pt \vspace{3mm}

% DB question_id: 1c600307-059b-4475-80f9-b8356ec5f044
\needspace{5\baselineskip}
\noindent\textbf{\color{ink}832.}
\par
Quyidagi ta'rif qaysi atamaga tegishli? Kodlashda qatnashadigan turli belgilar yoki holatlar soni.
\begin{enumerate}[label=\Alph*),leftmargin=9mm]
\item Minimal butun kod uzunligi
\item Alifbo quvvati
\item kB
\item KiB
\end{enumerate}
\begin{tcolorbox}[enhanced,breakable,colback=teal!4,colframe=teal!50!black,boxrule=0.5pt,arc=1.1mm,left=3mm,right=3mm,top=1.7mm,bottom=1.7mm,before skip=2mm,after skip=2mm]
\textbf{\color{teal}To'g'ri javob:} \textbf{B) Alifbo quvvati}
\par\smallskip
\textbf{Izoh:} Ta'rif Alifbo quvvati tushunchasiga tegishli.
\end{tcolorbox}
\vspace{2mm}\hrule height 0.2pt \vspace{3mm}

% DB question_id: 343bd62c-efa6-5ec2-9846-90df7352aec4
\needspace{5\baselineskip}
\noindent\textbf{\color{ink}833.}
\par
384\ensuremath{\times}640 pikselli siqilmagan rastr tasvir 90 KiB joy egallaydi. Har piksel bir xil butun bitlarda kodlansa, tasvir ko'pi bilan nechta rangni ifodalay oladi?
\begin{enumerate}[label=\Alph*),leftmargin=9mm]
\item 8 ta
\item 16 ta
\item 3 ta
\item 256 ta
\end{enumerate}
\begin{tcolorbox}[enhanced,breakable,colback=teal!4,colframe=teal!50!black,boxrule=0.5pt,arc=1.1mm,left=3mm,right=3mm,top=1.7mm,bottom=1.7mm,before skip=2mm,after skip=2mm]
\textbf{\color{teal}To'g'ri javob:} \textbf{A) 8 ta}
\par\smallskip
\textbf{Izoh:} Bit chuqurligi 90\ensuremath{\cdot}1024\ensuremath{\cdot}8/(384\ensuremath{\cdot}640)=3 bit/piksel; 2\textsuperscript{3}=8 rang.
\end{tcolorbox}
\vspace{2mm}\hrule height 0.2pt \vspace{3mm}

% DB question_id: 34c1c9e8-d379-4d2d-9e17-38d4c6b99107
\needspace{5\baselineskip}
\noindent\textbf{\color{ink}834.}
\par
Rasm uchun W \ensuremath{\times} H \ensuremath{\times} rang chuqurligi hisoblandi. Bu vaziyatdagi asosiy tushunchani aniqlang.
\begin{enumerate}[label=\Alph*),leftmargin=9mm]
\item Siqilgan hajm
\item Ko'p kanalli hisob
\item Xom hajm
\item Kadrlar soni
\end{enumerate}
\begin{tcolorbox}[enhanced,breakable,colback=teal!4,colframe=teal!50!black,boxrule=0.5pt,arc=1.1mm,left=3mm,right=3mm,top=1.7mm,bottom=1.7mm,before skip=2mm,after skip=2mm]
\textbf{\color{teal}To'g'ri javob:} \textbf{C) Xom hajm}
\par\smallskip
\textbf{Izoh:} Vaziyat Xom hajmga misol: siqish va xizmat ma'lumotlarisiz nazariy axborot hajmi.
\end{tcolorbox}
\vspace{2mm}\hrule height 0.2pt \vspace{3mm}

% DB question_id: 404fa955-3e79-4fcd-a7f4-27a29bddbb14
\needspace{5\baselineskip}
\noindent\textbf{\color{ink}835.}
\par
Quyidagi ta'rif qaysi atamaga tegishli? Har bir mustaqil kanal hajmini umumiy formulaga kiritish.
\begin{enumerate}[label=\Alph*),leftmargin=9mm]
\item Ko'p kanalli hisob
\item Kadrlar soni
\item Teskari masala
\item Mantiqiy tekshiruv
\end{enumerate}
\begin{tcolorbox}[enhanced,breakable,colback=teal!4,colframe=teal!50!black,boxrule=0.5pt,arc=1.1mm,left=3mm,right=3mm,top=1.7mm,bottom=1.7mm,before skip=2mm,after skip=2mm]
\textbf{\color{teal}To'g'ri javob:} \textbf{A) Ko'p kanalli hisob}
\par\smallskip
\textbf{Izoh:} Ta'rif Ko'p kanalli hisob tushunchasiga tegishli.
\end{tcolorbox}
\vspace{2mm}\hrule height 0.2pt \vspace{3mm}

% DB question_id: 428d427e-45ac-51c9-bf8c-bf7ab3ae5dcb
\needspace{5\baselineskip}
\noindent\textbf{\color{ink}836.}
\par
Mono audio 16 bitli kvantlash chuqurligi va 3200 Hz namuna olish chastotasida 3 minut yozildi. Siqilmagan fayl hajmini KiB da toping.
\begin{enumerate}[label=\Alph*),leftmargin=9mm]
\item 1125 KiB
\item 1152 KiB
\item 900 KiB
\item 9216 KiB
\end{enumerate}
\begin{tcolorbox}[enhanced,breakable,colback=teal!4,colframe=teal!50!black,boxrule=0.5pt,arc=1.1mm,left=3mm,right=3mm,top=1.7mm,bottom=1.7mm,before skip=2mm,after skip=2mm]
\textbf{\color{teal}To'g'ri javob:} \textbf{A) 1125 KiB}
\par\smallskip
\textbf{Izoh:} 16\ensuremath{\cdot}3200\ensuremath{\cdot}180=9 216 000 bit=1 152 000 bayt=1125 KiB.
\end{tcolorbox}
\vspace{2mm}\hrule height 0.2pt \vspace{3mm}

% DB question_id: 439688ab-1119-5e11-8a6c-c177d616981a
\needspace{5\baselineskip}
\noindent\textbf{\color{ink}837.}
\par
Meteorologik stansiyada namlikning har bir o'lchovi 0 dan 100 gacha bo'lgan butun son va eng kam bir xil uzunlikdagi ikkilik kod bilan yoziladi. 80 ta o'lchovning axborot hajmini baytda toping.
\begin{enumerate}[label=\Alph*),leftmargin=9mm]
\item 10 bayt
\item 70 bayt
\item 80 bayt
\item 50 bayt
\end{enumerate}
\begin{tcolorbox}[enhanced,breakable,colback=teal!4,colframe=teal!50!black,boxrule=0.5pt,arc=1.1mm,left=3mm,right=3mm,top=1.7mm,bottom=1.7mm,before skip=2mm,after skip=2mm]
\textbf{\color{teal}To'g'ri javob:} \textbf{B) 70 bayt}
\par\smallskip
\textbf{Izoh:} 0-100 oralig'ida 101 ta qiymat bor; \texttt{2\textsuperscript{6}\textless{}101\ensuremath{\leq}2\textsuperscript{7}}, demak har o'lchov 7 bit. 80\ensuremath{\times}7=560 bit=70 bayt.
\end{tcolorbox}
\vspace{2mm}\hrule height 0.2pt \vspace{3mm}

% DB question_id: 481b6956-3ca6-5ad6-924f-f8e02de06314
\needspace{5\baselineskip}
\noindent\textbf{\color{ink}838.}
\par
32 768 Hz namunalash chastotasi va 16 bit chuqurlikda 2 daqiqa davomida yozilgan siqilmagan mono audiofayl hajmini MiB da toping.
\begin{enumerate}[label=\Alph*),leftmargin=9mm]
\item 10 MiB
\item 12 MiB
\item 7,5 MiB
\item 0,125 MiB
\end{enumerate}
\begin{tcolorbox}[enhanced,breakable,colback=teal!4,colframe=teal!50!black,boxrule=0.5pt,arc=1.1mm,left=3mm,right=3mm,top=1.7mm,bottom=1.7mm,before skip=2mm,after skip=2mm]
\textbf{\color{teal}To'g'ri javob:} \textbf{C) 7,5 MiB}
\par\smallskip
\textbf{Izoh:} 32 768\ensuremath{\cdot}16\ensuremath{\cdot}120/8=7 864 320 bayt; 7 864 320/1 048 576=7,5 MiB. Mono uchun kanal soni 1.
\end{tcolorbox}
\vspace{2mm}\hrule height 0.2pt \vspace{3mm}

% DB question_id: 489627e2-c8b6-5880-8a93-efdc09b347a2
\needspace{5\baselineskip}
\noindent\textbf{\color{ink}839.}
\par
1080\ensuremath{\times}1024 pikselli siqilmagan rastr tasvirda 126 xil rang ishlatiladi. Har piksel eng kam butun bitlar sonida kodlansa, hajmni KiB da toping.
\begin{enumerate}[label=\Alph*),leftmargin=9mm]
\item 945 KiB
\item 1080 KiB
\item 806,4 KiB
\item 7560 KiB
\end{enumerate}
\begin{tcolorbox}[enhanced,breakable,colback=teal!4,colframe=teal!50!black,boxrule=0.5pt,arc=1.1mm,left=3mm,right=3mm,top=1.7mm,bottom=1.7mm,before skip=2mm,after skip=2mm]
\textbf{\color{teal}To'g'ri javob:} \textbf{A) 945 KiB}
\par\smallskip
\textbf{Izoh:} 126 rang uchun 7 bit/piksel. 1080\ensuremath{\cdot}1024\ensuremath{\cdot}7/(8\ensuremath{\cdot}1024)=945 KiB.
\end{tcolorbox}
\vspace{2mm}\hrule height 0.2pt \vspace{3mm}

% DB question_id: 49ec057a-4c5d-4fee-a62d-81530c955ebd
\needspace{5\baselineskip}
\noindent\textbf{\color{ink}840.}
\par
Siqilmagan stereo audio fayl 48 MB. Shu davomiylikdagi qayta yozuv mono bo'ldi, bit chuqurligi 1,5 baravar oshdi, namuna olish chastotasi 3 baravar kamaydi. Yangi fayl hajmini toping.
\begin{enumerate}[label=\Alph*),leftmargin=9mm]
\item 24 MB
\item 12 MB
\item 18 MB
\item 16 MB
\end{enumerate}
\begin{tcolorbox}[enhanced,breakable,colback=teal!4,colframe=teal!50!black,boxrule=0.5pt,arc=1.1mm,left=3mm,right=3mm,top=1.7mm,bottom=1.7mm,before skip=2mm,after skip=2mm]
\textbf{\color{teal}To'g'ri javob:} \textbf{B) 12 MB}
\par\smallskip
\textbf{Izoh:} Hajm kanallar\ensuremath{\cdot}bit chuqurligi\ensuremath{\cdot}chastotaga proporsional: (1/2)\ensuremath{\cdot}1,5\ensuremath{\cdot}(1/3)=1/4. 48/4=12 MB.
\end{tcolorbox}
\vspace{2mm}\hrule height 0.2pt \vspace{3mm}

% DB question_id: 5741b67a-1bee-5fb1-95f9-5f016ee900c0
\needspace{5\baselineskip}
\noindent\textbf{\color{ink}841.}
\par
999 ta turli ismga bir xil minimal uzunlikdagi ikkilik kod berildi. Shu uzunlikdagi foydalanilmay qolgan kodlar soni nechta?
\begin{enumerate}[label=\Alph*),leftmargin=9mm]
\item 5
\item 32
\item 25
\item 256
\end{enumerate}
\begin{tcolorbox}[enhanced,breakable,colback=teal!4,colframe=teal!50!black,boxrule=0.5pt,arc=1.1mm,left=3mm,right=3mm,top=1.7mm,bottom=1.7mm,before skip=2mm,after skip=2mm]
\textbf{\color{teal}To'g'ri javob:} \textbf{C) 25}
\par\smallskip
\textbf{Izoh:} 2\textsuperscript{9}=512\textless{}999\ensuremath{\leq}1024=2\textsuperscript{1}\textsuperscript{0}, demak 10 bitli 1024 kod mavjud. 1024-999=25 tasi ortib qoladi.
\end{tcolorbox}
\vspace{2mm}\hrule height 0.2pt \vspace{3mm}

% DB question_id: 60fe54ad-b5f0-5b3c-913c-8b18cfc6454a
\needspace{5\baselineskip}
\noindent\textbf{\color{ink}842.}
\par
Mono audio 64 bitli kvantlash va 32 kHz chastotada 10 soniya yozildi. Siqilmagan hajm necha KiB?
\begin{enumerate}[label=\Alph*),leftmargin=9mm]
\item 2500 KiB
\item 2560 KiB
\item 20 000 KiB
\item 312,5 KiB
\end{enumerate}
\begin{tcolorbox}[enhanced,breakable,colback=teal!4,colframe=teal!50!black,boxrule=0.5pt,arc=1.1mm,left=3mm,right=3mm,top=1.7mm,bottom=1.7mm,before skip=2mm,after skip=2mm]
\textbf{\color{teal}To'g'ri javob:} \textbf{A) 2500 KiB}
\par\smallskip
\textbf{Izoh:} 64\ensuremath{\cdot}32 000\ensuremath{\cdot}10=20 480 000 bit=2 560 000 bayt=2500 KiB.
\end{tcolorbox}
\vspace{2mm}\hrule height 0.2pt \vspace{3mm}

% DB question_id: 6102a94c-bbb2-5671-8d3a-d42a1b123334
\needspace{5\baselineskip}
\noindent\textbf{\color{ink}843.}
\par
Bir xil piksel soniga ega siqilmagan rastr fayl 266 rangli formatdan 56 rangli formatga o'tkazilganda hajmi 40 KiB ga kamaydi. Dastlabki hajmni toping.
\begin{enumerate}[label=\Alph*),leftmargin=9mm]
\item 120 KiB
\item 80 KiB
\item 160 KiB
\item 360 KiB
\end{enumerate}
\begin{tcolorbox}[enhanced,breakable,colback=teal!4,colframe=teal!50!black,boxrule=0.5pt,arc=1.1mm,left=3mm,right=3mm,top=1.7mm,bottom=1.7mm,before skip=2mm,after skip=2mm]
\textbf{\color{teal}To'g'ri javob:} \textbf{A) 120 KiB}
\par\smallskip
\textbf{Izoh:} 266 rang uchun 9 bit, 56 rang uchun 6 bit kerak. 3 bitlik kamayish 40 KiB bo'lsa, dastlabki 9 bitlik hajm 40\ensuremath{\cdot}9/3=120 KiB.
\end{tcolorbox}
\vspace{2mm}\hrule height 0.2pt \vspace{3mm}

% DB question_id: 6630c760-3d14-4c36-92b4-dbfa17fe0bac
\needspace{5\baselineskip}
\noindent\textbf{\color{ink}844.}
\par
Masalada piksel, rang chuqurligi va kadr soni berilgani uchun video formulasi tanlandi. Bu vaziyatdagi asosiy tushunchani aniqlang.
\begin{enumerate}[label=\Alph*),leftmargin=9mm]
\item Formula tanlash
\item Birliklarni normallashtirish
\item Xom hajm
\item Siqilgan hajm
\end{enumerate}
\begin{tcolorbox}[enhanced,breakable,colback=teal!4,colframe=teal!50!black,boxrule=0.5pt,arc=1.1mm,left=3mm,right=3mm,top=1.7mm,bottom=1.7mm,before skip=2mm,after skip=2mm]
\textbf{\color{teal}To'g'ri javob:} \textbf{A) Formula tanlash}
\par\smallskip
\textbf{Izoh:} Vaziyat Formula tanlashga misol: berilgan axborot turiga mos hajm formulasini aniqlash bosqichi.
\end{tcolorbox}
\vspace{2mm}\hrule height 0.2pt \vspace{3mm}

% DB question_id: 6c456e90-35b0-4530-a797-44730d917838
\needspace{5\baselineskip}
\noindent\textbf{\color{ink}845.}
\par
40 MB xom fayl 4:1 koeffitsiyent bilan 10 MB bo'ldi. Bu vaziyatdagi asosiy tushunchani aniqlang.
\begin{enumerate}[label=\Alph*),leftmargin=9mm]
\item Ko'p kanalli hisob
\item Kadrlar soni
\item Teskari masala
\item Siqilgan hajm
\end{enumerate}
\begin{tcolorbox}[enhanced,breakable,colback=teal!4,colframe=teal!50!black,boxrule=0.5pt,arc=1.1mm,left=3mm,right=3mm,top=1.7mm,bottom=1.7mm,before skip=2mm,after skip=2mm]
\textbf{\color{teal}To'g'ri javob:} \textbf{D) Siqilgan hajm}
\par\smallskip
\textbf{Izoh:} Vaziyat Siqilgan hajmga misol: siqish algoritmi qo'llangandan keyingi fayl yoki oqim hajmi.
\end{tcolorbox}
\vspace{2mm}\hrule height 0.2pt \vspace{3mm}

% DB question_id: 71c6d49c-52c9-5247-b5b3-e99b5add8930
\needspace{5\baselineskip}
\noindent\textbf{\color{ink}846.}
\par
Har bir lampochka uch holatdan birida bo'ladi: yoniq, o'chiq yoki miltillovchi. Kamida 575 xil signalni kodlash uchun eng kamida nechta mustaqil lampochka kerak?
\begin{enumerate}[label=\Alph*),leftmargin=9mm]
\item 6
\item 5
\item 4
\item 729
\end{enumerate}
\begin{tcolorbox}[enhanced,breakable,colback=teal!4,colframe=teal!50!black,boxrule=0.5pt,arc=1.1mm,left=3mm,right=3mm,top=1.7mm,bottom=1.7mm,before skip=2mm,after skip=2mm]
\textbf{\color{teal}To'g'ri javob:} \textbf{A) 6}
\par\smallskip
\textbf{Izoh:} n ta lampochka 3\textsuperscript{n} holat beradi. 3\textsuperscript{5}=243\textless{}575\ensuremath{\leq}729=3\textsuperscript{6}; demak 6 ta lampochka kerak.
\end{tcolorbox}
\vspace{2mm}\hrule height 0.2pt \vspace{3mm}

% DB question_id: 8155c6d9-6999-5e6c-8350-b7eae9bf5902
\needspace{5\baselineskip}
\noindent\textbf{\color{ink}847.}
\par
Mono audio 32 bitli kvantlash va 16 kHz chastotada 2 minut yozildi. Siqilmagan hajm necha KiB?
\begin{enumerate}[label=\Alph*),leftmargin=9mm]
\item 7500 KiB
\item 7680 KiB
\item 60 000 KiB
\item 3750 KiB
\end{enumerate}
\begin{tcolorbox}[enhanced,breakable,colback=teal!4,colframe=teal!50!black,boxrule=0.5pt,arc=1.1mm,left=3mm,right=3mm,top=1.7mm,bottom=1.7mm,before skip=2mm,after skip=2mm]
\textbf{\color{teal}To'g'ri javob:} \textbf{A) 7500 KiB}
\par\smallskip
\textbf{Izoh:} 32\ensuremath{\cdot}16 000\ensuremath{\cdot}120=61 440 000 bit=7 680 000 bayt=7500 KiB.
\end{tcolorbox}
\vspace{2mm}\hrule height 0.2pt \vspace{3mm}

% DB question_id: 81a33864-cec1-5367-bef4-7db23b88fe84
\needspace{5\baselineskip}
\noindent\textbf{\color{ink}848.}
\par
Kvadrat shakldagi siqilmagan rastr tasvir 56 KiB joy egallaydi. Har piksel 128 rangdan biri bilan kodlanadi. Kvadratning tomoni necha piksel?
\begin{enumerate}[label=\Alph*),leftmargin=9mm]
\item 256 piksel
\item 128 piksel
\item 512 piksel
\item 64 piksel
\end{enumerate}
\begin{tcolorbox}[enhanced,breakable,colback=teal!4,colframe=teal!50!black,boxrule=0.5pt,arc=1.1mm,left=3mm,right=3mm,top=1.7mm,bottom=1.7mm,before skip=2mm,after skip=2mm]
\textbf{\color{teal}To'g'ri javob:} \textbf{A) 256 piksel}
\par\smallskip
\textbf{Izoh:} 128 rang uchun 7 bit. Piksellar soni 56\ensuremath{\cdot}1024\ensuremath{\cdot}8/7=65 536; kvadrat tomoni \ensuremath{\sqrt{\vphantom{x}}}65 536=256 piksel.
\end{tcolorbox}
\vspace{2mm}\hrule height 0.2pt \vspace{3mm}

% DB question_id: 84b6e75b-93b8-4750-8a55-49919bf48284
\needspace{5\baselineskip}
\noindent\textbf{\color{ink}849.}
\par
Siqilmagan stereo audio 64 kHz namuna chastotasi va 24 bit chuqurlikda yozilgan. Fayl hajmi 120 MB. Hisobda 1 MB=1024\textsuperscript{2} bayt deb oling. Yozuv davomiyligiga daqiqalarda eng yaqin, 5 ga bo'linadigan butun sonni toping.
\begin{enumerate}[label=\Alph*),leftmargin=9mm]
\item 5
\item 10
\item 12
\item 6
\end{enumerate}
\begin{tcolorbox}[enhanced,breakable,colback=teal!4,colframe=teal!50!black,boxrule=0.5pt,arc=1.1mm,left=3mm,right=3mm,top=1.7mm,bottom=1.7mm,before skip=2mm,after skip=2mm]
\textbf{\color{teal}To'g'ri javob:} \textbf{A) 5}
\par\smallskip
\textbf{Izoh:} Bit tezligi 64 000\ensuremath{\cdot}24\ensuremath{\cdot}2=3 072 000 bit/s. Hajm 120\ensuremath{\cdot}1024\textsuperscript{2}\ensuremath{\cdot}8 bit, shuning uchun vaqt 327,68 s\ensuremath{\approx}5,46 min. Eng yaqin 5 ga karrali butun son 5.
\end{tcolorbox}
\vspace{2mm}\hrule height 0.2pt \vspace{3mm}

% DB question_id: 87cc89db-f7d0-4f63-9d54-29685edce0bc
\needspace{5\baselineskip}
\noindent\textbf{\color{ink}850.}
\par
Har bir lampochka qizil, yashil, sariq yoki oq rangda yonishi mumkin. Kamida 200 xil signalni kodlash uchun eng kamida nechta lampochka kerak?
\begin{enumerate}[label=\Alph*),leftmargin=9mm]
\item 3 ta
\item 5 ta
\item 8 ta
\item 4 ta
\end{enumerate}
\begin{tcolorbox}[enhanced,breakable,colback=teal!4,colframe=teal!50!black,boxrule=0.5pt,arc=1.1mm,left=3mm,right=3mm,top=1.7mm,bottom=1.7mm,before skip=2mm,after skip=2mm]
\textbf{\color{teal}To'g'ri javob:} \textbf{D) 4 ta}
\par\smallskip
\textbf{Izoh:} n ta lampochka 4\textasciicircum{}n holat beradi. 4\textasciicircum{}3=64\textless{}200, 4\textasciicircum{}4=256\ensuremath{\geq}200; demak 4 ta lampochka kerak. [Yutuqli\_test, test kodi 12; Telegram xabarlar 6446 va 6556.]
\end{tcolorbox}
\vspace{2mm}\hrule height 0.2pt \vspace{3mm}

% DB question_id: 8c269694-fcab-4de1-bb4d-6586a7d002ec
\needspace{5\baselineskip}
\noindent\textbf{\color{ink}851.}
\par
250\ensuremath{\times}300 pikselli siqilmagan rastr tasvir uchun 40 KiB xotira ajratilgan. Har piksel bir xil butun bit chuqurligida kodlansa, palitrada ko'pi bilan nechta rang bo'lishi mumkin?
\begin{enumerate}[label=\Alph*),leftmargin=9mm]
\item 8
\item 32
\item 18
\item 16
\end{enumerate}
\begin{tcolorbox}[enhanced,breakable,colback=teal!4,colframe=teal!50!black,boxrule=0.5pt,arc=1.1mm,left=3mm,right=3mm,top=1.7mm,bottom=1.7mm,before skip=2mm,after skip=2mm]
\textbf{\color{teal}To'g'ri javob:} \textbf{D) 16}
\par\smallskip
\textbf{Izoh:} 40 KiB\ensuremath{\cdot}8=327680 bit. 327680/(250\ensuremath{\cdot}300)\ensuremath{\approx}4,37, demak sig'adigan eng katta butun chuqurlik 4 bit/piksel; 2\textsuperscript{4}=16 rang.
\end{tcolorbox}
\vspace{2mm}\hrule height 0.2pt \vspace{3mm}

% DB question_id: 9197ef2a-3866-5804-a7da-674b6bf13e28
\needspace{5\baselineskip}
\noindent\textbf{\color{ink}852.}
\par
248\ensuremath{\times}360 pikselli siqilmagan rastr tasvir 62 xil rangdan foydalanadi. Har piksel bir xil eng kam butun bitlar soni bilan kodlansa, tasvir hajmi necha bayt?
\begin{enumerate}[label=\Alph*),leftmargin=9mm]
\item 66 960 bayt
\item 535 680 bayt
\item 11 160 bayt
\item 89 280 bayt
\end{enumerate}
\begin{tcolorbox}[enhanced,breakable,colback=teal!4,colframe=teal!50!black,boxrule=0.5pt,arc=1.1mm,left=3mm,right=3mm,top=1.7mm,bottom=1.7mm,before skip=2mm,after skip=2mm]
\textbf{\color{teal}To'g'ri javob:} \textbf{A) 66 960 bayt}
\par\smallskip
\textbf{Izoh:} 62 rang uchun \ensuremath{\lceil}log\textsubscript{2}62\ensuremath{\rceil}=6 bit/piksel. 248\ensuremath{\cdot}360\ensuremath{\cdot}6=535 680 bit=66 960 bayt.
\end{tcolorbox}
\vspace{2mm}\hrule height 0.2pt \vspace{3mm}

% DB question_id: 9468ee3a-509c-4371-b8a8-9e193d8dbdc3
\needspace{5\baselineskip}
\noindent\textbf{\color{ink}853.}
\par
Quyidagi ta'rif qaysi atamaga tegishli? Asl hajmning siqilgan hajmga yoki kamayish foiziga nisbati.
\begin{enumerate}[label=\Alph*),leftmargin=9mm]
\item Kodek
\item Transkodlash
\item Zaxira nusxasi
\item Siqish koeffitsiyenti
\end{enumerate}
\begin{tcolorbox}[enhanced,breakable,colback=teal!4,colframe=teal!50!black,boxrule=0.5pt,arc=1.1mm,left=3mm,right=3mm,top=1.7mm,bottom=1.7mm,before skip=2mm,after skip=2mm]
\textbf{\color{teal}To'g'ri javob:} \textbf{D) Siqish koeffitsiyenti}
\par\smallskip
\textbf{Izoh:} Ta'rif Siqish koeffitsiyenti tushunchasiga tegishli.
\end{tcolorbox}
\vspace{2mm}\hrule height 0.2pt \vspace{3mm}

% DB question_id: 9e5ae3e1-7d19-465f-9292-d02135622e0f
\needspace{5\baselineskip}
\noindent\textbf{\color{ink}854.}
\par
Quyidagi ta'rif qaysi atamaga tegishli? N holat uchun log\textsubscript{2}n yuqoriga yaxlitlangan bitlar soni.
\begin{enumerate}[label=\Alph*),leftmargin=9mm]
\item kB
\item KiB
\item Minimal butun kod uzunligi
\item Xotira sig'imi
\end{enumerate}
\begin{tcolorbox}[enhanced,breakable,colback=teal!4,colframe=teal!50!black,boxrule=0.5pt,arc=1.1mm,left=3mm,right=3mm,top=1.7mm,bottom=1.7mm,before skip=2mm,after skip=2mm]
\textbf{\color{teal}To'g'ri javob:} \textbf{C) Minimal butun kod uzunligi}
\par\smallskip
\textbf{Izoh:} Ta'rif Minimal butun kod uzunligi tushunchasiga tegishli.
\end{tcolorbox}
\vspace{2mm}\hrule height 0.2pt \vspace{3mm}

% DB question_id: a0c22bce-46f2-52f7-809d-29da1783d596
\needspace{5\baselineskip}
\noindent\textbf{\color{ink}855.}
\par
128\ensuremath{\times}128 pikselli siqilmagan rastr tasvir 16 KB xotira egallaydi. Palitradagi maksimal ranglar sonini toping.
\begin{enumerate}[label=\Alph*),leftmargin=9mm]
\item 8 ta
\item 16 ta
\item 256 ta
\item 128 ta
\end{enumerate}
\begin{tcolorbox}[enhanced,breakable,colback=teal!4,colframe=teal!50!black,boxrule=0.5pt,arc=1.1mm,left=3mm,right=3mm,top=1.7mm,bottom=1.7mm,before skip=2mm,after skip=2mm]
\textbf{\color{teal}To'g'ri javob:} \textbf{C) 256 ta}
\par\smallskip
\textbf{Izoh:} Tasvir hajmi \texttt{16\ensuremath{\cdot}1024\ensuremath{\cdot}8=131072} bit, piksel soni \texttt{128\ensuremath{\cdot}128=16384}. Har pikselga 8 bit to'g'ri keladi; \texttt{2\textsuperscript{8}=256} xil rang kodlanadi.
\end{tcolorbox}
\vspace{2mm}\hrule height 0.2pt \vspace{3mm}

% DB question_id: a329626c-1acd-5dd1-abb5-c98969351551
\needspace{5\baselineskip}
\noindent\textbf{\color{ink}856.}
\par
Mono audio 32 bitli kvantlash va 32 kHz chastotada 4 minut yozildi. \texttt{1 MiB = 1 048 576 bayt}. Siqilmagan hajmni eng yaqin butun MiB gacha yaxlitlang.
\begin{enumerate}[label=\Alph*),leftmargin=9mm]
\item 29 MiB
\item 30 MiB
\item 240 MiB
\item 32 MiB
\end{enumerate}
\begin{tcolorbox}[enhanced,breakable,colback=teal!4,colframe=teal!50!black,boxrule=0.5pt,arc=1.1mm,left=3mm,right=3mm,top=1.7mm,bottom=1.7mm,before skip=2mm,after skip=2mm]
\textbf{\color{teal}To'g'ri javob:} \textbf{A) 29 MiB}
\par\smallskip
\textbf{Izoh:} 32\ensuremath{\cdot}32 000\ensuremath{\cdot}240=245 760 000 bit=30 720 000 bayt\ensuremath{\approx}29,297 MiB; yaxlitlanganda 29 MiB.
\end{tcolorbox}
\vspace{2mm}\hrule height 0.2pt \vspace{3mm}

% DB question_id: a345a0c6-54db-5a97-a2af-1d6e8f2cb908
\needspace{5\baselineskip}
\noindent\textbf{\color{ink}857.}
\par
640\ensuremath{\times}480 pikselli siqilmagan rastr tasvir 24 xil rangdan foydalanadi. Har piksel eng kam butun bitlarda kodlansa, hajm necha KiB?
\begin{enumerate}[label=\Alph*),leftmargin=9mm]
\item 187,5 KiB
\item 192 KiB
\item 1536 KiB
\item 150 KiB
\end{enumerate}
\begin{tcolorbox}[enhanced,breakable,colback=teal!4,colframe=teal!50!black,boxrule=0.5pt,arc=1.1mm,left=3mm,right=3mm,top=1.7mm,bottom=1.7mm,before skip=2mm,after skip=2mm]
\textbf{\color{teal}To'g'ri javob:} \textbf{A) 187,5 KiB}
\par\smallskip
\textbf{Izoh:} 24 rang uchun 5 bit/piksel. 640\ensuremath{\cdot}480\ensuremath{\cdot}5/(8\ensuremath{\cdot}1024)=187,5 KiB.
\end{tcolorbox}
\vspace{2mm}\hrule height 0.2pt \vspace{3mm}

% DB question_id: a3783d65-c1c9-402e-94b3-db522943c7dc
\needspace{5\baselineskip}
\noindent\textbf{\color{ink}858.}
\par
Quyidagi ta'rif qaysi atamaga tegishli? Siqish va xizmat ma'lumotlarisiz nazariy axborot hajmi.
\begin{enumerate}[label=\Alph*),leftmargin=9mm]
\item Siqilgan hajm
\item Ko'p kanalli hisob
\item Xom hajm
\item Kadrlar soni
\end{enumerate}
\begin{tcolorbox}[enhanced,breakable,colback=teal!4,colframe=teal!50!black,boxrule=0.5pt,arc=1.1mm,left=3mm,right=3mm,top=1.7mm,bottom=1.7mm,before skip=2mm,after skip=2mm]
\textbf{\color{teal}To'g'ri javob:} \textbf{C) Xom hajm}
\par\smallskip
\textbf{Izoh:} Ta'rif Xom hajm tushunchasiga tegishli.
\end{tcolorbox}
\vspace{2mm}\hrule height 0.2pt \vspace{3mm}

% DB question_id: a38042dd-c365-4295-ba42-dba470fa42d5
\needspace{5\baselineskip}
\noindent\textbf{\color{ink}859.}
\par
16 ta teng ehtimolli belgi uchun bitta belgining kod uzunligi topildi. Bu vaziyatdagi asosiy tushunchani aniqlang.
\begin{enumerate}[label=\Alph*),leftmargin=9mm]
\item Xartli formulasi
\item Alifbo quvvati
\item Minimal butun kod uzunligi
\item kB
\end{enumerate}
\begin{tcolorbox}[enhanced,breakable,colback=teal!4,colframe=teal!50!black,boxrule=0.5pt,arc=1.1mm,left=3mm,right=3mm,top=1.7mm,bottom=1.7mm,before skip=2mm,after skip=2mm]
\textbf{\color{teal}To'g'ri javob:} \textbf{A) Xartli formulasi}
\par\smallskip
\textbf{Izoh:} Vaziyat Xartli formulasiga misol: teng ehtimolli N holatdan birini kodlash uchun i = log\textsubscript{2}N munosabatini beradigan formula.
\end{tcolorbox}
\vspace{2mm}\hrule height 0.2pt \vspace{3mm}

% DB question_id: a3cdf164-b0a1-50a7-b722-14588d6c8b5f
\needspace{5\baselineskip}
\noindent\textbf{\color{ink}860.}
\par
Mono audio 8 bitli kvantlash chuqurligi va 25 000 Hz chastotada 4 minut yozildi. Siqilmagan hajmni eng yaqin butun MiB gacha yaxlitlang. \texttt{1 MiB = 1 048 576 bayt}.
\begin{enumerate}[label=\Alph*),leftmargin=9mm]
\item 6 MiB
\item 5 MiB
\item 48 MiB
\item 6000 MiB
\end{enumerate}
\begin{tcolorbox}[enhanced,breakable,colback=teal!4,colframe=teal!50!black,boxrule=0.5pt,arc=1.1mm,left=3mm,right=3mm,top=1.7mm,bottom=1.7mm,before skip=2mm,after skip=2mm]
\textbf{\color{teal}To'g'ri javob:} \textbf{A) 6 MiB}
\par\smallskip
\textbf{Izoh:} 8\ensuremath{\cdot}25 000\ensuremath{\cdot}240=48 000 000 bit=6 000 000 bayt\ensuremath{\approx}5,722 MiB; eng yaqin butun qiymat 6 MiB.
\end{tcolorbox}
\vspace{2mm}\hrule height 0.2pt \vspace{3mm}

% DB question_id: aa7e9cf2-8fb8-5eb3-8709-ac61fea3d047
\needspace{5\baselineskip}
\noindent\textbf{\color{ink}861.}
\par
Mono audio 16 bitli kvantlash va 32 kHz namuna olish chastotasida 3 minut yozildi. \texttt{1 MiB = 1 048 576 bayt}. Siqilmagan hajmni eng yaqin butun MiB gacha yaxlitlang.
\begin{enumerate}[label=\Alph*),leftmargin=9mm]
\item 11 MiB
\item 90 MiB
\item 12 MiB
\item 10 MiB
\end{enumerate}
\begin{tcolorbox}[enhanced,breakable,colback=teal!4,colframe=teal!50!black,boxrule=0.5pt,arc=1.1mm,left=3mm,right=3mm,top=1.7mm,bottom=1.7mm,before skip=2mm,after skip=2mm]
\textbf{\color{teal}To'g'ri javob:} \textbf{A) 11 MiB}
\par\smallskip
\textbf{Izoh:} 16\ensuremath{\cdot}32 000\ensuremath{\cdot}180=92 160 000 bit=11 520 000 bayt\ensuremath{\approx}10,986 MiB; yaxlitlanganda 11 MiB.
\end{tcolorbox}
\vspace{2mm}\hrule height 0.2pt \vspace{3mm}

% DB question_id: b29a2951-1d61-49f1-a94c-181bf8a9b97e
\needspace{5\baselineskip}
\noindent\textbf{\color{ink}862.}
\par
Quyidagi ta'rif qaysi atamaga tegishli? Natijaning birlik, tartib va shartga mosligini baholash bosqichi.
\begin{enumerate}[label=\Alph*),leftmargin=9mm]
\item Formula tanlash
\item Birliklarni normallashtirish
\item Xom hajm
\item Mantiqiy tekshiruv
\end{enumerate}
\begin{tcolorbox}[enhanced,breakable,colback=teal!4,colframe=teal!50!black,boxrule=0.5pt,arc=1.1mm,left=3mm,right=3mm,top=1.7mm,bottom=1.7mm,before skip=2mm,after skip=2mm]
\textbf{\color{teal}To'g'ri javob:} \textbf{D) Mantiqiy tekshiruv}
\par\smallskip
\textbf{Izoh:} Ta'rif Mantiqiy tekshiruv tushunchasiga tegishli.
\end{tcolorbox}
\vspace{2mm}\hrule height 0.2pt \vspace{3mm}

% DB question_id: b5775f39-8e3b-5493-a5da-6a223295e5cb
\needspace{5\baselineskip}
\noindent\textbf{\color{ink}863.}
\par
Matnli faylning har bir sahifasida 32 qator, har qatorda 64 ta belgi bor. 80 sahifali faylda har bir belgi 1 baytli ASCII kodida saqlansa, faylning axborot hajmini KiB da toping.
\begin{enumerate}[label=\Alph*),leftmargin=9mm]
\item 80 KiB
\item 160 KiB
\item 320 KiB
\item 40 KiB
\item 120 KiB
\end{enumerate}
\begin{tcolorbox}[enhanced,breakable,colback=teal!4,colframe=teal!50!black,boxrule=0.5pt,arc=1.1mm,left=3mm,right=3mm,top=1.7mm,bottom=1.7mm,before skip=2mm,after skip=2mm]
\textbf{\color{teal}To'g'ri javob:} \textbf{B) 160 KiB}
\par\smallskip
\textbf{Izoh:} 32\ensuremath{\cdot}64\ensuremath{\cdot}80=163 840 bayt. 163 840/1024=160 KiB.
\end{tcolorbox}
\vspace{2mm}\hrule height 0.2pt \vspace{3mm}

% DB question_id: bdb9b49f-0d61-4966-8b94-5feff092d6fc
\needspace{5\baselineskip}
\noindent\textbf{\color{ink}864.}
\par
Quyidagi ta'rif qaysi atamaga tegishli? Siqish algoritmi qo'llangandan keyingi fayl yoki oqim hajmi.
\begin{enumerate}[label=\Alph*),leftmargin=9mm]
\item Ko'p kanalli hisob
\item Kadrlar soni
\item Teskari masala
\item Siqilgan hajm
\end{enumerate}
\begin{tcolorbox}[enhanced,breakable,colback=teal!4,colframe=teal!50!black,boxrule=0.5pt,arc=1.1mm,left=3mm,right=3mm,top=1.7mm,bottom=1.7mm,before skip=2mm,after skip=2mm]
\textbf{\color{teal}To'g'ri javob:} \textbf{D) Siqilgan hajm}
\par\smallskip
\textbf{Izoh:} Ta'rif Siqilgan hajm tushunchasiga tegishli.
\end{tcolorbox}
\vspace{2mm}\hrule height 0.2pt \vspace{3mm}

% DB question_id: ca9bc578-659b-5c2c-8cfd-79e7f9885055
\needspace{5\baselineskip}
\noindent\textbf{\color{ink}865.}
\par
192\ensuremath{\times}160 pikselli siqilmagan rastr tasvir 30 KiB joy egallaydi. Har piksel bir xil butun bitlarda kodlansa, tasvir ko'pi bilan nechta rangni ifodalay oladi?
\begin{enumerate}[label=\Alph*),leftmargin=9mm]
\item 256 ta
\item 8 ta
\item 128 ta
\item 65 536 ta
\end{enumerate}
\begin{tcolorbox}[enhanced,breakable,colback=teal!4,colframe=teal!50!black,boxrule=0.5pt,arc=1.1mm,left=3mm,right=3mm,top=1.7mm,bottom=1.7mm,before skip=2mm,after skip=2mm]
\textbf{\color{teal}To'g'ri javob:} \textbf{A) 256 ta}
\par\smallskip
\textbf{Izoh:} 30\ensuremath{\cdot}1024\ensuremath{\cdot}8/(192\ensuremath{\cdot}160)=8 bit/piksel; 2\textsuperscript{8}=256 rang.
\end{tcolorbox}
\vspace{2mm}\hrule height 0.2pt \vspace{3mm}

% DB question_id: cb4b033a-aea2-4631-b5ba-2dbce927d776
\needspace{5\baselineskip}
\noindent\textbf{\color{ink}866.}
\par
Kod jadvalida 32 xil belgi mavjudligi sanaldi. Bu vaziyatdagi asosiy tushunchani aniqlang.
\begin{enumerate}[label=\Alph*),leftmargin=9mm]
\item Minimal butun kod uzunligi
\item Alifbo quvvati
\item kB
\item KiB
\end{enumerate}
\begin{tcolorbox}[enhanced,breakable,colback=teal!4,colframe=teal!50!black,boxrule=0.5pt,arc=1.1mm,left=3mm,right=3mm,top=1.7mm,bottom=1.7mm,before skip=2mm,after skip=2mm]
\textbf{\color{teal}To'g'ri javob:} \textbf{B) Alifbo quvvati}
\par\smallskip
\textbf{Izoh:} Vaziyat Alifbo quvvatiga misol: kodlashda qatnashadigan turli belgilar yoki holatlar soni.
\end{tcolorbox}
\vspace{2mm}\hrule height 0.2pt \vspace{3mm}

% DB question_id: d942db05-c061-53ac-8484-26898ad5cc41
\needspace{5\baselineskip}
\noindent\textbf{\color{ink}867.}
\par
\texttt{Yoshlikda olingan bilim toshga bitilgan yozuv kabi o'chmasdir!} satrida bo'sh joy, apostrof va undov belgisi ham hisobga olinadi. Har bir belgi 8 bit bo'lsa, axborot hajmini o'n oltilik sanoq sistemasida toping.
\begin{enumerate}[label=\Alph*),leftmargin=9mm]
\item 1B7\textsubscript{1}\textsubscript{6}
\item 1B8\textsubscript{1}\textsubscript{6}
\item 1F0\textsubscript{1}\textsubscript{6}
\item 1B0\textsubscript{1}\textsubscript{6}
\end{enumerate}
\begin{tcolorbox}[enhanced,breakable,colback=teal!4,colframe=teal!50!black,boxrule=0.5pt,arc=1.1mm,left=3mm,right=3mm,top=1.7mm,bottom=1.7mm,before skip=2mm,after skip=2mm]
\textbf{\color{teal}To'g'ri javob:} \textbf{C) 1F0\textsubscript{1}\textsubscript{6}}
\par\smallskip
\textbf{Izoh:} Qo'shtirnoqlarsiz satrda 62 belgi bor. 62\ensuremath{\cdot}8=496 bit; 496\textsubscript{1}\textsubscript{0}=1F0\textsubscript{1}\textsubscript{6}.
\end{tcolorbox}
\vspace{2mm}\hrule height 0.2pt \vspace{3mm}

% DB question_id: dc513599-1d75-5bf3-bcd0-abede856493b
\needspace{5\baselineskip}
\noindent\textbf{\color{ink}868.}
\par
Mono audio 64 bitli kvantlash chuqurligi va 32 000 Hz chastotada 10 soniya yozildi. Siqilmagan hajmni eng yaqin butun MiB gacha yaxlitlang. \texttt{1 MiB = 1 048 576 bayt}.
\begin{enumerate}[label=\Alph*),leftmargin=9mm]
\item 2 MiB
\item 3 MiB
\item 20 MiB
\item 2,56 MiB
\end{enumerate}
\begin{tcolorbox}[enhanced,breakable,colback=teal!4,colframe=teal!50!black,boxrule=0.5pt,arc=1.1mm,left=3mm,right=3mm,top=1.7mm,bottom=1.7mm,before skip=2mm,after skip=2mm]
\textbf{\color{teal}To'g'ri javob:} \textbf{A) 2 MiB}
\par\smallskip
\textbf{Izoh:} 64\ensuremath{\cdot}32 000\ensuremath{\cdot}10=20 480 000 bit=2 560 000 bayt\ensuremath{\approx}2,441 MiB; eng yaqin butun qiymat 2 MiB.
\end{tcolorbox}
\vspace{2mm}\hrule height 0.2pt \vspace{3mm}

% DB question_id: df6a5b28-2ee4-4e0a-bc9b-bd0981106efc
\needspace{5\baselineskip}
\noindent\textbf{\color{ink}869.}
\par
Quyidagi ta'rif qaysi atamaga tegishli? Teng ehtimolli n holatdan birini kodlash uchun i = log\textsubscript{2}n munosabatini beradigan formula.
\begin{enumerate}[label=\Alph*),leftmargin=9mm]
\item Xartli formulasi
\item Alifbo quvvati
\item Minimal butun kod uzunligi
\item kB
\end{enumerate}
\begin{tcolorbox}[enhanced,breakable,colback=teal!4,colframe=teal!50!black,boxrule=0.5pt,arc=1.1mm,left=3mm,right=3mm,top=1.7mm,bottom=1.7mm,before skip=2mm,after skip=2mm]
\textbf{\color{teal}To'g'ri javob:} \textbf{A) Xartli formulasi}
\par\smallskip
\textbf{Izoh:} Ta'rif Xartli formulasi tushunchasiga tegishli.
\end{tcolorbox}
\vspace{2mm}\hrule height 0.2pt \vspace{3mm}

% DB question_id: e02ffe84-8b7f-58b8-9e6b-be5d744ca26d
\needspace{5\baselineskip}
\noindent\textbf{\color{ink}870.}
\par
Parol 9 ta belgi (\texttt{A}-\texttt{I}) alifbosidan tuzilgan 25 ta belgidan iborat. Har belgi bir xil eng kam bitlarda, har bir parol esa alohida eng kichik butun baytlar sonida saqlanadi. 15 ta parol uchun qancha bayt kerak?
\begin{enumerate}[label=\Alph*),leftmargin=9mm]
\item 195 bayt
\item 188 bayt
\item 200 bayt
\item 375 bayt
\end{enumerate}
\begin{tcolorbox}[enhanced,breakable,colback=teal!4,colframe=teal!50!black,boxrule=0.5pt,arc=1.1mm,left=3mm,right=3mm,top=1.7mm,bottom=1.7mm,before skip=2mm,after skip=2mm]
\textbf{\color{teal}To'g'ri javob:} \textbf{A) 195 bayt}
\par\smallskip
\textbf{Izoh:} 9 belgi uchun 4 bit kerak. Bir parol 25\ensuremath{\cdot}4=100 bit, ya'ni alohida saqlanganda 13 bayt. 15\ensuremath{\cdot}13=195 bayt.
\end{tcolorbox}
\vspace{2mm}\hrule height 0.2pt \vspace{3mm}

% DB question_id: e660b2c5-a42e-5c9d-b4ab-c317a0bec72e
\needspace{5\baselineskip}
\noindent\textbf{\color{ink}871.}
\par
Bir xil piksel sonli siqilmagan rastr fayl 64 rangli formatdan 4 rangli formatga o'tkazilganda hajmi 24 KiB ga kamaydi. Dastlabki hajm necha KiB edi?
\begin{enumerate}[label=\Alph*),leftmargin=9mm]
\item 36 KiB
\item 24 KiB
\item 48 KiB
\item 72 KiB
\end{enumerate}
\begin{tcolorbox}[enhanced,breakable,colback=teal!4,colframe=teal!50!black,boxrule=0.5pt,arc=1.1mm,left=3mm,right=3mm,top=1.7mm,bottom=1.7mm,before skip=2mm,after skip=2mm]
\textbf{\color{teal}To'g'ri javob:} \textbf{A) 36 KiB}
\par\smallskip
\textbf{Izoh:} 64 rang=6 bit, 4 rang=2 bit. 4 bitlik kamayish 24 KiB bo'lsa, dastlabki 6 bitlik hajm 24\ensuremath{\cdot}6/4=36 KiB.
\end{tcolorbox}
\vspace{2mm}\hrule height 0.2pt \vspace{3mm}

% DB question_id: ef5700de-6c07-5613-a74a-1eb585446028
\needspace{5\baselineskip}
\noindent\textbf{\color{ink}872.}
\par
Bir xil piksel sonli siqilmagan rastr fayl 128 rangli formatdan 8 rangli formatga o'tkazilganda hajmi 32 KiB ga kamaydi. Dastlabki hajm necha KiB edi?
\begin{enumerate}[label=\Alph*),leftmargin=9mm]
\item 56 KiB
\item 32 KiB
\item 24 KiB
\item 96 KiB
\end{enumerate}
\begin{tcolorbox}[enhanced,breakable,colback=teal!4,colframe=teal!50!black,boxrule=0.5pt,arc=1.1mm,left=3mm,right=3mm,top=1.7mm,bottom=1.7mm,before skip=2mm,after skip=2mm]
\textbf{\color{teal}To'g'ri javob:} \textbf{A) 56 KiB}
\par\smallskip
\textbf{Izoh:} 128 rang=7 bit, 8 rang=3 bit. 4 bitlik kamayish 32 KiB bo'lsa, dastlabki 7 bitlik hajm 32\ensuremath{\cdot}7/4=56 KiB.
\end{tcolorbox}
\vspace{2mm}\hrule height 0.2pt \vspace{3mm}

% DB question_id: efb56701-8923-5c4e-9a1b-dfeac1685c86
\needspace{5\baselineskip}
\noindent\textbf{\color{ink}873.}
\par
1024\ensuremath{\times}640 pikselli siqilmagan rastr tasvir 56 xil rangdan foydalanadi. Har piksel eng kam butun bitlarda kodlansa, hajm necha KiB bo'ladi?
\begin{enumerate}[label=\Alph*),leftmargin=9mm]
\item 480 KiB
\item 3840 KiB
\item 109,375 KiB
\item 640 KiB
\end{enumerate}
\begin{tcolorbox}[enhanced,breakable,colback=teal!4,colframe=teal!50!black,boxrule=0.5pt,arc=1.1mm,left=3mm,right=3mm,top=1.7mm,bottom=1.7mm,before skip=2mm,after skip=2mm]
\textbf{\color{teal}To'g'ri javob:} \textbf{A) 480 KiB}
\par\smallskip
\textbf{Izoh:} 56 rang uchun 6 bit/piksel. 1024\ensuremath{\cdot}640\ensuremath{\cdot}6/(8\ensuremath{\cdot}1024)=480 KiB.
\end{tcolorbox}
\vspace{2mm}\hrule height 0.2pt \vspace{3mm}

% DB question_id: f501f9ee-e7d1-5394-b596-9377d35f1a0e
\needspace{5\baselineskip}
\noindent\textbf{\color{ink}874.}
\par
O'quvchi minutiga 160 ta so'z o'qiydi; bir so'z o'rtacha 6 ta belgi va har belgi 1 bayt. 4 soatda ko'zdan kechirilgan axborot hajmini KiBda toping. 1 KiB=1024 bayt.
\begin{enumerate}[label=\Alph*),leftmargin=9mm]
\item 225 KiB
\item 220 KiB
\item 120 KiB
\item 240 KiB
\end{enumerate}
\begin{tcolorbox}[enhanced,breakable,colback=teal!4,colframe=teal!50!black,boxrule=0.5pt,arc=1.1mm,left=3mm,right=3mm,top=1.7mm,bottom=1.7mm,before skip=2mm,after skip=2mm]
\textbf{\color{teal}To'g'ri javob:} \textbf{A) 225 KiB}
\par\smallskip
\textbf{Izoh:} 4 soat=240 minut. 160\ensuremath{\cdot}6\ensuremath{\cdot}240=230 400 bayt; 230 400/1024=225 KiB.
\end{tcolorbox}
\vspace{2mm}\hrule height 0.2pt \vspace{3mm}

% DB question_id: f5ce0bbc-0699-4644-8573-bd35e89df86c
\needspace{5\baselineskip}
\noindent\textbf{\color{ink}875.}
\par
Quyidagi ta'rif qaysi atamaga tegishli? Hajm va boshqa parametrlar ma'lum bo'lganda noma'lum parametrni formuladan ajratib topish.
\begin{enumerate}[label=\Alph*),leftmargin=9mm]
\item Mantiqiy tekshiruv
\item Formula tanlash
\item Teskari masala
\item Birliklarni normallashtirish
\end{enumerate}
\begin{tcolorbox}[enhanced,breakable,colback=teal!4,colframe=teal!50!black,boxrule=0.5pt,arc=1.1mm,left=3mm,right=3mm,top=1.7mm,bottom=1.7mm,before skip=2mm,after skip=2mm]
\textbf{\color{teal}To'g'ri javob:} \textbf{C) Teskari masala}
\par\smallskip
\textbf{Izoh:} Ta'rif Teskari masala tushunchasiga tegishli.
\end{tcolorbox}
\vspace{2mm}\hrule height 0.2pt \vspace{3mm}

% DB question_id: f6dc74cd-2eb6-4b48-b51b-204c8b2040d1
\needspace{5\baselineskip}
\noindent\textbf{\color{ink}876.}
\par
Quyidagi ta'rif qaysi atamaga tegishli? Fps va davomiylik ko'paytmasi orqali topiladigan video tasvirlari soni.
\begin{enumerate}[label=\Alph*),leftmargin=9mm]
\item Teskari masala
\item Kadrlar soni
\item Mantiqiy tekshiruv
\item Formula tanlash
\end{enumerate}
\begin{tcolorbox}[enhanced,breakable,colback=teal!4,colframe=teal!50!black,boxrule=0.5pt,arc=1.1mm,left=3mm,right=3mm,top=1.7mm,bottom=1.7mm,before skip=2mm,after skip=2mm]
\textbf{\color{teal}To'g'ri javob:} \textbf{B) Kadrlar soni}
\par\smallskip
\textbf{Izoh:} Ta'rif Kadrlar soni tushunchasiga tegishli.
\end{tcolorbox}
\vspace{2mm}\hrule height 0.2pt \vspace{3mm}

% DB question_id: f80a6c80-f33d-497a-989c-75da1e558a62
\needspace{5\baselineskip}
\noindent\textbf{\color{ink}877.}
\par
Bir xil belgilar ketma-ketligi dastlab har belgi uchun 16 bitli kodda, keyin har belgi uchun 8 bitli kodda saqlandi. Fayl hajmi 40 baytga kamaydi. Dastlabki hajm qancha edi?
\begin{enumerate}[label=\Alph*),leftmargin=9mm]
\item 40 bayt
\item 60 bayt
\item 120 bayt
\item 80 bayt
\end{enumerate}
\begin{tcolorbox}[enhanced,breakable,colback=teal!4,colframe=teal!50!black,boxrule=0.5pt,arc=1.1mm,left=3mm,right=3mm,top=1.7mm,bottom=1.7mm,before skip=2mm,after skip=2mm]
\textbf{\color{teal}To'g'ri javob:} \textbf{D) 80 bayt}
\par\smallskip
\textbf{Izoh:} Har bir belgida 8 bit=1 bayt tejaladi. 40 bayt kamayish 40 belgini bildiradi; dastlab 40\ensuremath{\cdot}2=80 bayt.
\end{tcolorbox}
\vspace{2mm}\hrule height 0.2pt \vspace{3mm}

% DB question_id: f93e112c-f247-593b-acf3-5077f0d85959
\needspace{5\baselineskip}
\noindent\textbf{\color{ink}878.}
\par
Siqilmagan rastr tasvirning eni 384 piksel, hajmi 90 KiB va palitrasi 8 rangli. Tasvir bo'yi necha piksel? \texttt{1 KiB = 1024 bayt}.
\begin{enumerate}[label=\Alph*),leftmargin=9mm]
\item 640 piksel
\item 320 piksel
\item 480 piksel
\item 720 piksel
\end{enumerate}
\begin{tcolorbox}[enhanced,breakable,colback=teal!4,colframe=teal!50!black,boxrule=0.5pt,arc=1.1mm,left=3mm,right=3mm,top=1.7mm,bottom=1.7mm,before skip=2mm,after skip=2mm]
\textbf{\color{teal}To'g'ri javob:} \textbf{A) 640 piksel}
\par\smallskip
\textbf{Izoh:} 8 rang uchun 3 bit/piksel. Balandlik=(90\ensuremath{\cdot}1024\ensuremath{\cdot}8)/(384\ensuremath{\cdot}3)=640 piksel.
\end{tcolorbox}
\vspace{2mm}\hrule height 0.2pt \vspace{3mm}

% DB question_id: fed76e50-829c-4790-acd6-0946de029b53
\needspace{5\baselineskip}
\noindent\textbf{\color{ink}879.}
\par
5 daqiqa 300 soniyaga, 8 Mbit/s esa bit/s ga aylantirildi. Bu vaziyatdagi asosiy tushunchani aniqlang.
\begin{enumerate}[label=\Alph*),leftmargin=9mm]
\item Xom hajm
\item Birliklarni normallashtirish
\item Siqilgan hajm
\item Ko'p kanalli hisob
\end{enumerate}
\begin{tcolorbox}[enhanced,breakable,colback=teal!4,colframe=teal!50!black,boxrule=0.5pt,arc=1.1mm,left=3mm,right=3mm,top=1.7mm,bottom=1.7mm,before skip=2mm,after skip=2mm]
\textbf{\color{teal}To'g'ri javob:} \textbf{B) Birliklarni normallashtirish}
\par\smallskip
\textbf{Izoh:} Vaziyat Birliklarni normallashtirishga misol: hisoblashdan oldin bit-bayt, soniya-daqiqa va prefikslarni bir xil tizimga keltirish.
\end{tcolorbox}
\vspace{2mm}\hrule height 0.2pt \vspace{3mm}

% DB question_id: 0026e08f-3288-559a-8799-eb975cecca8f
\needspace{5\baselineskip}
\noindent\textbf{\color{ink}880.}
\par
154403 ta belgidan iborat matnda har bir belgi 8 bit bilan kodlangan. Matn hajmini baytda toping.
\begin{enumerate}[label=\Alph*),leftmargin=9mm]
\item 154403 bayt
\item 1235224 bayt
\item 77201.5 bayt
\item 308806 bayt
\end{enumerate}
\begin{tcolorbox}[enhanced,breakable,colback=teal!4,colframe=teal!50!black,boxrule=0.5pt,arc=1.1mm,left=3mm,right=3mm,top=1.7mm,bottom=1.7mm,before skip=2mm,after skip=2mm]
\textbf{\color{teal}To'g'ri javob:} \textbf{A) 154403 bayt}
\par\smallskip
\textbf{Izoh:} 154403 \ensuremath{\times} 8 = 1235224 bit; 8 ga bo'lsak 154403 bayt.
\end{tcolorbox}
\vspace{2mm}\hrule height 0.2pt \vspace{3mm}

% DB question_id: 0aa658b9-cbc1-47d3-8962-876abe165861
\needspace{5\baselineskip}
\noindent\textbf{\color{ink}881.}
\par
192\ensuremath{\times}160 pikselli tasvir 30 Kbayt xotira egallaydi. Siqishsiz bitta piksel uchun nechta rangni kodlash mumkin?
\begin{enumerate}[label=\Alph*),leftmargin=9mm]
\item 312
\item 256
\item 194
\item 128
\end{enumerate}
\begin{tcolorbox}[enhanced,breakable,colback=teal!4,colframe=teal!50!black,boxrule=0.5pt,arc=1.1mm,left=3mm,right=3mm,top=1.7mm,bottom=1.7mm,before skip=2mm,after skip=2mm]
\textbf{\color{teal}To'g'ri javob:} \textbf{B) 256}
\par\smallskip
\textbf{Izoh:} 30720 piksel va 30720 bayt, har piksel 8 bit: 256 rang.
\end{tcolorbox}
\vspace{2mm}\hrule height 0.2pt \vspace{3mm}

% DB question_id: 0ccda2c0-f94a-565c-b27e-4c523e99e34f
\needspace{5\baselineskip}
\noindent\textbf{\color{ink}882.}
\par
Siqilmagan rastrda ranglar soni 1 048 576 tadan 32 768 tagacha kamaytirildi va fayl hajmi 1 MiB ga qisqardi. Piksel soni o'zgarmagan, sarlavha hisobga olinmaydi. Dastlabki fayl hajmini toping. 1 MiB=2\textsuperscript{2}\textsuperscript{0} bayt.
\begin{enumerate}[label=\Alph*),leftmargin=9mm]
\item 32 MiB
\item 5 MiB
\item 4 MiB
\item 2 MiB
\end{enumerate}
\begin{tcolorbox}[enhanced,breakable,colback=teal!4,colframe=teal!50!black,boxrule=0.5pt,arc=1.1mm,left=3mm,right=3mm,top=1.7mm,bottom=1.7mm,before skip=2mm,after skip=2mm]
\textbf{\color{teal}To'g'ri javob:} \textbf{C) 4 MiB}
\par\smallskip
\textbf{Izoh:} 2\textsuperscript{2}\textsuperscript{0} rang uchun 20 bit/piksel, 2\textsuperscript{1}\textsuperscript{5} rang uchun 15 bit/piksel kerak. 5 bit/piksellik kamayish 1 MiB bo'lsa, dastlabki 20 bit/piksellik hajm 20/5\ensuremath{\cdot}1=4 MiB.
\end{tcolorbox}
\vspace{2mm}\hrule height 0.2pt \vspace{3mm}

% DB question_id: 0e94c1c4-b628-5299-a06e-43715ca5a8e3
\needspace{5\baselineskip}
\noindent\textbf{\color{ink}883.}
\par
64706 ta belgidan iborat matnda har bir belgi 16 bit bilan kodlangan. Matn hajmini baytda toping.
\begin{enumerate}[label=\Alph*),leftmargin=9mm]
\item 64706 bayt
\item 258824 bayt
\item 1035296 bayt
\item 129412 bayt
\end{enumerate}
\begin{tcolorbox}[enhanced,breakable,colback=teal!4,colframe=teal!50!black,boxrule=0.5pt,arc=1.1mm,left=3mm,right=3mm,top=1.7mm,bottom=1.7mm,before skip=2mm,after skip=2mm]
\textbf{\color{teal}To'g'ri javob:} \textbf{D) 129412 bayt}
\par\smallskip
\textbf{Izoh:} 64706 \ensuremath{\times} 16 = 1035296 bit; 8 ga bo'lsak 129412 bayt.
\end{tcolorbox}
\vspace{2mm}\hrule height 0.2pt \vspace{3mm}

% DB question_id: 0efddcfa-85f0-464f-a481-55a875f833e0
\needspace{5\baselineskip}
\noindent\textbf{\color{ink}884.}
\par
<<ATTESTATIYA>> degan 11 ta ASCII belgidan iborat satrni 824 marta bo'shliqsiz yozib saqladik. Har belgi 8 bit, boshqa xizmat bayti yo'q. Hajmni KiBda hisoblab eng yaqin butun songa yaxlitlang.
\begin{enumerate}[label=\Alph*),leftmargin=9mm]
\item 7 KiB
\item 8 KiB
\item 9 KiB
\item 10 KiB
\end{enumerate}
\begin{tcolorbox}[enhanced,breakable,colback=teal!4,colframe=teal!50!black,boxrule=0.5pt,arc=1.1mm,left=3mm,right=3mm,top=1.7mm,bottom=1.7mm,before skip=2mm,after skip=2mm]
\textbf{\color{teal}To'g'ri javob:} \textbf{C) 9 KiB}
\par\smallskip
\textbf{Izoh:} <<ATTESTATIYA>> 11 ta lotin harfidan iborat. 824 marta yozilganda 11\ensuremath{\times}824=9064 bayt bo'ladi. 9064/1024\ensuremath{\approx}8,852 KiB; eng yaqin butun son 9 KiB.
\end{tcolorbox}
\vspace{2mm}\hrule height 0.2pt \vspace{3mm}

% DB question_id: 11142174-4632-5e34-b51a-8c8e07591c41
\needspace{5\baselineskip}
\noindent\textbf{\color{ink}885.}
\par
161364 ta belgidan iborat matnda har bir belgi 8 bit bilan kodlangan. Matn hajmini baytda toping.
\begin{enumerate}[label=\Alph*),leftmargin=9mm]
\item 161364 bayt
\item 80682 bayt
\item 322728 bayt
\item 1290912 bayt
\end{enumerate}
\begin{tcolorbox}[enhanced,breakable,colback=teal!4,colframe=teal!50!black,boxrule=0.5pt,arc=1.1mm,left=3mm,right=3mm,top=1.7mm,bottom=1.7mm,before skip=2mm,after skip=2mm]
\textbf{\color{teal}To'g'ri javob:} \textbf{A) 161364 bayt}
\par\smallskip
\textbf{Izoh:} 161364 \ensuremath{\times} 8 = 1290912 bit; 8 ga bo'lsak 161364 bayt.
\end{tcolorbox}
\vspace{2mm}\hrule height 0.2pt \vspace{3mm}

% DB question_id: 11e20434-1e21-4a01-89e0-a62a243ab911
\needspace{5\baselineskip}
\noindent\textbf{\color{ink}886.}
\par
Har bir lampochka qizil yoki yashil yonadi. Kamida 200 xil signalni kodlash uchun nechta lampochka zarur?
\begin{enumerate}[label=\Alph*),leftmargin=9mm]
\item 7 ta
\item 9 ta
\item 8 ta
\item 10 ta
\end{enumerate}
\begin{tcolorbox}[enhanced,breakable,colback=teal!4,colframe=teal!50!black,boxrule=0.5pt,arc=1.1mm,left=3mm,right=3mm,top=1.7mm,bottom=1.7mm,before skip=2mm,after skip=2mm]
\textbf{\color{teal}To'g'ri javob:} \textbf{C) 8 ta}
\par\smallskip
\textbf{Izoh:} 2\textsuperscript{7}=128\textless{}200\ensuremath{\leq}256=2\textsuperscript{8}, minimal 8 ta. [2025 kuzgi demo YECHIMLARI, sahifa 4.]
\end{tcolorbox}
\vspace{2mm}\hrule height 0.2pt \vspace{3mm}

% DB question_id: 1479ecbe-f163-4b52-af62-bb64d3b81726
\needspace{5\baselineskip}
\noindent\textbf{\color{ink}887.}
\par
640 \ensuremath{\times} 480 pikselli, 24 bitli rasmning siqilmagan hajmi qancha?
\begin{enumerate}[label=\Alph*),leftmargin=9mm]
\item 921 600 B
\item 307 200 B
\item 614 400 B
\item 7 372 800 B
\end{enumerate}
\begin{tcolorbox}[enhanced,breakable,colback=teal!4,colframe=teal!50!black,boxrule=0.5pt,arc=1.1mm,left=3mm,right=3mm,top=1.7mm,bottom=1.7mm,before skip=2mm,after skip=2mm]
\textbf{\color{teal}To'g'ri javob:} \textbf{A) 921 600 B}
\par\smallskip
\textbf{Izoh:} 640 \ensuremath{\times} 480 \ensuremath{\times} 24 \ensuremath{\div} 8 = 921 600 bayt.
\end{tcolorbox}
\vspace{2mm}\hrule height 0.2pt \vspace{3mm}

% DB question_id: 1c82bbb9-1109-4aff-9425-ed98ddc1bee4
\needspace{5\baselineskip}
\noindent\textbf{\color{ink}888.}
\par
2048 baytli bir baytli kodlangan matnda nechta belgi bo'lishi mumkin?
\begin{enumerate}[label=\Alph*),leftmargin=9mm]
\item 1024
\item 4096
\item 16384
\item 2048
\end{enumerate}
\begin{tcolorbox}[enhanced,breakable,colback=teal!4,colframe=teal!50!black,boxrule=0.5pt,arc=1.1mm,left=3mm,right=3mm,top=1.7mm,bottom=1.7mm,before skip=2mm,after skip=2mm]
\textbf{\color{teal}To'g'ri javob:} \textbf{D) 2048}
\par\smallskip
\textbf{Izoh:} Har belgi 1 bayt bo'lsa, belgilar soni baytlar soniga teng.
\end{tcolorbox}
\vspace{2mm}\hrule height 0.2pt \vspace{3mm}

% DB question_id: 400c9d7d-fc85-56d8-b07f-a3efba6c2344
\needspace{5\baselineskip}
\noindent\textbf{\color{ink}889.}
\par
58323 ta belgidan iborat matnda har bir belgi 32 bit bilan kodlangan. Matn hajmini baytda toping.
\begin{enumerate}[label=\Alph*),leftmargin=9mm]
\item 233292 bayt
\item 1866336 bayt
\item 466584 bayt
\item 116646 bayt
\end{enumerate}
\begin{tcolorbox}[enhanced,breakable,colback=teal!4,colframe=teal!50!black,boxrule=0.5pt,arc=1.1mm,left=3mm,right=3mm,top=1.7mm,bottom=1.7mm,before skip=2mm,after skip=2mm]
\textbf{\color{teal}To'g'ri javob:} \textbf{A) 233292 bayt}
\par\smallskip
\textbf{Izoh:} 58323 \ensuremath{\times} 32 = 1866336 bit; 8 ga bo'lsak 233292 bayt.
\end{tcolorbox}
\vspace{2mm}\hrule height 0.2pt \vspace{3mm}

% DB question_id: 41ef9a1e-872c-5e21-8848-6bf3845855fb
\needspace{5\baselineskip}
\noindent\textbf{\color{ink}890.}
\par
72245 ta belgidan iborat matnda har bir belgi 8 bit bilan kodlangan. Matn hajmini baytda toping.
\begin{enumerate}[label=\Alph*),leftmargin=9mm]
\item 36122.5 bayt
\item 72245 bayt
\item 577960 bayt
\item 144490 bayt
\end{enumerate}
\begin{tcolorbox}[enhanced,breakable,colback=teal!4,colframe=teal!50!black,boxrule=0.5pt,arc=1.1mm,left=3mm,right=3mm,top=1.7mm,bottom=1.7mm,before skip=2mm,after skip=2mm]
\textbf{\color{teal}To'g'ri javob:} \textbf{B) 72245 bayt}
\par\smallskip
\textbf{Izoh:} 72245 \ensuremath{\times} 8 = 577960 bit; 8 ga bo'lsak 72245 bayt.
\end{tcolorbox}
\vspace{2mm}\hrule height 0.2pt \vspace{3mm}

% DB question_id: 4334de3f-b754-5f4e-b143-1b18fcd41e9d
\needspace{5\baselineskip}
\noindent\textbf{\color{ink}891.}
\par
23338 ta belgidan iborat matnda har bir belgi 16 bit bilan kodlangan. Matn hajmini baytda toping.
\begin{enumerate}[label=\Alph*),leftmargin=9mm]
\item 373408 bayt
\item 93352 bayt
\item 46676 bayt
\item 23338 bayt
\end{enumerate}
\begin{tcolorbox}[enhanced,breakable,colback=teal!4,colframe=teal!50!black,boxrule=0.5pt,arc=1.1mm,left=3mm,right=3mm,top=1.7mm,bottom=1.7mm,before skip=2mm,after skip=2mm]
\textbf{\color{teal}To'g'ri javob:} \textbf{C) 46676 bayt}
\par\smallskip
\textbf{Izoh:} 23338 \ensuremath{\times} 16 = 373408 bit; 8 ga bo'lsak 46676 bayt.
\end{tcolorbox}
\vspace{2mm}\hrule height 0.2pt \vspace{3mm}

% DB question_id: 47902b9d-ec28-586b-b23c-d7d796f312a9
\needspace{5\baselineskip}
\noindent\textbf{\color{ink}892.}
\par
Kitobda \texttt{615\textsubscript{9}} ta sahifa, har sahifada \texttt{203\textsubscript{4}} ta qator va har qatorda \texttt{3C\textsubscript{1}\textsubscript{6}} ta bir baytli belgi bor. Kitobdagi axborot hajmini baytda toping.
\begin{enumerate}[label=\Alph*),leftmargin=9mm]
\item 1 050 000 bayt
\item 1 008 000 bayt
\item 35 000 bayt
\item 3 150 000 bayt
\end{enumerate}
\begin{tcolorbox}[enhanced,breakable,colback=teal!4,colframe=teal!50!black,boxrule=0.5pt,arc=1.1mm,left=3mm,right=3mm,top=1.7mm,bottom=1.7mm,before skip=2mm,after skip=2mm]
\textbf{\color{teal}To'g'ri javob:} \textbf{A) 1 050 000 bayt}
\par\smallskip
\textbf{Izoh:} 615\textsubscript{9}=500, 203\textsubscript{4}=35, 3C\textsubscript{1}\textsubscript{6}=60. Shuning uchun 500\ensuremath{\cdot}35\ensuremath{\cdot}60=1 050 000 bayt.
\end{tcolorbox}
\vspace{2mm}\hrule height 0.2pt \vspace{3mm}

% DB question_id: 5364e149-664a-58ae-93b1-9d3732cf4d06
\needspace{5\baselineskip}
\noindent\textbf{\color{ink}893.}
\par
2852 ta belgidan iborat matnda har bir belgi 16 bit bilan kodlangan. Matn hajmini baytda toping.
\begin{enumerate}[label=\Alph*),leftmargin=9mm]
\item 5704 bayt
\item 45632 bayt
\item 2852 bayt
\item 11408 bayt
\end{enumerate}
\begin{tcolorbox}[enhanced,breakable,colback=teal!4,colframe=teal!50!black,boxrule=0.5pt,arc=1.1mm,left=3mm,right=3mm,top=1.7mm,bottom=1.7mm,before skip=2mm,after skip=2mm]
\textbf{\color{teal}To'g'ri javob:} \textbf{A) 5704 bayt}
\par\smallskip
\textbf{Izoh:} 2852 \ensuremath{\times} 16 = 45632 bit; 8 ga bo'lsak 5704 bayt.
\end{tcolorbox}
\vspace{2mm}\hrule height 0.2pt \vspace{3mm}

% DB question_id: 567f4acc-cf91-463f-8a5c-c9b3255f85fb
\needspace{5\baselineskip}
\noindent\textbf{\color{ink}894.}
\par
300 ta belgi, har biri 2 baytdan kodlansa, matnning nazariy hajmi qancha?
\begin{enumerate}[label=\Alph*),leftmargin=9mm]
\item 300 B
\item 450 B
\item 2400 B
\item 600 B
\end{enumerate}
\begin{tcolorbox}[enhanced,breakable,colback=teal!4,colframe=teal!50!black,boxrule=0.5pt,arc=1.1mm,left=3mm,right=3mm,top=1.7mm,bottom=1.7mm,before skip=2mm,after skip=2mm]
\textbf{\color{teal}To'g'ri javob:} \textbf{D) 600 B}
\par\smallskip
\textbf{Izoh:} 300 \ensuremath{\times} 2 = 600 bayt.
\end{tcolorbox}
\vspace{2mm}\hrule height 0.2pt \vspace{3mm}

% DB question_id: 5768b271-a0a9-5d33-9a8d-2219d181703d
\needspace{5\baselineskip}
\noindent\textbf{\color{ink}895.}
\par
20\ensuremath{\times}15 cm tasvir har ikki yo'nalishda 16 ppi zichlikda olinadi. Har piksel 55 rangdan biri bilan eng kam butun bitlarda kodlanadi. \texttt{1 dyuym = 2,5 cm}. Siqilmagan hajm necha KiB?
\begin{enumerate}[label=\Alph*),leftmargin=9mm]
\item 9 KiB
\item 55 KiB
\item 6 KiB
\item 72 KiB
\end{enumerate}
\begin{tcolorbox}[enhanced,breakable,colback=teal!4,colframe=teal!50!black,boxrule=0.5pt,arc=1.1mm,left=3mm,right=3mm,top=1.7mm,bottom=1.7mm,before skip=2mm,after skip=2mm]
\textbf{\color{teal}To'g'ri javob:} \textbf{A) 9 KiB}
\par\smallskip
\textbf{Izoh:} O'lcham 128\ensuremath{\times}96 piksel; 55 rang uchun 6 bit. 128\ensuremath{\cdot}96\ensuremath{\cdot}6/(8\ensuremath{\cdot}1024)=9 KiB.
\end{tcolorbox}
\vspace{2mm}\hrule height 0.2pt \vspace{3mm}

% DB question_id: 5995b501-f10a-5587-8a7c-8a02f2aa8fa7
\needspace{5\baselineskip}
\noindent\textbf{\color{ink}896.}
\par
195772 ta belgidan iborat matnda har bir belgi 8 bit bilan kodlangan. Matn hajmini baytda toping.
\begin{enumerate}[label=\Alph*),leftmargin=9mm]
\item 195772 bayt
\item 97886 bayt
\item 391544 bayt
\item 1566176 bayt
\end{enumerate}
\begin{tcolorbox}[enhanced,breakable,colback=teal!4,colframe=teal!50!black,boxrule=0.5pt,arc=1.1mm,left=3mm,right=3mm,top=1.7mm,bottom=1.7mm,before skip=2mm,after skip=2mm]
\textbf{\color{teal}To'g'ri javob:} \textbf{A) 195772 bayt}
\par\smallskip
\textbf{Izoh:} 195772 \ensuremath{\times} 8 = 1566176 bit; 8 ga bo'lsak 195772 bayt.
\end{tcolorbox}
\vspace{2mm}\hrule height 0.2pt \vspace{3mm}

% DB question_id: 5b01daf8-ccfc-4c7f-bfc0-f70569afaf26
\needspace{5\baselineskip}
\noindent\textbf{\color{ink}897.}
\par
30 FPS video 2 daqiqa davom etsa, jami nechta kadr bor?
\begin{enumerate}[label=\Alph*),leftmargin=9mm]
\item 1 800
\item 7 200
\item 60
\item 3 600
\end{enumerate}
\begin{tcolorbox}[enhanced,breakable,colback=teal!4,colframe=teal!50!black,boxrule=0.5pt,arc=1.1mm,left=3mm,right=3mm,top=1.7mm,bottom=1.7mm,before skip=2mm,after skip=2mm]
\textbf{\color{teal}To'g'ri javob:} \textbf{D) 3 600}
\par\smallskip
\textbf{Izoh:} 30 \ensuremath{\times} 120 = 3600 kadr.
\end{tcolorbox}
\vspace{2mm}\hrule height 0.2pt \vspace{3mm}

% DB question_id: 5c328b71-ec09-5d6b-a484-9276776a799f
\needspace{5\baselineskip}
\noindent\textbf{\color{ink}898.}
\par
16955 ta belgidan iborat matnda har bir belgi 32 bit bilan kodlangan. Matn hajmini baytda toping.
\begin{enumerate}[label=\Alph*),leftmargin=9mm]
\item 33910 bayt
\item 542560 bayt
\item 135640 bayt
\item 67820 bayt
\end{enumerate}
\begin{tcolorbox}[enhanced,breakable,colback=teal!4,colframe=teal!50!black,boxrule=0.5pt,arc=1.1mm,left=3mm,right=3mm,top=1.7mm,bottom=1.7mm,before skip=2mm,after skip=2mm]
\textbf{\color{teal}To'g'ri javob:} \textbf{D) 67820 bayt}
\par\smallskip
\textbf{Izoh:} 16955 \ensuremath{\times} 32 = 542560 bit; 8 ga bo'lsak 67820 bayt.
\end{tcolorbox}
\vspace{2mm}\hrule height 0.2pt \vspace{3mm}

% DB question_id: 630dcc43-fd17-54fb-b7fa-764c810ebbdd
\needspace{5\baselineskip}
\noindent\textbf{\color{ink}899.}
\par
190793 ta belgidan iborat matnda har bir belgi 32 bit bilan kodlangan. Matn hajmini baytda toping.
\begin{enumerate}[label=\Alph*),leftmargin=9mm]
\item 381586 bayt
\item 1526344 bayt
\item 6105376 bayt
\item 763172 bayt
\end{enumerate}
\begin{tcolorbox}[enhanced,breakable,colback=teal!4,colframe=teal!50!black,boxrule=0.5pt,arc=1.1mm,left=3mm,right=3mm,top=1.7mm,bottom=1.7mm,before skip=2mm,after skip=2mm]
\textbf{\color{teal}To'g'ri javob:} \textbf{D) 763172 bayt}
\par\smallskip
\textbf{Izoh:} 190793 \ensuremath{\times} 32 = 6105376 bit; 8 ga bo'lsak 763172 bayt.
\end{tcolorbox}
\vspace{2mm}\hrule height 0.2pt \vspace{3mm}

% DB question_id: 64d96315-6935-40fe-94d0-505152f4ca80
\needspace{5\baselineskip}
\noindent\textbf{\color{ink}900.}
\par
Ekran o'lchami 256\ensuremath{\times}128 piksel, tasvir 40 Kbayt xotira egallaydi. Har piksel ko'pi bilan nechta rangni aks ettirishi mumkin? (Siqish qo'llanmagan.)
\begin{enumerate}[label=\Alph*),leftmargin=9mm]
\item 768
\item 1024
\item 512
\item 1280
\end{enumerate}
\begin{tcolorbox}[enhanced,breakable,colback=teal!4,colframe=teal!50!black,boxrule=0.5pt,arc=1.1mm,left=3mm,right=3mm,top=1.7mm,bottom=1.7mm,before skip=2mm,after skip=2mm]
\textbf{\color{teal}To'g'ri javob:} \textbf{B) 1024}
\par\smallskip
\textbf{Izoh:} 40\ensuremath{\times}1024\ensuremath{\times}8/(256\ensuremath{\times}128)=10 bit/piksel, 2\textsuperscript{1}\textsuperscript{0}=1024 rang. [2025 kuzgi demo YECHIMLARI, sahifa 5.]
\end{tcolorbox}
\vspace{2mm}\hrule height 0.2pt \vspace{3mm}

% DB question_id: 6c3ea812-c36e-503f-9e57-93fe248a8a6b
\needspace{5\baselineskip}
\noindent\textbf{\color{ink}901.}
\par
25\ensuremath{\times}30 cm tasvir har ikki yo'nalishda 300 dpi aniqlikda skanerlandi. Har nuqta 3 baytda kodlanadi. \texttt{1 dyuym = 2,5 cm}, \texttt{1 MiB = 1 048 576 bayt}. Siqilmagan hajmni eng yaqin butun MiB gacha yaxlitlang.
\begin{enumerate}[label=\Alph*),leftmargin=9mm]
\item 31 MiB
\item 30 MiB
\item 32 MiB
\item 24 MiB
\end{enumerate}
\begin{tcolorbox}[enhanced,breakable,colback=teal!4,colframe=teal!50!black,boxrule=0.5pt,arc=1.1mm,left=3mm,right=3mm,top=1.7mm,bottom=1.7mm,before skip=2mm,after skip=2mm]
\textbf{\color{teal}To'g'ri javob:} \textbf{A) 31 MiB}
\par\smallskip
\textbf{Izoh:} O'lcham 3000\ensuremath{\times}3600 nuqta. Hajm 3000\ensuremath{\cdot}3600\ensuremath{\cdot}3=32 400 000 bayt\ensuremath{\approx}30,899 MiB; eng yaqin butun qiymat 31 MiB.
\end{tcolorbox}
\vspace{2mm}\hrule height 0.2pt \vspace{3mm}

% DB question_id: 789bbe66-b467-5bf5-9861-69db75107c5c
\needspace{5\baselineskip}
\noindent\textbf{\color{ink}902.}
\par
80609 ta belgidan iborat matnda har bir belgi 8 bit bilan kodlangan. Matn hajmini baytda toping.
\begin{enumerate}[label=\Alph*),leftmargin=9mm]
\item 80609 bayt
\item 644872 bayt
\item 161218 bayt
\item 40304.5 bayt
\end{enumerate}
\begin{tcolorbox}[enhanced,breakable,colback=teal!4,colframe=teal!50!black,boxrule=0.5pt,arc=1.1mm,left=3mm,right=3mm,top=1.7mm,bottom=1.7mm,before skip=2mm,after skip=2mm]
\textbf{\color{teal}To'g'ri javob:} \textbf{A) 80609 bayt}
\par\smallskip
\textbf{Izoh:} 80609 \ensuremath{\times} 8 = 644872 bit; 8 ga bo'lsak 80609 bayt.
\end{tcolorbox}
\vspace{2mm}\hrule height 0.2pt \vspace{3mm}

% DB question_id: 7bf5cef4-a085-588f-b30c-a7fb3880a3f4
\needspace{5\baselineskip}
\noindent\textbf{\color{ink}903.}
\par
113613 ta belgidan iborat matnda har bir belgi 8 bit bilan kodlangan. Matn hajmini baytda toping.
\begin{enumerate}[label=\Alph*),leftmargin=9mm]
\item 227226 bayt
\item 56806.5 bayt
\item 113613 bayt
\item 908904 bayt
\end{enumerate}
\begin{tcolorbox}[enhanced,breakable,colback=teal!4,colframe=teal!50!black,boxrule=0.5pt,arc=1.1mm,left=3mm,right=3mm,top=1.7mm,bottom=1.7mm,before skip=2mm,after skip=2mm]
\textbf{\color{teal}To'g'ri javob:} \textbf{C) 113613 bayt}
\par\smallskip
\textbf{Izoh:} 113613 \ensuremath{\times} 8 = 908904 bit; 8 ga bo'lsak 113613 bayt.
\end{tcolorbox}
\vspace{2mm}\hrule height 0.2pt \vspace{3mm}

% DB question_id: 8365d280-c304-40a3-92d0-0ac86646b1b5
\needspace{5\baselineskip}
\noindent\textbf{\color{ink}904.}
\par
128 betli kitobning har betida 48 qator, har qatorda 36 belgi bor. Har belgi 1 bayt bo'lsa, kitob hajmi necha Kbayt?
\begin{enumerate}[label=\Alph*),leftmargin=9mm]
\item 256
\item 196
\item 228
\item 216
\end{enumerate}
\begin{tcolorbox}[enhanced,breakable,colback=teal!4,colframe=teal!50!black,boxrule=0.5pt,arc=1.1mm,left=3mm,right=3mm,top=1.7mm,bottom=1.7mm,before skip=2mm,after skip=2mm]
\textbf{\color{teal}To'g'ri javob:} \textbf{D) 216}
\par\smallskip
\textbf{Izoh:} 128\ensuremath{\times}48\ensuremath{\times}36/1024=216 Kbayt.
\end{tcolorbox}
\vspace{2mm}\hrule height 0.2pt \vspace{3mm}

% DB question_id: 83aeecea-8d84-4845-962b-aeaa9ee8a8ee
\needspace{5\baselineskip}
\noindent\textbf{\color{ink}905.}
\par
25 FPS va 20 soniya uchun 500 kadr hisoblandi. Vaziyatni tahlil qilib, unga eng aniq mos tushunchani toping.
\begin{enumerate}[label=\Alph*),leftmargin=9mm]
\item Teskari masala
\item Mantiqiy tekshiruv
\item Formula tanlash
\item Kadrlar soni
\end{enumerate}
\begin{tcolorbox}[enhanced,breakable,colback=teal!4,colframe=teal!50!black,boxrule=0.5pt,arc=1.1mm,left=3mm,right=3mm,top=1.7mm,bottom=1.7mm,before skip=2mm,after skip=2mm]
\textbf{\color{teal}To'g'ri javob:} \textbf{D) Kadrlar soni}
\par\smallskip
\textbf{Izoh:} Vaziyat Kadrlar soniga misol: FPS va davomiylik ko'paytmasi orqali topiladigan video tasvirlari soni.
\end{tcolorbox}
\vspace{2mm}\hrule height 0.2pt \vspace{3mm}

% DB question_id: 843bc08e-b811-5503-b921-08fd50acc1a4
\needspace{5\baselineskip}
\noindent\textbf{\color{ink}906.}
\par
78627 ta belgidan iborat matnda har bir belgi 32 bit bilan kodlangan. Matn hajmini baytda toping.
\begin{enumerate}[label=\Alph*),leftmargin=9mm]
\item 314508 bayt
\item 157254 bayt
\item 629016 bayt
\item 2516064 bayt
\end{enumerate}
\begin{tcolorbox}[enhanced,breakable,colback=teal!4,colframe=teal!50!black,boxrule=0.5pt,arc=1.1mm,left=3mm,right=3mm,top=1.7mm,bottom=1.7mm,before skip=2mm,after skip=2mm]
\textbf{\color{teal}To'g'ri javob:} \textbf{A) 314508 bayt}
\par\smallskip
\textbf{Izoh:} 78627 \ensuremath{\times} 32 = 2516064 bit; 8 ga bo'lsak 314508 bayt.
\end{tcolorbox}
\vspace{2mm}\hrule height 0.2pt \vspace{3mm}

% DB question_id: 844e10f2-b8e8-5ab3-ba44-b1aad473541b
\needspace{5\baselineskip}
\noindent\textbf{\color{ink}907.}
\par
39241 ta belgidan iborat matnda har bir belgi 8 bit bilan kodlangan. Matn hajmini baytda toping.
\begin{enumerate}[label=\Alph*),leftmargin=9mm]
\item 78482 bayt
\item 19620.5 bayt
\item 313928 bayt
\item 39241 bayt
\end{enumerate}
\begin{tcolorbox}[enhanced,breakable,colback=teal!4,colframe=teal!50!black,boxrule=0.5pt,arc=1.1mm,left=3mm,right=3mm,top=1.7mm,bottom=1.7mm,before skip=2mm,after skip=2mm]
\textbf{\color{teal}To'g'ri javob:} \textbf{D) 39241 bayt}
\par\smallskip
\textbf{Izoh:} 39241 \ensuremath{\times} 8 = 313928 bit; 8 ga bo'lsak 39241 bayt.
\end{tcolorbox}
\vspace{2mm}\hrule height 0.2pt \vspace{3mm}

% DB question_id: 86aaf2e2-045b-54b4-843f-1bca9606048d
\needspace{5\baselineskip}
\noindent\textbf{\color{ink}908.}
\par
73649 ta belgidan iborat matnda har bir belgi 8 bit bilan kodlangan. Matn hajmini baytda toping.
\begin{enumerate}[label=\Alph*),leftmargin=9mm]
\item 36824.5 bayt
\item 147298 bayt
\item 589192 bayt
\item 73649 bayt
\end{enumerate}
\begin{tcolorbox}[enhanced,breakable,colback=teal!4,colframe=teal!50!black,boxrule=0.5pt,arc=1.1mm,left=3mm,right=3mm,top=1.7mm,bottom=1.7mm,before skip=2mm,after skip=2mm]
\textbf{\color{teal}To'g'ri javob:} \textbf{D) 73649 bayt}
\par\smallskip
\textbf{Izoh:} 73649 \ensuremath{\times} 8 = 589192 bit; 8 ga bo'lsak 73649 bayt.
\end{tcolorbox}
\vspace{2mm}\hrule height 0.2pt \vspace{3mm}

% DB question_id: 8b676960-96bf-54c9-a2e0-921767074a67
\needspace{5\baselineskip}
\noindent\textbf{\color{ink}909.}
\par
108056 ta belgidan iborat matnda har bir belgi 8 bit bilan kodlangan. Matn hajmini baytda toping.
\begin{enumerate}[label=\Alph*),leftmargin=9mm]
\item 54028 bayt
\item 108056 bayt
\item 864448 bayt
\item 216112 bayt
\end{enumerate}
\begin{tcolorbox}[enhanced,breakable,colback=teal!4,colframe=teal!50!black,boxrule=0.5pt,arc=1.1mm,left=3mm,right=3mm,top=1.7mm,bottom=1.7mm,before skip=2mm,after skip=2mm]
\textbf{\color{teal}To'g'ri javob:} \textbf{B) 108056 bayt}
\par\smallskip
\textbf{Izoh:} 108056 \ensuremath{\times} 8 = 864448 bit; 8 ga bo'lsak 108056 bayt.
\end{tcolorbox}
\vspace{2mm}\hrule height 0.2pt \vspace{3mm}

% DB question_id: 8fdf47b9-a54b-5c9b-8c90-e7afa3b864c4
\needspace{5\baselineskip}
\noindent\textbf{\color{ink}910.}
\par
Harflar \texttt{A=00}, \texttt{B=11}, \texttt{C=10}, \texttt{D=010}, \texttt{E=0110} kodlari bilan yoziladi. Quyidagi bit satrlaridan faqat bittasi boshidan oxirigacha berilgan kodlarga ajraladi. Uni toping.
\begin{enumerate}[label=\Alph*),leftmargin=9mm]
\item 001101011001110
\item 000110100010110
\item 011010110011100
\item 011001011110100
\end{enumerate}
\begin{tcolorbox}[enhanced,breakable,colback=teal!4,colframe=teal!50!black,boxrule=0.5pt,arc=1.1mm,left=3mm,right=3mm,top=1.7mm,bottom=1.7mm,before skip=2mm,after skip=2mm]
\textbf{\color{teal}To'g'ri javob:} \textbf{A) 001101011001110}
\par\smallskip
\textbf{Izoh:} Birinchi satr 00|11|010|11|00|11|10 ko'rinishida A B D B A B C ga to'liq ajraladi. Qolgan satrlarda biror joyda mos kodsiz qoldiq qoladi.
\end{tcolorbox}
\vspace{2mm}\hrule height 0.2pt \vspace{3mm}

% DB question_id: 941a7395-c467-51b9-a5c3-50c1eafc4274
\needspace{5\baselineskip}
\noindent\textbf{\color{ink}911.}
\par
Har bir identifikator 158 ta belgidan iborat. Alifboda 10 ta raqam va 2022 ta maxsus belgi bor. Har bir belgi bir xil va eng kam bit bilan, har bir identifikator esa mumkin bo'lgan eng kichik butun baytlar sonida saqlanadi. 15360 ta identifikator uchun qancha KB xotira kerak? \texttt{1 KB = 1024 bayt}.
\begin{enumerate}[label=\Alph*),leftmargin=9mm]
\item 3270
\item 3255
\item 3258
\item 3259
\end{enumerate}
\begin{tcolorbox}[enhanced,breakable,colback=teal!4,colframe=teal!50!black,boxrule=0.5pt,arc=1.1mm,left=3mm,right=3mm,top=1.7mm,bottom=1.7mm,before skip=2mm,after skip=2mm]
\textbf{\color{teal}To'g'ri javob:} \textbf{A) 3270}
\par\smallskip
\textbf{Izoh:} Alifbo hajmi \texttt{2032}, shuning uchun bir belgi uchun \texttt{ceil(log\textsubscript{2}2032)=11} bit. Bir identifikator \texttt{158\ensuremath{\cdot}11=1738} bit, ya'ni \texttt{ceil(1738/8)=218} bayt. Jami \texttt{15360\ensuremath{\cdot}218/1024=3270 KB}.
\end{tcolorbox}
\vspace{2mm}\hrule height 0.2pt \vspace{3mm}

% DB question_id: 9769b009-1ec0-451c-83a1-b6e06bd0c179
\needspace{5\baselineskip}
\noindent\textbf{\color{ink}912.}
\par
128 \ensuremath{\times} 128 piksel, 32 bitli rasm hajmi qancha?
\begin{enumerate}[label=\Alph*),leftmargin=9mm]
\item 16 384 B
\item 65 536 B
\item 524 288 B
\item 4 096 B
\end{enumerate}
\begin{tcolorbox}[enhanced,breakable,colback=teal!4,colframe=teal!50!black,boxrule=0.5pt,arc=1.1mm,left=3mm,right=3mm,top=1.7mm,bottom=1.7mm,before skip=2mm,after skip=2mm]
\textbf{\color{teal}To'g'ri javob:} \textbf{B) 65 536 B}
\par\smallskip
\textbf{Izoh:} 128 \ensuremath{\times} 128 \ensuremath{\times} 32 \ensuremath{\div} 8 = 65 536 bayt.
\end{tcolorbox}
\vspace{2mm}\hrule height 0.2pt \vspace{3mm}

% DB question_id: 9803392b-c484-4905-b1c9-ede4474b1611
\needspace{5\baselineskip}
\noindent\textbf{\color{ink}913.}
\par
16 bit rang chuqurligi nazariy jihatdan nechta rangni ifodalaydi?
\begin{enumerate}[label=\Alph*),leftmargin=9mm]
\item 16
\item 65 536
\item 256
\item 32 768
\end{enumerate}
\begin{tcolorbox}[enhanced,breakable,colback=teal!4,colframe=teal!50!black,boxrule=0.5pt,arc=1.1mm,left=3mm,right=3mm,top=1.7mm,bottom=1.7mm,before skip=2mm,after skip=2mm]
\textbf{\color{teal}To'g'ri javob:} \textbf{B) 65 536}
\par\smallskip
\textbf{Izoh:} Ranglar soni 2\textsuperscript{1}\textsuperscript{6} = 65 536.
\end{tcolorbox}
\vspace{2mm}\hrule height 0.2pt \vspace{3mm}

% DB question_id: a14b8369-0d23-494a-8500-836ceef16c8d
\needspace{5\baselineskip}
\noindent\textbf{\color{ink}914.}
\par
1200 ta belgi 8 bitli kodda saqlansa, hajm qancha bo'ladi?
\begin{enumerate}[label=\Alph*),leftmargin=9mm]
\item 1200 B
\item 150 B
\item 2400 B
\item 9600 B
\end{enumerate}
\begin{tcolorbox}[enhanced,breakable,colback=teal!4,colframe=teal!50!black,boxrule=0.5pt,arc=1.1mm,left=3mm,right=3mm,top=1.7mm,bottom=1.7mm,before skip=2mm,after skip=2mm]
\textbf{\color{teal}To'g'ri javob:} \textbf{A) 1200 B}
\par\smallskip
\textbf{Izoh:} 8 bit = 1 bayt; 1200 belgi = 1200 B.
\end{tcolorbox}
\vspace{2mm}\hrule height 0.2pt \vspace{3mm}

% DB question_id: a5ee576e-4795-50e1-bb20-8ee03892060f
\needspace{5\baselineskip}
\noindent\textbf{\color{ink}915.}
\par
106652 ta belgidan iborat matnda har bir belgi 8 bit bilan kodlangan. Matn hajmini baytda toping.
\begin{enumerate}[label=\Alph*),leftmargin=9mm]
\item 53326 bayt
\item 106652 bayt
\item 853216 bayt
\item 213304 bayt
\end{enumerate}
\begin{tcolorbox}[enhanced,breakable,colback=teal!4,colframe=teal!50!black,boxrule=0.5pt,arc=1.1mm,left=3mm,right=3mm,top=1.7mm,bottom=1.7mm,before skip=2mm,after skip=2mm]
\textbf{\color{teal}To'g'ri javob:} \textbf{B) 106652 bayt}
\par\smallskip
\textbf{Izoh:} 106652 \ensuremath{\times} 8 = 853216 bit; 8 ga bo'lsak 106652 bayt.
\end{tcolorbox}
\vspace{2mm}\hrule height 0.2pt \vspace{3mm}

% DB question_id: a748bc67-7566-412f-83b3-d785ed294e6c
\needspace{5\baselineskip}
\noindent\textbf{\color{ink}916.}
\par
Kitob hajmi 1024 KiB. Har sahifada 64 qator, har qatorda 32 belgi bor. Har belgi 1 bayt, boshqa xizmat baytlari yo'q. Kitob necha sahifadan iborat?
\begin{enumerate}[label=\Alph*),leftmargin=9mm]
\item 128
\item 256
\item 512
\item 1024
\end{enumerate}
\begin{tcolorbox}[enhanced,breakable,colback=teal!4,colframe=teal!50!black,boxrule=0.5pt,arc=1.1mm,left=3mm,right=3mm,top=1.7mm,bottom=1.7mm,before skip=2mm,after skip=2mm]
\textbf{\color{teal}To'g'ri javob:} \textbf{C) 512}
\par\smallskip
\textbf{Izoh:} Bir sahifada 64\ensuremath{\times}32=2048 belgi=2048 bayt. Kitob 1024\ensuremath{\times}1024=1048576 bayt. 1048576/2048=512 sahifa.
\end{tcolorbox}
\vspace{2mm}\hrule height 0.2pt \vspace{3mm}

% DB question_id: aa5dcce0-1952-584c-b1ba-38d70e35fb23
\needspace{5\baselineskip}
\noindent\textbf{\color{ink}917.}
\par
156385 ta belgidan iborat matnda har bir belgi 16 bit bilan kodlangan. Matn hajmini baytda toping.
\begin{enumerate}[label=\Alph*),leftmargin=9mm]
\item 625540 bayt
\item 156385 bayt
\item 2502160 bayt
\item 312770 bayt
\end{enumerate}
\begin{tcolorbox}[enhanced,breakable,colback=teal!4,colframe=teal!50!black,boxrule=0.5pt,arc=1.1mm,left=3mm,right=3mm,top=1.7mm,bottom=1.7mm,before skip=2mm,after skip=2mm]
\textbf{\color{teal}To'g'ri javob:} \textbf{D) 312770 bayt}
\par\smallskip
\textbf{Izoh:} 156385 \ensuremath{\times} 16 = 2502160 bit; 8 ga bo'lsak 312770 bayt.
\end{tcolorbox}
\vspace{2mm}\hrule height 0.2pt \vspace{3mm}

% DB question_id: abb4a6ff-6d93-50a1-a0fc-7528d495bc29
\needspace{5\baselineskip}
\noindent\textbf{\color{ink}918.}
\par
30876 ta belgidan iborat matnda har bir belgi 32 bit bilan kodlangan. Matn hajmini baytda toping.
\begin{enumerate}[label=\Alph*),leftmargin=9mm]
\item 247008 bayt
\item 61752 bayt
\item 123504 bayt
\item 988032 bayt
\end{enumerate}
\begin{tcolorbox}[enhanced,breakable,colback=teal!4,colframe=teal!50!black,boxrule=0.5pt,arc=1.1mm,left=3mm,right=3mm,top=1.7mm,bottom=1.7mm,before skip=2mm,after skip=2mm]
\textbf{\color{teal}To'g'ri javob:} \textbf{C) 123504 bayt}
\par\smallskip
\textbf{Izoh:} 30876 \ensuremath{\times} 32 = 988032 bit; 8 ga bo'lsak 123504 bayt.
\end{tcolorbox}
\vspace{2mm}\hrule height 0.2pt \vspace{3mm}

% DB question_id: b22c75e0-1561-5797-821e-00dc137f7262
\needspace{5\baselineskip}
\noindent\textbf{\color{ink}919.}
\par
Har bir foydalanuvchi identifikatori 101 ta belgidan iborat. Har bir belgi 10 ta o'nlik raqamdan yoki 4090 ta maxsus belgidan biri bo'lishi mumkin. Barcha belgilar bir xil eng kam bitlar soni bilan, har bir identifikator esa mumkin bo'lgan eng kichik butun baytlar sonida saqlanadi. 2048 ta identifikator uchun qancha KiB xotira kerak? \texttt{1 KiB = 1024 bayt}.
\begin{enumerate}[label=\Alph*),leftmargin=9mm]
\item 306 KiB
\item 302 KiB
\item 304 KiB
\item 330 KiB
\end{enumerate}
\begin{tcolorbox}[enhanced,breakable,colback=teal!4,colframe=teal!50!black,boxrule=0.5pt,arc=1.1mm,left=3mm,right=3mm,top=1.7mm,bottom=1.7mm,before skip=2mm,after skip=2mm]
\textbf{\color{teal}To'g'ri javob:} \textbf{D) 330 KiB}
\par\smallskip
\textbf{Izoh:} Alifbo hajmi 10+4090=4100; 2\textasciicircum{}12=4096\textless{}4100\ensuremath{\leq}8192=2\textasciicircum{}13, shuning uchun har bir belgi 13 bit. Bitta identifikator 101\ensuremath{\cdot}13=1313 bit, ya'ni \ensuremath{\lceil}1313/8\ensuremath{\rceil}=165 bayt. 2048\ensuremath{\cdot}165/1024=330 KiB. Asl tahlil 10 ta raqamni 4090 ga qo'shmagan va 304 ni belgilagan; variantlarda bo'lmagan 330 KiB to'g'ri javobi qo'shildi.
\end{tcolorbox}
\vspace{2mm}\hrule height 0.2pt \vspace{3mm}

% DB question_id: b6d65fd2-7bb8-5b1b-ba25-2e91cb335a9f
\needspace{5\baselineskip}
\noindent\textbf{\color{ink}920.}
\par
148021 ta belgidan iborat matnda har bir belgi 16 bit bilan kodlangan. Matn hajmini baytda toping.
\begin{enumerate}[label=\Alph*),leftmargin=9mm]
\item 592084 bayt
\item 148021 bayt
\item 296042 bayt
\item 2368336 bayt
\end{enumerate}
\begin{tcolorbox}[enhanced,breakable,colback=teal!4,colframe=teal!50!black,boxrule=0.5pt,arc=1.1mm,left=3mm,right=3mm,top=1.7mm,bottom=1.7mm,before skip=2mm,after skip=2mm]
\textbf{\color{teal}To'g'ri javob:} \textbf{C) 296042 bayt}
\par\smallskip
\textbf{Izoh:} 148021 \ensuremath{\times} 16 = 2368336 bit; 8 ga bo'lsak 296042 bayt.
\end{tcolorbox}
\vspace{2mm}\hrule height 0.2pt \vspace{3mm}

% DB question_id: b97c44a4-542b-461e-88dd-7a2b1431afb0
\needspace{5\baselineskip}
\noindent\textbf{\color{ink}921.}
\par
A, E, I, S, T, Y harflaridan tuzilgan barcha 5 belgili parollarda harflar takrorlanishi mumkin. Parollar alifbo tartibida 1 dan boshlab raqamlangan. T harfi bilan boshlanuvchi birinchi parol qaysi o'rinda turadi?
\begin{enumerate}[label=\Alph*),leftmargin=9mm]
\item 5184
\item 1296
\item 5185
\item 1297
\end{enumerate}
\begin{tcolorbox}[enhanced,breakable,colback=teal!4,colframe=teal!50!black,boxrule=0.5pt,arc=1.1mm,left=3mm,right=3mm,top=1.7mm,bottom=1.7mm,before skip=2mm,after skip=2mm]
\textbf{\color{teal}To'g'ri javob:} \textbf{C) 5185}
\par\smallskip
\textbf{Izoh:} T dan oldin A, E, I, S harflari turadi. Har biri bilan boshlanadigan parollar soni 6\textsuperscript{4}, chunki har bir o'ringa 6 xil harf tanlanadi va takrorlash mumkin. 4 \ensuremath{\times} 6\textsuperscript{4} = 5184; birinchi TAAAA paroli 5185-o'rinda.
\end{tcolorbox}
\vspace{2mm}\hrule height 0.2pt \vspace{3mm}

% DB question_id: bfb2102f-61f4-42ff-b67c-85b7af71595b
\needspace{5\baselineskip}
\noindent\textbf{\color{ink}922.}
\par
8 Mbit/s videoning 5 daqiqalik hajmi taxminan qancha? MB = 10\textsuperscript{6} B.
\begin{enumerate}[label=\Alph*),leftmargin=9mm]
\item 30 MB
\item 150 MB
\item 300 MB
\item 2 400 MB
\end{enumerate}
\begin{tcolorbox}[enhanced,breakable,colback=teal!4,colframe=teal!50!black,boxrule=0.5pt,arc=1.1mm,left=3mm,right=3mm,top=1.7mm,bottom=1.7mm,before skip=2mm,after skip=2mm]
\textbf{\color{teal}To'g'ri javob:} \textbf{C) 300 MB}
\par\smallskip
\textbf{Izoh:} 8 000 000 \ensuremath{\times} 300 \ensuremath{\div} 8 = 300 000 000 B = 300 MB.
\end{tcolorbox}
\vspace{2mm}\hrule height 0.2pt \vspace{3mm}

% DB question_id: c2210821-15fa-48c1-8c91-ec88c3eea075
\needspace{5\baselineskip}
\noindent\textbf{\color{ink}923.}
\par
256\ensuremath{\times}200 pikselli siqilmagan tasvir 56,25 KiB xotira egallaydi. Har bir piksel ko'pi bilan nechta rangni ifodalay oladi?
\begin{enumerate}[label=\Alph*),leftmargin=9mm]
\item 1024 ta
\item 1280 ta
\item 768 ta
\item 512 ta
\end{enumerate}
\begin{tcolorbox}[enhanced,breakable,colback=teal!4,colframe=teal!50!black,boxrule=0.5pt,arc=1.1mm,left=3mm,right=3mm,top=1.7mm,bottom=1.7mm,before skip=2mm,after skip=2mm]
\textbf{\color{teal}To'g'ri javob:} \textbf{D) 512 ta}
\par\smallskip
\textbf{Izoh:} 56,25\ensuremath{\cdot}1024\ensuremath{\cdot}8/(256\ensuremath{\cdot}200)=9 bit/piksel. 2\textasciicircum{}9=512 rang. [Yutuqli\_test, test kodi 12; Telegram xabarlar 6446 va 6556.]
\end{tcolorbox}
\vspace{2mm}\hrule height 0.2pt \vspace{3mm}

% DB question_id: c9e649b5-65e5-5ad3-8739-6ad979fdcd4a
\needspace{5\baselineskip}
\noindent\textbf{\color{ink}924.}
\par
10\ensuremath{\times}15 cm tasvir 300 ppi zichlikda olinib, 256 rang bilan kodlandi. \texttt{1 dyuym = 2,5 cm}. Siqilmagan hajmni baytda toping.
\begin{enumerate}[label=\Alph*),leftmargin=9mm]
\item 2 160 000 bayt
\item 17 280 000 bayt
\item 270 000 bayt
\item 21 600 000 bayt
\end{enumerate}
\begin{tcolorbox}[enhanced,breakable,colback=teal!4,colframe=teal!50!black,boxrule=0.5pt,arc=1.1mm,left=3mm,right=3mm,top=1.7mm,bottom=1.7mm,before skip=2mm,after skip=2mm]
\textbf{\color{teal}To'g'ri javob:} \textbf{A) 2 160 000 bayt}
\par\smallskip
\textbf{Izoh:} O'lchamlar 1200\ensuremath{\times}1800 piksel; 256 rang=8 bit=1 bayt/piksel. Hajm 1200\ensuremath{\cdot}1800=2 160 000 bayt.
\end{tcolorbox}
\vspace{2mm}\hrule height 0.2pt \vspace{3mm}

% DB question_id: ce983991-2657-490e-94fd-6212e8098567
\needspace{5\baselineskip}
\noindent\textbf{\color{ink}925.}
\par
Stereo audio formulasida kanal soni 2 ga ko'paytirildi. Vaziyatni tahlil qilib, unga eng aniq mos tushunchani toping.
\begin{enumerate}[label=\Alph*),leftmargin=9mm]
\item Kadrlar soni
\item Teskari masala
\item Ko'p kanalli hisob
\item Mantiqiy tekshiruv
\end{enumerate}
\begin{tcolorbox}[enhanced,breakable,colback=teal!4,colframe=teal!50!black,boxrule=0.5pt,arc=1.1mm,left=3mm,right=3mm,top=1.7mm,bottom=1.7mm,before skip=2mm,after skip=2mm]
\textbf{\color{teal}To'g'ri javob:} \textbf{C) Ko'p kanalli hisob}
\par\smallskip
\textbf{Izoh:} Vaziyat Ko'p kanalli hisobga misol: har bir mustaqil kanal hajmini umumiy formulaga kiritish.
\end{tcolorbox}
\vspace{2mm}\hrule height 0.2pt \vspace{3mm}

% DB question_id: cede8aa5-87e6-5e0e-9d0f-8d0c429ca53b
\needspace{5\baselineskip}
\noindent\textbf{\color{ink}926.}
\par
119996 ta belgidan iborat matnda har bir belgi 32 bit bilan kodlangan. Matn hajmini baytda toping.
\begin{enumerate}[label=\Alph*),leftmargin=9mm]
\item 479984 bayt
\item 3839872 bayt
\item 239992 bayt
\item 959968 bayt
\end{enumerate}
\begin{tcolorbox}[enhanced,breakable,colback=teal!4,colframe=teal!50!black,boxrule=0.5pt,arc=1.1mm,left=3mm,right=3mm,top=1.7mm,bottom=1.7mm,before skip=2mm,after skip=2mm]
\textbf{\color{teal}To'g'ri javob:} \textbf{A) 479984 bayt}
\par\smallskip
\textbf{Izoh:} 119996 \ensuremath{\times} 32 = 3839872 bit; 8 ga bo'lsak 479984 bayt.
\end{tcolorbox}
\vspace{2mm}\hrule height 0.2pt \vspace{3mm}

% DB question_id: d07183f9-bfc8-4cec-b509-ef6ef61366be
\needspace{5\baselineskip}
\noindent\textbf{\color{ink}927.}
\par
800 ta belgi 16 bitli kodda saqlansa, xom hajm qancha?
\begin{enumerate}[label=\Alph*),leftmargin=9mm]
\item 1 600 B
\item 800 B
\item 6 400 B
\item 12 800 B
\end{enumerate}
\begin{tcolorbox}[enhanced,breakable,colback=teal!4,colframe=teal!50!black,boxrule=0.5pt,arc=1.1mm,left=3mm,right=3mm,top=1.7mm,bottom=1.7mm,before skip=2mm,after skip=2mm]
\textbf{\color{teal}To'g'ri javob:} \textbf{A) 1 600 B}
\par\smallskip
\textbf{Izoh:} 800 \ensuremath{\times} 16 \ensuremath{\div} 8 = 1600 bayt.
\end{tcolorbox}
\vspace{2mm}\hrule height 0.2pt \vspace{3mm}

% DB question_id: d0cf9413-3101-547d-b210-964b08b84ca9
\needspace{5\baselineskip}
\noindent\textbf{\color{ink}928.}
\par
149425 ta belgidan iborat matnda har bir belgi 16 bit bilan kodlangan. Matn hajmini baytda toping.
\begin{enumerate}[label=\Alph*),leftmargin=9mm]
\item 2390800 bayt
\item 298850 bayt
\item 597700 bayt
\item 149425 bayt
\end{enumerate}
\begin{tcolorbox}[enhanced,breakable,colback=teal!4,colframe=teal!50!black,boxrule=0.5pt,arc=1.1mm,left=3mm,right=3mm,top=1.7mm,bottom=1.7mm,before skip=2mm,after skip=2mm]
\textbf{\color{teal}To'g'ri javob:} \textbf{B) 298850 bayt}
\par\smallskip
\textbf{Izoh:} 149425 \ensuremath{\times} 16 = 2390800 bit; 8 ga bo'lsak 298850 bayt.
\end{tcolorbox}
\vspace{2mm}\hrule height 0.2pt \vspace{3mm}

% DB question_id: d12b0576-b255-4153-8c81-92019daab0f8
\needspace{5\baselineskip}
\noindent\textbf{\color{ink}929.}
\par
64 xil teng ehtimolli belgini kodlash uchun har bir belgiga necha bit kerak?
\begin{enumerate}[label=\Alph*),leftmargin=9mm]
\item 5
\item 7
\item 6
\item 8
\end{enumerate}
\begin{tcolorbox}[enhanced,breakable,colback=teal!4,colframe=teal!50!black,boxrule=0.5pt,arc=1.1mm,left=3mm,right=3mm,top=1.7mm,bottom=1.7mm,before skip=2mm,after skip=2mm]
\textbf{\color{teal}To'g'ri javob:} \textbf{C) 6}
\par\smallskip
\textbf{Izoh:} log\textsubscript{2}64 = 6.
\end{tcolorbox}
\vspace{2mm}\hrule height 0.2pt \vspace{3mm}

% DB question_id: d14673e4-239b-4940-807e-a641a0267b1d
\needspace{5\baselineskip}
\noindent\textbf{\color{ink}930.}
\par
Namuna chastotasi ikki marta oshib, boshqa omillar o'zgarmasa, xom audio hajmi qanday o'zgaradi?
\begin{enumerate}[label=\Alph*),leftmargin=9mm]
\item yarim bo'ladi
\item o'zgarmaydi
\item ikki marta oshadi
\item to'rt marta oshadi
\end{enumerate}
\begin{tcolorbox}[enhanced,breakable,colback=teal!4,colframe=teal!50!black,boxrule=0.5pt,arc=1.1mm,left=3mm,right=3mm,top=1.7mm,bottom=1.7mm,before skip=2mm,after skip=2mm]
\textbf{\color{teal}To'g'ri javob:} \textbf{C) ikki marta oshadi}
\par\smallskip
\textbf{Izoh:} Hajm namuna chastotasiga to'g'ri proporsional.
\end{tcolorbox}
\vspace{2mm}\hrule height 0.2pt \vspace{3mm}

% DB question_id: d2d6f94c-5a1a-57b3-ab20-c70ea4403841
\needspace{5\baselineskip}
\noindent\textbf{\color{ink}931.}
\par
Stereo audio 16 kHz namuna olish chastotasi va har kanal uchun 64 bitli kvantlashda yozildi. Siqilmagan fayl hajmi 20 MiB. Yozuv davomiyligini eng yaqin butun minutgacha yaxlitlang.
\begin{enumerate}[label=\Alph*),leftmargin=9mm]
\item 1 minut
\item 2 minut
\item 82 minut
\item 20 minut
\end{enumerate}
\begin{tcolorbox}[enhanced,breakable,colback=teal!4,colframe=teal!50!black,boxrule=0.5pt,arc=1.1mm,left=3mm,right=3mm,top=1.7mm,bottom=1.7mm,before skip=2mm,after skip=2mm]
\textbf{\color{teal}To'g'ri javob:} \textbf{A) 1 minut}
\par\smallskip
\textbf{Izoh:} Tezlik 2\ensuremath{\cdot}16 000\ensuremath{\cdot}64=2 048 000 bit/s=256 000 bayt/s. 20\ensuremath{\cdot}1 048 576/256 000=81,92 s\ensuremath{\approx}1,365 minut; eng yaqin butun qiymat 1 minut.
\end{tcolorbox}
\vspace{2mm}\hrule height 0.2pt \vspace{3mm}

% DB question_id: d367a949-125a-5f49-9c6f-4af21e71de5a
\needspace{5\baselineskip}
\noindent\textbf{\color{ink}932.}
\par
10\ensuremath{\times}15 cm tasvir har ikki yo'nalishda 32 ppi zichlikda olinadi. Har piksel 16 rangdan biri bilan kodlanadi. \texttt{1 dyuym = 2,5 cm} bo'lsa, siqilmagan hajm necha KiB?
\begin{enumerate}[label=\Alph*),leftmargin=9mm]
\item 12 KiB
\item 96 KiB
\item 48 KiB
\item 24 KiB
\end{enumerate}
\begin{tcolorbox}[enhanced,breakable,colback=teal!4,colframe=teal!50!black,boxrule=0.5pt,arc=1.1mm,left=3mm,right=3mm,top=1.7mm,bottom=1.7mm,before skip=2mm,after skip=2mm]
\textbf{\color{teal}To'g'ri javob:} \textbf{A) 12 KiB}
\par\smallskip
\textbf{Izoh:} O'lcham 128\ensuremath{\times}192 piksel, rang chuqurligi 4 bit. 128\ensuremath{\cdot}192\ensuremath{\cdot}4/(8\ensuremath{\cdot}1024)=12 KiB.
\end{tcolorbox}
\vspace{2mm}\hrule height 0.2pt \vspace{3mm}

% DB question_id: d8b76803-98da-5c16-a98a-000452b9cf07
\needspace{5\baselineskip}
\noindent\textbf{\color{ink}933.}
\par
32280 ta belgidan iborat matnda har bir belgi 32 bit bilan kodlangan. Matn hajmini baytda toping.
\begin{enumerate}[label=\Alph*),leftmargin=9mm]
\item 1032960 bayt
\item 129120 bayt
\item 258240 bayt
\item 64560 bayt
\end{enumerate}
\begin{tcolorbox}[enhanced,breakable,colback=teal!4,colframe=teal!50!black,boxrule=0.5pt,arc=1.1mm,left=3mm,right=3mm,top=1.7mm,bottom=1.7mm,before skip=2mm,after skip=2mm]
\textbf{\color{teal}To'g'ri javob:} \textbf{B) 129120 bayt}
\par\smallskip
\textbf{Izoh:} 32280 \ensuremath{\times} 32 = 1032960 bit; 8 ga bo'lsak 129120 bayt.
\end{tcolorbox}
\vspace{2mm}\hrule height 0.2pt \vspace{3mm}

% DB question_id: d91822a4-01ac-5c5b-a3d6-a6ab74ab6bef
\needspace{5\baselineskip}
\noindent\textbf{\color{ink}934.}
\par
37259 ta belgidan iborat matnda har bir belgi 16 bit bilan kodlangan. Matn hajmini baytda toping.
\begin{enumerate}[label=\Alph*),leftmargin=9mm]
\item 74518 bayt
\item 596144 bayt
\item 37259 bayt
\item 149036 bayt
\end{enumerate}
\begin{tcolorbox}[enhanced,breakable,colback=teal!4,colframe=teal!50!black,boxrule=0.5pt,arc=1.1mm,left=3mm,right=3mm,top=1.7mm,bottom=1.7mm,before skip=2mm,after skip=2mm]
\textbf{\color{teal}To'g'ri javob:} \textbf{A) 74518 bayt}
\par\smallskip
\textbf{Izoh:} 37259 \ensuremath{\times} 16 = 596144 bit; 8 ga bo'lsak 74518 bayt.
\end{tcolorbox}
\vspace{2mm}\hrule height 0.2pt \vspace{3mm}

% DB question_id: dda80548-9f5e-5ad5-89d0-50c2354720eb
\needspace{5\baselineskip}
\noindent\textbf{\color{ink}935.}
\par
44220 ta belgidan iborat matnda har bir belgi 16 bit bilan kodlangan. Matn hajmini baytda toping.
\begin{enumerate}[label=\Alph*),leftmargin=9mm]
\item 88440 bayt
\item 44220 bayt
\item 176880 bayt
\item 707520 bayt
\end{enumerate}
\begin{tcolorbox}[enhanced,breakable,colback=teal!4,colframe=teal!50!black,boxrule=0.5pt,arc=1.1mm,left=3mm,right=3mm,top=1.7mm,bottom=1.7mm,before skip=2mm,after skip=2mm]
\textbf{\color{teal}To'g'ri javob:} \textbf{A) 88440 bayt}
\par\smallskip
\textbf{Izoh:} 44220 \ensuremath{\times} 16 = 707520 bit; 8 ga bo'lsak 88440 bayt.
\end{tcolorbox}
\vspace{2mm}\hrule height 0.2pt \vspace{3mm}

% DB question_id: e0e4d535-e677-5132-a717-09fbcecd5ae0
\needspace{5\baselineskip}
\noindent\textbf{\color{ink}936.}
\par
115017 ta belgidan iborat matnda har bir belgi 16 bit bilan kodlangan. Matn hajmini baytda toping.
\begin{enumerate}[label=\Alph*),leftmargin=9mm]
\item 1840272 bayt
\item 230034 bayt
\item 460068 bayt
\item 115017 bayt
\end{enumerate}
\begin{tcolorbox}[enhanced,breakable,colback=teal!4,colframe=teal!50!black,boxrule=0.5pt,arc=1.1mm,left=3mm,right=3mm,top=1.7mm,bottom=1.7mm,before skip=2mm,after skip=2mm]
\textbf{\color{teal}To'g'ri javob:} \textbf{B) 230034 bayt}
\par\smallskip
\textbf{Izoh:} 115017 \ensuremath{\times} 16 = 1840272 bit; 8 ga bo'lsak 230034 bayt.
\end{tcolorbox}
\vspace{2mm}\hrule height 0.2pt \vspace{3mm}

% DB question_id: e31fdafc-5799-55e5-a84f-4a9bdad5328d
\needspace{5\baselineskip}
\noindent\textbf{\color{ink}937.}
\par
128\ensuremath{\times}144 pikselli siqilmagan video 32 768 rang, 15 kadr/s va 3 minut davomiylikka ega. Video hajmini baytda toping.
\begin{enumerate}[label=\Alph*),leftmargin=9mm]
\item 93 312 000 bayt
\item 746 496 000 bayt
\item 6 220 800 bayt
\item 31 104 000 bayt
\end{enumerate}
\begin{tcolorbox}[enhanced,breakable,colback=teal!4,colframe=teal!50!black,boxrule=0.5pt,arc=1.1mm,left=3mm,right=3mm,top=1.7mm,bottom=1.7mm,before skip=2mm,after skip=2mm]
\textbf{\color{teal}To'g'ri javob:} \textbf{A) 93 312 000 bayt}
\par\smallskip
\textbf{Izoh:} 32 768 rang uchun 15 bit. 128\ensuremath{\cdot}144\ensuremath{\cdot}15\ensuremath{\cdot}15\ensuremath{\cdot}180/8=93 312 000 bayt.
\end{tcolorbox}
\vspace{2mm}\hrule height 0.2pt \vspace{3mm}

% DB question_id: e38c3b09-c9bd-45fa-8ca1-c79029999e14
\needspace{5\baselineskip}
\noindent\textbf{\color{ink}938.}
\par
100 \ensuremath{\times} 200 pikselli, 8 bit rang chuqurligidagi siqilmagan rasm hajmi qancha?
\begin{enumerate}[label=\Alph*),leftmargin=9mm]
\item 20 000 B
\item 10 000 B
\item 160 000 B
\item 200 000 B
\end{enumerate}
\begin{tcolorbox}[enhanced,breakable,colback=teal!4,colframe=teal!50!black,boxrule=0.5pt,arc=1.1mm,left=3mm,right=3mm,top=1.7mm,bottom=1.7mm,before skip=2mm,after skip=2mm]
\textbf{\color{teal}To'g'ri javob:} \textbf{A) 20 000 B}
\par\smallskip
\textbf{Izoh:} 100 \ensuremath{\times} 200 \ensuremath{\times} 8 \ensuremath{\div} 8 = 20 000 bayt.
\end{tcolorbox}
\vspace{2mm}\hrule height 0.2pt \vspace{3mm}

% DB question_id: e6d08dae-ebd7-5aaa-8127-c177f42f9e6a
\needspace{5\baselineskip}
\noindent\textbf{\color{ink}939.}
\par
154981 ta belgidan iborat matnda har bir belgi 16 bit bilan kodlangan. Matn hajmini baytda toping.
\begin{enumerate}[label=\Alph*),leftmargin=9mm]
\item 154981 bayt
\item 309962 bayt
\item 2479696 bayt
\item 619924 bayt
\end{enumerate}
\begin{tcolorbox}[enhanced,breakable,colback=teal!4,colframe=teal!50!black,boxrule=0.5pt,arc=1.1mm,left=3mm,right=3mm,top=1.7mm,bottom=1.7mm,before skip=2mm,after skip=2mm]
\textbf{\color{teal}To'g'ri javob:} \textbf{B) 309962 bayt}
\par\smallskip
\textbf{Izoh:} 154981 \ensuremath{\times} 16 = 2479696 bit; 8 ga bo'lsak 309962 bayt.
\end{tcolorbox}
\vspace{2mm}\hrule height 0.2pt \vspace{3mm}

% DB question_id: e7e168ab-5f38-58cd-a8c5-d569ae420134
\needspace{5\baselineskip}
\noindent\textbf{\color{ink}940.}
\par
189389 ta belgidan iborat matnda har bir belgi 32 bit bilan kodlangan. Matn hajmini baytda toping.
\begin{enumerate}[label=\Alph*),leftmargin=9mm]
\item 378778 bayt
\item 757556 bayt
\item 6060448 bayt
\item 1515112 bayt
\end{enumerate}
\begin{tcolorbox}[enhanced,breakable,colback=teal!4,colframe=teal!50!black,boxrule=0.5pt,arc=1.1mm,left=3mm,right=3mm,top=1.7mm,bottom=1.7mm,before skip=2mm,after skip=2mm]
\textbf{\color{teal}To'g'ri javob:} \textbf{B) 757556 bayt}
\par\smallskip
\textbf{Izoh:} 189389 \ensuremath{\times} 32 = 6060448 bit; 8 ga bo'lsak 757556 bayt.
\end{tcolorbox}
\vspace{2mm}\hrule height 0.2pt \vspace{3mm}

% DB question_id: e9615f10-8333-4ef5-bf8b-fe89d9651d24
\needspace{5\baselineskip}
\noindent\textbf{\color{ink}941.}
\par
640 \ensuremath{\times} 360, 24 bit, 25 FPS videoning siqilmagan bir soniyalik hajmi qancha?
\begin{enumerate}[label=\Alph*),leftmargin=9mm]
\item 34 560 000 B
\item 172 800 000 B
\item 17 280 000 B
\item 5 760 000 B
\end{enumerate}
\begin{tcolorbox}[enhanced,breakable,colback=teal!4,colframe=teal!50!black,boxrule=0.5pt,arc=1.1mm,left=3mm,right=3mm,top=1.7mm,bottom=1.7mm,before skip=2mm,after skip=2mm]
\textbf{\color{teal}To'g'ri javob:} \textbf{C) 17 280 000 B}
\par\smallskip
\textbf{Izoh:} 640 \ensuremath{\times} 360 \ensuremath{\times} 24 \ensuremath{\div} 8 \ensuremath{\times} 25 = 17 280 000 bayt.
\end{tcolorbox}
\vspace{2mm}\hrule height 0.2pt \vspace{3mm}

% DB question_id: eeb6664a-345e-5f7a-b4b4-a69fdd4fab11
\needspace{5\baselineskip}
\noindent\textbf{\color{ink}942.}
\par
460\ensuremath{\times}240 pikselli siqilmagan video 156 rang, 28 kadr/s va 35 soniya davomiylikka ega. Video hajmini baytda toping.
\begin{enumerate}[label=\Alph*),leftmargin=9mm]
\item 108 192 000 bayt
\item 13 524 000 bayt
\item 103 bayt
\item 865 536 000 bayt
\end{enumerate}
\begin{tcolorbox}[enhanced,breakable,colback=teal!4,colframe=teal!50!black,boxrule=0.5pt,arc=1.1mm,left=3mm,right=3mm,top=1.7mm,bottom=1.7mm,before skip=2mm,after skip=2mm]
\textbf{\color{teal}To'g'ri javob:} \textbf{A) 108 192 000 bayt}
\par\smallskip
\textbf{Izoh:} 156 rang uchun 8 bit=1 bayt/piksel. 460\ensuremath{\cdot}240\ensuremath{\cdot}28\ensuremath{\cdot}35=108 192 000 bayt.
\end{tcolorbox}
\vspace{2mm}\hrule height 0.2pt \vspace{3mm}

% DB question_id: ef4aeb2b-a063-52c2-9966-9acf48aa85cf
\needspace{5\baselineskip}
\noindent\textbf{\color{ink}943.}
\par
65284 ta belgidan iborat matnda har bir belgi 8 bit bilan kodlangan. Matn hajmini baytda toping.
\begin{enumerate}[label=\Alph*),leftmargin=9mm]
\item 130568 bayt
\item 32642 bayt
\item 65284 bayt
\item 522272 bayt
\end{enumerate}
\begin{tcolorbox}[enhanced,breakable,colback=teal!4,colframe=teal!50!black,boxrule=0.5pt,arc=1.1mm,left=3mm,right=3mm,top=1.7mm,bottom=1.7mm,before skip=2mm,after skip=2mm]
\textbf{\color{teal}To'g'ri javob:} \textbf{C) 65284 bayt}
\par\smallskip
\textbf{Izoh:} 65284 \ensuremath{\times} 8 = 522272 bit; 8 ga bo'lsak 65284 bayt.
\end{tcolorbox}
\vspace{2mm}\hrule height 0.2pt \vspace{3mm}

% DB question_id: f00dc636-4a24-5804-8af0-f2f9189473df
\needspace{5\baselineskip}
\noindent\textbf{\color{ink}944.}
\par
197754 ta belgidan iborat matnda har bir belgi 32 bit bilan kodlangan. Matn hajmini baytda toping.
\begin{enumerate}[label=\Alph*),leftmargin=9mm]
\item 791016 bayt
\item 6328128 bayt
\item 1582032 bayt
\item 395508 bayt
\end{enumerate}
\begin{tcolorbox}[enhanced,breakable,colback=teal!4,colframe=teal!50!black,boxrule=0.5pt,arc=1.1mm,left=3mm,right=3mm,top=1.7mm,bottom=1.7mm,before skip=2mm,after skip=2mm]
\textbf{\color{teal}To'g'ri javob:} \textbf{A) 791016 bayt}
\par\smallskip
\textbf{Izoh:} 197754 \ensuremath{\times} 32 = 6328128 bit; 8 ga bo'lsak 791016 bayt.
\end{tcolorbox}
\vspace{2mm}\hrule height 0.2pt \vspace{3mm}

% DB question_id: f61c9108-d92e-4291-9961-fdf31cfcc943
\needspace{5\baselineskip}
\noindent\textbf{\color{ink}945.}
\par
44,1 kHz, 16 bit, stereo audio uchun xom bit tezligi qancha?
\begin{enumerate}[label=\Alph*),leftmargin=9mm]
\item 705,6 kbit/s
\item 1 411,2 kbit/s
\item 2 822,4 kbit/s
\item 44 100 kbit/s
\end{enumerate}
\begin{tcolorbox}[enhanced,breakable,colback=teal!4,colframe=teal!50!black,boxrule=0.5pt,arc=1.1mm,left=3mm,right=3mm,top=1.7mm,bottom=1.7mm,before skip=2mm,after skip=2mm]
\textbf{\color{teal}To'g'ri javob:} \textbf{B) 1 411,2 kbit/s}
\par\smallskip
\textbf{Izoh:} 44 100 \ensuremath{\times} 16 \ensuremath{\times} 2 = 1 411 200 bit/s.
\end{tcolorbox}
\vspace{2mm}\hrule height 0.2pt \vspace{3mm}

% DB question_id: f70aabb0-743e-4f61-93c5-dd09cd318d03
\needspace{5\baselineskip}
\noindent\textbf{\color{ink}946.}
\par
8 kHz, 8 bit, mono, 10 soniyalik PCM audio hajmi qancha?
\begin{enumerate}[label=\Alph*),leftmargin=9mm]
\item 8 000 B
\item 80 000 B
\item 40 000 B
\item 640 000 B
\end{enumerate}
\begin{tcolorbox}[enhanced,breakable,colback=teal!4,colframe=teal!50!black,boxrule=0.5pt,arc=1.1mm,left=3mm,right=3mm,top=1.7mm,bottom=1.7mm,before skip=2mm,after skip=2mm]
\textbf{\color{teal}To'g'ri javob:} \textbf{B) 80 000 B}
\par\smallskip
\textbf{Izoh:} 8000 \ensuremath{\times} 8 \ensuremath{\times} 1 \ensuremath{\times} 10 \ensuremath{\div} 8 = 80 000 bayt.
\end{tcolorbox}
\vspace{2mm}\hrule height 0.2pt \vspace{3mm}

% DB question_id: f9d9e87e-8a53-42ab-8adc-65affbe18d54
\needspace{5\baselineskip}
\noindent\textbf{\color{ink}947.}
\par
16 kHz, 16 bit, stereo, 5 soniyalik PCM audio hajmi qancha?
\begin{enumerate}[label=\Alph*),leftmargin=9mm]
\item 320 000 B
\item 160 000 B
\item 640 000 B
\item 1 280 000 B
\end{enumerate}
\begin{tcolorbox}[enhanced,breakable,colback=teal!4,colframe=teal!50!black,boxrule=0.5pt,arc=1.1mm,left=3mm,right=3mm,top=1.7mm,bottom=1.7mm,before skip=2mm,after skip=2mm]
\textbf{\color{teal}To'g'ri javob:} \textbf{A) 320 000 B}
\par\smallskip
\textbf{Izoh:} 16 000 \ensuremath{\times} 16 \ensuremath{\times} 2 \ensuremath{\times} 5 \ensuremath{\div} 8 = 320 000 bayt.
\end{tcolorbox}
\vspace{2mm}\hrule height 0.2pt \vspace{3mm}

% DB question_id: fa48d0ee-26a8-5c86-9813-fcba4324e8b3
\needspace{5\baselineskip}
\noindent\textbf{\color{ink}948.}
\par
20\ensuremath{\times}25 cm o'lchamdagi siqilmagan rastr tasvirning zichligi har ikki yo'nalishda 28 ppi. Har piksel 24 xil rangdan biri bilan eng kam butun bitlarda kodlanadi. Hajmni KiB da toping. \texttt{1 dyuym = 2,5 cm}, \texttt{1 KiB = 1024 bayt}.
\begin{enumerate}[label=\Alph*),leftmargin=9mm]
\item 38,28125 KiB
\item 76,5625 KiB
\item 31,36 KiB
\item 39,2 KiB
\end{enumerate}
\begin{tcolorbox}[enhanced,breakable,colback=teal!4,colframe=teal!50!black,boxrule=0.5pt,arc=1.1mm,left=3mm,right=3mm,top=1.7mm,bottom=1.7mm,before skip=2mm,after skip=2mm]
\textbf{\color{teal}To'g'ri javob:} \textbf{A) 38,28125 KiB}
\par\smallskip
\textbf{Izoh:} O'lchamlar 8\ensuremath{\cdot}28=224 va 10\ensuremath{\cdot}28=280 piksel. 24 rang uchun 5 bit; 224\ensuremath{\cdot}280\ensuremath{\cdot}5/(8\ensuremath{\cdot}1024)=38,28125 KiB.
\end{tcolorbox}
\vspace{2mm}\hrule height 0.2pt \vspace{3mm}

% DB question_id: fec1f1dd-e479-4e14-91c6-91777b4db011
\needspace{5\baselineskip}
\noindent\textbf{\color{ink}949.}
\par
20 xil holatni kodlashda 4 bit yetmagani sabab 5 bit tanlandi. Vaziyatni tahlil qilib, unga eng aniq mos tushunchani toping.
\begin{enumerate}[label=\Alph*),leftmargin=9mm]
\item Minimal butun kod uzunligi
\item kB
\item KiB
\item Xotira sig'imi
\end{enumerate}
\begin{tcolorbox}[enhanced,breakable,colback=teal!4,colframe=teal!50!black,boxrule=0.5pt,arc=1.1mm,left=3mm,right=3mm,top=1.7mm,bottom=1.7mm,before skip=2mm,after skip=2mm]
\textbf{\color{teal}To'g'ri javob:} \textbf{A) Minimal butun kod uzunligi}
\par\smallskip
\textbf{Izoh:} Vaziyat Minimal butun kod uzunligiga misol: N holat uchun log\textsubscript{2}N yuqoriga yaxlitlangan bitlar soni.
\end{tcolorbox}
\vspace{2mm}\hrule height 0.2pt \vspace{3mm}

% DB question_id: 1fd4d393-45f4-4cdb-8bd7-8de6e572b7bc
\needspace{5\baselineskip}
\noindent\textbf{\color{ink}950.}
\par
Audio hajmi berilib, davomiylik hisoblandi. Yaqin tushunchalarni farqlab, aynan qaysi atama vaziyatni eng aniq ifodalashini toping.
\begin{enumerate}[label=\Alph*),leftmargin=9mm]
\item Mantiqiy tekshiruv
\item Formula tanlash
\item Birliklarni normallashtirish
\item Teskari masala
\end{enumerate}
\begin{tcolorbox}[enhanced,breakable,colback=teal!4,colframe=teal!50!black,boxrule=0.5pt,arc=1.1mm,left=3mm,right=3mm,top=1.7mm,bottom=1.7mm,before skip=2mm,after skip=2mm]
\textbf{\color{teal}To'g'ri javob:} \textbf{D) Teskari masala}
\par\smallskip
\textbf{Izoh:} Vaziyat Teskari masalaga misol: hajm va boshqa parametrlar ma'lum bo'lganda noma'lum parametrni formuladan ajratib topish.
\end{tcolorbox}
\vspace{2mm}\hrule height 0.2pt \vspace{3mm}

% DB question_id: 2c9f585b-76f0-4189-81c2-7911b2625f43
\needspace{5\baselineskip}
\noindent\textbf{\color{ink}951.}
\par
320 \ensuremath{\times} 240, 24 bit, 10 FPS videoning 2 soniyalik xom hajmi qancha?
\begin{enumerate}[label=\Alph*),leftmargin=9mm]
\item 46 080 000 B
\item 4 608 000 B
\item 460 800 B
\item 9 216 000 B
\end{enumerate}
\begin{tcolorbox}[enhanced,breakable,colback=teal!4,colframe=teal!50!black,boxrule=0.5pt,arc=1.1mm,left=3mm,right=3mm,top=1.7mm,bottom=1.7mm,before skip=2mm,after skip=2mm]
\textbf{\color{teal}To'g'ri javob:} \textbf{B) 4 608 000 B}
\par\smallskip
\textbf{Izoh:} 320 \ensuremath{\times} 240 \ensuremath{\times} 24 \ensuremath{\div} 8 \ensuremath{\times} 10 \ensuremath{\times} 2 = 4 608 000 bayt.
\end{tcolorbox}
\vspace{2mm}\hrule height 0.2pt \vspace{3mm}

% DB question_id: 2e04e7c2-c5e0-4379-97bc-d24b518ac613
\needspace{5\baselineskip}
\noindent\textbf{\color{ink}952.}
\par
30 xil holatni bir xil uzunlikdagi kod bilan ifodalash uchun minimal necha bit kerak?
\begin{enumerate}[label=\Alph*),leftmargin=9mm]
\item 4
\item 6
\item 30
\item 5
\end{enumerate}
\begin{tcolorbox}[enhanced,breakable,colback=teal!4,colframe=teal!50!black,boxrule=0.5pt,arc=1.1mm,left=3mm,right=3mm,top=1.7mm,bottom=1.7mm,before skip=2mm,after skip=2mm]
\textbf{\color{teal}To'g'ri javob:} \textbf{D) 5}
\par\smallskip
\textbf{Izoh:} 2\textsuperscript{4} = 16 yetmaydi, 2\textsuperscript{5} = 32 yetadi.
\end{tcolorbox}
\vspace{2mm}\hrule height 0.2pt \vspace{3mm}

% DB question_id: 56798430-a0d7-49a9-b0c7-88f421ee1155
\needspace{5\baselineskip}
\noindent\textbf{\color{ink}953.}
\par
5 sm kenglikdagi tasvir 300 PPI bilan tayyorlanadi. Taxminiy piksel kengligi qancha? 1 dyuym = 2,54 sm.
\begin{enumerate}[label=\Alph*),leftmargin=9mm]
\item 300
\item 591
\item 500
\item 762
\end{enumerate}
\begin{tcolorbox}[enhanced,breakable,colback=teal!4,colframe=teal!50!black,boxrule=0.5pt,arc=1.1mm,left=3mm,right=3mm,top=1.7mm,bottom=1.7mm,before skip=2mm,after skip=2mm]
\textbf{\color{teal}To'g'ri javob:} \textbf{B) 591}
\par\smallskip
\textbf{Izoh:} 5 \ensuremath{\div} 2,54 \ensuremath{\times} 300 \ensuremath{\approx} 590,55, ya'ni 591 piksel.
\end{tcolorbox}
\vspace{2mm}\hrule height 0.2pt \vspace{3mm}

% DB question_id: 64da09e7-e798-4eb8-904b-e3d368548a70
\needspace{5\baselineskip}
\noindent\textbf{\color{ink}954.}
\par
320 kbit/s bit tezlikdagi 5 daqiqalik siqilgan audio taxminan necha MB? MB = 10\textsuperscript{6} B.
\begin{enumerate}[label=\Alph*),leftmargin=9mm]
\item 6 MB
\item 24 MB
\item 96 MB
\item 12 MB
\end{enumerate}
\begin{tcolorbox}[enhanced,breakable,colback=teal!4,colframe=teal!50!black,boxrule=0.5pt,arc=1.1mm,left=3mm,right=3mm,top=1.7mm,bottom=1.7mm,before skip=2mm,after skip=2mm]
\textbf{\color{teal}To'g'ri javob:} \textbf{D) 12 MB}
\par\smallskip
\textbf{Izoh:} 320 000 \ensuremath{\times} 300 \ensuremath{\div} 8 = 12 000 000 B = 12 MB.
\end{tcolorbox}
\vspace{2mm}\hrule height 0.2pt \vspace{3mm}

% DB question_id: 8ad2fc2d-ed8f-4836-88b0-e35fb0099b61
\needspace{5\baselineskip}
\noindent\textbf{\color{ink}955.}
\par
64 MiB xotiraga 1,5 MiB hajmli fayldan ko'pi bilan nechta to'liq fayl sig'adi?
\begin{enumerate}[label=\Alph*),leftmargin=9mm]
\item 42
\item 41
\item 43
\item 44
\end{enumerate}
\begin{tcolorbox}[enhanced,breakable,colback=teal!4,colframe=teal!50!black,boxrule=0.5pt,arc=1.1mm,left=3mm,right=3mm,top=1.7mm,bottom=1.7mm,before skip=2mm,after skip=2mm]
\textbf{\color{teal}To'g'ri javob:} \textbf{A) 42}
\par\smallskip
\textbf{Izoh:} 64 \ensuremath{\div} 1,5 = 42,66; faqat 42 ta to'liq fayl.
\end{tcolorbox}
\vspace{2mm}\hrule height 0.2pt \vspace{3mm}

% DB question_id: aaf03d49-8933-4697-87be-6e2cbf04d49a
\needspace{5\baselineskip}
\noindent\textbf{\color{ink}956.}
\par
40 MB fayl 10 MB bo'lganda 4:1 nisbat qayd etildi. Vaziyatni tahlil qilib, unga eng aniq mos tushunchani toping.
\begin{enumerate}[label=\Alph*),leftmargin=9mm]
\item Kodek
\item Siqish koeffitsiyenti
\item Transkodlash
\item Zaxira nusxasi
\end{enumerate}
\begin{tcolorbox}[enhanced,breakable,colback=teal!4,colframe=teal!50!black,boxrule=0.5pt,arc=1.1mm,left=3mm,right=3mm,top=1.7mm,bottom=1.7mm,before skip=2mm,after skip=2mm]
\textbf{\color{teal}To'g'ri javob:} \textbf{B) Siqish koeffitsiyenti}
\par\smallskip
\textbf{Izoh:} Vaziyat Siqish koeffitsiyentiga misol: asl hajmning siqilgan hajmga yoki kamayish foiziga nisbati.
\end{tcolorbox}
\vspace{2mm}\hrule height 0.2pt \vspace{3mm}

% DB question_id: b8c0a041-bb69-4c43-8358-5c16249fb4ed
\needspace{5\baselineskip}
\noindent\textbf{\color{ink}957.}
\par
22 kHz, 8 bit, mono, 60 soniyalik siqilmagan audio hajmi qancha?
\begin{enumerate}[label=\Alph*),leftmargin=9mm]
\item 10 560 000 B
\item 165 000 B
\item 1 320 000 B
\item 2 640 000 B
\end{enumerate}
\begin{tcolorbox}[enhanced,breakable,colback=teal!4,colframe=teal!50!black,boxrule=0.5pt,arc=1.1mm,left=3mm,right=3mm,top=1.7mm,bottom=1.7mm,before skip=2mm,after skip=2mm]
\textbf{\color{teal}To'g'ri javob:} \textbf{C) 1 320 000 B}
\par\smallskip
\textbf{Izoh:} 22 000 \ensuremath{\times} 8 \ensuremath{\times} 1 \ensuremath{\times} 60 \ensuremath{\div} 8 = 1 320 000 bayt.
\end{tcolorbox}
\vspace{2mm}\hrule height 0.2pt \vspace{3mm}

% DB question_id: c1fae011-49ba-456a-8e58-3d3756846579
\needspace{5\baselineskip}
\noindent\textbf{\color{ink}958.}
\par
60 MB fayl 25\% ga kamaytirildi. Yangi hajm qancha?
\begin{enumerate}[label=\Alph*),leftmargin=9mm]
\item 15 MB
\item 35 MB
\item 45 MB
\item 75 MB
\end{enumerate}
\begin{tcolorbox}[enhanced,breakable,colback=teal!4,colframe=teal!50!black,boxrule=0.5pt,arc=1.1mm,left=3mm,right=3mm,top=1.7mm,bottom=1.7mm,before skip=2mm,after skip=2mm]
\textbf{\color{teal}To'g'ri javob:} \textbf{C) 45 MB}
\par\smallskip
\textbf{Izoh:} 60 \ensuremath{\times} (1 - 0,25) = 45 MB.
\end{tcolorbox}
\vspace{2mm}\hrule height 0.2pt \vspace{3mm}

% DB question_id: d9ab6356-005b-4bc0-a314-5b4e1aecbf85
\needspace{5\baselineskip}
\noindent\textbf{\color{ink}959.}
\par
Kadr o'lchami va rang chuqurligi o'zgarmay, FPS 25 dan 50 ga oshsa, siqilmagan hajm qanday o'zgaradi?
\begin{enumerate}[label=\Alph*),leftmargin=9mm]
\item yarim bo'ladi
\item o'zgarmaydi
\item to'rt marta oshadi
\item ikki marta oshadi
\end{enumerate}
\begin{tcolorbox}[enhanced,breakable,colback=teal!4,colframe=teal!50!black,boxrule=0.5pt,arc=1.1mm,left=3mm,right=3mm,top=1.7mm,bottom=1.7mm,before skip=2mm,after skip=2mm]
\textbf{\color{teal}To'g'ri javob:} \textbf{D) ikki marta oshadi}
\par\smallskip
\textbf{Izoh:} Hajm FPSga to'g'ri proporsional.
\end{tcolorbox}
\vspace{2mm}\hrule height 0.2pt \vspace{3mm}

% DB question_id: 073c348d-47b4-4f5e-a95a-b74b8acdb949
\needspace{5\baselineskip}
\noindent\textbf{\color{ink}960.}
\par
Stereo audio mono audiodan kichik chiqib qolgan natija qayta tekshirildi. Yaqin tushunchalarni farqlab, aynan qaysi atama vaziyatni eng aniq ifodalashini toping.
\begin{enumerate}[label=\Alph*),leftmargin=9mm]
\item Mantiqiy tekshiruv
\item Formula tanlash
\item Birliklarni normallashtirish
\item Xom hajm
\end{enumerate}
\begin{tcolorbox}[enhanced,breakable,colback=teal!4,colframe=teal!50!black,boxrule=0.5pt,arc=1.1mm,left=3mm,right=3mm,top=1.7mm,bottom=1.7mm,before skip=2mm,after skip=2mm]
\textbf{\color{teal}To'g'ri javob:} \textbf{A) Mantiqiy tekshiruv}
\par\smallskip
\textbf{Izoh:} Vaziyat Mantiqiy tekshiruvga misol: natijaning birlik, tartib va shartga mosligini baholash bosqichi.
\end{tcolorbox}
\vspace{2mm}\hrule height 0.2pt \vspace{3mm}

% DB question_id: 3d70d38a-6b2e-4cc2-aa56-b7a3aa24f772
\needspace{5\baselineskip}
\noindent\textbf{\color{ink}961.}
\par
64 MiB kartaga nechta 1,5 MiB fayl sig'ishi hisoblandi. Yaqin tushunchalarni farqlab, aynan qaysi atama vaziyatni eng aniq ifodalashini toping.
\begin{enumerate}[label=\Alph*),leftmargin=9mm]
\item Bit
\item Bayt
\item Xotira sig'imi
\item Xartli formulasi
\end{enumerate}
\begin{tcolorbox}[enhanced,breakable,colback=teal!4,colframe=teal!50!black,boxrule=0.5pt,arc=1.1mm,left=3mm,right=3mm,top=1.7mm,bottom=1.7mm,before skip=2mm,after skip=2mm]
\textbf{\color{teal}To'g'ri javob:} \textbf{C) Xotira sig'imi}
\par\smallskip
\textbf{Izoh:} Vaziyat Xotira sig'imiga misol: tashuvchida saqlanishi mumkin bo'lgan axborotning umumiy hajmi.
\end{tcolorbox}
\vspace{2mm}\hrule height 0.2pt \vspace{3mm}

% DB question_id: 560cfd5d-6ec7-4e7f-b6d1-319ccc9fd18d
\needspace{5\baselineskip}
\noindent\textbf{\color{ink}962.}
\par
16 bit mono yozuv stereo qilindi, qolgan parametrlar o'zgarmadi. Hajmga ta'siri qanday?
\begin{enumerate}[label=\Alph*),leftmargin=9mm]
\item yarimlashadi
\item o'zgarmaydi
\item ikki marta oshadi
\item to'rt marta oshadi
\end{enumerate}
\begin{tcolorbox}[enhanced,breakable,colback=teal!4,colframe=teal!50!black,boxrule=0.5pt,arc=1.1mm,left=3mm,right=3mm,top=1.7mm,bottom=1.7mm,before skip=2mm,after skip=2mm]
\textbf{\color{teal}To'g'ri javob:} \textbf{C) ikki marta oshadi}
\par\smallskip
\textbf{Izoh:} Kanal soni 1 dan 2 ga oshgani uchun hajm ikki baravar bo'ladi.
\end{tcolorbox}
\vspace{2mm}\hrule height 0.2pt \vspace{3mm}

% DB question_id: 7bfdf143-64bb-47a3-a099-5fd7aea097db
\needspace{5\baselineskip}
\noindent\textbf{\color{ink}963.}
\par
Video 6 Mbit/s va audio 192 kbit/s. Umumiy oqim tezligi taxminan qancha?
\begin{enumerate}[label=\Alph*),leftmargin=9mm]
\item 5,808 Mbit/s
\item 6,000 Mbit/s
\item 192 Mbit/s
\item 6,192 Mbit/s
\end{enumerate}
\begin{tcolorbox}[enhanced,breakable,colback=teal!4,colframe=teal!50!black,boxrule=0.5pt,arc=1.1mm,left=3mm,right=3mm,top=1.7mm,bottom=1.7mm,before skip=2mm,after skip=2mm]
\textbf{\color{teal}To'g'ri javob:} \textbf{D) 6,192 Mbit/s}
\par\smallskip
\textbf{Izoh:} 192 kbit/s = 0,192 Mbit/s; jami 6,192 Mbit/s.
\end{tcolorbox}
\vspace{2mm}\hrule height 0.2pt \vspace{3mm}

% DB question_id: 8839e247-d355-4182-9e8c-2518b3cf1b84
\needspace{5\baselineskip}
\noindent\textbf{\color{ink}964.}
\par
Siqilgan trek 192 kbit/s qiymat bilan kodlandi. Yaqin tushunchalarni farqlab, aynan qaysi atama vaziyatni eng aniq ifodalashini toping.
\begin{enumerate}[label=\Alph*),leftmargin=9mm]
\item Analog tovush
\item Audio bit tezligi
\item ADC
\item Namuna chastotasi
\end{enumerate}
\begin{tcolorbox}[enhanced,breakable,colback=teal!4,colframe=teal!50!black,boxrule=0.5pt,arc=1.1mm,left=3mm,right=3mm,top=1.7mm,bottom=1.7mm,before skip=2mm,after skip=2mm]
\textbf{\color{teal}To'g'ri javob:} \textbf{B) Audio bit tezligi}
\par\smallskip
\textbf{Izoh:} Vaziyat Audio bit tezligiga misol: bir soniyada uzatiladigan yoki saqlanadigan audio bitlari soni.
\end{tcolorbox}
\vspace{2mm}\hrule height 0.2pt \vspace{3mm}
\section{Attestatsiyadagi nozik farqlar va diagnostika}

% DB question_id: e56f8fbe-98eb-5ca1-a25e-14a4995e8cdf
\needspace{5\baselineskip}
\noindent\textbf{\color{ink}965.}
\par
Boshqalarning hujjat yoki ishlarini ruxsatsiz, manbasini ko'rsatmasdan o'ziniki kabi ishlatish nima deyiladi?
\begin{enumerate}[label=\Alph*),leftmargin=9mm]
\item Reference
\item Plagiat
\item Compose
\item Piktogramma
\end{enumerate}
\begin{tcolorbox}[enhanced,breakable,colback=teal!4,colframe=teal!50!black,boxrule=0.5pt,arc=1.1mm,left=3mm,right=3mm,top=1.7mm,bottom=1.7mm,before skip=2mm,after skip=2mm]
\textbf{\color{teal}To'g'ri javob:} \textbf{B) Plagiat}
\par\smallskip
\textbf{Izoh:} Boshqa muallif ishini manbasiz o'zlashtirish plagiat hisoblanadi.
\end{tcolorbox}
\vspace{2mm}\hrule height 0.2pt \vspace{3mm}
\section{Bir qarashda: formulalar va birliklar}

% DB question_id: b22b43c9-22f0-5611-9f2f-bbee183b361d
\needspace{5\baselineskip}
\noindent\textbf{\color{ink}966.}
\par
Ikkilik birliklardan foydalaning: 2048 KiB + 3 MiB necha MiB bo'ladi?
\begin{enumerate}[label=\Alph*),leftmargin=9mm]
\item 7,5 MiB
\item 5 KiB
\item 5 MiB
\item 2051 KiB
\item 2051 MiB
\end{enumerate}
\begin{tcolorbox}[enhanced,breakable,colback=teal!4,colframe=teal!50!black,boxrule=0.5pt,arc=1.1mm,left=3mm,right=3mm,top=1.7mm,bottom=1.7mm,before skip=2mm,after skip=2mm]
\textbf{\color{teal}To'g'ri javob:} \textbf{C) 5 MiB}
\par\smallskip
\textbf{Izoh:} 1024 KiB=1 MiB, shuning uchun 2048 KiB=2 MiB; 2+3=5 MiB.
\end{tcolorbox}
\vspace{2mm}\hrule height 0.2pt \vspace{3mm}
\section{Axborot uzatish tezligi: hajm, tezlik va vaqt}

% DB question_id: 009f6fd6-fc68-4cd5-8be1-4b59fb60c2ad
\needspace{5\baselineskip}
\noindent\textbf{\color{ink}967.}
\par
"Uzatish tezligi" atamasiga eng aniq ta'rifni toping.
\begin{enumerate}[label=\Alph*),leftmargin=9mm]
\item axborot hajmining amaliy uzatish tezligiga nisbatidan topiladigan vaqt
\item vaqt birligida uzatilgan axborot hajmi
\item kanalning nazariy yoki belgilangan maksimal o'tkazish quvvati
\item foydalanuvchi amalda oladigan foydali ma'lumot uzatish tezligi
\end{enumerate}
\begin{tcolorbox}[enhanced,breakable,colback=teal!4,colframe=teal!50!black,boxrule=0.5pt,arc=1.1mm,left=3mm,right=3mm,top=1.7mm,bottom=1.7mm,before skip=2mm,after skip=2mm]
\textbf{\color{teal}To'g'ri javob:} \textbf{B) vaqt birligida uzatilgan axborot hajmi}
\par\smallskip
\textbf{Izoh:} Uzatish tezligi - vaqt birligida uzatilgan axborot hajmi.
\end{tcolorbox}
\vspace{2mm}\hrule height 0.2pt \vspace{3mm}

% DB question_id: 11441353-5294-4c5e-8501-494c95967b61
\needspace{5\baselineskip}
\noindent\textbf{\color{ink}968.}
\par
"Uzatish vaqti" atamasiga eng aniq ta'rifni toping.
\begin{enumerate}[label=\Alph*),leftmargin=9mm]
\item kanalning nazariy yoki belgilangan maksimal o'tkazish quvvati
\item foydalanuvchi amalda oladigan foydali ma'lumot uzatish tezligi
\item axborot hajmining amaliy uzatish tezligiga nisbatidan topiladigan vaqt
\item bir soniyada uzatiladigan signal simvollari soni
\end{enumerate}
\begin{tcolorbox}[enhanced,breakable,colback=teal!4,colframe=teal!50!black,boxrule=0.5pt,arc=1.1mm,left=3mm,right=3mm,top=1.7mm,bottom=1.7mm,before skip=2mm,after skip=2mm]
\textbf{\color{teal}To'g'ri javob:} \textbf{C) axborot hajmining amaliy uzatish tezligiga nisbatidan topiladigan vaqt}
\par\smallskip
\textbf{Izoh:} Uzatish vaqti - axborot hajmining amaliy uzatish tezligiga nisbatidan topiladigan vaqt.
\end{tcolorbox}
\vspace{2mm}\hrule height 0.2pt \vspace{3mm}

% DB question_id: 32166b06-8014-40c7-8dff-31a7dbf4d942
\needspace{5\baselineskip}
\noindent\textbf{\color{ink}969.}
\par
"Bandwidth" atamasiga eng aniq ta'rifni toping.
\begin{enumerate}[label=\Alph*),leftmargin=9mm]
\item foydalanuvchi amalda oladigan foydali ma'lumot uzatish tezligi
\item bir soniyada uzatiladigan signal simvollari soni
\item axborot faqat bir yo'nalishda uzatiladigan aloqa
\item kanalning nazariy yoki belgilangan maksimal o'tkazish quvvati
\end{enumerate}
\begin{tcolorbox}[enhanced,breakable,colback=teal!4,colframe=teal!50!black,boxrule=0.5pt,arc=1.1mm,left=3mm,right=3mm,top=1.7mm,bottom=1.7mm,before skip=2mm,after skip=2mm]
\textbf{\color{teal}To'g'ri javob:} \textbf{D) kanalning nazariy yoki belgilangan maksimal o'tkazish quvvati}
\par\smallskip
\textbf{Izoh:} Bandwidth - kanalning nazariy yoki belgilangan maksimal o'tkazish quvvati.
\end{tcolorbox}
\vspace{2mm}\hrule height 0.2pt \vspace{3mm}

% DB question_id: 58705ece-a3de-4cea-8f07-ae7f5858f68d
\needspace{5\baselineskip}
\noindent\textbf{\color{ink}970.}
\par
Quyidagi ta'rif qaysi atamaga tegishli? Vaqt birligida uzatilgan axborot hajmi.
\begin{enumerate}[label=\Alph*),leftmargin=9mm]
\item Uzatish vaqti
\item Uzatish tezligi
\item Bandwidth
\item Throughput
\end{enumerate}
\begin{tcolorbox}[enhanced,breakable,colback=teal!4,colframe=teal!50!black,boxrule=0.5pt,arc=1.1mm,left=3mm,right=3mm,top=1.7mm,bottom=1.7mm,before skip=2mm,after skip=2mm]
\textbf{\color{teal}To'g'ri javob:} \textbf{B) Uzatish tezligi}
\par\smallskip
\textbf{Izoh:} Ta'rif Uzatish tezligi tushunchasiga tegishli.
\end{tcolorbox}
\vspace{2mm}\hrule height 0.2pt \vspace{3mm}

% DB question_id: ae364ac1-4d46-4cb1-9dfe-6788eaf0b63c
\needspace{5\baselineskip}
\noindent\textbf{\color{ink}971.}
\par
"Baud" atamasiga eng aniq ta'rifni toping.
\begin{enumerate}[label=\Alph*),leftmargin=9mm]
\item axborot faqat bir yo'nalishda uzatiladigan aloqa
\item bir soniyada uzatiladigan signal simvollari soni
\item ikki yo'nalish mavjud, ammo bir vaqtda faqat bittasi ishlaydigan aloqa
\item axborot ikki yo'nalishda bir vaqtda uzatiladigan aloqa
\end{enumerate}
\begin{tcolorbox}[enhanced,breakable,colback=teal!4,colframe=teal!50!black,boxrule=0.5pt,arc=1.1mm,left=3mm,right=3mm,top=1.7mm,bottom=1.7mm,before skip=2mm,after skip=2mm]
\textbf{\color{teal}To'g'ri javob:} \textbf{B) bir soniyada uzatiladigan signal simvollari soni}
\par\smallskip
\textbf{Izoh:} Baud - bir soniyada uzatiladigan signal simvollari soni.
\end{tcolorbox}
\vspace{2mm}\hrule height 0.2pt \vspace{3mm}

% DB question_id: c16a0c71-930b-47a4-9b7e-a554dfebb12b
\needspace{5\baselineskip}
\noindent\textbf{\color{ink}972.}
\par
Quyidagi ta'rif qaysi atamaga tegishli? Axborot hajmining amaliy uzatish tezligiga nisbatidan topiladigan vaqt.
\begin{enumerate}[label=\Alph*),leftmargin=9mm]
\item Bandwidth
\item Throughput
\item Uzatish vaqti
\item Baud
\end{enumerate}
\begin{tcolorbox}[enhanced,breakable,colback=teal!4,colframe=teal!50!black,boxrule=0.5pt,arc=1.1mm,left=3mm,right=3mm,top=1.7mm,bottom=1.7mm,before skip=2mm,after skip=2mm]
\textbf{\color{teal}To'g'ri javob:} \textbf{C) Uzatish vaqti}
\par\smallskip
\textbf{Izoh:} Ta'rif Uzatish vaqti tushunchasiga tegishli.
\end{tcolorbox}
\vspace{2mm}\hrule height 0.2pt \vspace{3mm}

% DB question_id: cdba1dbc-d47c-43c9-b904-a0e76e0c8877
\needspace{5\baselineskip}
\noindent\textbf{\color{ink}973.}
\par
"Baud va bit/s farqi" atamasiga eng aniq ta'rifni toping.
\begin{enumerate}[label=\Alph*),leftmargin=9mm]
\item HTTPS transportni shifrlashi, sahifa mazmunining haqiqatligini kafolatlamasligini ajratish
\item ma'lumot qayd etilgan qiymat, axborot esa qabul qiluvchi uchun mazmunli ma'lumot ekanini ajratish
\item baud simvol tezligi, bit/s esa axborot bitlari tezligi ekanini ajratish
\item kodlash taqdim etish formatini, shifrlash esa maxfiylikni ko'zlashini ajratish
\end{enumerate}
\begin{tcolorbox}[enhanced,breakable,colback=teal!4,colframe=teal!50!black,boxrule=0.5pt,arc=1.1mm,left=3mm,right=3mm,top=1.7mm,bottom=1.7mm,before skip=2mm,after skip=2mm]
\textbf{\color{teal}To'g'ri javob:} \textbf{C) baud simvol tezligi, bit/s esa axborot bitlari tezligi ekanini ajratish}
\par\smallskip
\textbf{Izoh:} Baud va bit/s farqi - baud simvol tezligi, bit/s esa axborot bitlari tezligi ekanini ajratish.
\end{tcolorbox}
\vspace{2mm}\hrule height 0.2pt \vspace{3mm}

% DB question_id: d7b5f091-f756-4a5b-930b-6e63c81cd18d
\needspace{5\baselineskip}
\noindent\textbf{\color{ink}974.}
\par
"Throughput" atamasiga eng aniq ta'rifni toping.
\begin{enumerate}[label=\Alph*),leftmargin=9mm]
\item foydalanuvchi amalda oladigan foydali ma'lumot uzatish tezligi
\item bir soniyada uzatiladigan signal simvollari soni
\item axborot faqat bir yo'nalishda uzatiladigan aloqa
\item ikki yo'nalish mavjud, ammo bir vaqtda faqat bittasi ishlaydigan aloqa
\end{enumerate}
\begin{tcolorbox}[enhanced,breakable,colback=teal!4,colframe=teal!50!black,boxrule=0.5pt,arc=1.1mm,left=3mm,right=3mm,top=1.7mm,bottom=1.7mm,before skip=2mm,after skip=2mm]
\textbf{\color{teal}To'g'ri javob:} \textbf{A) foydalanuvchi amalda oladigan foydali ma'lumot uzatish tezligi}
\par\smallskip
\textbf{Izoh:} Throughput - foydalanuvchi amalda oladigan foydali ma'lumot uzatish tezligi.
\end{tcolorbox}
\vspace{2mm}\hrule height 0.2pt \vspace{3mm}

% DB question_id: 04bc41f0-b8d8-4142-b145-36a21c19d721
\needspace{5\baselineskip}
\noindent\textbf{\color{ink}975.}
\par
Internetdan 600 MB hajmli faylning 10\% qismi 8 Mbit/s, 60\% qismi 2 MB/s, qolgan qismi 24 Mbit/s tezlikda yuklandi. Umumiy yuklash vaqtini toping.
\begin{enumerate}[label=\Alph*),leftmargin=9mm]
\item 6 daqiqa
\item 4 daqiqa
\item 5 daqiqa
\item 7 daqiqa
\end{enumerate}
\begin{tcolorbox}[enhanced,breakable,colback=teal!4,colframe=teal!50!black,boxrule=0.5pt,arc=1.1mm,left=3mm,right=3mm,top=1.7mm,bottom=1.7mm,before skip=2mm,after skip=2mm]
\textbf{\color{teal}To'g'ri javob:} \textbf{C) 5 daqiqa}
\par\smallskip
\textbf{Izoh:} 60 MB / 1 MB/s = 60 s; 360 MB / 2 MB/s = 180 s; 180 MB / 3 MB/s = 60 s. Jami 300 s = 5 daqiqa.
\end{tcolorbox}
\vspace{2mm}\hrule height 0.2pt \vspace{3mm}

% DB question_id: 0fc903ec-6d92-4418-aa97-1a4e9bc4b92b
\needspace{5\baselineskip}
\noindent\textbf{\color{ink}976.}
\par
Quyidagi ta'rif qaysi atamaga tegishli? Kanalning nazariy yoki belgilangan maksimal o'tkazish quvvati.
\begin{enumerate}[label=\Alph*),leftmargin=9mm]
\item Throughput
\item Baud
\item Simpleks aloqa
\item Bandwidth
\end{enumerate}
\begin{tcolorbox}[enhanced,breakable,colback=teal!4,colframe=teal!50!black,boxrule=0.5pt,arc=1.1mm,left=3mm,right=3mm,top=1.7mm,bottom=1.7mm,before skip=2mm,after skip=2mm]
\textbf{\color{teal}To'g'ri javob:} \textbf{D) Bandwidth}
\par\smallskip
\textbf{Izoh:} Ta'rif Bandwidth tushunchasiga tegishli.
\end{tcolorbox}
\vspace{2mm}\hrule height 0.2pt \vspace{3mm}

% DB question_id: 24a8ad01-2717-51fc-8b99-27c5e57e9cea
\needspace{5\baselineskip}
\noindent\textbf{\color{ink}977.}
\par
750 sahifali kitobning har sahifasida 56 qator, har qatorda 32 ta bir baytli belgi bor. Kitob 25 Kibit/s tezlikda uzatilsa, uzatish necha minut davom etadi? \texttt{1 Kibit = 1024 bit}.
\begin{enumerate}[label=\Alph*),leftmargin=9mm]
\item 7 minut
\item 8,75 minut
\item 70 minut
\item 420 minut
\end{enumerate}
\begin{tcolorbox}[enhanced,breakable,colback=teal!4,colframe=teal!50!black,boxrule=0.5pt,arc=1.1mm,left=3mm,right=3mm,top=1.7mm,bottom=1.7mm,before skip=2mm,after skip=2mm]
\textbf{\color{teal}To'g'ri javob:} \textbf{A) 7 minut}
\par\smallskip
\textbf{Izoh:} Hajm 750\ensuremath{\cdot}56\ensuremath{\cdot}32\ensuremath{\cdot}8=10 752 000 bit. Tezlik 25\ensuremath{\cdot}1024=25 600 bit/s; vaqt 420 s=7 minut.
\end{tcolorbox}
\vspace{2mm}\hrule height 0.2pt \vspace{3mm}

% DB question_id: 26a55bb8-82a3-4fac-814e-eb849828c88d
\needspace{5\baselineskip}
\noindent\textbf{\color{ink}978.}
\par
Hisobot 504 bet, har bet 80 satr, har satr 70 belgidan iborat. 1 belgi=1 bayt. Hisobot 24 soat 30 daqiqada o'qildi. O'rtacha o'qish tezligini bayt/sekundda toping.
\begin{enumerate}[label=\Alph*),leftmargin=9mm]
\item 32 bit/s
\item 56 bit/s
\item 115 B/s
\item 32 B/s
\end{enumerate}
\begin{tcolorbox}[enhanced,breakable,colback=teal!4,colframe=teal!50!black,boxrule=0.5pt,arc=1.1mm,left=3mm,right=3mm,top=1.7mm,bottom=1.7mm,before skip=2mm,after skip=2mm]
\textbf{\color{teal}To'g'ri javob:} \textbf{D) 32 B/s}
\par\smallskip
\textbf{Izoh:} Hajm=504\ensuremath{\cdot}80\ensuremath{\cdot}70=2 822 400 bayt. Vaqt=(24\ensuremath{\cdot}60+30)\ensuremath{\cdot}60=88 200 s. Tezlik=2 822 400/88 200=32 B/s. [Yutuqli\_test, test kodi 12; Telegram xabarlar 6446 va 6556.]
\end{tcolorbox}
\vspace{2mm}\hrule height 0.2pt \vspace{3mm}

% DB question_id: 274e7b4c-2933-5c15-a715-db01f0897601
\needspace{5\baselineskip}
\noindent\textbf{\color{ink}979.}
\par
Har pikseli 3 baytda kodlangan 640\ensuremath{\times}480 pikselli siqilmagan tasvir 28 800 bit/s modem orqali yuboriladi. Uzatish necha soniya davom etadi?
\begin{enumerate}[label=\Alph*),leftmargin=9mm]
\item 256 soniya
\item 32 soniya
\item 25,6 soniya
\item 2048 soniya
\end{enumerate}
\begin{tcolorbox}[enhanced,breakable,colback=teal!4,colframe=teal!50!black,boxrule=0.5pt,arc=1.1mm,left=3mm,right=3mm,top=1.7mm,bottom=1.7mm,before skip=2mm,after skip=2mm]
\textbf{\color{teal}To'g'ri javob:} \textbf{A) 256 soniya}
\par\smallskip
\textbf{Izoh:} Hajm 640\ensuremath{\cdot}480\ensuremath{\cdot}3\ensuremath{\cdot}8=7 372 800 bit. 7 372 800/28 800=256 soniya.
\end{tcolorbox}
\vspace{2mm}\hrule height 0.2pt \vspace{3mm}

% DB question_id: 49df163f-a1d6-4a37-8715-e80d26766af8
\needspace{5\baselineskip}
\noindent\textbf{\color{ink}980.}
\par
Fayl hajmi bitga o'tkazilib, kanal tezligiga bo'lindi. Bu vaziyatdagi asosiy tushunchani aniqlang.
\begin{enumerate}[label=\Alph*),leftmargin=9mm]
\item Bandwidth
\item Throughput
\item Uzatish vaqti
\item Baud
\end{enumerate}
\begin{tcolorbox}[enhanced,breakable,colback=teal!4,colframe=teal!50!black,boxrule=0.5pt,arc=1.1mm,left=3mm,right=3mm,top=1.7mm,bottom=1.7mm,before skip=2mm,after skip=2mm]
\textbf{\color{teal}To'g'ri javob:} \textbf{C) Uzatish vaqti}
\par\smallskip
\textbf{Izoh:} Vaziyat Uzatish vaqtiga misol: axborot hajmining amaliy uzatish tezligiga nisbatidan topiladigan vaqt.
\end{tcolorbox}
\vspace{2mm}\hrule height 0.2pt \vspace{3mm}

% DB question_id: 4ef2e8f3-bc09-4a48-ab03-7c2efdfc278c
\needspace{5\baselineskip}
\noindent\textbf{\color{ink}981.}
\par
Kitob 128 sahifa, har sahifa 32 qator, har qator 60 belgi. Har bir belgi 1 bayt deb olinadi. Kitob 512 bit/s tezlikda uzatilsa qancha vaqt ketadi?
\begin{enumerate}[label=\Alph*),leftmargin=9mm]
\item 3600 sekund
\item 60 minut
\item 240 sekund
\item 64 minut
\end{enumerate}
\begin{tcolorbox}[enhanced,breakable,colback=teal!4,colframe=teal!50!black,boxrule=0.5pt,arc=1.1mm,left=3mm,right=3mm,top=1.7mm,bottom=1.7mm,before skip=2mm,after skip=2mm]
\textbf{\color{teal}To'g'ri javob:} \textbf{D) 64 minut}
\par\smallskip
\textbf{Izoh:} Hajm 128\ensuremath{\cdot}32\ensuremath{\cdot}60\ensuremath{\cdot}8=1 966 080 bit. 1 966 080/512=3840 sekund=64 minut.
\end{tcolorbox}
\vspace{2mm}\hrule height 0.2pt \vspace{3mm}

% DB question_id: 8cb590e9-e8de-40c1-9186-bdc35c947cb2
\needspace{5\baselineskip}
\noindent\textbf{\color{ink}982.}
\par
Quyidagi ta'rif qaysi atamaga tegishli? Baud simvol tezligi, bit/s esa axborot bitlari tezligi ekanini ajratish.
\begin{enumerate}[label=\Alph*),leftmargin=9mm]
\item Baud va bit/s farqi
\item HTTPS va ishonchlilik farqi
\item Axborot va ma'lumot farqi
\item Kodlash va shifrlash farqi
\end{enumerate}
\begin{tcolorbox}[enhanced,breakable,colback=teal!4,colframe=teal!50!black,boxrule=0.5pt,arc=1.1mm,left=3mm,right=3mm,top=1.7mm,bottom=1.7mm,before skip=2mm,after skip=2mm]
\textbf{\color{teal}To'g'ri javob:} \textbf{A) Baud va bit/s farqi}
\par\smallskip
\textbf{Izoh:} Ta'rif Baud va bit/s farqi tushunchasiga tegishli.
\end{tcolorbox}
\vspace{2mm}\hrule height 0.2pt \vspace{3mm}

% DB question_id: b9fbaea3-59fd-4650-b8ea-057e91ad6909
\needspace{5\baselineskip}
\noindent\textbf{\color{ink}983.}
\par
Kanal bir soniyada 20 million bit uzatdi. Bu vaziyatdagi asosiy tushunchani aniqlang.
\begin{enumerate}[label=\Alph*),leftmargin=9mm]
\item Uzatish vaqti
\item Uzatish tezligi
\item Bandwidth
\item Throughput
\end{enumerate}
\begin{tcolorbox}[enhanced,breakable,colback=teal!4,colframe=teal!50!black,boxrule=0.5pt,arc=1.1mm,left=3mm,right=3mm,top=1.7mm,bottom=1.7mm,before skip=2mm,after skip=2mm]
\textbf{\color{teal}To'g'ri javob:} \textbf{B) Uzatish tezligi}
\par\smallskip
\textbf{Izoh:} Vaziyat Uzatish tezligiga misol: vaqt birligida uzatilgan axborot hajmi.
\end{tcolorbox}
\vspace{2mm}\hrule height 0.2pt \vspace{3mm}

% DB question_id: c7acecf8-43b4-44a1-b14a-4f9265034de3
\needspace{5\baselineskip}
\noindent\textbf{\color{ink}984.}
\par
100 Mbit/s kanal sinovda 72 Mbit/s foydali tezlik berdi. Bu vaziyatdagi asosiy tushunchani aniqlang.
\begin{enumerate}[label=\Alph*),leftmargin=9mm]
\item Throughput
\item Baud
\item Simpleks aloqa
\item Yarim dupleks
\end{enumerate}
\begin{tcolorbox}[enhanced,breakable,colback=teal!4,colframe=teal!50!black,boxrule=0.5pt,arc=1.1mm,left=3mm,right=3mm,top=1.7mm,bottom=1.7mm,before skip=2mm,after skip=2mm]
\textbf{\color{teal}To'g'ri javob:} \textbf{A) Throughput}
\par\smallskip
\textbf{Izoh:} Vaziyat Throughputga misol: foydalanuvchi amalda oladigan foydali ma'lumot uzatish tezligi.
\end{tcolorbox}
\vspace{2mm}\hrule height 0.2pt \vspace{3mm}

% DB question_id: de19d837-d3f0-4555-b4bc-84dfbde1f2f9
\needspace{5\baselineskip}
\noindent\textbf{\color{ink}985.}
\par
Provayder ulanish uchun 100 Mbit/s yuqori chegarani ko'rsatdi. Bu vaziyatdagi asosiy tushunchani aniqlang.
\begin{enumerate}[label=\Alph*),leftmargin=9mm]
\item Throughput
\item Baud
\item Simpleks aloqa
\item Bandwidth
\end{enumerate}
\begin{tcolorbox}[enhanced,breakable,colback=teal!4,colframe=teal!50!black,boxrule=0.5pt,arc=1.1mm,left=3mm,right=3mm,top=1.7mm,bottom=1.7mm,before skip=2mm,after skip=2mm]
\textbf{\color{teal}To'g'ri javob:} \textbf{D) Bandwidth}
\par\smallskip
\textbf{Izoh:} Vaziyat Bandwidthga misol: kanalning nazariy yoki belgilangan maksimal o'tkazish quvvati.
\end{tcolorbox}
\vspace{2mm}\hrule height 0.2pt \vspace{3mm}

% DB question_id: ec73ce57-a668-4c56-8fb3-7686688bb04e
\needspace{5\baselineskip}
\noindent\textbf{\color{ink}986.}
\par
Quyidagi ta'rif qaysi atamaga tegishli? Bir soniyada uzatiladigan signal simvollari soni.
\begin{enumerate}[label=\Alph*),leftmargin=9mm]
\item Simpleks aloqa
\item Baud
\item Yarim dupleks
\item To'liq dupleks
\end{enumerate}
\begin{tcolorbox}[enhanced,breakable,colback=teal!4,colframe=teal!50!black,boxrule=0.5pt,arc=1.1mm,left=3mm,right=3mm,top=1.7mm,bottom=1.7mm,before skip=2mm,after skip=2mm]
\textbf{\color{teal}To'g'ri javob:} \textbf{B) Baud}
\par\smallskip
\textbf{Izoh:} Ta'rif Baud tushunchasiga tegishli.
\end{tcolorbox}
\vspace{2mm}\hrule height 0.2pt \vspace{3mm}

% DB question_id: f6c72aef-3263-496f-8699-21c2f8daa25e
\needspace{5\baselineskip}
\noindent\textbf{\color{ink}987.}
\par
Quyidagi ta'rif qaysi atamaga tegishli? Foydalanuvchi amalda oladigan foydali ma'lumot uzatish tezligi.
\begin{enumerate}[label=\Alph*),leftmargin=9mm]
\item Throughput
\item Baud
\item Simpleks aloqa
\item Yarim dupleks
\end{enumerate}
\begin{tcolorbox}[enhanced,breakable,colback=teal!4,colframe=teal!50!black,boxrule=0.5pt,arc=1.1mm,left=3mm,right=3mm,top=1.7mm,bottom=1.7mm,before skip=2mm,after skip=2mm]
\textbf{\color{teal}To'g'ri javob:} \textbf{A) Throughput}
\par\smallskip
\textbf{Izoh:} Ta'rif Throughput tushunchasiga tegishli.
\end{tcolorbox}
\vspace{2mm}\hrule height 0.2pt \vspace{3mm}

% DB question_id: 01e5a3b9-0b8a-5047-8d37-0efc973dac92
\needspace{5\baselineskip}
\noindent\textbf{\color{ink}988.}
\par
10179 MB fayl 216 Mbit/s tezlikda uzatilsa, ideal sharoitda uzatish qancha vaqt oladi?
\begin{enumerate}[label=\Alph*),leftmargin=9mm]
\item 3016 soniya
\item 377 soniya
\item 593 soniya
\item 188 soniya
\end{enumerate}
\begin{tcolorbox}[enhanced,breakable,colback=teal!4,colframe=teal!50!black,boxrule=0.5pt,arc=1.1mm,left=3mm,right=3mm,top=1.7mm,bottom=1.7mm,before skip=2mm,after skip=2mm]
\textbf{\color{teal}To'g'ri javob:} \textbf{B) 377 soniya}
\par\smallskip
\textbf{Izoh:} 10179 MB = 81432 Mbit. Vaqt = hajm / tezlik = 81432 / 216 = 377 soniya.
\end{tcolorbox}
\vspace{2mm}\hrule height 0.2pt \vspace{3mm}

% DB question_id: 03563cda-ce34-5d48-9379-8728bb23b77a
\needspace{5\baselineskip}
\noindent\textbf{\color{ink}989.}
\par
22032 MB fayl 432 Mbit/s tezlikda uzatilsa, ideal sharoitda uzatish qancha vaqt oladi?
\begin{enumerate}[label=\Alph*),leftmargin=9mm]
\item 408 soniya
\item 840 soniya
\item 3264 soniya
\item 204 soniya
\end{enumerate}
\begin{tcolorbox}[enhanced,breakable,colback=teal!4,colframe=teal!50!black,boxrule=0.5pt,arc=1.1mm,left=3mm,right=3mm,top=1.7mm,bottom=1.7mm,before skip=2mm,after skip=2mm]
\textbf{\color{teal}To'g'ri javob:} \textbf{A) 408 soniya}
\par\smallskip
\textbf{Izoh:} 22032 MB = 176256 Mbit. Vaqt = hajm / tezlik = 176256 / 432 = 408 soniya.
\end{tcolorbox}
\vspace{2mm}\hrule height 0.2pt \vspace{3mm}

% DB question_id: 068b996a-bcf8-5a59-939c-ac33e0e245a8
\needspace{5\baselineskip}
\noindent\textbf{\color{ink}990.}
\par
19084 MB fayl 416 Mbit/s tezlikda uzatilsa, ideal sharoitda uzatish qancha vaqt oladi?
\begin{enumerate}[label=\Alph*),leftmargin=9mm]
\item 183 soniya
\item 2936 soniya
\item 783 soniya
\item 367 soniya
\end{enumerate}
\begin{tcolorbox}[enhanced,breakable,colback=teal!4,colframe=teal!50!black,boxrule=0.5pt,arc=1.1mm,left=3mm,right=3mm,top=1.7mm,bottom=1.7mm,before skip=2mm,after skip=2mm]
\textbf{\color{teal}To'g'ri javob:} \textbf{D) 367 soniya}
\par\smallskip
\textbf{Izoh:} 19084 MB = 152672 Mbit. Vaqt = hajm / tezlik = 152672 / 416 = 367 soniya.
\end{tcolorbox}
\vspace{2mm}\hrule height 0.2pt \vspace{3mm}

% DB question_id: 08277968-3cc8-570e-82e8-38b25821017f
\needspace{5\baselineskip}
\noindent\textbf{\color{ink}991.}
\par
1836 MB fayl 432 Mbit/s tezlikda uzatilsa, ideal sharoitda uzatish qancha vaqt oladi?
\begin{enumerate}[label=\Alph*),leftmargin=9mm]
\item 34 soniya
\item 272 soniya
\item 17 soniya
\item 466 soniya
\end{enumerate}
\begin{tcolorbox}[enhanced,breakable,colback=teal!4,colframe=teal!50!black,boxrule=0.5pt,arc=1.1mm,left=3mm,right=3mm,top=1.7mm,bottom=1.7mm,before skip=2mm,after skip=2mm]
\textbf{\color{teal}To'g'ri javob:} \textbf{A) 34 soniya}
\par\smallskip
\textbf{Izoh:} 1836 MB = 14688 Mbit. Vaqt = hajm / tezlik = 14688 / 432 = 34 soniya.
\end{tcolorbox}
\vspace{2mm}\hrule height 0.2pt \vspace{3mm}

% DB question_id: 095b7d15-ecd0-5869-9b62-8b57c2d335c7
\needspace{5\baselineskip}
\noindent\textbf{\color{ink}992.}
\par
31388 MB fayl 472 Mbit/s tezlikda uzatilsa, ideal sharoitda uzatish qancha vaqt oladi?
\begin{enumerate}[label=\Alph*),leftmargin=9mm]
\item 266 soniya
\item 532 soniya
\item 4256 soniya
\item 1004 soniya
\end{enumerate}
\begin{tcolorbox}[enhanced,breakable,colback=teal!4,colframe=teal!50!black,boxrule=0.5pt,arc=1.1mm,left=3mm,right=3mm,top=1.7mm,bottom=1.7mm,before skip=2mm,after skip=2mm]
\textbf{\color{teal}To'g'ri javob:} \textbf{B) 532 soniya}
\par\smallskip
\textbf{Izoh:} 31388 MB = 251104 Mbit. Vaqt = hajm / tezlik = 251104 / 472 = 532 soniya.
\end{tcolorbox}
\vspace{2mm}\hrule height 0.2pt \vspace{3mm}

% DB question_id: 1558990b-04e5-574a-85da-5f03619a3c05
\needspace{5\baselineskip}
\noindent\textbf{\color{ink}993.}
\par
13202 MB fayl 328 Mbit/s tezlikda uzatilsa, ideal sharoitda uzatish qancha vaqt oladi?
\begin{enumerate}[label=\Alph*),leftmargin=9mm]
\item 161 soniya
\item 2576 soniya
\item 650 soniya
\item 322 soniya
\end{enumerate}
\begin{tcolorbox}[enhanced,breakable,colback=teal!4,colframe=teal!50!black,boxrule=0.5pt,arc=1.1mm,left=3mm,right=3mm,top=1.7mm,bottom=1.7mm,before skip=2mm,after skip=2mm]
\textbf{\color{teal}To'g'ri javob:} \textbf{D) 322 soniya}
\par\smallskip
\textbf{Izoh:} 13202 MB = 105616 Mbit. Vaqt = hajm / tezlik = 105616 / 328 = 322 soniya.
\end{tcolorbox}
\vspace{2mm}\hrule height 0.2pt \vspace{3mm}

% DB question_id: 1752699f-24d5-5d7d-ae53-16d304d587a8
\needspace{5\baselineskip}
\noindent\textbf{\color{ink}994.}
\par
121 MB fayl 8 Mbit/s tezlikda uzatilsa, ideal sharoitda uzatish qancha vaqt oladi?
\begin{enumerate}[label=\Alph*),leftmargin=9mm]
\item 968 soniya
\item 121 soniya
\item 129 soniya
\item 60 soniya
\end{enumerate}
\begin{tcolorbox}[enhanced,breakable,colback=teal!4,colframe=teal!50!black,boxrule=0.5pt,arc=1.1mm,left=3mm,right=3mm,top=1.7mm,bottom=1.7mm,before skip=2mm,after skip=2mm]
\textbf{\color{teal}To'g'ri javob:} \textbf{B) 121 soniya}
\par\smallskip
\textbf{Izoh:} 121 MB = 968 Mbit. Vaqt = hajm / tezlik = 968 / 8 = 121 soniya.
\end{tcolorbox}
\vspace{2mm}\hrule height 0.2pt \vspace{3mm}

% DB question_id: 1981b70e-84ed-5b14-9e89-e7a1217cae51
\needspace{5\baselineskip}
\noindent\textbf{\color{ink}995.}
\par
12122 MB fayl 232 Mbit/s tezlikda uzatilsa, ideal sharoitda uzatish qancha vaqt oladi?
\begin{enumerate}[label=\Alph*),leftmargin=9mm]
\item 209 soniya
\item 650 soniya
\item 418 soniya
\item 3344 soniya
\end{enumerate}
\begin{tcolorbox}[enhanced,breakable,colback=teal!4,colframe=teal!50!black,boxrule=0.5pt,arc=1.1mm,left=3mm,right=3mm,top=1.7mm,bottom=1.7mm,before skip=2mm,after skip=2mm]
\textbf{\color{teal}To'g'ri javob:} \textbf{C) 418 soniya}
\par\smallskip
\textbf{Izoh:} 12122 MB = 96976 Mbit. Vaqt = hajm / tezlik = 96976 / 232 = 418 soniya.
\end{tcolorbox}
\vspace{2mm}\hrule height 0.2pt \vspace{3mm}

% DB question_id: 1aaaccb0-15a3-53f6-a707-b5f6e1b13838
\needspace{5\baselineskip}
\noindent\textbf{\color{ink}996.}
\par
486 MB fayl 24 Mbit/s tezlikda uzatilsa, ideal sharoitda uzatish qancha vaqt oladi?
\begin{enumerate}[label=\Alph*),leftmargin=9mm]
\item 81 soniya
\item 186 soniya
\item 162 soniya
\item 1296 soniya
\end{enumerate}
\begin{tcolorbox}[enhanced,breakable,colback=teal!4,colframe=teal!50!black,boxrule=0.5pt,arc=1.1mm,left=3mm,right=3mm,top=1.7mm,bottom=1.7mm,before skip=2mm,after skip=2mm]
\textbf{\color{teal}To'g'ri javob:} \textbf{C) 162 soniya}
\par\smallskip
\textbf{Izoh:} 486 MB = 3888 Mbit. Vaqt = hajm / tezlik = 3888 / 24 = 162 soniya.
\end{tcolorbox}
\vspace{2mm}\hrule height 0.2pt \vspace{3mm}

% DB question_id: 1d11477d-e403-4169-99cd-7727ff62c803
\needspace{5\baselineskip}
\noindent\textbf{\color{ink}997.}
\par
160 betli kitobda har betda 48 satr, har satrda 64 belgi bor. Uni o'qishga 21 soat 20 daqiqa ketdi. Bir belgi 1 bayt deb olinsa, o'rtacha o'qish tezligi necha bayt/sekund?
\begin{enumerate}[label=\Alph*),leftmargin=9mm]
\item 7,2
\item 6,4
\item 5,5
\item 8
\end{enumerate}
\begin{tcolorbox}[enhanced,breakable,colback=teal!4,colframe=teal!50!black,boxrule=0.5pt,arc=1.1mm,left=3mm,right=3mm,top=1.7mm,bottom=1.7mm,before skip=2mm,after skip=2mm]
\textbf{\color{teal}To'g'ri javob:} \textbf{B) 6,4}
\par\smallskip
\textbf{Izoh:} 160\ensuremath{\times}48\ensuremath{\times}64=491520 bayt; 21\ensuremath{\times}3600+20\ensuremath{\times}60=76800 sekund; 491520/76800=6,4. [2025 kuzgi demo YECHIMLARI, sahifa 13.]
\end{tcolorbox}
\vspace{2mm}\hrule height 0.2pt \vspace{3mm}

% DB question_id: 21fc7868-3f28-5f04-9944-cad40edaf8bd
\needspace{5\baselineskip}
\noindent\textbf{\color{ink}998.}
\par
375 MB fayl 120 Mbit/s tezlikda uzatilsa, ideal sharoitda uzatish qancha vaqt oladi?
\begin{enumerate}[label=\Alph*),leftmargin=9mm]
\item 12 soniya
\item 200 soniya
\item 145 soniya
\item 25 soniya
\end{enumerate}
\begin{tcolorbox}[enhanced,breakable,colback=teal!4,colframe=teal!50!black,boxrule=0.5pt,arc=1.1mm,left=3mm,right=3mm,top=1.7mm,bottom=1.7mm,before skip=2mm,after skip=2mm]
\textbf{\color{teal}To'g'ri javob:} \textbf{D) 25 soniya}
\par\smallskip
\textbf{Izoh:} 375 MB = 3000 Mbit. Vaqt = hajm / tezlik = 3000 / 120 = 25 soniya.
\end{tcolorbox}
\vspace{2mm}\hrule height 0.2pt \vspace{3mm}

% DB question_id: 24ec3dbc-be4b-5ccf-b985-fa047c633b9c
\needspace{5\baselineskip}
\noindent\textbf{\color{ink}999.}
\par
2831 MB fayl 152 Mbit/s tezlikda uzatilsa, ideal sharoitda uzatish qancha vaqt oladi?
\begin{enumerate}[label=\Alph*),leftmargin=9mm]
\item 301 soniya
\item 149 soniya
\item 74 soniya
\item 1192 soniya
\end{enumerate}
\begin{tcolorbox}[enhanced,breakable,colback=teal!4,colframe=teal!50!black,boxrule=0.5pt,arc=1.1mm,left=3mm,right=3mm,top=1.7mm,bottom=1.7mm,before skip=2mm,after skip=2mm]
\textbf{\color{teal}To'g'ri javob:} \textbf{B) 149 soniya}
\par\smallskip
\textbf{Izoh:} 2831 MB = 22648 Mbit. Vaqt = hajm / tezlik = 22648 / 152 = 149 soniya.
\end{tcolorbox}
\vspace{2mm}\hrule height 0.2pt \vspace{3mm}

% DB question_id: 27b1241a-be82-540d-878e-51f0c8122081
\needspace{5\baselineskip}
\noindent\textbf{\color{ink}1000.}
\par
24750 MB fayl 440 Mbit/s tezlikda uzatilsa, ideal sharoitda uzatish qancha vaqt oladi?
\begin{enumerate}[label=\Alph*),leftmargin=9mm]
\item 450 soniya
\item 225 soniya
\item 3600 soniya
\item 890 soniya
\end{enumerate}
\begin{tcolorbox}[enhanced,breakable,colback=teal!4,colframe=teal!50!black,boxrule=0.5pt,arc=1.1mm,left=3mm,right=3mm,top=1.7mm,bottom=1.7mm,before skip=2mm,after skip=2mm]
\textbf{\color{teal}To'g'ri javob:} \textbf{A) 450 soniya}
\par\smallskip
\textbf{Izoh:} 24750 MB = 198000 Mbit. Vaqt = hajm / tezlik = 198000 / 440 = 450 soniya.
\end{tcolorbox}
\vspace{2mm}\hrule height 0.2pt \vspace{3mm}

% DB question_id: 363bc602-18ff-5e05-ac47-a74352d173e3
\needspace{5\baselineskip}
\noindent\textbf{\color{ink}1001.}
\par
4256 MB fayl 224 Mbit/s tezlikda uzatilsa, ideal sharoitda uzatish qancha vaqt oladi?
\begin{enumerate}[label=\Alph*),leftmargin=9mm]
\item 152 soniya
\item 1216 soniya
\item 376 soniya
\item 76 soniya
\end{enumerate}
\begin{tcolorbox}[enhanced,breakable,colback=teal!4,colframe=teal!50!black,boxrule=0.5pt,arc=1.1mm,left=3mm,right=3mm,top=1.7mm,bottom=1.7mm,before skip=2mm,after skip=2mm]
\textbf{\color{teal}To'g'ri javob:} \textbf{A) 152 soniya}
\par\smallskip
\textbf{Izoh:} 4256 MB = 34048 Mbit. Vaqt = hajm / tezlik = 34048 / 224 = 152 soniya.
\end{tcolorbox}
\vspace{2mm}\hrule height 0.2pt \vspace{3mm}

% DB question_id: 456e17d4-81e1-4362-9c25-37ab9421fe14
\needspace{5\baselineskip}
\noindent\textbf{\color{ink}1002.}
\par
Fayl hajmi va uzatish tezligidan uzatish vaqtini topishda bajariladigan amallarni ish tartibida joylashtiring. Hisoblashga tayyorgarlikdan yakuniy natijani tekshirishgacha bo'lgan izchillikni aniqlang.
\begin{enumerate}[label=\Alph*),leftmargin=9mm]
\item t = I / v formulasi bo'yicha vaqtni hisoblash
\item Fayl hajmi va tezlik birliklarini aniqlash
\item Natija birligi va mantiqiyligini tekshirish
\item Fayl hajmini tezlik bilan bir xil bit birliklariga keltirish
\end{enumerate}
\begin{tcolorbox}[enhanced,breakable,colback=teal!4,colframe=teal!50!black,boxrule=0.5pt,arc=1.1mm,left=3mm,right=3mm,top=1.7mm,bottom=1.7mm,before skip=2mm,after skip=2mm]
\textbf{\color{teal}To'g'ri javob:} \textbf{B -- D -- A -- C}
\par\smallskip
\textbf{Izoh:} Avval birliklar aniqlanadi va moslashtiriladi, keyin formula qo'llanib natija tekshiriladi.
\end{tcolorbox}
\vspace{2mm}\hrule height 0.2pt \vspace{3mm}

% DB question_id: 46e71e39-f0a4-526b-83a5-5689c2e32ce5
\needspace{5\baselineskip}
\noindent\textbf{\color{ink}1003.}
\par
496 MB fayl 64 Mbit/s tezlikda uzatilsa, ideal sharoitda uzatish qancha vaqt oladi?
\begin{enumerate}[label=\Alph*),leftmargin=9mm]
\item 126 soniya
\item 62 soniya
\item 31 soniya
\item 496 soniya
\end{enumerate}
\begin{tcolorbox}[enhanced,breakable,colback=teal!4,colframe=teal!50!black,boxrule=0.5pt,arc=1.1mm,left=3mm,right=3mm,top=1.7mm,bottom=1.7mm,before skip=2mm,after skip=2mm]
\textbf{\color{teal}To'g'ri javob:} \textbf{B) 62 soniya}
\par\smallskip
\textbf{Izoh:} 496 MB = 3968 Mbit. Vaqt = hajm / tezlik = 3968 / 64 = 62 soniya.
\end{tcolorbox}
\vspace{2mm}\hrule height 0.2pt \vspace{3mm}

% DB question_id: 499c6376-43a0-5969-84b1-1da954eb40af
\needspace{5\baselineskip}
\noindent\textbf{\color{ink}1004.}
\par
18630 MB fayl 368 Mbit/s tezlikda uzatilsa, ideal sharoitda uzatish qancha vaqt oladi?
\begin{enumerate}[label=\Alph*),leftmargin=9mm]
\item 773 soniya
\item 405 soniya
\item 202 soniya
\item 3240 soniya
\end{enumerate}
\begin{tcolorbox}[enhanced,breakable,colback=teal!4,colframe=teal!50!black,boxrule=0.5pt,arc=1.1mm,left=3mm,right=3mm,top=1.7mm,bottom=1.7mm,before skip=2mm,after skip=2mm]
\textbf{\color{teal}To'g'ri javob:} \textbf{B) 405 soniya}
\par\smallskip
\textbf{Izoh:} 18630 MB = 149040 Mbit. Vaqt = hajm / tezlik = 149040 / 368 = 405 soniya.
\end{tcolorbox}
\vspace{2mm}\hrule height 0.2pt \vspace{3mm}

% DB question_id: 5902ac4e-1738-57f9-a77f-c216a9cb8271
\needspace{5\baselineskip}
\noindent\textbf{\color{ink}1005.}
\par
9141 MB fayl 264 Mbit/s tezlikda uzatilsa, ideal sharoitda uzatish qancha vaqt oladi?
\begin{enumerate}[label=\Alph*),leftmargin=9mm]
\item 138 soniya
\item 277 soniya
\item 2216 soniya
\item 541 soniya
\end{enumerate}
\begin{tcolorbox}[enhanced,breakable,colback=teal!4,colframe=teal!50!black,boxrule=0.5pt,arc=1.1mm,left=3mm,right=3mm,top=1.7mm,bottom=1.7mm,before skip=2mm,after skip=2mm]
\textbf{\color{teal}To'g'ri javob:} \textbf{B) 277 soniya}
\par\smallskip
\textbf{Izoh:} 9141 MB = 73128 Mbit. Vaqt = hajm / tezlik = 73128 / 264 = 277 soniya.
\end{tcolorbox}
\vspace{2mm}\hrule height 0.2pt \vspace{3mm}

% DB question_id: 6c281a46-0317-534c-ae80-4d99f4d772ba
\needspace{5\baselineskip}
\noindent\textbf{\color{ink}1006.}
\par
60 MB faylning birinchi yarmi 256 Kbit/s, ikkinchi yarmi esa 128 Kbit/s tezlikda yuklandi. \texttt{1 MB=1024 KB} va \texttt{1 bayt=8 bit} deb oling. Fayl necha daqiqada yuklanadi?
\begin{enumerate}[label=\Alph*),leftmargin=9mm]
\item 32 daqiqa
\item 40 daqiqa
\item 48 daqiqa
\item 55 daqiqa
\end{enumerate}
\begin{tcolorbox}[enhanced,breakable,colback=teal!4,colframe=teal!50!black,boxrule=0.5pt,arc=1.1mm,left=3mm,right=3mm,top=1.7mm,bottom=1.7mm,before skip=2mm,after skip=2mm]
\textbf{\color{teal}To'g'ri javob:} \textbf{C) 48 daqiqa}
\par\smallskip
\textbf{Izoh:} Har bir yarmi \texttt{30\ensuremath{\cdot}1024\ensuremath{\cdot}8=245760 Kbit}. Birinchi qism \texttt{245760/256=960 s=16 min}, ikkinchisi \texttt{245760/128=1920 s=32 min}. Jami 48 daqiqa. Asl PDFdagi 47 varianti noto'g'ri edi.
\end{tcolorbox}
\vspace{2mm}\hrule height 0.2pt \vspace{3mm}

% DB question_id: 6d831df0-23fc-5eea-9c35-db59cdd43c7a
\needspace{5\baselineskip}
\noindent\textbf{\color{ink}1007.}
\par
30141 MB fayl 408 Mbit/s tezlikda uzatilsa, ideal sharoitda uzatish qancha vaqt oladi?
\begin{enumerate}[label=\Alph*),leftmargin=9mm]
\item 4728 soniya
\item 295 soniya
\item 999 soniya
\item 591 soniya
\end{enumerate}
\begin{tcolorbox}[enhanced,breakable,colback=teal!4,colframe=teal!50!black,boxrule=0.5pt,arc=1.1mm,left=3mm,right=3mm,top=1.7mm,bottom=1.7mm,before skip=2mm,after skip=2mm]
\textbf{\color{teal}To'g'ri javob:} \textbf{D) 591 soniya}
\par\smallskip
\textbf{Izoh:} 30141 MB = 241128 Mbit. Vaqt = hajm / tezlik = 241128 / 408 = 591 soniya.
\end{tcolorbox}
\vspace{2mm}\hrule height 0.2pt \vspace{3mm}

% DB question_id: 6e848bcb-1b39-5472-bb44-36c57c70bd4d
\needspace{5\baselineskip}
\noindent\textbf{\color{ink}1008.}
\par
3486 MB fayl 112 Mbit/s tezlikda uzatilsa, ideal sharoitda uzatish qancha vaqt oladi?
\begin{enumerate}[label=\Alph*),leftmargin=9mm]
\item 1992 soniya
\item 124 soniya
\item 361 soniya
\item 249 soniya
\end{enumerate}
\begin{tcolorbox}[enhanced,breakable,colback=teal!4,colframe=teal!50!black,boxrule=0.5pt,arc=1.1mm,left=3mm,right=3mm,top=1.7mm,bottom=1.7mm,before skip=2mm,after skip=2mm]
\textbf{\color{teal}To'g'ri javob:} \textbf{D) 249 soniya}
\par\smallskip
\textbf{Izoh:} 3486 MB = 27888 Mbit. Vaqt = hajm / tezlik = 27888 / 112 = 249 soniya.
\end{tcolorbox}
\vspace{2mm}\hrule height 0.2pt \vspace{3mm}

% DB question_id: 79e1e625-4ab4-5621-9234-3c05fb80c81d
\needspace{5\baselineskip}
\noindent\textbf{\color{ink}1009.}
\par
126 MB fayl 48 Mbit/s tezlikda uzatilsa, ideal sharoitda uzatish qancha vaqt oladi?
\begin{enumerate}[label=\Alph*),leftmargin=9mm]
\item 10 soniya
\item 21 soniya
\item 168 soniya
\item 69 soniya
\end{enumerate}
\begin{tcolorbox}[enhanced,breakable,colback=teal!4,colframe=teal!50!black,boxrule=0.5pt,arc=1.1mm,left=3mm,right=3mm,top=1.7mm,bottom=1.7mm,before skip=2mm,after skip=2mm]
\textbf{\color{teal}To'g'ri javob:} \textbf{B) 21 soniya}
\par\smallskip
\textbf{Izoh:} 126 MB = 1008 Mbit. Vaqt = hajm / tezlik = 1008 / 48 = 21 soniya.
\end{tcolorbox}
\vspace{2mm}\hrule height 0.2pt \vspace{3mm}

% DB question_id: 7f2a46e2-60dc-46de-96dc-dff1389f5567
\needspace{5\baselineskip}
\noindent\textbf{\color{ink}1010.}
\par
2400 baud signalda har bir simvol 4 bit tashisa, bit tezligi qancha?
\begin{enumerate}[label=\Alph*),leftmargin=9mm]
\item 600 bit/s
\item 9600 bit/s
\item 2400 bit/s
\item 4800 bit/s
\end{enumerate}
\begin{tcolorbox}[enhanced,breakable,colback=teal!4,colframe=teal!50!black,boxrule=0.5pt,arc=1.1mm,left=3mm,right=3mm,top=1.7mm,bottom=1.7mm,before skip=2mm,after skip=2mm]
\textbf{\color{teal}To'g'ri javob:} \textbf{B) 9600 bit/s}
\par\smallskip
\textbf{Izoh:} 2400 simvol/s \ensuremath{\times} 4 bit/simvol = 9600 bit/s.
\end{tcolorbox}
\vspace{2mm}\hrule height 0.2pt \vspace{3mm}

% DB question_id: 965ec9e1-7dde-53eb-8e38-e6e34257b47d
\needspace{5\baselineskip}
\noindent\textbf{\color{ink}1011.}
\par
2144 MB fayl 32 Mbit/s tezlikda uzatilsa, ideal sharoitda uzatish qancha vaqt oladi?
\begin{enumerate}[label=\Alph*),leftmargin=9mm]
\item 268 soniya
\item 4288 soniya
\item 568 soniya
\item 536 soniya
\end{enumerate}
\begin{tcolorbox}[enhanced,breakable,colback=teal!4,colframe=teal!50!black,boxrule=0.5pt,arc=1.1mm,left=3mm,right=3mm,top=1.7mm,bottom=1.7mm,before skip=2mm,after skip=2mm]
\textbf{\color{teal}To'g'ri javob:} \textbf{D) 536 soniya}
\par\smallskip
\textbf{Izoh:} 2144 MB = 17152 Mbit. Vaqt = hajm / tezlik = 17152 / 32 = 536 soniya.
\end{tcolorbox}
\vspace{2mm}\hrule height 0.2pt \vspace{3mm}

% DB question_id: 9d4a4a1c-87c5-5a30-884e-4ed9d22defbb
\needspace{5\baselineskip}
\noindent\textbf{\color{ink}1012.}
\par
4180 MB fayl 176 Mbit/s tezlikda uzatilsa, ideal sharoitda uzatish qancha vaqt oladi?
\begin{enumerate}[label=\Alph*),leftmargin=9mm]
\item 190 soniya
\item 95 soniya
\item 1520 soniya
\item 366 soniya
\end{enumerate}
\begin{tcolorbox}[enhanced,breakable,colback=teal!4,colframe=teal!50!black,boxrule=0.5pt,arc=1.1mm,left=3mm,right=3mm,top=1.7mm,bottom=1.7mm,before skip=2mm,after skip=2mm]
\textbf{\color{teal}To'g'ri javob:} \textbf{A) 190 soniya}
\par\smallskip
\textbf{Izoh:} 4180 MB = 33440 Mbit. Vaqt = hajm / tezlik = 33440 / 176 = 190 soniya.
\end{tcolorbox}
\vspace{2mm}\hrule height 0.2pt \vspace{3mm}

% DB question_id: a1c93aeb-f734-54c0-90f1-7d7f71a678d9
\needspace{5\baselineskip}
\noindent\textbf{\color{ink}1013.}
\par
27987 MB fayl 456 Mbit/s tezlikda uzatilsa, ideal sharoitda uzatish qancha vaqt oladi?
\begin{enumerate}[label=\Alph*),leftmargin=9mm]
\item 947 soniya
\item 491 soniya
\item 245 soniya
\item 3928 soniya
\end{enumerate}
\begin{tcolorbox}[enhanced,breakable,colback=teal!4,colframe=teal!50!black,boxrule=0.5pt,arc=1.1mm,left=3mm,right=3mm,top=1.7mm,bottom=1.7mm,before skip=2mm,after skip=2mm]
\textbf{\color{teal}To'g'ri javob:} \textbf{B) 491 soniya}
\par\smallskip
\textbf{Izoh:} 27987 MB = 223896 Mbit. Vaqt = hajm / tezlik = 223896 / 456 = 491 soniya.
\end{tcolorbox}
\vspace{2mm}\hrule height 0.2pt \vspace{3mm}

% DB question_id: a7e52c35-9bae-5619-a10a-6ddd607c465b
\needspace{5\baselineskip}
\noindent\textbf{\color{ink}1014.}
\par
16335 MB fayl 360 Mbit/s tezlikda uzatilsa, ideal sharoitda uzatish qancha vaqt oladi?
\begin{enumerate}[label=\Alph*),leftmargin=9mm]
\item 723 soniya
\item 181 soniya
\item 363 soniya
\item 2904 soniya
\end{enumerate}
\begin{tcolorbox}[enhanced,breakable,colback=teal!4,colframe=teal!50!black,boxrule=0.5pt,arc=1.1mm,left=3mm,right=3mm,top=1.7mm,bottom=1.7mm,before skip=2mm,after skip=2mm]
\textbf{\color{teal}To'g'ri javob:} \textbf{C) 363 soniya}
\par\smallskip
\textbf{Izoh:} 16335 MB = 130680 Mbit. Vaqt = hajm / tezlik = 130680 / 360 = 363 soniya.
\end{tcolorbox}
\vspace{2mm}\hrule height 0.2pt \vspace{3mm}

% DB question_id: b7242ba2-25fa-593b-9d4b-4ed4f7722bba
\needspace{5\baselineskip}
\noindent\textbf{\color{ink}1015.}
\par
20200 MB fayl 320 Mbit/s tezlikda uzatilsa, ideal sharoitda uzatish qancha vaqt oladi?
\begin{enumerate}[label=\Alph*),leftmargin=9mm]
\item 252 soniya
\item 4040 soniya
\item 825 soniya
\item 505 soniya
\end{enumerate}
\begin{tcolorbox}[enhanced,breakable,colback=teal!4,colframe=teal!50!black,boxrule=0.5pt,arc=1.1mm,left=3mm,right=3mm,top=1.7mm,bottom=1.7mm,before skip=2mm,after skip=2mm]
\textbf{\color{teal}To'g'ri javob:} \textbf{D) 505 soniya}
\par\smallskip
\textbf{Izoh:} 20200 MB = 161600 Mbit. Vaqt = hajm / tezlik = 161600 / 320 = 505 soniya.
\end{tcolorbox}
\vspace{2mm}\hrule height 0.2pt \vspace{3mm}

% DB question_id: ba605d77-81e5-5b3f-82ee-ce1915821068
\needspace{5\baselineskip}
\noindent\textbf{\color{ink}1016.}
\par
4640 MB fayl 128 Mbit/s tezlikda uzatilsa, ideal sharoitda uzatish qancha vaqt oladi?
\begin{enumerate}[label=\Alph*),leftmargin=9mm]
\item 290 soniya
\item 145 soniya
\item 2320 soniya
\item 418 soniya
\end{enumerate}
\begin{tcolorbox}[enhanced,breakable,colback=teal!4,colframe=teal!50!black,boxrule=0.5pt,arc=1.1mm,left=3mm,right=3mm,top=1.7mm,bottom=1.7mm,before skip=2mm,after skip=2mm]
\textbf{\color{teal}To'g'ri javob:} \textbf{A) 290 soniya}
\par\smallskip
\textbf{Izoh:} 4640 MB = 37120 Mbit. Vaqt = hajm / tezlik = 37120 / 128 = 290 soniya.
\end{tcolorbox}
\vspace{2mm}\hrule height 0.2pt \vspace{3mm}

% DB question_id: c6f9b0e2-2e5d-4713-8930-fdd71e5d44ec
\needspace{5\baselineskip}
\noindent\textbf{\color{ink}1017.}
\par
Telefon 20 daqiqada 400 KiB axborot oldi. 1 KiB=1024 bayt va 1 bayt=8 bit deb oling. O'rtacha tezlikni eng yaqin butun bit/sekundga yaxlitlang.
\begin{enumerate}[label=\Alph*),leftmargin=9mm]
\item 273 bit/s
\item 2730 bit/s
\item 2731 bit/s
\item 4096 bit/s
\end{enumerate}
\begin{tcolorbox}[enhanced,breakable,colback=teal!4,colframe=teal!50!black,boxrule=0.5pt,arc=1.1mm,left=3mm,right=3mm,top=1.7mm,bottom=1.7mm,before skip=2mm,after skip=2mm]
\textbf{\color{teal}To'g'ri javob:} \textbf{C) 2731 bit/s}
\par\smallskip
\textbf{Izoh:} 20 minut=1200 sekund. 400 KiB=400\ensuremath{\times}1024\ensuremath{\times}8=3276800 bit; 3276800/1200\ensuremath{\approx}2730,67 bit/s. Eng yaqin butun qiymat 2731 bit/s.
\end{tcolorbox}
\vspace{2mm}\hrule height 0.2pt \vspace{3mm}

% DB question_id: cc44b02b-9cd7-44fe-bc01-3af644ae8783
\needspace{5\baselineskip}
\noindent\textbf{\color{ink}1018.}
\par
25 MB fayl 10 Mbit/s tezlikda ideal holatda qancha vaqtda uzatiladi? MB = 10\textsuperscript{6} B.
\begin{enumerate}[label=\Alph*),leftmargin=9mm]
\item 20 s
\item 2,5 s
\item 10 s
\item 200 s
\end{enumerate}
\begin{tcolorbox}[enhanced,breakable,colback=teal!4,colframe=teal!50!black,boxrule=0.5pt,arc=1.1mm,left=3mm,right=3mm,top=1.7mm,bottom=1.7mm,before skip=2mm,after skip=2mm]
\textbf{\color{teal}To'g'ri javob:} \textbf{A) 20 s}
\par\smallskip
\textbf{Izoh:} 25 MB = 200 Mbit; 200 \ensuremath{\div} 10 = 20 soniya.
\end{tcolorbox}
\vspace{2mm}\hrule height 0.2pt \vspace{3mm}

% DB question_id: d040651f-80f2-5757-949d-d0bb78dea54a
\needspace{5\baselineskip}
\noindent\textbf{\color{ink}1019.}
\par
495 MB fayl 8 Mbit/s tezlikda uzatilsa, ideal sharoitda uzatish qancha vaqt oladi?
\begin{enumerate}[label=\Alph*),leftmargin=9mm]
\item 495 soniya
\item 3960 soniya
\item 503 soniya
\item 247 soniya
\end{enumerate}
\begin{tcolorbox}[enhanced,breakable,colback=teal!4,colframe=teal!50!black,boxrule=0.5pt,arc=1.1mm,left=3mm,right=3mm,top=1.7mm,bottom=1.7mm,before skip=2mm,after skip=2mm]
\textbf{\color{teal}To'g'ri javob:} \textbf{A) 495 soniya}
\par\smallskip
\textbf{Izoh:} 495 MB = 3960 Mbit. Vaqt = hajm / tezlik = 3960 / 8 = 495 soniya.
\end{tcolorbox}
\vspace{2mm}\hrule height 0.2pt \vspace{3mm}

% DB question_id: df4811aa-c0c3-47ca-b7f9-c7378ea1c282
\needspace{5\baselineskip}
\noindent\textbf{\color{ink}1020.}
\par
Modem bir soniyada 2400 ta signal holati yubordi. Vaziyatni tahlil qilib, unga eng aniq mos tushunchani toping.
\begin{enumerate}[label=\Alph*),leftmargin=9mm]
\item Simpleks aloqa
\item Yarim dupleks
\item Baud
\item To'liq dupleks
\end{enumerate}
\begin{tcolorbox}[enhanced,breakable,colback=teal!4,colframe=teal!50!black,boxrule=0.5pt,arc=1.1mm,left=3mm,right=3mm,top=1.7mm,bottom=1.7mm,before skip=2mm,after skip=2mm]
\textbf{\color{teal}To'g'ri javob:} \textbf{C) Baud}
\par\smallskip
\textbf{Izoh:} Vaziyat Baudga misol: bir soniyada uzatiladigan signal simvollari soni.
\end{tcolorbox}
\vspace{2mm}\hrule height 0.2pt \vspace{3mm}

% DB question_id: e2a85296-0528-50e6-a9a5-013c6a3adcef
\needspace{5\baselineskip}
\noindent\textbf{\color{ink}1021.}
\par
23478 MB fayl 344 Mbit/s tezlikda uzatilsa, ideal sharoitda uzatish qancha vaqt oladi?
\begin{enumerate}[label=\Alph*),leftmargin=9mm]
\item 546 soniya
\item 273 soniya
\item 4368 soniya
\item 890 soniya
\end{enumerate}
\begin{tcolorbox}[enhanced,breakable,colback=teal!4,colframe=teal!50!black,boxrule=0.5pt,arc=1.1mm,left=3mm,right=3mm,top=1.7mm,bottom=1.7mm,before skip=2mm,after skip=2mm]
\textbf{\color{teal}To'g'ri javob:} \textbf{A) 546 soniya}
\par\smallskip
\textbf{Izoh:} 23478 MB = 187824 Mbit. Vaqt = hajm / tezlik = 187824 / 344 = 546 soniya.
\end{tcolorbox}
\vspace{2mm}\hrule height 0.2pt \vspace{3mm}

% DB question_id: e557625a-76ab-517b-b2b5-d49579ab07c9
\needspace{5\baselineskip}
\noindent\textbf{\color{ink}1022.}
\par
7520 MB fayl 256 Mbit/s tezlikda uzatilsa, ideal sharoitda uzatish qancha vaqt oladi?
\begin{enumerate}[label=\Alph*),leftmargin=9mm]
\item 117 soniya
\item 235 soniya
\item 1880 soniya
\item 491 soniya
\end{enumerate}
\begin{tcolorbox}[enhanced,breakable,colback=teal!4,colframe=teal!50!black,boxrule=0.5pt,arc=1.1mm,left=3mm,right=3mm,top=1.7mm,bottom=1.7mm,before skip=2mm,after skip=2mm]
\textbf{\color{teal}To'g'ri javob:} \textbf{B) 235 soniya}
\par\smallskip
\textbf{Izoh:} 7520 MB = 60160 Mbit. Vaqt = hajm / tezlik = 60160 / 256 = 235 soniya.
\end{tcolorbox}
\vspace{2mm}\hrule height 0.2pt \vspace{3mm}

% DB question_id: ed0eb54d-901f-5658-bb1b-533d06e02cf3
\needspace{5\baselineskip}
\noindent\textbf{\color{ink}1023.}
\par
11480 MB fayl 328 Mbit/s tezlikda uzatilsa, ideal sharoitda uzatish qancha vaqt oladi?
\begin{enumerate}[label=\Alph*),leftmargin=9mm]
\item 140 soniya
\item 2240 soniya
\item 608 soniya
\item 280 soniya
\end{enumerate}
\begin{tcolorbox}[enhanced,breakable,colback=teal!4,colframe=teal!50!black,boxrule=0.5pt,arc=1.1mm,left=3mm,right=3mm,top=1.7mm,bottom=1.7mm,before skip=2mm,after skip=2mm]
\textbf{\color{teal}To'g'ri javob:} \textbf{D) 280 soniya}
\par\smallskip
\textbf{Izoh:} 11480 MB = 91840 Mbit. Vaqt = hajm / tezlik = 91840 / 328 = 280 soniya.
\end{tcolorbox}
\vspace{2mm}\hrule height 0.2pt \vspace{3mm}

% DB question_id: f018f0b9-4ca7-4960-8e42-b9492a5692d7
\needspace{5\baselineskip}
\noindent\textbf{\color{ink}1024.}
\par
20 Mbit/s tezlikda 100 Mbit axborot ideal holatda qancha vaqtda uzatiladi?
\begin{enumerate}[label=\Alph*),leftmargin=9mm]
\item 2 s
\item 5 s
\item 20 s
\item 2000 s
\end{enumerate}
\begin{tcolorbox}[enhanced,breakable,colback=teal!4,colframe=teal!50!black,boxrule=0.5pt,arc=1.1mm,left=3mm,right=3mm,top=1.7mm,bottom=1.7mm,before skip=2mm,after skip=2mm]
\textbf{\color{teal}To'g'ri javob:} \textbf{B) 5 s}
\par\smallskip
\textbf{Izoh:} t = I \ensuremath{\div} v = 100 \ensuremath{\div} 20 = 5 soniya.
\end{tcolorbox}
\vspace{2mm}\hrule height 0.2pt \vspace{3mm}

% DB question_id: f41c80e0-d712-5ffa-9e60-f1a082231cbc
\needspace{5\baselineskip}
\noindent\textbf{\color{ink}1025.}
\par
9321 MB fayl 312 Mbit/s tezlikda uzatilsa, ideal sharoitda uzatish qancha vaqt oladi?
\begin{enumerate}[label=\Alph*),leftmargin=9mm]
\item 551 soniya
\item 239 soniya
\item 119 soniya
\item 1912 soniya
\end{enumerate}
\begin{tcolorbox}[enhanced,breakable,colback=teal!4,colframe=teal!50!black,boxrule=0.5pt,arc=1.1mm,left=3mm,right=3mm,top=1.7mm,bottom=1.7mm,before skip=2mm,after skip=2mm]
\textbf{\color{teal}To'g'ri javob:} \textbf{B) 239 soniya}
\par\smallskip
\textbf{Izoh:} 9321 MB = 74568 Mbit. Vaqt = hajm / tezlik = 74568 / 312 = 239 soniya.
\end{tcolorbox}
\vspace{2mm}\hrule height 0.2pt \vspace{3mm}

% DB question_id: f5f972c2-463c-4bb6-815a-96eece3739ca
\needspace{5\baselineskip}
\noindent\textbf{\color{ink}1026.}
\par
Internet orqali 15 MB hajmdagi fayl 2 Mbit/s tezlikda uzatilmoqda. Faylni to'liq yuklab olish uchun qancha vaqt ketadi?
\begin{enumerate}[label=\Alph*),leftmargin=9mm]
\item 60 soniya
\item 30 soniya
\item 7,5 soniya
\item 15 soniya
\end{enumerate}
\begin{tcolorbox}[enhanced,breakable,colback=teal!4,colframe=teal!50!black,boxrule=0.5pt,arc=1.1mm,left=3mm,right=3mm,top=1.7mm,bottom=1.7mm,before skip=2mm,after skip=2mm]
\textbf{\color{teal}To'g'ri javob:} \textbf{A) 60 soniya}
\par\smallskip
\textbf{Izoh:} 15 MB \ensuremath{\times} 8 = 120 Mbit, 120/2 = 60 soniya (protokol xarajatlari hisobga olinmaganda).
\end{tcolorbox}
\vspace{2mm}\hrule height 0.2pt \vspace{3mm}

% DB question_id: e42a1d99-4eb5-4323-87a1-76fc1080b2b2
\needspace{5\baselineskip}
\noindent\textbf{\color{ink}1027.}
\par
42 MiB fayl 10 Mbit/s kanalda ideal holatda taxminan qancha vaqtda uzatiladi?
\begin{enumerate}[label=\Alph*),leftmargin=9mm]
\item 4,2 s
\item 33,6 s
\item 42 s
\item 35,2 s
\end{enumerate}
\begin{tcolorbox}[enhanced,breakable,colback=teal!4,colframe=teal!50!black,boxrule=0.5pt,arc=1.1mm,left=3mm,right=3mm,top=1.7mm,bottom=1.7mm,before skip=2mm,after skip=2mm]
\textbf{\color{teal}To'g'ri javob:} \textbf{D) 35,2 s}
\par\smallskip
\textbf{Izoh:} 42 \ensuremath{\times} 2\textsuperscript{2}\textsuperscript{0} \ensuremath{\times} 8 \ensuremath{\div} 10\textsuperscript{7} \ensuremath{\approx} 35,2 soniya.
\end{tcolorbox}
\vspace{2mm}\hrule height 0.2pt \vspace{3mm}

% DB question_id: eb78d4c2-f4e2-46fb-80a2-26f4d271635a
\needspace{5\baselineskip}
\noindent\textbf{\color{ink}1028.}
\par
Nazariy tezlik 100 Mbit/s, amaliy throughput 80 Mbit/s. Samaradorlik qancha?
\begin{enumerate}[label=\Alph*),leftmargin=9mm]
\item 20\%
\item 100\%
\item 80\%
\item 125\%
\end{enumerate}
\begin{tcolorbox}[enhanced,breakable,colback=teal!4,colframe=teal!50!black,boxrule=0.5pt,arc=1.1mm,left=3mm,right=3mm,top=1.7mm,bottom=1.7mm,before skip=2mm,after skip=2mm]
\textbf{\color{teal}To'g'ri javob:} \textbf{C) 80\%}
\par\smallskip
\textbf{Izoh:} 80 \ensuremath{\div} 100 \ensuremath{\times} 100\% = 80\%.
\end{tcolorbox}
\vspace{2mm}\hrule height 0.2pt \vspace{3mm}

% DB question_id: 9ce14b1f-ac9b-45bf-8d2f-591cc8c86693
\needspace{5\baselineskip}
\noindent\textbf{\color{ink}1029.}
\par
10 Mbit/s kanal orqali 75 MB fayl ideal sharoitda necha soniyada uzatiladi? MB = 10\textsuperscript{6} B.
\begin{enumerate}[label=\Alph*),leftmargin=9mm]
\item 7,5 s
\item 60 s
\item 75 s
\item 600 s
\end{enumerate}
\begin{tcolorbox}[enhanced,breakable,colback=teal!4,colframe=teal!50!black,boxrule=0.5pt,arc=1.1mm,left=3mm,right=3mm,top=1.7mm,bottom=1.7mm,before skip=2mm,after skip=2mm]
\textbf{\color{teal}To'g'ri javob:} \textbf{B) 60 s}
\par\smallskip
\textbf{Izoh:} 75 MB \ensuremath{\times} 8 = 600 Mbit; 600 \ensuremath{\div} 10 = 60 soniya.
\end{tcolorbox}
\vspace{2mm}\hrule height 0.2pt \vspace{3mm}

% DB question_id: c6f11758-6103-436d-9ea6-446a018f088b
\needspace{5\baselineskip}
\noindent\textbf{\color{ink}1030.}
\par
2400 baud va 2400 bit/s qachon son jihatdan teng bo'lishi mumkin?
\begin{enumerate}[label=\Alph*),leftmargin=9mm]
\item har simvol 2 bit tashiganda
\item faqat optik tolada
\item har simvol 1 bit tashiganda
\item har doim
\end{enumerate}
\begin{tcolorbox}[enhanced,breakable,colback=teal!4,colframe=teal!50!black,boxrule=0.5pt,arc=1.1mm,left=3mm,right=3mm,top=1.7mm,bottom=1.7mm,before skip=2mm,after skip=2mm]
\textbf{\color{teal}To'g'ri javob:} \textbf{C) har simvol 1 bit tashiganda}
\par\smallskip
\textbf{Izoh:} Bit tezligi = baud \ensuremath{\times} bit/simvol; bit/simvol = 1 bo'lsa teng.
\end{tcolorbox}
\vspace{2mm}\hrule height 0.2pt \vspace{3mm}

% DB question_id: ca2bdc08-7d36-4203-9f70-81b8c685c26f
\needspace{5\baselineskip}
\noindent\textbf{\color{ink}1031.}
\par
Bitta simvol 4 bit tashigani sabab bit tezligi bauddan to'rt marta katta bo'ldi. Yaqin tushunchalarni farqlab, aynan qaysi atama vaziyatni eng aniq ifodalashini toping.
\begin{enumerate}[label=\Alph*),leftmargin=9mm]
\item HTTPS va ishonchlilik farqi
\item Axborot va ma'lumot farqi
\item Baud va bit/s farqi
\item Kodlash va shifrlash farqi
\end{enumerate}
\begin{tcolorbox}[enhanced,breakable,colback=teal!4,colframe=teal!50!black,boxrule=0.5pt,arc=1.1mm,left=3mm,right=3mm,top=1.7mm,bottom=1.7mm,before skip=2mm,after skip=2mm]
\textbf{\color{teal}To'g'ri javob:} \textbf{C) Baud va bit/s farqi}
\par\smallskip
\textbf{Izoh:} Vaziyat Baud va bit/s farqiga misol: baud simvol tezligi, bit/s esa axborot bitlari tezligi ekanini ajratish.
\end{tcolorbox}
\vspace{2mm}\hrule height 0.2pt \vspace{3mm}
\clearpage
\section{Darslikdan aniqlangan qo'shimcha nazariya va atamalar}
\textbf{Bu bo'lim nega kiritildi?}
2026-yilgi spetsifikatsiyada axborot, kodlash, axborotning hajmi va
raqamli madaniyat asoslari belgilangan. Darsliklar mazmuni bilan
mavjud 1 031 ta savol solishtirilganda, ayrim tushuncha nomlari
savollarning asosiy matnida aniq uchramadi yoki faqat qisman qamrab olindi.
Bunda atama so'zma-so'z topilmasligi uning mohiyati boshqa savolda
hech qachon so'ralmaganini isbotlamaydi.
\subsection{Axborot sifati: to'g'rilik, to'liqlik va moslik}
\textbf{Nazariya.} Axborotni baholashda uning to'g'riligi, maqsadga mosligi, yetarlicha batafsilligi va to'liqligi farqlanadi. Xabar to'g'ri bo'lsa-da, joriy vazifaga daxlsiz bo'lishi mumkin. Yetishmayotgan muhim tafsilot ham xulosani buzishi mumkin.
\par\textbf{Manba:} Cambridge 10--11-sinf, PDF 15--16-betlar, 1.03-bo'lim
\par\textbf{Savol matnlari bo'yicha qamrov:} 'axborot sifati' iborasi aynan savol matnida 0 marta uchradi.
\subsection{Dolzarblik va tafsilotlar darajasi}
\textbf{Nazariya.} Dolzarblik - axborotning uni ishlatish maqsadiga aloqadorligi. Juda ko'p tafsilot kerakli fikrni yashirishi mumkin; juda oz tafsilot esa asosli qaror qabul qilishga xalaqit qiladi. Bu ikki holatni amaliy misollar bilan farqlab o'rganing.
\par\textbf{Manba:} Cambridge 10--11-sinf, PDF 16-bet
\par\textbf{Savol matnlari bo'yicha qamrov:} 'dolzarblik' iborasi aynan savol matnida 0 marta uchradi.
\subsection{Raqamli obro' va raqamli izni farqlash}
\textbf{Nazariya.} Raqamli iz - insonning internet faoliyatidan qolgan sharh, foto, yozuv va boshqa dalillar. Raqamli obro' - boshqalarning shaxs haqida internetda shakllantirgan fikri va taassuroti. Raqamli izlar odamning kelajakdagi obro'siga ta'sir qilishi mumkin.
\par\textbf{Manba:} Cambridge 7-sinf, PDF 117 va 142-betlar
\par\textbf{Savol matnlari bo'yicha qamrov:} 'raqamli obro' atamasi aynan savol matnida 0 marta, 'raqamli iz' 2 marta uchradi.
\subsection{Morze kodi: qisqa va uzun signal}
\textbf{Nazariya.} Morze tizimida belgilar qisqa va uzun signallar (nuqta va tire) ketma-ketligi orqali ifodalanadi. Bu tarixiy kodlash usuli. Morze signallarini bit va baytlarning fizik ta'rifi bilan tenglashtirmaslik kerak.
\par\textbf{Manba:} ICT 5-sinf (2020), PDF 18-bet, 4-dars
\par\textbf{Savol matnlari bo'yicha qamrov:} 'Morze' atamasi aynan savol matnida 0 marta uchradi.
\subsection{Qidiruv natijasidagi reklama va ishonchlilik}
\textbf{Nazariya.} Qidiruv tizimidagi eng yuqoridagi natija doim ham eng yaxshi yoki eng ishonchli manba emas. Ayrim natijalar pul to'langan reklama sifatida yuqori turadi. Ma'lumotni muallif, manba, dolzarblik va dalil orqali baholang.
\par\textbf{Manba:} Cambridge 5-sinf, PDF 95-bet, 5-bob
\par\textbf{Savol matnlari bo'yicha qamrov:} 'sponsor' va 'birinchi natija' iboralari savol matnlarida 0 marta uchradi; umumiy reklama savollari bor.
\subsection{Veb-manbaga havola berish va plagiatdan saqlanish}
\textbf{Nazariya.} Internetdan matn, dalil yoki fikr olganda uning kelib chiqish manbasini ko'rsatish kerak. O'z so'zlari bilan qayta ifodalash ham olingan g'oyaning manbasini ko'rsatishga bo'lgan ehtiyojni yo'qotmaydi. Manbasiz so'zma-so'z ko'chirish plagiat muammosini tug'diradi.
\par\textbf{Manba:} ICT 6-sinf (2021), PDF 87-bet
\par\textbf{Savol matnlari bo'yicha qamrov:} 'manba ko'rsatish' iborasi 1 savolda uchradi, lekin 'havola ko'rsatish' aynan 0 marta uchradi.
\subsection{Kodlash va shifrlashning farqi}
\textbf{Nazariya.} Kodlash ma'lumotni belgilangan formatga o'tkazadi; shifrlash uni tegishli kalitsiz anglab bo'lmas holga keltirish uchun o'zgartiradi. Dasturiy kod yozish esa kompyuter buyruqlarini yaratish bilan bog'liq. Uchala tushuncha bir narsa emas.
\par\textbf{Manba:} Cambridge 10--11-sinf, PDF 17-bet, 1.04-bo'lim
\par\textbf{Savol matnlari bo'yicha qamrov:} Mazmuniy qamrov qisman: kodlash bo'yicha testlar ko'p; aynan uchalasini bir-biridan ajratadigan savollar to'liq tekshirilmagan.
\subsection{Elektron xavfsizlik va nomaqbul kontent}
\textbf{Nazariya.} Elektron xavfsizlik faqat parol emas: nomaqbul kontent, mualliflik huquqi buzilishi, plagiat, onlayn tahqirlash va muloqotdagi xavf-xatarlarni ham o'z ichiga oladi. Ularni manba va vaziyat asosida tanib olish lozim.
\par\textbf{Manba:} Cambridge 9-sinf, PDF 115-bet, 8.02-bo'lim
\par\textbf{Savol matnlari bo'yicha qamrov:} Ushbu elementlarning ayrimlari testlarda mavjud. Har bir kichik xavf alohida savolda sinovdan o'tgan deb xulosa qilinmaydi.
\section{Darsliklar bo'yicha tekshiruv reyestri}
Ko'rsatilgan raqamlar darslikning \textbf{PDF sahifalari}, bosma sahifalar
raqami ba'zi kitoblarda farqlanishi mumkin. Quyidagi reyestr ushbu mavzuga
aloqador yoki yaqin boblar ustida bajarilgan \emph{tanlangan sahifalar}
tekshiruvini ko'rsatadi; 13 ta darslikning har bir sahifasi to'liq
ekspert tahlilidan o'tgani haqidagi da'vo emas.
\begin{longtable}{@{}p{60mm}p{21mm}p{75mm}@{}}
\toprule
\textbf{Darslik}&\textbf{Sahifa}&\textbf{Tekshirilgan yo'nalish}\\
\midrule\endfirsthead
\toprule \textbf{Darslik}&\textbf{Sahifa}&\textbf{Yo'nalish}\\\midrule\endhead
5-sinf (Cambridge) & PDF 95 & Manba ishonchliligi va reklama natijalari\\
6-sinf (Cambridge) & PDF 29 & Maqsadli kontent tanlash\\
7-sinf (Cambridge) & PDF 117, 142 & Raqamli iz va obro'\\
8-sinf (Cambridge) & PDF 107 & Xavfsizlik terminlari (qo'shni mavzu)\\
9-sinf (Cambridge) & PDF 115 & Elektron xavfsizlik\\
10--11-sinf (Cambridge) & PDF 15--17 & Axborot sifati, kodlash va shifrlash\\
5-sinf (ICT, 2020) & PDF 18--21 & Morze, ikkilik kodlash, bit/bayt\\
6-sinf (ICT, 2021) & PDF 87--88 & Veb-manbaga havola va plagiat\\
7-sinf (ICT, 2021) & PDF 5--7 & Axborot kanali, analog va diskret shakllar\\
8-sinf (ICT, 2020) & PDF 4--6 & Ijtimoiy media; asosan boshqa mavzuga tegishli\\
9-sinf (ICT, 2020) & PDF 4--6 & Mantiq, 4-mavzuga oid\\
10-sinf (ICT, 2021) & PDF 14, 21 & Saqlash va mobil qurilmalar, 2-mavzu bilan tutash\\
11-sinf (ICT, 2021) & PDF 14, 136 & Ma'lumotlar bazasi va Internet/intranet, keyingi mavzular\\
\bottomrule\end{longtable}
\section*{Qamrov va nazorat bayonnomasi}
\addcontentsline{toc}{section}{Qamrov va nazorat bayonnomasi}
1 031 ta faol savol, 1 031 ta javob va 1 031 ta izoh asosiy qismga kiritildi.
Y1: 1 027, Y2: 2, Y3: 2. Asl bankdagi 147 ta archived versiya
alohida texnik audit reyestrida qayd qilingan; asosiy testlar soniga qo'shilmagan.
Barcha testlar manba bazadagi holati bilan bog'langan.
Darsliklardan yangi tuzilgan izohli nazariyalar mavzuning qo'shimcha
o'rganish qismi hisoblanadi. Ekspert tomonidan barcha fan faktlari
mustaqil tekshirilganligi da'vo qilinmaydi.
\end{document}